\documentclass[11pt,a4paper]{article}
\usepackage{amsmath,amssymb,amsthm,mathtools,mathrsfs}
\usepackage[margin=23mm]{geometry}
\usepackage{hyperref}
\allowdisplaybreaks
\newtheorem{theorem}{Theorem}
\newtheorem{proposition}[theorem]{Proposition}
\newcommand{\C}{\mathbb C}
\newcommand{\R}{\mathbb R}
\newcommand{\Tr}{\operatorname{Tr}}
\newcommand{\diag}{\operatorname{diag}}
\title{The joint Schur return in the original divided jets\newline and its signed arithmetic receiving bounds}
\author{Split-Zero research programme}
\date{14 September 2026}
\begin{document}
\maketitle

\section{The original objects and both coordinate maps}

Fix the original complete quartet with
$0<\delta<1/2$, $\gamma>2$, multiplicity $m\geq1$,
$k=4l+1$, $l\geq1$, and retain
\[
\begin{gathered}
 c=k/2,\quad e_k=1+k(m-1),\quad q=(k+1)^2e_k,\\
 \lambda_{ab}=c+(2a-k)\delta+i(2b-k)\gamma,\qquad
 \chi(S)=\prod_{a,b=0}^k(S-\lambda_{ab})^{e_k},\\
 t_j=2j+\tfrac12,\quad\alpha_j=c+t_j,\quad
 f(S)=\prod_{j=0}^{l-1}(S-\alpha_j),\qquad
 E=\C[S]/(\chi).
\end{gathered}\tag{JSR1}
\]
All remainders use the increasing coefficient frame
$1,S,\ldots,S^{q-1}$. Since $k$ is odd,
$\operatorname{Im}\lambda_{ab}=(2b-k)\gamma\ne0$,
whereas each $\alpha_j$ is real. Thus $f$ and $\chi$ are coprime.
The original multiplication $M_f:E\to E$ is invertible, with
inverse its Bezout remainder. No original Taylor unit, ordered
tensor inclusion, source map or boundary observation is altered.

The two original reference measures and their masses are
\[
\begin{gathered}
 d\sigma(y)=\frac{|\Gamma(1/4+iy/2)|^2}{2\pi}\,dy,
 \quad\int d\sigma=\sqrt{2\pi},\quad S=c+iy,\\
 dm_{0,k}=\beta_k|f(c+iy)|^2d\sigma,
 \quad\int dm_{0,k}=(2\pi)^{k/2},\\
 \beta_k=\frac{(2\pi)^{k/2}}{\sqrt2\,\Gamma(k/2)}
        =\frac{(\sqrt{2\pi})^k}{c_{k/4}}4^{-l},
 \qquad c_a=2^{1-2a}\Gamma(2a).
\end{gathered}\tag{JSR2}
\]
The last equality follows because
$4^lc_{k/4}=2^{2l+1-k/2}\Gamma(k/2)=\sqrt2\Gamma(k/2)$.
Consequently the new note's $C_k$ is exactly GCR's $\beta_k$.
Gamma recurrence gives the measure identity with the original
coordinate and both masses retained.

For $J_N:\mathcal P_N\to E$ the original remainder map, let
$\mathsf E:E\to\C^q$ be the divided Hermite jet map, with
rows $(a,b,r)$ in the original order:
\[
 (\mathsf E)_{(a,b,r),n}=\binom nr\lambda_{ab}^{n-r},\qquad
 \mathsf D=\diag_{a,b,\,0\leq r<e_k}(r!).
 \tag{JSR3}
\]
Entries with $n<r$ are zero. The raw jet map in the new note is
exactly $\mathsf E^{\rm raw}=\mathsf D\mathsf E$, and hence
its inverse is $\mathsf E^{-1}\mathsf D^{-1}$.
The confluent Vandermonde determinant of $\mathsf E$ is
$\prod_{\lambda<\mu}(\mu-\lambda)^{e_k^2}$; the raw map
retains in addition $\det\mathsf D=\prod_{a,b,r}r!$.
On the complete primary jets let $\mathsf U_f$ be multiplication
by $f$, with entries $f^{[r-s]}(\lambda)$ for $r\geq s$.
Then
\[
 \mathsf E M_f=\mathsf U_f\mathsf E,\qquad
 \mathsf F_\chi^{\rm raw}=\mathsf D\mathsf U_f\mathsf D^{-1},
 \qquad
 \det M_f=\det\mathsf U_f=\prod_{a,b}f(\lambda_{ab})^{e_k}.
 \tag{JSR4}
\]
Indeed the $(r,s)$ entry after this conjugation is
$r!f^{(r-s)}(\lambda)/(s!(r-s)!)
=\binom rs f^{(r-s)}(\lambda)$, exactly the raw Leibniz matrix.
This proves the coordinate dictionary on every nilpotent
direction, including all derivative factorials.

\section{The two metrized small matrices and the full mixed determinant}

For $N\geq q-1$ put $M=N+l$. Write all covariances in the
original coefficient remainder frame. With
$Z_MP=(P(\alpha_j))_{0\leq j<l}$ and the full sigma source
Gram $H_M^\sigma$, define
\[
 \binom{J_M}{Z_M}(H_M^\sigma)^{-1}
       \begin{pmatrix}J_M^*&Z_M^*\end{pmatrix}
 =\begin{pmatrix}K_M&C_M\\C_M^*&A_M\end{pmatrix},\quad
 G_M^\sigma=K_M^{-1}.
 \tag{JSR5}
\]
The joint observation is onto: division by the coprime monic
polynomial $f\chi$ realizes every joint remainder at degree
$q+l-1\leq M$. Its kernel is $f\chi\mathcal P_{N-q}$.
Thus the joint covariance is positive definite and both
$S_M=K_M-C_MA_M^{-1}C_M^*$ and
$L_M=A_M-C_M^*K_M^{-1}C_M$ are positive definite.
Multiplication $P\mapsto fP$ maps the original fibre
$J_NP=v$ bijectively onto the joint fibre $(M_fv,0)$.
Its inverse is polynomial division on $\ker Z_M$.
Minimization in this exact fibre and JSR2 give
\[
 G_N^0=\beta_k M_f^*S_M^{-1}M_f.
 \tag{JSR6}
\]
This proof retains the actual cross covariance $C_M$.

Let $p_n$ be the original monic sigma polynomials in $S$.
Their norms are exactly
$\gamma_n=\sqrt{2\pi}(2n)!/4^n
=\sqrt{2\pi}\,n!(1/2)_n$.
In the new note's notation put
\[
\begin{gathered}
 U_N=[Jp_{N+1},\ldots,Jp_{N+l}],\quad
 \Omega_N=\diag(\gamma_{N+1},\ldots,\gamma_{N+l}),\\
 X_N=\Omega_N^{-1}U_N^*G_N^\sigma U_N,\qquad
 Y_N=A_M^{-1}C_M^*G_M^\sigma C_M,\\
 \widehat X_N=\Omega_N^{1/2}X_N\Omega_N^{-1/2},\qquad
 \widehat Y_N=A_M^{1/2}Y_NA_M^{-1/2}.
\end{gathered}\tag{JSR7}
\]
The maps $x\mapsto\Omega_N^{1/2}x$ and
$y\mapsto A_M^{1/2}y$ are isometries from their displayed
positive metrics to the standard coordinate metric.
Thus $\widehat X_N=\Omega_N^{-1/2}U_N^*G_N^\sigma
U_N\Omega_N^{-1/2}\succeq0$ and
$\widehat Y_N=A_M^{-1/2}C_M^*G_M^\sigma C_MA_M^{-1/2}
\succeq0$. Moreover
$I-\widehat Y_N=A_M^{-1/2}L_MA_M^{-1/2}>0$.
In particular $X_N$ is self-adjoint in $\Omega_N$, and
$Y_N$ is self-adjoint in $A_M$ with eigenvalues in $[0,1)$.

Use tildes to denote precisely the GCR divided-jet matrices,
so that no coefficient covariance is identified with a jet
covariance without its map. Direct substitution of JSR3 gives
\[
\begin{gathered}
 \widetilde K_N=\mathsf E K_N\mathsf E^*,\qquad
 \widetilde U_N=\mathsf E U_N\Omega_N^{-1/2},\qquad
 \widetilde V_N=\mathsf E C_MA_M^{-1/2},\\
 \widetilde W_N=(\widetilde U_N\ \widetilde V_N),\quad
 J_l=\diag(I_l,-I_l),\qquad
 \widetilde S_N=\mathsf E S_M\mathsf E^*.
\end{gathered}\tag{JSR8}
\]
The real-node observation has order zero, so no factorial
coordinate change occurs on its $l$-dimensional target.
All cross kernels transform by the left factor $\mathsf E$.
The original finite source kernel sum gives
$K_M=K_N+U_N\Omega_N^{-1}U_N^*$; hence
$\widetilde S_N=\widetilde K_N+
\widetilde U_N\widetilde U_N^*-
\widetilde V_N\widetilde V_N^*$, exactly GCR22.
The GCR24 positive factors are therefore precisely
\[
\begin{gathered}
 \mathcal A_N^{\rm GCR}=I+\widehat X_N
      =\Omega_N^{1/2}(I+X_N)\Omega_N^{-1/2},\\
 \mathcal B_N^{\rm GCR}=I-\widehat Y_N
      =A_M^{1/2}(I-Y_N)A_M^{-1/2},\\
 \mathfrak d_N:=\frac{\det\widetilde S_N}{\det\widetilde K_N}
 =\det\bigl(I_{2l}+J_l\widetilde W_N^*
                  \widetilde K_N^{-1}\widetilde W_N\bigr)\\
 =\det(I+X_N)\det(I-Y_N)>0.
\end{gathered}\tag{JSR9}
\]
For completeness, the determinant of
$\left(\begin{smallmatrix}I&A\\-B&I\end{smallmatrix}\right)$
computed by either diagonal block proves
$\det(I+AB)=\det(I+BA)$, giving the $2l$ determinant.
Alternatively factor $\det S_M/\det K_N$ as
$(\det K_M/\det K_N)(\det S_M/\det K_M)$.
The two factors are respectively $\det(I+X_N)$ and
$\det(I-Y_N)$ by the same identity and the full joint
block determinant. This proves JSR9 without requiring
the two differently metrized endomorphisms to commute.
In the GCR matrices the second positive factor retains
the entire mixed term
\[
\begin{aligned}
 \mathcal B_N^{\rm GCR}={}&I-
 \widetilde V_N^*\widetilde K_N^{-1}\widetilde V_N\\
 &+\widetilde V_N^*\widetilde K_N^{-1}\widetilde U_N
 (I+\widetilde U_N^*\widetilde K_N^{-1}\widetilde U_N)^{-1}
 \widetilde U_N^*\widetilde K_N^{-1}\widetilde V_N.
\end{aligned}
\]
Multiplying the resulting inverse by
$\widetilde K_N+\widetilde U_N\widetilde U_N^*$ verifies
this equality with the displayed order of every factor.

The note's scalar $d_N^{\rm note}=\log\det(I+X_N)$ is
a logarithmic update, whereas GCR23's $d_N$ is the positive
determinant $\mathfrak d_N$ in JSR9. Their exact relation is
\[
 \log\mathfrak d_N=d_N^{\rm note}-\beta_N^{\rm note},
 \qquad\beta_N^{\rm note}=-\log\det(I-Y_N).
 \tag{JSR10}
\]
This also identifies both objects with the original source
and relation determinants, as follows.

\section{Exact source and relation determinant dictionary}

Write $H_N^f$ for the coefficient Gram in $|f|^2d\sigma$.
The complete Gamma monic norms give the positive source ratio
\[
 D_{l,N}=\frac{\det H_N^f}{\det H_N^\sigma}
     =\prod_{a=0}^N\prod_{r=0}^{2l-1}(a+r+\tfrac12).
 \tag{JSR11}
\]
Indeed the ratio of the monic norms at index $a$ is
$\gamma_{k,a}/(\beta_k\gamma_a)
=\Gamma(a+2l+1/2)/\Gamma(a+1/2)$, the displayed finite product.
Products of all monic norms equal the coefficient determinant
because the monic triangular change of frame has determinant one.
For $r_N=N+1-q$ define the full relation ratio
\[
 R_N=\frac{\det\operatorname{Gram}_\sigma
                 (f\chi,f\chi S,\ldots,f\chi S^{r_N-1})}
             {\det\operatorname{Gram}_\sigma
                 (\chi,\chi S,\ldots,\chi S^{r_N-1})}.
 \tag{JSR12}
\]
At $r_N=0$ both empty determinants are one.
The coefficient frame
$(1,\ldots,S^{q-1},\chi,\ldots,\chi S^{N-q})$
has determinant one. Block elimination along the entire
relation kernel gives $\det H_N=\det H_{\rm rel}\det G_N$.
Applying this twice yields
$\det G_N^f/\det G_N^\sigma=D_{l,N}/R_N$.
On the other hand JSR6 and JSR9 give that ratio as
$|\det M_f|^2/\mathfrak d_N$. Consequently the exact
dictionary is
\[
 \boxed{\quad
 \mathfrak d_N=|\det M_f|^2\frac{R_N}{D_{l,N}},\qquad
 d_N^{\rm note}-\beta_N^{\rm note}
 =2\log|\det M_f|+\log R_N-\log D_{l,N}.
 \quad}\tag{JSR13}
\]
This is an equality on the original remainder and relation
spaces. Every full-unit multiplication is retained by JSR4.
For the original boundary $\Lambda:E\to\operatorname{im}\Lambda$,
the product-jet observation is exactly
$(0,Y)\mapsto\Lambda\mathsf E^{-1}\mathsf U_f^{-1}Y$.
It sends $(0,\mathsf U_f\mathsf E v)$ to $\Lambda v$,
so the source comparison also retains the actual boundary.

Put $(N_0,N_1,N_2,N_3)=(q-1,q,2q-1,2q)$ and
$(\epsilon_0,\epsilon_1,\epsilon_2,\epsilon_3)=(1,1,-1,-1)$.
Define the same four-volume functional for each original
quotient metric and retain
$\mathfrak T_k=F_{0,k}-F_{\sigma,k}$.
Only after taking those four metrics and signs do
$q\log\beta_k$ and $2\log|\det M_f|$ cancel. Thus
\[
\begin{gathered}
 \mathcal H_k:=\mathcal R_k^0-\mathcal R_k^\sigma
 =\sum_{a=0}^3\epsilon_a\log\mathfrak d_{N_a}
 =-\mathfrak T_k,\\
 \mathcal H_k=\sum_{a=0}^3\epsilon_a
       \{\log\det(I+X_{N_a})+\log\det(I-Y_{N_a})\}.
\end{gathered}\tag{JSR14}
\]
In the precise GSR low-endpoint notation,
$R_{q-1}=1$ and
$R_q=\mathfrak r_\chi=
\frac{\int|f\chi|^2d\sigma}{\int|\chi|^2d\sigma}$.
Therefore
\[
\begin{gathered}
 C_k^{\rm rel}=\log D_{l,q-1}+\log D_{l,q}
   -\log\mathfrak r_\chi-\log D_{l,2q-1}-\log D_{l,2q},\\
 \mathcal H_k=-C_k^{\rm rel}-\log R_{2q-1}-\log R_{2q}.
\end{gathered}\tag{JSR15}
\]
The source mass $\beta_k$ is not $C_k^{\rm rel}$.
Every logarithm is of a positive scalar, and the two low
relation dimensions remain exactly zero and one.

\section{The full original relation return through the universal Hankel source}

Retain the entire original packet, $S=c+iy$, all original maps
in JSR1--15, and put $d=\delta^2$, $g=\gamma^2$ and
$e=e_k=1+k(m-1)$, so $q=e(k+1)^2$. The following full-source
calculation applies on the explicit original domain
\[
 k=4l+1,\quad l\geq1,\quad
 k\geq2048\sqrt{d+g},\quad
 \epsilon=2^{-10},\qquad r\in\{q,q+1\}.
 \tag{BHR1}
\]
Indeed $R_*/q=k\sqrt{d+g}/q\leq\epsilon/2$ and
$t_{l-1}/q=(2l-3/2)/q\leq\epsilon$ there. These are the
BRD22 hypotheses, including a strict density expansion
radius. Below this threshold the original exact jet,
relation, and finite trace formulas remain in use.

The complete constants are
\[
\begin{gathered}
 \alpha_k^\chi=\frac{k(k+2)(d-g)}{3q},\quad
 \tau_\epsilon=\frac{k^2(d+g)}{q^2\epsilon^2}\leq\frac14,
 \quad\delta_\chi=\frac{q\tau_\epsilon^2}{2(1-\tau_\epsilon)},\\
 E_\chi=l\{ |\alpha_k^\chi|(\epsilon^{-2}-64^{-2})
                                   +2\delta_\chi\},\qquad
 \Delta_f=\frac{S_2}{q^2\epsilon^2},\quad
 S_2=\sum_{j=0}^{l-1}(2j+\tfrac12)^2
       =\frac{l(16l^2-12l-1)}{12},\\
 \kappa_{I,r}=\frac{2\epsilon r^2}{\sqrt{\cos1}}
          \left(\frac43\,2^{-13}e^{2\pi}\right)^q,
 \quad\kappa_F=\frac{512q}{59\sqrt{\cos1}}e^{-19q},\quad
 L_r^{\rm tail}=r\log(1+\kappa_{I,r}+\kappa_F).
\end{gathered}\tag{BHR2}
\]
Here $\alpha_k^\chi$ is exactly CTR23's $a_k$, given a
superscript to preserve its distinction from the arithmetic
allowances $a_k^-,a_k^+$ in JSR20. Likewise
$L_r^{\rm tail}$ is CTR38's $L_r$ and is distinct from
the finite trace endpoints. The complete BRD and CTR bodies
accompany this calculation with these constants unchanged.

Let $H_{\chi,0}$ and $H_{\chi,f}$ be the original Grams of
$(\chi S^j)_{j<r}$ and $(f\chi S^j)_{j<r}$, respectively.
Their exact ratio is the GSR relation determinant
$R_{q+r-1}$. The new full moment matrices are
\[
 (\mathsf M_{a,r})_{ij}
 =\int_{\R}y^{2a+i+j}\sigma(y)\,dy,\qquad
 0\leq i,j<r,\qquad
 \rho_r=\log\frac{\det\mathsf M_{q+l,r}}
                         {\det\mathsf M_{q,r}}.
 \tag{BHR3}
\]
They are positive definite: the scalar weight is positive
away from zero and a nonzero polynomial cannot vanish on
an interval. Their finite moments retain the unscaled
sigma density of mass $\sqrt{2\pi}$.
The exact coefficient map $P(S)\mapsto P(c+iy)$ is
\[
 J_{hj}=\binom jh c^{j-h}(\sqrt{-1})^h,
 \quad 0\leq h\leq j<r,\qquad |\det J|=1.
 \tag{BHR4}
\]
Thus the full reference source form for
$|y|^{2(q+s)}|P(c+iy)|^2\sigma(y)dy$ is
$J^*\mathsf M_{q+s,r}J$, $s=0,l$.
This proves the precise coordinate morphism and its
determinant cancellation. It introduces no additional
$2lr\log q$ into the full-$y$ ratio $\rho_r$.

Here is the full comparison, with the original tails and
signs retained before estimation. Partition the line into
$I=\{|y|<q\epsilon\}$, $B=\{q\epsilon\leq|y|\leq64q\}$
and $F=\{|y|>64q\}$. For any of the four original or
reference forms $H$ define
\[
 t_H=\log\det\!\left(I_r+(H^B)^{-1/2}
                    (H^I+H^F)(H^B)^{-1/2}\right),\qquad
 \log\det H=\log\det H^B+t_H.
 \tag{BHR5}
\]
BRD22--38 bounds the two original tail matrices, and
CTR39--44 proves the same bounds for both reference
matrices by their actual Gamma/Legendre integrals.
Every eigenvalue of these positive matrices is at most
$\kappa_{I,r}+\kappa_F$, so in each case
$0\leq t_H\leq L_r^{\rm tail}$.

For clarity, the degree-$l$ density transport concerns the
original polynomial spaces $\mathcal P=\C[S]_{\leq r-1}$,
$f\mathcal P$, inside $\C[S]_{\leq r+l-1}$.
Their intersection is exactly
$f\C[S]_{\leq r-l-1}$, by division by the original
monic $f$. Its kernel and quotient maps are those in
the complete exact sequence CTR4.
On the bulk, the density logarithm for replacing
$|y|^{2q}$ by $|\chi(c+iy)|^2$ is
\[
 \psi_\chi(z)=\alpha_k^\chi/z^2+E_\chi^{\rm dens}(qz),
 \quad |E_\chi^{\rm dens}(qz)|\leq\delta_\chi,
 \quad\epsilon\leq|z|\leq64.
 \tag{BHR6}
\]
This is the absolutely convergent expansion of the full
root product, with the odd power sums cancelled by the
original root permutation. Its second coefficient is
$qk(k+2)(d-g)/3$ before the substitution $y=qz$.
The difference of the two actual orthogonal projectors
has inertia $(l,l,r-l)$, trace zero and trace norm at
most $2l$; CTR8--14 constructs its full residual frame
with its original cross Gram. Differentiating the two
restriction determinants along $e^{t\psi_\chi}$ gives
the trace of this projector difference against the
compressed multiplication by $\psi_\chi$. In its signed
eigenframe, each paired expectation differs by at most
$\operatorname{osc}\psi_\chi$ and the positive weights
sum to at most $l$. Integration therefore gives the
actual bulk-ratio error $d_\chi$ with
$|d_\chi|\leq l\operatorname{osc}\psi_\chi\leq E_\chi$.
All source coordinates, fixed inclusions and density
Jacobians in this application are exactly CTR22--29.

The remaining numerator factor satisfies on the bulk
$1\leq|f(c+iy)|^2/|y|^{2l}\leq e^{\Delta_f}$.
Integration on each original polynomial and multiplication
of the $r$ relative eigenvalues therefore gives a
one-sided correction $0\leq d_f\leq r\Delta_f$.
Denote the reference forms from BHR4 by $H_{\rm pow,s}$.
The determinant factorization BHR5 now proves
\[
\begin{aligned}
 e_r:=\log R_{q+r-1}-\rho_r
 ={}&d_\chi+d_f+t_{H_{\chi,f}}-t_{H_{\chi,0}}\\
    &-t_{H_{\rm pow,l}}+t_{H_{\rm pow,0}},\\
 -E_\chi-2L_r^{\rm tail}\leq e_r
 &\leq E_\chi+r\Delta_f+2L_r^{\rm tail}.
\end{aligned}\tag{BHR7}
\]
In particular the two numerator-density constants and all
four nonnegative full-line tail corrections have their
exact places in this equality. The upper-only
$r\Delta_f$ has not been reflected into its lower side.
This is the finite full-source CTR47--48 result returned
through the exact original relation injection.

Its quantitative error is $O(k/e)+o(1)=o(q)$ on the
original degree. Indeed CTR31 gives
\[
 |\alpha_k^\chi|\leq\frac{|d-g|}{3e},\qquad
 \delta_\chi\leq\frac{2(d+g)^2}{3e^3k^2\epsilon^4},\qquad
 r\Delta_f\leq\frac{k}{8e\epsilon^2},
 \tag{BHR8}
\]
using $r\leq2q$ and $S_2\leq k^3/16$.
The tail bound is at most
$r(\kappa_{I,r}+\kappa_F)$, exponentially small in $q$
times its displayed polynomial. Thus the source error
is $O(k)$ for $m=1$ and $O(1)$ for fixed $m\geq2$.
Every bound applies to the full original forms, not only
their restriction to the bulk.

The full moment inputs in BHR3 are also exact. Retain
$c_\sigma=\sqrt{2\pi}$ and let $c_0^{\rm mom}=1$,
\[
 c_n^{\rm mom}=-\sum_{j=1}^{n-1}\binom{2n-1}{2j-1}c_j^{\rm mom}
     -\frac12\sum_{j=0}^{n-1}\binom{2n-1}{2j}c_j^{\rm mom}.
 \tag{BHR9}
\]
The Euler-integral/Fourier calculation CTR50 proves the
original characteristic function
$c_\sigma(\cosh t)^{-1/2}$.
Its derivative equation
$h'\cosh t+h\sinh t/2=0$ proves BHR9 by comparing
the coefficient of $t^{2n-1}$. The Gamma exponential
bound permits differentiation of every finite order
under the integral. Hence
$\mu_{2n+1}=0$, $\mu_{2n}=c_\sigma(-1)^nc_n^{\rm mom}$
and $(\mathsf M_{a,r})_{ij}=\mu_{2a+i+j}$.
Every $c_n^{\rm mom}$ is rational by the finite recurrence.
Thus each ratio in BHR3 is the logarithm of a positive
rational determinant ratio, with $c_\sigma^r$ present
in both original determinants before cancellation.

The original low endpoints remain exactly $R_{q-1}=1$
and $R_q=\mathfrak r_\chi$ with its JSR24b coefficient
formula. Define the same complete low constant and the
full-line reference centre by
\[
\begin{gathered}
 C_{\rm rel}=\log D_{l,q-1}+\log D_{l,q}
 -\log\mathfrak r_\chi-\log D_{l,2q-1}-\log D_{l,2q},\\
 H_k^*=-C_{\rm rel}-\rho_q-\rho_{q+1},\qquad
 L_\Sigma=L_q^{\rm tail}+L_{q+1}^{\rm tail},\\
 b_k^-=2E_\chi+(2q+1)\Delta_f+2L_\Sigma,\qquad
 b_k^+=2E_\chi+2L_\Sigma,\\
 H_k^*-b_k^-\leq\mathcal H_k
       =\mathcal R_k^0-\mathcal R_k^\sigma=-\mathfrak T_k
                     \leq H_k^*+b_k^+.
\end{gathered}\tag{BHR10}
\]
This is a direct signed substitution, not a different
four-volume: JSR15 gives
$\mathcal H_k=H_k^*-e_q-e_{q+1}$.
Use the two upper sides of BHR7 for its lower endpoint
and the two lower sides for its upper endpoint.
The numerator factor $\Delta_f$ therefore occurs only
in $b_k^-$. Both $b_k^\pm=o(q)$ with the explicit
$O(k)$ and $O(1)$ branches just proved.

For comparison the prior BRD bulk calculation remains
an additional proved interval. With
\[
\begin{gathered}
 (K_{s,r})_{ij}=\int_{\epsilon\leq|z|\leq64}
 z^{i+j}|z|^{2q+2s}e^{\alpha_k^\chi/z^2}\sigma(qz)\,dz,\\
 H_k^B=-C_{\rm rel}-2l(2q+1)\log q
     -\sum_{r=q}^{q+1}\log\frac{\det K_{l,r}}{\det K_{0,r}},\\
 b_{B,k}^-=2(2q+1)\delta_\chi+(2q+1)\Delta_f+L_\Sigma,
 \quad b_{B,k}^+=2(2q+1)\delta_\chi+L_\Sigma,
\end{gathered}\tag{BHR11}
\]
BRD41a gives
$H_k^B-b_{B,k}^-\leq\mathcal H_k\leq H_k^B+b_{B,k}^+$.
The proof uses the same two high relation ranks and
the same reversal of all four signs. The explicit
$2l(2q+1)\log q$ is present in BHR11 because its
integrals use $y=qz$. It is already incorporated in
the full-$y$ determinants in BHR3 and must not be
added to BHR10 a second time.

\section{The complete universal low-and-high reference in the same receiver}

On exactly the finite original-parameter domain BHR1, retain
all constants, both high ranks, and the full original low ratio
$\mathfrak r_\chi$ from JSR24b. Write
$\mu_j=\int_{\R}y^j\sigma(y)\,dy$, with the literal power
of the original real coordinate as its index. In particular
$\mu_0=\sqrt{2\pi}$, $\mu_{2n+1}=0$, and every even
moment is strictly positive. LET uses $m_j=\mu_j$ for
this literal moment index; the PHT
even-moment sequence has $\mu_n^{\rm even}=m_{2n}=\mu_{2n}$.
These are exact index dictionaries for the same original
measure and coordinate, with no change of source mass.
The scalar source comparison in
the complete LET1--18 proof gives the following finite quantities:
\[
\begin{gathered}
 E_0=|\alpha_k^\chi|(\epsilon^{-2}-64^{-2})+2\delta_\chi,
 \qquad E_\chi=lE_0,\\
 L_1^{\rm tail}=\log(1+\kappa_{I,1}+\kappa_F),\quad
 \kappa_{I,1}=\frac{2\epsilon}{\sqrt{\cos1}}
               \left(\frac43\,2^{-13}e^{2\pi}\right)^q,\\
 a_0=E_0+2L_1^{\rm tail},\qquad b_0=a_0+\Delta_f,\\
 e_0:=\log\mathfrak r_\chi-
                 \log\frac{\mu_{2q+2l}}{\mu_{2q}}
                    \in[-a_0,b_0].
\end{gathered}\tag{WRC1}
\]
The symbols $a_0,b_0$ denote scalar low-source errors;
the baseline arithmetic volume remains $\mathcal B_k^0$.
The scalar tail estimate here is LET12--16, whose proof
works directly with the constant coefficient. No statement
requiring the high-rank condition $r\geq l$ is applied at
rank one.

For completeness, the exact scalar source maps are
$u\mapsto(iy)^q u$ and $u\mapsto\chi(c+iy)u$ on the
original bulk Hilbert space. Their multiplication bridge
is $h\mapsto\chi(c+iy)h/(iy)^q$, with its bounded inverse
on the stated bulk; it commutes with the same original
$f$ multiplication. If $w$ is the bulk power measure,
positivity of both scalar integrals gives
$|\log R_{e^\psi w}(|f|^2)-\log R_w(|f|^2)|
\leq\operatorname{osc}\psi\leq E_0$.
The original factor $|f|^2/|y|^{2l}$ contributes a
one-sided logarithm in $[0,\Delta_f]$.
All four full scalar integrals are their own bulk integrals
multiplied by $e^{t_H}$, where $0\leq t_H\leq L_1^{\rm tail}$.
Consequently the exact low error has the same four-tail
decomposition
\[
 e_0=d_\chi^{\rm low}+d_f^{\rm low}
  +t_{\chi,f}^{\rm low}-t_{\chi,0}^{\rm low}
  -t_{{\rm pow},l}^{\rm low}+t_{{\rm pow},0}^{\rm low}.
 \tag{WRC2}
\]
This proves both sides of WRC1 with their original signs.
The complete LET proof supplies every scalar inner and far
constant and its Gamma source bound, and is retained as a
whole provider.

Define the full universal reference by
\[
\begin{gathered}
 P_k=\log D_{l,q-1}+\log D_{l,q}
                     -\log D_{l,2q-1}-\log D_{l,2q},\\
 W_k=P_k-\log\frac{\mu_{2q+2l}}{\mu_{2q}}
                                  +\rho_q+\rho_{q+1},\\
 C_{\rm rel}=P_k-\log\mathfrak r_\chi,\qquad
 H_k^*=-W_k+e_0,\\
 \mathfrak T_k=W_k-e_0+e_q+e_{q+1},\qquad
 \mathcal H_k=-\mathfrak T_k=-W_k+e_0-e_q-e_{q+1}.
\end{gathered}\tag{WRC3}
\]
These are equalities obtained by substitution into the
existing original low/high determinant identity. The
two $\rho_r$ are the full-$y$ Hankel ratios of BHR3;
no additional $q$-coordinate power is inserted.
Every factor in $W_k$ depends only on the actual integers
$q,l$ and the same fixed Gamma density. Its completely
specified positive exponential is
\[
 e^{W_k}=
 \frac{D_{l,q-1}D_{l,q}}{D_{l,2q-1}D_{l,2q}}
 \frac{\mu_{2q}}{\mu_{2q+2l}}
 \frac{\det\mathsf M_{q+l,q}}{\det\mathsf M_{q,q}}
 \frac{\det\mathsf M_{q+l,q+1}}{\det\mathsf M_{q,q+1}}.
 \tag{WRC4}
\]
Each moment and determinant is computed by the complete
finite moment recurrence in CTR50--53, with all original
mass factors retained before cancellation. In particular
$\mathsf M_{a,r}=\sqrt{2\pi}\,\mathsf Q_{a,r}$ in the
original entries, and both determinants in each ratio
have the factor $(\sqrt{2\pi})^r$.
The same mass occurs in both scalar moments.
Thus WRC4 is an exact positive rational expression.
The complete PHT/HCT providers give its positive moment
and parity-Hankel expression. The complete RWB and GEL/EIQ
proofs now evaluate this exact expression: GRI1--2 below proves
$lq/64\leq W_k\leq6lq$ and
$W_k/(lq)\to\mathcal C_\Gamma>0$ with every original
low moment, D-product and high parity factor retained.

Use precisely the existing high-endpoint constants
$b_k^+=2E_\chi+2L_\Sigma$ and
$b_k^-=b_k^++(2q+1)\Delta_f$ from BHR10.
Since $e_q+e_{q+1}\in[-b_k^+,b_k^-]$, WRC1--3 prove
\[
 -W_k-b_k^--a_0\leq\mathcal H_k
                     \leq-W_k+b_k^++b_0.
 \tag{WRC5}
\]
The low error contributes $+e_0$, and the two high errors
contribute $-e_q-e_{q+1}$, as in WRC3. In particular
the positive multiplication term occurs in different
sides for the low subtraction and the two high terms.

This fourth interval is compared exactly with the retained
earlier $H_k^*$ interval, rather than assumed to improve it.
The typed scalar translation $H_k^*=-W_k+e_0$ proves
\[
 -W_k-b_k^--a_0\leq H_k^*-b_k^-,\qquad
 H_k^*+b_k^+\leq-W_k+b_k^++b_0.
 \tag{WRC6}
\]
Indeed their respective differences are $e_0+a_0\geq0$
and $b_0-e_0\geq0$. Thus the new interval contains the
previous exact-low interval at every stated finite input.
Its addition as a fourth maximum/minimum entry leaves
the existing exact current endpoints unchanged. Its gain
is the fully universal centre WRC4, whose growth is now
proved by RWB and GEL and returned in GRI1--7 below.
The exact low moment remains an input to the narrower
previous finite interval.

The same original arithmetic allowance now has the
additional universal-centred presentation
\[
\begin{gathered}
 \mathcal B_k^0-W_k-b_k^--a_0-a_k^-
 \leq\mathcal B_k^{\rm ar}
 \leq\mathcal B_k^0-W_k+b_k^++b_0+a_k^+,\\
 \mathcal B_k^0-W_k-b_k^--a_0-D_k^-
 \leq\mathcal B_k^{\rm ar}
 \leq\mathcal B_k^0-W_k+b_k^++b_0+D_k^+.
\end{gathered}\tag{WRC7}
\]
Add the unchanged intervals for $\delta_k^\sigma$ to
WRC5 and use $\mathcal B_k^{\rm ar}=\mathcal B_k^0+
\mathcal H_k+\delta_k^\sigma$. The original full Gamma
return therefore obeys
$\Delta_k^\Gamma+W_k=\delta_k^\sigma+e_0-e_q-e_{q+1}$.
Its arithmetic constants have not improved: the new
universal Gamma errors are still $o(q)$, being $O(k)$
for $m=1$ and $O(1)$ for fixed $m\geq2$; LET18 proves
$a_0,b_0=O(1)$ and $O(k^{-1})$ in those respective regimes.
Thus the already proved signed $q\log k$ constants for
$m\geq2$ and the original simple-zero $5/7$ rate on
$kq$ transport to this centre by direct addition.
Subtracting $4q\log(D_hk)$ from WRC7 and intersecting
its lower endpoint with zero gives the same original
nonnegative HC residual. No original arithmetic allowance,
boundary, nilpotent direction or source mass is changed.

\subsection{The evaluated original baseline and its finite arithmetic receiver}

The canonical source still has its original order $k=4l+1$,
$e=1+k(m-1)$, $q=e(k+1)^2=2n$, and centre $c=k/2$.
Its multiplying polynomial is the original $f_l$, and its full
baseline remains $\mathcal B_k^0$. The following calculation uses
the complete BSL1--22 and HBL1--18 proofs included below.

Write $\gamma_j=\sqrt{2\pi}(2j)!/4^j$ for the original monic
Gamma norms and $\mu_q^{\rm even}=\int_{\mathbb R}y^{2q}d\sigma(y)$.
This last scalar is BSL's $\mu_q$ and the literal-index moment
$m_{2q}$ of LET. With the complete square pushforward
$\nu=(y\mapsto y^2)_*\sigma$, retain
\[
\begin{gathered}
 H_s^{(a)}=\det\left[\int_0^\infty t^{a+i+j}d\nu(t)\right]_{i,j<s},
 \qquad H_0^{(a)}=1,\\
 S_q=-\log\gamma_q-2\sum_{j=q+1}^{2q-1}\log\gamma_j
                                             -\log\gamma_{2q},\\
 Z_q=H_n^{(q)}H_{n+1}^{(q)}(H_n^{(q+1)})^2,
 \qquad V_q=S_q+\log Z_q-\log\mu_q^{\rm even}.
\end{gathered}\tag{BRI1}
\]
The original relation Grams are
$A_{\chi,r}=\operatorname{Gram}_\sigma[\chi S^j]_{j<r}$;
the power Grams are
$(M_{q,r})_{ij}=\int y^{2q+i+j}d\sigma(y)$.
Set $\eta_r=\log\det A_{\chi,r}-\log\det M_{q,r}$ for
$r=1,q,q+1$. The exact remainder map $J_N$ and its monic
relation frame give $\det G_N^\sigma=\det H_N^\sigma/
\det A_{\chi,N+1-q}$ at the four original ranks $0,1,q,q+1$.
Evenness of $\sigma$ splits the two high power Grams into precisely
the four factors in $Z_q$, while their scalar low Gram is
$\mu_q^{\rm even}$. Thus
\[
 \mathcal B_k^\sigma=V_q+\eta_q+\eta_{q+1}-\eta_1,
 \qquad \mathcal H_k=\mathcal B_k^\sigma-\mathcal B_k^0.
 \tag{BRI2}
\]

The exact source transition here now has its original observation
inputs explicitly calculated. With $F_K^{(s)}$ and $F_B^{(s)}$
given by the section-corrected norm and period-column determinants
in MRI4a, the same equality is
\[
 \mathcal H_k
   =\bigl(F_K^{(1)}+F_B^{(1)}\bigr)
       -\bigl(F_K^{(k)}+F_B^{(k)}\bigr).
 \tag{BRI2a}
\]

The original transition now has a computed individual return
on the proved simple-quartet period domain:
\[
\begin{gathered}
 \mathcal H_k=(F_K^{(1)}-F_K^{(k)})+(F_B^{(1)}-F_B^{(k)}),\\
 m=1,\ |u|\ge R_*:\quad
 \frac{F_K^{(1)}-F_K^{(k)}}{kq}\to0,\qquad
 \frac{F_B^{(1)}-F_B^{(k)}}{kq}\to-\frac{\mathcal C_\Gamma}{4},\\
 \mathcal B_k^{\rm ar}=F_K^{\rm ar}+F_B^{\rm ar},\quad
 \frac{F_K^{\rm ar}-F_K^{(k)}}{kq}\to0,\qquad
 \frac{F_B^{\rm ar}-F_B^{(k)}}{kq}\to-\frac{\mathcal C_\Gamma}{4}.
\end{gathered}\tag{BRI2b}
\]
The full metric comparison, exact ranks and signed deficits
are proved at MRI1d, MRI3b and MRI5a--d below. They give
$O_h(k^2\log(q+1))$ for the kernel Gamma difference and
$O_h(k^2)$ for the arithmetic-to-fixed-Gamma kernel difference.
Their division by $kq$ tends to zero since $q=(k+1)^2$.
Subtracting them from the original total transitions
$\mathcal H_k$ and $\mathcal H_k+\delta_k^\sigma$ proves
BRI2b with the stated negative sign. All finite intervals
for $\mathcal H_k$ and $\delta_k^\sigma$ are retained.
Outside the displayed specialization, the full-rank finite
maps and intervals MRI3b and MRI5a--c retain general $m$.

For each source, MRI1a identifies its observed columns with
the original $\Lambda$, MRI2a transports the retained kernel
through the preceding-degree minimum section, and MRI3a keeps
the original two factors. Their product gives the original
quotient increment, so the four-endpoint sum is exactly
$\mathcal B_k^{(1)}=\mathcal B_k^\sigma$ or
$\mathcal B_k^{(k)}=\mathcal B_k^0$. Substitution in BRI2
proves BRI2a, including its negative tensor-source sign.

The second equality is the exact original source transition
$\mathcal B_k^0=\mathcal B_k^\sigma+\mathfrak T_k$ and
$\mathcal H_k=-\mathfrak T_k$. It identifies the scalar to
which the baseline bounds must be returned.

On the original full-source threshold
$k\geq2048\sqrt{d+g}$, retain the BSL constants
\[
\begin{gathered}
 a_k=\frac{k(k+2)(d-g)}{3q}<0,\quad
 t_\epsilon=\frac{k^2(d+g)}{q^2\epsilon^2}\leq\tfrac14,
 \quad d_\chi=\frac{qt_\epsilon^2}{2(1-t_\epsilon)},\\
 b_-=a_k\epsilon^{-2}-d_\chi,
 \qquad b_+=a_k64^{-2}+d_\chi,
 \quad L_*=L_q^{\rm tail}+L_{q+1}^{\rm tail}+L_1^{\rm tail},\\
 \mathsf L_k^\sigma=V_q+(2q+1)b_- -b_+-L_*,\quad
 \mathsf U_k^\sigma=V_q+(2q+1)b_+ -b_-+L_*,\\
 \mathsf L_k^\sigma\leq\mathcal B_k^\sigma
                                  \leq\mathsf U_k^\sigma.
\end{gathered}\tag{BRI3}
\]
Here $d=\delta^2$, $g=\gamma^2$, $\epsilon=2^{-10}$ and
all three $L_r^{\rm tail}$ are the original complete tail constants
already displayed above; $b_-,b_+$ are the packet density bounds,
while $b_k^-,b_k^+$ keep their different original high-error meaning.
The exact multiplication map between the bulk sources is
$P(c+iy)(iy)^q\mapsto P(c+iy)\chi(c+iy)$.
Its squared density ratio is between $\exp(b_-)$ and
$\exp(b_+)$, by the full paired-root logarithm in BSL6--8.
Congruence on the same relation coefficient space and both positive
full-source tails gives $rb_- -L_r^{\rm tail}\leq\eta_r
\leq rb_++L_r^{\rm tail}$. Taking the two high positive signs
and the scalar low negative sign proves BRI3 from BRI2.

By the second equality in BRI2 the interval
$[\mathsf L_k^\sigma-\mathcal B_k^0,
\mathsf U_k^\sigma-\mathcal B_k^0]$ contains the same original
$\mathcal H_k$ as every previously retained Gamma interval.
It can therefore be included in their exact maximum/minimum
intersection. No change is made to any original source or
arithmetic allowance. Independently of which finite interval is
tightest, the original equality
$\mathcal B_k^{\rm ar}=\mathcal B_k^\sigma+\delta_k^\sigma$
gives the explicit finite baseline receiver
\[
\begin{gathered}
 \mathsf L_k^\sigma-a_k^-\leq\mathcal B_k^{\rm ar}
                                   \leq\mathsf U_k^\sigma+a_k^+,\\
 \mathsf L_k^\sigma-D_k^-\leq\mathcal B_k^{\rm ar}
                                   \leq\mathsf U_k^\sigma+D_k^+,\\
 \max\{0,\mathsf L_k^\sigma-D_k^- -4q\log(D_hk)\}
 \leq\mathcal R_k
 \leq\mathsf U_k^\sigma+D_k^+ -4q\log(D_hk).
\end{gathered}\tag{BRI4}
\]
The first line adds the unchanged
$\delta_k^\sigma\in[-a_k^-,a_k^+]$ to BRI3; the second
uses $a_k^\pm\leq D_k^\pm$. The last subtracts the unchanged
HC threshold and intersects the lower endpoint with the original
$\mathcal R_k\geq0$, on the combined original source and HC domains.

The complete HBL13--18 coefficient is
\[
 C_B=9-8\log2+\mathcal F(2,\pi),\qquad
 C_B>\frac{769}{17010}>0,\qquad
 \frac{V_q}{q^2}\longrightarrow C_B.
 \tag{BRI5}
\]
Here $\mathcal F(2,\pi)$ is the supremum of the single-kernel
energy for $\pi\sqrt x-2\log x$, attained by the exact positive
density $\rho_{2,\pi}$ already specified in GRI2. HBL14--16
gives its endpoint-integral expression, and HBL17--18 proves
the strict rational lower bound by its explicit arcsine competitor.
The determinant calculation retains the complete factors
\[
\begin{split}
 S_q&=-6q^2\log q+(9-8\log2)q^2+O(q\log q),\\
 \log Z_q&=(6q^2+2q)\log q+q^2\mathcal F(2,\pi)+o(q^2),\qquad
 \log\mu_q^{\rm even}=O(q\log q).
\end{split}
\]
Their $q^2\log q$ terms cancel only in BRI1. BSL6--10 gives
$\eta_q+\eta_{q+1}-\eta_1=O(q/e)=o(q^2)$, while BRI3
retains its finite asymmetric enclosure.

Consequently the full original baselines, arithmetic volume and
HC residual have the evaluated common leading value
\[
\begin{gathered}
 \frac{\mathcal B_k^\sigma}{q^2}\longrightarrow C_B,\qquad
 \frac{\mathcal B_k^0}{q^2}\longrightarrow C_B,\qquad
 \frac{\mathcal B_k^{\rm ar}}{q^2}\longrightarrow C_B,\qquad
 \frac{\mathcal R_k}{q^2}\longrightarrow C_B,\\
 \frac{\mathcal B_k^0-\mathcal B_k^{\rm ar}}{kq}
                \longrightarrow\frac{\mathcal C_\Gamma}{4},
 \qquad \frac{\mathcal B_k^0}{kq}\longrightarrow+\infty.
\end{gathered}\tag{BRI6}
\]

For either original source order $s=1,k$, the first-line baseline
limit also has the following fully specified mixed presentation:
\[
\begin{split}
 \frac1{q^2}\left\{
 \mathcal W_q\left(\log\left[1+
   \frac{(\vartheta_{q+j,j}^{(s)})^*
        H_{q+j-1}^{(s),K}\vartheta_{q+j,j}^{(s)}}{1+\xi_j}
                                      \right]\right)
 +\log\frac{\det\mathcal M_{2q}^{(s)}\det\mathcal M_{2q+1}^{(s)}}
                 {\det\mathcal M_q^{(s)}\det\mathcal M_{q+1}^{(s)}}
                          \right\}\longrightarrow C_B.
\end{split}\tag{BRI6a}
\]

On the evaluated original period domain this combined
coefficient is now resolved into the actual individual terms:
\[
 m=1,\ |u|\ge R_*:\quad
 \frac{F_K^{(1)}}{q^2},\frac{F_K^{(k)}}{q^2},
 \frac{F_K^{\rm ar}}{q^2}\to0,\qquad
 \frac{F_B^{(1)}}{q^2},\frac{F_B^{(k)}}{q^2},
 \frac{F_B^{\rm ar}}{q^2}\to C_B.
 \tag{BRI6b}
\]
MRI1d computes $r=8k-16$ from the original full period
orbit. MRI4b--c proves the finite kernel upper bound
$2r\log Z_q+\widehat C_K^+=O_h(kq\log(q+1))$ and the
exact complement for the arithmetic boundary. Division by
$q^2=(k+1)^4$ makes this error tend to zero. The already
proved baseline limits therefore give BRI6b, uniformly in
the stated period domain and every logarithm branch.
The finer component differences are BRI2b and MRI5d;
the absolute common kernel divided by $kq$ is retained
as the next quantity to calculate from the original metrics.

% Replacement for the current BRI6c receiving paragraph.
The original absolute common kernel now has the sharper finite
control MRI4e--f, whose complete proof is AKS20--32 and AKJ1--4.
The all-multiplicity factorial bound MRI4d remains valid. On the
actual simple-quartet domain $m=1$, $|u|\ge R_*$, one has
\[
 \limsup\frac{F_K^a}{kq}
 \le\kappa_*=32\log\frac{25\log3}{4\pi},\qquad
 a\in\{\sigma,0,\mathrm{ar}\}.
 \tag{BRI6c}
\]
The finite estimate is the pointwise minimum of the original KAF
bound and the new confluent-Newton bound, retaining the first
compressed increment. MRI4f gives the correlated arithmetic kernel
interval $[L_{\rm ar},U_{\rm ar}]$ and its exact boundary complement.
The original signed mixed return has the exact splitting identity
\[
 F_K^{\rm ar}+F_B^{\rm ar}-F_B^{(k)}
   =F_K^{(k)}+\mathcal H_k+\delta_k^\sigma.
 \tag{BRI6d}
\]
Thus its finite lower and upper bounds are precisely
$L_k+\mathcal H_k+\delta_k^\sigma$ and
$U_k+\mathcal H_k+\delta_k^\sigma$, further intersected with the
unchanged source-difference intervals of AKS41--48. The identity
follows by subtracting
$F_K^{(k)}+F_B^{(k)}=\mathcal B_k^0$ from
$F_K^{\rm ar}+F_B^{\rm ar}=\mathcal B_k^{\rm ar}
 =\mathcal B_k^0+\mathcal H_k+\delta_k^\sigma$.
In particular its lower limiting coefficient is at least
$-\mathcal C_\Gamma/4$ and its upper limiting coefficient is at
most $\kappa_*-\mathcal C_\Gamma/4$, using the full proved signed
Gamma return and arithmetic error. The sourcewise differences
BRI2b and MRI5d are unchanged. The next calculation remains the
actual common kernel's compressed determinant with its original
period dependence.


\paragraph{The current conductor correction in the signed return.}
MRI4g--k now controls the actual conductor transport with its explicit
error $\epsilon_k^{(s)}=o(kq)$ for the fixed original period.
Substituting the exact MRI4j into the original BRI6d identity yields
\[
 \begin{split}
 F_K^{\rm ar}+F_B^{\rm ar}-F_B^{(k)}
 &=\mathcal T_k^{(k)}-F_{\Xi,k}^{(k)}
       +\mathcal H_k+\delta_k^\sigma+\mathcal E_{A,k}^{(k)},\\
 |\mathcal E_{A,k}^{(k)}|&\le\epsilon_k^{(k)}=o(kq).
 \end{split}\tag{BRI6e}
\]
The unchanged total identity proves the substitution exactly;
CEP46--49 proves the error estimate on its actual compressed
covariance. Thus the current finite interval has centre
$\mathcal T_k^{(k)}-F_{\Xi,k}^{(k)}+\mathcal H_k+\delta_k^\sigma$
and radius $\epsilon_k^{(k)}$, intersected with the preceding
intervals for this same signed return obtained from BRI, AKS and
the correlated arithmetic bounds through the unchanged identities.
The original Gamma return
and arithmetic scalar keep their exact signs and bounds.
The remaining shifted lower-grid determinants and rank-$8k-32$
observation quotient are explicit original matrices. The conductor
inverse cost has been controlled and is no longer an unevaluated
leading-scale error in this formula.


\paragraph{The outer-root and actual residual signed return.}
Substitute the exact MRI4m at source order $s=k$ into BRI6d.
With the complete positive residual sum MRI4l and the original
Gamma return and arithmetic scalar, the result is
\[
 \begin{gathered}
 F_K^{\rm ar}+F_B^{\rm ar}-F_B^{(k)}
   =\mathcal F_{\partial,k}^{(k)}-\Phi_{k,k}
        +\mathcal H_k+\delta_k^\sigma+\mathfrak e_{k,k},\\
 -e_{k,k}^-\le\mathfrak e_{k,k}\le e_{k,k}^+,\qquad
 e_{k,k}^\pm=o(kq).
 \end{gathered}\tag{BRI6f}
\]
OAR1--3 proves these finite error endpoints with every original
tail sign, and OAR11 proves the displayed substitution. Adding
$\mathcal H_k+\delta_k^\sigma$ to MRI4n gives the corresponding
finite interval for this same signed return. The more precise
full-sum interval in MRI4m is available with its actual residual
coefficients. Both quotient returns retain their original metrics
and all four cutoffs; their leading difference remains the next
explicit calculation.


\paragraph{The same-source moment formula in the signed return.}
On the explicit finite comparison domain of MRI4p, the original
signed-return identity further gives
\[
 \begin{gathered}
 F_K^{\rm ar}+F_B^{\rm ar}-F_B^{(k)}
   =\mathcal F_{0,k}^{(k)}-\Phi_{k,k}
       +\mathcal H_k+\delta_k^\sigma+\mathfrak e^0_{k,k},\\
 -e_{k,k}^{0,-}\le\mathfrak e^0_{k,k}\le e_{k,k}^{0,+},\qquad
 e_{k,k}^{0,\pm}=o(kq).
 \end{gathered}\tag{BRI6g}
\]
PMR7 proves this by substitution of the source-$k$ identity in
MRI4p into BRI6d, preserving all original signs. Adding
$\mathcal H_k+\delta_k^\sigma$ to MRI4q's source-$k$ interval
gives its finite signed interval. Every moment determinant in
$\mathcal F_{0,k}^{(k)}$ uses the original source order $k$,
and the actual residual observation remains in $\Phi_{k,k}$.


\paragraph{The evaluated outer term in the original signed return.}
On MRI4r's finite domain, OVR7 proves the exact refinement
\[
 \begin{gathered}
 \mathfrak S_k^{\rm mix}
   :=F_K^{\rm ar}+F_B^{\rm ar}-F_B^{(k)},\\
 \mathfrak S_k^{\rm mix}+\Phi_{k,k}
   =\Delta q\mathfrak C_\partial+\mathcal H_k
                     +\delta_k^\sigma+R_{k,k},\\
 b_{k,k}-e_{k,k}^-\le R_{k,k}\le b_{k,k}+e_{k,k}^+.
 \end{gathered}\tag{BRI6h}
\]
This substitutes MRI4s into the exact BRI6d identity; the original
Gamma return, arithmetic scalar and both finite error signs remain.
AMT48--49 gives
$(\mathcal H_k+\delta_k^\sigma)/(kq)\to-\mathcal C_\Gamma/4$.
Consequently OVR8 proves
\[
 \frac{\mathfrak S_k^{\rm mix}+\Phi_{k,k}}{kq}
       \longrightarrow16\mathfrak C_\partial-\mathcal C_\Gamma/4,
 \qquad
 \frac{\mathfrak S_k^{\rm mix}+\Phi_{k,1}}{kq}
       \longrightarrow16\mathfrak C_\partial-\mathcal C_\Gamma/4.
 \tag{BRI6i}
\]
The second limit also uses the proved source-transfer identity OVR4.
Both constants are the explicit original EIQ integrals retained in
OVR1 and OVR7--8. The complete actual residual Gamma quotient remains
in $\Phi$ through ROQ and the exact fixed-frame congruence MRI4v;
its leading value is the current analytic calculation.


The definitions of $\xi_j,\mathcal M_D^{(s)}$ and
$\vartheta_{D,j}^{(s)}$ are the exact original formulas MRI1a
and MRI2a. Equation MRI4a identifies the braces with
$F_K^{(s)}+F_B^{(s)}=\mathcal B_k^{(s)}$ before division by
$q^2$. The already proved BRI6 limit therefore proves BRI6a.
For general multiplicity the displayed coefficient remains on the
complete mixed sum. On the proved simple-quartet period domain,
BRI6b and MRI4c now evaluate all three actual individual boundary
coefficients and prove the vanishing kernel coefficients at this scale.

To prove these statements, BRI2--5 first gives the fixed Gamma
baseline. The original tensor return is
$\mathcal B_k^0=\mathcal B_k^\sigma+\mathfrak T_k$ with
$\mathfrak T_k/(lq)\to\mathcal C_\Gamma$ and $l/q\to0$.
The arithmetic return is
$\mathcal B_k^{\rm ar}=\mathcal B_k^\sigma+\delta_k^\sigma$,
where the retained simple and multiple branches prove
$\delta_k^\sigma/(kq)\to0$ and $k/q\to0$.
The original threshold satisfies $4\log(D_hk)/q\to0$.
These give the first line. The finer difference in the second line
is GRI5--7, unchanged. Finally $q/k\to\infty$ and $C_B>0$
prove the displayed divergence. In particular the previous necessary
lower bound at the $kq$ scale remains valid and is now strengthened
by the evaluated $q^2$ baseline; no finer remainder is assumed.

The original HC identity retains its actual components:
\[
 \mathcal B_k^{\rm ar}=4q\log\mathcal L_k-W_k^\sigma
                  +\mathcal S_k-\tau_k^{\rm win},\qquad
 \mathcal L_k=2\delta e(k+1)\lfloor(k+1)^2/4\rfloor.
 \tag{BRI7}
\]
The symbol $W_k^\sigma$ here is the original factorial term of
TWA5, and is distinct from the universal Gamma centre $W_k$.
Its complete formula gives $O(q\log q)$; TWA6 gives
$\tau_k^{\rm win}=O_h(k+\log q)$. Also
$\log\mathcal L_k=O_h(\log k)$ and $\log q=O_m(\log k)$.
Dividing the exact identity by $q^2$ and using BRI6 proves
\[
 \frac{\mathcal S_k}{q^2}\longrightarrow C_B>0.
 \tag{BRI8}
\]
Thus the original contraction, phase and relative-control sum has
an evaluated cost on this scale, with its individual terms retained.

Finally let $K_{\rm ker}=\ker\Lambda$ and retain the original
coefficient section $S_*:B_{\rm obs}\to E_\chi$.
For each of the same positive quotient metrics, the full frame map
$[I_{K_{\rm ker}},S_*]:K_{\rm ker}\oplus B_{\rm obs}\to E_\chi$
is an isomorphism, and the exact block determinant gives
\[
 \det(I_{K_{\rm ker}}^*G_NI_{K_{\rm ker}})\det Q_N^{B_{\rm obs}}
 =|\det[I_{K_{\rm ker}},S_*]|^2\det G_N.
 \tag{BRI9}
\]

For the original observation this same block identity has the
explicit boundary metric
\[
\begin{gathered}
 Q_N^{B_{\rm obs},(s)}=(\mathcal M_{N+1}^{(s)})^{-1},\qquad
 \det(I_{K_{\rm ker}}^*G_N^{(s)}I_{K_{\rm ker}})
 =|\det[I_{K_{\rm ker}},S_*]|^2\det G_N^{(s)}
                                  \det\mathcal M_{N+1}^{(s)},
 \quad N\ge q-1.
\end{gathered}\tag{BRI9a}
\]

For the original arithmetic source the same coefficient-frame
calculation has the complete evaluated Gram map
\[
\begin{gathered}
 \mathcal M_{N+1}^{\rm ar}:=
       \Lambda J_N(H_N^{\rm ar})^{-1}J_N^*\Lambda^*,\quad
 Q_N^{\rm ar}=(\mathcal M_{N+1}^{\rm ar})^{-1},\\
 \det H_{K,N}^{\rm ar}
 =|\det[I,S_*]|^2\det G_N^{\rm ar}
                                    \det\mathcal M_{N+1}^{\rm ar},\\
 F_B^{\rm ar}=\log\frac{
 \det\mathcal M_{2q}^{\rm ar}\det\mathcal M_{2q+1}^{\rm ar}}
 {\det\mathcal M_q^{\rm ar}\det\mathcal M_{q+1}^{\rm ar}}.
\end{gathered}\tag{BRI9b}
\]
The first line applies the unchanged remainder and observation
to the inverse original source Gram. Its positivity follows
from surjectivity and the full positive arithmetic source.
The minimum section in MRI3b is orthogonal to the actual
kernel; its frame differs from $[I,S_*]$ by a block triangular
matrix with identity diagonal. Taking its determinant proves
the second line. The same factor occurs at all four degrees,
so applying their original signs proves the third line.
Thus BRI6b and the signed component return BRI2b concern
these literal arithmetic and period-vector Grams, with all
off-diagonal entries and the original action defect retained.

Indeed the complete polynomial remainder map in orthonormal
source coordinates has columns $[u_0],\ldots,[u_N]$.
Its quotient inverse Gram is their column Gram, so applying
the original $\Lambda$ and its adjoint gives precisely
$\mathcal M_{N+1}^{(s)}$ by MRI1a. Inverting on the actual
image proves the first equality. Substituting it in BRI9
proves the second equality, retaining the original fixed
coefficient-section determinant. The four signs
$+,+,-,-$ cancel that same constant factor and give the
boundary ratio in MRI4a; the degree-dependent minimum-section
correction used for the individual kernel rows remains MRI2a.
Thus BRI9 and BRI6a refer to exactly the same full mixed sum,
with every original primary unit and period present.

The inverse splitting is $z\mapsto(z-S_*\Lambda z,\Lambda z)$;
orthogonal block elimination proves the determinant identity.
The frame factor cancels with signs $+,+,-,-$ at the same four
endpoints. BSL20--21 therefore transports the same result to the
actual combined mixed observation:
$(F_K(1)+F_B(1))/q^2\to C_B$. Each individual kernel and
boundary Gram, its off-diagonal entries and its Taylor unit
remain present. BRI6b and MRI4c give the proved individual
coefficients on the stated simple-quartet period domain, with the
complete source and actual period rank retained.


\subsection{The refined original low endpoint in the same complete return}

Keep the original $k=4l+1$, $q=e(4l+2)^2$ and both original
high ranks. The complete GDR1--22 and LER1--23 proofs give,
with all moment indices retained,
\[
 \log\frac{m_{2q+2l}}{m_{2q}}
 =F_{q,l}-K_{q,l}+r_{q,l}+\varepsilon_{q,l},\qquad
 -T_{q,l}\leq r_{q,l}\leq0,\quad
 |\varepsilon_{q,l}|\leq E_{q,l},
 \tag{LRI1}
\]
where the complete finite quantities are
\[
\begin{gathered}
 F_{q,l}=2l\log\frac{4q}{\pi}+\frac{l^2}{q}
 -\frac{l(16l^2-1)}{48q^2}
 +\frac{l^2(8l^2-1)}{48q^3},\\
 K_{q,l}=\frac{\pi^2}{64}\left\{
 \frac1{(2q-3/2)(2q-1/2)}
 -\frac1{(2q+2l-3/2)(2q+2l-1/2)}\right\}>0,\\
 T_{q,l}=\frac{l(768l^4-160l^2+7)}{7680q^4},\qquad
 E_{q,l}=\frac{51l}{16q^5}.
\end{gathered}\tag{LRI2}
\]
The literal-index ratio here equals LER's even-index ratio
$\mu_{q+l}/\mu_q$; the whole mass $\sqrt{2\pi}$ and both
half-axes remain present before cancellation. GDR proves the
global density identity with its explicit remainder, and LER3--14
integrates and telescopes it at every original integer exponent.
The exact cubic logarithm identity then gives the one-sided
$r_{q,l}$ in LER18. Thus no formal asymptotic expansion is
substituted for the finite equality LRI1.

The paired high determinants remain exactly
$\rho_q+\rho_{q+1}=\log(Z_{q+l}/Z_q)$, with the original
$Z_a=H_n^{(a)}H_{n+1}^{(a)}(H_n^{(a+1)})^2$ and $n=q/2$.
Define the refined, still complete high-determinant expression
\[
 \widehat W_k=P_k+\rho_q+\rho_{q+1}-F_{q,l}+K_{q,l}.
\]
Substitution in the original definition of $W_k$ gives
\[
 W_k=\widehat W_k-r_{q,l}-\varepsilon_{q,l},\qquad
 \widehat W_k-E_{q,l}\leq W_k
                  \leq\widehat W_k+T_{q,l}+E_{q,l}.
 \tag{LRI3}
\]
In particular the rational Gamma correction enters $W_k$
positively and the $l^2/q$ term enters negatively. Both signs
come from the original low-endpoint subtraction.

Use the same source errors $e_0\in[-a_0,b_0]$ and
$e_q+e_{q+1}\in[-b_k^+,b_k^-]$ in
$\mathcal H_k=-W_k+e_0-e_q-e_{q+1}$. This proves
\[
\begin{gathered}
 -\widehat W_k-T_{q,l}-E_{q,l}-b_k^--a_0
 \leq\mathcal H_k
 \leq-\widehat W_k+E_{q,l}+b_k^++b_0,\\
 \mathcal B_k^0-\widehat W_k-T_{q,l}-E_{q,l}-b_k^--a_0-a_k^-
 \leq\mathcal B_k^{\rm ar}\\
 \leq\mathcal B_k^0-\widehat W_k+E_{q,l}+b_k^++b_0+a_k^+.
\end{gathered}\tag{LRI4}
\]
For the second enclosure add the unchanged arithmetic interval;
its established outer version replaces $a_k^\pm$ by $D_k^\pm$.
Subtracting the original HC threshold and imposing the original
nonnegative lower constraint gives its corresponding HC enclosure.
For fixed exact $W_k$, LRI3 proves that LRI4 contains the earlier
WRC5 interval, so it does not remove the retained tighter finite
intersection. Its new content is the exact original low correction
and its finite signed error inside the complete high-determinant
expression. No high determinant or source unit has been removed.

Finally LER23 proves
$\log(m_{2q+2l}/m_{2q})-2l\log(4q/\pi)-l^2/q\to0$.
The term $l^2/q$ tends to $1/16$ at $m=1$ and to zero for
each fixed $m\geq2$. These are two original multiplicity
branches of this scalar endpoint, retained alongside the full
$W_k$ and both of its high relation ranks.


\subsection{The finite four-block estimate returned to the original signed maps}

The complete QLG1--32 and QGT1--28 proofs retain the original
$k=4l+1$, $q=[1+k(m-1)](k+1)^2=2n$ and the original Gamma
density. QGT15 is the exact sum $\mathcal C_{q,l}$ of the four
free-energy increments at ranks $n,n+1,n,n$ with multiplicities
$1,1,1,1$ and their actual quarter-shifted parameters. With its
proved finite constants $A=M/2+2C$, define
\[
\begin{gathered}
 \mathcal U_{q,l}=Al(3\sqrt n+\sqrt{n+1}),\\
 \mathcal V_{q,l}=\widetilde P_{q,l}+\mathcal C_{q,l}
                         +2l\log q-F_{q,l}+K_{q,l},\\
 \omega_{q,l}^-=\mathcal V_{q,l}-\mathcal U_{q,l}-E_{q,l},
 \qquad
 \omega_{q,l}^+=\mathcal V_{q,l}
                        +\mathcal U_{q,l}+T_{q,l}+E_{q,l}.
\end{gathered}\tag{QRI1}
\]
Here $\widetilde P_{q,l}=P_k+4lq\log q$ is the exact original
finite source sum, while $F,K,T,E$ are precisely LRI2. The
quantile-cell and circle-potential proof QLG, transported through
the exact square pushforward QGT3--11, gives a finite $Cs$ error
at each determinant rank $s$. The convex secant proof QGT12--17
then gives the exact full high ratio
\[
 \log(Z_{q+l}/Z_q)=(4q+2)l\log q+\mathcal C_{q,l}+\xi_{q,l},
 \qquad |\xi_{q,l}|\le\mathcal U_{q,l}.
\]
Substituting LRI1, with its original negative low-endpoint sign,
gives
$W_k=\mathcal V_{q,l}+\xi_{q,l}-r_{q,l}-\varepsilon_{q,l}$.
Thus $\omega_{q,l}^-\le W_k\le\omega_{q,l}^+$.

There is a second finite enclosure with no finite source sum in
its centre. Keep the explicit $p^{\rm c}_{q,l},p^-_{q,l},p^+_{q,l}$
and $\mathcal V^{\rm c}_{q,l}$ of QGT21, including all their
linear, quadratic and cubic product corrections. Set
\[
\begin{gathered}
 \omega_{q,l}^{\rm c,-}=\mathcal V^{\rm c}_{q,l}
                    -p^-_{q,l}-\mathcal U_{q,l}-E_{q,l},\\
 \omega_{q,l}^{\rm c,+}=\mathcal V^{\rm c}_{q,l}
                    +p^+_{q,l}+\mathcal U_{q,l}+T_{q,l}+E_{q,l}.
\end{gathered}\tag{QRI2}
\]
The proved original product interval
$-p^-_{q,l}\le\widetilde P_{q,l}-p^{\rm c}_{q,l}\le p^+_{q,l}$
implies
$\omega_{q,l}^{\rm c,-}\le\omega_{q,l}^-\le W_k
 \le\omega_{q,l}^+\le\omega_{q,l}^{\rm c,+}$.
This also proves the finite containment relation between these
two evaluated presentations.

For either pair $(\omega^-,\omega^+)$ in QRI1 or QRI2 the actual
source identity, with all original errors, is
\[
\begin{gathered}
 \mathcal H_k=-\mathfrak T_k=-W_k+e_0-e_q-e_{q+1},\\
 e_0\in[-a_0,b_0],\qquad
 e_q+e_{q+1}\in[-b_k^+,b_k^-],\\
 -\omega^+-b_k^--a_0\le\mathcal H_k
                            \le-\omega^-+b_k^++b_0.
\end{gathered}\tag{QRI3}
\]
The lower bound uses $-W_k\ge-\omega^+$, $e_0\ge-a_0$ and
$-(e_q+e_{q+1})\ge-b_k^-$; the upper bound uses the other three
endpoints. Thus every sign is fixed by the original map. This
interval contains the exact-$W_k$ interval WRC5. Its intersection
with the previously retained finite intervals therefore preserves
every previously proved tighter bound.

Adding the unchanged arithmetic scalar and then using its original
majorants gives, for either of the same pairs,
\[
\begin{split}
 \mathcal B_k^0-\omega^+-b_k^--a_0-a_k^-
 &\le\mathcal B_k^{\rm ar}
 \le\mathcal B_k^0-\omega^-+b_k^++b_0+a_k^+,\\
 \mathcal B_k^0-\omega^+-b_k^--a_0-D_k^-
 &\le\mathcal B_k^{\rm ar}
 \le\mathcal B_k^0-\omega^-+b_k^++b_0+D_k^+.
\end{split}\tag{QRI4}
\]
Indeed $\mathcal B_k^{\rm ar}=\mathcal B_k^0+\mathcal H_k
+\delta_k^\sigma$, $\delta_k^\sigma\in[-a_k^-,a_k^+]$ and
$0\le a_k^\pm\le D_k^\pm$. The original source and arithmetic
domains are intersected; no original degree or multiplicity
restriction is weakened. On the original HC domain the same map
$\mathcal R_k=\mathcal B_k^{\rm ar}-4q\log(D_hk)$ and the
proved $\mathcal R_k\ge0$ give
\[
\begin{split}
 \max\{0,\mathcal B_k^0-\omega^+-b_k^--a_0-D_k^-
                              -4q\log(D_hk)\}
 &\le\mathcal R_k\\
 &\le\mathcal B_k^0-\omega^-+b_k^++b_0+D_k^+
                              -4q\log(D_hk).
\end{split}\tag{QRI5}
\]

For the quantitative return around the previously evaluated coefficient,
let $\mathcal D_{q,l}$ be the full nonnegative finite expression QGT25
and write $\mathcal E_{q,l}^{\rm high}=\mathcal U_{q,l}+\mathcal D_{q,l}$.
QGT26 and the same source errors prove
\[
\begin{gathered}
 |W_k-lq\mathcal C_\Gamma|\le\mathcal E_{q,l}^{\rm high},\\
 -\mathcal E_{q,l}^{\rm high}-b_k^--a_0
 \le \mathcal H_k+lq\mathcal C_\Gamma
 \le \mathcal E_{q,l}^{\rm high}+b_k^++b_0.
\end{gathered}\tag{QRI6}
\]
The finite expression QGT25 and the exact packet relation give
$\mathcal E_{q,l}^{\rm high}=O(l\sqrt q+l^2)$.
The already proved original source errors are $O(k)$ for $m=1$
and $O(1)$ for each fixed $m\ge2$, and hence are $o(q)$ in both
branches. For each fixed $m\ge2$, $l/\sqrt q\to0$ and
$l^2/q\to0$; for $m=1$, $q=(4l+2)^2$ and both ratios remain
bounded. Dividing QRI6 by $q$ therefore proves
\[
 \begin{cases}
 (\mathcal H_k+lq\mathcal C_\Gamma)/q\longrightarrow0,
                   &m\ge2\text{ fixed},\\
 \mathcal H_k+lq\mathcal C_\Gamma=O(q),&m=1.
 \end{cases}
 \tag{QRI7}
\]
The first statement is equivalently
$(\mathfrak T_k-lq\mathcal C_\Gamma)/q\to0$ for the same fixed
multiple branch. The second keeps the entire finite uncertainty
in QRI3 and QRI6 and assigns no unproved limit to that quotient.

For the original arithmetic return the exact counterpart is
\[
\begin{gathered}
 \Delta_k^\Gamma+lq\mathcal C_\Gamma-\delta_k^\sigma
       =\mathcal H_k+lq\mathcal C_\Gamma,\\
 \mathcal B_k^{\rm ar}-\mathcal B_k^0+lq\mathcal C_\Gamma
       -\delta_k^\sigma=\mathcal H_k+lq\mathcal C_\Gamma,\\
 \mathcal R_k-\mathcal B_k^0+4q\log(D_hk)
       +lq\mathcal C_\Gamma-\delta_k^\sigma
       =\mathcal H_k+lq\mathcal C_\Gamma.
\end{gathered}\tag{QRI8}
\]
Each equality is the original definition followed by addition of
the displayed identical quantities. Each right-hand side obeys
QRI6 and both branches of QRI7. The arithmetic scalar remains
on the left with its full asymmetric bounds; in the multiple
branch the earlier $q\log k$ constants $64m+245$ and $128m+490$
remain unchanged. At $m=1$ its established $5/7$ rate also
remains unchanged. These equalities return the newly proved
$q$-scale Gamma information without assigning that scale to an
unestimated arithmetic contribution.

The original simple branch now has a more precise centre as well.
QGQ4--9 proves, from the complete QGT21 finite expression and its
actual four shifted equilibrium integrals,
\[
 \mathcal C_\Gamma^{[q,1]}
       =\frac{\partial_\alpha\mathcal L(2,\pi)}8
                            -\frac{\log2}{4},\qquad
 \frac{\mathcal V^{\rm c}_{q,l}-lq\mathcal C_\Gamma}{q}
                            \longrightarrow\mathcal C_\Gamma^{[q,1]}
 \quad(m=1).
\]
Its proof retains the Taylor remainder of every high block and
every original source and scalar correction. QGQ9 bounds the
remaining actual-$W_k$ error about that centre by $A/\sqrt2$
in the limit superior. Since
$(e_0-e_q-e_{q+1})/q\to0$ on the same original branch, the
identity in QRI3 gives the correctly signed return
\[
 \boxed{\displaystyle
 \limsup_{l\to\infty}
 \left|\frac{\mathcal H_k+lq\mathcal C_\Gamma}{q}
                         +\mathcal C_\Gamma^{[q,1]}\right|
                \le\frac A{\sqrt2}\qquad(m=1).}
 \tag{QRI9}
\]
Indeed the expression inside the absolute value equals the
negative of
$(W_k-lq\mathcal C_\Gamma)/q-\mathcal C_\Gamma^{[q,1]}$
plus the displayed vanishing source error. The triangle inequality
and QGQ9 prove QRI9. By each exact equality in QRI8 the same bound
applies to its left-hand side divided by $q$ and then increased
by $\mathcal C_\Gamma^{[q,1]}$. In particular the arithmetic scalar
is still subtracted there with its original sign. The finite
interval QRI3 remains valid before passing to this limit superior.

\subsection{The original Gamma source mapped into each mixed row}

The accepted GMB1--16 proof computes the original mixed rows from
the original quotient relation. Keep $s\in\{1,k\}$, the packet
centre $c=k/2$, its polynomial $\chi$, the fixed remainder frame
of $E_\chi=\mathbb C[S]/(\chi)$, and the same surjection
$\Lambda:E_\chi\to B_{\rm obs}$. The source order $s=1$ is
the original fixed Gamma comparison, and $s=k$ is the original
tensor Gamma source. The source norm is
$h_d^{(s)}=(2\pi)^{s/2}d!(s/2)_d$.

Let $T_\chi$ be multiplication by $S$ in the fixed remainder
frame, and put $Y_\chi=(T_\chi-cI)/i$. The finite original-source
recurrence is
\[
\begin{gathered}
 r_0=[1],\quad r_1=Y_\chi[1],\qquad
 r_{d+1}=Y_\chi r_d-d(s/2+d-1)r_{d-1},\\
 C_s=\left[\frac{r_0}{\sqrt{h_0^{(s)}}},\ldots,
          \frac{r_{q-1}}{\sqrt{h_{q-1}^{(s)}}}\right],\qquad
 v_j=\frac{C_s^{-1}r_{q+j}}{\sqrt{h_{q+j}^{(s)}}},\\
 K_j=I_q+\sum_{a<j}v_av_a^*,\qquad
 t_j=v_j^*K_j^{-1}v_j,\qquad R_s=\Lambda C_s,\\
 M_j^{\rm obs}=R_sK_jR_s^*,\qquad
 w_j=R_sv_j,\qquad
 \xi_j=w_j^*(M_j^{\rm obs})^{-1}w_j.
\end{gathered}\tag{MRI1}
\]

\paragraph{The actual observation columns in these coordinates.}
Write $\lambda_d^{(s)}=\Lambda r_d/\sqrt{h_d^{(s)}}$;
this is exactly OPG6's vector $\lambda_d$ with its source index
displayed. On the original image $B_{\rm obs}$ put
\[
\begin{gathered}
 \mathcal M_D^{(s)}=\sum_{d=0}^{D-1}
                \lambda_d^{(s)}(\lambda_d^{(s)})^*\quad(D\ge q),\\
 w_j=\lambda_{q+j}^{(s)},\qquad
 M_j^{\rm obs}=\mathcal M_{q+j}^{(s)},\qquad
 \xi_j=(\lambda_{q+j}^{(s)})^*
            (\mathcal M_{q+j}^{(s)})^{-1}\lambda_{q+j}^{(s)}.
\end{gathered}\tag{MRI1a}
\]
Indeed $R_sv_j=\Lambda C_sC_s^{-1}r_{q+j}/\sqrt{h_{q+j}^{(s)}}$
is the stated column. The identity term of $K_j$ gives the first
$q$ columns and the rank-one terms give exactly the remaining
columns in the displayed sum. Those first $q$ columns span the
actual image because $C_s$ is invertible and $\Lambda$ is onto.
Thus this matrix is positive definite on that image; the empty
image has its original empty inverse and determinant one.

Here these columns have the original period and amplitude input,
as a fully calculated formula. Retain $h,v_h=g/h$, $E_h,A_h,U$,
the original period map $\Pi$, its contour orientations and
parameter $u$, and the endpoint matrix $C$ of OPG1--5. The
one-factor dimension is $n_1=4m$ (OPG's local $n$), while the
high parity rank elsewhere in this receiver remains $n=q/2$.
The original average is
$P_k=(n_1+1)^{-1}\sum_{a=0}^{n_1}(C^{-a})^{\otimes k}$.
In particular $P_k$ here is OPG's cyclic projector, while the
scalar source product $P_k$ in QGT18 keeps its separately
specified meaning. For the image inclusion $\jmath$ and the
holomorphic branch $\arctan0=0$, OPG7--12 gives the exact map
\[
\begin{split}
 \jmath\sum_{d\ge0}\sqrt{h_d^{(s)}}\lambda_d^{(s)}\frac{z^d}{d!}
 &=\frac{(1+z^2)^{-s/4}}{n_1+1}
       \sum_{a=0}^{n_1}
          \bigl(C^{-a}\mathbf v(\arctan z)\bigr)^{\otimes k},\\
 \mathbf v(w)
 &=\sum_{\alpha\in\mathcal R}e^{w(\alpha-1/2)/i}
       \sum_{a=0}^{m-1}\mathbf b_{\alpha,a}
       \sum_{\nu=0}^{a}\frac{v_h^{(\nu)}(\alpha)}{\nu!}
                       \frac{(w/i)^{a-\nu}}{(a-\nu)!}.
\end{split}\tag{MRI1b}
\]

\paragraph{The computed original sector map and its actual rank.}
The full primary observation can now be evaluated in the literal
Fourier basis, with its inherited metric. Put $N_1=4m+1$,
$\zeta=\exp(2\pi i/N_1)$, $F_{jt}=\zeta^{jt}-1$ and
$(F^{-1})_{tj}=\zeta^{-tj}/N_1$, $1\le j,t<N_1$.
Let $S_\nu$ be the sum of all distinct ordered words with
$\nu_t$ copies of the column $f_t$ of $F$, and let
$\mathcal I_{k,0}=\{\nu\ge0:\sum_t\nu_t=k,
\sum_t t\nu_t\equiv0\pmod {N_1}\}$.
The map $\mathcal F_0$ sends the coordinate vector indexed by
$\nu\in\mathcal I_{k,0}$ to this complete ordered tensor $S_\nu$.
For the full sum primary basis $\mathcal H$ of OCS17, the actual
finite matrix is $\mathfrak A_{\nu;(a,b,D)}=i^DD!c_{\nu;a,b,D}$,
where OCS12--16 gives every $c$ by the finite multinomial sum of
the original contour integrals and all unit derivatives. Then
\[
\begin{gathered}
 L\mathcal H=\mathcal F_0\mathfrak A,\quad
 K=\mathcal H\ker\mathfrak A,\quad
 b=\operatorname{rank}\Lambda=\operatorname{rank}\mathfrak A,\quad r=q-b,\\
 \jmath\lambda_d^{(s)}=\mathcal F_0\widehat\lambda_d^{(s)},\qquad
 (\widehat\lambda_d^{(s)})_\nu
 =\frac1{\sqrt{h_d^{(s)}}}\sum_{a,b'=0}^k\sum_{D=0}^{e-1}
 c_{\nu;a,b',D}p_d^{(D)}(\omega_{a,b'}),\quad s\in\{1,k\}.
\end{gathered}\tag{MRI1c}
\]
Here $b'$ is the original second root index, while $b$ denotes
the actual boundary dimension. To prove the first identity,
expand $e^{wY}[1]$ in its full primary basis and compare the
coefficients of $w^De^{\omega_{a,b'}w}$ with the original
projected pure tensor. The primary expansion has coefficient
$i^{-D}/D!$, giving exactly the displayed matrix. Distinct
exponential polynomials are independent by applying the product
of their other constant-coefficient annihilators. Applying the
same comparison to $p_d(Y)[1]$ proves the second line, including
the cancellation of $i^DD!$ with its original Taylor coefficient.
The inverse of $F$ follows from the finite geometric sum;
the disjoint ordered word supports prove injectivity of
$\mathcal F_0$. Thus the kernel and rank identities concern
precisely the original observation. Its metric is the restriction
of OCS11, with $F^*F=N_1(I+\mathbf1\mathbf1^*)$; none of the
sector coordinates is declared orthonormal.

For $m=1$, the actual period quadric is
$Q_u(x)=Q_0(\mathsf H^{-1}x)$,
$Q_0(t)=t_{++}t_{--}-t_{+-}t_{-+}$ and
$\mathsf H=F^{-1}\Pi U\mathcal E_{\rm prim}$.
Its exact period entries and all deciding coefficients are
PQO6--17. RPC17, RPC21 and RPC23 define the positive finite
radius $R_*$ from the fixed original quartet and its actual unit;
the two cases are the computed $c_3\ne0$ and $c_3=0$,
with $c_1\ne0$ in the latter case. The complete estimates
RPC18--24 show a nonzero original commutator entry for every
$|u|\ge R_*$ and every logarithm branch. Therefore
\[
 m=1,\quad |u|\ge R_*,\quad k=4l+1,\ l\ge1
 \quad\Longrightarrow\quad
 b=k^2-6k+17,\qquad r=8k-16.\tag{MRI1d}
\]
Indeed the original quadric has five distinct prime orbit
factors there. An invariant polynomial divisible by $Q_u$
is divisible by each of those five factors; their product
$\mathcal H_Q$ is invariant and has degree ten. Cancellation
shows that the quotient is invariant. Hence its actual
invariant restriction rank is $d_{5,k,0}-d_{5,k-10,0}$.
The exact character formula
$d_{5,k,0}=(\binom{k+3}{3}+4\epsilon_k)/5$ has
$\epsilon_k=\mathbf1_{k\equiv0\ (5)}-\mathbf1_{k\equiv1\ (5)}$.
Subtracting the two cubics, with equal epsilon terms, gives
$k^2-6k+17$ for $k\ge10$; $k=5,9$ give respectively
$12,44$ directly with the negative-degree term zero.
Subtracting from $(k+1)^2$ proves MRI1d. Outside this proved
period domain the full original minors of MRI1c and the exact
orbit alternatives OQC10--13 remain the rank prescription.
For general $m$, MRI1c retains every primary direction.

This identity converges for $|z|<1$. The period vectors are the
original polynomial integrals
$(\mathbf b_{\alpha,a})_j=\int_{\Gamma_j}
e^{\Phi_h(z_1)/u}E_{\alpha,a}(z_1)\,dz_1$, where
$E_{\alpha,a}=\operatorname{rem}_h(h_\alpha t_\alpha
(z_1-\alpha)^a)$ and $t_\alpha$ is the degree $m-1$ Taylor
inverse of $h_\alpha=h/(z_1-\alpha)^m$ at $\alpha$.
The exact Gamma generating recurrence proves its scalar
prefactor $(1+z^2)^{-s/4}$, while the commuting tensor coordinate
actions factor its exponential with their original centres
$1/2$ and sum centre $c=k/2$. Taylor multiplication in each
primary algebra gives every derivative in the second line.
Applying the original polynomial period integrals and all
$n_1+1$ cyclic terms proves MRI1b. Its finite coefficient form
is OPG13--14, included in full below, with every ordered tuple,
nilpotent index and unit derivative retained.

The conjugate transpose is denoted by a star. The recurrence is
the remainder of the original monic Gamma orthogonal polynomial
recurrence; its determinant is
$\det C_s=i^{-q(q-1)/2}\prod_{d<q}(h_d^{(s)})^{-1/2}$.
It therefore gives an invertible map with its actual source
phase and mass. Positivity of $K_j$ and surjectivity of $R_s$
prove the stated inverses exist. For zero boundary dimension
the matrices are empty and $\xi_j=0$.

The original remainder map at degree $q+j-1$, in these orthonormal
source coordinates, is precisely
$C_s[I_q,v_0,\ldots,v_{j-1}]$. Its minimum norm lift of $C_sx$
is $[I_q,v_0,\ldots,v_{j-1}]^*K_j^{-1}x$. The orthogonality to
the full relation kernel proves the actual quotient metric
$G_{q+j-1}^{(s)}=C_s^{-*}K_j^{-1}C_s^{-1}$.
Furthermore
\[
\begin{gathered}
 r_j^{\rm ker}=v_j-K_jR_s^*(M_j^{\rm obs})^{-1}w_j,
 \qquad R_sr_j^{\rm ker}=0,\\
 t_j=\xi_j+(r_j^{\rm ker})^*K_j^{-1}r_j^{\rm ker},
 \qquad0\le\xi_j\le t_j.
\end{gathered}\tag{MRI2}
\]

\paragraph{The exact kernel vector supplied by the original recurrence.}
At $D=q+j$, keep the fixed coefficient section $S_*$ and
$I\kappa=I_E-S_*\Lambda$, and put
$e_D^{(s)}=r_D/\sqrt{h_D^{(s)}}$ and
$\zeta_D^{(s)}=\kappa e_D^{(s)}$. The required metric section
has the preceding degree $D-1$:
\[
\begin{gathered}
 S_{D-1}^{(s)}=C_sK_jC_s^*\Lambda^*
                          (\mathcal M_D^{(s)})^{-1},\qquad
 H_{D-1}^{(s),K}=I^*G_{D-1}^{(s)}I,\\
 \vartheta_{D,j}^{(s)}
       =\zeta_D^{(s)}-\kappa S_{D-1}^{(s)}\lambda_D^{(s)},\\
 C_sr_j^{\rm ker}
       =e_D^{(s)}-S_{D-1}^{(s)}\lambda_D^{(s)}
       =I\vartheta_{D,j}^{(s)},\\
 t_j-\xi_j=(\vartheta_{D,j}^{(s)})^*
                     H_{D-1}^{(s),K}\vartheta_{D,j}^{(s)}.
\end{gathered}\tag{MRI2a}
\]
The exact inverse Gram is
$(G_{D-1}^{(s)})^{-1}=C_sK_jC_s^*$, so MRI1a proves that
its boundary Gram is $(\mathcal M_D^{(s)})^{-1}$ and that
the displayed section obeys $\Lambda S_{D-1}^{(s)}=I$ and
$I^*G_{D-1}^{(s)}S_{D-1}^{(s)}=0$. Substituting MRI1a in
$C_sr_j^{\rm ker}$ gives the third line. Both $S_*$ and
$S_{D-1}^{(s)}$ are sections, so their difference lands in
$I(K)$; applying $\kappa$ gives exactly the correction on
the second line. Congruencing $G_{D-1}^{(s)}$ by $C_s$ gives
$K_j^{-1}$, which proves the last line from MRI2. Thus the
fixed-section kernel and the minimum-section kernel are connected
by the full original degree-dependent correction.

The vectors in MRI1a and MRI2a are also computed by the complete
coupled action. Let $\mathsf Y_B=\Lambda Y_\chi S_*$,
$\mathsf D=\Lambda Y_\chi I$, $\mathsf C=\kappa Y_\chi S_*$,
and $\mathsf Y_K=\kappa Y_\chi I$. To avoid changing the
original source-error names $a_0,b_0$, display OPG20's recurrence
coefficients here as
$\mathsf a_d^{(s)}=[(d+1)(s/2+d)]^{-1/2}$ and
$\mathsf b_d^{(s)}=[d(s/2+d-1)/((d+1)(s/2+d))]^{1/2}$,
with $\mathsf b_0^{(s)}=0$. The exact recurrence is
\[
\begin{split}
 \lambda_{d+1}^{(s)}
 &=\mathsf a_d^{(s)}(\mathsf Y_B\lambda_d^{(s)}
                         +\mathsf D\zeta_d^{(s)})
                           -\mathsf b_d^{(s)}\lambda_{d-1}^{(s)},\\
 \zeta_{d+1}^{(s)}
 &=\mathsf a_d^{(s)}(\mathsf C\lambda_d^{(s)}
                         +\mathsf Y_K\zeta_d^{(s)})
                           -\mathsf b_d^{(s)}\zeta_{d-1}^{(s)},\\
 \lambda_0^{(s)}&=\Lambda[1]/\sqrt{M_s},\qquad
 \zeta_0^{(s)}=\kappa[1]/\sqrt{M_s},\qquad
 \lambda_{-1}^{(s)}=\zeta_{-1}^{(s)}=0.
\end{split}\tag{MRI2b}
\]

\paragraph{The action defect after the computed sector and orbit maps.}
Let $J$ in the sum primary frame be the full matrix
$Je_{a,b,D}=\omega_{a,b}e_{a,b,D}+i^{-1}e_{a,b,D+1}$,
with the second term zero only at $D=e-1$. Then
\[
 Y\mathcal H=\mathcal HJ,\qquad
 \Lambda Y\mathcal Hx=\mathcal Q\mathfrak AJx\quad
 (x\in\ker\mathfrak A),\qquad
 \mathcal Q:\operatorname{im}\mathfrak A\xrightarrow{\sim}B,
 \quad x\mapsto\jmath^{-1}\mathcal F_0x.
 \tag{MRI2c}
\]
Multiplication by $(S-c)/i$ on each original primary power
proves the first formula. MRI1c proves the second and the
specified isomorphism, so this is the same defect
$\mathsf D=\Lambda YI$ of MRI2b with a computed original
coefficient matrix. The exact kernel undetected by every
iterate is $\mathcal H\ker[\mathfrak A;\mathfrak AJ;\ldots;
\mathfrak AJ^{q-1}]$: the degree-$q$ characteristic polynomial
expresses every subsequent power in these first $q$ powers.
Thus neither the rank calculation nor the cyclic quotient
removes the original action leaving the observed subspace.

On the simple-quartet five-orbit domain the degree-eight
product $\mathcal A_Q=\prod_{a=1}^4g^aQ_u$ gives the
exact invariant lift $[f]\mapsto[5\mathsf R(\mathcal A_Qf)]$
modulo $\mathcal H_Q$. Every nonidentity summand is divisible
by $Q_u$, so its restriction is $[\mathcal A_Qf]$.
This proves the inclusion of its image in the actual observed
dual space. Under the original restriction maps $\Theta_d$
the resulting injection $\mathfrak J_k:E_{k-8}^\vee\to B_k^\vee$
for $k\ge9$ satisfies the complete action law
\[
 Y_k^\vee\Lambda_k^\vee\mathfrak J_k
  -\Lambda_k^\vee\mathfrak J_kY_{k-8}^\vee
 =\Theta_kM_{[\mathcal D\mathcal A_Q]}\Theta_{k-8}^{-1},\qquad
 \mathcal D\mathcal A_Q=\sum_{a=1}^4(\mathcal D g^aQ_u)
                          \prod_{b\ne a,\ 1\le b\le4}g^bQ_u.
 \tag{MRI2d}
\]
Here $\mathcal D=(\mathsf H\Omega\mathsf H^{-1}x)\cdot\nabla$
is the original curve derivation. Differentiating the exact
curve restriction gives $\Theta_d\mathcal D=Y_d^\vee\Theta_d$,
and Leibniz's rule proves MRI2d. Its codomain is $E_k^\vee$;
restriction to the actual $K$ retains the full remaining term.
The degree eight comes from the four actual orbit quadrics.
The source orders in every metric remain $1,k$.

The exact norm ratio applied to the original Gamma polynomial
recurrence gives
$e_{d+1}^{(s)}=\mathsf a_d^{(s)}Y_\chi e_d^{(s)}
-\mathsf b_d^{(s)}e_{d-1}^{(s)}$. Splitting
$e_d^{(s)}=S_*\lambda_d^{(s)}+I\zeta_d^{(s)}$ and applying
$\Lambda$ and $\kappa$ proves both lines. The action defect is
the actual map
$\Lambda Y_\chi-\mathsf Y_B\Lambda=\mathsf D\kappa$;
in the original ambient tensor space OPG25--26 proves
$\jmath\mathsf D=P_k\mathcal T(I-P_k)\mathcal V I$.
Thus no invariant-kernel or invariant-boundary assumption enters
MRI2a or MRI2b. The observed-iterate quotient OPG23--24 retains
its original coefficients and its map to the original observation.

The two summands of $v_j$ are orthogonal in the $K_j^{-1}$
metric since their cross inner product contains $R_sr_j^{\rm ker}=0$.
The map $C_s$ restricts to a bijection
$\ker R_s\to\ker\Lambda$. Thus the original kernel and its
observation are both retained in the same computation.

Put $D=q+j$. The original monic relation has the exact squared
norm $a_{D-1}=h_D^{(s)}(1+t_j)$ and coefficient row
$\alpha_{D-1}=-i^{-D}\sqrt{h_D^{(s)}}\,
v_j^*K_j^{-1}C_s^{-1}$. The full section and block inverse
calculation in GMB11--13 gives the original two contractions
\[
\begin{gathered}
 \beta_{D-1}=h_D^{(s)}(1+t_j)\frac{t_j-\xi_j}{1+\xi_j},
 \qquad
 \nu_{D-1}=h_D^{(s)}(1+t_j)^2\frac{\xi_j}{1+\xi_j},\\
 1+\frac{\beta_{D-1}}{a_{D-1}}
          =\frac{1+t_j}{1+\xi_j},\qquad
 1+\frac{\nu_{D-1}}{a_{D-1}+\beta_{D-1}}=1+\xi_j.
\end{gathered}\tag{MRI3}
\]

With precisely the original observation input MRI1a and the
kernel vector MRI2a, these same ordered factors read
\[
\begin{gathered}
 1+\frac{\beta_{D-1}}{a_{D-1}}
   =1+\frac{(\vartheta_{D,j}^{(s)})^*
                  H_{D-1}^{(s),K}\vartheta_{D,j}^{(s)}}{1+\xi_j},\\
 1+\frac{\nu_{D-1}}{a_{D-1}+\beta_{D-1}}
   =1+(\lambda_D^{(s)})^*(\mathcal M_D^{(s)})^{-1}
                                              \lambda_D^{(s)}.
\end{gathered}\tag{MRI3a}
\]

\paragraph{The literal arithmetic source on these same mixed spaces.}
Write $H_N^{\rm ar}$ for the original polynomial Gram of
$m_{h,k}=w_h^{*k}$, $w_h(y)=|v_h(1/2+iy)|^2/(2\pi)$,
and $H_N^\sigma$ for the fixed Gamma Gram in the same
increasing $S$ coefficient frame. The complete AMT4 constants
$a_k,A_k,D_N$ are positive and give
$(a_k/D_N)H_N^\sigma\preceq H_N^{\rm ar}\preceq A_kH_N^\sigma$.
AMT5--6 proves this on every polynomial using the original
convolution envelopes, both Gamma density bounds, the full
weighted polynomial norm, and Cauchy--Schwarz; its exact
coordinate substitution retains the full coefficient vector
through an invertible map. For $a\in\{\sigma,\mathrm{ar}\}$ put
\[
\begin{gathered}
 G_N^a=(J_N(H_N^a)^{-1}J_N^*)^{-1},\quad
 V_N^a=(H_N^a)^{-1}J_N^*G_N^a,\quad H_{K,N}^a=I^*G_N^aI,\\
 Q_N^a=(\Lambda(G_N^a)^{-1}\Lambda^*)^{-1},\qquad
 S_N^a=(G_N^a)^{-1}\Lambda^*Q_N^a,\\
 (a_k/D_N)H_{K,N}^\sigma\preceq H_{K,N}^{\rm ar}
                                      \preceq A_kH_{K,N}^\sigma,\qquad
 (a_k/D_N)Q_N^\sigma\preceq Q_N^{\rm ar}\preceq A_kQ_N^\sigma.
\end{gathered}\tag{MRI3b}
\]
The same inequality holds for $G_N$. Its proof takes the
minimum over the identical affine fibre $J_NP=z$, then
restricts to $z=Ix$ or minimizes over $\Lambda z=y$.
Direct multiplication gives $J_NV_N^a=I$ and orthogonality
to $\ker J_N$, establishing the minimum formulas. The exact
lift bijection is $V_N^\sigma z\mapsto V_N^{\rm ar}z$,
whose inverse uses the same $z$; its difference belongs to
$\chi\mathcal P_{N-q}$. The boundary minimum lifts likewise
map $V_N^\sigma S_N^\sigma y\mapsto V_N^{\rm ar}S_N^{\rm ar}y$.
These prove the source-to-kernel and source-to-boundary maps
on the original observation, with its unchanged action defect.
Every formula retains general $m$ and the actual ranks $r,b$.

The tensor-source comparison retains the full divided Taylor
isomorphism $\mathsf E_\chi$ and the triangular multiplication
matrix $U_f$ of the original $f_l$. Its determinant is
$\prod_{a,b,j}[(2a-k)\delta-t_j+i(2b-k)\gamma]^e\ne0$,
since $k$ is odd and $t_j$ is real. Multiplication
$P\mapsto f_lP$ is a bijection from its original remainder
fibre to the full product-jet fibre
$(\mathsf E_fQ,\mathsf E_\chi Q)=(0,U_f\mathsf E_\chi z)$,
with inverse $Q\mapsto Q/f_l$. On these coordinates the
observation and its actual kernel are exactly
\[
 \widetilde\Lambda=\Lambda\mathsf E_\chi^{-1}U_f^{-1},\quad
 \widetilde\Lambda U_f\mathsf E_\chi=\Lambda,\quad
 \ker\widetilde\Lambda=U_f\mathsf E_\chi K,\quad
 \widetilde I=U_f\mathsf E_\chi I,\quad
 \widetilde S_*=U_f\mathsf E_\chi S_*.
 \tag{MRI3c}
\]
The inverse product-jet map proves both kernel inclusions.
All original unit and nilpotent coefficients remain in
$\Lambda$ and $U_f$. AMT27--28 computes the corresponding
positive product-kernel Schur complement before applying
these maps; this exact map carries that full joint metric
to the same kernel and boundary used here.

This is obtained by substituting $t_j-\xi_j$ from MRI2a and
$\xi_j$ from MRI1a into MRI3. It preserves the denominator
$1+\xi_j$, the monic relation norm, the phase in its row, and
the original order of elimination. In particular the kernel
norm is evaluated in degree $D-1$, while the original row
contractions in MRI3 retain their degree-$D$ block definitions.
The exact map OPR1--4 proves why those two degrees occur.

These are the original BRU factors in their original order. Their
product is $1+t_j$, equal to
$\det G_{D-1}^{(s)}/\det G_D^{(s)}$ by the exact rank-one
inverse update. The terms in each denominator and the phase in
the row remain part of the original map.

Define the actual four-endpoint weight
$\mathcal W_q(f)=f_0+2\sum_{j=1}^{q-1}f_j+f_q$. Telescoping
the preceding quotient and mixed determinants gives
\[
\begin{gathered}
 \mathcal B_k^{(s)}=\mathcal W_q\bigl(\log(1+t_j)\bigr),\\
 F_K^{(s)}=\mathcal W_q\left(\log\frac{1+t_j}{1+\xi_j}\right),
 \quad F_B^{(s)}=\mathcal W_q\bigl(\log(1+\xi_j)\bigr),\\
 F_K^{(s)}+F_B^{(s)}=\mathcal B_k^{(s)},\qquad
 0\le F_K^{(s)}\le\mathcal B_k^{(s)},\quad
 0\le F_B^{(s)}\le\mathcal B_k^{(s)}.
\end{gathered}\tag{MRI4}
\]

The boundary part of this same finite endpoint calculation is
now explicitly
\[
\begin{split}
 F_B^{(s)}
 &=\log\frac{\det\mathcal M_{2q}^{(s)}
                  \det\mathcal M_{2q+1}^{(s)}}
                 {\det\mathcal M_q^{(s)}
                  \det\mathcal M_{q+1}^{(s)}},\\
 F_K^{(s)}
 &=\mathcal W_q\left(\log\left[1+
       \frac{(\vartheta_{q+j,j}^{(s)})^*
                H_{q+j-1}^{(s),K}\vartheta_{q+j,j}^{(s)}}
            {1+(\lambda_{q+j}^{(s)})^*
                (\mathcal M_{q+j}^{(s)})^{-1}\lambda_{q+j}^{(s)}}
                                    \right]\right).
\end{split}\tag{MRI4a}
\]

\paragraph{Finite individual control on the evaluated period domain.}
For $m=1$ define the original finite quantities
\[
 D_*=\max\{3,\sqrt{\delta^2+\gamma^2}\},\quad
 Z_q=1+2(q+1)q^22^{6q}(2D_*q)^{4q},\quad
 C_*=12+4\log(2D_*).
\]
The exact remainder computation OKV7--15 retains every root
and the denominator $(q+j)!(s/2)_{q+j}$ in its estimate of
$\|v_j\|^2$. In particular $I\preceq K_j\preceq Z_qI$.
With the original $T_s=C_s^{-1}I$, congruence gives
$Z_q^{-1}T_s^*T_s\preceq T_s^*K_j^{-1}T_s\preceq T_s^*T_s$.
The four determinants retain their actual frame factor until
their signs $+,+,-,-$ cancel it. Hence for the actual rank $r$
\[
\begin{gathered}
 0\le F_K^{(s)}\le2r\log Z_q\le2rC_*q\log(q+1),\quad s=1,k,\\
 \max\{0,\mathcal B_k^{(s)}-2r\log Z_q\}
             \le F_B^{(s)}\le\mathcal B_k^{(s)}.
\end{gathered}\tag{MRI4b}
\]
Each determinant difference is nonnegative by the nested
minimum metric, and at most $r\log Z_q$ by the displayed
congruence, proving the first line. MRI4 proves the second.
The bound depends on the actual rank, and otherwise does not
depend on the position or conditioning of the period kernel.

For the literal arithmetic source, define
$L_N=\log(A_kD_N/a_k)$,
$\Xi_{k,N}=(N+1)\log A_k-\log(\det H_N^{\rm ar}/\det H_N^\sigma)$,
and $d_{E,N}=q\log A_k-\log(\det G_N^{\rm ar}/\det G_N^\sigma)$.
The exact source, relation, kernel and boundary decomposition
AMT11--14 proves $\Xi_{k,N}=d_{R,N}+d_{K,N}+d_{B,N}$,
where $d_{K,N}+d_{B,N}=d_{E,N}$ and every term is nonnegative.
Set $\widehat c_{K,N}=\min\{rL_N,\Xi_{k,N}\}$ and
$\widehat C_K^+=\widehat c_{K,2q-1}+\widehat c_{K,2q}$.
Then the complete arithmetic return at the same four degrees is
\[
\begin{gathered}
 0\le F_K^{\rm ar}\le U_K:=2r\log Z_q+\widehat C_K^+
              \le2rC_*q\log(q+1)+rE_k^-,\\
 \max\{0,\mathcal B_k^\sigma+\delta_k^\sigma-U_K\}
       \le F_B^{\rm ar}\le\mathcal B_k^\sigma+\delta_k^\sigma,\\
 m=1,\ |u|\ge R_*:\qquad
 \frac{F_K^{(1)}}{q^2},\frac{F_K^{(k)}}{q^2},
 \frac{F_K^{\rm ar}}{q^2}\longrightarrow0,\qquad
 \frac{F_B^{(1)}}{q^2},\frac{F_B^{(k)}}{q^2},
 \frac{F_B^{\rm ar}}{q^2}\longrightarrow C_B.
\end{gathered}\tag{MRI4c}
\]
Here $E_k^-=L_{2q-1}+L_{2q}$ is the original upper endpoint
width AMT19. The arithmetic kernel difference is
$-d_{K,q-1}-d_{K,q}+d_{K,2q-1}+d_{K,2q}$, so its upper
bound is $\widehat C_K^+$; this proves the first line from
MRI4b. The exact arithmetic block determinant splitting
proves the second. On MRI1d, $r=8k-16$ and $q=(k+1)^2$,
while $E_k^-=O_h(k+\log q)$. Thus $U_K/q^2\to0$.
BRI6 gives all three total baseline limits, proving the last
line. These limits are uniform over $|u|\ge R_*$ and all
original branches because every error bound just displayed is
independent of $u$. The finite source comparison MRI3b and
the actual-rank caps retain their general-multiplicity scope;
the simple-quartet specialization is used only where stated.

\paragraph{The retained factorial sharpens the same complete kernel bound.}
For every original multiplicity $m\ge1$ keep the full repeated
root list of $\chi$, $q=e(k+1)^2$, and its actual maximum
$R=k\sqrt{\delta^2+\gamma^2}$. Define the precise finite
constants of KAF6 and KAF15 by
\[
\begin{gathered}
 A_{k,m}=\max\{2\sqrt{1+k/(2q)},R/q\},\quad
 \mathcal A_{k,m}=6\log2+4+4\log A_{k,m},\\
 \eta_{k,m}=\frac{1-1/(4q)}{A_{k,m}^2},\qquad
 Z_{k,m}=1+q^22^{6q}\frac{(4qA_{k,m})^{4q}}{(4q)!}
                 \frac{1-\eta_{k,m}^{q+1}}{1-\eta_{k,m}}.
\end{gathered}
\]
Here $A_{k,m}$ differs from the original arithmetic constant
$A_k$; both are retained with their specified definitions.
The full source norm denominator gives the sharper estimate
\[
\begin{gathered}
 \|v_j^{(s)}\|^2\le
 q^22^{6q}\frac{(4qA_{k,m})^{2(q+j)}}{(2q+2j)!},\quad0\le j\le q,\\
 \|K_{q+1}^{(s)}\|\le Z_{k,m},\qquad
 \log Z_{k,m}\le q\mathcal A_{k,m}+2\log q,\\
 0\le F_K^{(s)}\le2r\log Z_{k,m},\qquad
 0\le F_K^{\rm ar}\le2r\log Z_{k,m}+\widehat C_K^+
                            \le2r\log Z_{k,m}+rE_k^-.
\end{gathered}\tag{MRI4d}
\]
These bounds retain general $m$, the actual rank $r$ and
the original source orders $s=1,k$. For a complete derivation,
the signed remainder formula KAF8--10 applies to the root
list with every multiplicity repeated. With $T=qA_{k,m}$,
the original recurrence bounds both low monomial norms by
$\sqrt{M_s}T^a$ and the full coefficient sum of $p_d$ by
$(2T)^d$. It follows that the squared remainder norm is
at most $M_sq^22^{6q}(2T)^{2d}$. Dividing by the original
$M_sd!(s/2)_d$ and using
$d!(s/2)_d\ge d!(1/2)_d=(2d)!/4^d$ proves the first line.
For $t_d=(4T)^{2d}/(2d)!$, its reversed successive ratio is
$(2d)(2d-1)/(16T^2)\le\eta_{k,m}<1/4$ for
$q+1\le d\le2q$. Summing all $q+1$ original updates gives
the precise geometric sum in $Z_{k,m}$. The integral bound
$\log(4q)!\ge4q\log(4q)-4q+1$ and that geometric sum
$\le4/3$ give the second line, as in KAF17. Applying the
same four original determinant comparisons used in MRI4b
proves the Gamma kernel bound. The actual arithmetic deficit
cap in MRI4c then proves the last bound without altering its
correlation with the boundary deficit.

% Replacement for the current MRI4e receiving paragraph in both full receivers.
For $m=1$, $|u|\ge R_*$ and all original branches, retain
$r=8k-16$, $q=(k+1)^2$. The complete AKS20--32 proof now gives
the sharper source bounds. With its exact covariance logarithms
$\mathcal L_{s,j}$ and first compressed increment $d_1^{(s)}$, set
\[
 \begin{gathered}
 L_s=d_1^{(s)},\qquad
 U_s=\min\{2r\log Z_{k,1},
       r(\mathcal L_{s,q}+\mathcal L_{s,q+1})-d_1^{(s)}\},\\
 L_s\le F_K^{(s)}\le U_s,\quad s\in\{1,k\},\qquad
 \kappa_*=32\log\frac{25\log3}{4\pi},\\
 0\le\liminf\frac{F_K^a}{kq}
 \le\limsup\frac{F_K^a}{kq}\le\kappa_*,\qquad
 a\in\{\sigma,0,\mathrm{ar}\}.
 \end{gathered}\tag{MRI4e}
\]
Here $\mathcal L_{s,j}$ is the fully evaluated finite expression
AKS3, proved through AKS24; $d_1^{(s)}$ is AKS29's actual minimum-section increment.
The value $\kappa_*\in(25.02078097649,25.02078097650)$ follows
from the exact logarithmic certificate. The old factorial estimate
remains in the minimum, since a smaller asymptotic constant alone
does not compare all finite parameters. AKS derives its second bound
by confluent Newton division with every repeated root and the exact
generating coefficients of the same Gamma polynomials. Inverting and
compressing the covariance gives
$F_K^{(s)}=d_q^{(s)}+d_{q+1}^{(s)}-d_1^{(s)}$;
each high term is at most $r\mathcal L_{s,j}$ and monotonicity gives
the lower term $d_1^{(s)}$. This proves the finite inequality.
AKS30--32 proves the stated constant from those explicit terms.

Retain the correlated caps from AMT18, with their signs. The
literal arithmetic source and its boundary satisfy the finer bounds
\[
 \begin{gathered}
 L_{\rm ar}=\max\{0,L_1-C_K^-,
                       L_1+\delta_k^\sigma-C_B^+\},\\
 U_{\rm ar}=\min\{\mathcal B_k^{\rm ar},
                  U_1+C_K^+,U_1+\delta_k^\sigma+C_B^-\},\\
 L_{\rm ar}\le F_K^{\rm ar}\le U_{\rm ar},\qquad
 \mathcal B_k^{\rm ar}-U_{\rm ar}
       \le F_B^{\rm ar}\le\mathcal B_k^{\rm ar}-L_{\rm ar}.
 \end{gathered}\tag{MRI4f}
\]
For proof, the two lower caps for $F_K^{\rm ar}-F_K^{(1)}$
are $-C_K^-$ and $\delta_k^\sigma-C_B^+$; the two upper caps
are $C_K^+$ and $\delta_k^\sigma+C_B^-$. Add the appropriate
endpoints of MRI4e and intersect with
$0\le F_K^{\rm ar}\le\mathcal B_k^{\rm ar}$, which follows
from the exact nonnegative kernel/boundary splitting. Subtracting
from that same total gives the second interval, retaining correlation.
The complete AMT source-difference bounds are $o(kq)$; hence the
arithmetic ratio has the same limit points as both original Gamma
ratios and the same upper constant. All uniformity statements concern
the specified $|u|\ge R_*$ domain with its original branches.

AKJ3 proves the exact equality between the AKS coefficient projector
determinant and the existing OKM conductor-frame determinant, including
the frame denominator and its cancellation at these four endpoints.
AKJ8--11 gives the residue-annihilator map in the complete original
primary coefficients. AKJ12--25 supplies the actual conductor's
degree-adapted basis and norm bound. These calculations act on the
same metric determinant; they do not assume a unique leading value.
AKP6--10 supplies the full corrected minimum-section Gram, matching
the four checked finite observation modules in PR32. The original
theta class and the primitive for a representative change remain
the explicit maps AKP1--5.


\paragraph{The conductor's exterior metric error is now controlled.}
The complete RUI and CEP calculations below sharpen the preceding
operator estimate at the level of the actual exterior determinant.
Keep the fixed original period, $m=1$, $k\ge9$, and all four original
source endpoints. With the explicit constants $E_N^{(s)}$ of
CEP29--36 set
\[
 \epsilon_k^{(s)}=2\sum_{N\in\{q-1,q,2q-1,2q\}}E_N^{(s)}.
 \tag{MRI4g}
\]
These constants retain every actual conductor coefficient and
$\mu_v$, the original centre shift, and source order $s\in\{1,k\}$.
CEP37 supplies the full graph-quotient Grams in the transported
original kernel frame. Its proof gives
\[
 \left|F_K^{(s)}-\widetilde F_K^{(s)}\right|
 \le\epsilon_k^{(s)}=o(kq).
 \tag{MRI4h}
\]
The determinant-one calculation RUI5--9 uses complementary singular
values to replace the earlier rank-multiplied inverse bound.
CEP33--39 proves the exact passage through the original source ideal
and its low-degree graph. Neither kernel invariance nor an assigned
source order enters this estimate.

There is also the exact compressed-covariance calculation CEP46--49.
In its original fixed frames define
\[
 \begin{gathered}
 \mathcal T_k^{(s)}=\sum_N\sigma_N
      (\log\det G_N^{(s)}-\log\det G_{N-v}^{\prime(s)}),\qquad
 F_{\Xi,k}^{(s)}=\sum_N\sigma_N\log\det Q_{\Xi,N}^{(s)},\\
 \mathcal E_{A,k}^{(s)}=\sum_N\sigma_N\log\det K_{A,N}^{(s)},
 \qquad\sigma=(1,1,-1,-1).
 \end{gathered}\tag{MRI4i}
\]
The precise lower grid, omitted-root products, $\Xi$ quotient,
compressed covariance, original masses and endpoint-dependent
minimum sections are the complete explicit maps CEP1--49. They give
\[
 F_K^{(s)}=\mathcal T_k^{(s)}-F_{\Xi,k}^{(s)}
                       +\mathcal E_{A,k}^{(s)},\qquad
 |\mathcal E_{A,k}^{(s)}|\le\epsilon_k^{(s)}=o(kq).
 \tag{MRI4j}
\]
Here the remaining observation quotient has the actual rank
$\Delta-r=8k-32$. Its metric is retained. The lower source uses
the original order $s$, with cutoffs $q-1-v,q-v,2q-1-v,2q-v$;
the original roots have lower tensor degree $k-8$.

All preceding finite intervals can now be intersected with this
explicit conductor-reduction interval. For example set
$X_k^{(1)}=\mathcal T_k^{(1)}-F_{\Xi,k}^{(1)}$ and retain the full
correlated caps of MRI4f. Define
\[
 \begin{gathered}
 \ell_{A,k}=X_k^{(1)}-\epsilon_k^{(1)}
                    +\max\{-C_K^-,\delta_k^\sigma-C_B^+\},\\
 u_{A,k}=X_k^{(1)}+\epsilon_k^{(1)}
                    +\min\{C_K^+,\delta_k^\sigma+C_B^-\},\\
 \max\{L_{\rm ar},\ell_{A,k}\}\le F_K^{\rm ar}
                \le\min\{U_{\rm ar},u_{A,k}\}.
 \end{gathered}\tag{MRI4k}
\]
Each inequality is obtained by adding the original source-difference
bound to the two endpoints in MRI4j and intersecting with MRI4f.
Subtracting this interval from the exact original arithmetic total
gives its boundary interval, preserving the same correlations.
The scale estimates here hold for fixed original period; any
moving-period assertion still requires control of the displayed
actual constants. The earlier period-uniform AKS ceiling remains
valid with its stated stronger uniformity scope.


\paragraph{The outer-root quotient and the actual residual observation.}
CTV and ROQ now give full formulas for both remaining matrices in
MRI4j. The outer quotient uses the complete polynomial
$P_{\partial,k}$ and the original multiplication-pullback measure
$|P_{k-8}(y)|^2dm_s(y)$ at degrees
$\Delta-1,\Delta,q+\Delta-1,q+\Delta$. Its rank is
$\Delta=16k-48$. The original source order remains $s\in\{1,k\}$.
The residual observation has the actual coefficient rows
$D=\mathsf F^{\mathsf T}$ and $\Xi\mathsf N^{\mathsf T}=Z^{\mathsf T}$
of ROQ4--7, with rank $j=8k-32$. Its full covariance, minimum
section, period-word metric and kernel-plus-ideal graph are proved
in ROQ9--27. The exact positive return is
\[
 \Phi_{k,s}=\log(1+\eta_{q-1})
        +2\sum_{N=q}^{2q-2}\log(1+\eta_N)
        +\log(1+\eta_{2q-1}),\qquad
 \lambda_s=\log(1+\eta_{q-1})\ge0.
 \tag{MRI4l}
\]
Here every $\eta_N$ is the original constrained covariance term
ROQ17. Retaining the signed tails gives the exact current formula
\[
 \begin{gathered}
 F_K^{(s)}=\mathcal F_{\partial,k}^{(s)}-\Phi_{k,s}+\mathfrak e_{k,s},\\
 -e_{k,s}^-\le\mathfrak e_{k,s}\le e_{k,s}^+,\qquad
 e_{k,s}^\pm=o(kq).
 \end{gathered}\tag{MRI4m}
\]
The two explicit errors are OAR1--3: the same conductor error
$\epsilon_{A,k}^{(s)}$ plus the four original CTV18 tail bounds,
with their positive and negative endpoint signs retained separately.
Their proof preserves fixed-period scope and all original constants.

With $V_s=j(\mathcal L_{s,q}+\mathcal L_{s,q+1})-\lambda_s$,
the new finite source intervals are
\[
 \begin{split}
 L_s^\partial&=\max\{L_s,
          \mathcal F_{\partial,k}^{(s)}-e_{k,s}^--V_s\},\\
 U_s^\partial&=\min\{U_s,
          \mathcal F_{\partial,k}^{(s)}+e_{k,s}^+-\lambda_s\},\\
 L_s^\partial&\le F_K^{(s)}\le U_s^\partial.
 \end{split}\tag{MRI4n}
\]
ROQ21 proves $\lambda_s\le\Phi_{k,s}\le V_s$; inserting it in
MRI4m and intersecting with MRI4e proves every bound in MRI4n.
The full value of MRI4l gives the still more precise interval
centred at $\mathcal F_{\partial,k}^{(s)}-\Phi_{k,s}$ with
the two errors $e_{k,s}^-,e_{k,s}^+$.

The arithmetic and boundary refinements preserve the original
correlations:
\[
 \begin{gathered}
 L_{\rm ar}^\partial=\max\{L_{\rm ar},L_1^\partial-C_K^-,
                     L_1^\partial+\delta_k^\sigma-C_B^+\},\\
 U_{\rm ar}^\partial=\min\{U_{\rm ar},U_1^\partial+C_K^+,
                     U_1^\partial+\delta_k^\sigma+C_B^-\},\\
 L_{\rm ar}^\partial\le F_K^{\rm ar}\le U_{\rm ar}^\partial,\\
 \mathcal B_k^{\rm ar}-U_{\rm ar}^\partial
       \le F_B^{\rm ar}\le\mathcal B_k^{\rm ar}-L_{\rm ar}^\partial.
 \end{gathered}\tag{MRI4o}
\]
Add the exact two-sided source-difference inequality in MRI4f
to MRI4n at $s=1$, then subtract from the same arithmetic total.
This proves MRI4o, as fully expanded in OAR8--10. All earlier
intervals remain available for intersection on their same target
quantities; no original coefficient or positive increment is omitted.


\paragraph{The original Gamma moment determinant comparison.}
OSP1--35 proves a full coefficient-space norm comparison, including
the inner real interval, for the complete original root polynomials.
Retain its finite domain $k\ge29$ and
$k\sqrt{\delta^2+\gamma^2}\le2^{-33}(k+1)^2$.
For every fixed original quartet this holds eventually. OSP32--35
defines the exact same-source moment expression
$\mathcal F_{0,k}^{(s)}$, with the full original Gamma mass;
OSP33 gives its error $c_k^\circ=4(q+\Delta+1)\mathcal E_k=O(q)$.
Thus MRI4m becomes
\[
 \begin{gathered}
 F_K^{(s)}=\mathcal F_{0,k}^{(s)}-\Phi_{k,s}+\mathfrak e^0_{k,s},\\
 -e_{k,s}^{0,-}\le\mathfrak e^0_{k,s}\le e_{k,s}^{0,+},\qquad
 e_{k,s}^{0,\pm}=e_{k,s}^\pm+c_k^\circ=o(kq).
 \end{gathered}\tag{MRI4p}
\]
PMR1--3 proves this identity and both error endpoints by adding
the exact outer-root comparison to the retained signed error.
The period-dependent residual sum is unchanged.

The complete source and arithmetic intervals transport as
\[
 \begin{gathered}
 L_s^0=\max\{L_s^\partial,\mathcal F_{0,k}^{(s)}-e_{k,s}^{0,-}-V_s\},\qquad
 U_s^0=\min\{U_s^\partial,\mathcal F_{0,k}^{(s)}+e_{k,s}^{0,+}-\lambda_s\},\\
 L_{\rm ar}^0=\max\{L_{\rm ar}^\partial,L_1^0-C_K^-,
                                    L_1^0+\delta_k^\sigma-C_B^+\},\\
 U_{\rm ar}^0=\min\{U_{\rm ar}^\partial,U_1^0+C_K^+,
                                    U_1^0+\delta_k^\sigma+C_B^-\},\\
 L_s^0\le F_K^{(s)}\le U_s^0,\qquad
 L_{\rm ar}^0\le F_K^{\rm ar}\le U_{\rm ar}^0,\\
 \mathcal B_k^{\rm ar}-U_{\rm ar}^0
      \le F_B^{\rm ar}\le\mathcal B_k^{\rm ar}-L_{\rm ar}^0.
 \end{gathered}\tag{MRI4q}
\]
Insert the exact residual caps from MRI4n into MRI4p, intersect
with MRI4n--o, and add the original correlated source differences.
Subtracting from the same arithmetic total proves the last line.
PMR4--6 gives the full proof. PMR8 proves that the displayed
endpoints equal the retained outer-root endpoints exactly:
the new moment comparison supplies an analytic calculation
target with controlled error. Earlier finite intervals cover all
smaller $k$ without using the additional OSP comparison.


\paragraph{The evaluated outer contribution and the actual residual term.}
OFV1--32 evaluates the complete seven-moment expression through the
original Gamma recurrence, both parity blocks, the low-rank determinant
bounds and the uniform finite partition estimate EIQ/QLG. Retain
$k\ge257$ and the original eventual domain of MRI4p; no earlier finite
bound needs that added restriction. Define
\[
 \mathfrak C_\partial=4\log(4/\pi)+2-2I_{2,\pi},\qquad
 I_{2,\pi}=\iint\log|x-y|\,d\mu_{2,\pi}(x)d\mu_{2,\pi}(y).
 \tag{MRI4r}
\]
Here the measure is exactly the proved EIQ density with its original
endpoints. OFV30 proves
$\mathcal F_{\partial,k}^{(s)}=\Delta q\mathfrak C_\partial+
O_{\delta,\gamma}(q\log q)$ for both $s=1,k$.
The full remainder and both finite signs of OVR2--3 give
\[
 \begin{gathered}
 b_{k,s}=\mathcal F_{\partial,k}^{(s)}-\Delta q\mathfrak C_\partial,
 \quad R_{k,s}=b_{k,s}+\mathfrak e_{k,s}=o(kq),\\
 F_K^{(s)}+\Phi_{k,s}=\Delta q\mathfrak C_\partial+R_{k,s},
 \quad b_{k,s}-e_{k,s}^-\le R_{k,s}\le b_{k,s}+e_{k,s}^+,\\
 \frac{\Phi_{k,k}-\Phi_{k,1}}{kq}\longrightarrow0.
 \end{gathered}\tag{MRI4s}
\]
The first identity substitutes the evaluated outer term into MRI4m,
preserving its entire original error. OVR4 proves the last limit by
subtracting these two equalities and using the proved original kernel
source difference. The matrices in $\Phi$ remain the full ROQ
constrained covariances; no residual leading value has been assigned.

For the original $d_k^{\rm ar}=F_K^{\rm ar}-F_K^{(1)}$, the
correlated finite bounds remain OAR8/MRI4o. In particular AMT38/46
gives $d_k^{\rm ar}=o(kq)$, and OVR5 proves the exact update
\[
 \begin{split}
 F_K^{\rm ar}
   &=\Delta q\mathfrak C_\partial-\Phi_{k,1}+R_{k,1}+d_k^{\rm ar},\\
 F_B^{\rm ar}
   &=\mathcal B_k^{\rm ar}-\Delta q\mathfrak C_\partial
                 +\Phi_{k,1}-R_{k,1}-d_k^{\rm ar}.
 \end{split}\tag{MRI4t}
\]
The second line subtracts the first from the same original arithmetic
total. This preserves the two errors and their correlation with the
boundary. OVR9--11 also proves the finite common cap
\[
 \Delta q\mathfrak C_\partial+R_{k,s}
 \le\Delta(\mathcal L_{s,q}+\mathcal L_{s,q+1})-d_1-\lambda_s,
 \tag{MRI4u}
\]
where $d_1$ and $\lambda_s$ are the actual first kernel and residual
increments. Indeed add AKS28 and ROQ21, then use $r+j=\Delta$.
With $\kappa_*=32\log(25\log3/(4\pi))$, both
$F_K^{(s)}/(kq)$ and $\Phi_{k,s}/(kq)$ have lower limit at least
$\max\{0,16\mathfrak C_\partial-\kappa_*\}$ and upper limit at
most $\min\{\kappa_*,16\mathfrak C_\partial\}$.
Their sum tends to $16\mathfrak C_\partial$, and hence
$0\le\mathfrak C_\partial\le\kappa_*/8$. OVR9--11 supplies the
complete finite-cap and limiting arguments, retaining both source orders.

\paragraph{Fixed original coefficient frames and their Gamma metric.}
FMC1--38 constructs the fixed rank-$625$ module over
$\mathbb C[x_1^5,x_2^5,x_3^5,x_4^5]$, its actual rank-$125$
invariant submodule, and the complete homogeneous relation reductions
for the original period coefficients. Its forward and inverse strip
frame maps have exterior logarithmic cost $O_{\rm actual}(k^2)=o(kq)$.
The exact maps into the earlier residual frame, period Gram and original
polynomial degree flag are FMC27--28 and FMC35--38.
In particular, with the actual primary frame
$J_k=\mathsf N_k^{\mathsf T}T_k^{-\mathsf T}$, they retain
\[
 \begin{gathered}
 \mathcal H_N=J_k^*G_NJ_k
   =\begin{pmatrix}H_{11,N}&H_{12,N}\\H_{21,N}&H_{22,N}\end{pmatrix},\\
 Q_{\Xi,N}^{\rm fix}=H_{11,N}-H_{12,N}H_{22,N}^{-1}H_{21,N},
 \qquad N=q-1,q,2q-1,2q.
 \end{gathered}\tag{MRI4v}
\]
Completing the square in the full unchanged Gamma form proves this
minimum; FMC29--34 gives every original recurrence, mass and coefficient.
The frame and inherited period-Gram changes cancel exactly in the
four-endpoint return. The separate primary-to-polynomial Vandermonde
Gram factor is retained in FMC35, and the degree-flag transition is
FMC38. Thus the new coefficient bounds come with the full original
metric morphism. The analytic quantity to evaluate next is the actual
compression in MRI4v, equivalently the full residual sum in MRI4s.


In the second line $\mathcal W_q$ acts on the displayed sequence
indexed by $j=0,\ldots,q$. For the first line, the exact update
$\mathcal M_{D+1}^{(s)}=\mathcal M_D^{(s)}
+\lambda_D^{(s)}(\lambda_D^{(s)})^*$ gives determinant ratio
$1+\xi_j$ at $D=q+j$. Applying the original weights
$1,2,\ldots,2,1$ leaves the two positive high logarithms and
the two negative low logarithms displayed. Applying the same
weights to MRI3a gives the second line. The full sum and every
nonnegative factor in MRI4 are consequently unchanged.

The weights are exactly those of degrees $q-1,q,2q-1,2q$;
the relation at step $j$ has degree $q+j$. Positivity in MRI2
proves each displayed inequality term by term. BRI6 and BRI9
therefore apply to the complete sum in MRI4 for each of the
two original source orders. The finite formulas MRI1--3 give
the individual factors. MRI4b--c now supplies their finite
actual-rank estimates and, on the proved simple-quartet period
domain, their individual $q^2$ coefficients for all three sources.

The arithmetic transition remains the exact original identity
\[
 \mathcal B_k^{\rm ar}
   =F_K^{(k)}+F_B^{(k)}+\mathcal H_k+\delta_k^\sigma
   =F_K^{(1)}+F_B^{(1)}+\delta_k^\sigma.
 \tag{MRI5}
\]

\paragraph{The full arithmetic split and its correlated finite return.}
Define $F_K^{\rm ar}$ and $F_B^{\rm ar}$ by the four
minimum metrics MRI3b with signs $+,+,-,-$. Their exact
coefficient-frame identity is
\[
 \mathcal B_k^{\rm ar}=F_K^{\rm ar}+F_B^{\rm ar},\qquad
 \Delta_X^\sigma=F_X^{\rm ar}-F_X^{(1)},\quad X=K,B,
 \qquad\Delta_K^\sigma+\Delta_B^\sigma=\delta_k^\sigma.
 \tag{MRI5a}
\]
For $r_K=r,r_B=b$, retain the full individual deficits
$d_{X,N}=r_X\log A_k-\log(\det H_{X,N}^{\rm ar}/
\det H_{X,N}^\sigma)$, where $H_{B,N}^a=Q_N^a$.
Put $c_{X,N}=\min\{r_XL_N,d_{E,N}\}$,
$C_X^-=c_{X,q-1}+c_{X,q}$ and
$C_X^+=c_{X,2q-1}+c_{X,2q}$. They retain the exact nested
caps $d_{X,N}\le c_{X,N}\le\min\{r_XL_N,\Xi_{k,N}\}$.
Substitution into each signed determinant gives
\[
\begin{gathered}
 \Delta_X^\sigma=-d_{X,q-1}-d_{X,q}+d_{X,2q-1}+d_{X,2q},\\
 \max\{-C_K^-,\delta_k^\sigma-C_B^+\}
  \le\Delta_K^\sigma\le
          \min\{C_K^+,\delta_k^\sigma+C_B^-\},\\
 \max\{-C_B^-,\delta_k^\sigma-C_K^+\}
  \le\Delta_B^\sigma\le
          \min\{C_B^+,\delta_k^\sigma+C_K^-\}.
\end{gathered}\tag{MRI5b}
\]
The determinant split in the common frame proves MRI5a;
four copies of $r_X\log A_k$ cancel to prove the first
line of MRI5b. The other lines intersect each component's
interval with the exact complement in MRI5a. This retains
all source deficits and their correlation. The symbols
$C_B^\pm$ here are the endpoint caps; the unsuperscripted
$C_B$ remains the baseline coefficient of BRI6.

The original full multiplier gives a second exact comparison
on these same spaces. Keep $f_l(S)=\prod_{j<l}(S-c-t_j)$,
$t_j=2j+1/2$, $\beta_k=(2\pi)^{k/2}/(\sqrt2\Gamma(k/2))$
and define
\[
 \Pi_l=\prod_{j<l}t_j^2,\quad
 U_{l,N}=\prod_{j<l}[4(N+j+1)^2+t_j^2],\quad
 \omega_N^\Gamma=\log(U_{l,N}/\Pi_l),\quad
 \Omega_{\rm low}^\Gamma=\omega_{q-1}^\Gamma+\omega_q^\Gamma,
 \quad\Omega_{\rm high}^\Gamma=\omega_{2q-1}^\Gamma+\omega_{2q}^\Gamma.
\]
The exact measure identity is
$dm_{0,k}=\beta_k|f_l(c+iy)|^2d\sigma$.
Every factor is $y^2+t_j^2$ on this line. Its lower bound
is $t_j^2$, while the full original multiplication recurrence
at degree $N+j$ gives its upper norm factor
$4(N+j+1)^2+t_j^2$. Iterating proves
$\beta_k\Pi_lH_N^\sigma\preceq H_N^0
\preceq\beta_kU_{l,N}H_N^\sigma$.
The same successive minimum and restriction maps as MRI3b
give this inequality for both actual mixed Grams.
Writing $\Delta_X^\Gamma=F_X^{(1)}-F_X^{(k)}$ therefore gives
\[
\begin{gathered}
 -r_X\Omega_{\rm low}^\Gamma\le\Delta_X^\Gamma
                              \le r_X\Omega_{\rm high}^\Gamma,
 \qquad\Delta_K^\Gamma+\Delta_B^\Gamma=\mathcal H_k,\\
 \mathcal H_k-r\Omega_{\rm high}^\Gamma
       \le\Delta_B^\Gamma\le
               \mathcal H_k+r\Omega_{\rm low}^\Gamma,\\
 \mathcal H_k+\delta_k^\sigma-r\widetilde E_k^-
       \le F_B^{\rm ar}-F_B^{(k)}\le
                  \mathcal H_k+\delta_k^\sigma+r\widetilde E_k^+.
\end{gathered}\tag{MRI5c}
\]
Here the original AMT31 widths are exactly
$\widetilde E_k^+=E_k^++\Omega_{\rm low}^\Gamma$ and
$\widetilde E_k^-=E_k^-+\Omega_{\rm high}^\Gamma$,
with $E_k^+=L_{q-1}+L_q$ and $E_k^-=L_{2q-1}+L_{2q}$.
For a direct sign check set
$z_{X,N}=\log(\det H_{X,N}^0/\det H_{X,N}^\sigma)
-r_X\log(\beta_k\Pi_l)$.
Then $0\le z_{X,N}\le r_X\omega_N^\Gamma$ and
$\Delta_X^\Gamma=-z_{X,q-1}-z_{X,q}+z_{X,2q-1}+z_{X,2q}$.
This proves the first line of MRI5c. The determinant split
gives its exact common sum and then the second line.
Combining the arithmetic and Gamma source comparisons gives
$-r\widetilde E_k^+\le F_K^{\rm ar}-F_K^{(k)}
\le r\widetilde E_k^-$. Subtraction from
$F_B^{\rm ar}-F_B^{(k)}=\mathcal H_k+\delta_k^\sigma
-(F_K^{\rm ar}-F_K^{(k)})$ proves the last line.

At $m=1$, $|u|\ge R_*$, the proved $r=8k-16$ and
$\omega_N^\Gamma\le l\log[1+16(N+l)^2]$ yield
\[
\begin{gathered}
 F_K^{(1)}-F_K^{(k)}=O_h(k^2\log(q+1))=o(kq),\\
 F_K^{\rm ar}-F_K^{(1)}=O_h(k^2)=o(kq),\qquad
 F_K^{\rm ar}-F_K^{(k)}=o(kq),\\
 \frac{F_B^{(1)}-F_B^{(k)}}{kq}\longrightarrow-\frac{\mathcal C_\Gamma}{4},
 \qquad\frac{F_B^{\rm ar}-F_B^{(k)}}{kq}
                                      \longrightarrow-\frac{\mathcal C_\Gamma}{4},\\
 \frac{F_B^{\rm ar}-F_B^{(1)}}{kq}\longrightarrow0,\qquad
 \frac{F_B^{(1)}-F_B^{(k)}}{\mathcal H_k},\quad
 \frac{F_B^{\rm ar}-F_B^{(k)}}{\mathcal H_k}\longrightarrow1.
\end{gathered}\tag{MRI5d}
\]
Indeed $q=(k+1)^2$, so $k\log(q+1)/q\to0$.
MRI5b is bounded by $rE_k^\pm=O_h(k^2)$, and MRI5c
is bounded by $r\Omega^\Gamma=O_h(k^2\log(q+1))$.
The already proved $\mathcal H_k/(kq)\to-\mathcal C_\Gamma/4$
and $\delta_k^\sigma/(kq)\to0$ then prove every boundary
limit from the exact sums. Since $\mathcal C_\Gamma>0$,
$\mathcal H_k$ is eventually nonzero and division proves
the last two ratios. Every bound is uniform over this entire
proved period domain and all original branches. The absolute
quantity $F_K^{(1)}/(kq)$ remains a common original quantity
up to the proved $o(1)$ differences; no limit value is assigned
to it. Its exact finite form is MRI4a, and its full source,
action and period data are MRI1--3.


In MRI5 each $F_B^{(s)}$ is now the actual period-column determinant
ratio MRI4a, and each $F_K^{(s)}$ is its displayed original
section-corrected kernel sum. The substitution is exact because
MRI1a and MRI2a identify both terms with the original MRI3
contractions at every degree. In particular the tensor-source
sum receives $\mathcal H_k+\delta_k^\sigma$ and the fixed-source
sum receives $\delta_k^\sigma$, with both transitions unchanged.
The complete coupled recurrence MRI2b and primary coefficient
formula OPG14 supply their original inputs. QRI1--9 continues
to bound the same signed Gamma transition and retains the
separate actual arithmetic scalar in QRI8.

Indeed MRI4 identifies the two sums with $\mathcal B_k^0$ and
$\mathcal B_k^\sigma$, respectively, while BRI2 gives
$\mathcal H_k=\mathcal B_k^\sigma-\mathcal B_k^0$.
This proves the exact map of the new row formulas to the same
arithmetic scalar and therefore to QRI4--8. Every primary
coefficient and Taylor unit in $\Lambda$ remains in $R_s$.
No new source order or multiplying factor is selected.

\subsection{The exterior Gamma law returned through its full relation mass}

The original input source remains $dm_{0,s}$ with $s\in\{1,k\}$,
mass $M_s=(2\pi)^{s/2}$, and parameter $b=s/2$. For the actual
relation degree $r\in\{1,q,q+1\}$ retain the exterior measure
\[
 d\lambda_{\chi,s,r}(y)=\frac1{r!}
       \prod_{i<j}(y_j-y_i)^2
       \prod_{i=1}^r|\chi(c+iy_i)|^2\,dm_{0,s}(y_i),
 \qquad \nu_{\chi,s,r}=(\Sigma_r)_*\lambda_{\chi,s,r},
 \quad \Sigma_r(y)=\sum_i y_i.
 \tag{IRI1}
\]
IGO7 proves that this is the squared modulus of the original
relation exterior vector, including its phase
$i^{r(r-1)/2}$ and factor $1/\sqrt{r!}$. Its total mass is the
original relation Gram determinant $Z_{s,r}$ of IFB17. At $r=0$
the space is the unit point, the measure is $\delta_0$, and the
mass is one.

For $r\ge1$ put $B=br+r(r-1)$. IGO9--11 constructs the exact
polynomial $D_{\chi,s,r}(z)=\sum_{j=0}^{qr}d_jz^{2j}$ by the
displayed differential recurrence from every coefficient of
$|\chi(c+iy)|^2$. Its full signed constituent and positive
polynomial formulas are
\[
\begin{gathered}
 a_h=\sum_{j=h}^{qr}(-1)^{j+h}\binom jh d_j,\qquad
 R_{\chi,s,r}(X)=\sum_{h=0}^{qr}a_h
       \frac{\Gamma(B)}{\Gamma(B+2h)}
       \prod_{j=0}^{h-1}[X+(B+2j)^2],\\
 d\nu_{\chi,s,r}(Y)
 =\sum_{h=0}^{qr}\frac{M_s^ra_h}{(2\pi)^{B+2h}}
                             dm_{0,2B+4h}(Y)
 =\frac{M_s^r}{(2\pi)^B}
                  R_{\chi,s,r}(Y^2)dm_{0,2B}(Y).
\end{gathered}\tag{IRI2}
\]
IGO15--21 proves the first equality by the finite binomial expansion
of $D(i\tanh\xi)$ in the exact characteristic transform; the Gamma
recurrence proves the second equality, with all powers of two
accounted for. The polynomial has degree $qr$ and is strictly
positive on $X\ge0$ by the original positive fibre integral.
The input source orders have not been selected anew. The produced
orders are $sr+2r(r-1)+4h$, with every $a_h$ retained. Each
constituent on its natural line $U_h=B+2h+iY$ is pulled back to
the actual sum line $T=rc+iY$ by
$U_h=T+(B+2h-rc)$. This translation and its inverse preserve every
polynomial degree and have coefficient determinant one.

The total mass, including all signed cancellations, is
\[
\begin{gathered}
 Z_{s,r}=M_s^rD_{\chi,s,r}(0)
        =M_s^r\sum_{h=0}^{qr}a_h,\\
 D_{\chi,s,r}(0)
   =\left(\prod_{j=0}^{r-1}(q+j)!(b)_{q+j}\right)\Omega_{s,r},
 \qquad
 \Omega_{s,r}=\det\left(I_q+\sum_{j<r}v_jv_j^*\right).
\end{gathered}\tag{IRI3}
\]
Here $v_j$ are precisely the original remainder coordinates MRI1.
To prove the second line, IFB9 identifies the full relation space
in orthonormal source coordinates as the graph $(-V_rz,z)$.
Its monic high-coordinate matrix is triangular with squared
determinant $\prod_{j<r}h_{q+j}^{(s)}$; the graph Gram is
$I_r+V_r^*V_r$. Taking determinants, using
$\det(I_r+V_r^*V_r)=\det(I_q+V_rV_r^*)$, and then substituting
$h_{q+j}^{(s)}=M_s(q+j)!(b)_{q+j}$ proves IRI3 from the first
line. Every source mass is present before this equality cancels it.
At $r=0$ both products and both determinants are one.

Return these actual masses to the same four original quotient
degrees. IFB30 and IRI3 give
\[
\begin{gathered}
 \mathcal B_k^{(s)}
   =S_q^{(s)}+\log Z_{s,q}+\log Z_{s,q+1}-\log Z_{s,1},\\
 S_q^{(s)}=-\log h_q^{(s)}
       -2\sum_{j=q+1}^{2q-1}\log h_j^{(s)}-\log h_{2q}^{(s)},\\
 \mathcal B_k^{(s)}
   =\log\Omega_{s,q}+\log\Omega_{s,q+1}-\log\Omega_{s,1}
   =\mathcal W_q\bigl(\log(1+t_j)\bigr).
\end{gathered}\tag{IRI4}
\]
The two high monic norm products divided by the scalar low norm
contain $h_q$ once, $h_{q+1},\ldots,h_{2q-1}$ twice, and
$h_{2q}$ once. These are exactly the factors with negative signs
in $S_q^{(s)}$; their full source mass powers are $2q$ and $-2q$.
This proves the last line before any asymptotic passage. The final
equality telescopes $\Omega_{s,j+1}/\Omega_{s,j}=1+t_j$ with
the same weight-one and weight-two endpoints as MRI4. Thus the
intrinsic exterior law returns to the actual baseline already
evaluated in BRI5--6, through the full original polynomial mass.
MRI3--5 then gives its individual mixed factors and arithmetic map.

For completeness the operation also retains its Hilbert-space
kernel. The pullback $U_rF=F\circ\Sigma_r$ is an isometry from
$L^2(\nu_{\chi,s,r})$ to $L^2(\lambda_{\chi,s,r})$. If
$w(y)$ is the complete Lebesgue density of IRI1 and
$p(Y)=\int_{\mathbb R^{r-1}}w(u,Y-\sum u)du$, its exact adjoint is
\[
 (U_r^*f)(Y)=\frac{\int_{\mathbb R^{r-1}}
                 f(u,Y-\sum u)w(u,Y-\sum u)\,du}{p(Y)},
 \qquad
 L^2(\lambda_{\chi,s,r})
       =U_rL^2(\nu_{\chi,s,r})\oplus\ker U_r^*.
 \tag{IRI5}
\]
IGO22--26 proves $p(Y)>0$ and the fibre Cauchy--Schwarz bound,
which proves the adjoint and the orthogonal decomposition. The
kernel is exactly the functions with zero displayed weighted fibre
integral almost everywhere. Multiplication by the original exterior
vector gives the isometry to the exterior Hilbert space and the
full adjoint in IGO27--28, including its original phase. For $r=1$
and $r=0$ the maps are the identities on the stated original spaces.
The original spectral sum $T=rc+iY$ intertwines on the exact domain
$\{F:\int |YF(Y)|^2d\nu(Y)<\infty\}$, as proved at the end of IGO.
These maps specify the original relation, the sum, and their retained
kernel throughout the receiving calculation.

\subsection{The current original-source interval receiver}
On BHR1 define the actual current receiving endpoints
by the intersection of the retained BHR10 and BHR11
intervals and the universal WRC5 interval with the
exact-low trace endpoints of JSR24a:
\[
\begin{gathered}
 L_k^{\rm cur}=\max\{\mathscr L_k,H_k^*-b_k^-,
 H_k^B-b_{B,k}^-,-W_k-b_k^--a_0,
                  \mathsf L_k^\sigma-\mathcal B_k^0\},\\
 U_k^{\rm cur}=\min\{\mathscr U_k,H_k^*+b_k^+,
 H_k^B+b_{B,k}^+,-W_k+b_k^++b_0,
                  \mathsf U_k^\sigma-\mathcal B_k^0\},\\
 L_k^{\rm cur}\leq\mathcal H_k\leq U_k^{\rm cur},
 \qquad0\leq U_k^{\rm cur}-L_k^{\rm cur}
                       \leq b_k^-+b_k^+=o(q).
\end{gathered}\tag{BHR12}
\]
Each interval contains the very same scalar by the
proved morphisms JSR8--15, BHR4--7 and WRC3--5.
WRC6 proves that the fourth lower entry is at most the
second, and the fourth upper entry is at least the second.
Thus the first four entries reproduce the previous exact
endpoints. BRI2--3 proves that the fifth interval contains the
same original scalar on BHR1. Its intersection preserves the
earlier width bound and may strengthen the actual finite endpoints.
The evaluated QRI3 and LRI4 intervals contain the exact-$W_k$
interval and retain every finite high and low error. Their explicit
values are available alongside this tighter intersection.
For tensor orders outside BHR1 define
$L_k^{\rm cur}=\mathscr L_k$ and
$U_k^{\rm cur}=\mathscr U_k$ instead; all original
finite formulas remain valid there.

This gives the following literal arithmetic return, with
the original source deficit quantities of JSR18--21:
\[
\begin{gathered}
 \mathcal B_k^0+L_k^{\rm cur}-a_k^-
 \leq\mathcal B_k^{\rm ar}
 \leq\mathcal B_k^0+U_k^{\rm cur}+a_k^+,\\
 \mathcal B_k^0+L_k^{\rm cur}-D_k^-
 \leq\mathcal B_k^{\rm ar}
 \leq\mathcal B_k^0+U_k^{\rm cur}+D_k^+,
 \qquad 0\leq a_k^\pm\leq D_k^\pm.
\end{gathered}\tag{BHR13}
\]
Add the original arithmetic interval to BHR12 and use
$\mathcal B_k^{\rm ar}=\mathcal B_k^0+
\mathcal H_k+\delta_k^\sigma$ to prove both lines.
The full Hankel centre gives in particular
\[
\begin{gathered}
 -b_k^--a_k^-\leq
 \Delta_k^\Gamma-H_k^*=\delta_k^\sigma-e_q-e_{q+1}
 \leq b_k^++a_k^+,\\
 -(64m+245)\leq\liminf\frac{\Delta_k^\Gamma-H_k^*}{q\log k}
 \leq\limsup\frac{\Delta_k^\Gamma-H_k^*}{q\log k}
 \leq128m+490\quad(m\geq2),\\
 \limsup\frac{(U_k^{\rm cur}+a_k^+)
                       -(L_k^{\rm cur}-a_k^-)}{q\log k}
 \leq192m+735\quad(m\geq2).
\end{gathered}\tag{BHR14}
\]
The same finite bounds hold with $a_k^\pm$ replaced by
$D_k^\pm$. Since $b_k^\pm=o(q)$, they vanish after
division by $q\log k$. The signed ODD limits in JSR21
then prove the second line. For the last, add the two
finite allowances and use BHR12; the limit of
$(D_k^-+D_k^+)/(q\log k)$ is $192m+735$.
For $m=1$, retain instead
$(\Delta_k^\Gamma-H_k^*)/(kq)
=O_h((\log k/k)^{5/7})$; the $O(k)$ Gamma comparison
error divided by $kq$ is $O(q^{-1})$, which is smaller
on this actual degree. No multiple-zero constant is
assigned to the simple-zero regime.

Finally the full original HC threshold returns as
\[
\begin{gathered}
 \max\{0,\mathcal B_k^0+L_k^{\rm cur}-D_k^-
                               -4q\log(D_hk)\}
 \leq\mathcal R_k\\
 \leq\mathcal B_k^0+U_k^{\rm cur}+D_k^+
                               -4q\log(D_hk).
\end{gathered}\tag{BHR15}
\]
Subtract its unchanged threshold from BHR13 and use
the original independent constraint $\mathcal R_k\geq0$.
This computes the full-source approximation error at
$o(q)$ while retaining the complete low moment ratio,
the universal Hankel determinants and the actual
arithmetic defect. The exact relation $H_k^*=-W_k+e_0$
and the now proved GEL18 limit give
$H_k^*/(lq)\to-\mathcal C_\Gamma$.
GRI4--7 below returns the evaluated finite bounds and
complete signed limit to the same arithmetic and HC objects,
including their necessary baseline relation.

\subsection{The evaluated full Gamma centre and its original arithmetic return}

Keep every original packet parameter, the original source line
$S=c+iy$, the divided-jet quotient and boundary maps, and the literal
moment index used above. The complete WGP, RWB, EIQ, GEL and WGR
proofs are included below as whole source bodies. They now evaluate
the growth of the very same $W_k$, before returning its sign through
$\mathcal H_k=-W_k+e_0-e_q-e_{q+1}$.

For all original integers $l,m\geq1$, set
\[
\begin{aligned}
 w^-_{q,l}=\max\bigg\{\frac{lq}{64},\;&
 4lq\log\frac{27}{8\pi}-4l^2-6l\log(3/2)-3l\\
 &-\frac{9l}{4q-1}-\frac{2l(l+1/4)}q\bigg\},\qquad
 w^+_{q,l}=6lq,\\
 &\frac{lq}{64}\leq w^-_{q,l}\leq W_k\leq w^+_{q,l}=6lq.
\end{aligned}\tag{GRI1}
\]
These are the finite original-polynomial bounds RWB20--27, including
their exact lower remainder. They need no large-$k$ source threshold.
The original source threshold $k\geq2048\sqrt{\delta^2+\gamma^2}$
is retained when adding the errors of the $\chi$-source comparison.

The growing limit has an explicit equilibrium integral. Define
\[
 K(\kappa)=\int_0^{\pi/2}(1-\kappa^2\sin^2\theta)^{-1/2}\,d\theta,
 \qquad
 E(\kappa)=\int_0^{\pi/2}(1-\kappa^2\sin^2\theta)^{1/2}\,d\theta.
\]
Let $0<u<v$ and $0<\kappa<1$ be the unique solution of
$uK(\kappa)=2$, $vE(\kappa)=4$, $\kappa^2=1-u^2/v^2$.
On $u^2<x<v^2$ put
\[
 \rho_{2,\pi}(x)=
 \frac{\sqrt{(v^2-x)(x-u^2)}}{4\pi x}
 \int_0^\infty
 \frac{\sqrt s\,ds}{(x+s)\sqrt{(u^2+s)(v^2+s)}},
\]
and put this density equal to zero elsewhere. The full EIQ1--33
proof establishes its positivity, mass one, unique endpoints,
variational inequality on the entire positive half-line and the
required logarithmic integral. For each fixed original packet
$(\delta,\gamma,m)$, GEL18--18a proves
\[
\begin{gathered}
 \frac{W_k}{lq}\longrightarrow\mathcal C_\Gamma,
 \qquad \mathcal C_\Gamma=
 2\int_{u^2}^{v^2}\log x\,\rho_{2,\pi}(x)\,dx
                                      -4(2\log2-1),\\
 4\log\frac{27}{8\pi}\leq\mathcal C_\Gamma\leq6,
 \qquad \mathcal C_\Gamma>0.
\end{gathered}\tag{GRI2}
\]
Here the elliptic modulus $\kappa$ is EIQ's local $k$, and its
midpoint $(u^2+v^2)/2$ is EIQ's local $m$; neither changes the
original packet order or multiplicity. GEL uses the ensemble
parameters $\alpha=a/n$, $\beta=\pi q/(2n)$ tending to $(2,\pi)$.
Its exact maps are the full pushforward $t=y^2$ and the literal
change $t=q^2x$. In particular GEL3 retains the factor
$q^{2an+2n(n-1)}$ in $H_n^{(a)}$, both halves of the source and
its complete Gamma mass. GEL14 retains all four parity blocks.
The original scalar moment ratio is $m_{2q+2l}/m_{2q}$, equal
to GEL's even-moment ratio $\mu_{q+l}/\mu_q$. Thus the finite
expression for $e^{W_k}$ above is exactly the object to which
GRI1--2 apply. Its leading remainder is $o(lq)$; the complete
WGP13 and GEL17b finite terms remain available at the smaller
$q$ scale, with no replacement of this remainder by $o(q)$.

For the original Gamma errors retain $A=b_k^+$ and $B=b_k^-$;
these are the positive error constants of WGR3, while the original
boundary target is $B_{\mathrm{obs}}=\operatorname{im}\Lambda$.
Adding the finite errors $e_0\in[-a_0,b_0]$ and
$e_q+e_{q+1}\in[-b_k^+,b_k^-]$ gives the evaluated return
\[
 -w^+_{q,l}-b_k^--a_0\leq\mathcal H_k
                        \leq-w^-_{q,l}+b_k^++b_0.
 \tag{GRI3}
\]
The proof is direct substitution in
$\mathcal H_k=-W_k+e_0-e_q-e_{q+1}$: the low error enters
positively and the entire high sum negatively. Since
$w^-_{q,l}\leq W_k\leq w^+_{q,l}$, this evaluated interval
contains the earlier interval
$[-W_k-b_k^--a_0,-W_k+b_k^++b_0]$. It therefore also
contains the existing finite four-interval intersection.
The current exact maximum and minimum remain unchanged;
GRI3 supplies explicit bounds requiring no unevaluated $W_k$.

The original quotient, kernel and boundary forms are transported
by precisely the remainder map $J_N:\mathbb C[S]_{\leq N}\to E$
and the fixed observation $\Lambda:E\to B_{\mathrm{obs}}$.
WGR7--9 proves in the common coefficient frame $[I_K,S_*]$ that
$\det H_N^K\det Q_N^B=|\det[I_K,S_*]|^2\det G_N$;
the common frame factor cancels between the two original sources.
It retains the signed low/high deficit identity
$\delta_k^\sigma=-b_{q-1}-b_q+b_{2q-1}+b_{2q}$.
Consequently the same original arithmetic allowances give
\[
\begin{gathered}
 \mathcal B_k^0-w^+_{q,l}-b_k^--a_0-a_k^-
 \leq\mathcal B_k^{\rm ar}
 \leq\mathcal B_k^0-w^-_{q,l}+b_k^++b_0+a_k^+,\\
 \mathcal B_k^0-w^+_{q,l}-b_k^--a_0-D_k^-
 \leq\mathcal B_k^{\rm ar}
 \leq\mathcal B_k^0-w^-_{q,l}+b_k^++b_0+D_k^+.
\end{gathered}\tag{GRI4}
\]
Add $\delta_k^\sigma\in[-a_k^-,a_k^+]$ to GRI3 and use
$\mathcal B_k^{\rm ar}=\mathcal B_k^0+\mathcal H_k+\delta_k^\sigma$
for the first line, then use the unchanged $a_k^\pm\leq D_k^\pm$
for the second. No source arithmetic allowance has been enlarged
or assumed to vanish at finite degree.

For fixed $m=1$, the original arithmetic estimate is
$\delta_k^\sigma/(kq)=O_h((\log k/k)^{5/7})$.
For each fixed $m\geq2$, its original signed constants are
$\liminf\delta_k^\sigma/(q\log k)\geq-(64m+245)$ and
$\limsup\delta_k^\sigma/(q\log k)\leq128m+490$.
These two finite bounds imply eventual boundedness at that
scale; multiplication by $\log k/k$ proves
$\delta_k^\sigma/(kq)\to0$. The simple-multiplicity rate
also tends to zero. The Gamma errors are $O(k)$ at $m=1$
and $O(1)$ at fixed $m\geq2$, hence their quotient by $kq$
tends to zero in both original branches. As $l/k\to1/4$,
GRI2 now gives the actual full signed limit
\[
 \frac{\Delta_k^\Gamma}{kq}
 =-\frac{l}{k}\frac{W_k}{lq}
       +\frac{e_0-e_q-e_{q+1}+\delta_k^\sigma}{kq}
 \longrightarrow-\frac{\mathcal C_\Gamma}{4}<0.
 \tag{GRI5}
\]
Thus this original Gamma/arithmetic difference is negative
at all sufficiently large original degrees, with every
retained boundary and kernel direction still in its exact map.

On the combined original source and HC domains, subtract the
same $4q\log(D_hk)$ and use the independent $\mathcal R_k\geq0$:
\[
\begin{aligned}
 \max\{0,\mathcal B_k^0-w^+_{q,l}-b_k^--a_0-D_k^-
                                      -4q\log(D_hk)\}
 &\leq\mathcal R_k\\
 &\leq\mathcal B_k^0-w^-_{q,l}+b_k^++b_0+D_k^+
                                      -4q\log(D_hk).
\end{aligned}\tag{GRI6}
\]
Here the unchanged $D_h>0$ is fixed on the original packet.
The identities $\mathcal B_k^0-\mathcal B_k^{\rm ar}=-\Delta_k^\Gamma$
and $\mathcal R_k-\mathcal B_k^0=\Delta_k^\Gamma-4q\log(D_hk)$
therefore prove
\[
\begin{gathered}
 \frac{\mathcal B_k^0-\mathcal B_k^{\rm ar}}{kq}
                  \longrightarrow\frac{\mathcal C_\Gamma}{4},\qquad
 \frac{\mathcal R_k-\mathcal B_k^0}{kq}
                  \longrightarrow-\frac{\mathcal C_\Gamma}{4},\\
 \liminf\frac{\mathcal B_k^0}{kq}\geq\frac{\mathcal C_\Gamma}{4}.
\end{gathered}\tag{GRI7}
\]
Indeed $4\log(D_hk)/k\to0$, and nonnegativity gives
$\mathcal B_k^0/(kq)\geq(\mathcal B_k^0-\mathcal R_k)/(kq)$
before the lower limit is taken. This is the necessary
baseline and residual relation of WGR15, with their original
objects retained. The next finite-scale calculation concerns
those objects and the full $q$-scale remainder; the proved
negative difference does not set either object equal to zero.


\section{The finite arithmetic majorant at each original degree}

The original arithmetic density is
$w_h(y)=|v_h(1/2+iy)|^2/(2\pi)$, $m_k=w_h^{*k}$.
Retain the proved TW/DMR positive constants and the original
source determinant deficit
\[
 \Xi_{k,N}=(N+1)\log A_k
       -\log\frac{\det H_{k,N}}{\det H_N^\sigma},\qquad
 A_k=\frac{C_h}{c_\Gamma}k^{p+1}M_\alpha^{k-1}.
 \tag{JSR16}
\]
Here $C_h$ is the original tilted-envelope constant. For the
finite bound all required constants are
\[
\begin{gathered}
 p=8m-\tfrac92,\quad B=42+8m,\quad M_0=21+4m,\quad
 \eta=(p-1)/p,\quad\alpha=\pi/2,\quad g_0=\log\sqrt{2\pi},\\
 \vartheta_h=\int_{-1}^1w_h(y)\,dy,\quad
 M_1=\int_{-1}^1e^{\alpha y}w_h(y)\,dy,\quad
 M_\alpha=\int_{\R}e^{\alpha y}w_h(y)\,dy,\\
 b_h=\frac{c_h}{C_\Gamma2^{3B/2}},\quad
 C_r=1+\frac{\alpha+\log(M_\alpha/\vartheta_h)+B}{\log2},\quad
 C_d=2C_r+M_0/3,\quad C_L=1+C_d,\\
 C_E=6C_r^2+(2M_0+4)C_d/3,\quad
 C_s=\left|\log\frac{C_h}{c_\Gamma b_h}+2\log M_\alpha\right|,
 \quad C_{\rm tail}=\frac{2C_h}{(p-1)M_1},\\
 T_k^{\rm deg}=\log(4m(C_L+2))+3\log(k+2),\quad
 z_k=\frac{kT_k^{\rm deg}}q,\quad R_k^*=\max(1,z_k^{-1/p}).
\end{gathered}\tag{JSR17}
\]
The letter $M_0$ denotes the original DMR exponent, and does
not replace the joint source degree $M=N+l$.
The Gamma-envelope constants remain
$c_\Gamma=\sqrt2e^{-7/6}$ and
$C_\Gamma=\sqrt{2\sqrt5}e^{2/3}$.
Every constant except $T_k^{\rm deg},z_k,R_k^*$ is fixed
as $k$ varies. The tilted finite mass $M_1$ and the
un-tilted finite mass $\vartheta_h$ remain their separate integrals.

For each $N\in\{q-1,q,2q-1,2q\}$ put
\[
\begin{aligned}
 \Omega_{k,N}^{\rm deg}={}&\frac{N+1}{q}
 \left\{C_{\rm tail}\min(2^{1-p},z_k^\eta)
       +\frac{(p+1)\log k+C_s}{k}\right\}
       +\frac{2M_0NT_k^{\rm deg}}{kq}\\
 &+C_Ez_k+C_d\max(z_k,z_k^\eta)+\frac{(g_0+2)C_d}{q},\\
 D_k^-={}&kq(\Omega_{k,q-1}^{\rm deg}+\Omega_{k,q}^{\rm deg}),\qquad
 D_k^+=kq(\Omega_{k,2q-1}^{\rm deg}+\Omega_{k,2q}^{\rm deg}).
\end{aligned}\tag{JSR18}
\]
This scalar $\Omega_{k,N}^{\rm deg}$ is not the positive
matrix $\Omega_N$ in JSR7. Their domains and definitions
are explicit. The full finite source statement is
\[
 0\leq\Xi_{k,N}\leq kq\Omega_{k,N}^{\rm deg}.
 \tag{JSR19}
\]
To verify its substitution from the already proved complete
DMR source determinant bound, keep the exact factors in ODD7:
\[
\begin{aligned}
 \frac{\Xi_{k,N}}{kq}\leq{}&\frac{N+1}{q}
 \left\{C_{\rm tail}(1+R)^{1-p}
               +\frac{(p+1)\log k+C_s}{k}\right\}
 +\frac{2M_0NT_k^{\rm deg}}{kq}\\
 &+\frac{kT_k^{\rm deg}}q(C_E+C_dR)
 +\frac{(g_0+2)C_d}{q},\qquad R\geq1.
\end{aligned}
\]
At $R=R_k^*$ one has
$z_kR_k^*=\max(z_k,z_k^\eta)$ and
$(1+R_k^*)^{1-p}\leq\min(2^{1-p},z_k^\eta)$.
For $z_k\leq1$ these follow from $R_k^*=z_k^{-1/p}$;
for $z_k\geq1$ they follow from $R_k^*=1$ and
$2^{1-p}\leq1\leq z_k^\eta$.
This proves JSR19 including the common boundary $z_k=1$.
The complete DMR proof and the unchanged ODD provider are
preserved with this source; the inequality is their exact
finite degree-dependent receiving substitution.

Write $b_N$ for the original quotient deficit and
$b_N^K,b_N^B$ for its fixed original kernel and boundary
deficits. The full source/quotient Schur decomposition gives
$0\leq b_N^K,b_N^B\leq b_N\leq\Xi_{k,N}$ and
$b_N^K+b_N^B=b_N$. To see the precise maps, for either
source use $G_N=(J_NH_N^{-1}J_N^*)^{-1}$,
$H_N^K=I_K^*G_NI_K$ and
$Q_N^B=(\Lambda G_N^{-1}\Lambda^*)^{-1}$.
The common coefficient frame $[I_K,S_*]$ has the same
determinant in both sources; its orthogonal lift differs
by a kernel column and a triangular matrix of determinant
one. Consequently
$\det H_N^K\det Q_N^B=|\det[I_K,S_*]|^2\det G_N$.
The quotient Schur factor and its kernel restriction are
positive contractions after the same full reference
congruence, proving the nonnegative deficit decomposition.
Empty factors at rank zero have determinant one.

The original four signs give the sharper finite interval
\[
\begin{gathered}
 a_k^-=\Xi_{k,q-1}+\Xi_{k,q}\leq D_k^-,\qquad
 a_k^+=\Xi_{k,2q-1}+\Xi_{k,2q}\leq D_k^+,\\
 -D_k^-\leq-a_k^-\leq X_k\leq a_k^+\leq D_k^+,\\
 \qquad X_k\in\{\delta_k^\sigma,\Delta F_K,\Delta F_B\}.
\end{gathered}\tag{JSR20}
\]
Indeed $X_k=-b_{q-1}-b_q+b_{2q-1}+b_{2q}$, using the
corresponding superscript for either original summand.
Drop the positive two terms for the lower bound and the
negative two for the upper bound, then use JSR19.
The original source inclusions into degree $2q$ give
$\Xi_{k,N}\leq\Xi_{k,2q}$, hence both $a_k^\pm$ are
at most $2\Xi_{k,2q}$. Thus the actual interval using
$a_k^-,a_k^+$ supersedes the note's symmetric interval.
The explicit majorants $D_k^-,D_k^+$ remain outer bounds
when smaller exact deficit sums are already known.

For $m\geq2$ let $C_m=p+1+6M_0=32m+245/2$.
At an original endpoint with $N/q\to s\in\{1,2\}$,
\[
\begin{gathered}
 \frac{k\Omega_{k,N}^{\rm deg}}{\log k}\longrightarrow sC_m,
 \quad\limsup\frac{\Xi_{k,N}}{q\log k}\leq sC_m,\\
 \frac{D_k^-}{q\log k}\longrightarrow64m+245,\qquad
 \frac{D_k^+}{q\log k}\longrightarrow128m+490,\\
 -(64m+245)\leq\liminf\frac{X_k}{q\log k}
 \leq\limsup\frac{X_k}{q\log k}\leq128m+490.
\end{gathered}\tag{JSR21}
\]
The proof retains $q=(m-1)k^3+O_m(k^2)$ and $\eta>1/2$.
Since $T_k^{\rm deg}/\log k\to3$,
$kz_k^\eta/\log k
=O_m(k^{1-2\eta}(\log k)^{\eta-1})\to0$,
$kz_k/\log k=O_m(k^{-1})$, and
$k/(q\log k)\to0$, the two surviving terms in JSR18
have limits $s(p+1)$ and $6M_0s$.
This proves the first line; summing the two low and two
high endpoints separately proves the rest.
For $m=1$ retain $q=(k+1)^2$ and $\eta=5/7$.
The same finite formulas instead give
$X_k/(kq)=O_h((\log k/k)^{5/7})$.
No $q\log k$ bound in JSR21 is assigned to that branch.

\section{The actual finite trace interval receives the asymmetric bounds}

Retain the note's common scalar and all eight metrized matrices
\[
 s_k=2+\sum_{a=0}^3\Tr X_{N_a},\qquad
 Z_{N_a}=I-(I+X_{N_a})/s_k,
 \qquad Y_{N_a}\quad(0\leq a\leq3).
 \tag{JSR22}
\]
Because each eigenvalue of $X_N$ is nonnegative and at
most its trace, $Z_N$ has eigenvalues in $[0,1)$ in
the original $\Omega_N$ metric. The same assertion for
$Y_N$ in its original $A_{N+l}$ metric was proved in JSR7.
For any one of these matrices $T$ define
\[
 b_p(T)=\sum_{n=1}^p\frac{\Tr T^n}{n},\quad
 a_p(T)=\Tr T^p,\quad
 u_p(T)=b_p(T)+\frac{b_{2p}(T)-b_p(T)}{1-a_p(T)}.
 \tag{JSR23}
\]
At every admitted positive integer order $p$ with $a_p(T)<1$,
$b_p(T)\leq-\log\det(I-T)\leq u_p(T)$.
Here is the finite-tail proof: for each eigenvalue
$0\leq x<1$, let $r_p(x)=\sum_{n>p}x^n/n$.
Then $r_{2p}(x)\leq x^pr_p(x)$ by shifting its index.
Summing the nonnegative terms gives
$\sum r_p\leq b_{2p}-b_p+a_p\sum r_p$,
which proves the upper bound. Since every eigenvalue is
strictly less than one, $a_p\to0$ and $u_p-b_p\to0$;
in particular every prescribed positive tolerance admits
a finite order.

For each endpoint take admitted orders $p_a,q_a$ and put
\[
\begin{gathered}
 L_a=b_{p_a}(Z_{N_a})+b_{q_a}(Y_{N_a}),\quad
 U_a=u_{p_a}(Z_{N_a})+u_{q_a}(Y_{N_a}),\\
 \mathscr L_k^{\rm note}=L_2+L_3-U_0-U_1,\qquad
 \mathscr U_k^{\rm note}=U_2+U_3-L_0-L_1.
\end{gathered}\tag{JSR24}
\]
The four original $l\log s_k$ terms cancel with their
four signs, so JSR14 is
$\mathcal H_k=-\sum_a\epsilon_a
\sum_{n\geq1}(\Tr Z_{N_a}^n+\Tr Y_{N_a}^n)/n$.
Inserting JSR23 with these signs proves
$\mathscr L_k^{\rm note}\leq\mathcal H_k
\leq\mathscr U_k^{\rm note}$.
This finite interval uses the same actual entries and
metrics as the note, GCR and GSR via JSR8--15.

The exact GSR low endpoints give a further contraction of
this particular trace interval. Define their complete sum
and the new finite endpoints by
\[
\begin{gathered}
 A_k^{\rm low}=4\log|\det M_f|-
 \log D_{l,q-1}-\log D_{l,q}+\log\mathfrak r_\chi,\\
 \mathscr L_k=A_k^{\rm low}-2l\log s_k+L_2+L_3,\qquad
 \mathscr U_k=A_k^{\rm low}-2l\log s_k+U_2+U_3,\\
 \mathscr L_k^{\rm note}\leq\mathscr L_k
 \leq\mathcal H_k\leq\mathscr U_k
 \leq\mathscr U_k^{\rm note},\\
 \mathscr U_k-\mathscr L_k=(U_2-L_2)+(U_3-L_3).
\end{gathered}\tag{JSR24a}
\]
To prove this, put
$Q_a=-\log\det(I-Z_{N_a})-\log\det(I-Y_{N_a})$.
Then $\log\mathfrak d_{N_a}=l\log s_k-Q_a$ and
$L_a\leq Q_a\leq U_a$. JSR13 with
$R_{q-1}=1$, $R_q=\mathfrak r_\chi$ proves
$\log\mathfrak d_{q-1}+\log\mathfrak d_q=A_k^{\rm low}$.
Consequently
$\mathcal H_k=A_k^{\rm low}-2l\log s_k+Q_2+Q_3$.
This proves its new endpoints. Also
$Q_0+Q_1=2l\log s_k-A_k^{\rm low}$ lies between
$L_0+L_1$ and $U_0+U_1$, which proves both outer
containments and the exact width formula.
Thus JSR24a is the substitution of the accepted GSR3--6
low-endpoint identities through the exact maps JSR8--13
into the note's actual trace interval. It is the same
original determinant in those three presentations.

Every low-endpoint quantity here has a finite expression
on the original polynomial source. In particular, let
$a_\chi$ be the full degree-$q$ coefficient column of
$\chi$, let $F_q$ be multiplication by $f$ from
$\mathcal P_q$ into $\mathcal P_{q+l}$, and let $P_d$
have as columns the monic polynomials $p_0,\ldots,p_d$
in the original $S$ coordinate. Then
\[
 H_d^\sigma=P_d^{-*}\diag(\gamma_0,\ldots,\gamma_d)P_d^{-1},
 \qquad
 \mathfrak r_\chi=
 \frac{a_\chi^*F_q^*H_{q+l}^\sigma F_qa_\chi}
      {a_\chi^*H_q^\sigma a_\chi}.
 \tag{JSR24b}
\]
The first identity follows from monic orthogonality and
the second from the exact coefficient multiplication map.
Both are finite, positive, and retain the full polynomial
$\chi$ with all multiplicities. Thus substituting these
low endpoints introduces no new unknown source measure.

The updated actual arithmetic receiver is therefore
\[
\boxed{\begin{aligned}
 \mathcal B_k^0+L_k^{\rm cur}-a_k^-
 &\leq\mathcal B_k^{\rm ar}
 \leq\mathcal B_k^0+U_k^{\rm cur}+a_k^+,\\
 \mathcal B_k^0+L_k^{\rm cur}-D_k^-
 &\leq\mathcal B_k^{\rm ar}
 \leq\mathcal B_k^0+U_k^{\rm cur}+D_k^+.
\end{aligned}}\tag{JSR25}
\]
Indeed the unchanged TWA equality is exactly
$\mathcal B_k^{\rm ar}=\mathcal B_k^0+
\mathcal H_k+\delta_k^\sigma$; add the signed original
arithmetic interval to the current Gamma intersection BHR12.
Its endpoints equal the exact-low trace endpoints below
the explicit full-source threshold and otherwise intersect
them with the three proved full-source enclosures.
The fourth total entry, contributed by WRC5, leaves
the prior endpoints exactly unchanged by WRC6; the
additional WRC7 presentation exposes the full universal centre.
The first line is contained in the preceding trace-only
interval because the Gamma lower endpoint increases and
its upper endpoint decreases. The same original
$a_k^-,a_k^+$ are retained. The second line evaluates
those unchanged arithmetic allowances by the same full
finite source constants $D_k^-,D_k^+$; the full-source
Gamma result does not improve the arithmetic deficit itself.
GRI1 and GRI4 now give the corresponding finite evaluated
arithmetic interval using $w^-_{q,l},w^+_{q,l}$, and GRI5
proves its complete signed growing limit. The exact current
finite endpoints remain unchanged because GRI3 contains
the retained exact-$W_k$ interval.
Equivalently these bounds apply to
$\Delta_k^\Gamma=\mathcal B_k^{\rm ar}-\mathcal B_k^0$.
In particular
\[
 \Delta_k^\Gamma-\mathcal H_k
 =\Delta_k^\Gamma+\mathfrak T_k=\delta_k^\sigma.
 \tag{JSR26}
\]
The constants in JSR21 apply to this exact difference.
The full-source calculation additionally gives
$\Delta_k^\Gamma-H_k^*=\delta_k^\sigma-e_q-e_{q+1}$,
with every finite signed error in BHR7 and BHR10.
Thus BHR14 transports the same arithmetic constants to
the explicitly specified Hankel centre with its $o(q)$
Gamma error. GEL18 and GRI5 now evaluate the full return:
$\Delta_k^\Gamma/(kq)\to-\mathcal C_\Gamma/4$ on both
original multiplicity branches. Neither formula omits
the Gamma return or changes its arithmetic allowance.

There is an entirely finite order choice that propagates
the constants to the width of the actual receiver.
For each of the four high-endpoint matrices choose the least positive
integer $p$ satisfying $a_p(T)<1$ and
$u_p(T)-b_p(T)\leq1/4$. The convergence proved after
JSR23 proves that each least order exists. Then
\[
\begin{gathered}
 0\leq\mathscr U_k-\mathscr L_k
   =\sum_{a=2}^3(U_a-L_a)\leq1,\\
 \limsup\frac{
   (U_k^{\rm cur}+a_k^+)-(L_k^{\rm cur}-a_k^-)}{q\log k}
   \leq192m+735\quad(m\geq2),\\
 \lim\frac{
   (U_k^{\rm cur}+D_k^+)-(L_k^{\rm cur}-D_k^-)}{q\log k}
   =192m+735\quad(m\geq2).
\end{gathered}\tag{JSR27}
\]
The current Gamma interval is contained in this exact-low
trace interval, so the finite arithmetic widths are still
bounded by $1+a_k^-+a_k^+$ and $1+D_k^-+D_k^+$ respectively.
Independently of this adaptive trace-order choice, BHR12
already gives $o(q)$ for the Gamma width on the original
growing-family domain, from its complete integral bounds.
The exact width identity in the first line and the
two limits in JSR21 prove both claims. This order choice
requires actual matrix traces; it inserts no unproved
uniform bound for them as the packet grows. When numerical
entries are used, their certified errors must also enter
the finite intervals before these claims are invoked.

Finally the original HC residual retains its constant,
sign and nonnegative constraint:
\[
\begin{gathered}
 \mathcal R_k=\mathcal B_k^{\rm ar}-4q\log(D_hk)\geq0,\\
 \max\{0,\mathcal B_k^0+L_k^{\rm cur}-D_k^-
                         -4q\log(D_hk)\}
 \leq\mathcal R_k\\
 \leq\mathcal B_k^0+U_k^{\rm cur}+D_k^+
                         -4q\log(D_hk).
\end{gathered}\tag{JSR28}
\]
Subtract the same original threshold from JSR25 and
intersect its lower endpoint with the proved nonnegative
constraint. Both the complete arithmetic source and the
entire signed joint-Schur return remain in this formula.
The full signed Gamma value is now evaluated by GRI5.
GRI6 gives the finite RWB substitution into this same HC
formula, and GRI7 proves the two exact difference limits
and $\liminf\mathcal B_k^0/(kq)\geq\mathcal C_\Gamma/4$.
The original baseline and residual now have the evaluated
$q^2$ coefficient $C_B>769/17010$ by BRI5--8, with the same
finite asymmetric packet bounds. QRI3--8 supplies the actual
four-block finite Gamma return and its $q$-scale multiplicity
branches. MRI1--5 now returns the computed actual period rank
and full arithmetic mixed metrics. On the proved simple-quartet
period domain their three kernel coefficients vanish at scale
$q^2$, their three boundary coefficients equal $C_B$, and both
boundary returns to the tensor source have coefficient
$-\mathcal C_\Gamma/4$ at scale $kq$. The sourcewise kernel
differences are $o(kq)$; the absolute common kernel retains
its full finite metric expression for the next calculation.

\paragraph{Exact source dependencies.}
This calculation uses the received \emph{Tau Gamma Joint
Schur Return} note, the complete GCR1--44 source, the complete
GSR1--33/JQR1--39 source, ODD1--14, the accepted full
HMR successor including HMR44a--44c, the complete BRD1--42
proof including BRD41a, the complete CTR1--54 proof,
the complete LET1--21 scalar-source proof, and the complete
PHT1--23/HCT1--39 positive moment and Hankel calculation.
WRC1--7 proves their receiving substitution and exact interval
containment. The complete WGP1--14, RWB1--27, EIQ1--33,
GEL1--19 including GEL17a--b and GEL18a, and WGR1--15
proof bodies are included below. GRI1--7 installs their exact
finite evaluated and asymptotic consequences in this original
receiver, retaining the smaller-scale terms. BRI1--9, LRI1--4,
QRI1--9, MRI1--5 and IRI1--5 now return the complete canonical BSL/HBL,
GDR/LER, QLG/QGT and IGO/IFB/GMB proofs, all included in full
below. No independently chosen Gamma order enters these formulas.
The companion JSON
records their complete source hashes and unchanged copies.
The new argument above proves the coordinate dictionary,
the finite signed receiving substitutions and their
original-degree constants; it does not reattribute the
complete inherited DMR determinant proof to this note.

\clearpage
\part*{Complete proofs of the growing Gamma return}
The following are the complete accepted proof bodies, included without
mathematical alteration. Their exact source pins, the complete cut20
predecessor and all earlier source providers are retained alongside
this complete successor.

% BEGIN RX INLINE 32; SHA256 e6122af0f0cc54028ceee5383c97070e7738d9d8ec1a5bc5481b349761e5fbfe
% Complete proof fragment. Original Gamma source and endpoint roles retained.
\section{The original four-product cancellation with a finite remainder}
\label{sec:WGP}

Retain $k=4l+1$, $l\ge1$, $m\ge1$, $e=1+k(m-1)$ and
$q=e(k+1)^2=2n$. The symbols $e$ and $n$ here are the original
packet multiplicity and the exact half of its degree. Define the original
source products and their signed combination by
\[
 D_{l,N}=\prod_{a=0}^{N}\prod_{j=0}^{2l-1}(a+j+1/2),\qquad
 P_k=\log D_{l,q-1}+\log D_{l,q}
       -\log D_{l,2q-1}-\log D_{l,2q}.
 \tag{WGP1}
\]
The present calculation estimates these literal products, including both
distinct endpoint ranges. In particular it does not replace the two ranges
by a common interval and omit their endpoint correction.

Cancelling identical factors in the four specified finite products gives
\[
 P_k=-\sum_{j=0}^{2l-1}\left\{
 \sum_{a=q}^{2q-1}\log(a+j+1/2)
 +\sum_{a=q+1}^{2q}\log(a+j+1/2)\right\}.
 \tag{WGP2}
\]
Put $x_b=(b+1/2)/q$ for $q\le b\le2q-1$. The change of index
$b=a-1$ in the second sum retains its full endpoints and yields
\[
 \widetilde P_{q,l}:=P_k+4lq\log q
 =-\sum_{b=q}^{2q-1}\sum_{j=0}^{2l-1}
  \left\{\log(x_b+j/q)+\log(x_b+(j+1)/q)\right\}.
 \tag{WGP3}
\]
This defines an explicitly recorded scalar correction to $P_k$; the
original products in WGP1 are still the objects to which it refers.

\paragraph{An elementary quadrature bound with its sign.}
For a twice continuously differentiable real $F$ on $[1,2]$, divide
that interval into the $q$ cells of width $h=1/q$ with centres $x_b$.
The identity
\[
 F(x_b+u)-F(x_b)-uF'(x_b)
 =\begin{cases}
 \displaystyle\int_0^u(u-v)F''(x_b+v)\,dv,&u\ge0,\\
 \displaystyle\int_u^0(v-u)F''(x_b+v)\,dv,&u<0
 \end{cases}
\]
follows by two integrations of $F''$. Its absolute value is at most
$\|F''\|_\infty u^2/2$. Integration over each symmetric cell kills the
linear term, giving
\[
 \left|\sum_b F(x_b)-q\int_1^2F(x)\,dx\right|
 \le\frac{\|F''\|_\infty}{24q}.
 \tag{WGP4}
\]
The remainder has the sign of $F''$ when this sign is constant; hence
the midpoint sum is below the integral multiplied by $q$ for convex $F$
and above it for concave $F$. Applying the formula to the three literal
functions $\log x,x^{-1},x^{-2}$ gives
\[
 \begin{aligned}
 \sum_b\log x_b&=q(2\log2-1)+\varepsilon_0,
       &0\le\varepsilon_0\le(24q)^{-1},\\
 \sum_b x_b^{-1}&=q\log2-\varepsilon_1,
       &0\le\varepsilon_1\le(12q)^{-1},\\
 \sum_b x_b^{-2}&=q/2-\varepsilon_2,
       &0\le\varepsilon_2\le(4q)^{-1}.
 \end{aligned}\tag{WGP5}
\]
Here the second derivative absolute bounds are respectively $1,2,6$.
The integrals are evaluated by their displayed elementary antiderivatives.

For every $u\ge0$, integrating
$1/(1+u)=1-u+u^2/(1+u)$ gives
\[
 \log(1+u)=u-u^2/2+r(u),\qquad
 r(u)=\int_0^u\frac{t^2}{1+t}\,dt,\qquad 0\le r(u)\le u^3/3.
 \tag{WGP6}
\]
Use this exact formula at $u=j/(qx_b)$ and $(j+1)/(qx_b)$.
The finite sums of their numerators are
\[
 \begin{aligned}
 \sum_{j<2l}(j+j+1)&=4l^2,\\
 \sum_{j<2l}(j^2+(j+1)^2)&=(16l^3+2l)/3,\\
 \sum_{j<2l}(j^3+(j+1)^3)&=8l^4+2l^2.
 \end{aligned}\tag{WGP7}
\]
These identities follow by adding the two shifted ranges; they also follow
directly from telescoping the differences of the powers $j^2,j^3,j^4$.
Consequently WGP3 is the exact identity
\[
 \widetilde P_{q,l}
 =-4l\sum_b\log x_b-\frac{4l^2}{q}\sum_bx_b^{-1}
   +\frac{8l^3+l}{3q^2}\sum_bx_b^{-2}-\mathcal R_{q,l},
 \tag{WGP8}
\]
where the original sum of the $2(2l)q$ nonnegative remainders satisfies
\[
 0\le\mathcal R_{q,l}
 \le\frac{8l^4+2l^2}{3q^3}\sum_bx_b^{-3}
 \le\frac{8l^4+2l^2}{3q^2}.
 \tag{WGP9}
\]
The last inequality uses $x_b\ge1$ and the exact number $q$ of cells.
There is no smallness hypothesis on $j/q$ in WGP6--9.

Substituting WGP5 gives the finite expansion
\[
 \boxed{\begin{aligned}
 P_k={}&-4lq\log q-4lq(2\log2-1)-4l^2\log2
             +\frac{8l^3+l}{6q}+\mathcal E_{q,l},\\
 -\mathcal E^-_{q,l}\le{}&\mathcal E_{q,l}\le\mathcal E^+_{q,l},\\
 \mathcal E^-_{q,l}={}&\frac{l}{6q}
       +\frac{8l^3+l}{12q^3}
       +\frac{8l^4+2l^2}{3q^2},\qquad
 \mathcal E^+_{q,l}=\frac{l^2}{3q^2}.
 \end{aligned}}\tag{WGP10}
\]
Indeed the exact error is
$-4l\varepsilon_0+(4l^2/q)\varepsilon_1
 -(8l^3+l)\varepsilon_2/(3q^2)-\mathcal R_{q,l}$.
Every coefficient is positive before its displayed sign is applied;
the inequalities in WGP10 follow separately from WGP5 and WGP9.
Since $q\ge(4l+2)^2$, this error is $O(1)$ at fixed $m=1$.
For any fixed $m\ge2$ it is $O(l^{-2})$ as $l\to\infty$, because
$e=1+(4l+1)(m-1)$ and hence $q$ is of order $l^3$.
In particular the explicit $-4l^2\log2$ term survives on the $q$ scale
at $m=1$ and has been retained. The next term is also displayed with
its sign and original endpoints rather than absorbed into that scale.

\paragraph{Exact interface to the positive original Gamma recurrence.}
The Gamma measure is still
$d\sigma(y)=|\Gamma(1/4+iy/2)|^2dy/(2\pi)$.
For $a,j\ge0$ let $H_j^{(a)}$ and $c_j^{(a)}$ have exactly their
HCT definitions: $H_j^{(a)}=\det[\int y^{2(a+i+r)}d\sigma]_{i,r<j}$
and $c_j^{(a)}=H_{j+1}^{(a+1)}H_j^{(a)}/
(H_{j+1}^{(a)}H_j^{(a+1)})$. All original mass factors are present
in the four determinants in this expression. Define only the scalar
notation $\lambda_j^{(a)}=\log(c_j^{(a)}/q^2)$, keeping the same $q$
for every shift $a$. HCT37 and the exact low identity PHT22 give
\[
 \boxed{\begin{aligned}
 W_k=\widetilde P_{q,l}+\sum_{s=0}^{l-1}\bigg[
  &2\sum_{j=0}^{n-1}\lambda_j^{(q+s)}+\lambda_n^{(q+s)}\\
  &+2\sum_{j=0}^{n-1}\lambda_j^{(q+s+1)}
       -\lambda_0^{(q+s)}\bigg],\quad n=q/2.
 \end{aligned}}\tag{WGP11}
\]
To verify the entire scale cancellation, each bracket has
$2n+1+2n-1=4n=2q$ copies of $2\log q$ before this scalar
notation is introduced. Across the $l$ original shifts this is
$4lq\log q$, exactly the term recorded in WGP3. The two high ranks,
the doubled odd block and the subtracted low endpoint remain visible
in WGP11. The finite remainders in WGP10 therefore enter the same
$W_k$ as in LET19--21, with coefficient $+1$; in the arithmetic
receiver $B_{\rm ar}=B_0+\delta_\sigma-\mathfrak T_k$ their signs
reverse because the coefficient of $W_k$ is $-1$.

\paragraph{The constant-scale term obtained during the independent check.}
The following refinement retains the same source products and all their
endpoints. Integrating $1/(1+u)=1-u+u^2-u^3/(1+u)$ gives
\[
 \log(1+u)=u-u^2/2+u^3/3-r_4(u),\qquad
 r_4(u)=\int_0^u\frac{t^3}{1+t}\,dt,\quad0\le r_4(u)\le u^4/4.
\]
Write
\[
 A_3=8l^4+2l^2,\qquad
 A_4=\sum_{j<2l}(j^4+(j+1)^4)=\frac{192l^5+80l^3-2l}{15}.
 \tag{WGP12}
\]
To verify the last polynomial, set $N=2l$ and subtract its value
$2N^5/5+2N^3/3-N/15$ at $N$ from its value at $N+1$.
Expansion gives exactly $N^4+(N+1)^4$; at $N=0$ both sides vanish.
This proves the finite sum by induction. WGP4 at $x^{-3}$ and
the convex midpoint inequality at $x^{-4}$ give
\[
 \sum_bx_b^{-3}=3q/8-\varepsilon_3,\quad
 0\le\varepsilon_3\le(2q)^{-1},\qquad
 \sum_bx_b^{-4}\le7q/24.
\]
The derivatives and integrals used here are $(x^{-3})''=12x^{-5}$,
$\int_1^2x^{-3}dx=3/8$, and $\int_1^2x^{-4}dx=7/24$.
Thus the nonnegative quartic remainder $Q_4$ obeys
$0\le Q_4\le7A_4/(96q^3)$, and the exact refinement is
\[
 \boxed{\begin{aligned}
 P_k={}&-4lq\log q-4lq(2\log2-1)-4l^2\log2
       +\frac{8l^3+l}{6q}-\frac{l^4+l^2/4}{q^2}
       +\mathcal E^{(3)}_{q,l},\\
 -\frac l{6q}-\frac{8l^3+l}{12q^3}
 \le{}&\mathcal E^{(3)}_{q,l}\le
 \frac{l^2}{3q^2}+\frac{A_3}{6q^4}+\frac{7A_4}{96q^3}.
 \end{aligned}}\tag{WGP13}
\]
Indeed its exact error is
\[
 \mathcal E^{(3)}_{q,l}=-4l\varepsilon_0
 +\frac{4l^2}{q}\varepsilon_1
 -\frac{8l^3+l}{3q^2}\varepsilon_2
 +\frac{A_3}{3q^3}\varepsilon_3+Q_4.
 \tag{WGP14}
\]
The cubic term has coefficient $-A_3(3q/8)/(3q^3)
=-(l^4+l^2/4)/q^2$, retaining its negative sign.
The displayed bounds give $\mathcal E^{(3)}_{q,l}=O(l^{-1})$
for $m=1$ and $O(l^{-2})$ for every fixed $m\ge2$.
For $m=1$ the earlier exact remainder therefore satisfies
$\mathcal E_{q,l}=-1/256+O(l^{-1})$, because $q=(4l+2)^2$.
WGP13 supersedes the coarser numerical expansion WGP10 when this
constant-scale information is used; WGP10 remains a valid independently
proved enclosure. Both versions enter the identical WGP11 and LET21
maps, with their already specified opposite arithmetic sign.

% END RX INLINE 32


% BEGIN RX INLINE 33; SHA256 580f282cb2be5225383da4913dec33810d5482a5d02097ed594173f0bc70b87f
% New complete growing-degree bound. No historical source is modified.
\section{A positive growing-degree bound for the complete Gamma return}
\label{sec:RWB}

Retain exactly the measure, monic polynomials, norms and Christoffel
coefficients of HCT1--24:
\[
 d\sigma(y)=\frac{|\Gamma(1/4+iy/2)|^2}{2\pi}\,dy,
 \quad \nu=(y\mapsto y^2)_*\sigma,
 \quad d\nu_a(t)=t^a\,d\nu(t),
 \quad h_j^{(a)}=\int_0^\infty(P_j^{(a)}(t))^2d\nu_a(t),
 \quad c_j^{(a)}=\frac{h_j^{(a+1)}}{h_j^{(a)}}.
 \tag{RWB1}
\]
Every $P_j^{(a)}$ is monic of degree $j$. The original mass
$\sqrt{2\pi}$ remains in the norms; the last quotient is the exact HCT13
identity. Throughout this section $a,j$ are nonnegative integers, and
the estimates containing $4a-1$ require $a\ge1$.

\paragraph{The logarithmic derivative of the original Gamma density.}
Write $\sigma(y)$ also for the displayed density and set
$L(y)=-\sigma'(y)/\sigma(y)$. For every $y\in\mathbb R$,
\[
 -\frac12\le yL(y)-\frac\pi2|y|\le\frac32,
 \qquad
 \left|y\left(L(y)-\frac\pi2\operatorname{sgn}y\right)\right|
 \le\frac32.
 \tag{RWB2}
\]
The expressions at $y=0$ have their continuous value zero. Here is a
derivation retaining the exact parameter $1/4$. For $\Re z>0$,
integration by parts in the beta integral and the change $u=Nx$ give
\[
 \frac{N!N^z}{z(z+1)\cdots(z+N)}
 =\int_0^N u^{z-1}(1-u/N)^N\,du\longrightarrow\Gamma(z).
 \tag{RWB3}
\]
The inequality $(1-u/N)^N\le e^{-u}$ proves locally uniform convergence
on every compact subset of $\Re z>0$: near zero use the least real part
on that compact set, and above one use its greatest real part. The same
majorants with a factor $|\log u|$ justify differentiated convergence.
Nonvanishing at each $z=\alpha+iv$, $\alpha>0$, follows directly from
the same limits, since
\[
 \frac{|\Gamma(\alpha+iv)|}{\Gamma(\alpha)}
 =\lim_{N\to\infty}\prod_{r=0}^N
       \left(1+\frac{v^2}{(r+\alpha)^2}\right)^{-1/2}>0.
\]
The infinite product has a positive limit because
$\sum_r\log(1+v^2/(r+\alpha)^2)$ converges by comparison with
$\sum_r v^2/(r+\alpha)^2$. Thus logarithmic differentiation, with
$z=1/4+iv$, gives the convergent imaginary-part series
\[
 \operatorname{Im}\frac{\Gamma'(1/4+iv)}{\Gamma(1/4+iv)}
 =\sum_{r=0}^{\infty}\frac{v}{(r+1/4)^2+v^2},
 \qquad L(y)=\operatorname{Im}\frac{\Gamma'(1/4+iy/2)}{
                                  \Gamma(1/4+iy/2)}.
 \tag{RWB4}
\]
Only an imaginary part is taken, so the real $\log N$ term disappears.
For $v>0$ let $F_v(x)=v^2/(x^2+v^2)$, a positive decreasing function
on $[0,\infty)$. Integral comparison on each unit interval proves
\[
 \int_{1/4}^{\infty}F_v(x)\,dx
 \le\sum_{r\ge0}F_v(r+1/4)
 \le F_v(1/4)+\int_{1/4}^{\infty}F_v(x)\,dx.
\]
Since $\int_0^\infty F_v(x)\,dx=\pi v/2$, the difference between
the sum and this last integral is at least $-1/4$ and at most
\[
 F_v(1/4)-\int_0^{1/4}F_v(x)\,dx
 \le\frac34F_v(1/4)\le\frac34.
\]
For $y=2v>0$, RWB4 gives $yL(y)=2\sum_r F_v(r+1/4)$.
Multiplying the preceding interval by two proves RWB2 at positive $y$.
Evenness of $\sigma$ makes $yL(y)$ even, proving the assertion on the
whole original real line.

All integrations by parts below are justified by the exponential bound
proved from the Euler contour in HCT1. Formula RWB4 shows that $L$ is
continuous at zero; RWB2 bounds it at infinity, so the differentiated
integrands have the required integrability as well.

\paragraph{An inverse-square estimate with its exact polynomial space.}
For $a\ge1$ fix $Q=P_j^{(a+1)}$ and put
\[
 J_{a,j}=
 \frac{\int_0^\infty Q(t)^2\,d\nu_a(t)}
      {\int_0^\infty Q(t)^2\,d\nu_{a+1}(t)}.
\]
Then
\[
 0<J_{a,j}\le\left(\frac{\pi}{4a-1}\right)^2.
 \tag{RWB5}
\]
To prove this, the exact density of $\nu_{a+1}$ from HCT1 is
$t^{a+1/2}\sigma(\sqrt t)$ on $(0,\infty)$. Integrate the derivative
of $Q(t)^2t^{a+1/2}\sigma(\sqrt t)$ over that interval. Its boundary
values vanish: at zero $a+1/2>0$, and at infinity the factor has a
polynomial times exponential bound. The term $2QQ'$ integrates to
zero because $\deg Q'<j$ and $Q$ is orthogonal in the exact measure
$\nu_{a+1}$. Dividing by $h_j^{(a+1)}$ therefore gives
\[
 (a+1/2)J_{a,j}
 =\frac12\,\frac{\int_0^\infty
                 Q(t)^2 L(\sqrt t)/\sqrt t\,d\nu_{a+1}(t)}
                    {h_j^{(a+1)}}.
 \tag{RWB6}
\]
RWB2 implies $L(\sqrt t)/\sqrt t\le
\pi/(2\sqrt t)+3/(2t)$. Cauchy--Schwarz in the same positive measure
gives an upper bound $\sqrt{J_{a,j}}$ for the ratio obtained by
inserting $t^{-1/2}$. Consequently
\[
 (a-1/4)J_{a,j}\le\frac\pi4\sqrt{J_{a,j}}.
\]
The denominator is positive and $J_{a,j}>0$, so division proves RWB5.
No assertion concerning an unspecified inverse operator is used: both
integrals defining $J_{a,j}$ are the original polynomial integrals.

\paragraph{A lower bound including both the degree and the shift.}
For every $a\ge1,j\ge0$,
\[
 \boxed{\displaystyle
 \sqrt{c_j^{(a)}}\ \ge\
 \frac{2(4a-1)(2a+2j+1)}{\pi(4a+2)}
 =\frac4\pi\,(a+j+1/2)\frac{4a-1}{4a+2}.}
 \tag{RWB7}
\]
Here the square root is the positive one. Work in the original
real-line measure $dw_a(y)=y^{2a}d\sigma(y)$ and define
\[
 R(y)=P_j^{(a)}(y^2),\qquad
 O(y)=yP_j^{(a+1)}(y^2).
 \tag{RWB8}
\]
Evenness of $w_a$ and the defining orthogonality on the square
pushforward show that $R,O$ are the monic orthogonal polynomials of
degrees $2j,2j+1$ in $L^2(w_a)$. Their norms are exactly
$h_j^{(a)}$ and $h_j^{(a+1)}$; both parity sectors have been retained.
In particular
\[
 \int O'R\,dw_a=(2j+1)h_j^{(a)},\qquad
 \int OR'\,dw_a=0,\qquad
 \int\frac{OR}{y}\,dw_a=h_j^{(a)}.
 \tag{RWB9}
\]
The last integrand is a polynomial: $O/y=P_j^{(a+1)}(y^2)$ is monic
of the same degree as $R$, so their difference has lower degree.
Integrating $(ORy^{2a}\sigma)'$ and using RWB9 proves the exact identity
\[
 2a+2j+1=\frac{\int OR L\,dw_a}{h_j^{(a)}}.
 \tag{RWB10}
\]
At zero the boundary term is zero because $O$ is odd, and the apparent
$1/y$ term has just been exhibited as a polynomial. At both infinite
ends exponential decay proves vanishing.

Write $L(y)=(\pi/2)\operatorname{sgn}y+\rho(y)$.
Cauchy--Schwarz bounds the first contribution in RWB10 by
$(\pi/2)\sqrt{c_j^{(a)}}$. For the second, $|y\rho(y)|\le3/2$
from RWB2 and $O/y=P_j^{(a+1)}(y^2)$ give the bound
\[
 \frac32\sqrt{
 \frac{\int(P_j^{(a+1)}(y^2))^2\,dw_a}{h_j^{(a)}}}
 =\frac32\sqrt{c_j^{(a)}J_{a,j}}
 \le\frac{3\pi}{2(4a-1)}\sqrt{c_j^{(a)}}.
\]
Thus $2a+2j+1\le(\pi/2)(4a+2)/(4a-1)\sqrt{c_j^{(a)}}$,
which is RWB7, including its original degree and constants.

\paragraph{The low coefficient with the sign needed in the return.}
For $p\ge0$ write $I_p=\int_0^\infty y^p\sigma(y)\,dy>0$.
Integration of $(y^{p+1}\sigma(y))'$ and RWB2 give
\[
 \frac{2p-1}{\pi}\le\frac{I_{p+1}}{I_p}
 \le\frac{2p+3}{\pi}.
 \tag{RWB11}
\]
For $p<1/2$ the lower bound is simply a valid negative bound; it will
not be used there. Indeed the integrated identity is
$(p+1)I_p=\int y^{p+1}L(y)\sigma(y)\,dy$; substituting the interval
in RWB2 proves both signs in RWB11. Since $P_0^{(a)}=1$, evenness
gives $c_0^{(a)}=I_{2a+2}/I_{2a}$. Multiplying two upper bounds in
RWB11 proves
\[
 \boxed{\displaystyle c_0^{(a)}\le
 \frac{(4a+3)(4a+5)}{\pi^2}.}
 \tag{RWB12}
\]
This is the upper bound appropriate to the negative low-endpoint term;
RWB7 will be used only in the positively signed high terms.

\paragraph{The full four-endpoint functional at the original ranks.}
Retain the packet integers
\[
 k=4l+1,\quad l\ge1,\quad m\ge1,\quad
 e=1+k(m-1),\quad q=e(k+1)^2=2n.
 \tag{RWB13}
\]
Let $W_k$ be exactly LET19. With
\[
 G=\prod_{b=0}^{2l-1}
 \left\{\prod_{u=q}^{2q-1}(u+b+1/2)
         \prod_{u=q+1}^{2q}(u+b+1/2)\right\},
\]
HCT37 and the low moment ratio give the exact identity
\[
 e^{W_k}=\frac1G\prod_{s=0}^{l-1}
 \frac{\displaystyle
       \left(\prod_{j=0}^{n-1}c_j^{(q+s)}\right)^2
       c_n^{(q+s)}
       \left(\prod_{j=0}^{n-1}c_j^{(q+s+1)}\right)^2}
      {c_0^{(q+s)}}.
 \tag{RWB14}
\]
Here $P_k=-\log G$ follows by cancelling the overlapping literal
products $D_{l,N}$ in LET19, with both endpoint intervals as displayed.
The source masses cancel only in these prescribed equal-dimension
ratios. There are $4lq$ original linear factors in $G$. In each $s$
factor the net number of $c$ coefficients is $2q$. Therefore division
of each linear factor by $q$ and each $c$ by $q^2$ cancels exactly
$4lq\log q$ on the two sides and yields
\begin{align}
 W_k={}&-\sum_{b=0}^{2l-1}
 \left[\sum_{u=q}^{2q-1}\log\frac{u+b+1/2}{q}
       +\sum_{u=q+1}^{2q}\log\frac{u+b+1/2}{q}\right]\notag\\
 &+\sum_{s=0}^{l-1}\left[
 2\sum_{j=0}^{n-1}\log\frac{c_j^{(q+s)}}{q^2}
 +\log\frac{c_n^{(q+s)}}{q^2}
 +2\sum_{j=0}^{n-1}\log\frac{c_j^{(q+s+1)}}{q^2}
 -\log\frac{c_0^{(q+s)}}{q^2}\right].
 \tag{RWB15}
\end{align}
This is an exact equality of the original quantities, with every scale
factor accounted for, rather than replacement of the original objects.

Define $r_q=(4q-1)/(4q+2)$. The function
$a\mapsto(4a-1)/(4a+2)$ is increasing for $a\ge1$. RWB7 thus implies,
for both $a=q+s$ and $a=q+s+1$,
\[
 \log\frac{c_j^{(a)}}{q^2}
 \ge2\log\frac4\pi+2\log(1+j/q)+2\log r_q.
 \tag{RWB16}
\]
RWB12 and $s\le l-1$ give the simultaneous upper bound
\[
 \frac{c_0^{(q+s)}}{q^2}\le
 U_{q,l}:=\frac{(4q+4l-1)(4q+4l+1)}{\pi^2q^2}
 \le\left(\frac4\pi\right)^2
          \left(1+\frac{l+1/4}{q}\right)^2.
 \tag{RWB17}
\]
For the increasing function $x\mapsto\log(1+x)$, comparison of its
left and right sums on $[0,1/2]$ gives
\[
 \sum_{j=0}^{n-1}\log(1+j/q)
 \ge q\int_0^{1/2}\log(1+x)\,dx-\log(3/2)
 =q\left[\frac32\log(3/2)-\frac12\right]-\log(3/2).
 \tag{RWB18}
\]
Inserting RWB16--18 into the $s$ bracket of RWB15, counting its
$2q+1$ positive coefficients, gives the lower bound
\[
 4q\left[\log(4/\pi)+3\log(3/2)-1\right]
 +2\log(4/\pi)+(4q+2)\log r_q
 -6\log(3/2)-\log U_{q,l}.
 \tag{RWB19}
\]

For the first line of RWB15 use
$\log((u+b+1/2)/q)\le\log(u/q)+(b+1/2)/u
\le\log(u/q)+(b+1/2)/q$. The sum of the latter corrections over
both $u$ intervals and all $b$ is exactly $4l^2$.
Concavity of $\log x$ shows that its composite trapezoidal sum on
$[1,2]$ is at most its integral: each chord lies below the graph.
Consequently
\[
 \sum_{u=q}^{2q-1}\log(u/q)+\sum_{u=q+1}^{2q}\log(u/q)
 \le2q\int_1^2\log x\,dx=2q(2\log2-1).
\]
The first line of RWB15 is therefore at least
$-4lq(2\log2-1)-4l^2$. Add $l$ copies of RWB19 and use
\[
 (4q+2)\log r_q\ge-\frac{3(4q+2)}{4q-1}
 =-3-\frac9{4q-1},\qquad
 \log(1+x)\le x\quad(x\ge0).
\]
The first inequality follows by integrating
$1/(1-t)\le1/(1-x)$ on $0\le t\le x<1$.
This proves the completely explicit bound
\[
 \boxed{\displaystyle
 W_k\ge4lq\log\frac{27}{8\pi}
       -4l^2-6l\log(3/2)-3l
       -\frac{9l}{4q-1}-\frac{2l(l+1/4)}q.}
 \tag{RWB20}
\]
Every term is at the original finite $q,l$, including the negatively
signed low coefficient. In particular no large-rank value of a Hankel
determinant or a recurrence coefficient has been assumed.

\paragraph{Uniform positivity and the actual packet limit.}
The preceding finite bound implies
\[
 \boxed{\displaystyle W_k\ge\frac{lq}{64}>0
 \quad\text{for every }l\ge1,\ m\ge1.}
 \tag{RWB21}
\]
Here are rational verifications of the constants. The identity
\[
 \frac{22}{7}-\pi=\int_0^1
                \frac{x^4(1-x)^4}{1+x^2}\,dx>0
\]
follows by polynomial division and $\int_0^1(1+x^2)^{-1}dx=\pi/4$.
Thus $27/(8\pi)>189/176$. For $x\ge0$, differentiation proves
$\log(1+x)\ge2x/(2+x)$, and hence
$4\log(27/(8\pi))>104/365$. Integrating
$1/(1+t)\le1-t+t^2$ over $0\le t\le1/2$ gives
$\log(3/2)\le5/12$.
RWB13 gives $q\ge(4l+2)^2\ge36l$; the second inequality is
$(4l+2)^2-36l=4(4l-1)(l-1)\ge0$. Dividing RWB20 by $lq$ now yields
\[
 \frac{W_k}{lq}\ge
 \frac{104}{365}-\frac19-\frac{11}{72}-\frac1{572}
                    -\frac5{2592}
 =\frac{2349349}{135289440}
 =\frac1{64}+\frac{470903}{270578880}>\frac1{64}.
 \tag{RWB22}
\]
The four subtracted rational terms bound respectively
$4l/q$, $(6\log(3/2)+3)/q$, $9/[q(4q-1)]$, and
$2(l+1/4)/q^2$. This proves RWB21 on the entire stated packet domain.

Finally, along the original $k=4l+1\to\infty$ and any integer choices
$m=m(l)\ge1$, RWB13 gives $l/q\to0$ and $1/q\to0$.
The exact finite remainder in RWB20 therefore proves
\[
 \boxed{\displaystyle
 \liminf_{l\to\infty}\frac{W_{4l+1}}{lq}
 \ge4\log\frac{27}{8\pi}>0.}
 \tag{RWB23}
\]
This is a lower bound, with its finite remainder RWB20. No equality
or numerical asymptotic value is asserted. The exact original receiving
identities LET20--21 may now use this bound with their already proved
signs and errors; the arithmetic contribution in those identities is
left exactly as defined there.

\paragraph{A matching order bound from the original polynomial operator.}
For completeness, the exact original real-line recurrence HCT6 gives,
for $a\ge1,j\ge0$ and $D=a+2j$,
\[
 c_j^{(a)}\le
 \left[\sqrt{(D+1)(D+1/2)}+\sqrt{D(D-1/2)}\right]^2
 \le(2a+4j+1/2)^2.
 \tag{RWB24}
\]
Indeed the norm of a monic orthogonal polynomial is the minimum norm
among all monic polynomials of that degree: subtracting the orthogonal
polynomial leaves a lower-degree polynomial perpendicular to it.
Consequently
\[
 h_j^{(a+1)}\le\int_0^\infty
                  t(P_j^{(a)}(t))^2\,d\nu_a(t)
 =\|yF\|_{L^2(\sigma)}^2,\qquad
 F(y)=y^aP_j^{(a)}(y^2),\quad
 \|F\|_{L^2(\sigma)}^2=h_j^{(a)}.
\]
This is the explicit multiplication map on the original polynomials of
degree at most $D$, with image degree at most $D+1$.
For the exact orthonormal basis $e_d=p_d/\sqrt{\gamma_0d!(1/2)_d}$
from HCT5, its raising and lowering parts have coefficients
$\sqrt{(d+1)(d+1/2)}$ and $\sqrt{d(d-1/2)}$, respectively.
Their norms on degrees at most $D$ are at most the two terms in
RWB24; this follows by summing their mutually orthogonal images
separately. The triangle inequality proves the first estimate.
Since $D\ge1$, the arithmetic--geometric mean bounds on the two
square roots are $D+3/4$ and $D-1/4$, proving the second.

Combining the two lower bounds in RWB11, at $p=2a,2a+1$, also proves
\[
 c_0^{(a)}\ge\frac{(4a-1)(4a+1)}{\pi^2}\quad(a\ge1).
 \tag{RWB25}
\]
At the indices in RWB15, $a\le q+l$ and $j\le q/2$, so RWB24 gives
$c_j^{(a)}/q^2\le[4+(2l+1/2)/q]^2$. RWB25 gives
$c_0^{(q+s)}/q^2\ge(16-q^{-2})/\pi^2>1$, using $q\ge36$
and $\pi<22/7$. Every logarithm in the first line of RWB15 is
positive. Hence
\[
 W_k\le l(4q+2)\log\left(4+\frac{2l+1/2}{q}\right)
 \le6lq.
 \tag{RWB26}
\]
For the last numerical bound, $q\ge36l$ gives
$4+(2l+1/2)/q\le293/72$.
The positive exponential series gives
\[
 e^{10/7}\ge\sum_{r=0}^4\frac{(10/7)^r}{r!}
 =\frac{29593}{7203}>\frac{293}{72},
\]
where the rational inequality follows by cross multiplication.
Thus the logarithm is less than $10/7$ and
$(4q+2)10/7\le6q$ for $q\ge10$.
Together RWB21 and RWB26 prove the finite two-sided order estimate
\[
 \boxed{\displaystyle \frac{lq}{64}\le W_k\le6lq
 \qquad(l\ge1,\ m\ge1).}
 \tag{RWB27}
\]
The narrower lower remainder RWB20, rather than the coarse lower
constant in RWB27, remains available for propagation into each
original receiving formula.

% END RX INLINE 33

\paragraph{Exact dictionary for the equilibrium provider.}
The elliptic modulus $k$ used locally in EIQ is the $\kappa$ of GEL,
and EIQ's local $m=(u^2+v^2)/2$ is an endpoint midpoint. They are
distinctly indexed coordinates from the original packet
$k_{\rm packet}=4l+1$ and multiplicity $m_{\rm packet}$.
GEL3 gives their exact parameter map
$\alpha=a/n$, $\beta=\pi q/(2n)$, tending to $(2,\pi)$ on
the original sequences. GEL's constant $C_\Gamma=\Gamma(1/4)^2/(2\pi)$
in its source bound retains that definition; the leading coefficient
is the separately defined $\mathcal C_\Gamma$ in GEL18.

% BEGIN RX INLINE 34; SHA256 a0a0346a19d0e035a186590dbd69881cdf7a30e8fef602d47bba0932f175e731
% Independent proof. No external one-cut or equilibrium theorem is used.
% May be included as a section in the cumulative paper.
\section{Exact equilibrium for the square-root logarithmic potential}
\label{sec:independent-sqrt-log-equilibrium}

Throughout this section $\alpha>0$ and $\beta>0$.  On $(0,\infty)$ set
\begin{equation}\tag{EIQ1}
 V_{\alpha,\beta}(x)=\beta\sqrt{x}-\alpha\log x,
 \qquad
 \mathcal E_{\alpha,\beta}(\mu)
 =\int V_{\alpha,\beta}(x)\,d\mu(x)
       -\iint\log|x-y|\,d\mu(x)d\mu(y).
\end{equation}
The measures considered initially are probability measures for which
$\int V_{\alpha,\beta}\,d\mu$ is finite.  Because this potential is bounded
below and tends to $+\infty$ at both ends, this implies
$\int(\sqrt{x}+|\log x|)\,d\mu<\infty$.
The positive part of $\log|x-y|$ is then integrable, by
$\log^+|x-y|\leq\log(1+x)+\log(1+y)$.
Thus the energy is well defined with value in $\mathbb R\cup\{+\infty\}$.
All comparisons below also apply to the usual extended energy defined by
integrating the kernel
$\tfrac12V_{\alpha,\beta}(x)+\tfrac12V_{\alpha,\beta}(y)-\log|x-y|$:
the kernel is bounded below, and finite energy for this kernel implies the
same moment conditions.  To see the latter assertion without cancellation,
put $g(x)=\tfrac12V_{\alpha,\beta}(x)-\log(1+x)$.
The kernel is at least $g(x)+g(y)$; $g$ is bounded below and controls positive
constant multiples of both $\sqrt{x}$ near infinity and $|\log x|$ near zero.

\subsection{Existence and uniqueness of the two endpoints}
For $0<k<1$ define the complete elliptic integrals by their actual real
integrals
\begin{equation}\tag{EIQ2}
 K(k)=\int_0^{\pi/2}\frac{d\theta}
            {\sqrt{1-k^2\sin^2\theta}},\qquad
 E(k)=\int_0^{\pi/2}\sqrt{1-k^2\sin^2\theta}\,d\theta.
\end{equation}
There is exactly one $k\in(0,1)$ satisfying
\begin{equation}\tag{EIQ3}
 \frac{\sqrt{1-k^2}K(k)}{E(k)}=\frac{\alpha}{\alpha+2}.
\end{equation}
For this $k$ put
\begin{equation}\tag{EIQ4}
 r=\sqrt{1-k^2},\qquad
 v=\frac{(\alpha+2)\pi}{\beta E(k)},\qquad u=rv,
 \qquad a=u^2,\quad b=v^2.
\end{equation}
In particular $0<a<b<\infty$, and the endpoint equations, with all
constants retained, are
\begin{equation}\tag{EIQ5}
 uK(k)=\frac{\alpha\pi}{\beta},\qquad
 vE(k)=\frac{(\alpha+2)\pi}{\beta},\qquad
 k^2=1-\frac{u^2}{v^2}.
\end{equation}

Here is a proof of the assertion of uniqueness. Differentiation of the
integrals in (EIQ2), and integration of the derivative of
$\sin\theta\cos\theta/\sqrt{1-k^2\sin^2\theta}$ between its two zero
endpoint values, give
\[
 E'(k)=\frac{E(k)-K(k)}{k},\qquad
 K'(k)=\frac{E(k)}{k(1-k^2)}-\frac{K(k)}{k}.
\]
For completeness the second identity is equivalent, after differentiation
under the integral sign, to
\[
 \int_0^{\pi/2}
 \left\{\frac{k^2\sin^2\theta}{(1-k^2\sin^2\theta)^{3/2}}
 -\frac{1-k^2\sin^2\theta}{(1-k^2)\sqrt{1-k^2\sin^2\theta}}
 +\frac1{\sqrt{1-k^2\sin^2\theta}}\right\}d\theta=0;
\]
the expression in braces is
$-k^2(1-k^2)^{-1}\,d(\sin\theta\cos\theta/
\sqrt{1-k^2\sin^2\theta})/d\theta$.
With $t=E(k)/K(k)$, the integral definitions show strictly
$r^2<t<1$.  Consequently, if $f(k)=rK(k)/E(k)$, then
\begin{equation}\tag{EIQ6}
 \frac{kf'(k)}{f(k)}
 =\frac{t}{r^2}+\frac1t-1-\frac1{r^2}
 =\frac{(t-1)(t-r^2)}{r^2t}<0.
\end{equation}
At $k=0$ the continuous limiting value is $f(0)=1$.
As $k\uparrow1$, the integrand of $rK(k)$ tends to zero at every
$\theta<\pi/2$ and is bounded by $1$, while $E(k)\to1$ by dominated
convergence.  Hence $f(k)\to0$.  This proves (EIQ3), including existence.
It also proves smooth dependence of $k,u,v,a,b$ on $(\alpha,\beta)$ in
$(0,\infty)^2$, by the nonzero derivative in (EIQ6).

Let $m=(a+b)/2$, $R_s=\sqrt{(a+s)(b+s)}$ for $s\geq0$, and
$R_0=uv$.  Substitution
$t=a\cos^2\theta+b\sin^2\theta$ gives, directly from (EIQ5),
\begin{equation}\tag{EIQ7}
 \int_a^b\frac{V'_{\alpha,\beta}(t)}
                  {\sqrt{(b-t)(t-a)}}\,dt=0,
 \qquad
 \int_a^b\frac{tV'_{\alpha,\beta}(t)}
                  {\sqrt{(b-t)(t-a)}}\,dt=2\pi.
\end{equation}
Indeed the three integrals with numerators $t^{-1/2},t^{-1},t^{1/2}$
are respectively $2K(k)/v$, $\pi/(uv)$, and $2vE(k)$, and
$V'(t)=\beta/(2\sqrt t)-\alpha/t$.

\subsection{A positive density and its exact Cauchy transform}
The elementary Stieltjes integral, obtained by $s=xz^2$, is
\begin{equation}\tag{EIQ8}
 \int_0^\infty\frac{ds}{\sqrt s(x+s)}=\frac\pi{\sqrt x}
 \quad(x>0).
\end{equation}
Using this identity in (EIQ7) and applying Tonelli to its positive
constituents gives
\begin{equation}\tag{EIQ9}
 \frac\alpha{R_0}=
 \frac\beta{2\pi}\int_0^\infty\frac{ds}{\sqrt s R_s},
 \qquad
 \int_0^\infty\frac{1-s/R_s}{\sqrt s}\,ds=2vE(k).
\end{equation}
Both integrals converge; their integrands are $O(s^{-1/2})$ at zero
and $O(s^{-3/2})$ at infinity.  The second formula also follows by
averaging $\sqrt t=\pi^{-1}\int_0^\infty
s^{-1/2}t/(t+s)\,ds$ against the arcsine probability measure on $[a,b]$.

Define for every $x>0$
\begin{align}\tag{EIQ10}
 h(x)
 &=\frac\alpha{xR_0}
   -\frac\beta{2\pi}
       \int_0^\infty\frac{ds}{\sqrt s(x+s)R_s} \\
 &=\frac\beta{2\pi x}
       \int_0^\infty\frac{\sqrt s}{(x+s)R_s}\,ds>0.\notag
\end{align}
The second equality is exactly (EIQ9), with no change of the potential.
All integrals in (EIQ10) converge absolutely and locally uniformly in
$x>0$.  The candidate density on the original positive half-line is
\begin{equation}\tag{EIQ11}
 \rho_{\alpha,\beta}(x)=
 \begin{cases}
 \displaystyle
 \frac{\sqrt{(b-x)(x-a)}}{2\pi}\,h(x),&a<x<b,\\
 0,&x\notin(a,b).
 \end{cases}
\end{equation}
It is strictly positive on $(a,b)$ and vanishes continuously with the
square-root factor at each endpoint.  Positivity here follows from the
last integral in (EIQ10), not from a one-interval ansatz.

The following elementary transform computes its mass and resolvent.
Choose $R(z)=\sqrt{(z-a)(z-b)}$ on $\mathbb C\setminus[a,b]$ with
$R(z)/z\to1$ at infinity. Thus $R(x)>0$ for $x>b$,
$R(x)<0$ for $x<a$, and $R(x+i0)=i\sqrt{(b-x)(x-a)}$ on $(a,b)$.
Direct substitution $t=m+(b-a)\cos\theta/2$, followed by
$w=\tan(\theta/2)$, proves
\[
 \frac1\pi\int_a^b\frac{dt}{(z-t)\sqrt{(b-t)(t-a)}}=\frac1{R(z)}.
\]
For real $z>b$, the substitution reduces the integral to
$\pi^{-1}\int_0^\infty2\,dw/[(z-b)+(z-a)w^2]$;
the identity elsewhere follows by analytic continuation.
Polynomial division then gives
\begin{equation}\tag{EIQ12}
 \frac1\pi\int_a^b\frac{\sqrt{(b-t)(t-a)}}{z-t}\,dt
     =z-m-R(z),\qquad
 \frac1\pi\int_a^b\frac{\sqrt{(b-t)(t-a)}}{t+s}\,dt
     =m+s-R_s.
\end{equation}
For example the first identity follows from
$(b-t)(t-a)=-(t-z)^2+2(m-z)(t-z)-(z-a)(z-b)$,
the preceding arcsine transform, and the arcsine first moment $m$.
Partial fractions in these two identities yield
\begin{equation}\tag{EIQ13}
 H_s(z):=\frac1\pi\int_a^b
       \frac{\sqrt{(b-t)(t-a)}}{(t+s)(z-t)}\,dt
       =1-\frac{R_s+R(z)}{z+s}.
\end{equation}
The formula at $z=-s$ is interpreted by its removable limit.

The mass of (EIQ11), computed using (EIQ12), is
\begin{align}\tag{EIQ14}
 \int_a^b\rho_{\alpha,\beta}(t)\,dt
 &=\frac\alpha{2R_0}(m-R_0)
 -\frac\beta{4\pi}\int_0^\infty
       \frac{(m+s)/R_s-1}{\sqrt s}\,ds\\
 &=-\frac\alpha2+
   \frac\beta{4\pi}\int_0^\infty\frac{1-s/R_s}{\sqrt s}\,ds
 =-\frac\alpha2+\frac{\beta vE(k)}{2\pi}=1.\notag
\end{align}
In particular this is a probability density.

Define its Cauchy transform with the explicit sign convention
$G(z)=\int_a^b\rho_{\alpha,\beta}(t)/(z-t)\,dt$.
For $z$ outside both $[a,b]$ and $(-\infty,0]$, formulas (EIQ10) and
(EIQ13) give
\begin{align}\tag{EIQ15}
 G(z)
 &=\frac\alpha{2R_0}H_0(z)
       -\frac\beta{4\pi}
          \int_0^\infty\frac{H_s(z)}{\sqrt sR_s}\,ds\\
 &=-\frac\alpha{2z}+\frac\beta{4\sqrt z}
       -\frac{R(z)h(z)}2
 =\frac{V'_{\alpha,\beta}(z)-R(z)h(z)}2.\notag
\end{align}
Here the positive square root is continued with its negative-real-axis
cut, and $h(z)$ is the first integral formula of (EIQ10).  Every integral
exchange is absolutely convergent, uniformly on compact subsets of this
domain: near zero the auxiliary integrand is $O(s^{-1/2})$, and near
infinity the original integral defining $H_s$ is $O(s^{-1})$.
The constant terms on the second line cancel by (EIQ9), and (EIQ8)
supplies the term $\beta/(4\sqrt z)$.

\subsection{The Euler equality, its strict exterior inequality, and uniqueness}
Set
\begin{equation}\tag{EIQ16}
 U(x)=V_{\alpha,\beta}(x)
             -2\int_a^b\log|x-t|\rho_{\alpha,\beta}(t)\,dt
 \quad(x>0).
\end{equation}
This function is continuous on $(0,\infty)$.  The integral is continuous
because the density is bounded, compactly supported, and the integral
of $|\log|x-t||$ over an interval of length $\varepsilon$ tends to zero
uniformly in its centre.  In the interior of $[a,b]$, differentiation of
the logarithmic potential is the principal-value integral
$\operatorname{PV}\int\rho(t)/(x-t)\,dt$.
This follows by subtracting $\rho(x)$ on a symmetric interval about
$x$; the remaining numerator is $O(|t-x|)$ because the density is
continuously differentiable locally in the interior.  The real part of
(EIQ15) on the upper edge of the interval is precisely this principal
value.  Indeed subtracting $\rho(x)$ in
$\int\rho(t)/(x+i\varepsilon-t)\,dt$ gives the same limit, while the
constant term has real part equal to its elementary principal value.
Since $R(x+i0)$ is purely imaginary and $h(x)$ is real, (EIQ15) proves
$U'(x)=0$ on $(a,b)$.  There is therefore a real constant $\ell$ such that
\begin{equation}\tag{EIQ17}
 U(x)=\ell\quad(a\leq x\leq b).
\end{equation}
For positive $x$ outside this interval, (EIQ15) instead gives the exact
real derivative
\begin{equation}\tag{EIQ18}
 U'(x)=R(x)h(x)
 \begin{cases}<0,&0<x<a,\\>0,&x>b.\end{cases}
\end{equation}
Combining (EIQ17), continuity, and these signs proves
\begin{equation}\tag{EIQ19}
 U(x)>\ell\quad\hbox{for }x\in(0,a)\cup(b,\infty).
\end{equation}
This verifies the entire exterior inequality on the original domain,
including the interval next to zero.

Here is a complete positivity argument for the energy of the difference
of two measures.  If $\nu$ is a finite signed real measure of total mass
zero for which $|\log|x-y||$ is integrable against
$|\nu|\otimes|\nu|$, the scalar identity
\begin{equation}\tag{EIQ20}
 -\log d=\frac12\int_0^\infty
           \frac{e^{-td^2}-e^{-t}}t\,dt\quad(d>0)
\end{equation}
follows by differentiating in $d$ and evaluating at $d=1$.
One half of the integral over any finite subinterval of $(0,\infty)$
has the same sign as $-\log d$ and absolute value at most $|\log d|$.
Consequently dominated convergence and the zero total mass give
\begin{equation}\tag{EIQ21}
 -\iint\log|x-y|\,d\nu(x)d\nu(y)
 =\frac12\int_0^\infty\frac1t
       \iint e^{-t(x-y)^2}\,d\nu(x)d\nu(y)\,dt\geq0.
\end{equation}
The nonnegative sign is proved, rather than assumed, by the Gaussian
Fourier identity
\begin{equation}\tag{EIQ22}
 \iint e^{-t(x-y)^2}\,d\nu(x)d\nu(y)
 =\frac1{\sqrt{4\pi t}}\int_{\mathbb R}
          e^{-\xi^2/(4t)}|\widehat\nu(\xi)|^2\,d\xi,
 \qquad \widehat\nu(\xi)=\int e^{i\xi x}\,d\nu(x).
\end{equation}
The scalar Gaussian Fourier identity follows by completing the square
in the Gaussian integral (or by differentiating its Fourier transform
and its value at zero); its integration against the finite measures is
justified by absolute integrability.  Equality in (EIQ21) forces
$\widehat\nu=0$: the right side of (EIQ22) is strictly positive for every
$t>0$ if this continuous Fourier transform has a nonzero value.
Finally $\widehat\nu=0$ implies $\nu=0$ without an extra uniqueness
assumption. Convolve $\nu$ with a centred Gaussian approximate identity.
Its value at every point is zero by that same Fourier integral.
For every compactly supported continuous test function, its convolution
with these Gaussians converges uniformly to the test function; integration
against $\nu$ then proves that $\nu$ annihilates every such function.

Let $\mu$ be any probability measure of finite energy in (EIQ1), and
write $\mu_*(dx)=\rho_{\alpha,\beta}(x)\,dx$ and $\nu=\mu-\mu_*$.
The integrability needed for (EIQ21) holds.  The $\mu\otimes\mu$ term
has it by finite energy and the already established logarithmic moment;
the $\mu_*\otimes\mu_*$ term has it by bounded compact support; and the
cross term has it because the negative logarithmic singularity integrated
against the bounded density is uniformly bounded, while the positive
part is controlled by $\log(1+x)+\log(1+y)$.
Expanding the energy therefore yields exactly
\begin{equation}\tag{EIQ23}
 \mathcal E_{\alpha,\beta}(\mu)-
       \mathcal E_{\alpha,\beta}(\mu_*)
 =\int(U(x)-\ell)\,d\mu(x)
       -\iint\log|x-y|\,d\nu(x)d\nu(y)\geq0.
\end{equation}
The first term is nonnegative by (EIQ17)--(EIQ19), and the second is
strictly positive unless $\mu=\mu_*$. Infinite energy competitors cannot
improve this finite value. Thus (EIQ11) is the unique global minimizing
probability measure, with exactly the support $[a,b]$.

\subsection{An explicit endpoint integral for the logarithmic moment}
Retain $m,R_s,u,v$ as defined above, and define
\begin{equation}\tag{EIQ24}
 C=\frac{(u+v)^2}4,\qquad
 \eta=\frac{v-u}{v+u},\qquad
 \zeta_s=\frac{b-a}{2(m+s+R_s)}\quad(s\geq0).
\end{equation}
Thus $0<\zeta_s\leq\eta<1$ and $\zeta_0=\eta$.
For every $s\geq0$ put
\begin{equation}\tag{EIQ25}
 L_s=(m+s-R_s)\log C-(m-R_0)
                  -2R_s\log(1-\eta\zeta_s).
\end{equation}
Then the logarithmic moment of the equilibrium measure is exactly
\begin{equation}\tag{EIQ26}
 \mathcal L(\alpha,\beta):=\int_a^b\log x\rho_{\alpha,\beta}(x)\,dx
 =\frac\beta{4\pi}\int_0^\infty
                    \frac{L_0-L_s}{\sqrt s R_s}\,ds.
\end{equation}
This is a convergent one-dimensional endpoint integral of explicitly
given real elementary functions; its only endpoints are those specified
uniquely in (EIQ3)--(EIQ5).

Here is a direct derivation, including the constants in (EIQ25).
Set $h_0=(b-a)/2$ and substitute $t=m+h_0\cos\theta$ in arcsine
integrals.  The identities
\[
 t=C(1+2\eta\cos\theta+\eta^2),\qquad
 t+s=\frac{m+s+R_s}2
                (1+2\zeta_s\cos\theta+\zeta_s^2)
\]
give the absolutely convergent Fourier series
\[
 \log t=\log C+
       2\sum_{j\geq1}\frac{(-1)^{j+1}\eta^j}j\cos(j\theta),
 \qquad
 \frac1{t+s}=\frac1{R_s}
       \left(1+2\sum_{j\geq1}(-\zeta_s)^j\cos(j\theta)\right).
\]
Orthogonality on $[0,\pi]$ gives
\begin{align}\tag{EIQ27}
 \frac1\pi\int_a^b\frac{\log t}{\sqrt{(b-t)(t-a)}}\,dt
 &=\log C,\\
 \frac1\pi\int_a^b\frac{t\log t}{\sqrt{(b-t)(t-a)}}\,dt
 &=m\log C+(m-R_0),\notag\\
 \frac1\pi\int_a^b\frac{\log t}{(t+s)\sqrt{(b-t)(t-a)}}\,dt
 &=\frac{\log C+2\log(1-\eta\zeta_s)}{R_s}.\notag
\end{align}
The second identity uses $h_0\eta=m-R_0$.
Polynomial division
\[
 \frac{(b-t)(t-a)}{t+s}
       =-t+2m+s-\frac{R_s^2}{t+s}
\]
now proves, exactly,
\begin{equation}\tag{EIQ28}
 L_s=\frac1\pi\int_a^b
               \frac{\log t\sqrt{(b-t)(t-a)}}{t+s}\,dt.
\end{equation}
Using the positive integral expression for $h$ in (EIQ10), and the
identity $1/[t(t+s)]=(1/t-1/(t+s))/s$, yields (EIQ26).
The exchanges are absolutely convergent: $\log t$ is bounded on
$[a,b]$, $L_0-L_s=O(s)$ at zero by (EIQ28), and $L_s=O(s^{-1})$
at infinity.  Thus the integrand in (EIQ26) is $O(s^{1/2})$ at zero
and $O(s^{-3/2})$ at infinity.

\subsection{Parameter continuity and the variational derivative}
The density and its logarithmic moment depend smoothly on
$(\alpha,\beta)\in(0,\infty)^2$ in the following precise sense.
On a compact parameter neighbourhood, (EIQ3)--(EIQ6) put $a,b$ in a
fixed compact subset of $(0,\infty)$, with $b-a$ bounded below by a
positive number.  The integral for $h$ and each of its parameter and
$x$ derivatives converge uniformly for $x$ in this fixed compact set:
their absolute integrands are bounded by constants times $s^{1/2}$
near zero and $s^{-3/2}$ near infinity. Parameter derivatives of the
endpoints are bounded on a smaller compact neighbourhood.
Under the fixed-interval coordinate
$x=a+(b-a)y$, $0\leq y\leq1$, the measure is exactly
\begin{equation}\tag{EIQ29}
 \rho_{\alpha,\beta}(x)\,dx
 =\frac{(b-a)^2}{2\pi}\sqrt{y(1-y)}
                  h\bigl(a+(b-a)y\bigr)\,dy.
\end{equation}
All parameter derivatives of its coefficient are bounded by a constant
times $\sqrt{y(1-y)}$. The same is true after multiplication by
$\log(a+(b-a)y)$ or by $\sqrt{a+(b-a)y}$.
Differentiation under this integral proves smoothness, in particular
continuity, of both moments on the entire positive parameter quadrant.

Define $E_{\min}(\alpha,\beta)=\mathcal E_{\alpha,\beta}(\mu_{\alpha,\beta})$,
where the measure is the unique minimizer proved above.
Testing the minimizer for a parameter pair against the neighbouring
potential, and then reversing the two tests, gives for $h>0$
\[
 -\mathcal L(\alpha+h,\beta)
 \leq\frac{E_{\min}(\alpha+h,\beta)-E_{\min}(\alpha,\beta)}h
 \leq-\mathcal L(\alpha,\beta).
\]
The corresponding inequalities for $h<0$ and the proved continuity
therefore give
\begin{equation}\tag{EIQ30}
 \partial_\alpha E_{\min}(\alpha,\beta)=-\mathcal L(\alpha,\beta),\qquad
 \partial_\beta E_{\min}(\alpha,\beta)
        =\int\sqrt{x}\,d\mu_{\alpha,\beta}(x).
\end{equation}
No unproved differentiability of a minimizing selection is being used.
For the opposite convention
$F(\alpha,\beta)=\sup_\mu\{\iint\log|x-y|\,d\mu(x)d\mu(y)
-\int V_{\alpha,\beta}\,d\mu\}=-E_{\min}(\alpha,\beta)$,
the exact derivative is therefore
$\partial_\alpha F(\alpha,\beta)=\mathcal L(\alpha,\beta)$.

The second moment in (EIQ30) has an exact value.  Push the minimizing
measure forward by the explicit dilation $x\mapsto c^2x$, $c>0$.
Writing $M=\int\sqrt{x}\,d\mu_{\alpha,\beta}(x)$, its energy minus
the energy at $c=1$ is precisely
$\beta(c-1)M-2(\alpha+1)\log c$.
Its derivative at the minimizing value $c=1$ is zero. Hence
\begin{equation}\tag{EIQ31}
 \int\sqrt{x}\,d\mu_{\alpha,\beta}(x)
       =\frac{2(\alpha+1)}\beta.
\end{equation}
The same explicit dilation, or (EIQ3)--(EIQ5), proves that the map
$x\mapsto(\pi/\beta)^2x$ pushes $\mu_{\alpha,\pi}$ to
$\mu_{\alpha,\beta}$, including the density and its two endpoints.
Consequently
\begin{equation}\tag{EIQ32}
 \mathcal L(\alpha,\beta)
       =\mathcal L(\alpha,\pi)+2\log(\pi/\beta).
\end{equation}

For the particular original limiting potential
$V(x)=\pi\sqrt{x}-2\log x$, the substitutions are exactly
$\alpha=2$, $\beta=\pi$. Equations (EIQ5), (EIQ10), and (EIQ26) become
\begin{equation}\tag{EIQ33}
 uK(k)=2,\qquad vE(k)=4,\qquad
 h(x)=\frac1{2x}\int_0^\infty\frac{\sqrt s}{(x+s)R_s}\,ds,
 \qquad
 \mathcal L(2,\pi)=\frac14\int_0^\infty
                         \frac{L_0-L_s}{\sqrt sR_s}\,ds.
\end{equation}
In particular the preceding proof includes positivity, unit mass,
the full Euler inequality, the unique global minimizer, and the
parameter continuity needed at these actual parameters.

% END RX INLINE 34


% BEGIN RX INLINE 35; SHA256 2014e7fefdcd963b70b00fb89d597e24c0826582a2541dd1ff44d63051979fc1
% Original-source large-degree calculation, 14 September 2026.
% Input fragment. Equilibrium provider:
% independent/EXACT_SQRT_LOG_EQUILIBRIUM.tex (EIQ1--33).
\section{The exact Gamma ensemble and the leading growing return}
\label{sec:GEL}

\paragraph{GEL1. Source, exponents, and the literal coordinate map.}
The source throughout is
\[
 d\sigma(y)=\frac{|\Gamma(1/4+iy/2)|^2}{2\pi}\,dy,
 \qquad \sigma(\mathbb R)=\sqrt{2\pi},\qquad
 d\nu(t)=\frac{|\Gamma(1/4+i\sqrt t/2)|^2}{2\pi\sqrt t}\,dt.
 \tag{GEL1}
\]
Here \(\nu\) is the full pushforward under \(y\mapsto y^2\); both
half-axes are present in its density. For integers \(a,n\geq0\), retain
\[
 H_n^{(a)}=\det\left[\int_0^\infty t^{a+i+j}\,d\nu(t)\right]_{i,j<n}.
\]
The definition extends to real \(a>-1/2\); all its integrals converge.
For \(n\geq1\), expansion of the two Vandermonde determinants followed
by integration term by term gives the exact identity
\[
 H_n^{(a)}=\frac1{n!}\int_{(0,\infty)^n}
       \prod_{i<j}(t_i-t_j)^2\prod_{i=1}^n t_i^a\,d\nu(t_i).
 \tag{GEL2}
\]
Indeed expansion gives \(n!\) copies of the Gram determinant after one
of the two permutation indices is fixed. Absolute integrability follows
by expanding the squared polynomial and using the existing moments.

Fix a positive real coordinate scale \(q\), independently of \(a,n\),
and make the literal map \(t_i=q^2x_i\). Introduce the reference
probability measure
\[
 d\omega(x)=\tfrac12x^{-1/2}e^{-\sqrt x}\,dx,
 \qquad B_q(x)=2q e^{\sqrt x}\sigma(q\sqrt x).
\]
Its mass is one by \(x=u^2\). This reference measure does not replace
or rescale the source: the exact Radon--Nikodym identity is
\(d\nu(q^2x)=B_q(x)d\omega(x)\), including its whole mass. Consequently
\[
 \boxed{H_n^{(a)}=q^{2an+2n(n-1)}J_{n,q}(a),\quad
 J_{n,q}(a)=\frac1{n!}\int\prod_{i<j}(x_i-x_j)^2
                         \prod_i x_i^aB_q(x_i)\,d\omega(x_i).}
 \tag{GEL3}
\]
The powers in this identity come separately from the \(n\) factors
\(t_i^a\) and the \(n(n-1)/2\) squared differences. No determinant,
source mass, or Jacobian has been omitted.

\paragraph{GEL2. An elementary global Gamma bound with the exact exponential.}
Put \(\lambda=1/4\), \(b\geq0\), and
\(C_\Gamma=\Gamma(1/4)^2/(2\pi)\). Euler's finite limit gives
\[
 \frac{|\Gamma(\lambda+ib)|^2}{\Gamma(\lambda)^2}
 =\prod_{j=0}^\infty\left(1+\frac{b^2}{(j+\lambda)^2}\right)^{-1}.
\]
For clarity, the finite limit follows by applying the beta integral to
\(\int_0^1t^{z-1}(1-t)^Ndt=N!/[z(z+1)\cdots(z+N)]\), replacing
\(t\) by \(u/N\), and using dominated convergence for
\(\operatorname{Re}z>0\). Taking squared moduli and dividing by the
same limit at \(z=\lambda\) proves the displayed product. Its logarithm
converges absolutely because its summand is \(O((j+1)^{-2})\).

The decreasing function \(v\mapsto\log(1+b^2/v^2)\) gives, with
\(S=\sum_{j\geq0}\log(1+b^2/(j+\lambda)^2)\),
\[
 I\leq S\leq I+\log(1+b^2/\lambda^2),\qquad
 I=\int_\lambda^\infty\log(1+b^2/v^2)\,dv
  =\pi b-\lambda\log(1+b^2/\lambda^2)
                      -2b\arctan(\lambda/b).
\]
The integral is evaluated by differentiating
\(v\log(1+b^2/v^2)+2b\arctan(v/b)\); at \(b=0\) it is zero
by continuity. Since \(0\leq2b\arctan(\lambda/b)\leq2\lambda\),
substitution of \(b=|y|/2\) proves the global, finite inequalities
\[
 \boxed{C_\Gamma(1+4y^2)^{-3/4}e^{-\pi|y|/2}
 \leq\sigma(y)\leq
 C_\Gamma e^{1/2}(1+4y^2)^{1/4}e^{-\pi|y|/2}.}
 \tag{GEL4}
\]
Both inequalities include \(y=0\). In particular, for
\(C_0=\Gamma(1/4)^2/\pi\),
\[
 B_q(x)\leq C_0q e^{1/2}(1+2q)^{1/2}
            e^{-(\pi q/2-3/2)\sqrt x}\quad(x>0).
 \tag{GEL5}
\]
Here \((1+4q^2x)^{1/4}\leq(1+2q\sqrt x)^{1/2}
\leq(1+2q)^{1/2}(1+\sqrt x)^{1/2}
\leq(1+2q)^{1/2}e^{\sqrt x/2}\). On each fixed compact interval
\([A,B]\subset(0,\infty)\), the lower bound is
\[
 B_q(x)\geq C_0q(1+4q^2B)^{-3/4}e^{-\pi q\sqrt x/2}.
 \tag{GEL6}
\]
Thus the exact exponential, and the size and position of the remaining
source factor, have been proved directly from the original Gamma function.

\paragraph{GEL3. The energy object and its original parameter domain.}
For \(\alpha,\beta>0\) put
\[
 V_{\alpha,\beta}(x)=\beta\sqrt x-\alpha\log x,
 \quad K_{\alpha,\beta}(x,y)=\log|x-y|
           -\tfrac12V_{\alpha,\beta}(x)-\tfrac12V_{\alpha,\beta}(y),
\]
\[
 \mathcal F(\alpha,\beta)=\sup_{\mu\in\mathcal P((0,\infty))}
                   \iint K_{\alpha,\beta}(x,y)\,d\mu(x)d\mu(y).
 \tag{GEL7}
\]
The kernel is bounded above, so its integral has an unambiguous value
in \(\mathbb R\cup\{-\infty\}\) for every probability in this domain.
Whenever its value is finite it equals
\(\iint\log|x-y|\,d\mu^2-\int V_{\alpha,\beta}\,d\mu\).
Indeed \(K\leq r(x)+r(y)\), with
\(r=\log(1+x)-V_{\alpha,\beta}/2\) bounded above and tending to
\(-\infty\) at both endpoints; finite kernel integral therefore forces
the \(\sqrt x\) and \(|\log x|\) moments to be finite. The positive
logarithmic part is then integrable, and finite kernel integral forces
the negative part to be integrable as well. This defines the extended
energy without forming an indeterminate difference of infinite quantities.
The equilibrium calculation supplied with this section proves existence,
uniqueness, and the full exterior variational inequality for its density
\(\rho_{\alpha,\beta}\), with support \([u^2,v^2]\). Its exact endpoints
and positive density are
\[
 uK(\kappa)=\frac{\alpha\pi}{\beta},\quad
 vE(\kappa)=\frac{(\alpha+2)\pi}{\beta},\quad
 \kappa^2=1-u^2/v^2,
\]
\[
 \rho_{\alpha,\beta}(x)=
 \frac{\beta\sqrt{(v^2-x)(x-u^2)}}{4\pi^2x}
 \int_0^\infty
 \frac{\sqrt s\,ds}{(x+s)\sqrt{(u^2+s)(v^2+s)}}
       \quad(u^2<x<v^2).
 \tag{GEL8}
\]
The endpoint proof shows continuous parameter dependence, finite entropy
relative to \(\omega\), finite logarithmic energy, and
\[
 \partial_\alpha\mathcal F(\alpha,\beta)
       =\int\log x\,\rho_{\alpha,\beta}(x)\,dx.
 \tag{GEL9}
\]
In particular \(\mathcal F\) and this derivative are continuous on the
positive parameter quadrant. The following ensemble argument proves the
connection to GEL3; it does not presume a random-matrix asymptotic theorem.

\paragraph{GEL4. Complete free-energy limit from a compactified kernel.}
Let \(n\to\infty\), \(q/n\to\theta>0\), and let \(\alpha>0\) be fixed.
Then the actual integral in GEL3 satisfies
\[
 \boxed{\frac1{n^2}\log J_{n,q}(n\alpha)
       \longrightarrow\mathcal F(\alpha,\pi\theta/2).}
 \tag{GEL10}
\]
Here is a proof of both bounds. For the upper bound GEL5 removes the
factor \([C_0q e^{1/2}(1+2q)^{1/2}]^n/n!\), whose logarithm divided by
\(n^2\) tends to zero. The remaining exponent is
\[
 2\sum_{i<j}\log|x_i-x_j|-n\sum_iV_{\alpha,\beta_n}(x_i),
 \qquad \beta_n=\frac{\pi q}{2n}-\frac3{2n}.
\]
For \(n>1\) it is exactly
\(\sum_{i\ne j}K_{Q_n}(x_i,x_j)\), where
\[
 Q_n=\frac n{n-1}V_{\alpha,\beta_n},\qquad
 K_Q(x,y)=\log|x-y|-\tfrac12Q(x)-\tfrac12Q(y).
\]
Write \(\beta=\pi\theta/2\). Choose \(\varepsilon>0\) so that
\(\alpha-\varepsilon>0\) and \(\beta-3\varepsilon>0\).
For all sufficiently large \(n\), the two coefficients of \(Q_n\)
differ from \((\alpha,\beta)\) by at most \(\varepsilon\). It follows
pointwise that
\(Q_n\geq Q_-:=V_{\alpha-\varepsilon,\beta-3\varepsilon}\).
Indeed its difference has the form
\(b\sqrt x-a\log x\), where \(b\geq2\varepsilon\) and
\(0\leq a\leq2\varepsilon\); it is nonnegative for \(x<1\), and for
\(x\geq1\) use \(\log x\leq\sqrt x\).

Compactify \((0,\infty)\) by adjoining its two endpoints.
The function \(K_{Q_-}\), assigned value \(-\infty\) on the diagonal
and at either added endpoint, has continuous truncations
\(K_M=\max\{K_{Q_-},-M\}\). To verify the endpoint assertion uniformly
in the other variable, use
\[
 K_{Q_-}(x,y)\leq r(x)+r(y),\qquad
 r(x)=\log(1+x)-Q_-(x)/2.
\]
The function \(r\) is bounded above and tends to \(-\infty\) at both
endpoints. Continuity at a diagonal point in the interior follows from
\(\log|x-y|\to-\infty\) and local boundedness of \(Q_-\).
For the empirical probability \(L_n=n^{-1}\sum_i\delta_{x_i}\),
\[
 \frac1{n^2}\sum_{i\ne j}K_{Q_-}(x_i,x_j)
 \leq\iint K_M\,dL_n^2+\frac Mn
 \leq s_M+\frac Mn,
 \quad s_M=\max_{\mu}\iint K_M\,d\mu^2.
\]
Since \(\omega\) has mass one, integration preserves this upper bound.
Moreover \(s_M\downarrow\mathcal F(\alpha-\varepsilon,
\beta-3\varepsilon)\). For completeness, choose maximizing probabilities
\(\mu_M\). Compactness gives a convergent subsequence. For every fixed
\(L\), the inequality \(K_M\leq K_L\) when \(M\geq L\), followed by
continuity of the integral of \(K_L\), bounds the subsequential limit
of \(s_M\) by \(\iint K_Ld\mu^2\). Decrease \(L\)'s truncation to the
untruncated kernel. Monotone convergence after subtracting a common upper
bound gives \(\lim s_M\leq\iint K_{Q_-}d\mu^2\leq\mathcal F\).
The reverse inequality follows by testing every \(K_M\) on an arbitrary
probability of finite energy. Probabilities charging an added endpoint
have energy \(-\infty\), so they do not alter the supremum.
First take \(n\to\infty\), then \(M\to\infty\), and finally
\(\varepsilon\downarrow0\). Continuity proved in GEL9 yields the
upper bound in GEL10.

For the lower bound choose exactly
\(d\mu(x)=\rho_{\alpha,\beta}(x)dx\), supported on \([A,B]\), and use
GEL6. Write \(p=d\mu/d\omega\). Restricting the integral to this support
and changing measure to \(\mu^n\), Jensen's inequality gives
\begin{align*}
 \log J_{n,q}(n\alpha)\geq{}&-
 \log n!+n\log\{C_0q(1+4q^2B)^{-3/4}\}\\
 &+n(n-1)\iint\log|x-y|\,d\mu(x)d\mu(y)\\
 &-n^2\int V_{\alpha,\pi q/(2n)}\,d\mu
       -n\int\log p\,d\mu .
\end{align*}
All logarithms on the right are integrable: the equilibrium density is
bounded, has square-root decay at its finite positive endpoints, and
\(\omega\) has a positive smooth density there. The prefactor contributes
\(O(n\log q+n\log n)\), and the entropy contributes \(O(n)\).
Division by \(n^2\) therefore yields exactly the value of GEL7 at its
maximizer. This proves GEL10.

\paragraph{GEL5. Finite positive shifts by a proved convex squeeze.}
For each fixed \(n,q\), differentiation under the integral in GEL3 on
every compact subinterval of \(a>-1/2\) gives
\[
 \frac{d}{da}\log J_{n,q}(a)=
    \mathbb E_a\sum_i\log x_i,\qquad
 \frac{d^2}{da^2}\log J_{n,q}(a)=
    \operatorname{Var}_a\!\left(\sum_i\log x_i\right)\geq0.
 \tag{GEL11}
\]
The expectation is with respect to the entire positive integrand in
GEL3 divided by its integral. Extra logarithmic factors are integrable
at zero because \(a>-1/2\), and at infinity by GEL4. Hence the
functions \(F_{n,q}(z)=n^{-2}\log J_{n,q}(nz)\) are differentiable
and convex. Suppose now
\[
 q/n\longrightarrow2,\qquad a/n\longrightarrow2,\qquad
 l>0,\quad l/n\longrightarrow0.
\]
Set
\[
 \mathcal L=\int_{u^2}^{v^2}\log x\,\rho_{2,\pi}(x)\,dx,
 \qquad uK(\kappa)=2,\quad vE(\kappa)=4.
 \tag{GEL12}
\]
For fixed \(0<\eta<2\), both endpoints of the secant from \(a/n\) to
\((a+l)/n\) lie in \((2-\eta/2,2+\eta/2)\) for large \(n\).
Convexity bounds that secant from below by the secant from
\(2-\eta\) to \(2-\eta/2\), and from above by the secant from
\(2+\eta/2\) to \(2+\eta\). Apply GEL10 at these four fixed positive
arguments, then let \(\eta\downarrow0\) and use GEL9. This proves
\[
 \frac{\log J_{n,q}(a+l)-\log J_{n,q}(a)}{nl}
       \longrightarrow\mathcal L.
\]
Combining this with the exact source factors in GEL3 gives
\[
 \boxed{\log\frac{H_n^{(a+l)}}{H_n^{(a)}}
       =2nl\log q+nl\mathcal L+o(nl).}
 \tag{GEL13}
\]
This proof covers moving starting exponents \(a\); it never differentiates
the error in a free-energy asymptotic. The remainder assertion has the
precise scope \(o(nl)\) along the displayed sequences.

\paragraph{GEL6. Both original matrices and the low moment.}
Return to the actual packet
\[
 k=4l+1,\quad e=1+k(m-1),\quad q=e(k+1)^2=2n,
 \qquad m\geq1\text{ fixed},\qquad l\to\infty.
\]
Retain the exact parity identity from HCT35,
\(Z_a=H_n^{(a)}H_{n+1}^{(a)}(H_n^{(a+1)})^2\).
Use GEL13 at the four literal factors, with sizes \(n,n+1,n,n\)
and starting exponents \(q,q,q+1,q+1\), respectively. Every sequence
satisfies its stated domain. Addition gives
\[
 \boxed{\log\frac{Z_{q+l}}{Z_q}
   =2(2q+1)l\log q+(2q+1)l\mathcal L+o(ql).}
 \tag{GEL14}
\]
Thus the contribution from the two original full matrices has been
evaluated, with all four parity blocks still identifiable.

Write \(\mu_a=\int y^{2a}d\sigma(y)\). GEL4 gives, for integer \(a\geq1\),
\[
 \log\mu_a=2a\log a+O(a).
 \tag{GEL15}
\]
One elementary proof of the upper bound uses
\((1+4y^2)^{1/4}\leq1+2y\) for \(y\geq0\); its integrals are
\((2a)!/(\pi/2)^{2a+1}\) and
\(2(2a+1)!/(\pi/2)^{2a+2}\), multiplied by the exact constant
\(2C_\Gamma e^{1/2}\). The bound \(N!\leq N^N\) suffices.
For the lower bound restrict the positive-axis integral to \([a,a+1]\):
there \((1+4y^2)^{-3/4}\geq5^{-3/4}(a+1)^{-3/2}\),
\(y^{2a}\geq a^{2a}\), and
\(e^{-\pi y/2}\geq e^{-\pi(a+1)/2}\). These prove both sides of GEL15
with constants independent of \(a\).

The real-exponent function \(a\mapsto\log\mu_a\) is convex by the
same variance proof as GEL11. Since \(q\) is even and \(1\leq l\leq q\),
its secant on \([q,q+l]\) lies between the secants on \([q/2,q]\)
and \([2q,4q]\). GEL15 bounds both latter secants in absolute value
by \(O(\log q)\). Therefore
\[
 \left|\log\frac{\mu_{q+l}}{\mu_q}\right|=O(l\log q)=o(ql).
 \tag{GEL16}
\]
The original low moment is precisely
\(\mu_{q+l}/\mu_q=m_{2q+2l}/m_{2q}\), so no index is reinterpreted.

\paragraph{GEL7. The leading value of the unchanged signed centre.}
Keep the source product and its exact row intervals:
\[
 P_k=-\sum_{a=q}^{2q-1}\sum_{j=0}^{2l-1}\log(a+j+1/2)
     -\sum_{a=q+1}^{2q}\sum_{j=0}^{2l-1}\log(a+j+1/2).
\]
The two outer sums each have \(q\) terms. Since \(l/q\to0\), direct
Riemann sums, with \(\log\) Lipschitz on \([1,3]\), give
\[
 P_k=-4ql\log q-4ql\int_1^2\log x\,dx+o(ql)
     =-4ql\log q-4ql(2\log2-1)+o(ql).
 \tag{GEL17}
\]
Explicitly, omitting \((j+1/2)/q\) changes each logarithm by at most
\(2l/q\), hence changes the total by at most \(8l^2\); the ordinary
left and right Riemann sums together have error \(O(l)\).

The full finite product calculation WGP10 strengthens this leading
identity to the following enclosure for the very same \(P_k\):
\[
 \begin{split}
 P_k={}&-4lq\log q-4lq(2\log2-1)-4l^2\log2
              +\frac{8l^3+l}{6q}+\mathcal E_{q,l},\\
 -\left\{\frac l{6q}+\frac{8l^3+l}{12q^3}
                 +\frac{8l^4+2l^2}{3q^2}\right\}
 \leq{}&\mathcal E_{q,l}\leq\frac{l^2}{3q^2}.
 \end{split}
 \tag{GEL17a}
\]
Its proof uses the literal midpoint coordinates in WGP3, the signed
quadrature remainders WGP4--5, and the nonnegative cubic remainder
WGP6--9. The subsequent exact refinement WGP12--14 gives
\[
 \begin{split}
 P_k={}&-4lq\log q-4lq(2\log2-1)-4l^2\log2
 +\frac{8l^3+l}{6q}-\frac{l^4+l^2/4}{q^2}
          +\mathcal E^{(3)}_{q,l},\\
 -\frac l{6q}-\frac{8l^3+l}{12q^3}
 \leq{}&\mathcal E^{(3)}_{q,l}\leq
 \frac{l^2}{3q^2}+\frac{8l^4+2l^2}{6q^4}
                 +\frac{7(192l^5+80l^3-2l)}{1440q^3}.
 \end{split}
 \tag{GEL17b}
\]
These original endpoint contributions are retained when the finer
\(q\)-scale calculation is continued. In particular the term
\(-4l^2\log2\) has order \(q\) at \(m=1\), and the cubic term tends
to \(-1/256\) there; the leading limit does not remove either term.

The original quantity is exactly
\(W_k=P_k-\log(\mu_{q+l}/\mu_q)+\log(Z_{q+l}/Z_q)\).
Substitute GEL14, GEL16, and GEL17. The surviving
\(2l\log q\) is \(o(ql)\), so the leading growing-degree calculation is
\[
 \boxed{\frac{W_k}{ql}\longrightarrow
 \mathcal C_\Gamma:=2\mathcal L-4(2\log2-1),\qquad
 \mathcal L=\int_{u^2}^{v^2}\log x\,\rho_{2,\pi}(x)\,dx,\quad
 uK(\kappa)=2,\ vE(\kappa)=4.}
 \tag{GEL18}
\]
Every measure, exponent, matrix dimension, and source factor entering this
limit has been connected to its original object by GEL1--3 and GEL14.
The exact positive density and unique endpoint equations in GEL8 make
the constant an explicitly defined one-dimensional equilibrium integral.
This leading result has remainder \(o(ql)\). It supplies no claim of
an \(o(q)\) remainder, and hence does not erase the next finite-scale
calculation required by the original arithmetic receiver.

The fully proved original-polynomial recurrence estimate RWB20--27
also applies to this identical \(W_k\). It gives the finite bounds
\(lq/64\leq W_k\leq6lq\), and its narrower finite lower remainder
gives \(\liminf W_k/(lq)\geq4\log(27/(8\pi))\). Combining those
proved inequalities with the now proved existence of GEL18 yields
\[
 \boxed{4\log\frac{27}{8\pi}\leq\mathcal C_\Gamma\leq6,
              \qquad\mathcal C_\Gamma>0.}
 \tag{GEL18a}
\]
The strict sign is analytic: RWB22 proves
\(4\log(27/(8\pi))>104/365>0\) using rational integral bounds for
\(\pi\) and the elementary logarithm inequality. Numerical evaluation
of the equilibrium integral is unnecessary for this sign conclusion.

Finally LET20, whose error is \(O(k)=o(q)\) at fixed original
\(m,\delta,\gamma\), transfers the proved limit without changing its sign:
\[
 \frac{\mathfrak T_k}{ql}\longrightarrow\mathcal C_\Gamma,
 \qquad
 \frac{R_k^0-R_k^\sigma}{ql}\longrightarrow-\mathcal C_\Gamma.
 \tag{GEL19}
\]
The arithmetic identity remains
\(B_{\rm ar}=B_0+\delta_\sigma-\mathfrak T_k\), including its original
\(\delta_\sigma\) allowance. Thus GEL19 evaluates the Gamma summand
and its exact signed map into that identity; it makes no unsupported
growth assertion about either of the other two summands.

% END RX INLINE 35


% BEGIN RX INLINE 36; SHA256 9b42c91c8cbc5e8a01aaf130f9d5d5ac3855880bfe85dd72c8ac6add2bfb5479
% Complete new receiving proof; original predecessor files remain preserved.
\section{The evaluated Gamma term in the original arithmetic and cohomology receiver}
\label{sec:WGR}

Retain the actual packet $k=4l+1$, $e=1+k(m-1)$,
$q=e(k+1)^2$, $c=k/2$, $0<\delta<1/2$, $\gamma>2$,
and the original polynomials
\[
 \chi(S)=\prod_{a,b=0}^k
 [S-c-(2a-k)\delta-i(2b-k)\gamma]^e,\qquad
 f(S)=\prod_{j=0}^{l-1}(S-c-2j-1/2).
 \tag{WGR1}
\]
Here $E=\mathbb C[S]/(\chi)$, $J_N:\mathbb C[S]_{\le N}\to E$
is the literal remainder map, and the source line remains $S=c+iy$.
The parameters, primary multiplicities, and polynomial maps have not
been changed by the evaluation of $W_k$.

Let $\mathfrak r_\chi$ be the original scalar norm ratio at $N=q$,
and let $\mathfrak R_{2q-1},\mathfrak R_{2q}$ be the two original
relation-Gram ratios from JSR12. Write the exact scalar errors
\[
 \begin{aligned}
 e_0&=\log\mathfrak r_\chi-\log(m_{2q+2l}/m_{2q}),\\
 e_r&=\log\mathfrak R_{q+r-1}
       -\log\frac{\det\mathsf M_{q+l,r}}{\det\mathsf M_{q,r}},
       \qquad r=q,q+1,\\
 (\mathsf M_{a,r})_{ij}&=\int_{\mathbb R}y^{2a+i+j}
            \frac{|\Gamma(1/4+iy/2)|^2}{2\pi}\,dy,
            \qquad 0\le i,j<r .
 \end{aligned}\tag{WGR2}
\]
Thus the exponent index $m_j$ is the literal power $j$, and both
matrices use the full original mass $\sqrt{2\pi}$.
The exact coefficient map $P(S)\mapsto P(c+iy)$ has entries
$J_{hj}=\binom jh c^{j-h}i^h$ and determinant of absolute value one.
In particular the high reference Gram is its congruence
$J^*\mathsf M_{a,r}J$, not a different-coordinate determinant with
an unrecorded power of $q$. The bulk multiplication
$\chi(c+iy)/(iy)^q$ and the full positive inner and far forms are
the exact source maps and corrections proved in CTR and LET.

On the explicit domain $k\ge2048\sqrt{\delta^2+\gamma^2}$,
LET18 and CTR48 prove
\[
 e_0\in[-a_0,b_0],\qquad e_q+e_{q+1}\in[-A,B],
 \quad A=2lE_0+2(L_q+L_{q+1}),\quad
 B=A+(2q+1)\Delta_f,
 \tag{WGR3}
\]
where $a_0=E_0+2L_1$, $b_0=a_0+\Delta_f$, and every constant
$E_0,\Delta_f,L_r$ is the finite original integral constant printed
in LET7--16 and BHR2. In particular $A,B=O(k)$ for fixed $m=1$
and $O(1)$ for fixed $m\ge2$; $a_0,b_0=O(1)$ and $O(k^{-1})$
in those respective branches. These statements were proved on the
full source, including the low rank one, and are used within their
stated domain here.

The exact identity LET19, retaining the low minus sign, is
\[
 \mathfrak T_k=W_k-e_0+e_q+e_{q+1},\qquad
 \mathcal H_k:=R_k^0-R_k^\sigma=-\mathfrak T_k.
 \tag{WGR4}
\]
RWB20--27 now supplies an evaluated pair of finite numbers
\[
 \begin{aligned}
 w^-_{q,l}=\max\bigg\{\frac{lq}{64},\;&
 4lq\log\frac{27}{8\pi}-4l^2-6l\log(3/2)-3l\\
 &-\frac{9l}{4q-1}-\frac{2l(l+1/4)}q\bigg\},
 \qquad w^+_{q,l}=6lq,\qquad
 w^-_{q,l}\le W_k\le w^+_{q,l}.
 \end{aligned}\tag{WGR5}
\]
This assertion holds at every original $l,m\ge1$; the extra threshold
in WGR3 is required only when returning to the original $\chi$ source.
Substituting WGR5 and then the three signed errors into WGR4 proves
\[
 \boxed{-w^+_{q,l}-B-a_0\le\mathcal H_k
                 \le-w^-_{q,l}+A+b_0.}
 \tag{WGR6}
\]
Indeed $\mathcal H_k=-W_k+e_0-e_q-e_{q+1}$.
This proves both signs separately; no positive high error has been
assigned the sign of the negative low endpoint.

\paragraph{The complete retained arithmetic and mixed-state map.}
Use exactly the previous signed arithmetic difference
$\delta_k^\sigma$ and the full baseline $\mathcal B_k^0$ from JSR16--21.
For either original source, the quotient, kernel and boundary forms
are
\[
 G_N=(J_NH_N^{-1}J_N^*)^{-1},\qquad
 H_N^K=I_K^*G_NI_K,\qquad
 Q_N^B=(\Lambda G_N^{-1}\Lambda^*)^{-1}.
 \tag{WGR7}
\]
Here $\Lambda:E\to B_{\mathrm{obs}}=\operatorname{im}\Lambda$ is the
original boundary observation, $K=\ker\Lambda$, and $I_K:K\hookrightarrow E$
is the original inclusion. The chosen original lift
$S_*:B_{\mathrm{obs}}\to E$ satisfies $\Lambda S_*=I_{B_{\mathrm{obs}}}$.
Thus $[I_K,S_*]:K\oplus B_{\mathrm{obs}}\to E$ is the exact
coefficient isomorphism. The exact sequence is represented in this common
coefficient frame $[I_K,S_*]$. Orthogonal projection subtracts a kernel
column from each lift in $S_*$ and is a triangular change with
determinant one, proving
\[
 \det H_N^K\det Q_N^B
       =|\det[I_K,S_*]|^2\det G_N.
 \tag{WGR8}
\]
The same frame is used for both sources, so its determinant cancels
in their ratio. The source comparison gives the original nonnegative
deficits $b_N^K+b_N^B=b_N\le\Xi_{k,N}$ and the four signed sums
\[
 \begin{aligned}
 X_k&=-b_{q-1}-b_q+b_{2q-1}+b_{2q},\\
 a_k^-&=\Xi_{k,q-1}+\Xi_{k,q},\quad
 a_k^+=\Xi_{k,2q-1}+\Xi_{k,2q},\qquad
 -a_k^-\le X_k\le a_k^+ .
 \end{aligned}\tag{WGR9}
\]
The same formula applies to $\Delta F_K,\Delta F_B$ by using the
corresponding superscripts, and to $\delta_k^\sigma$ by using the
full quotient. The last inequality follows by dropping the two
positive terms for the lower bound and the two negative terms for
the upper bound. Thus the original mixed kernel and boundary
directions remain in the exact map, rather than being lost in a
scalar positivity assertion. The previous finite source estimates
$a_k^-\le D_k^-$ and $a_k^+\le D_k^+$, with all original constants
in JSR17--18, continue to apply separately.

The original arithmetic equality and the evaluated finite bounds are
therefore
\[
 \begin{aligned}
 \mathcal B_k^{\rm ar}&=\mathcal B_k^0+\mathcal H_k+\delta_k^\sigma,
       &\Delta_k^\Gamma&=\mathcal H_k+\delta_k^\sigma,\\
 \mathcal B_k^0-w^+_{q,l}-B-a_0-a_k^-
 &\le\mathcal B_k^{\rm ar}
       \le\mathcal B_k^0-w^-_{q,l}+A+b_0+a_k^+ .
 \end{aligned}\tag{WGR10}
\]
Add WGR9 to WGR6 to prove WGR10. Replacing $a_k^\pm$ by their
respective outer bounds $D_k^\pm$ gives the corresponding fully
estimated arithmetic interval. In particular the original nonnegative
HC residual has the finite enclosure
\[
 \begin{aligned}
 \max\{0,\mathcal B_k^0-w^+_{q,l}-B-a_0-D_k^-
                     -4q\log(D_hk)\}
 &\le\mathcal R_k:=\mathcal B_k^{\rm ar}-4q\log(D_hk)\\
 &\le\mathcal B_k^0-w^-_{q,l}+A+b_0+D_k^+
                     -4q\log(D_hk).
 \end{aligned}\tag{WGR11}
\]
Its constant $D_h$ and its threshold are the previous original ones.
For every finite exact $W_k$, WGR6 contains the sharper WRC5 interval.
Consequently it provides an evaluated bound with the original source
constants on the growing family; it does not claim to narrow an interval already using
the exact scalar $W_k$ as input.

\paragraph{A definite growing arithmetic contribution.}
Fix an original packet $(\delta,\gamma,m)$ and let $k=4l+1\to\infty$.
The already proved JSR18--21 finite source estimates give
\[
 \begin{cases}
 \delta_k^\sigma/(kq)=O_h((\log k/k)^{5/7}),&m=1,\\
 \displaystyle\liminf\frac{\delta_k^\sigma}{q\log k}\ge-(64m+245),
 \quad\displaystyle\limsup\frac{\delta_k^\sigma}{q\log k}
                                      \le128m+490,&m\ge2.
 \end{cases}\tag{WGR12}
\]
Thus $\delta_k^\sigma/(kq)\to0$ in both branches. In the second
branch multiply the finite upper and lower constants by $\log k/k\to0$;
in the first branch the displayed rate tends to zero. Also WGR3 implies
$(e_0-e_q-e_{q+1})/(kq)\to0$, retaining its original asymmetric bounds.
Since $l/k=(k-1)/(4k)\to1/4$, RWB23 and RWB27 yield
\[
 \boxed{-\frac32\le\liminf\frac{\Delta_k^\Gamma}{kq}
 \le\limsup\frac{\Delta_k^\Gamma}{kq}
 \le-\log\frac{27}{8\pi}<0.}
 \tag{WGR13}
\]
This uses $\Delta_k^\Gamma=-W_k+e_0-e_q-e_{q+1}+\delta_k^\sigma$
and the proved finite bounds, so it does not infer a limit from the
exploratory numerical recurrence. In particular the Gamma/arithmetic
difference is strictly negative for all sufficiently large original
degrees. It remains a term in the full arithmetic equality WGR10.

The preceding full-source equilibrium calculation GEL18--19 evaluates
this term more precisely, proving
\[
 \boxed{\frac{\Delta_k^\Gamma}{kq}\longrightarrow
       -\frac{\mathcal C_\Gamma}{4},\qquad
 4\log\frac{27}{8\pi}\le\mathcal C_\Gamma\le6.}
 \tag{WGR14}
\]
Indeed GEL18 proves $W_k/(lq)\to\mathcal C_\Gamma$, while the
two other summands divided by $kq$ vanish by WGR12 and WGR3.
The finite RWB bounds give the stated interval for the exact equilibrium
integral $\mathcal C_\Gamma$. Combining this result with WGR10--11
gives the evaluated relation to the original baseline and residual:
\[
 \frac{\mathcal B_k^0-\mathcal B_k^{\rm ar}}{kq}
       \longrightarrow\frac{\mathcal C_\Gamma}{4},\qquad
 \frac{\mathcal R_k-\mathcal B_k^0}{kq}
       \longrightarrow-\frac{\mathcal C_\Gamma}{4},\qquad
 \liminf\frac{\mathcal B_k^0}{kq}\ge\frac{\mathcal C_\Gamma}{4}.
 \tag{WGR15}
\]
The second assertion also uses $4\log(D_hk)/k\to0$ at fixed $D_h>0$;
the third follows from the independent original inequality
$\mathcal R_k\ge0$. These are evaluated differences and a necessary
baseline bound with the original arithmetic object retained. They do
not assert that the baseline or the nonnegative residual has vanished.
The finite $q$-scale terms in WGP13 and the finite allowances in
WGR10--12 remain available for the next calculation on those actual
objects. In particular the $o(lq)$ remainder in GEL18 has not been
replaced by an unproved $o(q)$ remainder.

% END RX INLINE 36



\clearpage
\part*{Complete proofs of the original baseline and its intrinsic mixed return}
The following complete source bodies prove the baseline, the original
finite low and high corrections, and the intrinsic source and mixed-row
maps used above. The packet is unchanged throughout: only the already
present source orders $s=1$ and $s=k$ enter these receiving statements.
The original masses, relation polynomial, observation, four endpoint
signs and fibre kernels are retained in the displayed exact maps.

% BEGIN RX INLINE 37; SHA256 f49ae1f85afb633a7a00ee9a74e643fef3a57e072508dc28b1ebbbd197992468



\section{Original source and the homogeneous remainder map}

Retain the original packet integers
\[
 k=4l+1,\quad l\ge1,\quad m\ge1,\quad
 e=1+k(m-1),\quad q=e(k+1)^2=2n,\quad c=k/2,
\]
with $m$ an integer. Define the homogeneous comparison polynomial and
its literal quotient by
\[
 \chi_0(S)=(S-c)^q,\qquad E_0=\C[S]/(\chi_0),\qquad
 J_N^0:\C[S]_{\le N}\longrightarrow E_0.
 \tag{HBL1}
\]
Here $J_N^0$ is polynomial remainder, and the original source line is
$S=c+iy$. The measure and its full mass remain
\[
 d\sigma(y)=\frac{|\Gamma(1/4+iy/2)|^2}{2\pi}\,dy,
 \qquad \gamma_0=\sigma(\R)=\sqrt{2\pi}.
 \tag{HBL2}
\]
For $N\ge q-1$, let $H_N^\sigma$ be the source Gram in the increasing
$S$ coefficient frame, and set
\[
 G_N^{\sigma,0}=(J_N^0(H_N^\sigma)^{-1}(J_N^0)^*)^{-1}.
 \tag{HBL3}
\]
The map is onto in this range, so this is a positive definite quotient
Gram. Define the four-endpoint homogeneous baseline with its original
signs by
\[
 F_\sigma^{\rm hom}(q)=
 \log\det G_{q-1}^{\sigma,0}+\log\det G_q^{\sigma,0}
 -\log\det G_{2q-1}^{\sigma,0}-\log\det G_{2q}^{\sigma,0}.
 \tag{HBL4}
\]
The result below evaluates precisely HBL4. Its comparison with a
displaced original polynomial $\chi$ requires the separate original
source comparison; that comparison is not asserted in this calculation.

\section{The exact source/relation determinant identity}

Put $r=N+1-q$ and retain the full real-coordinate relation matrix
\[
 (M_{q,r})_{ij}=\int_{\R}y^{2q+i+j}\,d\sigma(y),
 \qquad 0\le i,j<r,
 \qquad \det M_{q,0}=1.
\]
The frame
\[
 1,S-c,\ldots,(S-c)^{q-1},\quad
 \chi_0,\chi_0(S-c),\ldots,\chi_0(S-c)^{r-1}
\]
is the increasing translated monomial frame, so its determinant in
the original coefficient coordinates is one. Its first $q$ columns
represent a basis of $E_0$ and its last $r$ columns form the entire
relation kernel. In this frame the remainder map is $(I_q\ 0)$.
The Schur complement of the positive relation block is the quotient
Gram in HBL3: this follows either by multiplying the block inverse,
or by minimizing the source norm in a fixed remainder fibre.
The relation columns evaluate to $(iy)^{q+j}$, so their Gram is a
congruence of $M_{q,r}$ by the diagonal matrix with entries $i^{q+j}$.
Its determinant has absolute value one. Therefore
\[
 \boxed{\det G_N^{\sigma,0}
             =\frac{\det H_N^\sigma}{\det M_{q,N+1-q}}.}
 \tag{HBL5}
\]
This also proves HBL5 at $r=0$ with the empty determinant convention.

The original monic real-line polynomials $p_j(y)$ have squared norms
\[
 \gamma_j=\sqrt{2\pi}\,j!(1/2)_j
          =\sqrt{2\pi}\frac{(2j)!}{4^j}.
 \tag{HBL6}
\]
These are the complete HCT5 norms. The polynomials
$i^jp_j((S-c)/i)$ are monic in the original $S$ coordinate and have
the same norms on $S=c+iy$. Their monic triangular coefficient map
has determinant one. It follows that
$\det H_N^\sigma=\prod_{j=0}^N\gamma_j$ with the full mass retained.

Let $\mu_a=\int y^{2a}d\sigma(y)$ and
$H_s^{(a)}=\det[\mu_{a+i+j}]_{i,j<s}$, with $H_0^{(a)}=1$.
The exact even/odd permutation of the original moment matrices gives
\[
 \det M_{q,q}=H_n^{(q)}H_n^{(q+1)},\qquad
 \det M_{q,q+1}=H_{n+1}^{(q)}H_n^{(q+1)},\qquad n=q/2.
\]
The permutation is applied to rows and columns alike; its determinant
sign is squared. Also $\det M_{q,1}=\mu_q$. Define
\[
 Z_q=H_n^{(q)}H_{n+1}^{(q)}(H_n^{(q+1)})^2.
\]
Substitution into HBL4--6 proves the exact finite equality
\[
 \boxed{F_\sigma^{\rm hom}(q)=S_q+\log Z_q-\log\mu_q,}
 \tag{HBL7}
\]
where all endpoints in the source contribution are
\[
 S_q=-\sum_{j=q}^{2q-1}\log\gamma_j
     -\sum_{j=q+1}^{2q}\log\gamma_j
     =-\log\gamma_q-2\sum_{j=q+1}^{2q-1}\log\gamma_j
      -\log\gamma_{2q}.
 \tag{HBL8}
\]

\section{The source products at their exact degree scale}

For every integer $d\ge1$, monotonicity of $\log x$ on each unit
interval gives
\[
 d\log d-d+1\le\log(d!)
                \le d\log d-d+1+\log d.
\]
At $d=2j$, HBL6 consequently gives, uniformly for $q\le j\le2q$,
\[
 \log\gamma_j=2j\log j-2j+O(\log q).
\]
The constant $\log\sqrt{2\pi}$ is included in this finite error.
For $x=j/q$, the displayed main expression is
$2qx(\log q+\log x-1)$. The left and right sums in HBL8 have
$q$ terms each. Each differs from its integral times $q$ by
$O(q\log q)$, because the total variation of this increasing function
on $[1,2]$ has that order. The $2q$ individual norm errors also sum
to $O(q\log q)$. The literal integrals
\[
 \int_1^2x\,dx=\frac32,\qquad
 \int_1^2x\log x\,dx=2\log2-\frac34
\]
therefore prove
\[
 \boxed{S_q=-6q^2\log q+(9-8\log2)q^2+O(q\log q).}
 \tag{HBL9}
\]
In particular both distinct endpoints in HBL8 were used before the
error was estimated.

\section{The complete high relation determinant at scale $q^2$}

We use the already proved original-source free-energy theorem GEL3,
GEL7--10 and the EIQ equilibrium calculation, with their original
measure and coordinate factors. Their definitions are repeated here
to specify every parameter entering the limit. For $\alpha,\beta>0$,
\[
 \mathcal F(\alpha,\beta)=\sup_{\mu\in\mathcal P((0,\infty))}
 \left\{\iint\log|x-y|\,d\mu(x)d\mu(y)
            -\int(\beta\sqrt x-\alpha\log x)\,d\mu(x)\right\}.
 \tag{HBL10}
\]
The rigorous extended definition is the single bounded-above kernel
$\log|x-y|-(V(x)+V(y))/2$, so infinite terms are never subtracted.
Its unique equilibrium probability has compact support bounded away
from zero, and $\mathcal F$ is continuous on the positive quadrant.

For the exact scale $t=q^2x$, define
\[
 d\omega(x)=\tfrac12x^{-1/2}e^{-\sqrt x}\,dx,
 \quad B_q(x)=2qe^{\sqrt x}\sigma(q\sqrt x),
\]
\[
 J_{s,q}(a)=\frac1{s!}\int\prod_{i<j}(x_i-x_j)^2
                    \prod_i x_i^aB_q(x_i)\,d\omega(x_i).
\]
Here $d\nu(q^2x)=B_q(x)d\omega(x)$ is the exact density identity
including the whole original mass and Jacobian. GEL3 and GEL10 state
\[
 H_s^{(a)}=q^{2as+2s(s-1)}J_{s,q}(a),\qquad
 \frac{\log J_{s,q}(s\alpha)}{s^2}
       \longrightarrow\mathcal F(\alpha,\pi\theta/2)
       \quad(q/s\to\theta>0).
 \tag{HBL11}
\]
This is a previously proved theorem about these literal Gamma integrals;
the present calculation applies it, without replacing its source by an
unproved asymptotic density.

Its fixed-$\alpha$ limit also applies when $\alpha_s\to2$. Here are
the needed details. The functions
$f_s(\alpha)=s^{-2}\log J_{s,q}(s\alpha)$ are convex, since their
second derivatives are variances of the sum of the original logarithmic
coordinates, multiplied by a positive constant. For fixed $\eta>0$,
convexity bounds their values on $[2-\eta,2+\eta]$ above by the larger
of the endpoint values. If $\alpha_s\ge2$, the lower secant inequality is
\[
 f_s(\alpha_s)\ge f_s(2)
 +\frac{\alpha_s-2}{\eta}\big(f_s(2)-f_s(2-\eta)\big).
\]
For $\alpha_s\le2$ the corresponding lower bound is
\[
 f_s(\alpha_s)\ge f_s(2)
 +\frac{\alpha_s-2}{\eta}\big(f_s(2+\eta)-f_s(2)\big).
\]
The fixed-point convergences in HBL11 make their differences bounded,
so both lower bounds have liminf at least $\mathcal F(2,\pi)$ when
$q/s\to2$. The upper bound and continuity let $\eta$ tend to zero.
This proves the moving-argument limit required below.

For $Z_q$, use $(s,a)=(n,q),(n+1,q),(n,q+1),(n,q+1)$ exactly.
Each has $q/s\to2$ and $a/s\to2$. The four exact powers in HBL11 are
\[
 \frac32q^2-q,\qquad\frac32q^2+3q,
 \qquad\frac32q^2,\qquad\frac32q^2.
\]
Their sum is $6q^2+2q$. The sum of the four squared sizes divided
by $q^2$ tends to one. Thus
\[
 \boxed{\log Z_q=(6q^2+2q)\log q
                      +q^2\mathcal F(2,\pi)+o(q^2).}
 \tag{HBL12}
\]

The scalar low moment is negligible on this scale with its full mass
still present. The established original Gamma contour bound gives
$\sigma(y)\le C_Ue^{-|y|}$ for a fixed $C_U>0$, so
$\mu_q\le2C_U(2q)!$. Positivity on $[1,2]$ gives
$\mu_q\ge2\int_1^2\sigma(y)\,dy>0$. Hence
$\log\mu_q=O(q\log q)=o(q^2)$.
Combining this with HBL7, HBL9 and HBL12 proves the exact leading value
\[
 \boxed{\frac{F_\sigma^{\rm hom}(q)}{q^2}\longrightarrow
   \mathcal C_B:=9-8\log2+\mathcal F(2,\pi).}
 \tag{HBL13}
\]
This holds along the original packet degrees as they tend to infinity;
the homogeneous quantity depends only on $q$ after the explicit
unit-determinant translation by $c$. There is no error estimate of
order $o(q)$ claimed in HBL12 or HBL13.

\section{An exact endpoint-integral presentation of the coefficient}

To evaluate the constant in the same equilibrium coordinates as GEL18,
let $u,v$ be the unique positive endpoints with $0<u<v$ and
\[
 uK(\kappa)=2,\qquad vE(\kappa)=4,\qquad
 \kappa^2=1-u^2/v^2,
\]
where
$K(\kappa)=\int_0^{\pi/2}(1-\kappa^2\sin^2\theta)^{-1/2}d\theta$
and
$E(\kappa)=\int_0^{\pi/2}(1-\kappa^2\sin^2\theta)^{1/2}d\theta$.
The EIQ/GEL equilibrium density is exactly
\[
 \rho(x)=\frac{\sqrt{(v^2-x)(x-u^2)}}{4\pi x}
   \int_0^\infty
    \frac{\sqrt s\,ds}{(x+s)\sqrt{(u^2+s)(v^2+s)}},
 \qquad u^2<x<v^2.
 \tag{HBL14}
\]
It has mass one. Write
$\mathcal L=\int\log x\,\rho(x)dx$ and
$M=\int\sqrt x\,\rho(x)dx$.
Under the literal dilation $x\mapsto r^2x$ of this probability,
the value of the functional HBL10 changes by
$6\log r-\pi(r-1)M$. The equilibrium maximizes it at $r=1$;
differentiating this explicit scalar expression there gives $\pi M=6$.
Thus $\int V_{2,\pi}\rho=6-2\mathcal L$.

The EIQ variational identity says that
$V_{2,\pi}(x)-2\int\log|x-y|\rho(y)dy$ is a constant $\ell$ on
the full support, including its endpoints. Set $b=v^2$. At the right
endpoint the same identity is
\[
 \ell=\pi v-2\log b
            -2\int_{u^2}^{b}\log(b-y)\rho(y)\,dy.
\]
All these integrals are finite; the support is bounded away from zero
and the density has square-root endpoint vanishing. Integrating the
variational identity against $\rho$ gives
$\mathcal F(2,\pi)=-\ell/2-3+\mathcal L$. Consequently HBL13 has
the exact one-dimensional endpoint-integral presentation
\[
 \boxed{\begin{aligned}
 \mathcal C_B={}&6-8\log2+\mathcal L-\ell/2\\
 ={}&6-8\log2+\log b-\frac{\pi v}{2}
       +\int_{u^2}^{b}\big[\log x+\log(b-x)\big]\rho(x)\,dx.
 \end{aligned}}
 \tag{HBL15}
\]

\section{A strictly positive rational bound from an explicit competitor}

We prove the sign without numerical evaluation of HBL14--15. Use the
arcsine probability on the literal interval $[1,9]$,
\[
 d\eta(x)=\frac{1_{(1,9)}(x)}{\pi\sqrt{(x-1)(9-x)}}\,dx,
 \qquad x=5+4\cos\theta,\quad 0\le\theta\le\pi.
\]
The substitution gives the uniform measure $d\theta/\pi$ with the
orientation reversed when $x$ increases. Its three exact integrals are
\[
 \iint\log|x-y|\,d\eta(x)d\eta(y)=\log2,\quad
 \int\log x\,d\eta(x)=2\log2,\quad
 \int\sqrt x\,d\eta(x)=\frac6\pi E(\sqrt{8/9}).
 \tag{HBL16}
\]
Here are proofs retaining the constants. The double-angle identity and
the substitution $\theta\mapsto\pi/2-\theta$ give
$\int_0^\pi\log\sin\theta\,d\theta=-\pi\log2$; all endpoint
logarithms are integrable. Factoring
$\cos\theta-\cos\phi=-2\sin((\theta+\phi)/2)
\sin((\theta-\phi)/2)$, then joining the two substituted intervals,
gives
\[
 \frac1\pi\int_0^\pi\log|\cos\theta-\cos\phi|d\phi
 =\log2+\frac2\pi\int_{(\theta-\pi)/2}^{(\theta+\pi)/2}
                              \log|\sin t|\,dt=-\log2.
\]
The interval has length $\pi$ and the integrand has period $\pi$.
Adding $\log4$ proves the first HBL16 identity, and also proves finite
logarithmic energy of this competitor. For the second, write
$5+4\cos\theta=|2+e^{i\theta}|^2$ and expand
$\log(1+e^{i\theta}/2)$ in its uniformly absolutely convergent power
series. Every nonconstant cosine has integral zero on $[0,\pi]$,
leaving $2\log2$. Finally
$5+4\cos\theta=9(1-(8/9)\sin^2(\theta/2))$; the substitution
$t=\theta/2$ gives the third identity.

Since HBL10 is a supremum over all such probability measures,
the competitor proves
\[
 \mathcal F(2,\pi)\ge5\log2-6E(\sqrt{8/9}),\qquad
 \mathcal C_B\ge9-3\log2-6E(\sqrt{8/9}).
 \tag{HBL17}
\]
An elementary polynomial bound makes this strictly positive. For
$0\le z\le1$ put $P(z)=1-z/2-z^2/8-z^3/16$. It is positive there,
since $P(z)\ge P(1)=5/16$, and direct multiplication gives
\[
 P(z)^2-(1-z)=\frac{5z^4}{64}+\frac{z^5}{64}+\frac{z^6}{256}\ge0.
\]
Therefore $\sqrt{1-z}\le P(z)$. The integrations by parts
$\int_0^{\pi/2}\sin^{2r}t\,dt=
 (2r-1)/(2r)\int_0^{\pi/2}\sin^{2r-2}t\,dt$ give the exact
second, fourth and sixth power integrals $\pi/4,3\pi/16,5\pi/32$.
At $z=(8/9)\sin^2t$ the polynomial bound consequently gives
\[
 E(\sqrt{8/9})\le\frac\pi2
 \left(1-\frac29-\frac1{27}-\frac{10}{729}\right)
 =\frac{265\pi}{729}.
\]
The identity
$22/7-\pi=\int_0^1x^4(1-x)^4/(1+x^2)\,dx>0$
follows by polynomial division and the arctangent integral.
Also the positive exponential series gives
$e^{7/10}>1+7/10+(7/10)^2/2+(7/10)^3/6=12013/6000>2$,
so $\log2<7/10$. Thus HBL17 implies the entirely rational strict bound
\[
 \boxed{\mathcal C_B\ge9-3\log2-\frac{530\pi}{243}
       >9-\frac{21}{10}-\frac{11660}{1701}
       =\frac{769}{17010}>0.}
 \tag{HBL18}
\]
This is a bound on the exact coefficient in HBL13, proved by a finite
competitor calculation. In particular the homogeneous baseline has a
strictly positive leading term of order $q^2$. Its transfer to a
displaced packet, and any subsequent combination with the arithmetic
baseline or HC residual, must use the corresponding original maps and
source comparisons rather than identifying those objects with HBL4.

\paragraph{Complete proved dependencies used.}
The original Gamma monic norms are HCT5 in
\texttt{independent\_contiguous.tex}. The exact density map, free-energy
limit, and its original-source bounds are GEL1--10 in
\texttt{GAMMA\_ENSEMBLE\_LEADING\_RETURN.tex}. The compact equilibrium,
endpoint equations, positive density, and full variational identity are
proved in \texttt{EXACT\_SQRT\_LOG\_EQUILIBRIUM.tex}. This note supplies
the complete new quotient calculation, moving-argument specialization,
original source expansion, coefficient identity, and its analytic sign.


% END RX INLINE 37


% BEGIN RX INLINE 38; SHA256 a53703c284d816930291aa4b02373564fe8cd5cf7c93a1d0517354614f438c30
% New original-packet baseline proof; predecessor delivery21 is immutable.
\section{The original packet baseline and its complete receiving maps}
\label{sec:BSL}

Retain the original parameters and all roots, counted with their orders:
\[
 k=4l+1,\quad l,m\geq1,\quad e=1+k(m-1),\quad
 q=e(k+1)^2=2n,\quad c=k/2,\quad
 0<\delta<1/2,\quad\gamma>2,
\]
\[
 \chi(S)=\prod_{a,b=0}^{k}
 \bigl[S-c-(2a-k)\delta-i(2b-k)\gamma\bigr]^e,
 \qquad d\sigma(y)=\frac{|\Gamma(1/4+iy/2)|^2}{2\pi}\,dy.
 \tag{BSL1}
\]
The letter $e$ here is the positive integer order; exponential functions
below are written $\exp$. Every limit in this section holds with the
original $m,\delta,\gamma$ fixed and $l\to\infty$.

\subsection{The exact source, relation, and quotient determinants}
Let $\mathcal P_N$ be the polynomials in $S$ of degree at most $N$,
with the stated source inner product evaluated on $S=c+iy$. Set
$E_\chi=\mathbb C[S]/(\chi)$ and let
$J_N:\mathcal P_N\to E_\chi$ send a polynomial to its remainder.
For $N\geq q-1$, this is onto, and its kernel is the literal image of
\[
 I_{\chi,N}:\mathcal P_{N-q}\longrightarrow\mathcal P_N,
 \qquad Q\longmapsto\chi Q.
 \tag{BSL2}
\]
The domain is the zero vector space when $N=q-1$. Let $H_N^\sigma$
be the source Gram in $(1,S,\ldots,S^N)$, $G_N^\sigma$ the minimum
norm quotient Gram in the fixed remainder frame, and
\[
 A_{\chi,r}=\operatorname{Gram}_{\sigma}
       (\chi,\chi S,\ldots,\chi S^{r-1}),\qquad r=N+1-q.
\]
All these nonempty Grams are positive definite. A nonzero source
polynomial is nonzero on an open real interval, and $\sigma$ is
strictly positive there. The quotient minimum exists and is positive
because this is an onto map between finite dimensional spaces with a
positive source form. The same argument applies to every relation form.

The ordered change of frame
$(1,\ldots,S^{q-1},\chi,\ldots,\chi S^{r-1})$ has determinant one:
its column of degree $j$ has leading coefficient one and no higher
coefficient. In this frame, minimizing the quadratic form over its last
$r$ coefficients gives the Schur complement of $A_{\chi,r}$.
Indeed, for blocks $\left(\begin{smallmatrix}A&C\\C^*&D\end{smallmatrix}\right)$,
the minimizing last coefficient is $-D^{-1}C^*z$, its squared norm is
$z^*(A-CD^{-1}C^*)z$, and triangular block elimination has determinant
one. Consequently, with the empty determinant equal to one,
\[
 \det G_N^\sigma=\frac{\det H_N^\sigma}{\det A_{\chi,N+1-q}}.
 \tag{BSL3}
\]
This equality retains the complete relation space before forming a
quotient. It gives the same form as
$G_N^\sigma=(J_N(H_N^\sigma)^{-1}J_N^*)^{-1}$: the minimizer of
$P^*H_N^\sigma P$ subject to $J_NP=z$ is
$(H_N^\sigma)^{-1}J_N^*G_N^\sigma z$, as multiplication by $J_N$
and completing the square directly verify.

Let $\gamma_j=\sqrt{2\pi}(2j)!/4^j$ be the original monic Gamma
norms, including $\gamma_0=\sqrt{2\pi}$. Their complete derivation is
HCT1--9. Monic orthogonalization has determinant one, so
$\det H_N^\sigma=\prod_{j=0}^N\gamma_j$. Define
\[
 \begin{split}
 S_q&=-\log\gamma_q-2\sum_{j=q+1}^{2q-1}\log\gamma_j
                      -\log\gamma_{2q},\\
 \mathcal B_k^\sigma
 &=\log\frac{\det G_{q-1}^\sigma\det G_q^\sigma}
                   {\det G_{2q-1}^\sigma\det G_{2q}^\sigma}.
 \end{split}
\]
Substitution of all four ranks $0,1,q,q+1$ in BSL3 proves the exact
baseline formula
\[
 \boxed{\mathcal B_k^\sigma
 =S_q+\log\det A_{\chi,q}+\log\det A_{\chi,q+1}
                         -\log\det A_{\chi,1}.}
 \tag{BSL4}
\]
Both high relation spaces and the one dimensional low relation space
are present with their original signs.

\subsection{A finite comparison retaining every original root}
Define the literal power moment matrix and original coefficient map
\[
 (M_{q,r})_{ij}=\int_{\mathbb R}y^{2q+i+j}\,d\sigma(y),\quad
 T_{h,j}=\binom jh c^{j-h}i^h\quad(0\leq h\leq j<r),
 \tag{BSL5}
\]
with other entries zero. The matrix $T$ sends coefficients of $P(S)$
to those of $P(c+iy)$; its inverse comes from $y=(S-c)/i$.
Since $|\det T|=1$, the Gram of
$(S-c)^q(1,S,\ldots,S^{r-1})$ on the original source is
$T^*M_{q,r}T$ and has determinant $\det M_{q,r}$.
The phase $i^q$ has squared modulus one, as used in this equality.

Put $d=\delta^2$, $g=\gamma^2$, $\epsilon=2^{-10}$, and assume
the same explicit finite threshold as CTR23 and LET4,
$k\geq2048\sqrt{d+g}$. Partition the line into
$I=\{|y|<q\epsilon\}$, $B=\{q\epsilon\leq|y|\leq64q\}$,
and $F=\{|y|>64q\}$. Retain
\[
 a_k=\frac{k(k+2)(d-g)}{3q}<0,\qquad
 t_\epsilon=\frac{k^2(d+g)}{q^2\epsilon^2}\leq\frac14,\qquad
 d_\chi=\frac{qt_\epsilon^2}{2(1-t_\epsilon)},
\]
\[
 b_- =a_k\epsilon^{-2}-d_\chi,\qquad
 b_+ =a_k64^{-2}+d_\chi.
 \tag{BSL6}
\]
Here $d_\chi$ is a nonnegative Taylor remainder bound, not a
determinant or the earlier Christoffel correction. At $y=qz$ in $B$,
the exact product gives
\[
 \psi_\chi(qz):=\log\frac{|\chi(c+iqz)|^2}{|qz|^{2q}}
       =\frac{a_k}{z^2}+E_\chi(qz),\quad |E_\chi(qz)|\leq d_\chi,
 \qquad b_-\leq\psi_\chi(qz)\leq b_+.
 \tag{BSL7}
\]
To check every coefficient, write each factor as
$iqz(1+i\zeta_{ab}/(qz))$ with
$\zeta_{ab}=(2a-k)\delta+i(2b-k)\gamma$. The logarithmic Taylor
series is absolutely convergent because $|\zeta_{ab}/(qz)|\leq1/2$.
The permutation $(a,b)\mapsto(k-a,k-b)$ cancels its odd powers.
The second power is
$e\operatorname{Re}\sum_{a,b}\zeta_{ab}^2/(qz)^2=a_k/z^2$,
using $\sum_{a=0}^k(2a-k)^2=(k+1)k(k+2)/3$ and
$\sum_{a=0}^k(2a-k)=0$. The even remainder has absolute value at
most $q\sum_{j\geq2}t_\epsilon^j/j\leq d_\chi$.
The negative sign of $a_k$ determines which end of the interval for
$z^{-2}$ gives each bound in BSL6.

The exact source map on $L^2(B,d\sigma)$ is multiplication by
$\chi(c+iy)/(iy)^q$. It and its inverse are bounded there because
the numerator and denominator have no zero on this compact set.
Composing this map with $P\mapsto(iy)^qP(c+iy)$ gives precisely
$P\mapsto\chi(c+iy)P(c+iy)$. Taking these Hilbert Grams and using
BSL7 proves the full vector inequality
\[
 \exp(b_-)\,T^*M_{q,r}^BT
 \preceq A_{\chi,r}^B
 \preceq\exp(b_+)\,T^*M_{q,r}^BT.
 \tag{BSL8}
\]
No map between different remainder algebras has been presumed here;
BSL8 is proved on the specified common relation coefficient space.

For either full Gram $H=H^I+H^B+H^F$, keep its exact tail value
\[
 t_H=\log\det\left(I_r+(H^B)^{-1/2}
                       (H^I+H^F)(H^B)^{-1/2}\right),\qquad
 \log\det H=\log\det H^B+t_H.
\]
The complete contour and polynomial estimates CTR37--44, and the
separate scalar proof LET10--16 when $r=1$, prove
\[
 0\leq t_H\leq L_r,
 \quad L_r=r\log(1+\kappa_{I,r}+\kappa_F),
\]
\[
 \kappa_{I,r}=\frac{2\epsilon r^2}{\sqrt{\cos1}}
       \left(\frac43\,2^{-13}\exp(2\pi)\right)^q,
 \qquad \kappa_F=\frac{512q}{59\sqrt{\cos1}}\exp(-19q),
 \quad r\in\{1,q,q+1\}.
 \tag{BSL9}
\]
The first base is smaller than $\exp(-2)$; these are finite
positive bounds for each original tail, not omitted integrals.
Define the actual finite error
\[
 \eta_r=\log\det A_{\chi,r}-\log\det M_{q,r}.
\]
Congruence of BSL8 with the positive power Gram bounds each of its
$r$ eigenvalues. Multiplying them and adding the two original tail
values gives the exact expression and its enclosure
\[
 \begin{split}
 \eta_r={}&\log\frac{\det A_{\chi,r}^B}{\det M_{q,r}^B}
                 +t_{A_{\chi,r}}-t_{M_{q,r}},\\
 rb_- -L_r\leq{}&\eta_r\leq rb_++L_r.
 \end{split}
 \tag{BSL10}
\]
For a further exact representation of the first term, put
$A_r(t)=\operatorname{Gram}_{1_B|y|^{2q}\exp(t\psi_\chi)d\sigma}
(1,S,\ldots,S^{r-1})$. Its derivative exists entrywise since
$\psi_\chi$ is bounded on $B$, and all its matrices are positive.
Jacobi's determinant derivative, proved by differentiating the
multilinear determinant one column at a time, gives
\[
 \log\frac{\det A_{\chi,r}^B}{\det M_{q,r}^B}
    =\int_0^1\operatorname{tr}(A_r(t)^{-1}A_r'(t))\,dt.
 \tag{BSL11}
\]
Thus the entire polynomial location contribution has both its original
finite matrix and its exact interpolation map.

\subsection{The evaluated baseline coefficient}
Let $\mu_q=\int y^{2q}\,d\sigma(y)$, let $H_j^{(a)}$ denote the
positive half-line moment determinant GEL1, and retain
\[
 Z_q=H_n^{(q)}H_{n+1}^{(q)}(H_n^{(q+1)})^2,
 \qquad V_q=S_q+\log Z_q-\log\mu_q.
 \tag{BSL12}
\]
Grouping even and odd coordinate indices in $M_{q,q}$ gives
$\det M_{q,q}=H_n^{(q)}H_n^{(q+1)}$.
The identical permutation in $M_{q,q+1}$ gives
$\det M_{q,q+1}=H_{n+1}^{(q)}H_n^{(q+1)}$.
The odd cross entries vanish by evenness of $\sigma$; the simultaneous
row and column permutations have product sign one. Also
$\det M_{q,1}=\mu_q$. Thus BSL4 becomes exactly
\[
 \boxed{\mathcal B_k^\sigma=V_q+\eta_q+\eta_{q+1}-\eta_1.}
 \tag{BSL13}
\]
Writing $L_*=L_q+L_{q+1}+L_1$, BSL10 retains the full finite bounds
\[
 \boxed{V_q+(2q+1)b_- -b_+-L_*
       \leq\mathcal B_k^\sigma
       \leq V_q+(2q+1)b_+-b_-+L_*.}
 \tag{BSL14}
\]
In particular no low determinant is silently assigned the value one.

The independent homogeneous calculation HBL gives the complete result
\[
 \frac{V_q}{q^2}\longrightarrow C_B
 :=9-8\log2+\mathcal F(2,\pi),\qquad
 C_B>\frac{769}{17010}>0,
 \tag{BSL15}
\]
where $\mathcal F$ is the exact variational energy GEL7 with the
unique positive density EIQ11. All logarithmic source powers cancel
only after the identity BSL12 is substituted: the source contribution
is $-6q^2\log q+(9-8\log2)q^2+O(q\log q)$, while
$\log Z_q=(6q^2+2q)\log q+q^2\mathcal F(2,\pi)+o(q^2)$.
The low term has size $O(q\log q)$. HBL proves the exact coefficient,
its one dimensional equilibrium integral, and the strict rational
lower bound by an explicit arcsine competitor. Its full proof is
included with this section.

The original degree identities give
\[
 |a_k|\leq\frac{|d-g|}{3e},\qquad
 d_\chi\leq\frac{2(d+g)^2}{3e^3k^2\epsilon^4}.
\]
Consequently $\eta_q+\eta_{q+1}-\eta_1=O(q/e)$ by BSL10;
the displayed tails decay exponentially. Since $q/e=(k+1)^2$
and $q\to\infty$, this error divided by $q^2$ tends to zero.
The stronger signed bounds BSL14 remain in force. We have proved
the original packet limit
\[
 \boxed{\frac{\mathcal B_k^\sigma}{q^2}\longrightarrow C_B>0.}
 \tag{BSL16}
\]

\subsection{Original tensor, arithmetic, kernel, and boundary receivers}
The tensor-Gamma source and its mass are unchanged:
\[
 dm_{0,k}(y)=\beta_k|f_l(c+iy)|^2d\sigma(y),\quad
 f_l(S)=\prod_{j=0}^{l-1}(S-c-2j-\tfrac12),\quad
 \beta_k=\frac{(2\pi)^{k/2}}{\sqrt2\Gamma(k/2)}.
 \tag{BSL17}
\]
Its full quotient baseline is the original $\mathcal B_k^0$.
Multiplication by $f_l$ in the original remainder fibres, including
the full product Taylor jets and their invertible multiplication germ,
gives TWA29--30 and the complete GCR proof. In their signs these exact
identities are
\[
 \mathcal B_k^0=\mathcal B_k^\sigma+\mathfrak T_k,
 \quad\mathcal B_k^{\rm ar}=\mathcal B_k^\sigma+\delta_k^\sigma,
 \quad\Delta_k^\Gamma=\delta_k^\sigma-\mathfrak T_k.
 \tag{BSL18}
\]
All four contributions $q\log\beta_k$ and all four full resultant
logarithms in that proof are retained before their signed cancellation.
Here $\mathfrak T_k/(lq)\to\mathcal C_\Gamma>0$ is the proved
GEL19 result and $\delta_k^\sigma/(kq)\to0$ is the proved original
source determinant estimate TWA15. Since $l/q\to0$ and $k/q\to0$,
BSL16--18 give
\[
 \boxed{\frac{\mathcal B_k^0}{q^2}\longrightarrow C_B,\qquad
 \frac{\mathcal B_k^{\rm ar}}{q^2}\longrightarrow C_B,\qquad
 \frac{\mathcal R_k}{q^2}\longrightarrow C_B,}
 \qquad\mathcal R_k=\mathcal B_k^{\rm ar}-4q\log(D_hk).
 \tag{BSL19}
\]
The last limit uses the original fixed $D_h>0$ and
$\log(D_hk)/q\to0$. The full finite enclosures are obtained without
any asymptotic substitution: add the original interval for
$\mathfrak T_k$ to BSL14 for $\mathcal B_k^0$, add
$[-a_k^-,a_k^+]$ to BSL14 for $\mathcal B_k^{\rm ar}$, and subtract
the original $4q\log(D_hk)$ for $\mathcal R_k$.

To see exactly which mixed objects receive this result, retain
$\Lambda:E_\chi\to B_{\rm obs}=\operatorname{im}\Lambda$,
$K=\ker\Lambda$, its inclusion $I_K$, and the original section
$S_*:B_{\rm obs}\to E_\chi$ with $\Lambda S_*=I$.
For any original positive quotient Gram $G_N$, the square matrix
$[I_K,S_*]$ is invertible: every $z$ is uniquely
$(z-S_*\Lambda z)+S_*\Lambda z$ with the first term in $K$.
Its restricted kernel Gram is $I_K^*G_NI_K$; the boundary minimum
Gram is the Schur complement after this restriction. Block elimination
therefore proves, with the original frame factor,
\[
 \det(I_K^*G_NI_K)\det Q_N^{B_{\rm obs}}
       =|\det[I_K,S_*]|^2\det G_N.
 \tag{BSL20}
\]
Both determinant factors are positive, with the empty one equal to one
at boundary rank zero or full rank. In the same four endpoint signed
sum the frame contribution has coefficient $1+1-1-1=0$.
The complete BRU row update construction gives these two sums as
$F_K(1)$ and $F_B(1)$ and hence
\[
 \frac{F_K(1)+F_B(1)}{q^2}\longrightarrow C_B.
 \tag{BSL21}
\]
This proves the exact combined mixed-support consequence. It keeps
the individual row contractions, the full Taylor unit in $\Lambda$,
and their off-diagonal terms through BSL20; no individual contribution
is assigned the combined limit.

The same original HC identity TWA10 is
$\mathcal B_k^{\rm ar}=4q\log\mathcal L_k-W_k^\sigma
 +\mathcal S_k-\tau_k^{\rm win}$, where all nonnegative contraction,
phase, and relative-control terms occur separately in TWA9.
Here $\mathcal L_k=2\delta e(k+1)\lfloor(k+1)^2/4\rfloor$,
$W_k^\sigma=O(q\log q)$ by its literal factorial formula TWA5,
and $\tau_k^{\rm win}=O_h(k+\log q)$ by TWA6.
The explicit polynomial formula for $q$ gives
$\log\mathcal L_k=O_h(\log k)$ and $\log q=O_m(\log k)$.
Dividing this exact HC identity by $q^2$ now proves
\[
 \boxed{\frac{\mathcal S_k}{q^2}\longrightarrow C_B>0.}
 \tag{BSL22}
\]
Thus the growing source volume is connected all the way to the
original contraction and phase costs. The finer Gamma law
$(\mathcal B_k^0-\mathcal B_k^{\rm ar})/(kq)
 \to\mathcal C_\Gamma/4$ continues to hold simultaneously.
Its complete finite remainders and all packet locations remain
in BSL10--14 and WGR; none is removed by the leading $q^2$ limit.

% END RX INLINE 38


% BEGIN RX INLINE 39; SHA256 642fc0b0b60a3253ad577b588b2b93d6d0aa4357cae40b69e7f45098e0c10304
% New proof fragment; no historical source is modified.
% Derived directly from the original Gamma series in recurrence_bounds.tex,
% RWB3--4. The derivation of that series is reproduced below.
% Scope: explicit Gamma logarithmic-derivative refinement only.
\section{An explicit remainder for the original Gamma logarithmic derivative}
\label{sec:GDR}

Retain the original density and its logarithmic derivative:
\[
 \sigma(y)=\frac{|\Gamma(1/4+iy/2)|^2}{2\pi},\qquad
 L(y)=-\frac{\sigma'(y)}{\sigma(y)},\qquad y\in\mathbb R.
 \tag{GDR1}
\]
For every real $y\ne0$, define
\[
 \epsilon(y)=yL(y)-\frac\pi2|y|-\frac12.
 \tag{GDR2}
\]
The following two estimates hold on this entire punctured real line:
\[
 \boxed{\displaystyle
 |\epsilon(y)|\le
 \left(\frac18+\frac{\sqrt3}{9}\right)\frac1{y^2}
 <\frac1{3y^2},}
 \tag{GDR3}
\]
\[
 \boxed{\displaystyle
 \left|\epsilon(y)-\frac1{8y^2}\right|
 \le\left(\frac5{32}+\frac7{12\sqrt3}\right)\frac1{y^4}
 <\frac1{2y^4}.}
 \tag{GDR4}
\]
In particular, the coefficient of the inverse-square term is positive
and exactly $1/8$. Both bounds are global inequalities with explicit
constants; no unspecified asymptotic remainder is used. The density
and $yL(y)$ have their original values at zero, namely
$\sigma(0)=|\Gamma(1/4)|^2/(2\pi)$ and $yL(y)=0$ at $y=0$.
Equations GDR3--4 make no assertion at that excluded point.

\paragraph{The exact convergent series.}
For $\Re z>0$, repeated integration by parts in the beta integral gives
\[
 G_N(z):=\frac{N!N^z}{z(z+1)\cdots(z+N)}
 =\int_0^N u^{z-1}(1-u/N)^N\,du.
 \tag{GDR5}
\]
The factor $(1-u/N)^N$ is interpreted as zero above $N$ and is
bounded by $e^{-u}$ on $[0,N]$. On a compact subset of $\Re z>0$,
use its least real part below $u=1$ and its greatest real part above
$u=1$. The resulting integrable majorants, also multiplied by
$|\log u|$, prove locally uniform convergence of both $G_N$ and
$G_N'$ to $\Gamma$ and $\Gamma'$. For $\alpha>0$ and $v\in\mathbb R$,
\[
 \frac{|\Gamma(\alpha+iv)|}{\Gamma(\alpha)}
 =\lim_{N\to\infty}\prod_{r=0}^N
 \left(1+\frac{v^2}{(r+\alpha)^2}\right)^{-1/2}>0.
 \tag{GDR6}
\]
The strict positivity follows because the sum of the logarithms of
the positive factors converges absolutely, by comparison with
$\sum_r v^2/(r+\alpha)^2$. Logarithmic differentiation of GDR5,
taking its imaginary part and then the limit, therefore proves
\[
 \operatorname{Im}\frac{\Gamma'(\alpha+iv)}{\Gamma(\alpha+iv)}
 =\sum_{r=0}^{\infty}\frac{v}{(r+\alpha)^2+v^2}.
 \tag{GDR7}
\]
Indeed, the real term $\log N$ contributes zero imaginary part.
Differentiating the two original conjugate Gamma factors in GDR1
shows that $L(y)$ equals the left side of GDR7 with
$\alpha=1/4$ and $v=y/2$. Consequently, for $y=2v>0$,
\[
 yL(y)=2\sum_{r=0}^{\infty}F_v(r+1/4),\qquad
 F_v(x)=\frac{v^2}{x^2+v^2},\quad x\ge0.
 \tag{GDR8}
\]
In particular the original quarter-shift has been retained exactly.

\paragraph{Shifted summation with its signed remainder.}
We prove the needed summation identity instead of invoking an
asymptotic Gamma expansion. Define the following polynomials on
$[0,1]$:
\[
\begin{aligned}
 B_1(t)&=t-\tfrac12,\qquad
 B_2(t)=t^2-t+\tfrac16,\\
 B_3(t)&=t^3-\tfrac32t^2+\tfrac12t,\\
 B_4(t)&=t^4-2t^3+t^2-\tfrac1{30},\\
 B_5(t)&=t^5-\tfrac52t^4+\tfrac53t^3-\tfrac16t.
\end{aligned}
 \tag{GDR9}
\]
Direct differentiation and evaluation give
$B_j'=jB_{j-1}$ for $2\le j\le5$,
$B_j(1)=B_j(0)$ for $2\le j\le5$, and
$B_j(1-t)=(-1)^jB_j(t)$ for $1\le j\le5$.
For $0<\alpha<1$, let
\[
 P_j(x)=\frac{B_j(\{x-\alpha\})}{j!},\qquad x\ge0,
 \tag{GDR10}
\]
where $\{u\}=u-\lfloor u\rfloor$. The functions $P_j$ are
bounded; those with $j\ge2$ are continuous. Between successive
points $r+\alpha$, $P_1'=1$, and at every such point $P_1$ has
a downward jump of size one. For a continuously differentiable
function $f$ for which the expressions converge, integration on
each interval and addition of the jump contributions give
\[
 \sum_{r\ge0}f(r+\alpha)
 =\int_0^\infty f(x)\,dx
   +f(0)P_1(0)+\int_0^\infty f'(x)P_1(x)\,dx.
 \tag{GDR11}
\]
One can first make this calculation on $[0,M]$, where $M$ is an
integer. The right boundary term is $-f(M)P_1(M)$ and tends to zero
when $f(M)\to0$. The left boundary is not a jump, since
$0<\alpha<1$, and
$P_j(0)=B_j(1-\alpha)/j!=(-1)^jB_j(\alpha)/j!$.
The derivative identity for $P_j$ holds almost everywhere and
$P_j$ is locally absolutely continuous for $j\ge2$. Repeated
ordinary integration by parts in the last integral of GDR11
therefore yields, for either $p=3$ or $p=5$,
\[
 \sum_{r\ge0}f(r+\alpha)
 =\int_0^\infty f(x)\,dx
  -\sum_{j=1}^{p}\frac{B_j(\alpha)}{j!}f^{(j-1)}(0)
  +\frac{(-1)^{p-1}}{p!}
   \int_0^\infty B_p(\{x-\alpha\})f^{(p)}(x)\,dx.
 \tag{GDR12}
\]
For the actual function $f=F_v$, every derivative in question
tends to zero at infinity and is integrable on $[0,\infty)$.
Thus every boundary passage and integral in this proof is justified
for the function to which GDR12 will be applied. In particular,
the remainder sign is positive for both odd values of $p$ used here.

\paragraph{The reciprocal-square bound.}
Put $H(u)=(1+u^2)^{-1}$, so that
$F_v^{(r)}(x)=v^{-r}H^{(r)}(x/v)$. The necessary derivatives are
\[
 H''(u)=\frac{6u^2-2}{(1+u^2)^3},\qquad
 H'''(u)=\frac{24u(1-u^2)}{(1+u^2)^4}.
 \tag{GDR13}
\]
The second derivative increases from $-2$ to $1/2$ on $[0,1]$
and decreases from $1/2$ to $0$ on $[1,\infty)$. It follows exactly
that
\[
 \int_0^\infty|F_v'''(x)|\,dx=\frac3{v^2}.
 \tag{GDR14}
\]
Writing $s=t-1/2$ gives $B_3(t)=s(s^2-1/4)$. On
$-1/2\le s\le1/2$, differentiating this cubic shows that its
largest absolute value is attained at $s=\pm1/(2\sqrt3)$ and equals
\[
 \max_{0\le t\le1}|B_3(t)|=\frac1{12\sqrt3}
 =\frac{\sqrt3}{36}.
 \tag{GDR15}
\]
At the original shift $\alpha=1/4$,
$B_1(1/4)=-1/4$, $B_3(1/4)=3/64$, while
$F_v'(0)=0$ and $F_v''(0)=-2/v^2$. Since
$\int_0^\infty F_v(x)\,dx=\pi v/2$, GDR12 with $p=3$ becomes
\[
 \sum_{r\ge0}F_v(r+1/4)
 =\frac{\pi v}{2}+\frac14+\frac1{64v^2}+R_3(v),
 \qquad
 |R_3(v)|\le\frac{\sqrt3}{72v^2}.
 \tag{GDR16}
\]
Multiply this equality by two and put $v=y/2$ as in GDR8.
This proves the non-strict inequality in GDR3 for $y>0$.
The stated strict comparison of constants is equivalent to
$\sqrt3/9<5/24$, whose square is $64<75$ after multiplying
by the positive common factors.

\paragraph{The inverse-fourth remainder after the exact coefficient.}
For the same $H$, direct differentiation gives
\[
 H^{(4)}(u)=\frac{24(1-10u^2+5u^4)}{(1+u^2)^5},\qquad
 H^{(5)}(u)=\frac{-240u(3-10u^2+3u^4)}{(1+u^2)^6}.
 \tag{GDR17}
\]
The fourth derivative decreases on $[0,1/\sqrt3]$, increases on
$[1/\sqrt3,\sqrt3]$, and decreases on $[\sqrt3,\infty)$.
Its endpoint values in that order are
\[
 24,\qquad-\frac{81}{8},\qquad\frac38,\qquad0.
\]
Its total variation is therefore
$(24+81/8)+(81/8+3/8)+3/8=45$, and scaling back to $F_v$ yields
\[
 \int_0^\infty|F_v^{(5)}(x)|\,dx=\frac{45}{v^4}.
 \tag{GDR18}
\]
The exact factorization
\[
 B_5(t)=t(t-1)(t-1/2)(t^2-t-1/3)
       =s(s^2-1/4)(s^2-7/12),\qquad s=t-1/2,
 \tag{GDR19}
\]
implies, using the maximization already proved in GDR15,
\[
 |B_5(t)|
 =|s|(1/4-s^2)(7/12-s^2)
 \le\frac7{12}\,\frac1{12\sqrt3}
 =\frac7{144\sqrt3},\qquad 0\le t\le1.
 \tag{GDR20}
\]
The remaining original-shift values are
$B_5(1/4)=-25/1024$, $F_v'''(0)=0$, and
$F_v^{(4)}(0)=24/v^4$. Thus the exact $p=5$ equality from GDR12 is
\[
\begin{split}
 \sum_{r\ge0}F_v(r+1/4)
 &=\frac{\pi v}{2}+\frac14+\frac1{64v^2}
       +\frac5{1024v^4}+R_5(v),\\
 R_5(v)&=\frac1{120}\int_0^\infty
 B_5(\{x-1/4\})F_v^{(5)}(x)\,dx,\qquad
 |R_5(v)|\le\frac7{384\sqrt3\,v^4}.
\end{split}
 \tag{GDR21}
\]
Equivalently, the exact signed remainder in the original coordinate is
\[
 \epsilon(y)=\frac1{8y^2}+\frac5{32y^4}
 +\frac1{60}\int_0^\infty
  B_5(\{x-1/4\})F_{y/2}^{(5)}(x)\,dx,
 \qquad y>0.
 \tag{GDR22}
\]
The last integral, including its factor $1/60$, has absolute value
at most $7/(12\sqrt3\,y^4)$. This proves the non-strict inequality
in GDR4. For its strict comparison of constants, observe that
$7/(12\sqrt3)<11/32$ is equivalent to $56<33\sqrt3$; squaring
both positive sides gives $3136<3267$.

Finally, $\Gamma(1/4-iy/2)$ is the conjugate of
$\Gamma(1/4+iy/2)$ by the original defining integral. Hence $\sigma$
is even, $L$ is odd, and $yL(y)$ is even. This proves GDR3--4
on the full stated domain $y\in\mathbb R\setminus\{0\}$.
Every occurrence of $y/2$, the quarter-shift, the original density
constant $1/(2\pi)$, and the remainder orientation is retained.

% END RX INLINE 39


% BEGIN RX INLINE 40; SHA256 b3607ad6e0cd2480f917bdce057683cb8526d5243f583363eee297835403422e
% Complete new low-endpoint calculation. Historical sources are unchanged.
\section{The original scalar Gamma endpoint to fifth-order finite error}
\label{sec:LER}

Retain the original density and both literal moment indices:
\[
 \sigma(y)=\frac{|\Gamma(1/4+iy/2)|^2}{2\pi},\qquad
 I_p=\int_0^\infty y^p\sigma(y)\,dy\quad(p\ge0),\qquad
 \mu_a=\int_{\mathbb R}y^{2a}\sigma(y)\,dy=2I_{2a}.
 \tag{LER1}
\]
Thus $\mu_a=m_{2a}$ in LET's all-power notation, and
$\mu_0=\sqrt{2\pi}$. The factors two and the original mass are kept
until cancellation in the specified ratios. Let $\alpha_\Gamma=\pi/2$ and
$L(y)=-\sigma'(y)/\sigma(y)$. The complete accompanying derivation
\texttt{DERIVATIVE\_REMAINDER.tex}, based on the original Gamma
beta-limit series, proves
\[
 yL(y)=\alpha_\Gamma y+\frac12+\frac1{8y^2}+E(y),\qquad
 |E(y)|\le\frac1{2y^4}\quad(y>0).
 \tag{LER2}
\]
No asymptotic notation substitutes for this global finite estimate.

\paragraph{Exact moment recursion and its positive comparison data.}
For every integer $p\ge4$, integrate the derivative of
$y^{p+1}\sigma(y)$ on $(0,\infty)$. Its boundary values are zero
because the density is bounded at zero and has the exponential bound
proved in HCT1. The differentiated integrand is integrable by RWB2.
Substitution of LER2 gives the exact identity
\[
 \alpha_\Gamma I_{p+1}=(p+1/2)I_p-\frac18 I_{p-2}-e_p,
 \qquad e_p=\int_0^\infty y^p E(y)\sigma(y)\,dy,
 \qquad |e_p|\le\frac12 I_{p-4}.
 \tag{LER3}
\]
The lower powers on the right are integrable at zero for $p\ge4$.
We retain the sign of $e_p$; the bound on its absolute value is the
only replacement made in an inequality.

RWB11 already gives, with its complete proof from the original density,
\[
 \frac{2v-1}{\pi}\le\frac{I_{v+1}}{I_v}
 \le\frac{2v+3}{\pi}\quad(v\ge0).
 \tag{LER4}
\]
The lower bound is used only when it is positive. For $p\ge10$ define
\[
 U_{p,2}=\frac{\pi^2}{(2p-5)(2p-3)},\qquad
 U_{p,4}=\frac{\pi^4}{(2p-9)(2p-7)(2p-5)(2p-3)},
\]
\[
 b_p=\frac{U_{p,2}/8+U_{p,4}/2}{p+1/2}.
 \tag{LER5}
\]
Telescoping the respective two or four positive lower bounds in LER4
proves $I_{p-2}/I_p\le U_{p,2}$ and
$I_{p-4}/I_p\le U_{p,4}$.
Put
\[
 A_p=\frac{p+1/2}{\alpha_\Gamma},\qquad
 u_p=\frac{I_{p-2}/(8I_p)+e_p/I_p}{p+1/2}.
\]
Then LER3 gives the exact positive ratio
\[
 \frac{I_{p+1}}{I_p}=A_p(1-u_p),\qquad
 |u_p|\le b_p\le\frac7{4p^3}<\frac1{100}\quad(p\ge10).
 \tag{LER6}
\]
For the numerical estimate, $\pi<22/7$ implies $\pi^2<10$ and
$\pi^4<100$. Every factor $2p-9,2p-7,2p-5,2p-3$ is at least $p$
when $p\ge10$. Hence
$b_p\le5/(4p^3)+50/p^5\le7/(4p^3)$, proving every stated constant.

\paragraph{The logarithmic step with a fifth-order remainder.}
For $p\ge12$ define
\[
 V_p=\frac{\alpha_\Gamma^2}{(p-3/2)(p-1/2)},\qquad
 z_p=\frac{V_p}{8(p+1/2)}
     =\frac{\pi^2}{32(p-3/2)(p-1/2)(p+1/2)}.
 \tag{LER7}
\]
There is an exact real $\eta_p$ satisfying
\[
 \boxed{\displaystyle
 \log\frac{I_{p+1}}{I_p}
 =\log\frac{p+1/2}{\alpha_\Gamma}-z_p+\eta_p,
 \qquad |\eta_p|\le\frac{51}{p^5}.}
 \tag{LER8}
\]
To prove it without assuming a moment expansion, multiply the two
identities LER6 at $p-2,p-1$ and invert their positive product:
\[
 \frac{I_{p-2}}{I_p}
 =\frac{V_p}{(1-u_{p-2})(1-u_{p-1})}.
 \tag{LER9}
\]
For $|u|\le b\le1/100$, $|v|\le c\le1/100$, direct subtraction gives
\[
 \left|\frac1{(1-u)(1-v)}-1\right|
 \le\frac{b+c+bc}{(1-b)(1-c)}\le2(b+c).
\]
Indeed $bc\le(b+c)/200$ and $(1-b)(1-c)\ge9801/10000$;
$10050/9801<2$ proves the last bound. LER6 and LER9 thus imply
\[
 \left|\frac{I_{p-2}}{I_p}-V_p\right|
 \le2V_p(b_{p-2}+b_{p-1})\le\frac{7V_p}{(p-2)^3}.
 \tag{LER10}
\]
Consequently the exact $u_p$ satisfies
\[
 |u_p-z_p|\le
 \frac{7V_p/[8(p-2)^3]+U_{p,4}/2}{p+1/2}.
 \tag{LER11}
\]
For $|u|<1$, integration of the derivative of $\log(1-u)+u$ gives
\[
 |\log(1-u)+u|\le\frac{u^2}{2(1-|u|)}.
\]
Combining this with LER6 and LER11 bounds $|\eta_p|$ by
\[
 \frac{7V_p/[8(p-2)^3]+U_{p,4}/2}{p+1/2}
 +\frac{b_p^2}{2(1-b_p)}.
 \tag{LER12}
\]
Here are explicit estimates for the advertised simpler constant.
For $p\ge12$, $p-3/2\ge7p/8$, $p-1/2\ge23p/24$, and $\pi^2<10$
give $V_p<3/p^2$. Also $p-2\ge5p/6$. The first term in LER12 is
therefore bounded by
\[
 \frac{21}{8}\left(\frac65\right)^3p^{-6}+50p^{-5}
 <5p^{-6}+50p^{-5}.
\]
Using LER6, the second term is at most
$[49/(16p^6)](50/99)<2p^{-6}$.
Thus LER12 is less than
$50p^{-5}+7p^{-6}<51p^{-5}$, which proves LER8.

\paragraph{The original scalar endpoint and an exact telescoping correction.}
For integers $a\ge6$ and $l\ge1$ define
\[
 B(x)=\frac1{(x-3/2)(x-1/2)}\quad(x>3/2),\qquad
 K_{a,l}=\frac{\pi^2}{64}\bigl[B(2a)-B(2a+2l)\bigr]>0.
 \tag{LER13}
\]
Summing LER8 at every original integer $p=2a,\ldots,2a+2l-1$ gives
\[
 \boxed{\displaystyle
 \log\frac{\mu_{a+l}}{\mu_a}
 =\sum_{j=0}^{2l-1}\log\frac{2a+j+1/2}{\pi/2}
       -K_{a,l}+\varepsilon_{a,l},\qquad
 |\varepsilon_{a,l}|\le\frac{51l}{16a^5}.}
 \tag{LER14}
\]
In particular no square-pushforward index is substituted for an
all-power moment index: $\mu_{a+l}/\mu_a=I_{2a+2l}/I_{2a}$
is precisely the telescoped ratio. The rational term telescopes
because its exact partial fraction is
\[
 z_p=\frac{\pi^2}{64}\bigl[B(p)-B(p+1)\bigr].
 \tag{LER15}
\]
The sum of the error bounds in LER8 is at most
$51(2l)/(2a)^5=51l/(16a^5)$.
All constants in LER14 are therefore uniform over $a=q,\ldots,q+l$,
as well as over every other integer $a\ge6$.

At $l=1$, this gives the requested individual Christoffel endpoint:
\[
 \boxed{\displaystyle
 \log c_0^{(a)}
 =\log\frac{(4a+1)(4a+3)}{\pi^2}
 -\frac{\pi^2}{64}\bigl[B(2a)-B(2a+2)\bigr]
 +\varepsilon_{a,1},\qquad
 |\varepsilon_{a,1}|\le\frac{51}{16a^5}.}
 \tag{LER16}
\]
The full original low moment remains the product of these coefficients;
the corrections in LER16 add to the single term in LER14, with no
loss from bounding the exact telescoping contribution.

\paragraph{Explicit finite expansion along the original packet.}
Retain $k=4l+1$, $e=1+k(m-1)$ and
$q=e(k+1)^2=e(4l+2)^2$, with $l,m\ge1$. Thus $q\ge36$.
Set
\[
 F_{q,l}=2l\log\frac{4q}{\pi}
 +\frac{l^2}{q}-\frac{l(16l^2-1)}{48q^2}
 +\frac{l^2(8l^2-1)}{48q^3},
\]
\[
 T_{q,l}=\frac{l(768l^4-160l^2+7)}{7680q^4}>0,
 \qquad E_{q,l}=\frac{51l}{16q^5}.
 \tag{LER17}
\]
There is an exact $r_{q,l}$ with $-T_{q,l}\le r_{q,l}\le0$ such that
\[
 \boxed{\displaystyle
 \log\frac{\mu_{q+l}}{\mu_q}
 =F_{q,l}-K_{q,l}+r_{q,l}+\varepsilon_{q,l},\qquad
 |\varepsilon_{q,l}|\le E_{q,l}.}
 \tag{LER18}
\]
Indeed for every $x\ge0$,
\[
 \log(1+x)=x-\frac{x^2}{2}+\frac{x^3}{3}
              -\int_0^x\frac{t^3}{1+t}\,dt,
\]
and the last term belongs to $[-x^4/4,0]$.
Apply this exact formula to $x_j=(j+1/2)/(2q)$, $0\le j<2l$,
in LER14. The required literal power sums, for $M=2l$, are
\[
 \sum_{j=0}^{M-1}(j+1/2)=\frac{M^2}{2},\qquad
 \sum_{j=0}^{M-1}(j+1/2)^2=\frac{M(4M^2-1)}{12},
\]
\[
 \sum_{j=0}^{M-1}(j+1/2)^3=\frac{M^2(2M^2-1)}8,\qquad
 \sum_{j=0}^{M-1}(j+1/2)^4
 =\frac{M(48M^4-40M^2+7)}{240}.
 \tag{LER19}
\]
Each identity follows by expanding the left-hand summand and summing
integer powers; alternatively its right-hand difference at $M+1,M$
equals $(M+1/2)^r$ and its value at $M=0$ is zero, proving it by
induction. Substitution proves every coefficient and the sign in LER18.
The additional rational correction has the elementary bound
\[
 0<K_{q,l}=\sum_{p=2q}^{2q+2l-1}z_p
 \le\frac{3l}{32q^3},
 \tag{LER20}
\]
since the established $V_p<3/p^2$ gives $z_p<3/(8p^3)$.

\paragraph{Exact insertion into the entire Gamma functional.}
Keep the original paired high determinant
$Z_a=H_n^{(a)}H_{n+1}^{(a)}(H_n^{(a+1)})^2$ with $n=q/2$, and
the original product term $P_k$ from LET19. Write
\[
 \mathcal V_k^{\mathrm{high}}=P_k+\log\frac{Z_{q+l}}{Z_q},\qquad
 W_k=\mathcal V_k^{\mathrm{high}}-\log\frac{\mu_{q+l}}{\mu_q}.
 \tag{LER21}
\]
These are the literal original determinant terms, with all four
high blocks retained. LER18 proves the finite receiving interval
\[
 \boxed{\displaystyle
 \mathcal V_k^{\mathrm{high}}-F_{q,l}+K_{q,l}-E_{q,l}
 \ \le W_k\le\
 \mathcal V_k^{\mathrm{high}}-F_{q,l}+K_{q,l}+T_{q,l}+E_{q,l}.}
 \tag{LER22}
\]
In particular the Gamma rational correction is positive in $W_k$,
while the $l^2/q$ term is negative there. The one-sided Taylor error
is also positive after the original low-endpoint minus sign.
This formula updates the scalar contribution only; it leaves the
original high determinants, quotient maps, arithmetic contribution
and their stated bounds intact for their own calculations.

Finally LER17--20 prove, along the actual packet, the more precise
low-endpoint limit
\[
 \log\frac{\mu_{q+l}}{\mu_q}-2l\log\frac{4q}{\pi}
 -\frac{l^2}{q}\longrightarrow0.
 \tag{LER23}
\]
To check this directly, $q\ge(4l+2)^2$ bounds each remaining term by
a constant times respectively $l^{-1},l^{-2},l^{-3},l^{-5},l^{-9}$
as $l\to\infty$, with the displayed exact finite bounds available
in place of these orders. For $m=1$, $l^2/q\to1/16$; for every
fixed $m\ge2$, $l^2/q\to0$ since
$q=[1+(4l+1)(m-1)](4l+2)^2$ grows cubically in $l$.
These conclusions concern the exact original scalar moment ratio;
the complete $W_k$ continues to include both original high ranks.

% END RX INLINE 40


% BEGIN RX INLINE 41; SHA256 52256dd3671db7944c6758ed8a2f7129b7a226a08c433b1f32449985f2137b71
% Complete finite-n proof, 14 September 2026.
% Provider: equilibrium/independent/EXACT_SQRT_LOG_EQUILIBRIUM.tex, EIQ1--33.
% The Heine factor 1/n! and both reference measures are retained explicitly.
\section{A uniform finite partition estimate from the original equilibrium}
\label{sec:QLG}

\paragraph{QLG1. Exact measures and parameter domain.}
Fix the compact parameter rectangle
\[
 \mathcal K=[3/2,5/2]\times[\pi/2,3\pi/2],\qquad
 V_{\alpha,\beta}(x)=\beta\sqrt{x}-\alpha\log x\quad(x>0).
 \tag{QLG1}
\]
The equilibrium proved in EIQ1--33 has density
\[
 \rho(x)=\frac{\beta\sqrt{(b-x)(x-a)}}{4\pi^2x}
       \int_0^\infty
       \frac{\sqrt{s}\,ds}{(x+s)\sqrt{(a+s)(b+s)}}\quad(a<x<b),
 \qquad \rho(x)=0\quad(x\notin(a,b)),
 \tag{QLG2}
\]
where the exact endpoints are those of EIQ3--5. In particular
\(0<a<b<\infty\), the density is bounded and continuous, positive
in the interior, and has integral one. Write
\[
 d\mu(x)=\rho(x)\,dx,\quad
 P(z)=\int_a^b\log|z-t|\,d\mu(t)\quad(z\in\mathbb C),\quad
 I(\mu)=\iint\log|x-y|\,d\mu(x)d\mu(y),
\]
\[
 F(\alpha,\beta)=I(\mu)-\int V_{\alpha,\beta}\,d\mu.
 \tag{QLG3}
\]
The original Euler equality and full exterior inequality in EIQ17--19
give a real number \(\ell\) for which
\[
 D(x)=V_{\alpha,\beta}(x)-2P(x)-\ell\geq0\quad(x>0),
 \qquad D(x)=0\quad(a\leq x\leq b).
 \tag{QLG4}
\]
Integrating the equality on the support gives the exact sign identities
\[
 \ell=\int V_{\alpha,\beta}\,d\mu-2I(\mu),\qquad
 F(\alpha,\beta)=-\ell-I(\mu).
 \tag{QLG5}
\]
The two integration measures treated below are, separately,
\[
 d\lambda_0(x)=dx,\qquad
 d\lambda_\omega(x)=d\omega(x)=\tfrac12 x^{-1/2}e^{-\sqrt{x}}\,dx.
 \tag{QLG6}
\]
The map \(x=r^2\) gives \(\int d\omega=\int_0^\infty e^{-r}dr=1\).
For either of these measures, with its stated density \(w_\lambda\),
define the literal finite integral
\[
 J_n^\lambda(\alpha,\beta)=\frac1{n!}
 \int_{(0,\infty)^n}
 \exp\!\left(2\sum_{i<j}\log|x_i-x_j|
                 -n\sum_{i=1}^n V_{\alpha,\beta}(x_i)\right)
       \prod_{i=1}^n d\lambda(x_i),\qquad n\geq1.
 \tag{QLG7}
\]
No source measure, scale, factorial, or exponent has been absorbed into
\(F\). The unweighted integral without the Heine factor is exactly
\(n!J_n^\lambda\).

\paragraph{QLG2. Constants proved finite on the whole rectangle.}
Let
\[
 a_* =\min_{(\alpha,\beta)\in\mathcal K} a(\alpha,\beta)>0,
 \qquad B_* =\max_{(\alpha,\beta)\in\mathcal K} b(\alpha,\beta)<\infty,
 \qquad \beta_-=\pi/2,\quad\beta_+=3\pi/2.
 \tag{QLG8}
\]
These are explicit compact-parameter extrema of the uniquely specified
endpoint integrals EIQ3--5. Their asserted strict inequalities follow
from the continuous positive endpoints proved in EIQ6 and compactness
of \(\mathcal K\). Thus their definition requires no existence or
regularity assumption about an unspecified minimizing measure.
For \(x\in[a,b]\), the integral in QLG2 satisfies
\[
 \int_0^\infty\frac{\sqrt{s}\,ds}{(x+s)\sqrt{(a+s)(b+s)}}
 \leq\int_0^\infty\frac{\sqrt{s}\,ds}{(a_*+s)^2}
 =\frac{\pi}{2\sqrt{a_*}}.
\]
For the last equality substitute \(s=a_*r^2\), followed by
\(r=\tan\theta\); the resulting integral is
\(2a_*^{-1/2}\int_0^{\pi/2}\sin^2\theta\,d\theta\).
Since \(\sqrt{(b-x)(x-a)}\leq(b-a)/2\leq B_*/2\), it follows that
\[
 0\leq\rho(x)\leq M_*:=\frac{\beta_+B_*}{16\pi a_*^{3/2}}
       \quad\hbox{for all real }x\hbox{ and all parameters in }\mathcal K.
 \tag{QLG9}
\]
In particular \(B_*M_*\geq1\), by the mass-one identity and support
length at most \(B_*\).

The entropy relative to either stated measure is
\[
 H_\lambda(\mu)=\int_a^b
                 \rho(x)\log\frac{\rho(x)}{w_\lambda(x)}\,dx,
       \qquad 0\log0=0.
 \tag{QLG10}
\]
It is absolutely integrable: the function \(u|\log u|\) is bounded
on \([0,M_*]\), and \(\log w_\lambda\) is bounded on the fixed
compact interval \([a_*,B_*]\). In the two cases its upper bounds are
\[
 H_{\lambda_0}(\mu)\leq\log M_* ,\qquad
 H_{\lambda_\omega}(\mu)\leq
 \log M_*+\log2+\tfrac12\log B_*+\sqrt{B_*}.
 \tag{QLG11}
\]
All pair logarithmic energies on this support are absolutely
integrable. For the negative part, integration against a density
bounded by \(M_*\) is bounded above by
\(M_*\int_{-1}^1-\log|s|\,ds=2M_*\), uniformly in the centre.
For the positive part the bound \(\log^+ B_*\) suffices.

Put
\[
 T_\lambda=\int_0^\infty e^{-D(x)}\,d\lambda(x).
 \tag{QLG12}
\]
For \(\lambda=\lambda_\omega\), QLG4 and the mass of \(\omega\)
give \(0<T_{\lambda_\omega}\leq1\). For the Lebesgue measure a finite
uniform bound can also be displayed entirely. For every \(x>0\),
\[
 P(x)\leq\log(1+x)+\log(1+B_*),\qquad P(x)\geq-2M_*.
\]
The lower bound uses only the negative part estimate just proved.
Evaluating QLG4 at any point of \([a,b]\) therefore yields
\[
 \ell\leq L_*:=\beta_+\sqrt{B_*}
    +\tfrac52\max\{|\log a_*|,|\log B_*|\}+4M_*.
\]
Consequently
\[
 T_{\lambda_0}\leq e^{L_*}(1+B_*)^2
    \int_0^\infty (1+x)^2(1+x^{5/2})e^{-\beta_-\sqrt{x}}\,dx
    =:T_*<\infty.
 \tag{QLG13}
\]
Here \(x^\alpha\leq1+x^{5/2}\) for all parameters in \(\mathcal K\).
The integral in QLG13 equals, by \(x=r^2\) and elementary repeated
integration by parts,
\[
 2\left(\frac{1!}{\beta_-^2}
       +\frac{2\,3!}{\beta_-^4}
       +\frac{5!}{\beta_-^6}
       +\frac{6!}{\beta_-^7}
       +\frac{2\,8!}{\beta_-^9}
       +\frac{10!}{\beta_-^{11}}\right).
 \tag{QLG14}
\]

\paragraph{QLG3. The planar logarithmic comparison, with its proof.}
We use the plane only to construct an exact comparison of logarithmic
kernels; the original particles and potential remain on \((0,\infty)\).
If \(\eta\) is a finite signed measure of compact support on
\(\mathbb R^2\), has total mass zero, and
\(\iint|\log|z-w||\,d|\eta|(z)d|\eta|(w)<\infty\), then
\[
 I(\eta):=\iint\log|z-w|\,d\eta(z)d\eta(w)\leq0.
 \tag{QLG15}
\]
To prove this, for \(d>0\) integrate the elementary identity
\[
 -\log d=\frac12\int_0^\infty
                    \frac{e^{-td^2}-e^{-t}}{t}\,dt.
\]
It follows by differentiation in \(d\) and evaluation at \(d=1\).
On any finite subinterval of \((0,\infty)\), the integral on the
right has the same sign as \(-\log d\) and absolute value at most
\(|\log d|\). The stipulated logarithmic integrability therefore
permits dominated passage to the full integral against \(d\eta^2\).
The constant term vanishes by the zero mass. The elementary Gaussian
Fourier integral in each of the two real coordinates gives
\[
 -I(\eta)=\frac12\int_0^\infty\frac{dt}{t}
       \iint e^{-t|z-w|^2}\,d\eta(z)d\eta(w),
\]
\[
 \iint e^{-t|z-w|^2}\,d\eta(z)d\eta(w)
 =\frac1{4\pi t}\int_{\mathbb R^2}
       e^{-|\xi|^2/(4t)}
       \left|\int e^{i\xi\cdot z}\,d\eta(z)\right|^2d\xi\geq0.
 \tag{QLG16}
\]
For completeness the scalar Fourier identity is obtained by
differentiating \(\int e^{-u^2/(4t)}e^{isu}du\) with respect to
\(s\), integrating by parts in \(u\), and using its value
\(\sqrt{4\pi t}\) at zero. The product of the two scalar identities
has the displayed factor \((4\pi t)^{-1}\). All finite Gaussian
integrations against the measures are absolutely convergent. This
proves QLG15 without requiring an equilibrium theorem on the plane.

Let \(c_{x,\varepsilon}\) be probability arc measure on the circle
\(z=x+\varepsilon e^{i\theta}\), where the real centre \(x\) is
arbitrary and \(\varepsilon>0\). Its exact logarithmic potential is
\[
 \int\log|z-w|\,dc_{x,\varepsilon}(z)
           =\log\max\{\varepsilon,|x-w|\}.
 \tag{QLG17}
\]
Indeed, if \(|x-w|>\varepsilon\), expand
\(\log|1-re^{i\theta}|=-\sum_{j\geq1}r^j\cos(j\theta)/j\)
with \(r=\varepsilon/|x-w|<1\); all nonconstant terms average to
zero. If \(|x-w|<\varepsilon\), use the same expansion with
\(r=|x-w|/\varepsilon\). At equality the assertion is the elementary
integral \(\int_0^{2\pi}\log|1-e^{i\theta}|d\theta=0\).
To verify this last integral directly, write
\(|1-e^{i\theta}|=2\sin(\theta/2)\), and put
\(A=\int_0^{\pi/2}\log(\sin u)du\). Reflection gives the same
integral for \(\log\cos u\); the double-angle identity then gives
\(2A=-(\pi/2)\log2+A\), whence \(A=-(\pi/2)\log2\).
The endpoint logarithms are integrable, for example from
\(2u/\pi\leq\sin u\leq u\) on \([0,\pi/2]\).

For finite centres \(x,y\), the logarithmic kernel is absolutely
integrable against \(c_{x,\varepsilon}\otimes c_{y,\varepsilon}\).
Its positive part is bounded by a finite constant depending on the
centres and \(\varepsilon\). Applying QLG17 first to this bounded
positive part minus the negative part, Tonelli's theorem and the finite
value \(\log\max\{\varepsilon,|x-w|\}\) show that the negative
part also has finite double integral. This argument includes identical
and tangent circles. The same argument proves absolute integrability
of the circle--\(\mu\) cross kernel, using compactness of the support
of \(\mu\); the \(\mu\)--\(\mu\) term was proved above.
In particular QLG15 applies to a finite mixture of these circles minus
\(\mu\).

Two consequences of QLG17, retaining both self and cross terms, are
\[
 I(c_{x,\varepsilon})=\log\varepsilon,\qquad
 I(c_{x,\varepsilon},c_{y,\varepsilon})
 \geq\log\max\{\varepsilon,|x-y|\}\geq\log|x-y|.
 \tag{QLG18}
\]
The first equality applies QLG17 at \(w\) on its own circle.
For the second inequality, the first circle average is
\(\log\max\{\varepsilon,|x-w|\}\geq\log|x-w|\);
averaging the right side on the second circle and applying QLG17 again
gives its asserted value. When \(x=y\) the last lower bound is
understood in the extended sense.

\paragraph{QLG4. A finite upper estimate on the original particle domain.}
For fixed particles \(x_1,\ldots,x_n>0\), put
\(\nu_\varepsilon=n^{-1}\sum_i c_{x_i,\varepsilon}\).
By QLG18 and QLG15, respectively,
\[
 2\sum_{i<j}\log|x_i-x_j|
       \leq n^2I(\nu_\varepsilon)-n\log\varepsilon,
 \qquad
 I(\nu_\varepsilon)\leq2I(\nu_\varepsilon,\mu)-I(\mu).
 \tag{QLG19}
\]
There is no potential evaluated at a complex particle in these
identities. The exact comparison back to the original real potential is
\[
 0\leq\int P(z)\,dc_{x,\varepsilon}(z)-P(x)
 =\int_{|t-x|<\varepsilon}
               \log\frac{\varepsilon}{|t-x|}\rho(t)\,dt
 \leq2M_*\varepsilon.
 \tag{QLG20}
\]
The last constant is the explicit integral
\(M_*\int_{-\varepsilon}^{\varepsilon}
\log(\varepsilon/|s|)ds=2M_*\varepsilon\).
Combining QLG19--20 and then substituting QLG4--5 proves, for every
original particle configuration,
\[
 2\sum_{i<j}\log|x_i-x_j|-n\sum_iV_{\alpha,\beta}(x_i)
 \leq n^2F-n\sum_iD(x_i)-n\log\varepsilon
                                      +4M_*n^2\varepsilon.
 \tag{QLG21}
\]
The coefficient \(4M_*\) comes from twice the logarithmic potential
and the bound \(2M_*\varepsilon\); the coefficient of the self energy
is \(-n\log\varepsilon\).
Take \(\varepsilon=1/n\). Since \(D\geq0\) and \(n\geq1\),
\(\int e^{-nD}d\lambda\leq T_\lambda\). Integrating QLG21 gives
the finite positive upper bound
\[
 \log J_n^\lambda
 \leq n^2F+4M_*n+n\log T_\lambda+n\log n-\log(n!).
 \tag{QLG22}
\]
In particular the original integrals in QLG7 are finite. The elementary
inequality
\(\log(n!)\geq\int_1^n\log t\,dt=n\log n-n+1\)
also proves
\[
 \log J_n^\lambda
 \leq n^2F+n(1+4M_*+\log T_\lambda)-1.
 \tag{QLG23}
\]
The inequality for \(n!\) follows by monotonicity of \(\log t\)
on each interval \([j-1,j]\), for \(j=2,\ldots,n\), and is exact
at \(n=1\).

\paragraph{QLG5. Quantile lower estimate and the exact factorial cancellation.}
Because \(\rho\) is continuous and positive on \((a,b)\), there are
unique endpoints
\[
 a=t_0<t_1<\cdots<t_n=b,\quad
 \int_{t_{j-1}}^{t_j}\rho(x)\,dx=\frac1n,\quad
 \Delta_j=t_j-t_{j-1}>0,
 \quad p_j(x)=n\rho(x)\mathbf1_{(t_{j-1},t_j)}(x).
 \tag{QLG24}
\]
The endpoints have measure zero. Each \(p_j(x)dx\) is a probability
measure with bounded density, on a compact interval of positive length.
Every pair logarithmic integral below is therefore absolutely
integrable, including self integrals and integrals over touching
intervals, by the same bounded-density estimate used after QLG11.
Also \(\int p_j|\log(p_j/w_\lambda)|dx<\infty\), since
\(p_j\leq nM_*\) and \(\log w_\lambda\) is bounded on the support.

The product cell \(\prod_{j=1}^n(t_{j-1},t_j)\) and its \(n!\)
coordinate permutations are disjoint. Symmetry of the integrand and the
product integration measure in QLG7 implies that the Heine factor
\(1/n!\) cancels these \(n!\) equal integrals exactly. Restricting
to their union therefore gives
\[
 J_n^\lambda\geq
 \int_{\prod_j(t_{j-1},t_j)}
 \exp\!\left(2\sum_{i<j}\log|x_i-x_j|-n\sum_iV(x_i)\right)
                       \prod_j w_\lambda(x_j)\,dx_j.
 \tag{QLG25}
\]
On this cell use the probability density \(\prod_jp_j(x_j)\).
The logarithm of the integrand relative to this probability is
integrable in absolute value, by the preceding estimates and boundedness
of \(V\) on \([a,b]\). The elementary convex inequality
\(e^u\geq e^v(1+u-v)\), integrated with
\(v\) equal to the expectation of this logarithm, proves Jensen's
lower bound directly. Since \(\sum_jp_j=n\rho\), its pair and
potential contributions are exactly
\[
 \sum_{i\ne j}I(p_i,p_j)-n\sum_j\int Vp_j\,dx
   =n^2F-\sum_j I(p_j).
\]
Here \(I(p_i,p_j)=\iint\log|x-y|p_i(x)p_j(y)dxdy\) and
\(I(p_j)=I(p_j,p_j)\). Its entropy contribution is exactly
\[
 \sum_j\int p_j\log\frac{p_j}{w_\lambda}\,dx
    =n\log n+nH_\lambda(\mu).
 \tag{QLG26}
\]
Moreover \(|x-y|\leq\Delta_j\) on the \(j\)-th cell, whence
\(I(p_j)\leq\log\Delta_j\). Concavity of the logarithm, proved
for example by its negative second derivative, and
\(\sum_j\Delta_j=b-a\) now give
\[
 \sum_jI(p_j)\leq\sum_j\log\Delta_j
       \leq n\log\frac{b-a}{n}.
 \tag{QLG27}
\]
Thus the \(n\log n\) in QLG26 cancels that in QLG27, and the complete
lower bound is
\[
 \boxed{\log J_n^\lambda(\alpha,\beta)
       \geq n^2F(\alpha,\beta)
                -n\{\log(b-a)+H_\lambda(\mu)\}.}
 \tag{QLG28}
\]
This also proves strict positivity of QLG7. No lower bound on the
density at an endpoint or on a quantile interval length has been used.

\paragraph{QLG6. The uniform order-\(n\) theorem.}
For every \((\alpha,\beta)\in\mathcal K\), every integer \(n\geq1\),
and each of the two measures QLG6, QLG22 and QLG28 are the explicit
two-sided finite estimate
\[
 \boxed{
 -n\{\log(b-a)+H_\lambda(\mu)\}
 \leq\log J_n^\lambda-n^2F
 \leq 4M_*n+n\log T_\lambda+n\log n-\log(n!).}
 \tag{QLG29}
\]
For entirely parameter-uniform constants set
\[
 L_0=\log(B_*M_*),\qquad
 L_\omega=\log(B_*M_*)+\log2+\tfrac12\log B_*+\sqrt{B_*},
\]
\[
 C_0=\max\{0,L_0,1+4M_*+\log^+T_*\},\qquad
 C_\omega=\max\{0,L_\omega,1+4M_*\}.
 \tag{QLG30}
\]
These constants are finite by QLG8--14. Since \(b-a\leq B_*\),
QLG11 bounds the lower side of QLG29 by \(-nL_\lambda\); QLG12--14
and QLG23 bound its upper side by \(nC_\lambda\). Consequently
\[
 \boxed{\left|\log J_n^{\lambda_0}(\alpha,\beta)
                 -n^2F(\alpha,\beta)\right|\leq C_0n,
 \qquad
 \left|\log J_n^{\lambda_\omega}(\alpha,\beta)
                 -n^2F(\alpha,\beta)\right|\leq C_\omega n.}
 \tag{QLG31}
\]
For the integral \(Z_n^\lambda\) without the Heine factor, the exact
relation \(\log Z_n^\lambda=\log J_n^\lambda+\log(n!)\) gives
\[
 \left|\log Z_n^\lambda-n^2F-\log(n!)\right|\leq C_\lambda n,
 \qquad
 \left|\log Z_n^\lambda-n^2F\right|
                         \leq n\log n+C_\lambda n.
 \tag{QLG32}
\]
The theorem therefore retains the distinction between the original
Heine determinant integral and the integral with labelled particles.
The gain to order \(n\) in QLG31 follows from the proved circle
comparison and the quantile-cell cancellation, with the same exact
potential, positive-half-line particle domain, and equilibrium of
EIQ1--33 throughout.

% END RX INLINE 41


% BEGIN RX INLINE 42; SHA256 0b552f80782b86651d3876ecf2dfd71a7c3f2e2c04ef005d6cadecb15d9bb629
% Complete original-source quantitative return, 14 September 2026.
% Providers: GDR, QLG, EIQ, WGP and LER; their proof bodies accompany this source.
\section{The original Gamma Hankel return with a finite moving-shift error}
\label{sec:QGT}

\paragraph{Original objects and exact square pushforward.}
Retain the original packet integers and the original density:
\[
 k=4l+1,\quad e=1+k(m-1),\quad q=e(k+1)^2=2n,
 \quad l,m\ge1,\quad
 \sigma(y)=\frac{|\Gamma(1/4+iy/2)|^2}{2\pi}.
 \tag{QGT1}
\]
In particular $n\ge18$ and $n\ge18l$: indeed
$n/l\ge8l+8+2/l\ge18$, the last expression increasing for $l\ge1$.
For every real $a>-1/2$ define the actual even moments and their
Hankel determinant by
\[
 \mu_a=\int_{\mathbb R}|y|^{2a}\sigma(y)\,dy,
 \qquad H_s^{(a)}=\det[\mu_{a+i+j}]_{0\le i,j<s},\quad s\ge1.
 \tag{QGT2}
\]
For the integer values used in the programme this is precisely
$\mu_a=m_{2a}$, with $\mu_0=\sqrt{2\pi}$ retained. The extension to
real $a$ is only the displayed integral on the same original measure.
Evenness of $\sigma$ and $t=y^2$ give, including both real half-lines,
\[
 \mu_a=\int_0^\infty t^{a-1/2}\sigma(\sqrt t)\,dt,
 \quad
 H_s^{(a)}=\frac1{s!}\int_{(0,\infty)^s}
       \prod_{i<j}(t_i-t_j)^2
       \prod_{i=1}^s t_i^{a-1/2}\sigma(\sqrt{t_i})\,dt_i.
 \tag{QGT3}
\]
To verify the second identity, expand the two determinants
$\det[t_i^{j-1}]_{i,j=1}^s$ as finite permutation sums, multiply them,
and integrate each monomial against the displayed positive product
measure. Every integral is finite: the density is bounded at zero,
$a>-1/2$, and the global Gamma upper bound below gives exponential
decay in $\sqrt t$ at infinity. Relabeling the integration variables
in each term gives $s!$ copies of the defining moment determinant.
The Vandermonde identity follows by its alternating zeros and its
leading coefficient one. This proves QGT3 with its literal factorial.
It is strictly positive because its integrand is positive on each
open cell of distinct positive coordinates.

\paragraph{The actual Gamma density, including its prefactor.}
Write $\alpha_\Gamma=\pi/2$, $D=1/16+\sqrt3/18$, and set
\[
 C_+=\max\{\sigma(0)e^{\alpha_\Gamma},\sigma(1)e^{\alpha_\Gamma+D}\},
 \qquad r(y)=y^{1/2}e^{\alpha_\Gamma y}\sigma(y)\quad(y>0).
 \tag{QGT4}
\]
Integrating the proved original logarithmic derivative GDR3 establishes,
by the complete deduction below, on the full stated domains
\[
 \sigma(y)\ge\sigma(1)(1+y)^{-1/2}e^{-\alpha_\Gamma y}\quad(y\ge0),
 \qquad \sigma(y)\le C_+y^{-1/2}e^{-\alpha_\Gamma y}\quad(y>0).
 \tag{QGT5}
\]
Here is also the short full deduction needed for this use. GDR3 says
$L(y)=\alpha_\Gamma+1/(2y)+\epsilon(y)/y$ and
$|\epsilon(y)|\le(1/8+\sqrt3/9)y^{-2}$. Integrating from $1$ to
$y\ge1$ shows
\[
 \log r(y)=\log\sigma(1)+\alpha_\Gamma-\int_1^y\epsilon(u)\,du/u,
 \qquad \left|\int_1^y\epsilon(u)\,du/u\right|\le D.
\]
Since $\alpha_\Gamma>D$, this gives $r(y)\ge\sigma(1)$ for $y\ge1$ and the
upper in QGT4 there. GDR7 implies $L(y)>0$ for $y>0$, so for
$0\le y\le1$ one has $\sigma(1)\le\sigma(y)\le\sigma(0)$.
On that interval $(1+y)^{-1/2}e^{-\alpha_\Gamma y}\le1$ and
$y^{-1/2}e^{-\alpha_\Gamma y}\ge e^{-\alpha_\Gamma}$ when $y>0$. These facts prove both
inequalities in QGT5 and keep the density's finite value at zero.
The elementary inequality $\alpha_\Gamma>D$ follows, for example, from
$\pi>3$ and $\sqrt3<2$.

Use exactly the dilation $t_i=q^2x_i$ in QGT3 and put
\[
 \alpha=\frac{a-3/4}{s},\qquad \beta=\frac{\pi q}{2s},
 \quad
 \mathcal J_{s,q}(\alpha,\beta)=\frac1{s!}
 \int_{(0,\infty)^s}\prod_{i<j}(x_i-x_j)^2
 \prod_i e^{-s(\beta\sqrt{x_i}-\alpha\log x_i)}
                     r(q\sqrt{x_i})\,dx_i.
 \tag{QGT6}
\]
Then the exact original determinant is
\[
 \boxed{H_s^{(a)}
   =q^{2as+2s(s-1)+s/2}\mathcal J_{s,q}(\alpha,\beta).}
 \tag{QGT7}
\]
Indeed the Vandermonde contributes $q^{2s(s-1)}$, and each density
factor transforms as
\[
 (q^2x)^{a-1/2}\sigma(q\sqrt x)q^2dx
 =q^{2a+1/2}x^{a-3/4}e^{-\alpha_\Gamma q\sqrt x}r(q\sqrt x)dx.
\]
This accounts separately for the original Jacobian, the quarter shift,
and the factor $q^{s/2}$. None of them is incorporated in the
equilibrium free energy.

\paragraph{A uniform finite partition theorem for this exact source.}
Let $\mathcal K=[3/2,5/2]\times[\pi/2,3\pi/2]$. Use the uniquely
proved equilibrium $\mu_{\alpha,\beta}$ of EIQ on the positive
half-line, its actual endpoints $a(\alpha,\beta),b(\alpha,\beta)$,
and its energy
\[
 F(\alpha,\beta)=\iint\log|x-y|\,d\mu_{\alpha,\beta}(x)
                    d\mu_{\alpha,\beta}(y)
       -\int(\beta\sqrt x-\alpha\log x)\,d\mu_{\alpha,\beta}(x).
 \tag{QGT8}
\]
In this paragraph only, the endpoint $a(\alpha,\beta)$ carries
arguments to distinguish it from the moment exponent $a$.
Let $a_*,B_*,M_*,C_0$ be the finite constants explicitly defined and
proved in QLG8--14 and QLG30. Set
\[
 L_0=\log(B_*M_*),\qquad
 C_* =\sigma(1)
          \left(\frac{\sqrt{a_*}}{1+\sqrt{a_*}}\right)^{1/2}>0,
\]
\[
 C=\max\{0,\ C_0+\log C_+,\ L_0-\log C_*\}.
 \tag{QGT9}
\]
For every $q\ge1$, integer $s\ge1$ and $(\alpha,\beta)\in\mathcal K$,
\[
 \boxed{\left|\log\mathcal J_{s,q}(\alpha,\beta)
                         -s^2F(\alpha,\beta)\right|\le Cs.}
 \tag{QGT10}
\]
For the upper estimate QGT5 gives $r(q\sqrt x)\le C_+$ everywhere
on the original positive half-line. QLG31 for the Lebesgue reference
measure therefore proves
$\log\mathcal J\le s^2F+(C_0+\log C_+)s$.
For the lower estimate retain the precise quantile cells QLG24--28
for this same equilibrium and its same $s$. Every coordinate in
these cells lies in $[a_*,B_*]$, where QGT5 gives
\[
 r(q\sqrt x)\ge\sigma(1)
      \left(\frac{q\sqrt x}{1+q\sqrt x}\right)^{1/2}\ge C_*.
\]
The last inequality uses $q\ge1$, $x\ge a_*$ and monotonicity of
$u/(1+u)$. QLG25--28 bound the integral over precisely these cells,
with their $s!$ permutations cancelling the original Heine factor.
Multiplying its integrand by the lower factor $C_*^s$ yields
$\log\mathcal J\ge s^2F-(L_0-\log C_*)s$.
This proves QGT10. A global lower bound by $x^{-3/4}$ at zero was
not used; the original density and its exact positive-domain factor
$r(q\sqrt x)$ remain in QGT6.

Combining QGT7 and QGT10 gives the exact finite decomposition
\[
 \boxed{\log H_s^{(a)}=
 [2as+2s(s-1)+s/2]\log q
 +s^2F\left(\frac{a-3/4}{s},\frac{\pi q}{2s}\right)+E_{s,q}(a),
 \qquad |E_{s,q}(a)|\le Cs.}
 \tag{QGT11}
\]
The bound holds whenever the displayed parameter pair belongs to
$\mathcal K$. The quantity $E_{s,q}(a)$ denotes the exact difference
in this equality; no derivative estimate for it has been assumed.

\paragraph{A quantitative moving secant, proved from convexity.}
Write $\mathcal L(\alpha,\beta)=\partial_\alpha F(\alpha,\beta)$,
the smooth original logarithmic moment proved in EIQ26 and EIQ29--30.
The fixed-interval integral EIQ29 and its proved differentiated
majorants give a finite constant
\[
 M=\max_{(\alpha,\beta)\in\mathcal K}
                   |\partial_\alpha\mathcal L(\alpha,\beta)|,
 \qquad A=M/2+2C.
 \tag{QGT12}
\]
For the two original ranks $s=n,n+1$, every real initial exponent
$a\in[q,q+l+1]$, and the original positive integer shift $l$, one has
\[
 \boxed{\left|\log\frac{H_s^{(a+l)}}{H_s^{(a)}}-2sl\log q
 -s^2\left[F\left(\frac{a+l-3/4}{s},\frac{\pi q}{2s}\right)
           -F\left(\frac{a-3/4}{s},\frac{\pi q}{2s}\right)\right]
                         \right|\le A l\sqrt s.}
 \tag{QGT13}
\]
All parameter neighbourhoods used to prove this bound lie in the
stated rectangle. Indeed put $u=(a-3/4)/s$, $t=l/s$ and
$\eta=s^{-1/2}<1/4$. For either rank,
\[
 u\ge2-\frac{11}{4s}\ge2-\frac{11}{72},\qquad
 u+t\le2+\frac{2l+1/4}{n}\le2+\frac18.
\]
Here $n\ge18l$ and $n\ge18$ give
$(2l+1/4)/n\le1/9+1/72=1/8$. Thus
$[u-\eta,u+t+\eta]\subset(3/2,5/2)$.
The second coordinate is $\pi$ for $s=n$ and $\pi n/(n+1)$ for
$s=n+1$, both lying in $[\pi/2,3\pi/2]$.

For fixed $q,s,\beta$, define
$g(v)=s^{-2}\log\mathcal J_{s,q}(v,\beta)$ on this neighbourhood.
This function is twice continuously differentiable and convex.
To justify differentiation under its integral, on every compact
subinterval of $v>0$ the exponent at zero is a positive power of
$x$, and the Gamma upper in QGT5 bounds the tails by a polynomial
times $\exp(-s\beta\sqrt x)$. The finite Vandermonde polynomial,
multiplied by any of the finitely many $\log x_i$ or their products
arising in the first two derivatives, has an integrable majorant.
Division by the strictly positive integral is legitimate. Under its
exact probability density,
\[
 (\log\mathcal J_{s,q})''(v)
       =s^2\operatorname{Var}_v\!\left(\sum_i\log x_i\right),
 \qquad g''(v)=\operatorname{Var}_v\!\left(\sum_i\log x_i\right)\ge0.
\]
This proves convexity without a property of the remainder.

For each $v\in[u,u+t]$, convexity gives
\[
 \frac{g(v)-g(v-\eta)}{\eta}\le g'(v)
              \le\frac{g(v+\eta)-g(v)}{\eta}.
\]
QGT10 implies $|g-F|\le C/s$. The integral formula for a difference
quotient of $F$ and the bound in QGT12 give, for both signs,
\[
 |g'(v)-\mathcal L(v,\beta)|
          \le\frac{M\eta}{2}+\frac{2C}{s\eta}.
 \tag{QGT14}
\]
In detail, the upper quotient of $F$ is
$\eta^{-1}\int_v^{v+\eta}\mathcal L(w,\beta)dw$ and differs
from $\mathcal L(v,\beta)$ by at most
$\eta^{-1}\int_0^\eta Mr\,dr=M\eta/2$; the lower quotient
has the identical bound. Integrating QGT14 over $[u,u+t]$ and
multiplying by $s^2$ gives $Al\sqrt s$ because
$\eta=s^{-1/2}$ and $t=l/s$. In QGT7 the terms
$2s(s-1)\log q$ and $(s/2)\log q$ cancel exactly between the
two values of $a$, while $2as\log q$ contributes $2sl\log q$.
This proves QGT13, retaining the rank $n+1$ and the source prefactor.

\paragraph{All four original high blocks in one exact centre.}
Define the same positive high product as in LET and LER,
\[
 Z_a=H_n^{(a)}H_{n+1}^{(a)}(H_n^{(a+1)})^2.
\]
For the original pair of endpoints $a=q,q+l$ put
\[
\begin{split}
 \mathcal C_{q,l}={}&n^2\left[
 F\left(\frac{q+l-3/4}{n},\pi\right)
 -F\left(\frac{q-3/4}{n},\pi\right)\right]\\
 &+(n+1)^2\left[
 F\left(\frac{q+l-3/4}{n+1},\frac{\pi n}{n+1}\right)
 -F\left(\frac{q-3/4}{n+1},\frac{\pi n}{n+1}\right)\right]\\
 &+2n^2\left[
 F\left(\frac{q+l+1/4}{n},\pi\right)
 -F\left(\frac{q+1/4}{n},\pi\right)\right],
\end{split}
 \tag{QGT15}
\]
\[
 \mathcal U_{q,l}=Al(3\sqrt n+\sqrt{n+1}).
 \tag{QGT16}
\]
Adding QGT13 with the displayed coefficients $1,1,2$ yields the
finite exact decomposition
\[
 \boxed{\log\frac{Z_{q+l}}{Z_q}
       =(4q+2)l\log q+\mathcal C_{q,l}+\xi_{q,l},
              \qquad |\xi_{q,l}|\le\mathcal U_{q,l}.}
 \tag{QGT17}
\]
The coefficient of the logarithm is exactly
$2l[n+(n+1)+2n]=2l(2q+1)$. The two odd blocks, the different
high ranks, and the original quarter shifts are all explicit in QGT15.

\paragraph{Finite return with the low endpoint and every product correction.}
Use the original $P_k$ of WGP1 and the original scalar moment of LER1,
so that
\[
 W_k=P_k+\log(Z_{q+l}/Z_q)-\log(\mu_{q+l}/\mu_q).
 \tag{QGT18}
\]
Put $\widetilde P_{q,l}=P_k+4lq\log q$, the exact finite sum WGP3,
and retain $F_{q,l},K_{q,l},T_{q,l},E_{q,l}$ with exactly the
definitions LER13 and LER17. Define the following centre in terms of
the already proved original objects:
\[
 \mathcal V_{q,l}=\widetilde P_{q,l}+\mathcal C_{q,l}
                         +2l\log q-F_{q,l}+K_{q,l}.
 \tag{QGT19}
\]
QGT17 and the signed remainder LER18 give
\[
 \boxed{\mathcal V_{q,l}-\mathcal U_{q,l}-E_{q,l}
       \le W_k\le
       \mathcal V_{q,l}+\mathcal U_{q,l}+T_{q,l}+E_{q,l}.}
 \tag{QGT20}
\]
To check the signs, LER18 is
$\log(\mu_{q+l}/\mu_q)=F_{q,l}-K_{q,l}+r_{q,l}+\varepsilon_{q,l}$
with $-T_{q,l}\le r_{q,l}\le0$ and
$|\varepsilon_{q,l}|\le E_{q,l}$. Substitution in QGT18 gives
$W_k=\mathcal V_{q,l}+\xi_{q,l}-r_{q,l}-\varepsilon_{q,l}$,
which proves QGT20. The original logarithmic powers cancel exactly:
$-4lq+(4q+2)l-2l=0$, since the logarithmic part of $F_{q,l}$ is
$2l\log(4q/\pi)$.

For an entirely integral-and-rational centre instead of the finite
source sum, set
\[
 p^{\rm c}_{q,l}=-4lq(2\log2-1)-4l^2\log2
       +\frac{8l^3+l}{6q}-\frac{l^4+l^2/4}{q^2},
\]
\[
 p^-_{q,l}=\frac l{6q}+\frac{8l^3+l}{12q^3},\qquad
 p^+_{q,l}=\frac{l^2}{3q^2}
       +\frac{8l^4+2l^2}{6q^4}
       +\frac{7(192l^5+80l^3-2l)}{1440q^3},
\]
\[
 \mathcal V^{\rm c}_{q,l}
     =p^{\rm c}_{q,l}+\mathcal C_{q,l}
                        +2l\log q-F_{q,l}+K_{q,l}.
 \tag{QGT21}
\]
WGP13 is the proved finite statement
$-p^-_{q,l}\le\widetilde P_{q,l}-p^{\rm c}_{q,l}\le p^+_{q,l}$.
Consequently the full original Gamma functional satisfies
\[
 \boxed{\mathcal V^{\rm c}_{q,l}-p^-_{q,l}-\mathcal U_{q,l}-E_{q,l}
       \le W_k\le
       \mathcal V^{\rm c}_{q,l}+p^+_{q,l}
                         +\mathcal U_{q,l}+T_{q,l}+E_{q,l}.}
 \tag{QGT22}
\]
In particular this bound uses the exact Gamma source and both original
low and high endpoint contributions; its finite constants were
obtained from their integrals and the proved equilibrium.

\paragraph{An explicit bound around the previously proved leading coefficient.}
Write $\mathcal L_0=\mathcal L(2,\pi)$ and
$\mathcal C_\Gamma=2\mathcal L_0-4(2\log2-1)$, exactly the positive
coefficient proved in GEL. The actual dilation proved in EIQ32 gives
$\mathcal L(v,\pi n/(n+1))=
\mathcal L(v,\pi)+2\log(1+1/n)$.
For $d\in\{-3/4,-11/4,1/4\}$, the fundamental theorem of calculus
and QGT12 give
\[
 \left|s^2\int_{2+d/s}^{2+(d+l)/s}
           \mathcal L(v,\pi)\,dv-sl\mathcal L_0\right|
 \le M\left(|d|l+\frac{l^2}{2}\right).
 \tag{QGT23}
\]
Indeed set $v=2+(d+r)/s$, $0\le r\le l$; the integrand's
difference from $\mathcal L_0$ is at most $M|d+r|/s$,
and $|d+r|\le |d|+r$. Apply this respectively at ranks $n,n+1,n$
with coefficients $1,1,2$ in QGT15. Since
$3/4+11/4+2(1/4)=4$, the result is
\[
 \left|\mathcal C_{q,l}-(2q+1)l\mathcal L_0
                      -2(n+1)l\log(1+1/n)\right|
                    \le M(2l^2+4l).
 \tag{QGT24}
\]
Every $v$ in these integrals is in the rectangle already verified above.
Define the fully explicit finite quantity
\[
\begin{split}
 \mathcal D_{q,l}={}&M(2l^2+4l)+l|\mathcal L_0|
        +2(n+1)l\log(1+1/n)+4l^2\log2\\
 &+\frac{8l^3+l}{6q}+\frac{l^4+l^2/4}{q^2}
        +\max\{p^-_{q,l},p^+_{q,l}\}
        +2l|\log(4/\pi)|+\frac{l^2}{q}\\
 &+\frac{l(16l^2-1)}{48q^2}
        +\frac{l^2(8l^2-1)}{48q^3}
        +K_{q,l}+T_{q,l}+E_{q,l}.
\end{split}
 \tag{QGT25}
\]
All displayed terms are nonnegative for $l\ge1$. Combining
QGT18, QGT17, WGP13, LER17--18 and QGT24, with the triangle inequality
applied only after their signs have been recorded, proves
\[
 \boxed{|W_k-lq\mathcal C_\Gamma|
                     \le\mathcal U_{q,l}+\mathcal D_{q,l}.}
 \tag{QGT26}
\]
The sign-sensitive sharper statement remains QGT22; QGT26 is its
explicit estimate about the original leading coefficient.

Along the actual packet, $\mathcal D_{q,l}=O(l^2+l)$ with a constant
independent of $m$. This follows term by term from
$q\ge(4l+2)^2$, the displayed rational quantities, LER20, and
$\log(1+1/n)\le1/n$ (integrate $1/(1+t)\le1$ on $[0,1/n]$).
Furthermore $\mathcal U_{q,l}\le4Al\sqrt{n+1}$.
Thus the quantitative theorem proves, for the same original integers,
\[
 \boxed{W_k=lq\mathcal C_\Gamma+O(l\sqrt q+l^2).}
 \tag{QGT27}
\]
The finite bounds QGT22 and QGT26 specify this error whenever an
order symbol is insufficient.
For each fixed $m\ge2$, the exact formula
$q=[1+(4l+1)(m-1)](4l+2)^2$ gives
$l/\sqrt q\to0$ and $l^2/q\to0$. Dividing QGT26 by $q$ therefore
proves the stronger original-return limit
\[
 \boxed{\frac{W_k-lq\mathcal C_\Gamma}{q}\longrightarrow0
                        \quad(l\to\infty,\ m\ge2\text{ fixed}).}
 \tag{QGT28}
\]
For $m=1$, the same exact calculation gives a finite bound on
$(W_k-lq\mathcal C_\Gamma)/q$, since $q=(4l+2)^2$ and
$l/\sqrt q\to1/4$. The bound does not identify the limit of that
quotient: the retained high-block uncertainty
$\mathcal U_{q,l}/q$ need not tend to zero in this case.
This identifies precisely what this finite estimate resolves at the
original $q$ scale, with every packet relation unchanged.

% END RX INLINE 42


% BEGIN RX INLINE 43; SHA256 9803d6ddf301807a1df7b4678b325241f6785f6ac1f4778c036037607ddc7dc6
% Exact evaluation of the already proved finite QGT error at original m=1.
% Provider: QUANTITATIVE_ORIGINAL_GAMMA_RETURN.tex, QGT1--28.
\section{The finite original \(m=1\) return bound at the \(q\) scale}
\label{sec:QGQ}

Retain the same original measure, determinants, source products, and
constants $A,M,\mathcal C_\Gamma$ of QGT1--28. At the original value
$m=1$, the packet is exactly
\[
 k=4l+1,\qquad q=(4l+2)^2,\qquad n=q/2,\qquad l\ge1.
 \tag{QGQ1}
\]
The explicit two finite nonnegative bounds in QGT16 and QGT25 have
the following limits on this packet:
\[
 \frac{\mathcal U_{q,l}}q\longrightarrow\frac A{\sqrt2},
 \qquad
 \frac{\mathcal D_{q,l}}q\longrightarrow
                  \frac M8+\frac{\log2}{4}.
 \tag{QGQ2}
\]
For the first limit use the exact identity
\[
 \frac{\mathcal U_{q,l}}q
 = A\frac l{\sqrt q}
          \left(\frac3{\sqrt2}+\sqrt{\frac12+\frac1q}\right).
\]
Here $l/\sqrt q=l/(4l+2)\to1/4$, so the limit is
$A(1/4)(4/\sqrt2)$, as displayed.

For the second limit, extract from the actual QGT25 sum the two
terms $2Ml^2$ and $4l^2\log2$. Their quotients by $q$ tend to
$M/8$ and $(\log2)/4$, since $l^2/q\to1/16$.
Every other summand tends to zero after division by $q$.
Here is the term-by-term justification using the original summands.
The terms $4Ml$, $l|\mathcal L_0|$, and
$2l|\log(4/\pi)|$ are fixed constants times $l$; their quotients
by $q$ tend to zero. The actual rank correction satisfies
\[
 0\le2(n+1)l\log(1+1/n)
       \le2(n+1)l/n\le4l,
\]
because $\log(1+u)=\int_0^u(1+t)^{-1}dt\le u$ and $n\ge1$.
The two next rational terms of QGT25 have sizes $O(l)$ and $O(1)$,
respectively:
\[
 \frac{8l^3+l}{6q},\qquad\frac{l^4+l^2/4}{q^2}.
\]
To verify these orders directly, substitute
$q=l^2(4+2/l)^2$; the remaining denominators are bounded below
by positive constants for $l\ge1$.
The two nonnegative product bounds have orders
$p^-_{q,l}=O(l^{-1})$ and $p^+_{q,l}=O(l^{-1})$.
For $p^-$ its two terms have orders $l^{-1},l^{-3}$.
For $p^+$ its three displayed terms have orders
$l^{-2},l^{-4},l^{-1}$, respectively. Each follows by the same
literal substitution in QGT21; the maximum therefore also has order
$l^{-1}$.
The remaining three displayed polynomial low terms
\[
 \frac{l^2}{q},\qquad
 \frac{l(16l^2-1)}{48q^2},\qquad
 \frac{l^2(8l^2-1)}{48q^3}
\]
have orders $1,l^{-1},l^{-2}$, respectively.
Finally the actual LER bounds retained in QGT25 give
$0<K_{q,l}\le3l/(32q^3)=O(l^{-5})$,
$T_{q,l}=O(l^{-3})$, and $E_{q,l}=O(l^{-9})$.
These are all the terms of QGT25, and each of these remaining terms
divided by $q\ge16l^2$ tends to zero. This proves the second
limit in QGQ2.

Apply the finite original inequality QGT26 for every integer $l$,
divide it by the same $q$, and use both proved limits in QGQ2.
The result is the quantitative statement
\[
 \boxed{\displaystyle
 \limsup_{l\to\infty}
 \left|\frac{W_k-lq\mathcal C_\Gamma}{q}\right|
 \le\frac A{\sqrt2}+\frac M8+\frac{\log2}{4}
 \qquad(m=1).}
 \tag{QGQ3}
\]
All constants in this inequality are the finite constants of the
proved original-source QGT calculation. The two limits in QGQ2
evaluate its explicit error bounds. Inequality QGQ3 alone does not
determine a limit for the actual expression inside its absolute value;
the exact centre $\mathcal C_{q,l}$ in QGT15 and the signed finite
interval QGT22 remain available for that continuing calculation.

\paragraph{The exact \(q\)-scale limit of the proved finite centre.}
The centre in QGT22 can itself be evaluated one scale further.
Use the smooth original equilibrium logarithmic moment of EIQ29--30,
with its same parameter domain, and put
\[
 \mathcal L_1=\partial_\alpha\mathcal L(2,\pi),\qquad
 M_2=\max_{(\alpha,\beta)\in\mathcal K}
                 |\partial_\alpha^2\mathcal L(\alpha,\beta)|<\infty,
 \qquad
 \mathcal C_\Gamma^{[q,1]}=\frac{\mathcal L_1}{8}
                              -\frac{\log2}{4}.
 \tag{QGQ4}
\]
The maximum is finite by the smoothness proved from the actual
fixed-interval density integral in EIQ29 and compactness of the same
rectangle $\mathcal K$ used in QGT12. Thus $\mathcal L_1$ and $M_2$
are derivatives of that explicitly proved integral.

For the original triples
$(s,d)=(n,-3/4),(n+1,-11/4),(n,1/4)$, all intermediate parameter
points below belong to $\mathcal K$, by the domain calculation in
QGT13. Integrating the derivative of $\mathcal L$ twice gives
\[
 \left|\mathcal L(2+v,\pi)-\mathcal L_0-v\mathcal L_1\right|
                       \le\frac{M_2v^2}{2}
\]
whenever the interval joining $2$ and $2+v$ is in that rectangle.
Indeed the exact remainder is
$\int_0^v\int_0^u\partial_\alpha^2\mathcal L(2+w,\pi)\,dw\,du$;
using oriented integrals gives the same absolute bound for $v<0$.
Now use exactly $v=(d+r)/s$, $0\le r\le l$, in the integral
of $\mathcal L$ defining the difference of $F$. The result is
\[
\begin{split}
 s^2\!\left[F\left(2+\frac{d+l}s,\pi\right)
              -F\left(2+\frac ds,\pi\right)\right]
   &=sl\mathcal L_0+(dl+l^2/2)\mathcal L_1+\zeta_{s,d,l},\\
 |\zeta_{s,d,l}|
   &\le\frac{M_2}{6s}\bigl[(|d|+l)^3-|d|^3\bigr].
\end{split}
 \tag{QGQ5}
\]
For the bound, the exact integral of the remainder is bounded by
$M_2(2s)^{-1}\int_0^l(|d|+r)^2dr$, which is the quantity displayed.

Apply QGQ5 to the literal three terms of QGT15 with their weights
$1,1,2$. The rank-$n+1$ beta change contributes exactly
$2(n+1)l\log(1+1/n)$ by EIQ32. Therefore
\[
\begin{split}
 \mathcal C_{q,l}={}&(2q+1)l\mathcal L_0
 +(2l^2-3l)\mathcal L_1
 +2(n+1)l\log(1+1/n)+\zeta_{q,l},\\
 |\zeta_{q,l}|\le{}&\frac{M_2}{6}
 \left\{
 \frac{(l+3/4)^3-(3/4)^3}{n}
 +\frac{(l+11/4)^3-(11/4)^3}{n+1}
 +2\frac{(l+1/4)^3-(1/4)^3}{n}
 \right\}.
\end{split}
 \tag{QGQ6}
\]
The weighted rank sum is $n+(n+1)+2n=2q+1$.
The weighted starting shift is exactly
$-3/4-11/4+2(1/4)=-3$, and the weights have sum four;
these give the coefficient $2l^2-3l$ in this formula.

Substitute QGQ6 into the exact centre $\mathcal V^{\rm c}_{q,l}$
of QGT21 and use the full expression for $F_{q,l}$ of LER17.
After displaying all original terms, the resulting identity is
\[
\begin{split}
 \mathcal V^{\rm c}_{q,l}-lq\mathcal C_\Gamma
 ={}&(2\mathcal L_1-4\log2)l^2
       +l\mathcal L_0-3l\mathcal L_1
       +2(n+1)l\log(1+1/n)\\
 &+\frac{8l^3+l}{6q}-\frac{l^4+l^2/4}{q^2}
       -2l\log(4/\pi)-\frac{l^2}{q}\\
 &+\frac{l(16l^2-1)}{48q^2}
       -\frac{l^2(8l^2-1)}{48q^3}+K_{q,l}+\zeta_{q,l}.
\end{split}
 \tag{QGQ7}
\]
The equality retains the original source, rank correction, low endpoint,
and signed cubic terms. Its first coefficient comes from the actual
four-block high contribution $2l^2\mathcal L_1$ and the original
source contribution $-4l^2\log2$.

On the original packet QGQ1, the finite bound in QGQ6 gives
$\zeta_{q,l}=O(l)$, because $n=q/2\ge8l^2$ and each numerator
there is a cubic polynomial in $l$ with zero constant term.
Every term of QGQ7 after its first term is consequently $O(l)$
or smaller: the rank correction was bounded by $4l$ above, the
rational terms have the orders individually proved before QGQ3,
and $K_{q,l}\le3l/(32q^3)$. Dividing QGQ7 by the literal
$q=(4l+2)^2$ proves
\[
 \boxed{\displaystyle
 \frac{\mathcal V^{\rm c}_{q,l}-lq\mathcal C_\Gamma}{q}
             \longrightarrow\mathcal C_\Gamma^{[q,1]}
             =\frac{\partial_\alpha\mathcal L(2,\pi)}8
                                  -\frac{\log2}{4}.}
 \tag{QGQ8}
\]

Finally QGT22 gives, before taking any limit, the two signed errors
\[
 -p^-_{q,l}-\mathcal U_{q,l}-E_{q,l}
 \le W_k-\mathcal V^{\rm c}_{q,l}
 \le p^+_{q,l}+\mathcal U_{q,l}+T_{q,l}+E_{q,l}.
\]
The two product bounds and both low errors divided by $q$ tend to
zero by the estimates already proved above. The remaining bound
$\mathcal U_{q,l}/q$ tends to $A/\sqrt2$ by QGQ2.
Combining this finite interval with QGQ8 proves the sharper conclusion
\[
 \boxed{\displaystyle
 \limsup_{l\to\infty}
 \left|\frac{W_k-lq\mathcal C_\Gamma}{q}
                  -\mathcal C_\Gamma^{[q,1]}\right|
                           \le\frac A{\sqrt2}\qquad(m=1).}
 \tag{QGQ9}
\]
Here the centre is forced by the original finite formula. The remaining
width is the proved error of its high-block transport. The estimate
does not replace that width by zero or assign an unproved limit to
the actual original Gamma functional.

% END RX INLINE 43


% BEGIN RX INLINE 44; SHA256 a086f4f2ea260621f5e1a1ec5ca2d0ce8843b464634d35f9affba276a472920f
% Canonical intrinsic operation: no independent source order is inserted.
\section{The Gamma orders produced by the existing exterior and sum maps}
\label{sec:IGO}

The source order in this section is one of the two orders already present
in the construction. The determinant size is the actual exterior degree.
Every relation coefficient remains in the resulting density.

\subsection{Original coordinates and the exact source measure}
Retain
\[
 k=4l+1,\quad l,m\geq1,\quad e=1+k(m-1),\quad
 q=e(k+1)^2,\quad c=k/2,\quad
 0<\delta<1/2,\quad\gamma>2,
\]
\[
 \chi(S)=\prod_{a,d=0}^{k}
 [S-c-(2a-k)\delta-i(2d-k)\gamma]^e,
 \qquad s\in\{1,k\},\quad b=s/2,\quad M_s=(2\pi)^{s/2}.
 \tag{IGO1}
\]
Here $s=1$ is the fixed source $\sigma$ and $s=k$ is the original
tensor Gamma source. The original line throughout is $S=c+iy$.
The change of letter from the second root index to $d$ leaves every
root and its order unchanged.

For each real $v>0$ define a density on the real coordinate by
\[
 \rho_v(y)=\frac{2^v}{4\pi\Gamma(v)}
       \left|\Gamma(v/2+iy/2)\right|^2,
 \qquad dm_{0,s}(y)=M_s\rho_b(y)\,dy.
 \tag{IGO2}
\]
Thus the full mass $M_s$ is retained. In particular
$dm_{0,1}=|\Gamma(1/4+iy/2)|^2dy/(2\pi)=d\sigma$.
For the other existing source, the recurrence of the Gamma function
at $k/4=l+1/4$ gives
\[
 dm_{0,k}(y)=\beta_k\prod_{j=0}^{l-1}
           [y^2+(2j+\tfrac12)^2]d\sigma(y),\qquad
 \beta_k=\frac{(2\pi)^{k/2}}{\sqrt2\Gamma(k/2)}.
\]
Indeed the squared Gamma recurrence contributes $4^{-l}$ times
the displayed product. The coefficient in IGO2, multiplied by
$4^{-l}$ and divided by the coefficient $1/(2\pi)$ of $\sigma$,
is $M_k2^{k/2-1-2l}/\Gamma(k/2)=\beta_k$.
This proves agreement with the original tensor source including
its full polynomial factor and mass.
For real $|t|<\pi/2$ and real $\xi$, the exact transforms are
\[
 \int_{\mathbb R}e^{ty}\rho_v(y)\,dy=(\cos t)^{-v},\qquad
 \int_{\mathbb R}e^{i\xi y}\rho_v(y)\,dy=(\cosh\xi)^{-v},
 \qquad \phi_s(t):=\int e^{ty}dm_{0,s}=M_s(\cos t)^{-b}.
 \tag{IGO3}
\]

Here is a direct justification of the source identities and the
integrability used below. The beta integral with the substitution
$u=\exp(2x)$ gives
\[
 \frac{\Gamma(v/2+iy/2)\Gamma(v/2-iy/2)}{\Gamma(v)}
  =2\int_{\mathbb R} e^{iyx}(2\cosh x)^{-v}\,dx.
\]
Consequently $\rho_v(y)=(2\pi)^{-1}
\int e^{iyx}(\cosh x)^{-v}dx$. Fourier inversion applies: the
function and its first two derivatives are integrable, and two
integrations by parts make its transform integrable. This proves
the characteristic identity and mass one in IGO3, with the displayed
Fourier convention. More strongly, for every $0<\theta<\pi/2$, move
the $x$ contour to $x+i\theta\operatorname{sgn}(y)$. The analytic
branch of $(\cosh x)^{-v}$ is fixed by its positive values on the
real axis; it exists in this strip because $\cosh x$ has no zero
there. Its modulus decays as $\exp(-v|\operatorname{Re}x|)$ on
each closed smaller strip. The vertical sides of the rectangle
therefore tend to zero and
$\rho_v(y)\leq C_{v,\theta}\exp(-\theta|y|)$.
This proves absolute convergence of every polynomial moment with
an exponential factor $e^{ty}$ when $|t|<\pi/2$. Holomorphic
continuation of the characteristic identity through this strip
then gives the first identity in IGO3. Strict positivity follows
from the zero-free Euler Gamma product in $\operatorname{Re}z>0$:
\[
 \frac1{\Gamma(z)}=z\exp(\gamma_{\!E}z)
       \prod_{n=1}^{\infty}(1+z/n)\exp(-z/n).
\]
Indeed the logarithms of the factors after their linear terms are
removed converge locally uniformly, and none of the factors vanishes
there. This product follows by taking the beta-integral limit
$\Gamma(z)=\lim_{n\to\infty}n!n^z/[z(z+1)\cdots(z+n)]$;
the linear sum is absorbed by
$\gamma_{\!E}=\lim_n(\sum_{j=1}^n1/j-\log n)$.

The natural line of the source of order $s$ is $U=b+iy$.
Its presentation on the original line is the exact translation
\[
 U=S+(b-c),\qquad
 F(S)\longmapsto F(U-(b-c)).
 \tag{IGO4}
\]
This is an algebra isomorphism, with inverse given by the opposite
translation, and its coefficient matrix on each degree filtration
is triangular with diagonal one. In particular, using $s=1$ in
IGO2 retains the original line $c+i\mathbb R$ through IGO4.
The original tensor identification is exact in the real coordinate:
\[
 dm_{0,k}=(dm_{0,1})^{*k}.
\]
Indeed the characteristic function of the convolution on the right
is $[M_1(\cosh\xi)^{-1/2}]^k
=M_k(\cosh\xi)^{-k/2}$, which is the characteristic function
of the left side by IGO3. The finite-measure Fourier uniqueness
argument given after IGO17 proves equality. The natural individual
coordinates are $1/2+iy_j$, whose sum is $c+i\sum_jy_j$.
If the individual sources are first written on the fixed original
line as $S_j=c+iy_j$, the same sum is obtained by the exact map
$\sum_jS_j\mapsto\sum_jS_j+k(1/2-c)$, applying IGO4 to all
$k$ factors. This retains the full tensor mass $M_1^k=M_k$
and its original centre $c$.

\subsection{The actual relation exterior vector and its sum measure}
The relation multiplier is the full polynomial
\[
 W_\chi(y)=|\chi(c+iy)|^2
 =\prod_{a,d=0}^{k}
 \left([(2a-k)\delta]^2+[y-(2d-k)\gamma]^2\right)^e
 =\sum_{j=0}^{2q}w_jy^j.
 \tag{IGO5}
\]
All its original parameters occur in the displayed product.
It is real, even, strictly positive on $\mathbb R$, and monic of
degree $2q$. Evenness follows by the permutation $d\mapsto k-d$.
Strict positivity follows because $k$ is odd, so
$(2a-k)\delta\ne0$ for every $a$. Thus $w_{2q}=1$ and
$w_j=0$ for odd $j$; no other coefficient is removed.

For $r\geq1$ put
\[
 \Delta_r(y)=\prod_{1\leq i<j\leq r}(y_j-y_i),\qquad
 d\lambda_{\chi,s,r}(y_1,\ldots,y_r)
 =\frac1{r!}\Delta_r(y)^2
       \prod_{i=1}^r W_\chi(y_i)\,dm_{0,s}(y_i),
 \qquad \Sigma_r(y)=\sum_{i=1}^r y_i,
 \quad \nu_{\chi,s,r}=(\Sigma_r)_*\lambda_{\chi,s,r}.
 \tag{IGO6}
\]
The notation with subscript $\varnothing$ instead of $\chi$ means
that each factor $W_\chi$ is replaced by the literal factor $1$.
For $r=0$, the product space is the one-point space with mass one,
the sum is $0$, and both pushforward measures are $\delta_0$.

This is precisely the existing exterior measure. To check all
phases and factorials, write $S_i=c+iy_i$ and define the exterior
vector in the product source Hilbert space by
\[
 \psi_{\chi,r}(y)=\frac1{\sqrt{r!}}
           \det[\chi(S_i)S_i^{j-1}]_{i,j=1}^r
   =\frac{i^{r(r-1)/2}}{\sqrt{r!}}
           \Delta_r(y)\prod_{i=1}^r\chi(c+iy_i).
 \tag{IGO7}
\]
Its squared modulus times $\prod dm_{0,s}(y_i)$ is exactly IGO6.
The phase in IGO7 is retained in the exterior map below.
Expanding both determinants and relabelling integration variables
shows that $\|\psi_{\chi,r}\|^2$ equals the Gram determinant of
$\chi,\chi S,\ldots,\chi S^{r-1}$ in the original source.
The $r!$ identical permutation contributions cancel the displayed
$1/r!$. This also proves the following full Laplace identity:
\[
 g_{\chi,s}(t)=\int e^{ty}W_\chi(y)\,dm_{0,s}(y)
       =W_\chi(\partial_t)\phi_s(t),\qquad
 H_{\chi,s,r}(t):=\int e^{tY}\,d\nu_{\chi,s,r}(Y)
       =\det[g_{\chi,s}^{(i+j)}(t)]_{i,j=0}^{r-1}.
 \tag{IGO8}
\]
All expansions and differentiations here are absolutely convergent
by the exponential estimate after IGO3, since the remaining
factors are fixed polynomials. For real $|t|<\pi/2$ the matrix
on the right is the Gram matrix of $1,y,\ldots,y^{r-1}$ for the
strictly positive measure $e^{ty}W_\chi(y)dm_{0,s}(y)$.
It is positive definite: a nonzero polynomial is nonzero on a real
open interval. Thus $H_{\chi,s,r}(t)>0$.

\subsection{The exact differential polynomial and its degree}
Put $z=\tan t$ and define polynomials recursively by
\[
 P_0(z)=1,\qquad
 P_{j+1}(z)=bzP_j(z)+(1+z^2)P_j'(z),\qquad
 Q_0(z)=\sum_{j=0}^{2q}w_jP_j(z),\qquad
 Q_{j+1}(z)=bzQ_j(z)+(1+z^2)Q_j'(z).
 \tag{IGO9}
\]
The product rule proves
$\phi_s^{(j)}(t)=\phi_s(t)P_j(\tan t)$ and
$g_{\chi,s}^{(j)}(t)=\phi_s(t)Q_j(\tan t)$.
These are exact polynomial recursions in all the $w_j$, with no
limit or discarded term.

For polynomials $F_0,\ldots,F_{r-1}$ write
$\operatorname{Wr}_z(F_0,\ldots,F_{r-1})
=\det[\partial_z^iF_j]_{i,j=0}^{r-1}$. Then
\[
 B=br+r(r-1),\qquad
 D_{\chi,s,r}(z)=\operatorname{Wr}_z(Q_0,\ldots,Q_{r-1}),
 \qquad
 H_{\chi,s,r}(t)=M_s^r(\cos t)^{-B}D_{\chi,s,r}(\tan t).
 \tag{IGO10}
\]
To prove this without losing a determinant factor, the matrix in
IGO8 is the $t$-Wronskian of $g,g',\ldots,g^{(r-1)}$.
In a Wronskian, multiplying every column function by a common
function $h$ multiplies the determinant by $h^r$:
the product-rule matrix on its rows is lower triangular with
diagonal $h$. Changing the independent variable from $t$ to $z$
multiplies its rows by a lower triangular matrix with diagonal
$1,z',\ldots,(z')^{r-1}$. Its determinant is
$(z')^{r(r-1)/2}$. Here $h=\phi_s$ and
$z'=1+z^2=(\cos t)^{-2}$. Multiplication of these two factors
gives exactly $M_s^r(\cos t)^{-br-r(r-1)}$ in IGO10.

The polynomial $P_j$ has parity $j$, degree $j$, and leading
coefficient $(b)_j=\prod_{h=0}^{j-1}(b+h)$; this follows directly
from IGO9. The polynomial $Q_j$ has parity $j$, degree $2q+j$,
and leading coefficient $(b)_{2q}(b+2q)_j$. Therefore
\[
 D_{\chi,s,r}(z)=\sum_{j=0}^{qr}d_jz^{2j},\qquad
 d_{qr}=((b)_{2q})^r
             \prod_{j=0}^{r-1}j!(b+2q)_j>0,
 \qquad D_{\chi,s,r}(z)>0\quad(z\in\mathbb R).
 \tag{IGO11}
\]
For the degree and leading coefficient, factor the leading
coefficient of each $Q_j$ from the Wronskian. If its degrees are
$n_j=2q+j$, the coefficient of its largest possible power of $z$
is $\det[(n_j)_{\underline i}]_{i,j=0}^{r-1}$.
The falling factorial in row $i$ is a monic polynomial of degree
$i$ in $n_j$; subtracting lower rows therefore changes it to the
ordinary Vandermonde determinant. Its value is
$\prod_{i<j}(n_j-n_i)=\prod_{j=0}^{r-1}j!$.
The power is $\sum_j(2q+j)-\sum_i i=2qr$, so the coefficient
cannot cancel. The parity of entry $(i,j)$ is $i+j$ modulo two;
each determinant term has even total parity
$\sum_i i+\sum_j j=r(r-1)$. Finally positivity of $D$ for
every real $z$ follows from the positive Gram in IGO8 by the
bijection $z=\tan t$ between $(-\pi/2,\pi/2)$ and $\mathbb R$.

\subsection{The Gamma order actually produced without relation factors}
For the literal exterior of the source alone, use $Q_0=1$ and
$Q_j=P_j$ in the same Wronskian calculation. Its degrees are
$j$, so its Wronskian is the positive constant
\[
 C_{b,r}=\prod_{j=0}^{r-1}j!(b)_j,\qquad
 H_{\varnothing,s,r}(t)=M_s^r C_{b,r}(\cos t)^{-B},
 \qquad
 d\nu_{\varnothing,s,r}(Y)=M_s^r C_{b,r}\rho_B(Y)\,dY.
 \tag{IGO12}
\]
The transform identity and Fourier uniqueness prove the density
formula, including its whole mass $M_s^r C_{b,r}$. The source
order in this sum law is forced by the actual exterior degree:
\[
 2B=sr+2r(r-1),\qquad
 d\nu_{\varnothing,s,r}
    =\frac{M_s^r C_{b,r}}{(2\pi)^B}\,dm_{0,2B}.
 \tag{IGO13}
\]
Here $dm_{0,a}(Y):=(2\pi)^{a/2}\rho_{a/2}(Y)dY$ names
the Gamma measure at the produced positive order $a=2B$;
the same notation below is used for the produced orders $2B+4h$.
The input order remains the original $s\in\{1,k\}$.
In the second expression $dm_{0,2B}$ is written in the real
coordinate $Y$. Its natural complex line is $U=B+iY$, whereas
the original sum coordinate is $T=rc+iY$. The exact translation
of these lines and their polynomial algebras is
\[
 U=T+(B-rc),\qquad
 F(T)\longmapsto F(U-(B-rc)).
 \tag{IGO14}
\]
It has inverse translation by $rc-B$, preserves every degree,
and has coefficient determinant one on every finite degree
filtration. If IGO4 is first applied to all $r$ factors, their
sum is $rb+iY$. The further displacement in IGO14 is exactly
$r(r-1)$, since $B-rb=r(r-1)$. Thus both the original sum centre
and the naturally produced Gamma centre have explicit maps.
Equations IGO12--IGO14 concern $r\geq1$; the unit case $r=0$
remains the point mass specified after IGO6.

\subsection{All relation terms: finite signed Gamma decomposition}
For the actual polynomial $W_\chi$, define every coefficient
\[
 a_h=\sum_{j=h}^{qr}(-1)^{j+h}\binom{j}{h}d_j,
 \qquad 0\leq h\leq qr.
 \tag{IGO15}
\]
Analytic continuation of IGO10 to $t=i\xi$, or the characteristic
integral in IGO8 directly, gives
\[
 \widehat\nu_{\chi,s,r}(\xi)
 =M_s^r(\cosh\xi)^{-B}D_{\chi,s,r}(i\tanh\xi)
 =M_s^r\sum_{h=0}^{qr}a_h(\cosh\xi)^{-(B+2h)}.
 \tag{IGO16}
\]
Indeed each term $d_j(i\tanh\xi)^{2j}$ is
$d_j(-1)^j[1-(\cosh\xi)^{-2}]^j$; expanding the finite
binomial power gives IGO15 with its stated signs. Thus
\[
 d\nu_{\chi,s,r}(Y)
  =M_s^r\sum_{h=0}^{qr}a_h\rho_{B+2h}(Y)\,dY
  =\sum_{h=0}^{qr}\frac{M_s^r a_h}{(2\pi)^{B+2h}}
                   \,dm_{0,\,2B+4h}(Y).
 \tag{IGO17}
\]
Fourier uniqueness applies to these finite signed measures. For
completeness, a finite signed measure with zero Fourier transform
has zero convolution with every Gaussian, by Fourier inversion
and Fubini. Letting the Gaussian widths decrease to zero and
testing against compactly supported continuous functions gives
zero measure. This proves the asserted equality of measures,
without a sign assumption on any $a_h$.

For each term in the second expression of IGO17 the original
sum line is still $T=rc+iY$. Its natural source line is
$U_h=B+2h+iY$, and its pullback is the explicit translation
$U_h=T+(B+2h-rc)$, with the polynomial map and inverse as in
IGO14. Thus IGO17 adds measures on the same original line; the
different natural centres have not been identified implicitly.

The top coefficient satisfies $a_{qr}=d_{qr}>0$. The largest
order forced by the full relation operation is consequently
\[
 2B+4qr=sr+2r(r-1)+4qr.
 \tag{IGO18}
\]
It equals $sq+6q^2-2q$ at $r=q$, and
$s(q+1)+6q(q+1)$ at $r=q+1$. These orders result from IGO6
and the leading coefficient of its original polynomial.
All the lower signed constituents in IGO17 remain part of the
same exact measure.

\subsection{A positive polynomial density with all masses retained}
The recurrence $\Gamma(z+1)=z\Gamma(z)$ gives, with no residual
power of two,
\[
 \frac{\rho_{B+2h}(Y)}{\rho_B(Y)}
 =\frac{\Gamma(B)}{\Gamma(B+2h)}
       \prod_{j=0}^{h-1}\left[Y^2+(B+2j)^2\right].
 \tag{IGO19}
\]
In detail, the ratio of the constants in IGO2 is
$2^{2h}\Gamma(B)/\Gamma(B+2h)$; the squared Gamma ratio is
$\prod_{j<h}[(B/2+j)^2+Y^2/4]$. Each of its $h$ factors
contributes $1/4$, cancelling $2^{2h}$ exactly. Define
\[
 R_{\chi,s,r}(X)=\sum_{h=0}^{qr}
   a_h\frac{\Gamma(B)}{\Gamma(B+2h)}
      \prod_{j=0}^{h-1}[X+(B+2j)^2].
 \tag{IGO20}
\]
Then the exact positive measure, in both conventions for the
underlying Gamma mass, is
\[
 \begin{split}
 d\nu_{\chi,s,r}(Y)
    &=M_s^r\rho_B(Y)R_{\chi,s,r}(Y^2)\,dY\\
    &=\frac{M_s^r}{(2\pi)^B}
           R_{\chi,s,r}(Y^2)\,dm_{0,2B}(Y),\\
 \deg R_{\chi,s,r}&=qr,\qquad
 [X^{qr}]R_{\chi,s,r}
       =d_{qr}\frac{\Gamma(B)}{\Gamma(B+2qr)}>0.
 \end{split}
 \tag{IGO21}
\]
The source on the second line has the centre map IGO14. In its
complex coordinate the multiplier is exactly
$R_{\chi,s,r}(-(U-B)^2)$; on the original sum line it is
$R_{\chi,s,r}(-(T-rc)^2)$.

To prove its strict positivity at every real $Y$, as opposed to
positivity only almost everywhere from the measure identity,
write the density of IGO6 relative to $dy_1\cdots dy_r$ as
\[
 w_{\chi,s,r}(y)
  =\frac{M_s^r}{r!}\Delta_r(y)^2
                  \prod_{i=1}^r W_\chi(y_i)\rho_b(y_i).
\]
For $r\geq2$ the coordinate change
$y_j=u_j$ for $j<r$ and $y_r=Y-\sum_{j<r}u_j$ has
absolute Jacobian one. It gives the actual fibre density
\[
 p_{\chi,s,r}(Y)=\int_{\mathbb R^{r-1}}
 w_{\chi,s,r}(u_1,\ldots,u_{r-1},Y-\textstyle\sum_{j<r}u_j)
                 \,du_1\cdots du_{r-1}.
 \tag{IGO22}
\]
It is finite and continuous. On a compact set of $Y$, the
exponential estimate following IGO3 bounds the integrand by
an integrable polynomial in $|u|$ times
$\exp(-\theta\sum_{j<r}|u_j|)$, with one fixed
$0<\theta<\pi/2$. This also justifies dominated convergence.
For each $Y$, choose any $t_0\ne0$ and distinct coordinates
$y_i=Y/r+t_0(i-(r+1)/2)$. Their sum is $Y$.
The Vandermonde is nonzero there, and all the other factors
are strictly positive. They remain positive in an open set
of the $u$ variables. Hence $p_{\chi,s,r}(Y)>0$.
For $r=1$ the same conclusion follows from
$p_{\chi,s,1}(Y)=M_s W_\chi(Y)\rho_b(Y)>0$.
The density in IGO21 is continuous and equals this density
almost everywhere by IGO17. Their equality therefore holds
everywhere. In particular
\[
 R_{\chi,s,r}(X)>0\quad(X\geq0),\qquad
 \nu_{\chi,s,r}(\mathbb R)
 =M_s^rD_{\chi,s,r}(0)
 =M_s^r\sum_{h=0}^{qr}a_h
 =M_s^r\int\rho_B(Y)R_{\chi,s,r}(Y^2)\,dY.
 \tag{IGO23}
\]
The coefficient sum equals $D(0)$ by the binomial identity in
IGO16 at $\xi=0$. Thus the full original relation determinant
is present as this mass. The Gamma order in IGO13, the full
polynomial in IGO20, and the signed constituents in IGO17 are
related by the proved identities IGO19--IGO23. In particular,
no choice of a larger source parameter is needed to form
this operation, and the polynomial coefficients are required
to recover its actual mass.

\subsection{The complete sum map, its adjoint, and the retained kernel}
Let $\mathcal H_{\lambda}=L^2(\mathbb R^r,\lambda_{\chi,s,r})$
and $\mathcal H_{\nu}=L^2(\mathbb R,\nu_{\chi,s,r})$, using
the full measures IGO6. The sum pullback
\[
 U_r:\mathcal H_{\nu}\longrightarrow\mathcal H_{\lambda},
 \qquad (U_rF)(y)=F(\Sigma_r y)
 \tag{IGO24}
\]
is an isometry, since its squared norm is exactly the defining
pushforward integral. Its range is closed. We use complex
inner products linear in the second variable. For $r\geq2$
its adjoint is the exact fibre integral
\[
 (U_r^*h)(Y)=\frac{1}{p_{\chi,s,r}(Y)}
  \int_{\mathbb R^{r-1}}
    h(u,Y-\textstyle\sum u_j)
    w_{\chi,s,r}(u,Y-\textstyle\sum u_j)\,du.
 \tag{IGO25}
\]
Every factor in this formula is specified in IGO22. The
denominator is the actual positive fibre mass; the source
measure and its total mass in IGO23 are unchanged. In the
Euclidean surface measure on $\Sigma_r^{-1}(Y)$, the same
numerator is $r^{-1/2}\int_{\Sigma_r^{-1}(Y)}hw\,d\mathcal H^{r-1}$.
Indeed $|\nabla\Sigma_r|=\sqrt r$, and the parametrization
in IGO22 has surface Jacobian $\sqrt r$.

For a proof of the adjoint formula and its domain, Cauchy--Schwarz
on each fibre gives
$|\int hw\,du|^2\leq p(Y)\int|h|^2w\,du$.
After division by $p(Y)$ and integration in $Y$, this gives
$\|U_r^*h\|_{\nu}\leq\|h\|_{\lambda}$, so the formula is
well defined for every equivalence class in $\mathcal H_\lambda$.
Fubini applied first to bounded integrable functions proves
$\langle U_rF,h\rangle_\lambda=\langle F,U_r^*h\rangle_\nu$;
the bound extends it to all $L^2$ functions. Taking $h=U_rF$
in IGO25 gives $U_r^*U_rF=F$. Consequently the exact orthogonal
projection and complete kernel are
\[
 \begin{split}
 \Pi_r&=U_rU_r^*,\qquad
 \mathcal H_\lambda=U_r\mathcal H_\nu\mathbin{\oplus}\ker U_r^*,\\
 \ker U_r^*&=\left\{h\in\mathcal H_\lambda:
    \int_{\mathbb R^{r-1}}h(u,Y-\textstyle\sum u_j)
     w_{\chi,s,r}(u,Y-\textstyle\sum u_j)\,du=0
      \text{ for almost every }Y\right\},\\
 \|h\|_\lambda^2&=\|U_r^*h\|_\nu^2+
                   \|h-U_rU_r^*h\|_\lambda^2.
 \end{split}
 \tag{IGO26}
\]
This identifies both the map and every discarded direction
in a passage to the sum. For $r=1$, $\Sigma_1$ is the identity,
$U_1$ and $U_1^*$ are the identity between the identical spaces,
and the kernel is zero. For $r=0$ the two Hilbert spaces are
the one-dimensional spaces at their unit points, the maps
are the identity under that identification, and the kernel
is again zero.

For $r\geq2$ the kernel remains nonzero even in the symmetric
subspace. To see this with its exact fibre dependence, let
$A(y)=\sum_i y_i^2$ and let
\[
 \overline A(Y)=\frac{1}{p(Y)}
     \int_{\mathbb R^{r-1}}A(u,Y-\textstyle\sum u_j)
                 w(u,Y-\textstyle\sum u_j)\,du.
\]
All its moments here exist by the same exponential bound.
The symmetric function $A-U_r\overline A$ lies in
$\ker U_r^*$ by IGO25. It is nonzero: for each fixed $Y$,
$A$ is a nonconstant quadratic function on the sum fibre,
and the weight is positive on the open set of distinct
coordinates. It therefore has strictly positive fibre
variance. This also records a concrete direction that the
sum coordinate alone does not reconstruct.

Finally, for $r\geq1$, the exterior Hilbert space has a complete typed map
from the same sum space. Realize $\bigwedge^rL^2(dm_{0,s})$
as the antisymmetric subspace of the product $L^2$ space with
the determinant convention in IGO7. Multiplication by
$\psi_{\chi,r}$ gives a unitary map from the symmetric
subspace of $\mathcal H_\lambda$ to that antisymmetric
subspace: its norm identity is IGO7, and its inverse is
division by $\psi_{\chi,r}$ away from the diagonal. The
diagonal has product source measure zero, and $\chi$ has
no zero on the original line, so this inverse exists almost
everywhere. Division takes an antisymmetric function to a
symmetric function and the norm identity proves its $L^2$
membership. Thus
\[
 \mathcal J_r:\mathcal H_\nu\longrightarrow
                  \bigwedge^r L^2(dm_{0,s}),\qquad
 \mathcal J_rF=\psi_{\chi,r}\,F(\Sigma_r y)
 \tag{IGO27}
\]
is an isometry with the original determinant phase retained.
For an antisymmetric vector $v$, its exact adjoint is
\[
 (\mathcal J_r^*v)(Y)=\frac1{p(Y)}
 \int_{\mathbb R^{r-1}}
   \overline{\psi_{\chi,r}(u,Y-\textstyle\sum u_j)}
   v(u,Y-\textstyle\sum u_j)
   \prod_{i=1}^r[M_s\rho_b(y_i)]\,du,
 \quad y=(u,Y-\textstyle\sum u_j).
 \tag{IGO28}
\]
Its kernel is precisely the vectors for which the numerator
in IGO28 vanishes for almost every $Y$; equivalently it is
$\psi_{\chi,r}$ times the symmetric part of the kernel in
IGO26. The proof is multiplication of IGO25 by the exact
identity $w=|\psi_{\chi,r}|^2\prod_i[M_s\rho_b(y_i)]$.
It also gives the full orthogonal decomposition with
projection $\mathcal J_r\mathcal J_r^*$. The nonzero
symmetric kernel vector just constructed supplies a nonzero
exterior kernel vector through this unitary multiplication.
At $r=0$, the empty determinant is $1$ and the zeroth exterior
space is $\mathbb C$; the exterior map and its adjoint are the
identity between the unit spaces, as for IGO24--IGO26.

For a bounded measurable function $f$, multiplication by
$f(Y)$ on $\mathcal H_\nu$ intertwines through IGO24 and
IGO27 with multiplication by $f(\Sigma_r y)$ on their
target spaces. For the unbounded original sum coordinate,
the domain on $\mathcal H_\nu$ is
$\{F:\int |YF(Y)|^2d\nu(Y)<\infty\}$; the identical
integral on the target proves the same intertwining on
exactly its transported domain. Adding the original
constant centre $rc$ and the factor $i$ gives the
intertwining for $T=rc+iY$ as well. Together with the
translations in IGO14 and after IGO17, these formulas
specify the spaces, measures, actions, adjoints, and kernels
through which the intrinsically produced Gamma density
enters the original spectral-sum construction.

% END RX INLINE 44


% BEGIN RX INLINE 45; SHA256 50fc30d31e2b6bdbf8d630f5029cf1091e10618de3e7bd3c22d229cc7ccf340a
% Original-source calculation. No auxiliary Gamma order is introduced.
\section{The intrinsic exterior mass and the original four-endpoint relation graph}
\label{sec:IFB}

The source order in this section is $s\in\{1,k\}$, with
$k=4l+1$, $l,m\geq1$, $e=1+k(m-1)$,
$q=e(k+1)^2$, and $c=k/2$.
Retain $0<\delta<1/2$, $\gamma>2$ and the complete original polynomial
\[
 \chi(S)=\prod_{a,b=0}^k
 [S-c-(2a-k)\delta-i(2b-k)\gamma]^e.
 \tag{IFB1}
\]
The repeated index $b$ in this product is local to the product; below
$b_s=s/2$ denotes the source parameter. We write
\[
 M_s=(2\pi)^{s/2},\qquad b_s=s/2,\qquad
 dm_{0,s}(y)=
 \frac{M_s2^{b_s-2}}{\pi\Gamma(b_s)}
 \left|\Gamma\left(\frac{b_s}{2}+\frac{iy}{2}\right)\right|^2dy.
 \tag{IFB2}
\]
For $s=1$ this is exactly $d\sigma$ in BSL1. For $s=k$ it is
exactly the original tensor-Gamma measure in BSL17: repeated use of
$\Gamma(z+1)=z\Gamma(z)$ gives
\[
 \left|\Gamma\left(l+\frac14+\frac{iy}{2}\right)\right|^2
 =4^{-l}\prod_{j=0}^{l-1}\bigl(y^2+(2j+\tfrac12)^2\bigr)
       |\Gamma(\tfrac14+\tfrac{iy}{2})|^2,
\]
and the coefficient in IFB2 times $4^{-l}$ equals
$M_k/(2\pi\sqrt2\Gamma(k/2))$. Thus no source has been selected
by matching its size to $q$. Both sources are the measures already
present in the original calculation. The original evaluation line in
both cases is $S=c+iy$.

\subsection{The original orthogonal coordinates and the complete quotient map}

The original source transform, with its full mass, is
\[
 \int_{\mathbb R}e^{ty}\,dm_{0,s}(y)
       =M_s(\cos t)^{-b_s},\qquad |t|<\pi/2.
 \tag{IFB3}
\]
Here is a direct check of its scale. With $\lambda=b_s/2$ and
$g(x)=\exp(\lambda x-e^x)$, Euler's integral gives
$\widehat g(\xi)=\Gamma(\lambda+i\xi)$ for
$\widehat g(\xi)=\int g(x)e^{i\xi x}dx$. Autocorrelation gives
\[
 \int e^{iu\xi}|\Gamma(\lambda+i\xi)|^2d\xi
 =2\pi\int g(x+u/2)g(x-u/2)dx
 =\frac{2\pi\Gamma(2\lambda)}{(2\cosh(u/2))^{2\lambda}}.
\]
The first equality follows by Fourier inversion of the integrable
autocorrelation. The function $g$ and every derivative decay
exponentially at $-\infty$ and faster than exponentially at $+\infty$,
so these Fourier operations and the displayed integrals converge.
The last equality follows from $v=e^x$. Substituting $y=2\xi$,
$u=2w$, and the coefficient in IFB2 gives
$M_s(\cosh w)^{-b_s}$. Rotation of Euler's integral through
$\theta\operatorname{sgn}(y)$, $0<\theta<\pi/2$, proves an
upper bound $C_{s,\theta}\exp(-\theta|y|)$ for the squared Gamma
factor. The small sector arc tends to zero because $\lambda>0$;
the large arc tends to zero because $\cos\theta>0$. This proves
analytic continuation of the characteristic transform to
$|\operatorname{Im}w|<\pi/2$ and hence IFB3. In particular all
polynomial moments below exist and $\int dm_{0,s}=M_s$.

Define real monic polynomials $p_n^{(s)}(y)$ by
\[
 \sum_{n=0}^\infty p_n^{(s)}(y)\frac{z^n}{n!}
 =(1+z^2)^{-b_s/2}\exp(y\arctan z).
 \tag{IFB4}
\]
Branches are analytic at zero. For sufficiently small real $z,w$,
IFB3 and the cosine addition formula show that the integral of the
product of these two generating functions is
$M_s(1-zw)^{-b_s}$. The exponential bound used for IFB3 permits
each derivative in a neighborhood of zero under the integral.
Coefficient comparison gives the exact norms and recurrence
\[
 h_n^{(s)}=M_s n!(b_s)_n,\qquad
 \int p_n^{(s)}p_j^{(s)}\,dm_{0,s}
       =\mathbf1_{n=j}h_n^{(s)},
\]
\[
 yp_n^{(s)}=p_{n+1}^{(s)}+n(b_s+n-1)p_{n-1}^{(s)},\quad
 p_0^{(s)}=1,\quad p_{-1}^{(s)}=0.
 \tag{IFB5}
\]
The recurrence also follows by differentiating IFB4 and multiplying
by $1+z^2$. Each $p_n^{(s)}$ has parity $n$. We henceforth retain
the superscript $s$ on $h_n$ where needed and suppress it on $p_n$.

Let $E_\chi=\mathbb C[S]/(\chi)$ with its fixed remainder frame
$(1,S,\ldots,S^{q-1})$, and define actual source polynomials
\[
 u_n(S)=\frac{p_n((S-c)/i)}{\sqrt{h_n^{(s)}}},\qquad
 C_s:\mathbb C^q\longrightarrow E_\chi,\quad
 C_s x=\sum_{n=0}^{q-1}x_n[u_n].
 \tag{IFB6}
\]
The positive square root is used. On the original line $u_n(c+iy)$
equals $p_n(y)/\sqrt{h_n^{(s)}}$. The source polynomials $u_n$ are
therefore an orthonormal frame, with all source factors explicitly
present. The leading coefficient of $u_n$ in $S$ is
$i^{-n}/\sqrt{h_n^{(s)}}$. Thus $C_s$ is triangular and invertible,
and
\[
 |\det C_s|^2=\prod_{n=0}^{q-1}(h_n^{(s)})^{-1}.
 \tag{IFB7}
\]
This is an exact change of coordinates on the original remainder
space, not a replacement of that space or of its frame.

For $j\geq0$ put
\[
 v_j=C_s^{-1}[u_{q+j}]\in\mathbb C^q,\qquad
 V_r=[v_0\ \cdots\ v_{r-1}],\qquad
 K_{s,r}=I_q+V_rV_r^*,\qquad \Omega_{s,r}=\det K_{s,r}.
 \tag{IFB8}
\]
The definitions for $r=0$ are $V_0$ the empty $q$ by zero matrix,
$K_{s,0}=I_q$, and $\Omega_{s,0}=1$. These vectors use the same
$C_s$ and the same remainder map for every endpoint. Fix
$r\geq0$ and $N=q+r-1$. The orthonormal source-coordinate map
\[
 U_{s,N}:\mathcal P_N\longrightarrow\mathbb C^q\oplus\mathbb C^r,
 \quad
 \sum_{n<q}x_nu_n+\sum_{j<r}z_ju_{q+j}\longmapsto(x,z)
\]
is a linear isometry. Polynomial division and IFB8 give the exact
commuting relation
\[
 J_NU_{s,N}^{-1}(x,z)=C_s(x+V_rz),\qquad
 U_{s,N}(\chi\mathcal P_{r-1})
       =\{(-V_rz,z):z\in\mathbb C^r\}.
 \tag{IFB9}
\]
To verify equality on the right, a vector in the displayed graph
has zero remainder by the first formula and hence is divisible by
the original monic $\chi$. Its degree is at most $q+r-1$, so its
quotient has degree at most $r-1$. Conversely every such multiple
has zero remainder, which forces its low coordinate to be $-V_rz$.
The map $z\mapsto U_{s,N}^{-1}(-V_rz,z)$ is consequently an exact
bijection from the high source coordinates to the original relation
space. No relation direction is projected out. When $r=0$, this
means the zero space on both sides.

For a fixed original remainder $a\in E_\chi$, write $w=C_s^{-1}a$.
The unique minimum source-norm lift has coordinates
\[
 (x,z)=\bigl(K_{s,r}^{-1}w, V_r^*K_{s,r}^{-1}w\bigr),
\qquad
 \|P_{\min}\|^2=w^*K_{s,r}^{-1}w.
 \tag{IFB10}
\]
Indeed $x+V_rz=K_{s,r}K_{s,r}^{-1}w=w$. The displayed vector
is orthogonal to every $(-V_rz',z')$, since its inner product with
that vector is zero. Every other lift differs by such a relation
vector, so the Pythagorean identity proves uniqueness and the
minimum formula. It follows in the original fixed remainder frame
that
\[
 G_{q+r-1}^{(s)}=C_s^{-*}K_{s,r}^{-1}C_s^{-1},\qquad
 \det G_{q+r-1}^{(s)}
       =\frac{\prod_{n=0}^{q-1}h_n^{(s)}}{\Omega_{s,r}}.
 \tag{IFB11}
\]
This gives the entire quotient Gram, including every off-diagonal
entry, before its determinant is taken.

\subsection{The relation determinant as a positive finite sum}

The integer $q$ is divisible by four. Pairing both original signed
root coordinates in IFB1 gives the exact real polynomial
\[
 \vartheta(y):=\chi(c+iy)
 =\prod_{A\in\{\delta,3\delta,\ldots,k\delta\}}
   \prod_{B\in\{\gamma,3\gamma,\ldots,k\gamma\}}
 \bigl[((y-B)^2+A^2)((y+B)^2+A^2)\bigr]^e.
 \tag{IFB12}
\]
Each quartet contributes two negative factors from its two opposite
real-root pairs, so their product has the positive sign displayed.
Each factor is strictly positive on $\mathbb R$ because $A>0$.
Thus $\vartheta$ is real, even, strictly positive on the real line,
and monic of degree $q$; its squared modulus is exactly
$W_\chi(y)=\vartheta(y)^2$. Write
\[
 \vartheta(y)=\sum_{a=0}^q c_a y^a,\qquad
 W_\chi(y)=\sum_{d=0}^{2q}w_dy^d,\qquad
 w_d=\sum_{a+a'=d}c_ac_{a'}.
 \tag{IFB13}
\]
These coefficients retain all roots and all repetitions. For an
explicit finite rule producing them, each factor of IFB12 is
$(y^4+2(A^2-B^2)y^2+(A^2+B^2)^2)^e$. Expand it as
\[
 \sum_{u+v+w=e}
 \frac{e!}{u!v!w!}\,[2(A^2-B^2)]^v(A^2+B^2)^{2w}
             y^{4u+2v},
 \tag{IFB14}
\]
and multiply these finite coefficient lists. In particular this rule
does not discard their signed middle coefficients.

Let $T_{n,h}$ be the coefficient of $p_h$ in $y^n$. The exact
finite recurrence following from IFB5 is
\[
 T_{0,0}=1,\quad T_{0,h}=0\ (h\ne0),\quad
 T_{n+1,h}=T_{n,h-1}+(h+1)(b_s+h)T_{n,h+1},
 \tag{IFB15}
\]
with indices outside $0\leq h\leq n$ interpreted as zero. Induction
using $yp_h=p_{h+1}+h(b_s+h-1)p_{h-1}$ proves the formula and
$y^n=\sum_hT_{n,h}p_h$. In particular $T_{n,n}=1$, and only
$h\equiv n\pmod2$ can occur. Define the full rectangular matrix
\[
 A_{h,j}=\sum_{a=0}^q c_aT_{a+j,h},\quad
 0\leq h\leq q+r-1,\quad 0\leq j<r,
 \qquad
 L_{h,j}=\sqrt{h_h^{(s)}}A_{h,j}.
 \tag{IFB16}
\]
Here $h_h^{(s)}$ means the source norm with index $h$; the matrix
letter $L$ denotes the full weighted coefficient matrix only in this
section. The definition says exactly that
$\vartheta(y)y^j=\sum_h A_{h,j}p_h(y)$.
Thus the Gram of these $r$ polynomials is $L^*L$.

In the original $S$-coefficient frame the relation vectors are
$(\chi,\chi S,\ldots,\chi S^{r-1})$. Its coefficient change to
the vectors $(\vartheta,\vartheta y,\ldots,\vartheta y^{r-1})$
is the exact triangular matrix
\[
 (T_c)_{h,j}=\binom jh c^{j-h}i^h\quad(0\leq h\leq j<r).
 \tag{IFB17}
\]
Its inverse is substitution $y=(S-c)/i$, and
$|\det T_c|=1$. Consequently the original relation determinant
$Z_{s,r}:=\det\operatorname{Gram}_{m_{0,s}}
(\chi,\ldots,\chi S^{r-1})$ equals $\det(L^*L)$.
We set $Z_{s,0}=1$.

Let $L_{\mathrm{low}}$ comprise rows $0,\ldots,q-1$ of $L$ and
$R$ comprise rows $q,\ldots,q+r-1$. Since $\vartheta y^j$ is
monic of degree $q+j$, the matrix $R$ is upper triangular with
diagonal $\sqrt{h_{q+j}^{(s)}}$. Hence it is invertible and
\[
 |\det R|^2=\prod_{j=0}^{r-1}h_{q+j}^{(s)},\qquad
 L_{\mathrm{low}}R^{-1}=-V_r.
 \tag{IFB18}
\]
For the second identity, the column space of $L$ is the relation
space in IFB9. Its high coordinates are precisely $R$ times its
coefficient vector; expressing a relation by its high coordinate
therefore gives the low graph matrix $L_{\mathrm{low}}R^{-1}$.
Uniqueness of the graph in IFB9 proves the asserted sign and every
matrix entry. Therefore
\[
 Z_{s,r}
 =\left(\prod_{j=0}^{r-1}h_{q+j}^{(s)}\right)
       \det(I_r+V_r^*V_r)
 =\left(\prod_{j=0}^{r-1}h_{q+j}^{(s)}\right)\Omega_{s,r}.
 \tag{IFB19}
\]
To justify the second equality without discarding zero singular
directions, eliminate either diagonal identity block in
$\left(\begin{smallmatrix}I_q&V_r\\-V_r^*&I_r\end{smallmatrix}\right)$.
The two block eliminations have determinant one and give
$\det(I_q+V_rV_r^*)=\det(I_r+V_r^*V_r)$, for all dimensions.

The full finite expansion of the same determinant is
\[
 \boxed{\displaystyle
 \Omega_{s,r}=
 \sum_{\substack{I\subset\{0,\ldots,q+r-1\}\\|I|=r}}
 |\det A_{I,\{0,\ldots,r-1\}}|^2
 \frac{\prod_{h\in I}h_h^{(s)}}
      {\prod_{j=0}^{r-1}h_{q+j}^{(s)}}.}
 \tag{IFB20}
\]
For completeness, expand each column of $L$ in the standard
orthonormal coordinates of $\mathbb C^{q+r}$. Their exterior
product has, in the exterior frame indexed by $r$-element subsets
$I$, coefficient $\det L_I$. The exterior frame is orthonormal,
so its squared norm is $\sum_I|\det L_I|^2$. Direct expansion of
that same squared norm gives $\det(L^*L)$ by the permutation
formula for a determinant. Substituting IFB16 and dividing by the
explicit positive factor in IFB19 proves IFB20. In particular the
subset $I=\{q,\ldots,q+r-1\}$ contributes exactly one. All other
terms are nonnegative; none of their signed coefficient calculations
has been replaced by a degree estimate. When $r=0$, the single
empty minor contributes one.

There is also a direct moment evaluation of exactly the same mass.
Define polynomials $P_n(z)$ by
\[
 P_0(z)=1,\qquad
 P_{n+1}(z)=b_szP_n(z)+(1+z^2)P_n'(z).
\]
Induction by differentiation gives
$\frac{d^n}{dt^n}(\cos t)^{-b_s}
=(\cos t)^{-b_s}P_n(\tan t)$. Differentiating IFB3 at zero
is justified by its exponential bound. It follows that
\[
 \boxed{\displaystyle
 Z_{s,r}=M_s^r\det\left[
       \sum_{d=0}^{2q}w_dP_{d+i+j}(0)
                          \right]_{i,j=0}^{r-1}.}
 \tag{IFB21}
\]
The moment degree required at the largest original endpoint $r=q+1$
is exactly $2q+2r-2=4q$. The finite recurrences IFB14--16 and
IFB21 therefore calculate every contribution from the original
parameters in finitely many additions, multiplications and divisions
by the displayed nonzero positive constants. IFB20 certifies the
positivity of the result even when coefficients in IFB21 cancel.

\subsection{The actual spectral-sum mass at the four ranks}

For $r\geq1$, let $\Delta(y)=\prod_{1\leq i<j\leq r}(y_j-y_i)$
and define the actual exterior measure and actual sum map
\[
 d\mathfrak m_{s,\chi,r}(y_1,\ldots,y_r)
 =\frac1{r!}\Delta(y)^2\prod_{j=1}^r
       W_\chi(y_j)\,dm_{0,s}(y_j),\qquad
 \Sigma_r(y_1,\ldots,y_r)=\sum_{j=1}^ry_j.
 \tag{IFB22}
\]
Set $\nu_{s,\chi,r}=(\Sigma_r)_*\mathfrak m_{s,\chi,r}$.
At $r=0$ the source space is a one-point space of mass one,
$\Sigma_0=0$, and $\nu_{s,\chi,0}=\delta_0$.
Expanding both Vandermonde determinants in IFB22 gives
\[
 \nu_{s,\chi,r}(\mathbb R)=Z_{s,r}.
 \tag{IFB23}
\]
To check the prefactor, write $\Delta=\det[y_i^{j-1}]$ and
expand it twice. Integrating a product gives a product of the
moments of $W_\chi dm_{0,s}$. Holding the first permutation fixed
and reindexing the second gives the determinant of its $r$ by $r$
moment matrix; the first permutation has $r!$ possible values.
Those $r!$ values cancel precisely the $1/r!$ in IFB22. Absolute
integrability follows from the polynomial degree and the exponential
source bound. The moment determinant agrees with the original
$S$-relation determinant by IFB17. This proves IFB23 with its
entire original mass.

Here is the precise retained fibre map. Write $\rho_s$ for the
density in IFB2 and, for $r\geq2$, put
\[
 y_r=Y-\sum_{a=1}^{r-1}y_a,\qquad
 F_r(Y;y_1,\ldots,y_{r-1})
 =\frac1{r!}\Delta(y)^2\prod_{a=1}^r W_\chi(y_a)\rho_s(y_a),
\]
\[
 g_r(Y)=\int_{\mathbb R^{r-1}}F_r(Y;y')\,dy'.
 \tag{IFB24}
\]
The change of variables $(y_1,\ldots,y_r)\mapsto(y_1,\ldots,
y_{r-1},Y)$ has Jacobian determinant one. Therefore
$d\nu_{s,\chi,r}(Y)=g_r(Y)dY$. These integrals are finite for every
$Y$: a polynomial in $|Y|+\sum|y_a|$ times
$\exp(-\theta\sum_{a<r}|y_a|)$ is an integrable dominating
function when $Y$ is fixed. Each $g_r(Y)>0$, because the source
densities and $W_\chi$ are positive and the distinct-coordinate
configurations form a nonempty open subset of the fibre.
For $r=1$, the formula is $g_1(Y)=W_\chi(Y)\rho_s(Y)>0$.

The map
\[
 \mathcal U_r:L^2(\nu_{s,\chi,r})\longrightarrow
                 L^2(\mathfrak m_{s,\chi,r}),\qquad
 (\mathcal U_r f)(y)=f(\Sigma_r y)
 \tag{IFB25}
\]
is an isometry, directly by the defining integral of a pushforward.
Its full adjoint, for $r\geq2$, is
\[
 (\mathcal U_r^*H)(Y)
 =\frac1{g_r(Y)}\int H(y_1,\ldots,y_{r-1},Y-\textstyle\sum y_a)
                         F_r(Y;y')\,dy'.
 \tag{IFB26}
\]
The integral exists for almost every $Y$ by Cauchy--Schwarz in the
fibre and Fubini. The same inequality shows its squared modulus,
multiplied by $g_r(Y)$, is bounded by the fibre integral of
$|H|^2F_r$; hence the adjoint is a contraction. Multiplying by a
test function $f(Y)$ and applying Fubini proves the adjoint identity
and thus IFB26. Its kernel is exactly the functions with weighted
fibre integral zero for almost every $Y$. For $r=1$, both maps are
the identity; at $r=0$ they are the identity of the one-dimensional
one-point $L^2$ space. Thus the fibre kernel and the full source
norm have specified domains and codomains; passing to a sum does
not silently identify the entire exterior space with its image.

The original complex spectral sum is
$\sum_{a=1}^r(c+iy_a)=rc+iY$. Accordingly the actual target line
of this pushforward is $rc+i\mathbb R$. If a Gamma density of
order $a$ is written on its natural line $a/2+i\mathbb R$, the
coordinate isomorphism for that constituent is
\[
 a/2+iY\longmapsto rc+iY,
 \quad\hbox{translation by }rc-a/2.
 \tag{IFB27}
\]
It preserves $dY$ and the complete density and mass. Different
constituents use their own stated translation; none changes the
original target centre $rc$.

Let $\mathcal D_{s,\chi,r}(z)$ denote the polynomial provided by
the intrinsic exterior/sum calculation, so its exact Laplace
transform is
\[
 \int e^{tY}\,d\nu_{s,\chi,r}(Y)
 =M_s^r(\cos t)^{-[b_sr+r(r-1)]}
             \mathcal D_{s,\chi,r}(\tan t).
 \tag{IFB28}
\]
This polynomial is completely determined by IFB21 as a differentiated
moment determinant, or by the full Wronskian construction in the
intrinsic source provider. To check this definition directly, let
$Q_0(z)=\sum_dw_dP_d(z)$ and
$Q_{j+1}(z)=b_szQ_j(z)+(1+z^2)Q_j'(z)$.
The Laplace version of IFB23 has moment entries
$M_s(\cos t)^{-b_s}Q_{i+j}(\tan t)$.
Thus one may take
\[
 \mathcal D_{s,\chi,r}(z)
 =(1+z^2)^{-r(r-1)/2}\det[Q_{i+j}(z)]_{i,j<r}.
\]
This expression is a polynomial: its determinant is the
$t$-Wronskian of the first $r$ derivatives of
$(\cos t)^{-b_s}Q_0(\tan t)$, and the following elementary two
identities show the factors explicitly.
For a common factor $f$, successive row subtractions in
$[\partial_t^i(fg_j)]$ give determinant $f^r$
times the Wronskian of the $g_j$; the subtraction coefficients come
from the product rule, and the triangular derivative matrix has
diagonal $f$. For a coordinate $z=z(t)$, the chain rule expresses
$\partial_t^i$ as $(z')^i\partial_z^i$ plus lower derivatives;
the determinant multiplier is $(z')^{r(r-1)/2}$.
Apply these identities with $f=(\cos t)^{-b_s}$,
$g_j=Q_j(z)$ and $z'=1+z^2$ to obtain
$\mathcal D_{s,\chi,r}(z)=\det[\partial_z^iQ_j(z)]_{i,j<r}$.
This also proves IFB28 without any assumed statement about the
degree or signs of this polynomial. Its full degree and Gamma
constituents are evaluated in the intrinsic source provider.

At the original four ranks its constant term is now calculated by
the positive original relation graph, not by size counting:
\[
 \boxed{\displaystyle
 \mathcal D_{s,\chi,r}(0)
 =\left(\prod_{j=0}^{r-1}(q+j)!(b_s)_{q+j}\right)
                       \Omega_{s,r},\quad r=0,1,q,q+1.}
 \tag{IFB29}
\]
Indeed IFB28 at $t=0$ gives $M_s^r\mathcal D(0)=Z_{s,r}$.
Substituting IFB19 and every factor
$h_{q+j}^{(s)}=M_s(q+j)!(b_s)_{q+j}$ proves IFB29. In particular
the full exterior-source mass appears to power $r$ on both sides
before its algebraic cancellation. Equation IFB29 is the exact
factor relating the intrinsic Gamma expression to the original
quotient determinant.

\subsection{The original four-endpoint identity and its positive update chain}

Write $\mathcal B_k^{(1)}=\mathcal B_k^\sigma$ and
$\mathcal B_k^{(k)}=\mathcal B_k^0$. The monic source norms in the
original $S$-basis are $h_n^{(s)}$: multiplication of $u_n$ by
$i^n\sqrt{h_n^{(s)}}$ makes it monic and leaves its squared norm
$h_n^{(s)}$. The triangular monic source/relation frame in BSL2--3
therefore gives, at the original four degrees, the full signed formula
\[
 \begin{split}
 \mathcal B_k^{(s)}
 &=S_q^{(s)}+\log Z_{s,q}+\log Z_{s,q+1}-\log Z_{s,1},\\
 S_q^{(s)}
 &=-\log h_q^{(s)}-2\sum_{j=q+1}^{2q-1}\log h_j^{(s)}
                                      -\log h_{2q}^{(s)}.
 \end{split}
 \tag{IFB30}
\]
The omitted rank-zero logarithm is exactly $\log Z_{s,0}=0$.
One can also prove this directly from IFB11 by taking the prescribed
$+,+,-,-$ signs at $r=0,1,q,q+1$. Substituting IFB19 into the
displayed formula cancels the monic norm factors as follows:
the two high products contain $h_q^2$, each
$h_{q+1},\ldots,h_{2q-1}$ twice, and $h_{2q}$ once; division by
the low factor $h_q$ leaves precisely the factors canceled by
$S_q^{(s)}$. In particular the full source mass exponent in the
relation terms is $q+(q+1)-1=2q$, while that in $S_q^{(s)}$ is
$-1-2(q-1)-1=-2q$. The exact result is
\[
 \boxed{\mathcal B_k^{(s)}
    =\log\Omega_{s,q}+\log\Omega_{s,q+1}-\log\Omega_{s,1}.}
 \tag{IFB31}
\]
Every original sign and all low and high polynomial factors have
now been retained and evaluated before this cancellation.

Define the finite, strictly positive update factors
\[
 d_{s,j}=1+v_j^*K_{s,j}^{-1}v_j,\qquad j\geq0.
 \tag{IFB32}
\]
Since $K_{s,j}=I_q+\sum_{a<j}v_av_a^*$ is positive definite, its
inverse exists and $d_{s,j}\geq1$. Multiplying
$K_{s,j+1}=K_{s,j}+v_jv_j^*$ by $K_{s,j}^{-1/2}$ on both sides
reduces its determinant ratio to that of $I+ww^*$,
$w=K_{s,j}^{-1/2}v_j$. In an orthonormal frame with first vector
$w/\|w\|$ when $w\ne0$, that operator has eigenvalues
$1+\|w\|^2,1,\ldots,1$; the case $w=0$ is the identity.
Consequently
\[
 \Omega_{s,r}=\prod_{j=0}^{r-1}d_{s,j},\qquad
 K_{s,j+1}^{-1}=K_{s,j}^{-1}
 -\frac{K_{s,j}^{-1}v_jv_j^*K_{s,j}^{-1}}{d_{s,j}}.
 \tag{IFB33}
\]
The inverse identity follows by multiplication by
$K_{s,j}+v_jv_j^*$; its two rank-one terms cancel exactly.
Starting with $K_{s,0}^{-1}=I_q$, it is a finite formula for every
inverse required in IFB32. The original four-endpoint expression
therefore equals the complete nonnegative sum
\[
 \boxed{\displaystyle
 \mathcal B_k^{(s)}
 =\log d_{s,0}+2\sum_{j=1}^{q-1}\log d_{s,j}+\log d_{s,q}.}
 \tag{IFB34}
\]
This is proved from the signed identity IFB31: each factor with
$1\leq j<q$ occurs in both high endpoints, $j=q$ occurs once,
and the rank-one low endpoint removes one of the two occurrences
of $j=0$. Formula IFB34 preserves all singular directions and all
off-diagonal interactions through the full inverse $K_{s,j}^{-1}$.
It is not an assertion that these interactions are independent.

\subsection{A finite companion calculation of every update vector}

Let $E_{\vartheta}=\mathbb C[y]/(\vartheta)$ in its monomial
frame and let $Y_\vartheta$ be the matrix of multiplication by $y$.
Thus its column $j<q-1$ is the vector with a one in row $j+1$,
and its last column is $(-c_0,\ldots,-c_{q-1})^t$.
The exact algebra isomorphism is
\[
 E_{\vartheta}\longrightarrow E_\chi,\qquad
 [P(y)]\longmapsto[P((S-c)/i)],
 \tag{IFB35}
\]
because substituting $y=(S-c)/i$ in
$\vartheta(y)=\chi(c+iy)$ gives $\chi(S)$, and the inverse
substitution is $S=c+iy$. Let $E_s$ be the coefficient matrix of
$p_0/\sqrt{h_0^{(s)}},\ldots,p_{q-1}/\sqrt{h_{q-1}^{(s)}}$
in the $y$-monomial frame; its diagonal is
$1/\sqrt{h_n^{(s)}}$. Hence $E_s$ is invertible. Define
\[
 Y_s=E_s^{-1}Y_\vartheta E_s,\qquad
 a_n=\sqrt{n(b_s+n-1)}\ (n\geq1),\quad a_0=0.
 \tag{IFB36}
\]
All entries are now determined by the finite original coefficients
in IFB14 and the finite original recurrence IFB5. For all $n\geq0$
let $z_n=C_s^{-1}[u_n]$. Its initial values are
$z_{-1}=0$ and $z_n$ the $n$th coordinate vector for $0\leq n<q$.
Dividing IFB5 by $\sqrt{h_n^{(s)}}$ and using
$h_n^{(s)}/h_{n-1}^{(s)}=n(b_s+n-1)$ gives
\[
 \boxed{\displaystyle
 z_{n+1}=\frac{Y_sz_n-a_nz_{n-1}}{a_{n+1}},\qquad
 v_j=z_{q+j}.}
 \tag{IFB37}
\]
Multiplication in the quotient is exactly $Y_s$ by IFB35--36, so
this is a recurrence for the actual remainders, including all of
their original polynomial relation coefficients. Every denominator
$a_{n+1}$ is positive. It suffices to run it through $n=2q-1$ to
obtain the last vector $v_q=z_{2q}$ used in IFB34. Combining
IFB14, IFB36--37, and IFB32--33 is thus a finite exact calculation
of the original four-endpoint quantity, using a fixed $q$-dimensional
space and its full matrices.

As an explicit coefficient check, summing the coefficients of
$y^{q-2}$ in the original quartets gives
\[
 c_{q-2}=\frac{qk(k+2)(\delta^2-\gamma^2)}6.
 \tag{IFB38}
\]
Indeed $\sum_{a\in\{1,3,\ldots,k\}}a^2=k(k+1)(k+2)/6$ and
there are $(k+1)/2$ choices of the other positive coordinate.
Each quartet contributes $2e(A^2-B^2)$, which proves IFB38.
The coefficient of $y^{n-2}$ in $p_n$ is
\[
 -A_n^{(s)},\qquad
 A_n^{(s)}=\sum_{j=1}^{n-1}j(b_s+j-1)
          =\frac{n(n-1)(3b_s+2n-4)}6.
 \tag{IFB39}
\]
The first equality follows successively from the coefficient of
$y^{n-1}$ in the recurrence for $p_{n+1}$, starting with $p_1=y$;
the second uses the finite sums of $j$ and $j^2$. Consequently
the coefficient of $p_{q-2}$ in the full original $\vartheta$
is exactly $c_{q-2}+A_q^{(s)}$. In IFB20 the minor obtained by
replacing the row $q$ in the highest-degree minor by $q-2$ is
therefore $c_{q-2}+A_q^{(s)}$, for every $r\geq1$.
To check that assertion, its first column has this entry in the
replacement row and is zero in all remaining rows, while the
remaining top-degree block is upper triangular with diagonal one.
It follows that
\[
 \Omega_{s,r}\geq
 1+\frac{(c_{q-2}+A_q^{(s)})^2h_{q-2}^{(s)}}{h_q^{(s)}}
 \quad(r\geq1).
 \tag{IFB40}
\]
For $r\geq2$, the analogous replacement of row $q+1$ by $q-1$
has squared minor $(c_{q-2}+A_{q+1}^{(s)})^2$: its two lowest
rows and first two columns have determinant
$-(c_{q-2}+A_{q+1}^{(s)})$, and the remaining block is upper
triangular. Making both replacements gives the product of the
two coefficients, because their opposite-parity cross entries
vanish. Keeping these four distinct nonnegative summands in
IFB20 proves the stronger finite bound
\[
 \Omega_{s,r}\geq
 \left(1+\frac{(c_{q-2}+A_q^{(s)})^2h_{q-2}^{(s)}}{h_q^{(s)}}\right)
 \left(1+\frac{(c_{q-2}+A_{q+1}^{(s)})^2h_{q-1}^{(s)}}
                         {h_{q+1}^{(s)}}\right),\quad r\geq2.
 \tag{IFB41}
\]
These are specific calculated terms in the full original polynomial
mass; IFB20 and IFB34 continue to retain every remaining term.
In particular they assert no asymptotic replacement of the full
quantity by these selected summands.

Finally $\vartheta$ is even and each $p_n$ has its original parity.
Division by this even monic polynomial preserves parity. Hence
$v_j$ has coordinates only in indices congruent to $q+j$, and
$K_{s,r}$ has zero entries between the two source parity subspaces.
This follows entrywise from IFB8, not by removing coordinates.
The full matrix $C_s$ in IFB11 carries this exact structure back
to the original $S$-remainder frame. The downstream GMB calculation
then composes that same map with the original mixed boundary map;
its contractions use the entire matrix $K_{s,r}$ and all its
original source phases.

% END RX INLINE 45


% BEGIN RX INLINE 46; SHA256 32c56427cb2f62c0d41bd2f6830813a84aa01b1bff21a39fd4291e9b54f35556
% Original source remainder coordinates and the full BRU mixed row.
\section{The exact map from Gamma remainders to the original mixed rows}
\label{sec:GMB}

This calculation uses only the two sources already present in the
original construction. It retains the packet order $k$, its centre
$c=k/2$, its monic polynomial $\chi$ of degree $q$, and the original
surjection $\Lambda:E_\chi\to B_{\rm obs}$, where
$E_\chi=\mathbb C[S]/(\chi)$. The source order is $s\in\{1,k\}$.
Write $b_s=s/2$ and $M_s=(2\pi)^{s/2}$. For the full source measure
$dm_{0,s}$, the monic real orthogonal polynomials $p_d(y)$ and their
squared norms are
\[
\begin{gathered}
 p_0=1,\quad p_1=y,\qquad
 p_{d+1}(y)=yp_d(y)-d(b_s+d-1)p_{d-1}(y),\\
 h_d^{(s)}=M_s d!(b_s)_d,\qquad
 u_d(S)=\frac{p_d((S-c)/i)}{\sqrt{h_d^{(s)}}}.
\end{gathered}\tag{GMB1}
\]
The recurrence and every source mass are those proved for the
original Gamma source. On $S=c+iy$ the $u_d$ are orthonormal;
their leading coefficients in $S$ are $i^{-d}/\sqrt{h_d^{(s)}}$.
Thus no phase or mass is removed by introducing these explicit
coordinates. Inner products below are conjugate linear in the
first variable.

Let $C_s:\mathbb C^q\to E_\chi$ have columns
$[u_0],\ldots,[u_{q-1}]$ in the fixed original remainder frame.
It is invertible, with the literal determinant
\[
 \det C_s=i^{-q(q-1)/2}\prod_{d=0}^{q-1}(h_d^{(s)})^{-1/2}.
 \tag{GMB2}
\]
For $j\geq0$ define the actual remainder coordinates and matrices
\[
\begin{gathered}
 v_j=C_s^{-1}[u_{q+j}],\qquad
 K_j=I_q+\sum_{a=0}^{j-1}v_av_a^*,\qquad
 t_j=v_j^*K_j^{-1}v_j,\qquad d_j^{\rm gain}=1+t_j,\\
 A_j=[I_q,v_0,\ldots,v_{j-1}]:\mathbb C^{q+j}\to\mathbb C^q.
\end{gathered}\tag{GMB3}
\]
The empty sum is zero. In particular $K_0=I_q$, $K_j=A_jA_j^*$
is positive definite, and $d_j^{\rm gain}\geq1$.
The original polynomial remainder map in these orthonormal source
coordinates is exactly $C_sA_j$. This follows by applying it to
each displayed basis vector, including all high degrees.

For $x\in\mathbb C^q$, the unique minimum norm source lift of $C_sx$
at degree $q+j-1$ has coefficient vector
$A_j^*K_j^{-1}x$. Indeed its remainder coordinates are $x$,
and it is orthogonal to $\ker A_j$. Every other lift differs by
an element of that kernel, so its squared norm is
$x^*K_j^{-1}x$ plus the squared norm of that difference. Hence the
original quotient Gram is
\[
 G_{q+j-1}=C_s^{-*}K_j^{-1}C_s^{-1}.
 \tag{GMB4}
\]
This proves the actual quotient map and its metric, not only its
determinant. The source kernel has its full dimension $j$.

\subsection{The monic relation, its original phase, and its row}
Put $D=q+j$. In degree at most $D$, with the last coordinate
corresponding to $u_D$, form
\[
 z_j=\begin{pmatrix}-A_j^*K_j^{-1}v_j\\1\end{pmatrix},\qquad
 \mathfrak d_{D-1}(S)
   =i^D\sqrt{h_D^{(s)}}\left(
       u_D(S)-\sum_{a=0}^{D-1}(A_j^*K_j^{-1}v_j)_a u_a(S)
                         \right).
 \tag{GMB5}
\]
The first vector has zero remainder because
$-A_jA_j^*K_j^{-1}v_j+v_j=0$. Thus $\mathfrak d_{D-1}$ is
divisible by the original $\chi$. Its leading coefficient in $S$
is exactly one by GMB1. Moreover $z_j$ is orthogonal to every
old relation: its old coordinates lie in $\operatorname{im}A_j^*$.
It follows that $\mathfrak d_{D-1}=\chi\,p_j^{\chi,s}$, where
$p_j^{\chi,s}$ is the original monic polynomial orthogonal to lower
degrees in the measure $|\chi(c+iy)|^2dm_{0,s}(y)$. To verify
uniqueness without any assumption on coordinates, the difference
of two such monic relations is an old relation, and its inner
product with that difference is zero. Positivity makes it zero.
Consequently GMB5 is exactly the relation line used in BRU and
HMR32, with its prescribed monic orientation.

Let $a_{D-1}=\|\mathfrak d_{D-1}\|^2$ and let
$\alpha_{D-1}$ be the row with entries
$\langle\mathfrak d_{D-1},S^a\rangle$, $0\leq a<q$.
GMB3--5 give the complete formulas
\[
\begin{gathered}
 a_{D-1}=h_D^{(s)}d_j^{\rm gain},\\
 \alpha_{D-1}
   =-i^{-D}\sqrt{h_D^{(s)}}\,v_j^*K_j^{-1}C_s^{-1}.
\end{gathered}\tag{GMB6}
\]
For the first equality, the squared norm of $z_j$ is
$1+v_j^*K_j^{-1}A_jA_j^*K_j^{-1}v_j=1+t_j$.
For the second, the degree $D$ orthogonal polynomial has zero
inner product with every $S^a$, $a<q$. A polynomial of degree
less than $q$ has its low orthonormal coordinates $C_s^{-1}[P]$,
so the old lift term in GMB5 has inner product
$v_j^*K_j^{-1}C_s^{-1}[P]$ with it. Conjugating the scalar
$i^D\sqrt{h_D^{(s)}}$ proves precisely the phase and minus sign
in GMB6.

Direct multiplication proves
\[
 (K_j+v_jv_j^*)^{-1}
 =K_j^{-1}-\frac{K_j^{-1}v_jv_j^*K_j^{-1}}{d_j^{\rm gain}}.
\]
Using GMB4 and GMB6 therefore gives the original BRU update,
including its complete coefficient row:
\[
 G_{D-1}=G_D+\frac{\alpha_{D-1}^*\alpha_{D-1}}{a_{D-1}},
 \qquad
 \frac{\det G_{D-1}}{\det G_D}=d_j^{\rm gain}.
 \tag{GMB7}
\]
The determinant equality follows either by taking determinants
in GMB4, or by conjugating $K_j+v_jv_j^*$ by $K_j^{-1/2}$:
the resulting rank-one update has eigenvalue $1+t_j$ in its
one possible nontrivial direction and eigenvalue one elsewhere.
This also proves it when $v_j=0$.

\subsection{The unchanged observation and both mixed contractions}
The observation in these explicit source coordinates is
\[
 R_s=\Lambda C_s:\mathbb C^q\twoheadrightarrow B_{\rm obs},
 \quad M_j^{\rm obs}=R_sK_jR_s^*,\quad w_j=R_sv_j,\quad
 \xi_j=w_j^*(M_j^{\rm obs})^{-1}w_j.
 \tag{GMB8}
\]
The notation $R_s$ here names a matrix of the fixed observation,
not a new observation. Every Taylor unit and primary coefficient
already in $\Lambda$ is still in this product. For zero boundary
rank the matrices are empty and $\xi_j=0$. Otherwise
$M_j^{\rm obs}$ is positive definite by surjectivity. Define
\[
 r_j^{\rm ker}=v_j-K_jR_s^*(M_j^{\rm obs})^{-1}w_j.
 \tag{GMB9}
\]
Then $R_sr_j^{\rm ker}=0$. Its two summands give the exact
$K_j^{-1}$-orthogonal decomposition
\[
 t_j=\xi_j+
       (r_j^{\rm ker})^*K_j^{-1}r_j^{\rm ker},\qquad
 0\leq\xi_j\leq t_j.
 \tag{GMB10}
\]
Indeed the cross inner product is
$(r_j^{\rm ker})^*R_s^*(M_j^{\rm obs})^{-1}w_j=0$, and the
squared norm of the second summand in GMB9 is $\xi_j$.
The map $C_s$ restricts to a bijection
$\ker R_s\to\ker\Lambda$, since
$\Lambda C_sx=R_sx$. This supplies the exact map of the retained
kernel, including its original coefficient embedding.

Let $I:\ker\Lambda\hookrightarrow E_\chi$ be that original
embedding. At the new degree $D$, retain the original forms and
minimum section
\[
\begin{gathered}
 H_D^{\rm ker}=I^*G_DI,\qquad
 Q_D=(\Lambda G_D^{-1}\Lambda^*)^{-1},\qquad
 S_D=G_D^{-1}\Lambda^*Q_D,\\
 \beta_{D-1}=(\alpha_{D-1}I)(H_D^{\rm ker})^{-1}
                                  (\alpha_{D-1}I)^*,\qquad
 \nu_{D-1}=(\alpha_{D-1}S_D)Q_D^{-1}
                                  (\alpha_{D-1}S_D)^*.
\end{gathered}\tag{GMB11}
\]
All inverses on a zero dimensional space have their unique empty
meaning. The section obeys $\Lambda S_D=I$ and $I^*G_DS_D=0$.
The map $[I,S_D]$ is a bijection: the inverse sends $z$ to its
kernel component $z-S_D\Lambda z$ and its boundary component
$\Lambda z$. In this frame the metric has diagonal blocks
$H_D^{\rm ker}$ and $Q_D$. Inverting the congruence proves
\[
 G_D^{-1}=I(H_D^{\rm ker})^{-1}I^*
                       +S_DQ_D^{-1}S_D^*.
 \tag{GMB12}
\]
This proves all decomposition identities used below while
retaining the original source, kernel, boundary and section maps.

Set $d=1+t_j$, $h=h_D^{(s)}$, and
$N=M_j^{\rm obs}+w_jw_j^*$. Equations GMB4 and GMB8 imply
$Q_D=N^{-1}$ and
$S_D=C_s(K_j+v_jv_j^*)R_s^*N^{-1}$. Direct multiplication yields
$N^{-1}w_j=(M_j^{\rm obs})^{-1}w_j/(1+\xi_j)$.
GMB6 now gives
\[
 \alpha_{D-1}S_D
 =-i^{-D}\sqrt h\,d\,w_j^*N^{-1},\qquad
 \nu_{D-1}=h d^2\frac{\xi_j}{1+\xi_j}.
\]
On the other hand GMB6 and GMB12 give
\[
 \beta_{D-1}+\nu_{D-1}
   =\alpha_{D-1}G_D^{-1}\alpha_{D-1}^*
   =h(t_j+t_j^2)=h d t_j.
\]
Subtracting the displayed expression for $\nu_{D-1}$, with no
sign assumption, proves the exact full mixed row formula
\[
\boxed{\begin{gathered}
 a_{D-1}=h d,\qquad
 \beta_{D-1}=h d\frac{t_j-\xi_j}{1+\xi_j},\qquad
 \nu_{D-1}=h d^2\frac{\xi_j}{1+\xi_j},\\
 1+\frac{\beta_{D-1}}{a_{D-1}}
      =\frac{1+t_j}{1+\xi_j},\qquad
 1+\frac{\nu_{D-1}}{a_{D-1}+\beta_{D-1}}=1+\xi_j.
\end{gathered}}\tag{GMB13}
\]
Both factors are at least one by GMB10, and their product is
$1+t_j$, the original complete increment in GMB7. These are the
individual BRU kernel and boundary factors, in their original
order. In particular the kernel factor has the denominator
$1+\xi_j$; it is not replaced by $1+t_j-\xi_j$.

\subsection{The four original endpoints and their finite calculation}
For a sequence $f_j$, put
\[
 \mathcal W_q(f)=f_0+2\sum_{j=1}^{q-1}f_j+f_q.
 \tag{GMB14}
\]
The original signed quotient volume is
$\log\det G_{q-1}+\log\det G_q
 -\log\det G_{2q-1}-\log\det G_{2q}$.
Telescoping GMB7 at these four precise degrees gives
\[
\begin{gathered}
 \mathcal B_k^{(s)}=\mathcal W_q\bigl(\log(1+t_j)\bigr),\\
 F_{K}^{(s)}=\mathcal W_q\left(
                    \log\frac{1+t_j}{1+\xi_j}\right),\qquad
 F_{B}^{(s)}=\mathcal W_q\bigl(\log(1+\xi_j)\bigr),\\
 F_{K}^{(s)}+F_{B}^{(s)}=\mathcal B_k^{(s)}.
\end{gathered}\tag{GMB15}
\]
To check the indices, the relation added at this step has degree
$D=q+j$ and is the BRU row with index $D-1$. Thus $j=0,q$
are exactly the two weight-one rows $q-1,2q-1$; all intervening
rows have weight two. GMB13 proves the individual factors, so
GMB15 is the full original mixed decomposition at $s=1$ or $s=k$.
It does not assign an arithmetic source to either Gamma order.
The established arithmetic comparison retains its own exact
transition scalar and its original maps.

For completeness this is a finite calculation from the original
polynomial, rather than a definition by an unspecified projection.
Let $T_\chi$ be multiplication by $S$ on the fixed remainder frame
and set $Y_\chi=(T_\chi-cI)/i$. Starting with
$r_0=[1]$, $r_1=Y_\chi[1]$, compute
\[
 r_{d+1}=Y_\chi r_d-d(b_s+d-1)r_{d-1},\qquad
 C_s=\left[\frac{r_0}{\sqrt{h_0^{(s)}}},\ldots,
                 \frac{r_{q-1}}{\sqrt{h_{q-1}^{(s)}}}\right],
 \qquad v_j=\frac{C_s^{-1}r_{q+j}}{\sqrt{h_{q+j}^{(s)}}}.
 \tag{GMB16}
\]
Polynomial division and GMB1 prove by induction that $r_d=[p_d(Y)]$.
Every entry of $T_\chi$ is a literal coefficient of the original
$\chi$, including its primary multiplicities. The matrices
$K_j$, $M_j^{\rm obs}$ and the scalars in GMB13--15 are then obtained
by the displayed finite operations. Their inverses exist by the
proved positivity, including the stated empty cases. This closes
the exact morphism from the intrinsic exterior mass calculation,
through the original polynomial remainders, to both original
mixed contractions. It introduces no new source order or factor.

% END RX INLINE 46



\paragraph{Current original observation input.}
MRI1a--b now computes the original observed columns from the complete
unit jets, period vectors and cyclic average. MRI2a--b relates their
full coupled recurrence to the preceding-degree minimum-section kernel
norm, preserving the actual action defect. MRI3a--4a returns both
ordered contractions and all four boundary determinants. The actual
BRI2, BRI6 and BRI9 sites now propagate the exact sourcewise
kernel differences and boundary return, including the literal
arithmetic source. MRI1d proves the original period rank, MRI4c
proves the individual $q^2$ coefficients on its stated domain,
and MRI5d gives the boundary differences at the $kq$ scale.
The absolute common kernel at that finer scale remains the
specified original metric quantity to calculate.

\clearpage
\part*{The complete original observation and its exact receiving maps}
The following full proofs retain the actual period contours, amplitude
unit, primary multiplicities and cyclic average. Their only source
orders are the existing $1,k$. The preceding full IGO and GMB inputs
already supply the source transform and original mixed-row metric.
The additional original period and observation proof bodies are
included here with their full source editions preserved alongside them.

% BEGIN RX INLINE 47; SHA256 d43c493eee0fc4e3e69f7383e7520b91900cdb183ed7c159bd9b2b473fdc5b04
% Exact original BRU observation, without a chosen replacement source order.
\section{The original period vectors and their Gamma generating map}
\label{sec:OPG}

The observation in this calculation is the literal BRU1--3 matrix. We
retain its entire amplitude unit, its primary multiplicities, its
period contours, its cyclic projector and its kernel. The only source
orders used are the already present orders $s\in\{1,k\}$. The letter
$s$ in a source-order subscript is an integer parameter; the coordinate
on the one-factor polynomial algebra will be denoted $z_1$ below to
keep it distinct. The generating variable is $z$.

\subsection{The original finite algebras and the observation}
Retain
\[
\begin{gathered}
 \rho=\tfrac12+\delta+i\gamma,\quad 0<\delta<\tfrac12,\quad\gamma>2,
 \quad\mathcal R=\{\rho,\bar\rho,1-\rho,1-\bar\rho\},\\
 h(z_1)=\prod_{\alpha\in\mathcal R}(z_1-\alpha)^m,
 \quad n=4m,\quad v_h=g/h,\quad g=2\xi,\quad m\geq1,\\
 E_h=\mathbb C[z_1]/(h),\qquad A_h=M_{z_1}:E_h\to E_h,
 \qquad U=M_{j_hv_h}:E_h\to E_h.
\end{gathered}\tag{OPG1}
\]
Here $m$ is the original complete zero order at each of the four
points. Consequently $v_h(\alpha)\ne0$ for every $\alpha\in\mathcal R$.
The entire-function remainder $j_hv_h$ records its derivatives of
orders $0,\ldots,m-1$ at every one of these points. In each primary
algebra it has nonzero constant term, and its inverse is the finite
geometric series for that constant plus a nilpotent. Thus $U$ is the
same invertible unit as in BRU, with none of its values prescribed.

For the original positive integer $k$, set
\[
\begin{gathered}
 c=k/2,\quad \ell_k=1+k(m-1),\quad q=(k+1)^2\ell_k,\\
 \chi(S)=\prod_{a,b=0}^k
   \bigl(S-c-(2a-k)\delta-i(2b-k)\gamma\bigr)^{\ell_k},\\
 E=\mathbb C[S]/(\chi),\quad A=M_S:E\to E,\quad
 Y=(A-cI)/i,\\
 A_\Sigma=\sum_{j=1}^k
      I_{E_h}^{\otimes(j-1)}\otimes A_h\otimes I_{E_h}^{\otimes(k-j)},
 \qquad
 \iota:E\hookrightarrow E_h^{\otimes k},\quad
 [P]\longmapsto[P(z_1+\cdots+z_k)].
\end{gathered}\tag{OPG2}
\]
For clarity, the stated map is injective with exactly this defining
polynomial. On an ordered primary tensor the eigenvalue of
$A_\Sigma$ is the sum of the ordered roots. Its nilpotent part is
the sum of the $k$ independent nilpotents of orders $m$.
Its power $k(m-1)$ has nonzero coefficient
$(k(m-1))!/((m-1)!)^k$ on their top monomial, and its next power
is zero. The distinct sums are exactly the $(k+1)^2$ displayed
roots of $\chi$, since the signs of the real and imaginary parts
can be selected independently in every factor. Distinct pairs
$(a,b)$ give distinct sums by $\delta,\gamma>0$. Their common
nilpotency index is $\ell_k$. These facts prove that the minimal
polynomial of $A_\Sigma$ is $\chi$, and therefore prove the
kernel of polynomial evaluation and the asserted injection.
In particular
\[
 \iota A=A_\Sigma\iota,\qquad
 \iota Y=\frac{A_\Sigma-cI}{i}\iota,\qquad
 \iota[1]=[1]^{\otimes k}.
 \tag{OPG3}
\]

Keep the original $u\ne0$, its logarithm branch, and $t=0$ in
the period construction. With $\Phi_h'=h$, $\Phi_h(0)=0$,
let $\ell_j$ be the ray of angle
$(\arg u+\pi+2\pi j)/(n+1)$ and let
$\Gamma_j=-\ell_0+\ell_j$, $1\leq j\leq n$, with the original
orientations. The original period map is
\[
 \Pi:E_h\to\mathbb C^n,\qquad
 (\Pi[P])_j=\int_{\Gamma_j}e^{\Phi_h(z_1)/u}
                    \operatorname{rem}_h P(z_1)\,dz_1.
 \tag{OPG4}
\]
The representative on the right has degree less than $n$. These
are exactly the original convergent polynomial period integrals;
their existence and full determinant are proved in PDE1--13 and
IPM1--7 in the retained source closure. This definition does not
assign a period to an arbitrary representative of a quotient class.
Let $C$ be the original matrix with rows corresponding to
$\Gamma_j\mapsto\Gamma_{j+1}-\Gamma_1$ for $j<n$ and
$\Gamma_n\mapsto-\Gamma_1$. The reduced endpoint basis gives
$C^{n+1}=I$. Hence the literal finite average below is idempotent:
\[
\begin{gathered}
 P_k=\frac1{n+1}\sum_{a=0}^{n}(C^{-a})^{\otimes k},\qquad
 L=P_k\Pi^{\otimes k}U^{\otimes k}\iota,\\
 B=\operatorname{im}L,\quad\jmath:B\hookrightarrow(\mathbb C^n)^{\otimes k},
 \qquad\Lambda:E\twoheadrightarrow B,\quad\jmath\Lambda=L,
 \qquad K=\ker\Lambda.
\end{gathered}\tag{OPG5}
\]
The basis inclusion $\jmath$, the inclusion $I:K\hookrightarrow E$,
and a coefficient section $S_*:B\to E$ with $\Lambda S_*=I_B$
are the fixed BRU choices. The full actual image is used at every
rank. Empty boundary matrices use the original empty conventions.

\subsection{An exact generating identity before and after observation}
Put $b_s=s/2$, $M_s=(2\pi)^{s/2}$, and retain the monic Gamma
polynomials and their full norms from GMB1:
\[
\begin{gathered}
 p_0(y)=1,\quad p_1(y)=y,\quad
 p_{d+1}(y)=yp_d(y)-d(b_s+d-1)p_{d-1}(y),\\
 h_d^{(s)}=M_s d!(b_s)_d,\quad
 r_d=p_d(Y)[1],\quad e_d=r_d/\sqrt{h_d^{(s)}},\quad
 \lambda_d=\Lambda e_d\in B.
\end{gathered}\tag{OPG6}
\]
The square root is the positive real square root. In particular
$e_d=[u_d]$ in the original notation, including the phase $i^{-d}$
in its leading coefficient. Define the holomorphic germs of
$\log(1+z^2)$ and $\arctan z$ by their values zero at $z=0$.
For $|z|<1$, the exact algebra-valued identity is
\[
 \sum_{d=0}^{\infty}r_d\frac{z^d}{d!}
 = (1+z^2)^{-s/4}\exp((\arctan z)Y)[1].
 \tag{OPG7}
\]
Indeed the right side is holomorphic in that disc, is $[1]$ at
zero, and satisfies
$(1+z^2)F'(z)=(Y-b_s zI)F(z)$.
Comparing its Taylor coefficients gives exactly the recurrence
in OPG6: the contribution of $z^2F'$ to coefficient $z^d/d!$
is $d(d-1)r_{d-1}$, and the contribution of $b_szF$ is
$b_sd r_{d-1}$. Thus its derivatives at zero are the $r_d$,
which proves the convergent identity and not merely a formal
coefficient equality.

All factors of $A_\Sigma$ act on different tensor coordinates
and therefore commute. Their matrix exponential power series
are absolutely convergent in finite dimension. The multinomial
expansion of every power consequently proves
\[
 \exp\left(w\frac{A_\Sigma-cI}{i}\right)[1]^{\otimes k}
 =\left(\exp\left(w\frac{A_h-\tfrac12I}{i}\right)[1]\right)^{\otimes k}.
 \tag{OPG8}
\]
The centre on the left is exactly $c=k/2$; each of the $k$
centres on the right is exactly $1/2$. Applying OPG3--5 to OPG7,
and then using OPG8, gives the original observed generating map
\[
\boxed{
 \jmath\sum_{d=0}^{\infty}\sqrt{h_d^{(s)}}\lambda_d\frac{z^d}{d!}
 = (1+z^2)^{-s/4}P_k
 \left[\Pi U\exp\left((\arctan z)
                \frac{A_h-\tfrac12I}{i}\right)[1]\right]^{\otimes k}.
}\tag{OPG9}
\]
The scalar outside $P_k$ has exponent $-s/4$. It has not acquired
an extra factor of $k$: OPG7 concerns one original Gamma source
on the sum coordinate, while OPG8 is the exact factorization of
that coordinate action in the already existing tensor algebra.
For $s=1$ and $s=k$ this gives the two stated source versions of
the same original observation. No additional order is selected.

\subsection{All primary coefficients and the original period integrals}
Here is a finite formula for every vector in OPG9. For
$\alpha\in\mathcal R$ write
\[
\begin{gathered}
 h_\alpha(z_1)=h(z_1)/(z_1-\alpha)^m,\qquad
 t_\alpha(z_1)=\sum_{j=0}^{m-1}
       \frac{(h_\alpha^{-1})^{(j)}(\alpha)}{j!}(z_1-\alpha)^j,\\
 \varepsilon_\alpha=[h_\alpha t_\alpha]\in E_h,
 \qquad E_{\alpha,a}(z_1)=
   \operatorname{rem}_h\bigl(h_\alpha(z_1)t_\alpha(z_1)
                                  (z_1-\alpha)^a\bigr),\quad0\leq a<m,\\
 \mathbf b_{\alpha,a}=\Pi[E_{\alpha,a}],\qquad
 (\mathbf b_{\alpha,a})_j=
   \int_{\Gamma_j}e^{\Phi_h(z_1)/u}E_{\alpha,a}(z_1)\,dz_1.
\end{gathered}\tag{OPG10}
\]
The Taylor inverse exists because $h_\alpha(\alpha)\ne0$.
Reduction modulo $(z_1-\alpha)^m$ gives
$\varepsilon_\alpha=1$ on its primary factor and zero on every
other primary factor. These observations prove
$\varepsilon_\alpha\varepsilon_\beta=0$ for $\alpha\ne\beta$,
$\varepsilon_\alpha^2=\varepsilon_\alpha$, and
$\sum_\alpha\varepsilon_\alpha=1$: their differences are divisible
by all four coprime factors and hence by $h$. They also prove that
$[E_{\alpha,a}]=\varepsilon_\alpha(z_1-\alpha)^a$ is a basis of
$E_h$. Every representative in OPG10 has degree less than $n$,
so each integral is precisely of the original polynomial type.

For $w\in\mathbb C$ put
$\mathbf v(w)=\Pi U\exp(w(A_h-\tfrac12I)/i)[1]$.
The primary nilpotent $N_\alpha=(z_1-\alpha)\varepsilon_\alpha$
satisfies $N_\alpha^m=0$. Taylor multiplication in that algebra
therefore proves, with all amplitude derivatives retained,
\[
\boxed{
 \mathbf v(w)=\sum_{\alpha\in\mathcal R}
 e^{w(\alpha-1/2)/i}\sum_{a=0}^{m-1}\mathbf b_{\alpha,a}
       \sum_{\nu=0}^{a}
       \frac{v_h^{(\nu)}(\alpha)}{\nu!}
       \frac{(w/i)^{a-\nu}}{(a-\nu)!}.
}\tag{OPG11}
\]
For a direct verification before integration, both
$U\exp(w(A_h-\tfrac12I)/i)[1]$ and
$j_h(v_h(z_1)e^{w(z_1-1/2)/i})$ have the same first $m$ Taylor
coefficients at every $\alpha$. Their difference is therefore
zero in $E_h$. The representative to which $\Pi$ is applied is
exactly the finite polynomial in the basis $E_{\alpha,a}$ whose
coefficients are displayed in OPG11. In particular this proof
does not replace the integrand in OPG10 by the unreduced entire
function $v_h(z_1)e^{w(z_1-1/2)/i}$. The precise map from that
entire function to the period vector is its finite Taylor
remainder $j_h$, followed by the degree-less-than-$n$ representative,
followed by the original polynomial period integral.

The cyclic projector is also fully evaluated as a finite operation:
\[
 \jmath\sum_{d\geq0}\sqrt{h_d^{(s)}}\lambda_d\frac{z^d}{d!}
 =\frac{(1+z^2)^{-s/4}}{n+1}
       \sum_{a=0}^{n}\bigl(C^{-a}\mathbf v(\arctan z)\bigr)^{\otimes k}.
 \tag{OPG12}
\]
This is obtained by applying each tensor power in OPG5 to a pure
tensor. It retains every term of the original average, its factor
$1/(n+1)$, and its inverse powers of $C$.

There is also a finite coefficient formula requiring no series
extraction. For an ordered $k$-tuple
$\boldsymbol\alpha\in\mathcal R^k$ set
$\omega_{\boldsymbol\alpha}=(\sum_i\alpha_i-c)/i$.
For $0\leq a_i<m$ and $0\leq\nu_i\leq a_i$ put
$D(\mathbf a,\boldsymbol\nu)=\sum_i(a_i-\nu_i)$ and
\[
 T_d(\boldsymbol\alpha,\mathbf a)=
 \sum_{\substack{0\leq\nu_i\leq a_i\\1\leq i\leq k}}
 \frac{i^{-D(\mathbf a,\boldsymbol\nu)}
        p_d^{(D(\mathbf a,\boldsymbol\nu))}(\omega_{\boldsymbol\alpha})}
        {\prod_{i=1}^k(a_i-\nu_i)!}
 \prod_{i=1}^k\frac{v_h^{(\nu_i)}(\alpha_i)}{\nu_i!}.
 \tag{OPG13}
\]
As usual the polynomial derivative is zero when its order exceeds
$d$. Multivariate Taylor expansion in the independent nilpotents
$(z_i-\alpha_i)$ proves the exact formula
\[
\boxed{
 \jmath\lambda_d=
 \frac1{\sqrt{h_d^{(s)}}}\frac1{n+1}
 \sum_{a=0}^{n}\ \sum_{\boldsymbol\alpha\in\mathcal R^k}
 \sum_{\substack{0\leq a_i<m\\1\leq i\leq k}}
 T_d(\boldsymbol\alpha,\mathbf a)
 \bigotimes_{i=1}^k C^{-a}\mathbf b_{\alpha_i,a_i}.
}\tag{OPG14}
\]
Indeed on that primary tensor, expand each $v_h(z_i)$ first.
For the remaining multi-degree $a_i-\nu_i$, Taylor expansion of
$p_d((\sum_i z_i-c)/i)$ supplies its derivative of the total
order $D$, the factor $i^{-D}$, and the factorials
$\prod_i(a_i-\nu_i)!$. These are exactly the coefficients in
OPG13. Apply the $k$ polynomial period maps and the literal cyclic
average to get OPG14. All sums are finite. Ordered tuples have
not been replaced by just their eigenvalue sum, so both the
amplitude data and every independent primary nilpotent survive.

\subsection{The actual observation increment without a coordinate inverse}
Retain the GMB maps
$C_s=[e_0,\ldots,e_{q-1}]:\mathbb C^q\to E$,
$v_j=C_s^{-1}e_{q+j}$,
$K_j=I_q+\sum_{a<j}v_av_a^*$, and $R_s=\Lambda C_s$.
The low-degree leading coefficients prove $C_s$ invertible as in
GMB2. On each column, applying the maps gives
\[
 R_s v_j=\lambda_{q+j},\qquad
 R_sK_jR_s^*=\sum_{d=0}^{q+j-1}\lambda_d\lambda_d^*.
 \tag{OPG15}
\]
For the second equality the identity term in $K_j$ gives exactly
the first $q$ columns of $R_s$, and each remaining rank-one term
gives the next displayed column. The low vectors $\lambda_d$,
$0\leq d<q$, span $B$ because $C_s$ is invertible and $\Lambda$
is onto. Thus, in the fixed original basis of $B$, the matrix
\[
 \mathcal M_D=\sum_{d=0}^{D-1}\lambda_d\lambda_d^*\quad(D\geq q)
 \tag{OPG16}
\]
is positive definite whenever $B$ has positive dimension:
for a nonzero covector its quadratic form is the sum of its
squared values on spanning vectors. Consequently the exact
original boundary increment of GMB8 is
\[
 \boxed{\displaystyle
 \xi_j=\lambda_{q+j}^*\mathcal M_{q+j}^{-1}\lambda_{q+j},\qquad
 1+\xi_j=\frac{\det\mathcal M_{q+j+1}}{\det\mathcal M_{q+j}}.}
 \tag{OPG17}
\]
The determinant equality follows by conjugating the rank-one
update by $\mathcal M_{q+j}^{-1/2}$; its sole possible nonzero
added eigenvalue is the displayed quadratic form. If $B=0$,
the vector is empty, $\xi_j=0$ and both determinants equal one.
Thus OPG11--14 compute the boundary increment from the original
period vectors, with no $C_s^{-1}$ in OPG17. This cancellation
is an identity of the proved maps, rather than a change of the
observation or the source measure.

At the four original endpoints the boundary contribution is
therefore the exact finite expression
\[
 F_B^{(s)}=
 \log\frac{\det\mathcal M_{2q}\det\mathcal M_{2q+1}}
                 {\det\mathcal M_q\det\mathcal M_{q+1}}
 =\log(1+\xi_0)+2\sum_{j=1}^{q-1}\log(1+\xi_j)
                         +\log(1+\xi_q).
 \tag{OPG18}
\]
Expanding the last sum of consecutive log determinant differences
leaves coefficient $-1$ at $\mathcal M_q,\mathcal M_{q+1}$ and
$+1$ at $\mathcal M_{2q},\mathcal M_{2q+1}$, proving every sign.
GMB13 proves that the other original factor is
$(1+t_j)/(1+\xi_j)$, with $t_j=v_j^*K_j^{-1}v_j$ and
$0\leq\xi_j\leq t_j$. Thus OPG18 is the actual BRU boundary
contribution, and the full original kernel contribution remains
$\mathcal W_q(\log((1+t_j)/(1+\xi_j)))$.

\subsection{A finite recurrence retaining the exact action defect}
The period formula also has a finite recurrence on the fixed
original coefficient spaces. Define the uniquely typed projection
$\kappa:E\to K$ by
$I\kappa=I_E-S_*\Lambda$. It exists because the right side has
image in $\ker\Lambda=\operatorname{im}I$, and $I$ is an
inclusion. It satisfies $\kappa I=I_K$ and $\kappa S_*=0$.
In particular $E=I(K)\oplus S_*(B)$ as coefficient spaces, without
an assertion that this splitting is preserved by $Y$.
Set
\[
\begin{gathered}
 \mathsf Y_B=\Lambda YS_*:B\to B,\quad
 \mathsf D=\Lambda YI:K\to B,\\
 \mathsf C=\kappa YS_*:B\to K,\quad
 \mathsf Y_K=\kappa YI:K\to K,\quad
 \zeta_d=\kappa e_d\in K.
\end{gathered}\tag{OPG19}
\]
All four maps are computable from OPG1--5 and the fixed BRU
section. For $d\geq0$ put
\[
 a_d=\frac1{\sqrt{(d+1)(b_s+d)}},\qquad
 b_d=\sqrt{\frac{d(b_s+d-1)}{(d+1)(b_s+d)}}\quad(d\geq1),
 \qquad b_0=0.
 \tag{OPG20}
\]
Set $\lambda_{-1}=0$ and $\zeta_{-1}=0$. Then the exact coupled
recurrence is
\[
\boxed{\begin{aligned}
 \lambda_{d+1}&=a_d(\mathsf Y_B\lambda_d+\mathsf D\zeta_d)
                                          -b_d\lambda_{d-1},\\
 \zeta_{d+1}&=a_d(\mathsf C\lambda_d+\mathsf Y_K\zeta_d)
                                          -b_d\zeta_{d-1},\\
 \lambda_0&=\Lambda[1]/\sqrt{M_s},\qquad
 \zeta_0=\kappa[1]/\sqrt{M_s}.
\end{aligned}}\tag{OPG21}
\]
To prove this, the norm ratio from OPG6 gives
$\sqrt{h_d^{(s)}/h_{d+1}^{(s)}}=a_d$ and
$d(b_s+d-1)\sqrt{h_{d-1}^{(s)}/h_{d+1}^{(s)}}=b_d$.
Hence $e_{d+1}=a_dYe_d-b_de_{d-1}$, including $d=0$ by its
initial values. Substituting $e_d=S_*\lambda_d+I\zeta_d$ and
applying $\Lambda$ and $\kappa$ proves both lines of OPG21.
It retains the kernel coordinate at every degree and supplies
the period vectors through degree $2q$, exactly the range used
in OPG17--18.

The failure of a putative boundary action has the precise map
\[
 \Lambda Y-\mathsf Y_B\Lambda=\mathsf D\kappa:E\to B,\qquad
 \jmath\mathsf D=
 \frac1iP_k\Pi^{\otimes k}U^{\otimes k}
                   (A_\Sigma-cI)\iota I:K\to(\mathbb C^n)^{\otimes k}.
 \tag{OPG22}
\]
The first equality follows from $I_E=S_*\Lambda+I\kappa$;
the second follows from OPG3--5. These equations preserve the
cyclic average and the full unit. They also prove exactly that
$Y(K)\subseteq K$ holds if and only if $\mathsf D=0$.
No such vanishing is used anywhere in OPG7--21. Thus every
possible boundary action defect is retained inside a proved
coupled action and inside the explicit original period formula.

For completeness the multiplication action does descend to a
canonical quotient that remembers every observed iterate. Define
\[
 \mathcal K_\infty=\bigcap_{a=0}^{q-1}\ker(\Lambda Y^a),\qquad
 \mathcal O:E\to B^{\oplus q},\quad
 x\longmapsto(\Lambda x,\Lambda Yx,\ldots,\Lambda Y^{q-1}x).
 \tag{OPG23}
\]
The polynomial $\chi(c+iT)$ annihilates $Y$ by OPG2, has degree
$q$ and leading coefficient $i^q\ne0$. It therefore expresses
$Y^q$ as a linear combination of $I,Y,\ldots,Y^{q-1}$, with
the exact coefficients of this original polynomial. If
$x\in\mathcal K_\infty$, each component of $\mathcal O(Yx)$
vanishes by this relation. Thus $\mathcal K_\infty$ is
$Y$-invariant and is precisely $\bigcap_{a\geq0}\ker(\Lambda Y^a)$.
Moreover any $Y$-invariant subspace contained in $K$ lies in
this intersection. Consequently the first isomorphism map
\[
\begin{gathered}
 E/\mathcal K_\infty\ \xrightarrow{\ \overline{\mathcal O}\ }
 \operatorname{im}\mathcal O,\quad[x]\longmapsto\mathcal O(x),\\
 \chi_j=[T^j]\chi(c+iT),\quad\chi_q=i^q,\\
 \mathsf{Shift}(b_0,\ldots,b_{q-1})=
 \left(b_1,\ldots,b_{q-1},-i^{-q}\sum_{j=0}^{q-1}\chi_j b_j\right).
\end{gathered}
 \tag{OPG24}
\]
is a bijection intertwining multiplication by $Y$ with the shift
of the $q$ components, whose last component is the stated exact
linear combination. It projects to the original observation by
the first component, with retained kernel $K/\mathcal K_\infty$.
This proves an action-preserving enlargement of the observed
data on the original algebra. It neither discards the original
kernel in OPG21 nor presumes a vanished defect in OPG22.

The original period isomorphism also locates that defect in the
cyclic projection. Its inverse exists by the retained period
determinant theorem used in OPG4. On the unchanged ambient tensor
space put
\[
 \mathcal V=\Pi^{\otimes k}U^{\otimes k}\iota,
 \qquad
 \mathcal T=\sum_{j=1}^k
 \left(\Pi\frac{A_h-\tfrac12I}{i}\Pi^{-1}\right)_j.
 \tag{OPG25}
\]
The subscript $j$ means that the displayed factor acts in position
$j$ and identities act in all other positions. Multiplication by
$j_hv_h$ commutes with $A_h$. Equations OPG3 and OPG25 therefore
give $\mathcal VY=\mathcal T\mathcal V$ exactly. Since
$L=P_k\mathcal V$ and $P_k^2=P_k$, splitting the ambient identity
as $P_k+(I-P_k)$ now proves
\[
\begin{gathered}
 LY=P_k\mathcal T L+P_k\mathcal T(I-P_k)\mathcal V,\\
 \jmath\mathsf D=P_k\mathcal T(I-P_k)\mathcal V I.
\end{gathered}\tag{OPG26}
\]
For the second equality, $P_k\mathcal V I=LI=0$, so the
first term of the first equality vanishes on $I(K)$. This is
the same defect as OPG22 on its original domain. The two ambient
terms in the first line sum to a vector in $\jmath(B)$; neither
individual term is asserted to preserve that image. OPG19--21
give the full recurrence within $K\oplus B$ without such an
assertion. Thus the preserved unit and period isomorphism carry
the multiplication action exactly, and OPG26 identifies precisely
the retained contribution of the original cyclic projection.

Equations OPG9, OPG11--14, OPG17 and OPG21 now give the same
original boundary increments by analytic generating functions,
finite primary coefficients, and finite coupled recurrences.
Their derivations establish the exact maps connecting these
calculations. In particular the new input into the original
mixed-control problem is the fully specified sequence
$\lambda_0,\ldots,\lambda_{2q}$ with its actual unit and periods,
rather than a freely selected observation or a source-order
replacement. The formulas themselves do not assign a growing-
parameter estimate to the individual weighted sums; they supply
their unchanged mathematical input for that calculation.

% END RX INLINE 47


% BEGIN RX INLINE 48; SHA256 b7de12c81133c7fb75e1bd3e46be3ca591988692ff6f6b8ef6ddecfbd4bfa92e
\section{The exact receiving map into the original mixed factors}
\label{sec:OPR}

The current receiver MRI1--5 is retained as a frozen predecessor.
This section proves its next observation input from OPG1--26.
Fix either original source order $s\in\{1,k\}$, put $D=q+j$,
and retain all the fixed coefficient maps of OPG and GMB.
In particular $e_D=[u_D]$, $\lambda_D=\Lambda e_D$, and
$\mathcal M_D=\sum_{d<D}\lambda_d\lambda_d^*$ on the original
boundary image. At the preceding polynomial degree $D-1$, GMB4
and OPG15 give
\[
\begin{gathered}
 G_{D-1}^{(s)}=C_s^{-*}K_j^{-1}C_s^{-1},\quad
 (G_{D-1}^{(s)})^{-1}=C_sK_jC_s^*,\\
 Q_{D-1}^{(s)}=\mathcal M_D^{-1},\qquad
 S_{D-1}^{(s)}=C_sK_jC_s^*\Lambda^*\mathcal M_D^{-1}:B\to E.
\end{gathered}\tag{OPR1}
\]
Indeed $\Lambda(G_{D-1}^{(s)})^{-1}\Lambda^*=
R_sK_jR_s^*=\mathcal M_D$. Direct multiplication therefore
proves $\Lambda S_{D-1}^{(s)}=I_B$ and
$I^*G_{D-1}^{(s)}S_{D-1}^{(s)}=0$. Thus this is precisely
the original minimum section for this degree and this source.
All statements use their unique empty meaning if $B=0$.

MRI2 uses the kernel vector
$r_j^{\rm ker}=v_j-K_jR_s^*\mathcal M_D^{-1}\lambda_D$
in the $C_s$ coordinates. The fixed-section kernel vector of
OPG21 is $\zeta_D=\kappa e_D$, with $I\kappa=I_E-S_*\Lambda$.
The exact morphism between these two specified vectors is
\[
\boxed{\begin{aligned}
 C_s r_j^{\rm ker}
   &=e_D-S_{D-1}^{(s)}\lambda_D\\
   &=I\zeta_D+(S_*-S_{D-1}^{(s)})\lambda_D
    =I\vartheta_{D,j}^{(s)},\\
 \vartheta_{D,j}^{(s)}
   &=\zeta_D-\kappa S_{D-1}^{(s)}\lambda_D\in K.
\end{aligned}}\tag{OPR2}
\]
For the first equality substitute $v_j=C_s^{-1}e_D$ and
$R_s^*=C_s^*\Lambda^*$ into MRI2. For the second substitute
$e_D=I\zeta_D+S_*\lambda_D$. The difference of the two
sections lands in $K$ because their observations are both $I_B$.
Applying $\kappa$ and using $\kappa S_*=0$ proves the last
line and hence the third equality. This proves the complete
section correction in the original coefficient space.

Retain the actual kernel Gram
$H_{D-1}^{(s),K}=I^*G_{D-1}^{(s)}I$. Congruencing OPR1 by
$C_s$ and then using OPR2 proves
\[
 t_j-\xi_j
 = (r_j^{\rm ker})^*K_j^{-1}r_j^{\rm ker}
 = (\vartheta_{D,j}^{(s)})^*
                    H_{D-1}^{(s),K}\vartheta_{D,j}^{(s)}.
 \tag{OPR3}
\]
The first equality is the proved MRI2/GMB10 orthogonal
decomposition. Thus the recurrence OPG21 supplies both the
actual boundary column and, through the complete section
correction OPR2, the exact kernel norm needed at each step.
The original ordered factors are consequently
\[
\begin{gathered}
 \xi_j=\lambda_D^*\mathcal M_D^{-1}\lambda_D,\qquad
 1+\frac{\nu_{D-1}}{a_{D-1}+\beta_{D-1}}=1+\xi_j,\\
 1+\frac{\beta_{D-1}}{a_{D-1}}
 =\frac{1+t_j}{1+\xi_j}
 =1+\frac{(\vartheta_{D,j}^{(s)})^*
                  H_{D-1}^{(s),K}\vartheta_{D,j}^{(s)}}{1+\xi_j}.
\end{gathered}\tag{OPR4}
\]
The first line uses OPG17 and GMB13. The second line follows
by inserting OPR3 into that same GMB13 identity. This preserves
the denominator $1+\xi_j$ and the original order of elimination.

The finite endpoint operation remains
$\mathcal W_q(f)=f_0+2\sum_{j=1}^{q-1}f_j+f_q$.
Applying it to the logarithms in OPR4 reproduces MRI4 on its
unchanged source. The exact arithmetic receiving equality is
therefore still
\[
\begin{aligned}
 \mathcal B_k^{\rm ar}
 &=F_K^{(k)}+F_B^{(k)}+\mathcal H_k+\delta_k^\sigma\\
 &=F_K^{(1)}+F_B^{(1)}+\delta_k^\sigma.
\end{aligned}\tag{OPR5}
\]
Its proof is substitution into the existing MRI5/BRI2 source
identity: $F_K^{(k)}+F_B^{(k)}=\mathcal B_k^0$,
$F_K^{(1)}+F_B^{(1)}=\mathcal B_k^\sigma$,
$\mathcal H_k=\mathcal B_k^\sigma-\mathcal B_k^0$, and
$\delta_k^\sigma=\mathcal B_k^{\rm ar}-\mathcal B_k^\sigma$.
No term is assigned a new sign or asymptotic coefficient.
The combined coefficient received by BRI6 and BRI9 remains
attached to this same sum. OPR1--4 provide its individual
original factors from the actual OPG observation vectors,
with all original source transitions present in OPR5.

% END RX INLINE 48


% BEGIN RX INLINE 49; SHA256 8d5957fb08f2e1d6448a8130dc4cc5ff530cfae2791ef15d24fd274fceaf4542
\section{The full norms of the two existing Gamma sources}
\label{sec:OGN}

The included canonical source proof IGO1--3 defines
\[
 dm_{0,s}(y)=M_s\frac{2^{b_s}}{4\pi\Gamma(b_s)}
       |\Gamma(b_s/2+iy/2)|^2\,dy,
 \qquad s\in\{1,k\},\quad b_s=s/2,\quad M_s=(2\pi)^{s/2},
 \tag{OGN1}
\]
and proves its exact Laplace transform
$\int e^{ty}dm_{0,s}=M_s(\cos t)^{-b_s}$ for real $|t|<\pi/2$.
Its proof includes exponential domination on each smaller strip,
so all the differentiated integrals below converge absolutely.
Every source mass and the original line $S=c+iy$, $c=k/2$, are
retained here.

Let $p_d$ be the recurrence polynomials of OPG6, and put
$G(z,y)=(1+z^2)^{-b_s/2}e^{y\arctan z}$ near zero. The
coefficient calculation in OPG7, with scalar $y$ in place of its
matrix, proves $G(z,y)=\sum_d p_d(y)z^d/d!$.
Each $p_d$ is real and monic by the recurrence and its initial
values. For sufficiently small real $z,w$, the elementary identity
\[
 \cos(\arctan z+\arctan w)
       =\frac{1-zw}{\sqrt{(1+z^2)(1+w^2)}}
\]
and the full Laplace transform give
\[
 \int_{\mathbb R}G(z,y)G(w,y)\,dm_{0,s}(y)
                  =M_s(1-zw)^{-b_s}.
 \tag{OGN2}
\]
Choose the neighbourhood so that
$|\arctan z|+|\arctan w|\leq a<\pi/2$.
Every fixed mixed derivative of the integrand is bounded on a
smaller neighbourhood by a constant times a polynomial in $|y|$
times $e^{a|y|}$. The exponential domination proved in IGO3,
with a strictly larger exponent smaller than $\pi/2$, integrates
this bound. Differentiation under the integral is therefore valid
at every fixed order. Applying $\partial_z^d\partial_w^e$ at zero
to OGN2 yields
\[
 \int p_d(y)p_e(y)\,dm_{0,s}(y)
  =\begin{cases}0,&d\ne e,\\
        M_s d!(b_s)_d,&d=e.
    \end{cases}
 \tag{OGN3}
\]
Indeed $(1-zw)^{-b_s}=\sum_{a\geq0}(b_s)_a(zw)^a/a!$,
so its $(d,e)$ derivative at zero vanishes off the diagonal and
equals $d!(b_s)_d$ on it. Since $b_s>0$, every diagonal entry
is strictly positive. Consequently
$u_d(S)=p_d((S-c)/i)/\sqrt{M_s d!(b_s)_d}$ is exactly the
orthonormal source polynomial used in GMB1 and OPG6, with
leading coefficient $i^{-d}/\sqrt{M_s d!(b_s)_d}$ in $S$.
This proves those source norms directly from the two retained
source measures and fixes the full phase and mass of each
observation column.

% END RX INLINE 49


% BEGIN RX INLINE 50; SHA256 517daf9230b256ed0fad88e5104ba9fd8f8d51e793998fea955c1a34bacf480d
\section{The boundary and its retained kernel under each original relation}
\label{GraphV3:reader-v3-BRU:sec:bru-original-relations}

This calculation continues the incoming boundary observation A1 and the
boundary/kernel metric A4--A18. It calculates those metrics and every
cross-degree overlap from the original relation rows. All maps below are
on the actual packet, before taking a completion. In particular the
boundary observation does not remove its kernel.

\subsection{The packet, observation, source and finite relation rows}

Retain an actual hypothetical complete-order quartet
\[
 \rho=\tfrac12+\delta+i\gamma,\quad 0<\delta<\tfrac12,\quad\gamma>2,
 \quad h(s)=\prod_{r\in\{\rho,\bar\rho,1-\rho,1-\bar\rho\}}(s-r)^m,
 \quad n=4m,
 \tag{BRU1}
\]
where $m\geq1$ is its complete zero order in $g=2\xi$.
Thus $v_h=g/h$ is the original entire amplitude and
$U=M_{j_hv_h}$ is its invertible multiplication matrix on
$E_h=\mathbb{C}[s]/(h)$. No value of this unit is assigned.
For $k\geq1$ the original cyclic coefficient space and action are
\[
 \begin{gathered}
 \ell_k=1+k(m-1),\qquad q=(k+1)^2\ell_k,\qquad E=\mathbb{C}[S]/(\chi),\\
 \chi(S)=\prod_{a,b=0}^k
 (S-k/2-(2a-k)\delta-i(2b-k)\gamma)^{\ell_k},\qquad
 A=M_S,\\
 \iota:E\longrightarrow E_h^{\otimes k},\qquad
 [P]\longmapsto[P(s_1+\cdots+s_k)],\qquad
 \eta=U^{\otimes k}\iota .
 \end{gathered}\tag{BRU2}
\]
The original cyclic inclusion is injective: on each ordered primary
tensor the sum of the nilpotents has nilpotency index
$1+k(m-1)$, since its power $k(m-1)$ has nonzero coefficient
$(k(m-1))!/((m-1)!)^k$ on the product of the top powers. Its next
power is zero by total degree. All sums in BRU2 are distinct for
different $(a,b)$ because $\delta,\gamma>0$. Each occurs among the
ordered factors, so their product with exactly these exponents is
the minimal polynomial of the tensor sum. This proves the stated
kernel of polynomial evaluation and injectivity of $\iota$.

Fix $u\ne0$ on the original logarithm branch and $t=0$. Let
$\Pi$ be the original period matrix with rays
$\Gamma_j=-\ell_0+\ell_j$, angles
$(\arg u+\pi+2\pi j)/(n+1)$, and entries
$\int_{\Gamma_j}e^{\Phi_h(s)/u}s^b\,ds$, where $\Phi_h'=h$ and
$\Phi_h(0)=0$. Its full determinant and convergence are proved in
PDE1--13 and IPM1--7, included in the source package. Let $C$ have
rows describing $\Gamma_j\mapsto\Gamma_{j+1}-\Gamma_1$ for $j<n$
and $\Gamma_n\mapsto-\Gamma_1$. Set
\[
 P_k=\frac1{n+1}\sum_{a=0}^{n}(C^{-a})^{\otimes k},\qquad
 L=P_k\Pi^{\otimes k}\eta:E\longrightarrow(\mathbb{C}^n)^{\otimes k}.
 \tag{BRU3}
\]
For the algebra used here, $C^{n+1}=I$ follows by applying the
cyclic permutation to the reduced endpoint basis $e_j-e_0$.
Multiplying the finite sum therefore gives $P_k^2=P_k$.
We use the literal coefficient map BRU3, without evaluating a
formal stationary series or identifying its action with Frobenius.
Choose a basis inclusion $\jmath:B=\operatorname{im}L\hookrightarrow
(\mathbb{C}^n)^{\otimes k}$, and define the surjection
$\Lambda:E\to B$ by $\jmath\Lambda=L$. Its kernel is $K$.
Fix a basis inclusion $I:K\hookrightarrow E$, and a coefficient
section $S_*$ with $\Lambda S_*=I_B$. Gaussian elimination on the
specified finite matrix constructs these maps, including at a
rank drop by using its actual image there. Keep these choices fixed
as the source metric varies. Write $r=\dim B$ and $d_K=q-r$.
Zero-dimensional factors use their unique empty inverse and determinant one.


% END RX INLINE 50


% BEGIN RX INLINE 51; SHA256 347a84b27a1b7cc9fee7f5c4a155400c1e3030a7a1f69c84400438e9fb21d772
\section{The actual entire theta amplitude and the finite period observation}
\label{sec:ipm-entire-source}

This calculation uses the actual packet
\(h(s)=\prod_\rho(s-\rho)^{m_\rho}\), with all complete orders,
\(d_h=\deg h\ge1\), \(\Phi_h'=h\), \(\Phi_h(0)=0\),
\(g=2\xi\), \(v_h=g/h\), and
\(\upsilon_h=j_hv_h\in E_h^\times\), where
\(E_h=\mathbb C[s]/(h)\). No root outside this actual packet is
introduced. The analytic statements have these data as their domain;
they do not assert that an off-critical packet has been found.
The finite-source map is the original injective map
\(\mathcal T_hP=P(D)F_h\), with Mellin transform \(v_hP\).
Thus all maps below on \(\mathcal T_h\mathcal P_N\) are defined
unambiguously through that injective source parametrization.

\subsection{An explicit growth bound proves the global integration domain}

Write
\[
 \psi(x)=\sum_{n\ge1}e^{-\pi n^2x},\qquad
 C_\theta=(1-e^{-3\pi})^{-1},\qquad
 \beta_R=(R+1)/2.
 \tag{IPM.1}
\]
The original theta representation, also recorded in DLMF 25.5.13--14,
gives the entire identity
\[
 g(s)=1+s(s-1)\int_1^\infty
       \bigl(x^{s/2}+x^{(1-s)/2}\bigr)\psi(x)\,\frac{dx}{x}.
 \tag{IPM.2}
\]
For completeness, the identity follows initially for \(\Re s>1\)
by integrating each Gaussian:
\(\int_0^\infty\psi(x)x^{s/2}dx/x
=\pi^{-s/2}\Gamma(s/2)\zeta(s)\).
Split the integral at one and use the original Poisson identity
\(\psi(x)=x^{-1/2}\psi(1/x)+(x^{-1/2}-1)/2\).
The elementary integral on \((0,1)\) is
\(1/(s-1)-1/s=1/(s(s-1))\). Multiplication by \(s(s-1)\)
proves (IPM.2) in that half-plane. Its right side is entire, since its
integral and every coefficient derivative converge uniformly on compact
sets by Gaussian decay. Analytic continuation proves the identity
everywhere, including \(s=0,1\).

Since \(n^2\ge1+3(n-1)\), for \(x\ge1\)
one has \(\psi(x)\le C_\theta e^{-\pi x}\). Consequently
\[
 \begin{gathered}
 |g(s)|\le\mathcal X(R):=
  1+2C_\theta R(R+1)\pi^{-\beta_R}\Gamma(\beta_R,\pi)\\
 \le 1+2C_\theta R(R+1)\pi^{-\beta_R}\Gamma(\beta_R),
 \qquad |s|\le R,\ R\ge1.
 \end{gathered}
 \tag{IPM.3}
\]
Indeed both powers of \(x\) in the integral have absolute value at most
\(x^{(R+1)/2}\), and \(|s(s-1)|\le R(R+1)\).
The upper incomplete gamma is the literal integral after \(z=\pi x\).
An elementary bound sufficient here is
\(\Gamma(\beta)\le2(2\beta)^\beta\) for \(\beta\ge1\).
To prove it, write its integrand as
\((t^{\beta-1}e^{-t/2})e^{-t/2}\), bound the first factor by
its maximum \([2(\beta-1)/e]^{\beta-1}\), and integrate the second.
At \(\beta=1\) use the value one for the zeroth power. This maximum
is at most \((2\beta)^\beta\). In particular
\(\log\mathcal X(R)=O(R\log(R+2))\); the explicit function
\(\mathcal X\), rather than a substituted constant, is retained in
all tail integrals.

Let \(R_h=\max_\rho|\rho|\). For \(|s|=r\ge\max(1,2R_h)\),
\[
 |h(s)|\ge(r/2)^{d_h},\qquad
 |v_h(s)|\le 2^{d_h}r^{-d_h}\mathcal X(r).
 \tag{IPM.4}
\]
The complete-divisor hypothesis makes \(v_h\) entire, so it is bounded
on the remaining compact disc. Cauchy's integral formula on radius-one
circles gives the same \(\exp(O(r\log(r+2)))\) growth for each fixed
derivative. The formula retains both the degree and every original root
in its compact-disc and exterior bounds.

Fix \(u\ne0\), \(t\in\mathbb C\), a branch of \(\arg u\), and
the original BC rays and oriented contours
\[
 \theta_j=\frac{\arg u+\pi+2\pi j}{d_h+1},\quad
 \ell_j(r)=re^{i\theta_j},\quad
 \Gamma_j=-\ell_0+\ell_j\quad(1\le j\le d_h),\quad
 f_t=\Phi_h-ts.
 \tag{IPM.5}
\]
Their leading real exponent is
\(-r^{d_h+1}/((d_h+1)|u|)\). If
\(f_t(s)=s^{d_h+1}/(d_h+1)+\sum_{\nu=0}^{d_h}c_\nu s^\nu\),
then, for \(r\ge1\), its full real exponent is at most
\[
 -\frac{r^{d_h+1}}{(d_h+1)|u|}
       +\frac{\sum_{\nu=0}^{d_h}|c_\nu|r^\nu}{|u|}.
 \tag{IPM.6}
\]
Combining (IPM.3)--(IPM.6) proves absolute convergence of
\[
 \mathcal I_{u,t}(a)=
  \left(\int_{\Gamma_j}e^{f_t(s)/u}a(s)\,ds\right)_{j=1}^{d_h}
 \tag{IPM.7}
\]
for \(a=v_hP\), every polynomial \(P\), every derivative used below,
and every entire quotient by \(h\) constructed below. On a compact
parameter set with \(u\ne0\), the same domination proves coefficient
differentiation and the needed contour deformations inside the fixed
decaying sectors. This is a proof of the domain for the actual theta
amplitude. It makes no assertion for arbitrary undeclared rays.


% END RX INLINE 51


% BEGIN RX INLINE 52; SHA256 4df10819db4933332fc6d72992b26b1c9a9f0da11e1d8c813c1c0611e31a9f9c
% Exact scalar dependency used in PDE9--11 and IPM.20--21.
\section{The positive Gamma grid product in the original period determinant}
\label{sec:PGCC}

The full original PDE determinant proof reduces its constant to
\[
 \prod_{r=1}^{d-1}\Gamma(r/d)
       =\frac{(2\pi)^{(d-1)/2}}{\sqrt d},\qquad d=n+1\geq2.
 \tag{PGCC1}
\]
Here is a complete proof of this particular scalar identity from the
positive defining Gamma integral, retaining its positive square root.
For a real $0<a<1$, Tonelli's theorem applies to the nonnegative
integrand defining the product of two Gamma integrals. Set
$x=st/(1+t)$ and $y=s/(1+t)$, with $s,t>0$. This is a bijective
change of variables on the positive quadrant with Jacobian
$s/(1+t)^2$, and direct multiplication gives
\[
 \begin{split}
 \Gamma(a)\Gamma(1-a)
 &=\int_0^\infty\int_0^\infty
        e^{-x-y}x^{a-1}y^{-a}\,dx\,dy\\
 &=\int_0^\infty e^{-s}\,ds
       \int_0^\infty\frac{t^{a-1}}{1+t}\,dt
  =\int_0^\infty\frac{t^{a-1}}{1+t}\,dt.
 \end{split}\tag{PGCC2}
\]
The last integral converges at zero because $a>0$ and at infinity
because $a<1$.

To evaluate it with its exact sign, use
$z^{a-1}=\exp((a-1)\operatorname{Log}z)$ with $0<\arg z<2\pi$.
The positively oriented keyhole contour about the positive real axis
has an upper segment from $\varepsilon$ to $R$, an outer circle
counterclockwise, a lower segment from $R$ to $\varepsilon$, and an
inner circle clockwise. Take $0<\varepsilon<1<R$ and let the banks
approach the cut. The only pole of $z^{a-1}/(1+z)$ inside is $z=-1$,
with residue $e^{\pi i(a-1)}$. The two straight segments sum to
\[
 (1-e^{2\pi ia})\int_\varepsilon^R
                       \frac{t^{a-1}}{1+t}\,dt.
 \tag{PGCC3}
\]
The outer circular contribution has absolute value at most
$2\pi R^a/(R-1)$, which tends to zero as $R\to\infty$.
The inner contribution has absolute value at most
$2\pi\varepsilon^a/(1-\varepsilon)$, which tends to zero as
$\varepsilon\to0$. The residue theorem and these bounds give
\[
 \begin{split}
 (1-e^{2\pi ia})\int_0^\infty\frac{t^{a-1}}{1+t}\,dt
 &=2\pi i e^{\pi i(a-1)},\\
 \Gamma(a)\Gamma(1-a)&=\frac{\pi}{\sin(\pi a)}.
 \end{split}\tag{PGCC4}
\]
In the second line the quotient sign follows explicitly from
$1-e^{2\pi ia}=-2i e^{\pi ia}\sin(\pi a)$ and
$e^{\pi i(a-1)}=-e^{\pi ia}$.

Now put $\zeta=e^{2\pi i/d}$. Factoring $X^d-1$ over its distinct
roots and evaluating the quotient by $X-1$ at $X=1$ proves
\[
 d=\prod_{r=1}^{d-1}(1-\zeta^r),\qquad
 d=\prod_{r=1}^{d-1}|1-\zeta^r|
   =2^{d-1}\prod_{r=1}^{d-1}\sin(\pi r/d).
 \tag{PGCC5}
\]
The sines in this product are all positive. Apply PGCC4 at
$a=r/d$ and multiply over $1\leq r<d$. The involution
$r\mapsto d-r$ permutes the same factors, so
\[
 \left(\prod_{r=1}^{d-1}\Gamma(r/d)\right)^2
 =\prod_{r=1}^{d-1}\frac{\pi}{\sin(\pi r/d)}
 =\frac{(2\pi)^{d-1}}d.
 \tag{PGCC6}
\]
Each Gamma integral is strictly positive. Taking the positive square
root proves PGCC1, including the case $d=2$.

This evaluates the precise scalar used in PDE9--11 without replacing
the original Fourier determinant or its phase. In fact the polynomial
period isomorphism already follows before this evaluation: PDE9 gives
nonzero positive Gamma factors and fixed nonzero powers of $u$, while
PDE10 gives the nonzero Fourier Vandermonde determinant. PDE6--8 then
propagate that nonzero determinant through the full coefficient family.
PGCC1 additionally supplies its displayed value $(2\pi)^{n/2}$,
with all original powers and phases calculated in PDE9--11 unchanged.

% END RX INLINE 52


% BEGIN RX INLINE 53; SHA256 bc8cc467be1b3858772131152cb2823968f886d9aae34baa845e71483cec97bb
\section{The full polynomial period deformation and its entire-amplitude forcing}
\label{GraphV3:sec:pde-full-deformation}

Retain the original nonempty complete packet $h(s)$, its degree $n\geq1$,
its primitive $\Phi_h(0)=0$, its arithmetic coefficient space
$E_h=\mathbb{C}[s]/(h)$ in the increasing frame $1,s,\ldots,s^{n-1}$,
and the actual entire amplitude $v_h=2\xi/h$. All zero orders in $h$
are complete, so this last quotient is entire and its Taylor unit
$\upsilon_h=j_h(v_h)$ is invertible in $E_h$. Write
\[
 h(s)=s^n+\sum_{r=0}^{n-1}h_rs^r,\qquad
 f_t(s)=\Phi_h(s)-ts,\qquad
 A_h(t)=A_h+tR,\quad R=e_0e_{n-1}^{\mathsf T}.
 \tag{PDE1}
\]
Here $A_h$ is the original arithmetic multiplication by $s$ on $E_h$;
$A_h(t)$ is its displayed companion deformation on the same coefficient
frame. The parameter $t$ does not replace the original marked arithmetic
quotient by a different quotient. The matrix $R$ includes the case
$n=1$, where it is the identity.

Fix $u\ne0$ on a specified branch of $\log u$. Set $d=n+1$ and retain
the actual rays and their orientations
\[
 \theta_j=\frac{\arg u+\pi+2\pi j}{d},\quad
 \ell_j(r)=re^{i\theta_j}\ (r\geq0),\quad
 \Gamma_j=-\ell_0+\ell_j\quad(1\leq j\leq n).
 \tag{PDE2}
\]
For an admitted amplitude $a$, let $\mathcal I_t(a)$ be the column with
entries $\int_{\Gamma_j}e^{f_t/u}a(s)\,ds$. Every amplitude used below is
a polynomial times a derivative of the fixed $v_h$, or a polynomial.
IPM1--6 proves their convergence on these rays with the original theta
bound; the explicit derivative estimate below also proves the required
first-derivative convergence directly. Compact sets of parameters have
a common dominating function with negative leading term
$-r^d/(d|u|)$ and lower polynomial terms. Thus these integrals and their
displayed parameter derivatives are entire in $t$. The fixed rays have
no $t$-dependent endpoint or orientation.

\subsection{The exact determinant on the whole polynomial family}

For this calculation only, allow all lower phase coefficients to vary:
\[
 f_{\mathbf c}(s)=\frac{s^{n+1}}{n+1}+\sum_{j=0}^n c_js^j,
 \quad h_{\mathbf c}=f_{\mathbf c}',\quad
 L_{\mathbf c}=u\partial_s+h_{\mathbf c},\quad
 \Pi(\mathbf c,u)_{j,b}
       =\int_{\Gamma_j}e^{f_{\mathbf c}/u}s^b\,ds.
 \tag{PDE3}
\]
The row index is $1\leq j\leq n$ and the column index $0\leq b<n$.
The original phase $f_t$ is the specified coefficient slice of this
family. No coefficient of that slice will be suppressed.

Ordinary monic division gives, for $0\leq j\leq n$ and $0\leq b<n$,
\[
 s^{b+j}=h_{\mathbf c}q_{b,j}+r_{b,j},\qquad
 \deg r_{b,j}<n,\quad\deg q_{b,j}\leq b+j-n.
 \tag{PDE4}
\]
A negative bound means the zero polynomial. In particular
$\deg q_{b,j}'\leq b-1<n$, so the exact differential reduction is
\[
 s^{b+j}=L_{\mathbf c}q_{b,j}
                    +(r_{b,j}-u q_{b,j}').
 \tag{PDE5}
\]
The leading term of $L_{\mathbf c}P$ has degree $n+\deg P$ and the same
leading coefficient as $P$. Repeated subtraction therefore proves that
the cokernel of $L_{\mathbf c}$ has the unique increasing frame of
degrees less than $n$, for every $\mathbf c$ and $u$. Formula PDE5 is
its literal remainder, not a specialization at $u=0$.

Let $B_j(\mathbf c,u)$ have these remainder columns, and let
$M_{\mathbf c}$ be multiplication by $s$ on
$\mathbb{C}[s]/(h_{\mathbf c})$. Ordinary remainders form
$M_{\mathbf c}^j$. The correction column $-u q_{b,j}'$ has row index
strictly less than $b$. It is consequently strictly triangular and has
zero trace. Thus, retaining its full nonzero matrix,
\[
 B_j=M_{\mathbf c}^{j}-u(\operatorname{coeff}q_{b,j}')_{b=0}^{n-1},
 \qquad \operatorname{Tr}B_j=\operatorname{Tr}M_{\mathbf c}^{j}.
 \tag{PDE6}
\]
The boundary integral of $e^{f_{\mathbf c}/u}q_{b,j}$ vanishes at
the two outer ends and cancels at the common origin. Differentiating
the original integrals and applying PDE5 gives
\[
 \partial_{c_j}\Pi=\frac1u\Pi B_j\quad(0\leq j\leq n).
 \tag{PDE7}
\]

Define the polynomial in the original coefficients
\[
 \mathcal S(\mathbf c)
  =\operatorname{Tr} f_{\mathbf c}(M_{\mathbf c})
  =\frac{\operatorname{Tr}M_{\mathbf c}^{n+1}}{n+1}
                +\sum_{j=0}^n c_j\operatorname{Tr}M_{\mathbf c}^j.
 \tag{PDE8}
\]
Its derivative with respect to $c_j$ is
$\operatorname{Tr}M_{\mathbf c}^j$. Here is a proof including repeated
roots. On the open set where $h_{\mathbf c}$ has distinct roots
$r_1,\ldots,r_n$, they admit local holomorphic enumerations and
$\mathcal S=\sum_i f_{\mathbf c}(r_i)$. Differentiation gives
$\partial_{c_j}\mathcal S=\sum_i r_i^j$, since
$f_{\mathbf c}'(r_i)=0$ makes every root-derivative term exactly zero.
Both sides of the identity are polynomials in $\mathbf c$; the open
set is dense because the derivative $s^n-1$ occurs in this family
and has distinct roots. Equality extends as a polynomial identity to
every coefficient value. Equivalently, the trace at a repeated root
is its full multiplicity times the scalar value, since powers of its
nilpotent part have zero diagonal. No multiplicity is removed.

Put $D=\det\Pi$. Multilinearity of the determinant in its columns
gives $\partial_{c_j}D=u^{-1}(\operatorname{Tr}B_j)D$ without any
prior assumption that $\Pi$ is invertible. PDE6 and PDE8 imply that
every derivative of $D\exp(-\mathcal S/u)$ is zero on the connected
coefficient space. Its value is determined at $\mathbf c=0$.

For completeness, write $\zeta=\exp(2\pi i/d)$ and
$\beta_b=(b+1)/d$. Direct integration on the original two rays gives
\[
 \Pi(0,u)_{j,b}
  =u^{(b+1)/d}d^{\beta_b-1}
          e^{i\pi\beta_b}\Gamma(\beta_b)
                       (\zeta^{j(b+1)}-1).
 \tag{PDE9}
\]
The power of $u$ uses the branch fixed in PDE2. To compute the remaining
determinant let $V_{r,b}=\zeta^{rb}$, $0\leq r,b<d$. Subtract row zero
from each other row and expand at its first column. The resulting
minor is exactly $(\zeta^{j(b+1)}-1)_{1\leq j\leq n,0\leq b<n}$,
with no row or column reversal. The Vandermonde formula gives
\[
 \det V=\prod_{0\leq p<q<d}(\zeta^q-\zeta^p),\qquad
 \zeta^q-\zeta^p
       =2i e^{\pi i(p+q)/d}\sin\frac{\pi(q-p)}d.
 \tag{PDE10}
\]
Every displayed sine is positive. The phase of this product is
$\pi d(d-1)/4+\pi(d-1)^2/2
=\pi(d-1)(3d-2)/4$.
Its first term comes from the $d(d-1)/2$ factors $i$, and its second
from $\sum_{p<q}(p+q)=d(d-1)^2/2$.
Also $V^*V=dI$ by the finite geometric sum, so $|\det V|=d^{d/2}$.
The retained gamma product is
$\prod_{b=0}^{n-1}\Gamma((b+1)/d)=(2\pi)^{n/2}/\sqrt d$,
the multiplication identity used in IPM21, also recorded in DLMF
5.5.7 at \url{https://dlmf.nist.gov/5.5.E7}.
The powers of $d$ from PDE9 sum to $-n/2$ and cancel this product's
$d^{-1/2}$ against $d^{d/2}$. Its additional phase is
$\pi\sum_b\beta_b=\pi n/2$. The total phase is therefore
$3\pi n(n+1)/4$. Consequently the exact determinant is
\[
 \boxed{\displaystyle
 \det\Pi(\mathbf c,u)
   =(2\pi)^{n/2}u^{n/2}
       \exp\left(\frac{3\pi i n(n+1)}4
                         +\frac{\mathcal S(\mathbf c)}u\right).}
 \tag{PDE11}
\]
In particular it never vanishes for $u\ne0$, including coefficients
with repeated critical points or repeated critical values. This proves
invertibility rather than assuming it to use a determinant equation.
For $n=1$ the formula is the downward Gaussian for $u>0$ on
$\arg u=0$, continued on the specified branch:
$-i\sqrt{2\pi u}\exp((c_0-c_1^2/2)/u)$.

\subsection{The original full quartet and the exact complex-parameter transport}

For the original full quartet $n=4m$, retain
$\rho=1/2+\delta+i\gamma$, $0<\delta<1/2$, $\gamma>2$, and
$h=\prod_{r\in\{\rho,\bar\rho,1-\rho,1-\bar\rho\}}(s-r)^m$.
Its reflection is $h(1-s)=h(s)$. With the original primitive constant
$A=\Phi_h(1/2)$, differentiation and evaluation at $s=1/2$ prove
\[
 \Phi_h(1-s)=2A-\Phi_h(s),\qquad
 f_t(1-s)=2A-t-f_t(s),\qquad
 \mathcal S(f_t)=n(A-t/2).
 \tag{PDE12}
\]
For the trace assertion, reflection pairs the roots of $h(s)-t$ with
their full multiplicities. Each pair has the displayed sum of phase
values. A root fixed by reflection has $s=1/2$ and phase value
$A-t/2$ itself. Summing over all $n$ roots, or taking the trace of the
corresponding algebra involution, proves the identity even at every
multiple-root parameter. Hence on this original quartet
\[
 \det\Pi_h(t,u)
  =(2\pi)^{n/2}u^{n/2}
       e^{3\pi i n(n+1)/4}\,e^{n(A-t/2)/u}.
 \tag{PDE13}
\]
Neither the original $A$ nor the $t$ term has been discarded.


% END RX INLINE 53



\clearpage
\part*{Complete original cyclic, period and arithmetic mixed control}
These complete provider editions are preserved byte-for-byte and
inserted in full. Their earlier finite estimates remain valid.
The later AMT proof sharpens OMV11's source-difference interval
and completes OMV10's literal arithmetic metric transfer; the
current mathematical receiving statements are MRI4c and MRI5d.
The exact invariant lift in OQC retains its whole action defect,
and OQX gives its sharper original rank after the period test
is evaluated in PQO and RPC. General multiplicities retain
the complete primary matrix and sourcewise finite metric maps.

% BEGIN RX INLINE 54; SHA256 db666f23d2b15878cfe4ea3618236cfea436031cefd0d4a64a260514136c126a
% New original-observation calculation; the sealed cut22 is not edited.
\section{The literal cyclic sectors and the original mixed kernel}
\label{sec:OCS}

This calculation evaluates the cyclic projector already present in
the original period observation. Its coefficients retain the original
period integrals, amplitude unit, and all primary nilpotents. The
result includes an exact finite matrix for the actual rank, as well
as unconditional rank bounds obtained from the original symmetric
tensor geometry.

\subsection{Original data and the Fourier basis of the literal matrix}
Use the following data of OPG1--5, including the same contours and
their orientations:
\[
 \begin{gathered}
 \alpha_{\epsilon,\eta}=\tfrac12+\epsilon\delta+i\eta\gamma,
 \quad\epsilon,\eta\in\{+1,-1\},\quad0<\delta<\tfrac12,\quad\gamma>2,\\
 h(z_1)=\prod_{\epsilon,\eta}(z_1-\alpha_{\epsilon,\eta})^m,
 \quad n=4m,\quad N=n+1,\quad E_h=\mathbb C[z_1]/(h),\\
 k=4l+1,\quad l,m\geq1,\quad c=k/2,
 \quad e=1+k(m-1),\quad q=e(k+1)^2,\\
 \beta_{a,b}=c+(2a-k)\delta+i(2b-k)\gamma,
 \quad\chi(S)=\prod_{a,b=0}^k(S-\beta_{a,b})^e,
 \quad E=\mathbb C[S]/(\chi).
 \end{gathered}\tag{OCS1}
\]
Here $A_h$ and $A$ are multiplication by $z_1$ and $S$ on their
respective algebras, and $Y=(A-cI)/i$. The original amplitude
unit is $U=M_{j_hv_h}$, with $v_h=(2\xi)/h$ and every original
$v_h(\alpha)\ne0$. For the specified nonzero $u$ and its branch,
$\Phi_h'=h$, $\Phi_h(0)=0$, and
\[
 \Gamma_j=-\ell_0+\ell_j,\qquad
 \arg\ell_j=\frac{\arg u+\pi+2\pi j}{N},\qquad
 (\Pi[P])_j=\int_{\Gamma_j}e^{\Phi_h(z_1)/u}
                       \operatorname{rem}_hP(z_1)\,dz_1.
 \tag{OCS2}
\]
The representative has degree less than $n$. The convergence and
invertibility of this original period map are the complete period
theorems retained in OPG4. No period value or unit derivative is
chosen in the present calculation. The original injection and
observation are
\[
 \iota[P]=[P(z_1+\cdots+z_k)],\quad
 \mathcal V=\Pi^{\otimes k}U^{\otimes k}\iota,
 \quad P_k=\frac1N\sum_{a=0}^{N-1}(C^{-a})^{\otimes k},
 \quad L=P_k\mathcal V=\jmath\Lambda,
 \quad B=\operatorname{im}L,\quad K=\ker\Lambda.
 \tag{OCS3}
\]
The map $\jmath:B\hookrightarrow(\mathbb C^n)^{\otimes k}$
is the original image inclusion. The injection $\iota$ has precisely
the defining polynomial in OCS1: on each ordered primary tensor
the sum of its $k$ independent nilpotents has index
$e=1+k(m-1)$, its top coefficient being
$(k(m-1))!/((m-1)!)^k\ne0$, and its distinct eigenvalues are
the displayed $\beta_{a,b}$. Thus $\mathcal V$ is injective.

On a period-coordinate column $x\in\mathbb C^{N-1}$ the literal
matrix in OPG4--5 is
\[
 (Cx)_j=x_{j+1}-x_1\quad(1\leq j<N-1),\qquad
 (Cx)_{N-1}=-x_1.
 \tag{OCS4}
\]
Set $\zeta=\exp(2\pi i/N)$ and, for $1\leq t\leq N-1$, let
\[
 (f_t)_j=\zeta^{jt}-1,\qquad
 F=[f_1\ \cdots\ f_{N-1}],\qquad
 (F^{-1}x)_t=\frac1N\sum_{j=1}^{N-1}\zeta^{-tj}x_j.
 \tag{OCS5}
\]
The last expression is the exact inverse. Indeed
$\sum_{j=1}^{N-1}\zeta^{-tj}(\zeta^{uj}-1)=N\mathbf1_{t=u}$
for $t,u\ne0$ modulo $N$, by the finite geometric sum.
The first $N-2$ rows of OCS4 give
$\zeta^{(j+1)t}-\zeta^t=\zeta^t(\zeta^{jt}-1)$, and the
last row gives $1-\zeta^t=\zeta^t(\zeta^{(N-1)t}-1)$.
Consequently $Cf_t=\zeta^tf_t$, proving the orientation of
every inverse power in OCS3, as well as $C^N=I$.
The original Euclidean Gram of this basis is
\[
 F^*F=N(I_{N-1}+\mathbf1\mathbf1^*).
 \tag{OCS6}
\]
To verify it, expand
$\sum_j(\zeta^{-jt}-1)(\zeta^{ju}-1)$ and use the same
geometric sums; its value is $N(\mathbf1_{t=u}+1)$.
Thus the complete coefficient and metric transformations are
both specified.

\subsection{The ordered-to-symmetric map and the exact sector projector}
Let
\[
 \mathcal A_k=\{\nu\in\mathbb Z_{\geq0}^{N-1}:\sum_t\nu_t=k\},
 \quad\operatorname{wt}(\nu)=\sum_{t=1}^{N-1}t\nu_t\pmod N,
 \quad\mathcal I_{k,\tau}=\{\nu\in\mathcal A_k:
                                      \operatorname{wt}(\nu)=\tau\}.
 \tag{OCS7}
\]
For each $\nu$ define the following unscaled symmetric tensor:
\[
 S_\nu=\sum_{\substack{(t_1,\ldots,t_k)\in\{1,\ldots,N-1\}^k\\
                   \#\{j:t_j=t\}=\nu_t\ \text{for each }t}}
                    f_{t_1}\otimes\cdots\otimes f_{t_k}.
 \tag{OCS8}
\]
Every distinct ordered word occurs once. Since the ordered tensors
in the $f_t$ form a basis, their disjoint word supports show that
the $S_\nu$ are linearly independent. A permutation-invariant tensor
has equal coefficients on each word orbit, so they span exactly
$\operatorname{Sym}^k(\mathbb C^{N-1})$, realized as that invariant
subspace of the full ordered tensor product. In particular, for
$x=(x_1,\ldots,x_{N-1})$,
\[
 (Fx)^{\otimes k}=\sum_{\nu\in\mathcal A_k}x^\nu S_\nu,
 \quad
 (C^{-a})^{\otimes k}S_\nu=
                         \zeta^{-a\operatorname{wt}(\nu)}S_\nu,
 \quad
 P_k(Fx)^{\otimes k}=\sum_{\nu\in\mathcal I_{k,0}}x^\nu S_\nu.
 \tag{OCS9}
\]
The finite average in OCS3 is exactly one on weight zero and
zero on every other weight, by its geometric sum. No multinomial
coefficient is omitted: the multiplicity $k!/\prod_t\nu_t!$
is already the number of summands in OCS8.

Write $\mathcal F_k:\mathbb C^{\mathcal A_k}\to
\operatorname{Sym}^k(\mathbb C^{N-1})$ for the basis isomorphism
$e_\nu\mapsto S_\nu$. Let $J_0$ be the coordinate inclusion of
$\mathbb C^{\mathcal I_{k,0}}$ into $\mathbb C^{\mathcal A_k}$,
$\pi_0$ the coordinate restriction in the opposite direction,
and let $\mathcal F_0=\mathcal F_kJ_0$, with its target the original
ordered tensor space. The exact projector factorization is
\[
 P_k\mathcal F_k=\mathcal F_0\pi_0.
 \tag{OCS10}
\]
All images of $\mathcal V$ are symmetric: polynomial evaluation
at a sum is invariant under every factor permutation, and the
same map $\Pi U$ is applied in all factors. Therefore OCS10
applies on the entire original image $\mathcal V(E)$.

For completeness, the full inherited metric is retained in these
coordinates. Put $G=F^*F$ as in OCS6. The Gram of the basis OCS8 is
\[
 \mathcal G_{\nu,\mu}=
 \sum_{\substack{a_{tu}\geq0\\\sum_u a_{tu}=\nu_t\ ,\
                                    \sum_t a_{tu}=\mu_u}}
       \frac{k!}{\prod_{t,u}a_{tu}!}\prod_{t,u}G_{tu}^{a_{tu}}.
 \tag{OCS11}
\]
Indeed the coefficient of $\bar x^\nu y^\mu$ in
$(\bar x^TGy)^k$ is the right side by the multinomial theorem.
Taking the inner product of the two pure tensors in OCS9 gives
the same coefficient as $\langle S_\nu,S_\mu\rangle$.
This proves OCS11, including its off-diagonal entries; it is
positive definite because OCS8 is a basis. On the full symmetric
coordinates the projector matrix is $J_0\pi_0$, and its adjoint
for this exact Gram is
$\mathcal G^{-1}(J_0\pi_0)^*\mathcal G$.
Thus the Fourier-sector computation makes no unproved
orthogonality or norm assertion about the original cyclic average.

\subsection{Every primary period coefficient of the projected curve}
For each original root $\alpha$ let $E_{\alpha,a}$, $0\leq a<m$,
be the degree-less-than-$4m$ primary representatives of OPG10.
Explicitly, if $h_\alpha=h/(z_1-\alpha)^m$, put
\[
 E_{\alpha,a}(z_1)=\operatorname{rem}_h\left[
 h_\alpha(z_1)\sum_{j=0}^{m-1}
       \frac{(h_\alpha^{-1})^{(j)}(\alpha)}{j!}
                  (z_1-\alpha)^{j+a}\right],\qquad
 b_{j;\alpha,a}=\int_{\Gamma_j}e^{\Phi_h(z_1)/u}
                                    E_{\alpha,a}(z_1)\,dz_1.
 \tag{OCS12}
\]
The four primary factors are coprime. The Taylor inverse gives
constant one on its own factor and zero on the other three;
thus the classes $[E_{\alpha,a}]$ form the full primary basis
with all $m$ nilpotent degrees at every root.
Define the exact finite Fourier and unit coefficients
\[
 \widehat b_{t;\alpha,a}=
     \frac1N\sum_{j=1}^{N-1}\zeta^{-tj}b_{j;\alpha,a},
 \qquad
 f_{t;\alpha,d}=\frac{i^{-d}}{d!}
       \sum_{a=d}^{m-1}\widehat b_{t;\alpha,a}
                         \frac{v_h^{(a-d)}(\alpha)}{(a-d)!},
 \quad0\leq d<m.
 \tag{OCS13}
\]
Let $\omega_\alpha=(\alpha-1/2)/i$ and set
\[
 \mathbf v(w)=\Pi U\exp\bigl(w(A_h-\tfrac12I)/i\bigr)[1],
 \quad x(w)=F^{-1}\mathbf v(w),\qquad
 x_t(w)=\sum_{\alpha}e^{\omega_\alpha w}
                          \sum_{d=0}^{m-1}f_{t;\alpha,d}w^d.
 \tag{OCS14}
\]
To verify the final equality, in the primary algebra at $\alpha$
multiply the Taylor series of $v_h$ by the terminating series of
$\exp(w(z_1-\alpha)/i)$. The coefficient on
$[E_{\alpha,a}]$ is
$e^{\omega_\alpha w}\sum_{d=0}^a
v_h^{(a-d)}(\alpha)(w/i)^d/((a-d)!d!)$.
Apply the original polynomial period map and the inverse in OCS5;
this gives precisely OCS13--14. In particular, the unreduced
entire function has passed through its full primary remainder
before any period is integrated.

For $\nu\in\mathcal I_{k,0}$, $0\leq a,b\leq k$ and
$0\leq D<e$, define $c_{\nu;a,b,D}$ by the following finite
multinomial rule. Sum over integers $n_{t,\alpha,d}\geq0$ satisfying
\[
 \begin{gathered}
 \sum_{\alpha,d}n_{t,\alpha,d}=\nu_t\quad\text{for each }t,\qquad
 \sum_{t,\alpha,d}d\,n_{t,\alpha,d}=D,\\
 \sum_{t,d,\alpha:\operatorname{Re}\alpha=1/2+\delta}
                                 n_{t,\alpha,d}=a,\qquad
 \sum_{t,d,\alpha:\operatorname{Im}\alpha=+\gamma}
                                 n_{t,\alpha,d}=b.
 \end{gathered}
\]
For these indices put
\[
 c_{\nu;a,b,D}=
 \sum_{\text{stated indices}}
   \left(\prod_{t=1}^{N-1}
             \frac{\nu_t!}{\prod_{\alpha,d}n_{t,\alpha,d}!}\right)
        \prod_{t,\alpha,d}f_{t;\alpha,d}^{\,n_{t,\alpha,d}}.
 \tag{OCS15}
\]
Empty index sets give zero and zero exponents give factors one.
Every original primary coefficient is present in this finite rule.
The ordinary multinomial expansion of OCS14 now proves
\[
 x(w)^\nu=\sum_{a,b=0}^k e^{\omega_{a,b}w}
                   \sum_{D=0}^{e-1}c_{\nu;a,b,D}w^D,
 \qquad
 \omega_{a,b}=(\beta_{a,b}-c)/i
                 =(2b-k)\gamma-i(2a-k)\delta.
 \tag{OCS16}
\]
Indeed $D\leq k(m-1)=e-1$ and the counts of each original
sign give exactly the displayed exponent. The sums with different
$(a,b)$ are distinct since their real and imaginary parts are
independently distinct.

\subsection{The actual finite observation matrix and its complete kernel}
Define the primary coefficient isomorphism
$\mathcal H:\mathbb C^q\to E$ by sending coordinate $(a,b,D)$
to $\varepsilon_{a,b}(S-\beta_{a,b})^D$, where
$\varepsilon_{a,b}$ is obtained from $\chi$ by the same truncated
Taylor-inverse construction as OCS12, now of order $e$.
The proof by the coprime factors is identical and shows that this
is a basis of all $q$ primary directions. Set
\[
 \mathfrak A_{\nu;(a,b,D)}=i^D D!\,c_{\nu;a,b,D},
 \qquad
 \mathfrak A:\mathbb C^q\longrightarrow
                            \mathbb C^{\mathcal I_{k,0}}.
 \tag{OCS17}
\]
Then, on the original spaces and with the original ordered target,
\[
 \boxed{\ L\mathcal H=\mathcal F_0\mathfrak A,\qquad
 K=\mathcal H(\ker\mathfrak A),\qquad
 \operatorname{rank}\Lambda=\operatorname{rank}\mathfrak A. }
 \tag{OCS18}
\]
Here is a proof giving every factorial and phase. In the primary
sum algebra,
\[
 e^{wY}[1]=\sum_{a,b}e^{\omega_{a,b}w}
          \sum_{D=0}^{e-1}\frac{i^{-D}w^D}{D!}
                  \varepsilon_{a,b}(S-\beta_{a,b})^D.
\]
The exact original tensor identity of OPG8 gives
$\mathcal V e^{wY}[1]=\mathbf v(w)^{\otimes k}$: the centre
$c=k/2$ splits into the $k$ original centres $1/2$, while
$U$ commutes with multiplication by $z_1$.
Applying the literal projector and then OCS9 and OCS16 gives
the right-hand curve with coefficients
$\mathcal F_0c_{a,b,D}$.
Exponential polynomials $w^De^{\omega_{a,b}w}$ with these distinct
exponents are linearly independent. To prove this directly,
apply $\prod_{(a',b')\ne(a,b)}(\partial_w-\omega_{a',b'})^e$
to any zero linear combination. It kills all the other exponents.
On the selected polynomial factor, each remaining operator is
$\partial_w+(\omega_{a,b}-\omega_{a',b'})$, invertible on the
polynomials of degree less than $e$ because its constant term
is nonzero and differentiation is nilpotent there. The surviving
polynomial must therefore vanish. Doing this for every pair proves
the asserted independence. Equating coefficients now gives
$L\mathcal H e_{a,b,D}=\mathcal F_0(i^DD!c_{a,b,D})$.
This proves the first equality in OCS18 on a basis. The other
two follow from injectivity of $\mathcal F_0$ and $\jmath$.

In particular the isomorphism from the actual sector image to
the original boundary is
\[
 \mathcal Q:\operatorname{im}\mathfrak A\xrightarrow{\ \sim\ }B,
 \qquad \mathcal Qx=\jmath^{-1}\mathcal F_0x,
 \qquad
 E/K\xrightarrow{\ \sim\ }\operatorname{im}\mathfrak A,
 \quad[z]\longmapsto\mathfrak A\mathcal H^{-1}z.
 \tag{OCS19}
\]
Both inverses exist on precisely their displayed images by OCS18.
These are full coefficient maps; the inherited Gram of
$\mathcal F_0$ is $J_0^*\mathcal GJ_0$ from OCS11.

For a completely explicit finite rank prescription, define
\[
 \Delta_r(\mathfrak A)=
   \sum_{\substack{I\subseteq\mathcal I_{k,0},\ J\subseteq\{(a,b,D)\}\\
                   |I|=|J|=r}}|\det\mathfrak A_{I,J}|^2,
 \qquad\Delta_0=1.
 \tag{OCS20}
\]
Every entry of every minor is the finite expression OCS12--17
in the actual periods and amplitude derivatives. Gaussian
elimination proves that a matrix has rank at least $r$ exactly
when one of its $r$ by $r$ minors is nonzero: pivot rows and
columns supply such a minor, and a nonzero minor supplies $r$
independent columns. Hence the actual rank is exactly
$\max\{r:\Delta_r(\mathfrak A)>0\}$. This evaluates its
finite deciding data without prescribing generic values or
asserting that any unexamined maximal minor is nonzero.

The Gamma vectors used in the existing endpoint increments have
the equally explicit coordinates
\[
 \jmath\lambda_d=\mathcal F_0\widehat\lambda_d,
 \qquad
 (\widehat\lambda_d)_\nu=
 \frac1{\sqrt{h_d^{(s)}}}\sum_{a,b=0}^k\sum_{D=0}^{e-1}
               c_{\nu;a,b,D}p_d^{(D)}(\omega_{a,b}),
 \quad h_d^{(s)}=(2\pi)^{s/2}d!(s/2)_d,
 \quad s\in\{1,k\}.
 \tag{OCS21}
\]
Taylor expansion of $p_d(Y)[1]$ in the sum primary basis gives
the coefficient $i^{-D}p_d^{(D)}(\omega_{a,b})/D!$.
Multiplication by OCS17 cancels exactly these factorials and
phases, proving OCS21. This also proves that the generating
formula of OPG9 becomes
$(1+z^2)^{-s/4}\mathcal F_0
 (x(\arctan z)^\nu)_{\nu\in\mathcal I_{k,0}}$,
with the original scalar exponent and both existing sources.

\subsection{Exact invariant dimensions and unconditional forced losses}
Put $d_{N,k,\tau}=|\mathcal I_{k,\tau}|$.
By OCS7--10 this is the dimension of the full symmetric sector,
and its exact finite character formula is
\[
 d_{N,k,0}=[T^k]\frac1N
       \sum_{h\mid N}\varphi(h)\frac{1-T}{(1-T^h)^{N/h}}.
 \tag{OCS22}
\]
To prove it, the trace generating series for $C^a$ on symmetric
tensors is
$\prod_{t=1}^{N-1}(1-T\zeta^{at})^{-1}$, by summing the
monomial basis OCS8. If $\zeta^a$ has order $h$, the list
$\zeta^{at}$ with $0\leq t<N$ contains every $h$th root
$N/h$ times. Their product is
$\prod_{t=0}^{N-1}(1-T\zeta^{at})=(1-T^h)^{N/h}$.
Removing the $t=0$ factor gives the series
$(1-T)/(1-T^h)^{N/h}$. There are exactly $\varphi(h)$
indices $a$ of that order, so averaging the traces of the
original finite projector proves OCS22. The inverse powers
give the same average by $a\mapsto-a$.
In ordinary finite binomial terms, for $k\geq1$ this is
\[
 d_{N,k,0}=\frac1N\sum_{h\mid N}\varphi(h)
 \left\{
 \mathbf1_{h\mid k}\binom{k/h+N/h-1}{N/h-1}
 -\mathbf1_{h\mid(k-1)}
                  \binom{(k-1)/h+N/h-1}{N/h-1}
 \right\}.
 \tag{OCS23}
\]
The binomial expansion of $(1-T^h)^{-N/h}$ follows by
counting weak compositions, proving the coefficient extraction
used here. It follows from the actual matrix OCS17 that
\[
 \operatorname{rank}\Lambda\leq\min(q,d_{N,k,0}),
 \qquad \dim K\geq\max(0,q-d_{N,k,0}).
 \tag{OCS24}
\]
These inequalities apply to the original map, since its full
image factorization was proved in OCS18; the ambient dimension
has not been equated with its actual rank.
There is also an unconditional nonzero observation for the
original $k\geq5$. Each entire function $x_t(w)$ in OCS14 is
nonzero. Otherwise its nonzero coordinate covector would
annihilate every derivative at zero of
$\Pi U\exp(w(A_h-1/2)/i)[1]$. Those derivatives span the full
space, since $[1]$ is cyclic for multiplication by $z_1$ and
$\Pi U$ is invertible. A product of finitely many nonzero entire
functions is nonzero. The zero sector is nonempty: take the
weights $1,1,N-2$ and then $(k-3)/2$ pairs $1,N-1$.
Their sum is a multiple of $N$. Its monomial curve in OCS9
is nonzero, so $\operatorname{rank}\Lambda\geq1$.

For $m=1$, $N=5$, and every $k\geq0$, write
\[
 D_k=\binom{k+3}{3},\qquad
 \epsilon_k=\begin{cases}1&k\equiv0\pmod5,\\-1&k\equiv1\pmod5,\\
                         0&\text{otherwise}.
             \end{cases}
 \qquad
 d_{5,k,0}=\frac{D_k+4\epsilon_k}{5},\quad
 d_{5,k,\tau}=\frac{D_k-\epsilon_k}{5}\quad(\tau\ne0).
 \tag{OCS25}
\]
For each of the four nonidentity powers of $C$, the trace
series is $(1-T)/(1-T^5)$, whose coefficient is $\epsilon_k$.
The weight-$\tau$ projector introduces the Fourier multiplier;
the sum of its four nontrivial fifth roots is $4$ when
$\tau=0$ and $-1$ otherwise. This proves both formulas.
For the requested original orders the exact counts and resulting
kernel lower bounds are
\[
\begin{array}{c|r|r|r|r}
 k&D_k&d_{5,k,0}&q=(k+1)^2&\max(0,q-d_{5,k,0})\\\hline
 5&56&12&36&24\\
 9&220&44&100&56\\
 13&560&112&196&84\\
 17&1140&228&324&96\\
 21&2024&404&484&80\\
 25&3276&656&676&20\\
 29&4960&992&900&0
\end{array}\tag{OCS26}
\]
These are integer evaluations of OCS25. The strict inequality
$d_{5,k,0}<q$ holds for $k=4l+1$, $l\geq1$, exactly for the
first six orders in the table. They exhaust $k\leq25$ in this
progression. For $k\geq29$,
\[
 D_k-5(k+1)^2=\frac{(k+1)(k^2-25k-24)}6\geq460>4,
\]
where $k^2-25k-24$ is positive and increasing from $k=29$.
Even the most negative correction $4\epsilon_k\geq-4$
therefore gives $d_{5,k,0}>q$. This proves the threshold
without extrapolating the finite table.

\subsection{The actual simple-quartet quadric and a stronger rank bound}
Continue with $m=1$. Order the four roots as
$(++),(+-),(-+),(--)$ and let
$\mathcal E_{\rm prim}:\mathbb C^4\to E_h$ be the primary
basis isomorphism whose columns are the four $[E_{\alpha,0}]$.
Define the actual invertible matrix and the original exponential
quartet by
\[
 \mathsf H=F^{-1}\Pi U\mathcal E_{\rm prim},\qquad
 t_{\epsilon,\eta}(w)=e^{(\eta\gamma-i\epsilon\delta)w},
 \qquad x(w)=\mathsf Ht(w).
 \tag{OCS27}
\]
Every column of $\mathsf H$ is OCS13 with $d=0$, namely the
Fourier transform of its original primary period times its
actual nonzero $v_h(\alpha)$. Each displayed map is invertible,
so $\mathsf H$ is invertible without a genericity assumption.
Put $\mathcal R_k=\mathbb C[x_1,x_2,x_3,x_4]_k$, the homogeneous
polynomials of degree $k$, and define
\[
 Q_0(t)=t_{++}t_{--}-t_{+-}t_{-+},\qquad
 Q(x)=Q_0(\mathsf H^{-1}x).
 \tag{OCS28}
\]
This is the literal transformed quadratic relation of the
original curve; its coefficients are determined by the periods
and unit through OCS27. In particular $Q\ne0$.

The kernel of homogeneous restriction to that curve is exactly
\[
 \ker\{\mathcal R_k\to\operatorname{Hol}(\mathbb C),
                         f\mapsto f(x(w))\}
            =Q\mathcal R_{k-2}\quad(k\geq2).
 \tag{OCS29}
\]
Here is a complete verification. First work in the $t$ coordinates.
A degree-$k$ monomial has exponent
$(2b-k)\gamma-i(2a-k)\delta$, where $a$ is its count of
first sign $+$ and $b$ its count of second sign $+$.
These $(k+1)^2$ complex exponents are distinct: equality of
their real parts forces $b$ equal and equality of their
imaginary parts forces $a$ equal. Equivalently, the two
frequencies $2\gamma$ and $-2i\delta$ admit no nonzero
integer relation, by their real and imaginary parts. Thus
the exponentials have no extra polynomial relation of this
homogeneous type; their linear independence follows, for
example, by differentiating $0,\ldots,(k+1)^2-1$ times and
using the Vandermonde determinant of the distinct exponents.
Every pair $(a,b)$ occurs: the count $r_{++}$ may be any integer
between $\max(0,a+b-k)$ and $\min(a,b)$, a nonempty interval,
and the other three counts are $a-r_{++}$, $b-r_{++}$,
and $k-a-b+r_{++}$.

Within any fixed pair $(a,b)$, two monomial exponent vectors
differ by an integer multiple of $(1,-1,-1,1)$. After their
common monomial is factored out, their difference is therefore
the difference of equal powers of $t_{++}t_{--}$ and
$t_{+-}t_{-+}$. The identity
$A^r-B^r=(A-B)\sum_{j=0}^{r-1}A^{r-1-j}B^j$ shows that
this difference is divisible by $Q_0$. Any polynomial
vanishing on the curve has coefficient sum zero within
each pair $(a,b)$, by the proved exponential independence.
It is thus a sum of these monomial differences and is
divisible by $Q_0$. Its quotient is homogeneous of degree
$k-2$. Conversely $Q_0(t(w))=0$, since each of its two
products is one. This proves the assertion in $t$
coordinates, and the invertible substitution OCS27 proves
OCS29. In particular
$\dim(\mathcal R_k/Q\mathcal R_{k-2})=(k+1)^2=q$;
no information about the original curve has been replaced
by a freely chosen quadric.

Let $\mathcal I_k\subseteq\mathcal R_k$ be the span of the
weight-zero monomials $x^\nu$, using the weights $1,2,3,4$
from the exact Fourier basis. The coordinate functions of
the original projected curve in OCS9 are precisely these
monomials restricted to $x(w)$. The span of the vector
curve $Le^{wY}[1]$ equals $\operatorname{im}L$: its
derivatives at zero contain $LY^d[1]$, $0\leq d<q$, and
$Y^d[1]$ is a basis because $Y$ is an invertible affine
change of multiplication by $S$ in the degree filtration.
The rank of that vector curve equals the dimension of
the span of its coordinate functions, by the equality of
row and column ranks of its finite exponential coefficient
matrix. Hence OCS29 proves the exact equality
\[
 \operatorname{rank}\Lambda
   =d_{5,k,0}-\dim(\mathcal I_k\cap Q\mathcal R_{k-2}).
 \tag{OCS30}
\]

Fix any multiplicative monomial order, for instance
lexicographic order, and write $x^\mu$ for the leading
monomial of the actual $Q$, with $|\mu|=2$. Choose a
row-echelon basis of the finite-dimensional intersection
in OCS30 in descending monomial order. Its leading
monomials are distinct, all have weight zero, and all
are divisible by $x^\mu$: the leading monomial of
$Qf$ is the product of those of $Q$ and $f$ because
the order is multiplicative and the coefficient field
has no zero divisors. Thus each such leading monomial
is $x^\mu x^\eta$ for a distinct degree-$k-2$ monomial
of weight $-\operatorname{wt}(\mu)$. Counting this set
and using OCS30 gives
\[
 \boxed{\quad
 d_{5,k,0}-d_{5,k-2,-\operatorname{wt}(\mu)}
 \ \leq\ \operatorname{rank}\Lambda
 \ \leq\ \min(q,d_{5,k,0}).\quad}
 \tag{OCS31}
\]
The left side depends only on the actual leading monomial
of the actual period quadric, and remains valid whatever
its other coefficients are. A uniform consequence retaining
the strongest count independent of that monomial is
\[
 r_k^-:=d_{5,k,0}-\max_{\tau\in\mathbb Z/5}d_{5,k-2,\tau}
       =\left\lfloor\frac{(k+1)^2}{5}\right\rfloor
                                  +\mathbf1_{5\mid k}
       \leq\operatorname{rank}\Lambda.
 \tag{OCS32}
\]
For the equality, $D_k-D_{k-2}=(k+1)^2$, as direct
subtraction of the two cubic binomials shows. Formula
OCS25 then gives
\[
 r_k^-=\frac{(k+1)^2+4\epsilon_k
                   -\max(4\epsilon_{k-2},-\epsilon_{k-2})}{5}.
\]
For $k$ congruent respectively to $0,1,2,3,4$ modulo five,
its numerator correction is respectively $4,-4,-4,-1,0$.
The corresponding residues of $(k+1)^2$ are $1,4,4,1,0$.
These five cases prove OCS32 for every $k\geq2$.
For the original orders in OCS26 the complete unconditional
intervals are therefore
\[
\begin{array}{c|r|r|r|r}
 k&r_k^-&\min(q,d_{5,k,0})&\max(0,q-d_{5,k,0})&q-r_k^-\\\hline
 5&8&12&24&28\\
 9&20&44&56&80\\
 13&39&112&84&157\\
 17&64&228&96&260\\
 21&96&404&80&388\\
 25&136&656&20&540\\
 29&180&900&0&720
\end{array}\tag{OCS33}
\]
The middle two rank columns bound $\operatorname{rank}\Lambda$;
the last two columns bound $\dim K$ from below and above,
respectively. No endpoint of either interval has been assigned
as the actual value without its original minors.

\subsection{A second exact finite matrix for the simple-quartet kernel}
The quadric calculation also gives the actual kernel a
polynomial quotient model. Write $Q(x)=\sum_{|\mu|=2}q_\mu x^\mu$.
Let $T_Q$ be the full multiplication matrix
$\mathcal R_{k-2}\to\mathcal R_k$ in the monomial bases:
\[
 (T_Q)_{\nu,\eta}=q_{\nu-\eta},
 \quad |\nu|=k,\quad|\eta|=k-2,
 \qquad q_{\nu-\eta}=0\text{ if an entry of }\nu-\eta<0.
\]
Let $T_{\rm out}$ be its rows of nonzero cyclic weight.
Every coefficient $q_\mu$ is explicitly fixed by
$Q_0(\mathsf H^{-1}x)$, and the inverse matrix may be
computed as $\operatorname{adj}\mathsf H/\det\mathsf H$,
with nonzero denominator proved in OCS27. Multiplication
by $Q$ is injective. The condition that $Qf$ belongs
to $\mathcal I_k$ is exactly $T_{\rm out}f=0$.
Thus OCS30 gives the further exact rank identities
\[
 \begin{split}
 \operatorname{rank}\Lambda
   &=d_{5,k,0}-D_{k-2}+\operatorname{rank}T_{\rm out},\\
 \dim K&=D_k-d_{5,k,0}-\operatorname{rank}T_{\rm out}.
 \end{split}\tag{OCS34}
\]
The ranks of this concrete matrix are decided by its finite
minors, by the same elimination proof as OCS20. This is
an alternative exact calculation from the actual periods,
not an assumption about the quadric's coefficients.

To give the typed map behind the kernel equality, let
$E^*$ denote the complex linear dual. Define
$\Theta:\mathcal R_k\to E^*$ as follows: $\Theta(f)$ is
the unique covector such that
\[
 \Theta(f)(e^{wY}[1])=f(x(w))\quad\text{for all }w.
 \tag{OCS35}
\]
For $m=1$ every function on the right is a linear
combination of the $q$ distinct exponentials in OCS16.
The primary expansion of $e^{wY}[1]$ gives exactly those
exponentials as a basis of covector functions. This proves
existence and uniqueness. The proof of OCS29 proves
surjectivity and $\ker\Theta=Q\mathcal R_{k-2}$.
Moreover $\Theta(\mathcal I_k)=\operatorname{im}L^\vee$,
where $L^\vee$ denotes the complex linear dual map:
the monomial coefficient functionals in the basis OCS8
pull back along $L$ to exactly OCS35 for $x^\nu$ of
weight zero. These functionals span the full dual of
the actual image, since every linear functional on a
subspace of a finite-dimensional space extends to the
whole space. Restriction of $E^*$ to $K$ has kernel
$\operatorname{im}L^\vee$: a covector vanishing on $K$
factors uniquely through $E/K\simeq\operatorname{im}L$.
Thus the exact maps give
\[
 \boxed{\quad
 \operatorname{coker}T_{\rm out}
 \simeq\mathcal R_k/(\mathcal I_k+Q\mathcal R_{k-2})
 \xrightarrow{\ \Theta|_K\ }K^*.
 \quad}\tag{OCS36}
\]
For the first isomorphism, use the direct monomial splitting
of $\mathcal R_k$ into its zero and nonzero sectors; after
quotienting by $\mathcal I_k$, multiplication by $Q$
becomes precisely $T_{\rm out}$. For the second, the
surjectivity and kernel were just proved. This identifies
the actual lost directions by a full quotient map.

\subsection{All primary action directions remain in the mixed calculation}
For arbitrary original $m\geq1$, retain the full primary
matrix $\mathfrak A$ in OCS17. In the basis $\mathcal H$
the original sum action is exactly
\[
 J e_{a,b,D}=\omega_{a,b}e_{a,b,D}
       +\begin{cases}i^{-1}e_{a,b,D+1}&D<e-1,\\0&D=e-1,
         \end{cases}
 \qquad Y\mathcal H=\mathcal HJ.
 \tag{OCS37}
\]
This follows by multiplying
$\varepsilon_{a,b}(S-\beta_{a,b})^D$ by
$(S-c)/i=\omega_{a,b}+(S-\beta_{a,b})/i$ in its original
primary algebra. Consequently the actual possible action
leaving the cyclic kernel is the fully specified map
\[
 \ker\mathfrak A\longrightarrow\operatorname{im}\mathfrak A,
 \qquad x\longmapsto\mathfrak A Jx,
 \qquad
 \Lambda Y\mathcal Hx=\mathcal Q\mathfrak A Jx.
 \tag{OCS38}
\]
The codomain assertion follows from $\mathfrak A Jx$ being
a value of the same matrix $\mathfrak A$ on $Jx$.
For the original inclusion $I:K\hookrightarrow E$, this
is the OPG22 defect transported through the exact kernel
isomorphism in OCS18. In particular none of the discarded
sectors has been declared invariant under $Y$.

The part of the kernel undetected by all original iterates
is exactly computed by the stacked matrix
\[
 \mathfrak O=\begin{bmatrix}
          \mathfrak A\\\mathfrak AJ\\\vdots\\\mathfrak AJ^{q-1}
          \end{bmatrix},\qquad
 \mathcal K_\infty=\mathcal H\ker\mathfrak O,
 \quad
 \dim\mathcal K_\infty=q-\operatorname{rank}\mathfrak O,
 \quad
 \dim(K/\mathcal K_\infty)=
        \operatorname{rank}\mathfrak O-\operatorname{rank}\mathfrak A.
 \tag{OCS39}
\]
Indeed $\mathfrak A J^a x$ is the exact sector-coordinate
value of $\Lambda Y^a\mathcal Hx$ by OCS18 and OCS37.
The polynomial $\chi(c+iT)$ annihilates $J$ and has
degree $q$ with leading coefficient $i^q$. It expresses
every higher power in the span of the first $q$, so the
stacked kernel is precisely the intersection over all
iterates and is $J$-invariant. Rank-nullity proves both
dimension formulas. Its quotient identifies the directions
that the original one-step observation loses and subsequent
original iterates detect.

Equations OCS17--21, OCS34--39 supply actual finite maps
and deciding minors for the original observation and its
kernel. The forced kernel losses in OCS26 and the rank
intervals OCS33 concern this unchanged projector and these
unchanged periods. The endpoint volume calculation still
uses the full positive sums of the actual vectors
$\lambda_d$ in OPG16--18 and their exact inherited
coordinates OCS11 and OCS19. No volume growth estimate
is deduced from these dimension bounds alone.

% END RX INLINE 54


% BEGIN RX INLINE 55; SHA256 cda01263392b8297f2516026434bbaf56ba063ed782200493615eee886d3b391
% Derived from the literal original period quadric and cyclic projector.
\section{The original quadric orbit and an explicit invariant lifting map}
\label{sec:OQC}

This calculation uses the original simple-quartet construction, with
$m=1$, $k=4l+1$, $l\geq1$, $q=(k+1)^2$, $0<\delta<1/2$, and
$\gamma>2$. The period parameter, its branch, the polynomial contours,
the complete amplitude unit, and the observation are those of OPG and
OCS. The degree-eight polynomial below is the product of four actual
orbit equations. Its degree and all of its coefficients are derived
from that orbit; no source order is selected to match an estimate.

\subsection{The actual coefficient ring and the original observation}
Order the original roots as $(++),(+-),(-+),(--)$, with
$\alpha_{\epsilon,\eta}=1/2+\epsilon\delta+i\eta\gamma$.
Let $\mathcal E_{\rm prim}$ have columns the four original primary
idempotents in $E_h=\mathbb C[z_1]/(h)$. For the original matrices
$\Pi,U,F$ in OCS2, OCS5, OCS27 put
\[
 \mathsf H=F^{-1}\Pi U\mathcal E_{\rm prim},\qquad
 t_{\epsilon,\eta}(w)=e^{(\eta\gamma-i\epsilon\delta)w},\qquad
 x(w)=\mathsf Ht(w),\qquad
 Q(x)=Q_0(\mathsf H^{-1}x),\quad
 Q_0(t)=t_{++}t_{--}-t_{+-}t_{-+}.
 \tag{OQC1}
\]
Every matrix in $\mathsf H$ is the specified original isomorphism;
in particular its nonzero amplitude factors are the actual
$v_h(\alpha)$, and its period entries are the original convergent
polynomial integrals. Thus $\det\mathsf H\ne0$ by the included
period-isomorphism and primary-basis proofs. Define
\[
 \mathscr P_d=\mathbb C[x_1,x_2,x_3,x_4]_d,\qquad
 R_d=\mathscr P_d/(Q\mathscr P_{d-2}),\qquad
 \mathscr P_d=R_d=0\quad(d<0).
 \tag{OQC2}
\]
The matrix of $Q_0(t)=t^TM_0t$ has entries
$(M_0)_{14}=(M_0)_{41}=1/2$ and
$(M_0)_{23}=(M_0)_{32}=-1/2$, all others zero. Hence
\[
 M=\mathsf H^{-T}M_0\mathsf H^{-1},\qquad
 \det M=\frac{1}{16(\det\mathsf H)^2}\ne0.
 \tag{OQC3}
\]
In particular $Q$ is irreducible. Indeed, a reducible homogeneous
quadratic over $\mathbb C$ is a product of two linear forms, whose
symmetric coefficient matrix is a sum of two rank-one matrices and
has rank at most two. This contradicts OQC3. The polynomial ring
over a field is a unique factorization domain, by the repeated
one-variable Gauss lemma, so its irreducible element $Q$ is prime.
Consequently $R=\mathbb C[x_1,x_2,x_3,x_4]/(Q)$ is an integral domain.

For clarity, the complete homogeneous restriction calculation has
the precise kernel $Q\mathscr P_{d-2}$. In the $t$ coordinates a
degree-$d$ monomial restricts to
\[
 \exp\bigl(((2b-d)\gamma-i(2a-d)\delta)w\bigr),
       \qquad 0\leq a,b\leq d.
 \tag{OQC4}
\]
Every pair occurs: the number of $(++)$ factors can range from
$\max(0,a+b-d)$ to $\min(a,b)$. The exponents are distinct by
their real and imaginary parts. Their exponentials are linearly
independent, since differentiating orders zero through $(d+1)^2-1$
gives a Vandermonde matrix with nonzero determinant. Two monomials
giving the same pair have exponent vectors differing by a multiple
of $(1,-1,-1,1)$. After their common monomial is removed, their
difference is divisible by $t_{++}t_{--}-t_{+-}t_{-+}$ through
$A^r-B^r=(A-B)\sum_{j=0}^{r-1}A^{r-1-j}B^j$. A vanishing
polynomial has coefficient sum zero within each pair by the
independence just proved, and is therefore a sum of these differences.
Conversely $Q_0(t(w))=0$. Transport by $\mathsf H$ proves the
stated kernel and
\[
                  \dim R_d=(d+1)^2\quad(d\geq0).
 \tag{OQC5}
\]

Write $\zeta=e^{2\pi i/5}$ and
$D=\operatorname{diag}(\zeta,\zeta^2,\zeta^3,\zeta^4)$.
The literal Fourier calculation gives $CF=FD$. On polynomials
define the original induced action and its average by
\[
 (g f)(x)=f(D^{-1}x),\qquad g^5=1,\qquad
 \mathsf R f=\frac15\sum_{a=0}^4g^af,
 \qquad \mathscr I_d=\mathscr P_d^{\langle g\rangle}.
 \tag{OQC6}
\]
Thus $\mathscr I_d$ consists of the weight-zero monomials, since
$gx^\nu=\zeta^{-\sum j\nu_j}x^\nu$. This is the dual of the
original average of $(C^{-a})^{\otimes d}$; both the inverse powers
and the factor $1/5$ are retained.

For the original algebra $E_d$ of the $d$-fold sum, define
$\Theta_d:R_d\xrightarrow{\sim}E_d^\vee$ by
\[
 \Theta_d([f])(e^{wY_d}[1])=f(x(w)),\qquad
       Y_d=(A_d-dI/2)/i.
 \tag{OQC7}
\]
Here $\vee$ denotes the complex-linear dual. The primary expansion
of the vector on the left gives all distinct exponentials OQC4;
therefore existence and uniqueness of this covector and bijectivity
follow from the preceding complete restriction calculation. The
original observation $\Lambda_d:E_d\to B_d$ obeys
\[
 J_d:=\operatorname{im}(\mathscr I_d\to R_d),\qquad
 \Theta_d(J_d)=\operatorname{im}\Lambda_d^\vee,
 \qquad \dim J_d=\operatorname{rank}\Lambda_d,
 \qquad R_d/J_d\xrightarrow{\sim}(\ker\Lambda_d)^\vee.
 \tag{OQC8}
\]
Indeed the coordinate functions of the projected pure tensor in
OCS9 are precisely $x(w)^\nu$ of weight zero; their pulled-back
covectors span the dual of its actual image. A covector vanishing
on $\ker\Lambda_d$ factors through that image, which proves the
last isomorphism as well. OQC8 specifies the exact original image;
it has not been replaced by the whole invariant tensor space.

\subsection{The complete orbit decision from the actual coefficients}
Put $Q_a=g^aQ$ for $0\leq a<5$. The orbit of the line $\mathbb C Q$
has size one or five because five is prime. Moreover a size-one
orbit has $gQ=Q$, with no other possible scalar. To prove this, write
$gQ=\lambda Q$. Then $\lambda^5=1$, while OQC3 and the determinant
of $D$, which is $\zeta^{10}=1$, imply
\[
 \det(D^{-T}MD^{-1})=\det M=\lambda^4\det M.
 \tag{OQC9}
\]
Thus $\lambda^4=1$ and $\lambda^5=1$, forcing $\lambda=1$.
If any nonidentity power fixed the line, its powers generate the
same cyclic group, so the preceding argument applies. This proves
the entire orbit classification.

There is an exact coefficient test, requiring no genericity premise.
Write the actual $Q=\sum_{|\mu|=2}q_\mu x^\mu$ and put
\[
 \mathfrak e(Q)=\sum_{\substack{|\mu|=2\\
                      \sum j\mu_j\not\equiv0\ (5)}}|q_\mu|^2.
 \tag{OQC10}
\]
All coefficients here are the entries of OQC3, including the full
inverse of the original period matrix. Since the monomials are a
basis and each summand is nonnegative, $\mathfrak e(Q)=0$ exactly
when $Q$ is invariant. The only degree-two weight-zero monomials
are $x_1x_4$ and $x_2x_3$. Hence the size-one case is exactly
$Q=a x_1x_4+b x_2x_3$ with $ab\ne0$; nonvanishing of the two
coefficients follows from OQC3. Positive $\mathfrak e(Q)$ gives
the five distinct irreducible orbit equations. These are exhaustive
outcomes of a test on the given periods, not assumptions on them.

In the size-one outcome, multiplication by $Q$ preserves every
weight. If $Qf$ is invariant, then $Q(gf-f)=0$ in the polynomial
ring, and $gf=f$. Therefore
\[
 \ker(\mathscr I_k\to R_k)=Q\mathscr I_{k-2},\qquad
 \operatorname{rank}\Lambda_k=d_{5,k,0}-d_{5,k-2,0},
 \quad \dim\ker\Lambda_k=q-d_{5,k,0}+d_{5,k-2,0}.
 \tag{OQC11}
\]
This also proves that $J_k=R_k^{\langle g\rangle}$: an invariant
class can be lifted to a polynomial, and averaging that lift gives
an invariant representative. The count itself follows from
\[
 d_{5,k,0}=\frac{\binom{k+3}{3}+4\epsilon_k}{5},\qquad
 \epsilon_k=\mathbf1_{k\equiv0\ (5)}-\mathbf1_{k\equiv1\ (5)}.
 \tag{OQC12}
\]
For a proof, the identity element has symmetric-power generating
series $(1-z)^{-4}$; each of the other four powers has series
$\prod_{j=1}^4(1-z\zeta^j)^{-1}=(1-z)/(1-z^5)$.
Averaging their coefficients gives OQC12. Subtraction and
$\binom{k+3}{3}-\binom{k+1}{3}=(k+1)^2$ evaluate the rank in
OQC11, for every $k\geq2$, as
\[
 \left\lfloor\frac{(k+1)^2}{5}\right\rfloor
                 +\mathbf1_{k\equiv0\text{ or }3\ (5)}.
 \tag{OQC13}
\]
The five numerator corrections before division by five are
$4,-4,-4,4,0$, respectively; the residues of $(k+1)^2$ are
$1,4,4,1,0$. This verifies every case in OQC13.

\subsection{An explicit invariant lift for the five-member orbit}
For the size-five outcome determined by OQC10, retain the full
actual products
\[
 \mathcal H_Q=\prod_{a=0}^4Q_a\quad(\deg\mathcal H_Q=10),
 \qquad \mathcal A_Q=\prod_{a=1}^4Q_a
                                      \quad(\deg\mathcal A_Q=8).
 \tag{OQC14}
\]
The action permutes all five factors cyclically, so
$g\mathcal H_Q=\mathcal H_Q$. None of the four factors of
$\mathcal A_Q$ is a scalar multiple of $Q$. Since each is
irreducible and $Q$ is prime, multiplication by $[\mathcal A_Q]$
is injective on the integral domain $R$.

For every homogeneous $f\in\mathscr P_d$ the exact trace identity is
\[
 \left[5\mathsf R(\mathcal A_Q f)\right]_{(Q)}
                          =[\mathcal A_Q f]_{(Q)}.
 \tag{OQC15}
\]
Indeed $g^a\mathcal A_Q$ is the product of all orbit factors
except $Q_a$. For $a\ne0$ it contains $Q_0=Q$. Those four
terms vanish in the displayed quotient, and the $a=0$ term is
exactly the right side. This proves OQC15 with its original
average and full factors, including every period-dependent
coefficient.

The lift also has a well-defined typed quotient map. Let
$\mathscr U=\mathbb C[x_1,x_2,x_3,x_4]/(\mathcal H_Q)$;
the invariant equation $\mathcal H_Q$ makes the cyclic action
well defined on this quotient. Then
\[
 \begin{gathered}
 \mathfrak L_d:R_d\longrightarrow
                      (\mathscr U_{d+8})^{\langle g\rangle},\qquad
 [f]_{(Q)}\longmapsto[5\mathsf R(\mathcal A_Q f)]_{(\mathcal H_Q)},\\
 \operatorname{res}\mathfrak L_d([f])=[\mathcal A_Q f]_{(Q)},
 \qquad \operatorname{res}:\mathscr U\to R.
 \end{gathered}\tag{OQC16}
\]
If $f$ changes by $Qb$, its proposed image changes by
$5\mathsf R(\mathcal H_Qb)=5\mathcal H_Q\mathsf Rb$, hence
by zero in $\mathscr U$. This proves well-definedness. The
image is invariant because it has an invariant polynomial
representative; OQC15 proves the second line. Injectivity follows
from the injectivity of multiplication by $[\mathcal A_Q]$ on
$R$. All maps in OQC16 are thus established, including their kernels.

The actual original invariant image consequently contains
\[
 [\mathcal A_Q]R_{k-8}\ \subseteq\ J_k\ \subseteq\ R_k.
 \tag{OQC17}
\]
The first inclusion follows by taking the invariant polynomial
representatives OQC15, and uses no dimension inference. Therefore,
for every $k\geq8$ in this outcome,
\[
 \operatorname{rank}\Lambda_k\geq(k-7)^2,\qquad
 \dim\ker\Lambda_k\leq(k+1)^2-(k-7)^2=16k-48.
 \tag{OQC18}
\]
These inequalities follow from the proved injection and OQC5.
They may be combined with the unconditional actual-matrix bounds
in OCS24, OCS31--33 by taking the stronger lower and upper bounds.
For the original degrees $k=9,13,17,21,25,29$, the displayed upper
bounds on the original kernel are respectively
$96,160,224,288,352,416$; at the first two degrees the earlier
unconditional bounds may be stronger. No such upper bound is
assigned to the size-one outcome, whose exact rank is OQC11.
Together OQC10--18 compute the complete orbit alternatives from
the original coefficients. Evaluating $\mathfrak e(Q)$ in the
original period family is a separate calculation on those explicit
coefficients; no outcome is postulated here.

\subsection{The induced original dual map and its complete action defect}
At $k\geq9$ the degree $d=k-8$ is positive, and the original
$d$-fold sum algebra is defined using the same $h,\Pi,U$.
For the size-five outcome, OQC7, OQC8 and OQC17 give the uniquely
typed injective map
\[
 \begin{gathered}
 \mathfrak J_k:E_{k-8}^\vee\hookrightarrow B_k^\vee,\\
 \Lambda_k^\vee\mathfrak J_k
        =\Theta_k\,M_{[\mathcal A_Q]}\,\Theta_{k-8}^{-1},\qquad
 \mathfrak J_k^\vee:B_k\twoheadrightarrow E_{k-8}.
 \end{gathered}\tag{OQC19}
\]
The right side of the second line has image in
$\operatorname{im}\Lambda_k^\vee$ by OQC17, and
$\Lambda_k^\vee$ is injective because $\Lambda_k$ is onto its
actual image. This gives existence and uniqueness. Each of the
three right-hand maps is injective, so $\mathfrak J_k$ is
injective. Finite-dimensional duality then proves the stated
surjection, using the canonical bidual identifications. The new
degree $k-8$ records contraction against the degree-eight
polynomial derived in OQC14. It does not replace either original
Gamma source of order $1$ or $k$.

The surviving action and its defect are fully determined. Put
\[
 \Omega=\operatorname{diag}(\gamma-i\delta,-\gamma-i\delta,
                              \gamma+i\delta,-\gamma+i\delta),
 \quad G=\mathsf H\Omega\mathsf H^{-1},\qquad
 (\mathcal Df)(x)=\sum_{j=1}^4(Gx)_j\partial_jf(x).
 \tag{OQC20}
\]
Here $x'(w)=Gx(w)$ by OQC1. In the original $t$ coordinates,
the weights of each product in $Q_0$ sum to zero, so
$\mathcal DQ=0$. Thus this derivation acts on each $R_d$.
Differentiating OQC7 proves exactly
\[
          \Theta_d\mathcal D=Y_d^\vee\Theta_d.
 \tag{OQC21}
\]
Both sides evaluated at $e^{wY_d}[1]$ are the derivative of
$f(x(w))$, and these vectors span $E_d$ by the distinct
exponentials; hence equality holds as maps, not only on a curve.
Leibniz's rule gives the complete defect
\[
 \begin{split}
 Y_k^\vee\Lambda_k^\vee\mathfrak J_k
       -\Lambda_k^\vee\mathfrak J_kY_{k-8}^\vee
 &=\Theta_k M_{[\mathcal D\mathcal A_Q]}\Theta_{k-8}^{-1},\\
 \mathcal D\mathcal A_Q
 &=\sum_{a=1}^4(\mathcal DQ_a)
                  \prod_{\substack{1\leq b\leq4\\b\ne a}}Q_b.
 \end{split}\tag{OQC22}
\]
The codomain in the first line is $E_k^\vee$. Its image has not
been asserted to lie in $\operatorname{im}\Lambda_k^\vee$.
Restriction to $\ker\Lambda_k$ supplies the exact remaining
map into $(\ker\Lambda_k)^\vee$, through OQC8. In particular
OQC19 is not used as an unproved intertwining of boundary actions.
OQC22 retains every term that can leave the observed dual subspace.

Finally, the failure of the cyclic average to commute with the
original action is explicitly calculable before any quotient:
\[
 ([\mathcal D,g^a]f)(x)
   =(\nabla f)(D^{-a}x)^T
                   (D^{-a}G-GD^{-a})x,
 \qquad
 [\mathcal D,\mathsf R]=\frac15\sum_{a=0}^4[\mathcal D,g^a].
 \tag{OQC23}
\]
The chain rule applied separately to $f(D^{-a}x)$ and to
$(\mathcal Df)(D^{-a}x)$ proves both identities. Thus all
action comparisons in the invariant lift have explicit original
coefficient maps. The rank estimates in OQC18 concern exactly
the kernel used in the original Gamma and arithmetic quotient
Grams, because those constructions use the same $\Lambda_k$.
An estimate for its volume still requires those original metrics;
no volume conclusion has been inferred solely from its dimension.

% END RX INLINE 55


% BEGIN RX INLINE 56; SHA256 d90037182769f7a7ab65b3443de6339fa28eea1aa6e4b4ac0615d9d94ae9a34b
% Exact completion of the orbit rank calculation. No period outcome is assumed.
\section{Exact original rank from the complete quadric orbit}
\label{sec:OQX}

Retain every original object and map in OQC1--23. In particular
$m=1$, $k=4l+1$, $l\geq1$, $q=(k+1)^2$,
$\mathscr P=\mathbb C[x_1,x_2,x_3,x_4]$, and
$g f(x)=f(D^{-1}x)$ for the original Fourier matrix
$D=\operatorname{diag}(\zeta,\zeta^2,\zeta^3,\zeta^4)$.
The polynomial $Q=Q_0(\mathsf H^{-1}x)$ has its actual original
period and amplitude coefficients. Write
$\mathscr I=\mathscr P^{\langle g\rangle}$,
$R=\mathscr P/(Q)$ and $J=\operatorname{im}(\mathscr I\to R)$,
with the grading and the actual observation identification of OQC8.
For negative indices put $\mathscr I_d=0$ and $d_{5,d,0}=0$.

The orbit decision is the exhaustive test OQC10. The size-one
case already has the exact rank OQC11--13. The following calculation
completes the other case using the same original coefficients. It
does not assert that a period parameter belongs to that case before
its coefficient test has been evaluated.

\subsection{Every invariant vanishing polynomial contains the full orbit}

In the size-five case, let $Q_a=g^aQ$ and retain exactly
$\mathcal H_Q=\prod_{a=0}^4Q_a$ from OQC14. Each $Q_a$ is
irreducible, because an invertible linear substitution preserves
factorizations and $Q$ is irreducible by OQC3. The five factors are
pairwise nonassociate by the proved orbit classification, and each
is prime in the original polynomial ring. Consequently
\[
 \boxed{\quad
 \mathscr I\cap Q\mathscr P=\mathcal H_Q\mathscr I,
 \qquad
 \ker(\mathscr I_d\longrightarrow R_d)
                       =\mathcal H_Q\mathscr I_{d-10}.
 \quad}\tag{OQX1}
\]
Here is the complete divisibility proof. If $f\in\mathscr I$ and
$f=Qb$, then for every $a=0,\ldots,4$ the exact equality
$g^af=f$ gives $f=Q_a g^ab$. Thus every one of the five original
prime factors divides $f$. Distinctness of the factors implies
that their product divides $f$: after factoring out $Q_0$,
primality of $Q_1$ and the fact that $Q_1$ does not divide $Q_0$
force $Q_1$ to divide the remaining quotient. Repeating this
argument with $Q_2,Q_3,Q_4$ proves
$f=\mathcal H_Q b_0$. The action permutes the five exact factors
cyclically, so $g\mathcal H_Q=\mathcal H_Q$. Applying $g$ to
this last equality and using $gf=f$ gives
$\mathcal H_Q(gb_0-b_0)=0$. The polynomial ring is an integral
domain and $\mathcal H_Q\ne0$, whence $gb_0=b_0$.
This proves the first inclusion in OQX1. Conversely an invariant
$b_0$ gives an invariant $\mathcal H_Qb_0$ divisible by $Q$,
which proves the reverse inclusion. The product is homogeneous
of degree ten. If its product with $b_0$ is homogeneous of degree
$d$, decomposition of $b_0$ into homogeneous parts shows that
only its part of degree $d-10$ can be nonzero: each other part
has zero product with $\mathcal H_Q$ and hence is zero.
This proves the graded kernel formula, including $d<10$.

\subsection{The exact invariant quotient and the restriction map}

Let $\mathscr U=\mathscr P/(\mathcal H_Q)$ as in OQC16. The
following are isomorphisms of graded complex algebras with their
specified maps:
\[
 \boxed{\quad
 \mathscr I/(\mathcal H_Q\mathscr I)
 \xrightarrow{\ [f]\mapsto[f]_{(\mathcal H_Q)}\ }
       \mathscr U^{\langle g\rangle}
 \xrightarrow{\ \operatorname{res}\ } J\subseteq R.
 \quad}\tag{OQX2}
\]
The action on $\mathscr U$ is defined because $\mathcal H_Q$ is
invariant. If $[f]_{(\mathcal H_Q)}$ is invariant, averaging its
representatives gives
$[\mathsf R f]_{(\mathcal H_Q)}=[f]_{(\mathcal H_Q)}$, since
$[g^af]=[f]$ for every $a$. Thus the first map is onto.
Its kernel before passage to the displayed domain is
$\mathscr I\cap\mathcal H_Q\mathscr P
=\mathcal H_Q\mathscr I$: the reverse inclusion is immediate,
and for $f=\mathcal H_Qb\in\mathscr I$, cancellation in
$\mathcal H_Q(gb-b)=0$ proves $b\in\mathscr I$.
This proves that the first map is injective as well.

The full quotient restriction
$\operatorname{res}:\mathscr U\to R$ is well defined because
$(\mathcal H_Q)\subseteq(Q)$; its kernel is exactly
$Q\mathscr P/\mathcal H_Q\mathscr P$.
For its invariant subspace, choose the invariant representative
just constructed. If its restriction vanishes, that representative
belongs to $\mathscr I\cap Q\mathscr P$, which is
$\mathcal H_Q\mathscr I$ by OQX1. Its class in $\mathscr U$ is
therefore zero. The invariant restriction is injective. Its image
is exactly $J$, since an invariant representative restricts to
$J$ by definition and every element of $J$ has such a representative.
Both maps are induced by polynomial quotient maps, so they preserve
products, sums, complex scalars and degrees. No action of $g$ on
$R$ is asserted in the size-five case. In particular the second
map retains precisely the invariant image of the original
observation, with its complete restriction kernel already computed.

\subsection{All original degrees and the exact kernel dimension}

The map $\Theta_k$ and the original observation in OQC8 identify
$\dim J_k$ with $\operatorname{rank}\Lambda_k$ and identify
$R_k/J_k$ with $(\ker\Lambda_k)^\vee$. Multiplication by the
nonzero polynomial $\mathcal H_Q$ is injective in $\mathscr P$.
Taking dimensions in OQX1 thus gives, for every nonnegative degree
$d$, the exact invariant restriction rank
$\dim J_d=d_{5,d,0}-d_{5,d-10,0}$. At the original degrees this
evaluates to
\[
 \boxed{\quad
 \operatorname{rank}\Lambda_k
      =d_{5,k,0}-d_{5,k-10,0}
      =k^2-6k+17,\qquad
 \dim\ker\Lambda_k=8k-16,
 \quad k=4l+1,\ l\geq1,
 \quad}\tag{OQX3}
\]
in the size-five outcome of the actual coefficient test.
To verify every constant and degree range, for $k\geq10$ the
complete character formula OQC12 gives
\[
 d_{5,k,0}-d_{5,k-10,0}
 =\frac15\left\{\binom{k+3}{3}-\binom{k-7}{3}
                         +4(\epsilon_k-\epsilon_{k-10})\right\}.
\]
The two $\epsilon$ values are equal because their indices differ
by ten. Expanding the remaining cubics gives
\[
 \binom{k+3}{3}-\binom{k-7}{3}
 =\frac{(k^3+6k^2+11k+6)
             -(k^3-24k^2+191k-504)}6
 =5(k^2-6k+17).
\]
The only original degrees below ten are five and nine. Their
negative-degree terms in OQX1 vanish, while the same exact
character formula gives
\[
 d_{5,5,0}=(56+4)/5=12=5^2-6\cdot5+17,
 \qquad
 d_{5,9,0}=220/5=44=9^2-6\cdot9+17.
\]
This proves the first equality of OQX3 for every original $l\geq1$.
Since $\dim E_k=\dim R_k=(k+1)^2$, rank-nullity gives
$(k+1)^2-(k^2-6k+17)=8k-16$, proving the second equality.
For the original degrees already tabulated in OCS26, the exact
five-orbit values are
\[
 \begin{array}{c|r|r}
 k&\operatorname{rank}\Lambda_k&\dim\ker\Lambda_k\\\hline
 5&12&24\\
 9&44&56\\
 13&108&88\\
 17&204&120\\
 21&332&152\\
 25&492&184\\
 29&684&216
 \end{array}
\]
These are exact values for this exhaustively classified orbit
outcome, replacing the less precise bounds OQC18 wherever that
outcome occurs. The full invariant lift and its action defect
OQC15--23 remain valid with all original coefficients. The
size-one outcome retains its different exact values OQC11--13.
No choice between the two outcomes has been made here without
evaluating the original period coefficient test, and no estimate
for a metric volume has been inferred from its dimension alone.

% END RX INLINE 56


% BEGIN RX INLINE 57; SHA256 5ef06d49d4b576f969a467fd4fb239b5af908ef188e8836d6dbff93836f3959d
\section{The actual period quadric and its cyclic continuation}
\label{sec:PQO}

This calculation uses the original simple quartet, $m=1$, and the
literal period matrix, amplitude unit and cyclic matrix of BRU and
OPG. Its formulas evaluate the transported quadric from those data.
All source orders elsewhere in the programme remain unchanged.

\subsection{Original roots, entire unit and centred phase}
Order the roots and keep their actual values as
\[
\begin{gathered}
 \alpha_1=\tfrac12+\delta+i\gamma,\quad
 \alpha_2=\tfrac12+\delta-i\gamma,\quad
 \alpha_3=\tfrac12-\delta+i\gamma,\quad
 \alpha_4=\tfrac12-\delta-i\gamma,\\
 0<\delta<\tfrac12,\quad\gamma>2,\quad
 h(s)=\prod_{j=1}^4(s-\alpha_j),\quad g=2\xi,\quad v_h=g/h.
\end{gathered}\tag{PQO1}
\]
The four zero orders of $g$ are complete and equal to one.
Thus $v_h$ is entire and has nonzero value at each displayed root.
Put $w=s-1/2$, and define the following original constants:
\[
\begin{gathered}
 z_+=\delta+i\gamma,\quad z_-=\delta-i\gamma,\quad
 \lambda=z_+^2,\quad\mu=z_-^2,\\
 A=2(\gamma^2-\delta^2),\qquad B=(\delta^2+\gamma^2)^2,\qquad
 \mathfrak a=\Phi_h(1/2)=\frac1{160}+\frac A{24}+\frac B2,\\
 h(1/2+w)=w^4+Aw^2+B,\qquad
 \Phi_h(1/2+w)=\mathfrak a+\frac{w^5}{5}+\frac{Aw^3}{3}+Bw.
\end{gathered}\tag{PQO2}
\]
The product of the two quadratics
$((w-\delta)^2+\gamma^2)((w+\delta)^2+\gamma^2)$ proves the
quartic identity. Integrating it and using $\Phi_h(0)=0$
proves the phase and constant in PQO2. In particular the translation
has not removed the original constant $\mathfrak a$.

The theta integral IPM.2 has real coefficients, real theta kernel,
and an expression unchanged by $s\mapsto1-s$. It therefore gives
$g(\bar s)=\overline{g(s)}$ and $g(1-s)=g(s)$ directly; the same
two identities hold for $h$. Taking their entire quotient gives
\[
 v_h(\alpha_1)=v_h(\alpha_4)=a,\qquad
 v_h(\alpha_2)=v_h(\alpha_3)=\bar a,\qquad a\ne0.
 \tag{PQO3}
\]
These are the full actual unit values. Their magnitude and phase
are not assigned. Set $t_+=a^{-2}$ and $t_-=\bar a^{-2}$.
For $j=0,1,2,3$ retain
\[
 \sigma_j=t_+\lambda^j-t_-\mu^j,\qquad
 \sigma_2=-A\sigma_1-B\sigma_0,\qquad
 \sigma_3=(A^2-B)\sigma_1+AB\sigma_0.
 \tag{PQO4}
\]
The recurrence follows because $\lambda+\mu=-A$ and
$\lambda\mu=B$. Each $\sigma_j$ is purely imaginary.
Furthermore $\sigma_0$ and $\sigma_1$ cannot both vanish:
the first would give $t_+=t_-\ne0$, and the second would then
give $t_+(\lambda-\mu)=0$, whereas
$\lambda-\mu=4i\delta\gamma\ne0$.

\subsection{The original periods as an absolutely convergent series}
Fix $u\ne0$ on its original logarithm branch. Let
$\ell_j(r)=r\exp(i(\arg u+\pi+2\pi j)/5)$,
$\Gamma_j=-\ell_0+\ell_j$, $1\leq j\leq4$.
Let $\Pi$ integrate the degree-less-than-four representative
against $e^{\Phi_h(s)/u}\,ds$ on these original oriented contours.
Write $F_{ja}=\zeta^{ja}-1$, $\zeta=e^{2\pi i/5}$,
$1\leq j,a\leq4$. Direct finite geometric sums prove
\[
 (F^{-1})_{aj}=\frac15\zeta^{-ja},\qquad
 CF=F D,\qquad D=\operatorname{diag}(\zeta,\zeta^2,\zeta^3,\zeta^4),
 \tag{PQO5}
\]
where $(Cx)_j=x_{j+1}-x_1$ for $j<4$ and $(Cx)_4=-x_1$.
For example multiplying the asserted inverse by $F$ gives
$\frac15\sum_{j=1}^4\zeta^{-ja}(\zeta^{jb}-1)=\delta_{ab}$;
substitution into each row gives $CF=FD$.

Let $T$ be the exact increasing-coefficient translation matrix
$T_{rb}=\binom br2^{-(b-r)}$ for $0\leq r\leq b\leq3$, and
zero otherwise. Thus it sends coefficients of $p(s)$ to those
of $p(1/2+w)$, with determinant one. Define the centred
Fourier-period matrix
\[
 P_{ar}(u)=\frac15\sum_{j=1}^4\zeta^{-ja}
 \int_{\Gamma_j}e^{(w^5/5+Aw^3/3+Bw)/u}w^r\,dw,
 \quad 1\leq a\leq4,\quad0\leq r\leq3.
 \tag{PQO6}
\]
The contours in this formula start at $w=0$. They give precisely
the translated original period by
\[
 F^{-1}\Pi=e^{\mathfrak a/u}P(u)T.
 \tag{PQO7}
\]
To justify this exact contour translation, the substitution
$w=s-1/2$ first gives rays starting at $-1/2$. Join each shifted
ray at a large radius to the corresponding unshifted ray. On the
joining segment their arguments differ by $O(1/r)$, and the
leading real phase is $-r^5/(5|u|)+O(r^4)$; all lower phase
terms and polynomial factors are dominated by this negative
leading term. Those outer joining integrals tend to zero.
The two common-origin joining integrals cancel in
$-\ell_0+\ell_j$ with their opposite orientations. Cauchy's
theorem for the entire integrand proves equality in the limit.
Substitution of the polynomial gives exactly the matrix $T$.

For every integer $L\geq1$ define, on the same branch,
\[
 G_L(u)=u^{L/5}5^{L/5-1}e^{\pi iL/5}\Gamma(L/5).
 \tag{PQO8}
\]
Integration of the leading phase on the two rays gives
$\int_{\Gamma_j}e^{w^5/(5u)}w^{L-1}dw
 =G_L(u)(\zeta^{jL}-1)$ by the real substitution
$r^5/(5|u|)$ on each ray, including its phase in $dw$.
Expanding the remaining two exponentials now gives the exact
entries
\[
\boxed{
 P_{ar}(u)=
 \sum_{\substack{p,q\geq0\\r+3p+q+1\equiv a\pmod5}}
 \frac{A^pB^q}{3^p p!q!u^{p+q}}\,
                 G_{r+3p+q+1}(u).
}\tag{PQO9}
\]
Terms with $r+3p+q+1\equiv0$ have zero two-ray integral
in every column and contribute zero before the Fourier map.
To prove every exchange, the sum of the absolute values of the
unintegrated series on either ray is bounded by a constant times
\[
 R^r\exp\left(-\frac{R^5}{5|u|}
             +\frac{|A|R^3}{3|u|}+\frac{|B|R}{|u|}\right).
\]
Here $R$ is the radial coordinate and $r$ is the fixed column
degree. This integrable majorant
proves absolute termwise integration. It is uniform on compact
coefficient and nonzero-$u$ sets on the chosen branch. The finite
Fourier sum in PQO5 then selects exactly the stated residue class.
Thus PQO9 is an actual period formula, with all coefficients and
all terms retained.

For additional finite computation, let $I_d$ denote the vector
of the four centred ray moments of $w^d$. Integration of the
derivative of $e^{(w^5/5+Aw^3/3+Bw)/u}w^d$, with the same
vanishing outer endpoints and cancelling inner endpoints, proves
\[
 I_{d+4}+A I_{d+2}+B I_d=-ud I_{d-1}\quad(d\geq0),
 \qquad u\partial_B I_d=I_{d+1},\quad
 3u\partial_A I_d=I_{d+3}.
 \tag{PQO10}
\]
The first right side is zero at $d=0$. Coefficient derivatives
are justified by the same compact domination. These are the
original polynomial period recurrences, and they compute every
higher moment from the first four and the actual coefficients.

The determinant is especially explicit in these centred
coordinates:
\[
 \det\Pi=-(2\pi)^2u^2e^{4\mathfrak a/u},\qquad
 \det F=-25\sqrt5,\qquad
 \boxed{\det P(u)=\frac{(2\pi)^2}{25\sqrt5}\,u^2.}
 \tag{PQO11}
\]
For the first equality use the full PDE1--13 determinant
calculation at $n=4$: the phase is $e^{15\pi i}=-1$ and
the critical-value sum is $4\mathfrak a$. This critical sum
also follows directly by pairing $w$ with $-w$ in the odd
centred phase. For $F$, subtract row zero from the other rows
in the five-by-five Fourier Vandermonde and expand its first
column. Its phase is $e^{13\pi i}=-1$ and its magnitude
$5^{5/2}$, by the finite Fourier orthogonality identity.
The last equality follows by taking determinants in PQO7 and
using $\det T=1$. The Gamma grid-product proof PGCC1--6 retains
the full evaluated constant in the first equality. In particular
every inverse below exists on the full stated domain.

\subsection{The actual transformed quadric and every deciding coefficient}
Let $V$ be the centred evaluation matrix with rows
$(1,z_j,z_j^2,z_j^3)$, where
$(z_1,z_2,z_3,z_4)=(z_+,z_-,-z_-,-z_+)$. Its inverse sends
the four primary values to their unique degree-less-than-four
polynomial. Thus the actual matrix from primary coordinates
to the Fourier-period space is
\[
 B_{\rm per}=F^{-1}\Pi UJ_{\rm prim}
       =e^{\mathfrak a/u}P(u)V^{-1}
                         \operatorname{diag}(a,\bar a,\bar a,a).
 \tag{PQO12}
\]
This follows either by applying all maps to the four CRT
idempotents or by evaluating the full unit at each root.
It agrees with OCS and OPG in the given original orientation.

For $Q_0(t)=t_1t_4-t_2t_3$, define
$Q_u(x)=Q_0(B_{\rm per}^{-1}x)$. On the unchanged original
period coordinate $\mathbf y\in\mathbb C^4$ the quadratic is
$Q_u^{\rm per}(\mathbf y)=Q_0((\Pi UJ_{\rm prim})^{-1}\mathbf y)
 =Q_u(F^{-1}\mathbf y)$, because $FB_{\rm per}=\Pi UJ_{\rm prim}$.
Thus the exact isomorphism $\mathbf y=Fx$ carries the following
quadric to the original period space. If
$p(w)=p_0+p_1w+p_2w^2+p_3w^3$, multiplication gives
\[
 p(w)p(-w)=p_0^2+(2p_0p_2-p_1^2)w^2
                     +(p_2^2-2p_1p_3)w^4-p_3^2w^6.
\]
Consequently the exact symmetric matrix in the centred
coefficient frame is
\[
 S=\begin{pmatrix}
 \sigma_0&0&\sigma_1&0\\
 0&-\sigma_1&0&-\sigma_2\\
 \sigma_1&0&\sigma_2&0\\
 0&-\sigma_2&0&-\sigma_3
 \end{pmatrix},\qquad
 \boxed{Q_u(x)=x^{\mathsf T}H(u)x,
 \quad H(u)=e^{-2\mathfrak a/u}P(u)^{-\mathsf T}S P(u)^{-1}.}
 \tag{PQO13}
\]
Indeed applying the preceding product formula at $z_+$ and
$z_-$ with multipliers $t_+$ and $-t_-$ gives exactly $S$.
This retains both actual unit values through every $\sigma_j$.
The two parity blocks have determinants
\[
 \sigma_0\sigma_2-\sigma_1^2
          =-t_+t_-(\lambda-\mu)^2,\qquad
 \sigma_1\sigma_3-\sigma_2^2
          =-t_+t_-B(\lambda-\mu)^2.
\]
Thus $S$ and $H(u)$ are nondegenerate, with the exact determinant
\[
 \det H(u)=e^{-8\mathfrak a/u}
       \frac{B(t_+t_-)^2(\lambda-\mu)^4}{(\det P(u))^2}\ne0.
 \tag{PQO14}
\]

Write $N(u)=\operatorname{adj}P(u)$, with entries $N_{ra}$,
$0\leq r\leq3$, $1\leq a\leq4$, and retain
$\Delta(u)=\det P(u)$ from PQO11. Then every coefficient is
given by the explicit minor formula
\[
\begin{aligned}
 H_{ab}(u)=\frac{e^{-2\mathfrak a/u}}{\Delta(u)^2}
 \{&\sigma_0N_{0a}N_{0b}
 +\sigma_1(N_{0a}N_{2b}+N_{2a}N_{0b}-N_{1a}N_{1b})\\
 &+\sigma_2(N_{2a}N_{2b}-N_{1a}N_{3b}-N_{3a}N_{1b})
                         -\sigma_3N_{3a}N_{3b}\}.
\end{aligned}\tag{PQO15}
\]
This is matrix multiplication in PQO13 using
$P^{-1}=N/\Delta$. Each $N_{ra}$ is the signed original
three-by-three period minor. Its entries are given by the
absolutely convergent series PQO9 or by the original integrals
PQO6. The polynomial coefficient of $x_a^2$ is $H_{aa}$;
that of $x_ax_b$, $a<b$, is $2H_{ab}$. Both factors are
retained in this specified quadratic form.

The exact eight deciding entries are
\[
 \mathcal E(u)=
 (H_{11},H_{22},H_{33},H_{44},H_{12},H_{13},H_{24},H_{34}).
 \tag{PQO16}
\]
The cyclic action on quadrics is the pullback
$Q_u\mapsto Q_u(D^{-1}x)$; its symmetric matrix is
$D^{-\mathsf T}H(u)D^{-1}$. Since $D^5=I$ and $\det D=1$,
the line spanned by a nondegenerate form can be stable only
with character one. In fact if its multiplier is $\eta$, then
$\eta^5=1$ by the group action, while taking determinants
gives $\eta^4=1$ because $\det H(u)\ne0$. Thus $\eta=1$.
Entrywise, invariance is precisely
$(\zeta^{-(a+b)}-1)H_{ab}=0$. The only pairs with
$a+b\equiv0$ are $(1,4)$ and $(2,3)$. Therefore
\[
\boxed{\begin{gathered}
 \text{The quadric is invariant exactly on }
       Z=\{u:\mathcal E(u)=0\};\\
 u\in Z\ \Longrightarrow\
 Q_u(x)=2H_{14}(u)x_1x_4+2H_{23}(u)x_2x_3,
                 \quad H_{14}(u)H_{23}(u)\ne0;\\
 \text{at every }u\notin Z\text{ its orbit consists of five distinct quadrics.}
\end{gathered}}\tag{PQO17}
\]
Every assertion follows from the entry test and the rank-four
determinant. The orbit size divides five; it is one exactly
when the line is stable. Thus no two-member orbit occurs.
Moreover $F^{-1}C^{-1}=D^{-1}F^{-1}$ by PQO5, so this is
exactly the orbit of $Q_u^{\rm per}(C^{-j}\mathbf y)$ in the
original period coordinates, with the same five indexed maps.
PQO15--17 determine the exact exceptional locus by original
period minors and original amplitude values, without assigning
generic values to them.

\subsection{The cyclic action is the actual original continuation in $u$}
Continue the original logarithm of $u$ by $2\pi i$. Its rays
advance from $\ell_j$ to $\ell_{j+1}$, and their oriented
differences advance by the exact original matrix $C$.
For the last row $\ell_5=\ell_0$, giving $-\Gamma_1$.
Equivalently every summand of PQO9 acquires the phase
\[
 e^{2\pi i((r+3p+q+1)/5-p-q)}
       =\zeta^{r+3p+q+1}=\zeta^a.
\]
The original constant $e^{\mathfrak a/u}$ is single-valued
under this continuation. Hence
\[
 P(\log u+2\pi i)=DP(\log u),\quad
 B_{\rm per}(\log u+2\pi i)=DB_{\rm per}(\log u),\quad
 H(\log u+2\pi i)=D^{-\mathsf T}H(\log u)D^{-1}.
 \tag{PQO18}
\]
The contour argument and the convergent series prove the same
map. Thus PQO17 describes the actual period continuation of
this original quadric. It introduces no freely chosen cyclic
action on its coefficients.

\subsection{The large-parameter geometry and a positive-ray obstruction}
Let $\varepsilon=u^{-1/5}$ on the same branch, and set
$D_\varepsilon=\operatorname{diag}(\varepsilon^{-1},
\varepsilon^{-2},\varepsilon^{-3},\varepsilon^{-4})$.
The exact series PQO9 gives
\[
 P(u)=\mathscr P(A\varepsilon^2,B\varepsilon^4)D_\varepsilon,
 \qquad
 \mathscr P(0,0)=\operatorname{diag}(g_1,g_2,g_3,g_4),
 \quad g_a=5^{a/5-1}e^{\pi ia/5}\Gamma(a/5).
 \tag{PQO19}
\]
The entries of $\mathscr P$ are entire in the two displayed
arguments, by the same absolute integrable series majorant.
The same determinant proof PQO11 applies to every complex pair
of lower odd-phase coefficients, since the centred phase remains
odd and its critical values pair to zero with full multiplicities.
Thus their determinant is the nonzero constant
$(2\pi)^2/(25\sqrt5)$, by PQO11 and $\det D_\varepsilon=u^2$.
Their inverse is therefore holomorphic everywhere, given by its
adjugate divided by that precise constant. Near $\varepsilon=0$
it is $\operatorname{diag}(g_a^{-1})+O(\varepsilon^2)$.
This bound follows from its convergent power series after
substituting $(A\varepsilon^2,B\varepsilon^4)$; it holds on
a full small complex disc, with the fixed original $A,B$.

If $\sigma_0\ne0$, PQO13 and PQO19 give
\[
 e^{2\mathfrak a/u}H_{11}(u)
              =\frac{\sigma_0}{g_1^2}\varepsilon^2
                                     +O(\varepsilon^4).
\]
If $\sigma_0=0$, then $\sigma_1\ne0$ by PQO4, and they give
\[
 e^{2\mathfrak a/u}H_{22}(u)
              =-\frac{\sigma_1}{g_2^2}\varepsilon^4
                                     +O(\varepsilon^6).
 \tag{PQO20}
\]
For verification of the orders, insert
$P^{-1}=D_\varepsilon^{-1}\mathscr P^{-1}$ into PQO13.
The middle matrix has entries $S_{rs}\varepsilon^{r+s+2}$.
Its first possible degree is two, from $S_{00}$, and its
next possible degree is four, from $S_{02},S_{20},S_{11}$.
The constant term of $\mathscr P^{-1}$ is diagonal, so the
indicated leading entry is exactly the displayed nonzero
coefficient. Every further contribution has the stated higher
degree. This also proves the expansion when some other
$\sigma_j$ vanish.

Consequently for every fixed original quartet and its full
actual amplitude there exists a finite radius outside which
the actual orbit has five distinct quadrics. This conclusion
holds on every original logarithm branch: the estimate is on
the full complex $\varepsilon$ disc, and its leading coefficient
is nonzero. The set $Z$ of PQO17 is a discrete subset of the
logarithmic $u$ surface. Indeed the decisive entry just proved
is holomorphic on that connected surface and is not identically
zero; its zeros are isolated by its local power series, and
$Z$ is contained in that zero set. The explicit finite-radius
bound from the companion period-connection calculation
strengthens this argument while retaining the same actual
coefficients.

Moreover PQO18 multiplies each deciding entry by its nonzero
root-of-unity factor, so the condition $\mathcal E(u)=0$ is
unchanged by every full logarithm turn. Therefore $Z$ is the
inverse image under $\log u\mapsto u$ of a well-defined set
$Z_{\rm abs}\subset\mathbb C^*$. On a sufficiently small disc
about any nonzero $u$, choose one logarithm branch. The preceding
isolated-zero proof shows that $Z_{\rm abs}$ is discrete and
closed on that disc. The explicit radius in the companion
calculation gives $|u|<R_*$ throughout $Z_{\rm abs}$.
Consequently its only possible accumulation point in the finite
$u$-plane is zero. In particular each compact annulus about
zero contains only finitely many such points: an infinite
subset of that compact set would have a nonzero accumulation
point, contradicting the proved local isolated-zero property.

For positive real $u$ retain the actual lift
$\log u=\ln u+2\pi i\ell_{\rm br}$, $\ell_{\rm br}\in\mathbb Z$.
The exact series has an additional real structure. Every
surviving index in PQO9 is $a+5L$, and its phase at the lift
$\ell_{\rm br}=0$ is $e^{\pi ia/5}(-1)^L$. Its remaining
coefficients are real. The actual lift multiplies row $a$ by
$\zeta^{\ell_{\rm br}a}$, by PQO18. Absolute convergence
therefore proves
\[
 P(u)=E_{\ell_{\rm br}} R(u),\qquad
 E_{\ell_{\rm br}}=D^{\ell_{\rm br}}
 \operatorname{diag}(e^{\pi i/5},e^{2\pi i/5},
                       e^{3\pi i/5},e^{4\pi i/5}),
 \qquad R(u)\text{ real and invertible}.
 \tag{PQO21}
\]
Let $M$ be multiplication by $w$ in the increasing centred
polynomial frame, so
\[
 M=\begin{pmatrix}0&0&0&-B\\1&0&0&0\\0&1&0&-A\\0&0&1&0\end{pmatrix},
 \quad c_3=\frac{t_+-t_-}{4i\delta\gamma},\quad
 c_1=\frac{t_++t_-}{2}-(\delta^2-\gamma^2)c_3.
 \tag{PQO22}
\]
Both $c_1,c_3$ are real by PQO3, and the polynomial
$c_1+c_3 w^2$ takes values $t_+,t_-,t_-,t_+$ at the four
original roots. Hence $c_1M+c_3M^3$ is precisely multiplication
by $w(j_hv_h)^{-2}$, with its entire unit retained.

The companion period-connection proof establishes the exact
identity between cyclic invariance of $Q_u$ and
\[
 [P(u)(c_1M+c_3M^3)P(u)^{-1},D]=0.
 \tag{PQO23}
\]
It proves this by the flat residue/reflection form on the
original polynomial periods and retains all its constants.
Using that proved identity, PQO21 gives a further precise
obstruction on the positive ray. Since $D$ has distinct
diagonal entries, a matrix commuting with it is diagonal.
The matrix in PQO23 is $E_{\ell_{\rm br}}$ times a real matrix
times $E_{\ell_{\rm br}}^{-1}$;
if diagonal, its diagonal entries must therefore be real.
Its eigenvalues are exactly
\[
 \frac{\delta+i\gamma}{a^2},\quad
 \frac{\delta-i\gamma}{\bar a^2},\quad
 -\frac{\delta-i\gamma}{\bar a^2},\quad
 -\frac{\delta+i\gamma}{a^2},
 \tag{PQO24}
\]
as follows by evaluation at the four distinct roots of $h$.
Thus every positive-real-$u$ point in $Z$ satisfies the explicit
original-unit equation
\[
 \operatorname{Im}\left(\frac{\delta+i\gamma}{a^2}\right)=0.
 \tag{PQO25}
\]
If that imaginary part is nonzero for the original unit, the
positive ray contains no point of $Z$. When it is zero, the
four displayed eigenvalues are the two nonzero real numbers
$r,-r$, each twice, and the complete period-minor tests
PQO15--17 still determine $Z$. The equation is not assumed
to hold or fail. Together the exact minor tests, the actual
period continuation, the proved large-parameter result and
this real-structure obstruction give control of the original
quadric family with every amplitude value and phase present.

% END RX INLINE 57


% BEGIN RX INLINE 58; SHA256 0345ee8ee8c2580663de6fabd76a88745c0531a5d4f531e3e2666251ef6783e0
% Independent mathematical provider. Original periods, amplitude and branch retained.
\section{The original period quadric as a residue commutator}
\label{sec:actual-period-residue-commutator}

Retain $0<\delta<1/2$, $\gamma>2$, the ordered simple quartet
$\alpha_1=1/2+\delta+i\gamma$, $\alpha_2=1/2+\delta-i\gamma$,
$\alpha_3=1/2-\delta+i\gamma$, $\alpha_4=1/2-\delta-i\gamma$,
and $h(s)=\prod_a(s-\alpha_a)$. Put
\[
 w=s-\tfrac12,\quad r_a=\alpha_a-\tfrac12,\quad
 \beta=2(\gamma^2-\delta^2),\quad \eta=(\delta^2+\gamma^2)^2,
 \quad H(w)=w^4+\beta w^2+\eta.
 \tag{RPC1}
\]
Thus $r_4=-r_1$, $r_3=-r_2$, $r_2=\overline{r_1}$ and
$r_1^2-r_2^2=4i\delta\gamma$. No root or original coordinate is
discarded by this translation. Its coefficient isomorphism is
$J_c:\mathbb C[w]_{<4}\to E_h$, $J_cp=[p(s-1/2)]$; its inverse
is substitution $s=w+1/2$ in the unique original representative.
The primitive with its original constant is
\[
 \Phi_h(s)=A+\phi(w),\quad
 \phi(w)=\frac{w^5}{5}+\frac{\beta w^3}{3}+\eta w,\quad
 A=\Phi_h(1/2)=\frac1{160}+\frac\beta{24}+\frac\eta2.
 \tag{RPC2}
\]
The last equality is forced by $\Phi_h(0)=0$ and follows by evaluating
the displayed odd polynomial at $w=-1/2$.

Let $u\ne0$ carry the specified original branch of its logarithm.
Let $\Pi$ be the original period map on the rays and orientations
of PDE2, and let $F_{ja}=\zeta^{ja}-1$, $\zeta=e^{2\pi i/5}$,
$1\leq j,a\leq4$. Its inverse has entries
$(F^{-1})_{aj}=\zeta^{-ja}/5$: multiplying the two matrices and using
$\sum_{j=1}^4\zeta^{jb}=4$ for $5\mid b$ and $-1$ otherwise proves
this entry by entry. The literal centered coefficient-to-Fourier map is
\[
 P=F^{-1}\Pi J_c,\qquad P_0=e^{-A/u}P.
 \tag{RPC3}
\]
These are invertible by the complete original period determinant
PDE3--13. The contours for $P_0$ may be taken to be the rays based at
$w=0$ with the original five angles. Indeed translating the old rays
gives rays based at $-1/2$. Join corresponding finite endpoints and
close each contour at radius $R$. The integrand is entire. On the
bounded-length outer joining curves its exponent is
$-R^5/(5|u|)+O(R^4)$, uniformly along the joining curves, so the
outer integrals tend to zero. The common finite joining path occurs
with opposite coefficients in $-\ell_0+\ell_j$ and cancels exactly.
Cauchy's integral theorem therefore proves this contour identity,
including its signs and the factor $e^{A/u}$.

\subsection{A flat alternating pairing with all constants}

On $\mathbb C[w]/(H)$ define the bilinear, rather than sesquilinear,
pairing
\[
 \mathcal J(p,q)=\sum_{H(r)=0}\frac{p(r)q(-r)}{H'(r)}.
 \tag{RPC4}
\]
The roots are simple at the original packet. The equivalent polynomial
formula is the coefficient of $w^3$ in the remainder of $p(w)q(-w)$
modulo $H$. To prove it, use Lagrange interpolation: the coefficient
of $w^3$ in the idempotent at $r$ is $1/H'(r)$. This polynomial
formula extends the pairing to every complex value of $\beta,\eta$,
including repeated roots, without a division by a repeated-root value.
In the increasing frame $1,w,w^2,w^3$, multiplication by $w$ and
this pairing are the exact matrices
\[
 M=\begin{pmatrix}0&0&0&-\eta\\1&0&0&0\\0&1&0&-\beta\\0&0&1&0\end{pmatrix},
 \qquad
 J=\begin{pmatrix}0&0&0&-1\\0&0&1&0\\0&-1&0&\beta\\1&0&-\beta&0\end{pmatrix}.
 \tag{RPC5}
\]
For the last matrix, the coefficient functional is zero through degree
two, one in degree three, zero in degree four and six, and $-\beta$
in degree five. Inserting the sign $(-1)^j$ from the second argument
gives every displayed entry. Thus $J^{\mathsf T}=-J$, $\det J=1$,
and direct multiplication gives $M^{\mathsf T}J+JM=0$.

Allow the two original odd-phase coefficients to vary complexly,
solely to calculate the existing period map. The ray domination
$e^{-r^5/(5|u|)+O(r^3+r)}$ on compact parameter sets justifies all
coefficient derivatives and integrations by parts below. Ordinary
division of $w^{j+3}$ by $H$, followed by differential reduction,
gives
\[
 \partial_\eta P_0=\frac1uP_0M,\qquad
 \partial_\beta P_0=\frac1{3u}P_0(M^3-uN),\qquad
 N=\begin{pmatrix}0&0&1&0\\0&0&0&2\\0&0&0&0\\0&0&0&0\end{pmatrix}.
 \tag{RPC6}
\]
For clarity, the quotient derivatives for columns $j=0,1,2,3$ are
$0,0,1,2w$, respectively; hence this is the full correction matrix.
The two outer boundary terms vanish, and the two values at the common
origin cancel. No differential remainder was replaced by an arithmetic
remainder in RPC6.

There is also a direct nonvanishing proof on this entire two-parameter
family. The displayed matrices satisfy $\operatorname{Tr}M=0$,
$\operatorname{Tr}M^3=0$ and $\operatorname{Tr}N=0$. The first and
third follow from their diagonals; the second follows by multiplying
RPC5 three times, or from the vanishing third Newton sum of the even
monic polynomial $H$. Multilinearity of the determinant in its columns
then gives $\partial_\eta\det P_0=\partial_\beta\det P_0=0$ from
RPC6, without first assuming an inverse exists. At the origin the
literal Gamma ray integrals give the diagonal matrix $G_u$ defined in
RPC8, whose determinant is nonzero. Hence $P_0$ is invertible at every
point of $\mathbb C^2$, justifying the inverses throughout this family.

The matrix identity $M^{\mathsf T}J+JM=0$ also gives
$(M^3)^{\mathsf T}J+JM^3=0$. Direct multiplication of RPC5--6 gives
\[
 \partial_\beta J+\frac13(N^{\mathsf T}J+JN)=0.
 \tag{RPC7}
\]
In particular $\partial_\beta J$ has entries $(2,3)=1$, $(3,2)=-1$
when indices begin at zero, while $N^{\mathsf T}J+JN$ has the
opposite threefold entries. Differentiate the inverse matrices in
$P_0^{-\mathsf T}JP_0^{-1}$. RPC6--7 show that both its $\beta$ and
its $\eta$ derivatives vanish. This matrix is consequently constant
on the connected coefficient space $\mathbb C^2$.

Define the unchanged positive-Gamma ray constants
\[
 g_a=5^{a/5-1}e^{\pi ia/5}\Gamma(a/5),\qquad
 G=\operatorname{diag}(g_1,g_2,g_3,g_4),\quad
 G_u=\operatorname{diag}(u^{a/5}g_a)_{a=1}^4.
 \tag{RPC8}
\]
At $\beta=\eta=0$, literal substitution in each original ray integral
gives $P_0=G_u$. Every $g_a$ is nonzero, since its Gamma factor is
the integral of a strictly positive function. We have therefore proved
the complete flat-pairing identity
\[
 \begin{split}
 \Omega_0&=P_0^{-\mathsf T}JP_0^{-1}
       =G_u^{-\mathsf T}J\big|_{\beta=0}G_u^{-1},\\
 (\Omega_0)_{14}&=-\frac1{u g_1g_4},\qquad
 (\Omega_0)_{23}=\frac1{u g_2g_3},\qquad
 (\Omega_0)_{ba}=-(\Omega_0)_{ab},\\
 \Omega&=P^{-\mathsf T}JP^{-1}=e^{-2A/u}\Omega_0.
 \end{split}
 \tag{RPC9}
\]
All entries not listed are zero; the indices in RPC9 begin at one.
This formulation retains the Gamma constants directly and requires
no unevaluated normalization of a period or source mass. For
$D=\operatorname{diag}(\zeta,\zeta^2,\zeta^3,\zeta^4)$, the two
paired sums of weights are five. It follows entry by entry that
\[
 D^{\mathsf T}\Omega D=\Omega,
 \qquad D^{\mathsf T}\Omega_0D=\Omega_0.
 \tag{RPC10}
\]

\subsection{The full arithmetic unit and the exact quadratic form}

Retain the original $v_h=2\xi/h$ and its complete simple-root unit
values. Reflection and conjugation of this original function give
\[
 v_h(\alpha_1)=v_h(\alpha_4)=a,\qquad
 v_h(\alpha_2)=v_h(\alpha_3)=\overline a,\qquad a\ne0.
 \tag{RPC11}
\]
Indeed both numerator and denominator are invariant under $s\mapsto1-s$
and conjugate under $s\mapsto\overline s$, first away from the roots
and then by the removable values. The complete zero orders give
nonzero values at these simple roots. Write
\[
 c_3=\frac{a^{-2}-\overline a^{-2}}{4i\delta\gamma},\qquad
 c_1=\frac{a^{-2}+\overline a^{-2}}2-(\delta^2-\gamma^2)c_3,
 \quad \mathfrak r(w)=c_1w+c_3w^3.
 \tag{RPC12}
\]
Both coefficients are real, and $(c_1,c_3)\ne(0,0)$. Directly,
$c_1+c_3r_1^2=a^{-2}$ and $c_1+c_3r_2^2=\overline a^{-2}$;
the same equalities hold at the respective opposite roots. Thus
$\mathfrak r(M)=M\,v_h(M+1/2)^{-2}$ in the original finite algebra, with the
entire function evaluated by its complete original arithmetic unit.
RPC12 is an exact interpolation of this specific multiplier.

For $Q_0(t)=t_1t_4-t_2t_3$, let $J_{\rm prim}$ be the original
Lagrange coordinate isomorphism and set
\[
 B=F^{-1}\Pi UJ_{\rm prim},\qquad Q(x)=Q_0(B^{-1}x),\qquad
 \mathcal H=P\mathfrak r(M)P^{-1}.
 \tag{RPC13}
\]
If $p=P^{-1}x$, the primary components of $B^{-1}x$ are exactly
$p(r_1)/a,p(r_2)/\overline a,p(r_3)/\overline a,p(r_4)/a$.
Since $H'(r)=4r(r^2+\gamma^2-\delta^2)$, its values give
$r/H'(r)=1/(8i\delta\gamma)$ on roots one and four, and
$r/H'(r)=-1/(8i\delta\gamma)$ on roots two and three. Inserting
these four values into RPC4 proves the exact identity
\[
 \boxed{\quad Q(x)=4i\delta\gamma\,
       p^{\mathsf T}\mathfrak r(M)^{\mathsf T}Jp
       =4i\delta\gamma\,x^{\mathsf T}\mathcal H^{\mathsf T}\Omega x.\quad}
 \tag{RPC14}
\]
There is no conjugate transpose in this quadratic identity. Because
$\mathfrak r$ is odd, $\mathfrak r(M)^{\mathsf T}J+J\mathfrak r(M)=0$; hence the matrices in
RPC14 are symmetric. This also proves
$\mathcal H^{\mathsf T}\Omega+\Omega\mathcal H=0$.
The four eigenvalues of $\mathfrak r(M)$ are exactly
$r_1/a^2,\overline{r_1/a^2},-\overline{r_1/a^2},-r_1/a^2$.
They are nonzero, although they need not be distinct. Both $B$ and
$Q_0$ are invertible as linear and quadratic matrices, respectively,
so $Q$ always has rank four, including the possible eigenvalue collisions.

By RPC10, substituting $Dx$ into RPC14 transports $\mathcal H$ to
$D^{-1}\mathcal H D$. The resulting symmetric quadratic matrices
are equal if and only if the corresponding $\mathcal H$ matrices
are equal, because $\Omega$ is invertible. Consequently
\[
 \boxed{\quad Q(Dx)=Q(x)\quad\Longleftrightarrow\quad
 [\mathcal H,D]=0\quad\Longleftrightarrow\quad
 \mathcal H_{ab}=0\text{ for every }a\ne b.\quad}
 \tag{RPC15}
\]
This is an exact test in the actual original period matrix and unit.
The final equivalence uses the four distinct diagonal entries of $D$.
It makes no assumption that their off-diagonal tests are independent.

For completeness, a rank-four quadratic semi-invariant of cyclic
weight $j\ne0$ is impossible: its coefficient in row $j$ could pair
only with the missing weight zero, so that row is zero. Thus any
semi-invariant rank-four quadratic must have weight zero. Since five
is prime, the projective orbit of $Q$ under $D$ has size one or five,
and size one is exactly RPC15. Every failure of any displayed
off-diagonal test proves size five.

\subsection{An explicit original-parameter radius forcing orbit five}

Here we prove that RPC15 fails for all sufficiently large $|u|$,
with a completely specified finite radius for the given original
quartet and unit. The argument evaluates the original periods; it
does not choose or insert a replacement amplitude.

Set $b=\beta u^{-2/5}$, $c=\eta u^{-4/5}$, and let $T(b,c)$
be the Fourier period matrix for
$z^5/5+bz^3/3+cz$ at $u=1$, in columns $1,z,z^2,z^3$ and on
the five rays of angles $(\pi+2\pi j)/5$. Substitution
$w=u^{1/5}z$, with exactly the original branch, gives
\[
 P_0=T(b,c)S_u,\quad S_u=\operatorname{diag}(u^{a/5})_{a=1}^4,
 \quad S_uMS_u^{-1}=u^{1/5}M(b,c),\quad T(0,0)=G.
 \tag{RPC16}
\]
Every column includes its original differential factor. The matrix
$M(b,c)$ is RPC5 with $\beta,\eta$ replaced by $b,c$; in particular
$N_0:=M(0,0)$ has its three subdiagonal entries equal to one.

The following positive constants are explicit finite real integrals
or finite expressions in the original Gamma constants:
\[
 \begin{split}
 E_j&=\frac85\int_0^\infty r^j(r^3/3+r)
                    e^{-r^5/5+r^3/3+r}\,dr\quad(0\le j\le3),\\
 C_T&=2\left(\sum_{j=0}^3E_j^2\right)^{1/2},\qquad
 g_+=\max_a|g_a|,\quad g_-^{-1}=\max_a|g_a|^{-1},\\
 \kappa&=\max_{1\le a<4}|g_{a+1}/g_a|,\qquad
 K_0=GN_0G^{-1},\\
 C_K&=4C_Tg_-^{-1}+2g_+g_-^{-1}
                       +2g_+g_-^{-2}C_T,\\
 C_{K,3}&=C_K(3\kappa^2+3\kappa C_K+C_K^2),\qquad
 \epsilon_0=\min\{1,(2C_Tg_-^{-1})^{-1}\}.
 \end{split}
 \tag{RPC17}
\]
Here $g_-^{-2}$ means $(\max_a|g_a|^{-1})^2$. The negative fifth
degree exponent proves convergence of every integral. All norms in
the following estimates are the Euclidean operator norm, with the
Frobenius norm used explicitly as an upper bound where stated.

Write $\epsilon=|b|+|c|$. When $|b|,|c|\le1$, on each ray
\[
 |e^{bz^3/3+cz}-1|
 \le \epsilon(r^3/3+r)e^{r^3/3+r}.
\]
Each contour has two rays, and each entry of $F^{-1}$ has absolute
value $1/5$ with four summands. The entry in column $j$ of
$T-G$ is therefore at most $E_j\epsilon$. Summing the squares
over all four rows proves
\[
 \|T-G\|\le\|T-G\|_{\rm F}\le C_T\epsilon.
 \tag{RPC18}
\]
For $\epsilon\le\epsilon_0$, the convergent inverse series for
$I+G^{-1}(T-G)$ gives
$\|T^{-1}\|\le2g_-^{-1}$ and
$\|T^{-1}-G^{-1}\|\le2g_-^{-2}C_T\epsilon$.
Also $\|M(b,c)-N_0\|\le\epsilon$, $\|M(b,c)\|\le2$,
and $\|N_0\|=1$. Subtract the products
$K:=T M(b,c)T^{-1}$ and $K_0$ in three terms, changing the left
factor, then the middle factor, then the inverse factor. The preceding
inequalities give, with the complete RPC17 constants,
\[
 \|K-K_0\|\le C_K\epsilon,\quad
 \|K\|\le\kappa+C_K\epsilon,\quad
 \|K^3-K_0^3\|\le C_{K,3}\epsilon.
 \tag{RPC19}
\]
For the last assertion use the exact noncommutative identity
$K^3-K_0^3=(K-K_0)K^2+K_0(K-K_0)K+K_0^2(K-K_0)$.
Its norm is at most $C_K\epsilon$ times
$(\kappa+C_K)^2+\kappa(\kappa+C_K)+\kappa^2$, as claimed.

The original multiplication and full unit in RPC13 now give exactly
\[
 \mathcal H=c_1u^{1/5}K+c_3u^{3/5}K^3,
 \qquad (K_0^3)_{41}=g_4/g_1,\quad (K_0)_{21}=g_2/g_1.
 \tag{RPC20}
\]
The scalar $e^{A/u}$ cancels from this conjugation, while it remains
in the quadratic form through RPC9. Put $B_*=\beta+\eta>0$.
For $|u|\ge1$, $\epsilon\le B_*|u|^{-2/5}$.

If $c_3\ne0$, define the completely specified radius
\[
 \begin{split}
 T_*&=\max\left\{1,\frac{B_*}{\epsilon_0},
 \frac{2\bigl(|c_3|C_{K,3}B_*+|c_1|(\kappa+C_K)\bigr)}
 {|c_3g_4/g_1|}\right\},\qquad R_*=T_*^{5/2}.
 \end{split}
 \tag{RPC21}
\]
For every $u$ on its original branch with $|u|\ge R_*$,
RPC19--20 imply
\[
 \left|u^{-3/5}\mathcal H_{41}-c_3g_4/g_1\right|
 \le |u|^{-2/5}\bigl(|c_3|C_{K,3}B_*
                           +|c_1|(\kappa+C_K)\bigr)
 \le\tfrac12|c_3g_4/g_1|.
 \tag{RPC22}
\]
In particular $\mathcal H_{41}\ne0$. In the quadratic matrix
RPC14 this is the nonzero $(1,1)$ entry
$4i\delta\gamma\,\mathcal H_{41}\Omega_{41}$, whose cyclic
weight is two.

If $c_3=0$, then $c_1\ne0$. In this case define
\[
 T_* =\max\left\{1,\frac{B_*}{\epsilon_0},
                   \frac{2C_KB_*}{|g_2/g_1|}\right\},
 \qquad R_*=T_*^{5/2}.
 \tag{RPC23}
\]
For every $|u|\ge R_*$, the same estimates give
\[
 \left|u^{-1/5}\mathcal H_{21}-c_1g_2/g_1\right|
       \le |c_1|C_KB_*|u|^{-2/5}
       \le\tfrac12|c_1g_2/g_1|.
 \tag{RPC24}
\]
Thus $\mathcal H_{21}\ne0$. The $(1,3)$ entry of the quadratic
matrix is then $4i\delta\gamma\mathcal H_{21}\Omega_{23}$,
which has cyclic weight four; symmetry retains its $(3,1)$ partner.
Both cases prove the unconditional conclusion for the actual original
parameters and the explicit radius belonging to those parameters:
\[
 \boxed{\quad |u|\ge R_*\quad\Longrightarrow\quad
       \#\{[Q(D^jx)]:0\le j<5\}=5.\quad}
 \tag{RPC25}
\]
The bracket denotes the projective class of the quadratic polynomial.
RPC21 or RPC23 supplies its radius for every nonzero original unit;
neither case assumes a condition on an uncomputed period coefficient.

Finally, the original periods are holomorphic on the connected
logarithm cover of $\mathbb C^*$, and on each branch domain: locally the contours may
be held in common decay sectors and differentiated under the dominated
integral, and contour deformation identifies these local definitions.
Their determinant never vanishes. Hence all entries of $\mathcal H$
are holomorphic there. RPC22 or RPC24 proves that its respective
decisive entry is not identically zero on the logarithm cover and,
by the identity theorem on that connected cover, on any open branch domain. The zeros
of that one entry are isolated, by its convergent local power series.
The exact exceptional set
\[
 \mathcal E=\{u:\mathcal H_{ab}(u)=0\text{ for all }a\ne b\}
 \tag{RPC26}
\]
is a subset of those isolated zeros and satisfies $|u|<R_*$.
Thus it has no accumulation inside the branch domain; accumulation
at $u=0$ or a branch boundary has not been excluded. RPC15 is the
full calculation specifying this set. No assertion that it is empty
at arbitrary finite complex $u$ is made or used.

% END RX INLINE 58


% BEGIN RX INLINE 59; SHA256 95bbd2cabc395d3f19a1f4f0c012d99524f1e7f2712e96c9913972a1baa77b0d
% Complete finite estimate on the original m=1 packet and its actual kernel.
% Original source and quotient providers: IFB1--11 and GMB1--16.
\section{A finite bound for the original kernel volume from its actual dimension}
\label{sec:OKV}

\paragraph{The original packet, sources, and exact coordinate map.}
Throughout this calculation the original multiplicity is $m=1$ and
\[
 k=4l+1,\quad l\ge1,\quad q=(k+1)^2\ge36,\quad c=k/2,
 \qquad 0<\delta<1/2,\quad\gamma>2.
 \tag{OKV1}
\]
Retain every original root and its two coordinates:
\[
\begin{split}
 \chi(S)&=\prod_{a,b=0}^k
  [S-c-(2a-k)\delta-i(2b-k)\gamma],\\
 \omega_{ab}&=(2b-k)\gamma-i(2a-k)\delta,\\
 P(y)&=i^{-q}\chi(c+iy)=\prod_{a,b=0}^k(y-\omega_{ab}).
\end{split}
 \tag{OKV2}
\]
The identity follows factor by factor from
$iy-(2a-k)\delta-i(2b-k)\gamma=i(y-\omega_{ab})$;
it retains the full phase $i^q$. The algebra map
\[
 \Phi:\mathbb C[y]/(P)\longrightarrow E_\chi:=\mathbb C[S]/(\chi),
 \qquad [r(y)]\longmapsto[r((S-c)/i)]
 \tag{OKV3}
\]
is well defined because $P((S-c)/i)=i^{-q}\chi(S)$.
Its inverse is $[f(S)]\mapsto[f(c+iy)]$, as verified by substitution.
Both maps preserve degrees and are inverse on each polynomial ring
before taking the quotients. In particular division by $P$ is the
original remainder operation through this exact coordinate map.

Use only the two original source orders $s\in\{1,k\}$ and write
\[
 b_s=s/2,\qquad M_s=(2\pi)^{s/2},\qquad
 dm_{0,s}(y)=\frac{M_s2^{b_s-2}}{\pi\Gamma(b_s)}
     \left|\Gamma\left(\frac{b_s}{2}+\frac{iy}{2}\right)\right|^2dy.
 \tag{OKV4}
\]
The exact source calculation IFB2--5 supplies mass $M_s$, the transform
$\int e^{ty}dm_{0,s}=M_s(\cos t)^{-b_s}$ for $|t|<\pi/2$, and
the original monic orthogonal polynomials and norms
\[
\begin{gathered}
 p_0=1,\quad p_1=y,\qquad
 p_{d+1}=yp_d-d(b_s+d-1)p_{d-1},\\
 h_d^{(s)}=\int|p_d|^2dm_{0,s}=M_s d!(b_s)_d,\qquad
 \varphi_d(y)=p_d(y)/\sqrt{h_d^{(s)}},\qquad
 u_d(S)=\varphi_d((S-c)/i).
\end{gathered}
 \tag{OKV5}
\]
All these integrals are on the same real line, with the original
source mass. For example the stated orthogonality and norms also
follow directly from its transform and the exact generating function
\[
 \sum_{d\ge0}p_d(y)\frac{z^d}{d!}
       =(1+z^2)^{-b_s/2}e^{y\arctan z}.
\]
Indeed the source integral of the product at sufficiently small real
$z,w$ is $M_s(1-zw)^{-b_s}$. The cosine addition identity and the
displayed transform give this equality; the source exponential bound
justifies each finite coefficient derivative under the integral.
The coefficients of $z^dw^e$ then give zero for $d\ne e$ and
$M_s d!(b_s)_d$ for $d=e$. The initial coefficients are $1,y$,
and differentiating the generating function gives precisely the
recurrence in OKV5. Thus the source norms used below retain their
original constants and all finite-degree factors.

Let $C_s:\mathbb C^q\to E_\chi$ have columns $[u_0],\ldots,[u_{q-1}]$
in the original frame $(1,S,\ldots,S^{q-1})$. Its exact leading
coefficients give
\[
 \det C_s=i^{-q(q-1)/2}\prod_{d=0}^{q-1}(h_d^{(s)})^{-1/2}\ne0.
\]
Retain the original GMB vectors and matrices, including the last row:
\[
 v_j=C_s^{-1}[u_{q+j}],\qquad
 K_j=I_q+\sum_{a=0}^{j-1}v_av_a^*,\qquad 0\le j\le q+1.
 \tag{OKV6}
\]
The vector $v_q$ has original polynomial degree $2q$.

\paragraph{Exact remainder coefficients in all original roots.}
Index the full list in OKV2 by $\omega_1,\ldots,\omega_q$, with
each pair $(a,b)$ occurring exactly once. Define the elementary and
complete homogeneous polynomials by their finite sums
\[
 e_h(\omega)=\sum_{1\le a_1<\cdots<a_h\le q}
                      \omega_{a_1}\cdots\omega_{a_h},\qquad
 h_t(\omega)=\sum_{n_1+\cdots+n_q=t}\prod_{a=1}^q\omega_a^{n_a},
 \tag{OKV7}
\]
with $e_0=h_0=1$, $e_h=0$ outside $0\le h\le q$, and
nonnegative integer $n_a$ in the second sum. The notation $h_t(\omega)$
in this paragraph is a polynomial in the roots; the original squared
norm continues to be $h_d^{(s)}$ of OKV5. Multiplication of formal
power series gives the exact identity
\[
 \left(\sum_{h=0}^q(-1)^he_h(\omega)z^h\right)
 \left(\sum_{t\ge0}h_t(\omega)z^t\right)=1,
\]
because each geometric series $(1-\omega_a z)^{-1}$ multiplies its
own factor $1-\omega_a z$. Consequently, for $D=q+j$ with $0\le j\le q$,
the quotient and remainder of $y^D$ upon division by $P$ are exactly
\[
\begin{split}
 Q_D(y)&=\sum_{t=0}^jh_t(\omega)y^{j-t},\qquad
 \operatorname{rem}_P(y^D)=y^D-P(y)Q_D(y),\\
 [y^a]\operatorname{rem}_P(y^D)
   &=-\sum_{t=0}^j h_t(\omega)(-1)^{D-t-a}e_{D-t-a}(\omega),
       \qquad 0\le a<q.
\end{split}
 \tag{OKV8}
\]
To verify the quotient assertion, the power-series identity makes
the coefficients of $y^D,y^{D-1},\ldots,y^q$ in $P Q_D$ equal
to those of $y^D$. The remaining coefficients are exactly the last
line of OKV8. Their degree is less than $q$, so uniqueness of monic
polynomial division proves both assertions. For $0\le D<q$ the
remainder is $y^D$ itself.

The actual maximal root modulus is
\[
 R=k\sqrt{\delta^2+\gamma^2},\qquad
 D_*:=\max\{3,\sqrt{\delta^2+\gamma^2}\},\qquad T_*:=D_*q.
 \tag{OKV9}
\]
Every $|\omega_{ab}|\le R$, and equality holds at each corner,
so this expression for $R$ is exact. In particular $R\le T_*$.
Counting the summands of OKV7 proves
\[
 |e_h(\omega)|\le\binom qh R^h,\qquad
 |h_t(\omega)|\le\binom{q+t-1}{t}R^t.
\]
The second count follows by placing $q-1$ separators among $t+q-1$
positions to encode $(n_1,\ldots,n_q)$. For fixed $j$, the binomial
$\binom{q+t-1}{t}$ is increasing in $0\le t\le j$, because its
successive ratio is $(q+t)/(t+1)\ge1$. In each nonzero summand
of OKV8 the total root exponent is $(D-t-a)+t=D-a$.
As $t$ varies the indices $D-t-a$ are distinct, so the sum of
their $\binom q{D-t-a}$ is at most $2^q$. It follows that
\[
 \left|[y^a]\operatorname{rem}_P(y^{q+j})\right|
 \le2^q\binom{q+j-1}{j}R^{q+j-a}
 \le2^{3q-1}R^{q+j-a},\qquad 0\le a<q,\quad0\le j\le q.
 \tag{OKV10}
\]
For the last inequality use
$\binom{q+j-1}{j}\le2^{q+j-1}\le2^{2q-1}$, by the binomial
expansion of $(1+1)^{q+j-1}$. These estimates retain the exact
coefficient formula OKV8, including every signed cancellation.

\paragraph{Original Gamma norms of the low monomials.}
The same recurrence and norms in OKV5 give, without changing the measure,
\[
 y\varphi_a=\sqrt{(a+1)(b_s+a)}\,\varphi_{a+1}
          +\sqrt{a(b_s+a-1)}\,\varphi_{a-1},
 \quad \varphi_{-1}=0.
 \tag{OKV11}
\]
For a polynomial of degree at most $q-2$, its orthonormal coefficient
vector is therefore sent by multiplication by $y$ to the sum of
one upward and one downward weighted shift, with output degree at most
$q-1$. Each shift has operator norm at most
$\sqrt{q(b_s+q)}$: the squared norm of its coefficient vector is
bounded by the largest squared weight times the input squared norm.
Since $s\le k<q$ gives $b_s\le q/2$, the triangle inequality yields
\[
 \|yf\|_{L^2(m_{0,s})}
       \le2\sqrt{q(b_s+q)}\,\|f\|_{L^2(m_{0,s})}
       \le3q\,\|f\|_{L^2(m_{0,s})}
       \le T_*\|f\|_{L^2(m_{0,s})}.
\]
Here $2\sqrt{3/2}=\sqrt6<3$. Starting from
$\|1\|_{L^2(m_{0,s})}=\sqrt{M_s}$ and applying this inequality
to $1,y,\ldots,y^{a-1}$ proves
\[
 \boxed{\|y^a\|_{L^2(m_{0,s})}\le\sqrt{M_s}\,T_*^a,
                      \qquad0\le a<q.}
 \tag{OKV12}
\]
All intermediate polynomial degrees satisfy the stated operator domain.

For any polynomial $f=\sum_a f_a y^a$, write
$\mathfrak c_{T_*}(f)=\sum_a|f_a|T_*^a$. This is a finite
coefficient sum, not a change of source or quotient norm.
The exact original recurrence gives
\[
 \mathfrak c_{T_*}(p_{d+1})
 \le T_*\mathfrak c_{T_*}(p_d)
       +d(b_s+d-1)\mathfrak c_{T_*}(p_{d-1}).
\]
For $1\le d<2q$ its actual positive coefficient is bounded by
$d(b_s+d-1)\le2q(b_s+2q)\le5q^2\le T_*^2$.
The initial values are $1,T_*$. Induction now proves
\[
 \boxed{\mathfrak c_{T_*}(p_d)\le(2T_*)^d,
                            \qquad0\le d\le2q.}
 \tag{OKV13}
\]
Indeed substituting the two preceding bounds in the recurrence gives
$T_*(2T_*)^d+T_*^2(2T_*)^{d-1}
 =\tfrac34(2T_*)^{d+1}\le(2T_*)^{d+1}$.

Combining all $q$ coefficients of OKV10 with OKV12, by the triangle
inequality in the original $L^2$ space, gives
\[
 \|\operatorname{rem}_P(y^D)\|_{L^2(m_{0,s})}
           \le\sqrt{M_s}\,q2^{3q}T_*^D\quad(0\le D\le2q).
\]
For $D\ge q$, each summand is bounded by
$\sqrt{M_s}2^{3q-1}R^{D-a}T_*^a
 \le\sqrt{M_s}2^{3q-1}T_*^D$ and there are $q$ summands.
For $D<q$, the assertion follows from OKV12 and $q2^{3q}\ge1$.
Linearity of the exact remainder map, followed by OKV13, proves
\[
 \boxed{\|\operatorname{rem}_P(p_d)\|_{L^2(m_{0,s})}
         \le\sqrt{M_s}\,q2^{3q}(2T_*)^d,
                           \qquad0\le d\le2q.}
 \tag{OKV14}
\]

\paragraph{The original remainder vectors and the full matrix norm.}
The polynomial $\operatorname{rem}_P(p_d)$ has degree less than $q$.
Its expansion in the orthonormal basis
$\varphi_0,\ldots,\varphi_{q-1}$ has squared coefficient norm equal
to its squared original $L^2(m_{0,s})$ norm, by expanding the finite
inner product and using OKV5. The exact quotient map OKV3 and
definition OKV6 therefore give, for $d=q+j$, $0\le j\le q$,
\[
 \|v_j\|_2^2
   =\frac{\|\operatorname{rem}_P(p_{q+j})\|_{L^2(m_{0,s})}^2}
                 {M_s(q+j)!(b_s)_{q+j}}
   \le\frac{q^22^{6q}(2T_*)^{2(q+j)}}{(q+j)!(b_s)_{q+j}}.
 \tag{OKV15}
\]
This equality is the actual map of the original quotient coordinates
to the remainder polynomial. The full source factor $M_s$ occurs in
both the numerator bound and its original norm denominator before
the displayed cancellation. No estimate on the conditioning of
$C_s$ or a companion matrix is used.

Since $b_s\ge1/2$, for each integer $d\ge1$ the first factor of
$(b_s)_d$ is at least $1/2$ and every subsequent factor is at least
one. Thus $d!(b_s)_d\ge1/2$. Also $T_*\ge1$ and $q+j\le2q$.
Consequently
\[
 \|v_j\|_2^2\le2q^22^{6q}(2D_*q)^{4q},\qquad0\le j\le q,
\]
and the original complete matrix obeys
\[
 \boxed{1\le\|K_{q+1}\|_2\le
   1+2(q+1)q^22^{6q}(2D_*q)^{4q}.}
 \tag{OKV16}
\]
To check its norm bound, $\|v_jv_j^*\|_2=\|v_j\|_2^2$
follows directly from Cauchy--Schwarz, with equality on a nonzero
$v_j$; it is zero for a zero vector. Add all $q+1$ summands of
OKV6 and the norm of $I_q$. Positivity also proves its lower bound.

Here is an explicit bound of the requested scale, with constants
depending only on the original two root coordinates:
\[
 \boxed{\log\|K_{q+1}\|_2
       \le C(\delta,\gamma)\,q\log(q+1),\qquad
 C(\delta,\gamma)=12+4\log(2D_*).}
 \tag{OKV17}
\]
For a direct verification, put
$U=2(q+1)q^22^{6q}(2D_*q)^{4q}\ge1$ and use
$\log(1+U)\le\log(2U)$. This upper bound is exactly
\[
 \log4+\log(q+1)+2\log q+6q\log2
       +4q\log(2D_*)+4q\log q.
\]
Write $L=\log(q+1)$. Since $q\ge36$, one has
$\log4\le L$, $\log q\le L$, $6\log2\le2L$, and $L\ge1$.
The last inequality follows already from
$\log4=\int_1^4dt/t>1$, by bounding the integrand below by
$1/2$ on $[1,2]$ and $1/4$ on $[2,4]$.
The preceding display is therefore at most
$4L+6qL+4\log(2D_*)qL$, which is bounded by the right side
of OKV17. The completely finite and sharper bound OKV16 remains
available before this logarithmic estimate.

\paragraph{The four original kernel Grams with their full coefficient frame.}
Retain the original surjection $\Lambda:E_\chi\to B_{\rm obs}$,
the actual subspace $\mathcal K_\Lambda=\ker\Lambda$, and its original
inclusion $I:\mathcal K_\Lambda\hookrightarrow E_\chi$. Put
$r_K=\dim\mathcal K_\Lambda$; no separate formula or upper bound
for this actual dimension is needed here. Use its fixed original
coefficient frame and write
\[
 T_s=C_s^{-1}I,\qquad
 G_{q+j-1}^{(s)}=C_s^{-*}K_j^{-1}C_s^{-1},\qquad
 H_j^{(s)}=I^*G_{q+j-1}^{(s)}I=T_s^*K_j^{-1}T_s,
             \qquad0\le j\le q+1.
 \tag{OKV18}
\]
For completeness, the quotient metric in this identity follows from
the original orthonormal source-coordinate map
$A_j=[I_q,v_0,\ldots,v_{j-1}]$.
It obeys $A_jA_j^*=K_j$. The minimum norm solution of $A_jz=x$
is $z=A_j^*K_j^{-1}x$: it has the required image and is orthogonal
to $\ker A_j$. Every other solution differs by such a kernel vector,
so the Pythagorean identity gives minimum squared norm $x^*K_j^{-1}x$.
The actual remainder is $C_sx$. Substitution proves the original
quotient Gram and then its full restriction in OKV18.

Let $\lambda_s=\|K_{q+1}\|_2$. The original update chain gives
$I_q=K_0\preceq K_j\preceq K_{q+1}\preceq\lambda_s I_q$.
For positive definite matrices, inversion reverses the order:
if $A\preceq B$, conjugate by $A^{-1/2}$ to obtain a matrix
with eigenvalues at least one, invert those eigenvalues and conjugate
back. Thus the full original kernel forms satisfy
\[
 \lambda_s^{-1}H_0^{(s)}\preceq H_j^{(s)}\preceq H_0^{(s)},
 \qquad H_0^{(s)}=T_s^*T_s,
 \qquad H_{j+1}^{(s)}\preceq H_j^{(s)}.
 \tag{OKV19}
\]
For $r_K>0$, the matrix $T_s$ is injective, so $H_0^{(s)}$ is
positive definite. Conjugation in OKV19 by its inverse square root
gives eigenvalues in $[\lambda_s^{-1},1]$. The determinant identity
for that exact congruence is
\[
 \det\bigl[(H_0^{(s)})^{-1/2}H_j^{(s)}(H_0^{(s)})^{-1/2}\bigr]
       =\frac{\det H_j^{(s)}}{\det(T_s^*T_s)}.
\]
Therefore
\[
 -r_K\log\lambda_s
 \le\log\det H_j^{(s)}-\log\det(T_s^*T_s)\le0.
 \tag{OKV20}
\]
In particular the determinant and conditioning of the actual
$T_s^*T_s$ are retained in this identity; no new kernel basis is
substituted for its original frame.

Define the original four-endpoint kernel volume, with its literal
degrees $q-1,q,2q-1,2q$, by
\[
 F_K^{(s)}=
 \log\det H_0^{(s)}+\log\det H_1^{(s)}
       -\log\det H_q^{(s)}-\log\det H_{q+1}^{(s)}.
 \tag{OKV21}
\]
Monotonicity in OKV19 makes each of
$\log\det H_0^{(s)}-\log\det H_q^{(s)}$ and
$\log\det H_1^{(s)}-\log\det H_{q+1}^{(s)}$ nonnegative.
Each is at most $r_K\log\lambda_s$ by OKV20. Consequently the
actual kernel volume satisfies the complete finite estimates
\[
 \boxed{\begin{split}
 0\le F_K^{(s)}
 &\le2r_K\log\|K_{q+1}\|_2\\
 &\le2r_K\log\left[1+2(q+1)q^22^{6q}(2D_*q)^{4q}\right]\\
 &\le2r_K C(\delta,\gamma)q\log(q+1),
                    \qquad s\in\{1,k\}.
 \end{split}}
 \tag{OKV22}
\]
For $r_K=0$, every kernel determinant has its original empty value
one, and $F_K^{(s)}=0$. Hence OKV22 includes this case.
The four copies of $\log\det(T_s^*T_s)$ in OKV21 have coefficients
$+1,+1,-1,-1$ and cancel exactly only after their presence in
OKV20 has been recorded.

\paragraph{The original boundary and arithmetic transitions.}
At each original degree set
\[
 Q_j^{(s)}=
   \left(\Lambda (G_{q+j-1}^{(s)})^{-1}\Lambda^*\right)^{-1},
 \qquad
 S_j^{(s)}=(G_{q+j-1}^{(s)})^{-1}\Lambda^*Q_j^{(s)}.
 \tag{OKV23}
\]
Their zero dimensional cases have their empty meaning. Surjectivity
of $\Lambda$ and positivity give the displayed inverses. Direct
multiplication proves $\Lambda S_j^{(s)}=I$ and
$I^*G_{q+j-1}^{(s)}S_j^{(s)}=0$.
For any fixed original section $S_*$ of $\Lambda$, the map $[I,S_*]$
is bijective: its inverse records $\Lambda z$ and the unique kernel
coordinate of $z-S_*\Lambda z$. The minimum section differs from
$S_*$ by a kernel column. The two frames therefore differ by a
block triangular matrix with identity diagonal and determinant one.
Taking the determinant of the full metric in the orthogonal frame
$[I,S_j^{(s)}]$ gives
\[
 \det H_j^{(s)}\det Q_j^{(s)}
       =|\det[I,S_*]|^2\det G_{q+j-1}^{(s)}.
 \tag{OKV24}
\]
This proves the exact relation with all kernel and boundary directions
retained. The same frame factor occurs at all four endpoints.
Their signed sum is consequently
\[
 F_K^{(s)}+F_B^{(s)}=\mathcal B_k^{(s)},\qquad
 F_B^{(s)}=\log\det Q_0^{(s)}+\log\det Q_1^{(s)}
               -\log\det Q_q^{(s)}-\log\det Q_{q+1}^{(s)}.
 \tag{OKV25}
\]
The original matrices $G_{q+j-1}^{(s)}$ decrease with $j$, so their
inverse forms increase. Compression by $\Lambda$ and inversion then
show that $Q_j^{(s)}$ decreases as well. Thus $F_B^{(s)}\ge0$,
and the finite companion bound is
\[
 \max\{0,\mathcal B_k^{(s)}-2r_K C(\delta,\gamma)q\log(q+1)\}
       \le F_B^{(s)}\le\mathcal B_k^{(s)}.
 \tag{OKV26}
\]
Here the finer middle bound in OKV22 can be used in place of its
last line, with the same proof.

Finally the original source identifications are
$\mathcal B_k^{(1)}=\mathcal B_k^\sigma$ and
$\mathcal B_k^{(k)}=\mathcal B_k^0$. Their already proved full
source transition is $\mathcal H_k=\mathcal B_k^\sigma-\mathcal B_k^0$,
and the original arithmetic transition is
$\mathcal B_k^{\rm ar}=\mathcal B_k^\sigma+\delta_k^\sigma$,
with its original finite signed allowance for $\delta_k^\sigma$.
Substitution of the proved exact splitting OKV25 gives, with every
source and arithmetic term still present,
\[
 \boxed{\mathcal B_k^{\rm ar}
   =F_K^{(k)}+F_B^{(k)}+\mathcal H_k+\delta_k^\sigma
   =F_K^{(1)}+F_B^{(1)}+\delta_k^\sigma.}
 \tag{OKV27}
\]
Thus OKV22 controls the original Gamma kernel summands in this
same arithmetic identity. It keeps the separate source transition
and arithmetic scalar with their established signs and bounds.
The result is finite for the actual $r_K=\dim\ker\Lambda$ and
uses no additional asserted dimension estimate or replacement source.

% END RX INLINE 59


% BEGIN RX INLINE 60; SHA256 be7ca05ac200e0f9bd7748e7bae9d01dd355300c10e50ce16b6d1317a3393f2f
\section{The actual simple-quartet kernel and boundary at the original volume scale}
\label{sec:OMV}

This section combines the complete original period calculation,
the exact cyclic-orbit rank calculation, and the finite original
Gamma kernel estimate. The hypotheses are the actual simple-quartet
data already used in the programme, and the period range below has
an explicit bound proved from those data. No orbit condition or
kernel-rank premise is added to the statement.

\subsection{The original domain and the proved period range}
Fix the original simple quartet $m=1$, its coordinates
$0<\delta<1/2$, $\gamma>2$, and its actual nonzero amplitude value
$a=v_h(1/2+\delta+i\gamma)$, where $v_h=2\xi/h$ and every
zero order in this quartet is complete. Reflection and conjugation
give the other three unit values exactly as in PQO3 and RPC11.
For every integer $l\geq1$ retain
\[
 k=4l+1,\qquad q=(k+1)^2,\qquad c=k/2,
 \qquad s\in\{1,k\}.
 \tag{OMV1}
\]
The polynomial $\chi$, the source measures $m_{0,s}$, their masses
$(2\pi)^{s/2}$, the original observation $\Lambda_k:E_\chi\to B_k$,
its kernel, and all four endpoint degrees are unchanged.

Let $R_*$ be the finite positive real number defined by RPC17 and
RPC21 when the actual $c_3$ of RPC12 is nonzero, and by RPC17 and
RPC23 when it is zero. Those two explicitly evaluated alternatives
exhaust the original nonzero unit, because $c_1,c_3$ cannot both
vanish. In particular $R_*$ depends on the fixed quartet and its
actual unit and is independent of $k$. Its constants are finite
positive Gamma values and the convergent positive integrals RPC17.
The full period proof RPC18--25 establishes
\[
 |u|\geq R_*\quad\Longrightarrow\quad
 \#\{[Q_u(D^j x)]:0\leq j<5\}=5
 \tag{OMV2}
\]
on every original logarithm branch. To recall the actual nonvanishing
calculation, RPC22 bounds the error in
$u^{-3/5}\mathcal H_{41}$ by half the nonzero value
$|c_3g_4/g_1|$ in the first alternative. RPC24 bounds the error
in $u^{-1/5}\mathcal H_{21}$ by half $|c_1g_2/g_1|$ in the
second. Each is a nonzero off-diagonal entry of the exact
$\mathcal H=P(c_1M+c_3M^3)P^{-1}$; the flat residue identity
RPC14--15 makes it a nonzero noninvariant coefficient of the actual
quadric. The degree-five cyclic action and its nondegenerate
quadratic matrix then give OMV2. All powers of $u$ use its original
branch, and the absolute error estimates hold for every such branch.

The exact invariant-ideal proof OQX1--3 now applies at every point
of this proved domain. It gives, for all the original degrees OMV1,
\[
 \boxed{\quad \dim B_k=k^2-6k+17,\qquad
                 r_K:=\dim\ker\Lambda_k=8k-16.\quad}
 \tag{OMV3}
\]
Specifically, an invariant polynomial vanishing on the original
quadric is divisible by all five distinct original orbit factors;
its quotient by their invariant product is invariant. Thus the
actual observed rank is $d_{5,k,0}-d_{5,k-10,0}$, with negative
degrees zero. The full character calculation OQX3 evaluates this
difference to $k^2-6k+17$, including the two original low degrees
$5,9$. Subtraction from $q=(k+1)^2$ proves the second equality.
This is the dimension of the kernel of the literal period observation,
not an assigned ambient dimension or a presumed maximal rank.

\subsection{A uniform finite bound on the actual kernel volume}
Define, from the original root coordinates only,
\[
 D_*=\max\{3,\sqrt{\delta^2+\gamma^2}\},\quad
 C_*=12+4\log(2D_*),\qquad
 \mathfrak M_k=(16k-32)C_*q\log(q+1).
 \tag{OMV4}
\]
Retain the exact original source-coordinate vectors
$v_j=C_s^{-1}[u_{q+j}]$, $K_j=I+\sum_{a<j}v_av_a^*$,
and the full original kernel frame $T_s=C_s^{-1}I$ of OKV6,18.
The original kernel metrics at degrees $q+j-1$ are
$H_j^{(s)}=T_s^*K_j^{-1}T_s$. Their actual four-endpoint volume
satisfies
\[
 \boxed{\begin{split}
 0\leq F_K^{(s)}(u)
 &\leq(16k-32)\log\!\left[
       1+2(q+1)q^22^{6q}(2D_*q)^{4q}\right]\\
 &\leq\mathfrak M_k,
 \qquad s\in\{1,k\},\quad |u|\geq R_*.
 \end{split}}\tag{OMV5}
\]
Here the source index and the actual period dependence of the kernel
are retained. The proof is direct substitution of the now proved
$r_K=8k-16$ into OKV22. Its finite coefficient proof retains
every original root in the signed remainder formula, all Gamma
recurrence coefficients through degree $2q$, and the full source
mass before its exact norm-ratio cancellation. Its metric proof
retains $\det(T_s^*T_s)$ at each of the four degrees, and cancels
only the four copies with coefficients $+1,+1,-1,-1$.
In particular neither the conditioning nor the changing position
of the original period kernel has been suppressed. The bound is
uniform in the entire period domain OMV2 because its right side
depends only on $k,\delta,\gamma$, and the dimension in OMV3 is
the same at every point of that domain.

Write the four original total quotient baselines as
$\mathcal B_k^{(1)}=\mathcal B_k^\sigma$ and
$\mathcal B_k^{(k)}=\mathcal B_k^0$. The full determinant identity
OKV24--25, with the original fixed coefficient section, gives
\[
 F_B^{(s)}(u)=\mathcal B_k^{(s)}-F_K^{(s)}(u),\qquad
 \max\{0,\mathcal B_k^{(s)}-\mathfrak M_k\}
           \leq F_B^{(s)}(u)\leq\mathcal B_k^{(s)}.
 \tag{OMV6}
\]
Both factors in the exact splitting are nonnegative by monotonicity
of their original minimum metrics. The literal boundary term is
the actual period-vector expression
\[
 F_B^{(s)}(u)=\log\frac{
       \det\mathcal M_{2q}^{(s)}(u)\det\mathcal M_{2q+1}^{(s)}(u)}{
       \det\mathcal M_q^{(s)}(u)\det\mathcal M_{q+1}^{(s)}(u)},
 \qquad
 \mathcal M_D^{(s)}(u)=\sum_{d=0}^{D-1}
       \lambda_d^{(s)}(u)(\lambda_d^{(s)}(u))^*.
 \tag{OMV7}
\]
The vectors are the complete original primary-period coefficients
OPG14 and OCS21 on the actual image of $\Lambda_k$. Formula OPG17
proves that each consecutive determinant ratio is $1+\xi_j$;
the unchanged weights $1,2,\ldots,2,1$ give exactly OMV7.
The companion kernel vectors retain the minimum-section correction
OPR2 and the action defect OPG22/OQC22. No invariance of that kernel
under the original sum action is needed for OMV5--7.

\subsection{The two individual leading coefficients}
The complete original baseline proofs BSL16--19 and HBL establish,
for fixed original $\delta,\gamma$ and $m=1$, both source limits
\[
 \frac{\mathcal B_k^{(s)}}{q^2}\longrightarrow
 C_B=9-8\log2+F(2,\pi)>\frac{769}{17010},
 \qquad s=1\text{ and }s=k.
 \tag{OMV8}
\]
Here $F(2,\pi)$ is exactly the energy functional evaluated in those
complete proofs; its original endpoint convention and coefficients
are retained. No value of it is redefined by the present comparison.
Since $q=(k+1)^2$,
\[
 \frac{\mathfrak M_k}{q^2}
   =C_*\frac{(16k-32)\log((k+1)^2+1)}{(k+1)^2}
               \longrightarrow0.
\]
For an elementary verification, $((k+1)^2+1)<(k+2)^2$, and
the last fraction is at most $32\log(k+2)/k$. The inequality
$\log x=\int_1^x dt/t\leq2(\sqrt x-1)$ for $x\geq1$
proves that this upper bound tends to zero. The finite estimate
OMV5 and the exact identity OMV6 therefore prove
\[
 \boxed{\quad
 \frac{F_K^{(1)}(u)}{q^2},\ \frac{F_K^{(k)}(u)}{q^2}
       \longrightarrow0,\qquad
 \frac{F_B^{(1)}(u)}{q^2},\ \frac{F_B^{(k)}(u)}{q^2}
       \longrightarrow C_B.
 \quad}\tag{OMV9}
\]
These limits are uniform for all $|u|\geq R_*$ and all original
branches. Indeed the kernel error is bounded by
$\mathfrak M_k/q^2$ uniformly, while the boundary error is at
most
$|\mathcal B_k^{(s)}/q^2-C_B|+\mathfrak M_k/q^2$.
The baseline is independent of $u$. This proves the asserted
uniformity, rather than only a conclusion for each separately
fixed observation. The original combined coefficient has now been
resolved into its two individual contributions on the proved
simple-quartet period domain.

\subsection{Exact propagation to the arithmetic identity and the finer scale}
Keep the original signed transitions
$\mathcal H_k=\mathcal B_k^\sigma-\mathcal B_k^0$ and
$\delta_k^\sigma=\mathcal B_k^{\rm ar}-\mathcal B_k^\sigma$.
Their complete finite allowances, including the QGQ simple-
multiplicity refinement, remain those of the full current receiver.
The exact arithmetic receiving equality becomes
\[
 \begin{split}
 \mathcal B_k^{\rm ar}-F_B^{(1)}(u)
       &=F_K^{(1)}(u)+\delta_k^\sigma,\\
 \mathcal B_k^{\rm ar}-F_B^{(k)}(u)
       &=F_K^{(k)}(u)+\mathcal H_k+\delta_k^\sigma.
 \end{split}\tag{OMV10}
\]
This is substitution in OKV27, so every source and sign is the
original one. In particular the finite upper and lower allowances
are obtained by adding $[0,\mathfrak M_k]$ to the actual interval
for $\delta_k^\sigma$, or to the actual interval for
$\mathcal H_k+\delta_k^\sigma$, respectively. The original proved
relations $\delta_k^\sigma/(kq)\to0$ and
$\mathcal H_k/(kq)\to-\mathcal C_\Gamma/4$, together with
$k/q\to0$ and OMV5, prove that both differences in OMV10,
divided by $q^2$, tend to zero uniformly on the stated period domain.
This is an exact comparison with the full arithmetic baseline.
The literal arithmetic kernel and boundary metrics retain their
own source until their corresponding metric comparison is applied.

The finer original Gamma return is still carried by an exact mixed
identity:
\[
 \begin{split}
 F_B^{(1)}(u)-F_B^{(k)}(u)
       &=\mathcal H_k-
                   \bigl(F_K^{(1)}(u)-F_K^{(k)}(u)\bigr),\\
 -\mathfrak M_k\leq F_K^{(1)}(u)-F_K^{(k)}(u)
       &\leq\mathfrak M_k,\\
 \mathcal H_k-\mathfrak M_k
       \leq F_B^{(1)}(u)-F_B^{(k)}(u)
       &\leq\mathcal H_k+\mathfrak M_k.
 \end{split}\tag{OMV11}
\]
Subtract OMV6 for the two actual source orders to obtain the first
line. Each kernel term lies in $[0,\mathfrak M_k]$, so their
difference lies in the stated interval, proving the other two lines.
The error bound here is of order $kq\log(q+1)$, whereas the
already proved signed return $\mathcal H_k$ has order $kq$.
No sign or limiting coefficient for either finer mixed difference
is assigned by OMV11. Its explicit remaining term is the difference
of the two original kernel metrics, on the exactly calculated
$8k-16$ dimensional period kernel. The full source multiplier and
its action on this actual kernel remain the input to that next
calculation, rather than a freely selected source order or kernel.

% END RX INLINE 60

\paragraph{The precise successor to the preceding finite interval.}
The preceding OMV11 interval is sharpened by the complete proof
AMT42--49 immediately below. That calculation retains the full
original polynomial factors and proves that all sourcewise kernel
differences are $o(kq)$ on the stated period domain. AMT36--41
also evaluates the literal arithmetic factors left to their source
metric comparison in OMV10. Both results have already been
propagated to the actual BRI and MRI receiving sites above.

% BEGIN RX INLINE 61; SHA256 f1a74298f3159df08e5833e96fd4ad46a51264b94e71656dd3883a95a2950c2e
% Exact original arithmetic mixed transfer. Historical sources remain unchanged.
\section{The original arithmetic kernel and boundary receive the full source comparison}
\label{sec:AMT}

This calculation uses the proved full-polynomial comparison TW12/TWA3,
the original observation OPG1--5, and the complete original Gamma
kernel estimate OKV1--27. It gives the exact maps between their source,
quotient, kernel and boundary metrics, before taking any determinant.
The subsequent simple-quartet conclusion uses the evaluated original
period test RPC1--26 and its exact rank calculation OQX1--3. Every
source order below is one of the already present orders $1,k$.

\subsection{The complete original source and the same observation}
Retain the original full-order quartet, with $0<\delta<1/2$,
$\gamma>2$, and positive integer complete zero order $m$. Thus
\[
\begin{gathered}
 \mathcal R=\{\tfrac12+\delta+i\gamma,\tfrac12+\delta-i\gamma,
              \tfrac12-\delta+i\gamma,\tfrac12-\delta-i\gamma\},\quad
 h(z)=\prod_{\alpha\in\mathcal R}(z-\alpha)^m,\quad v_h=2\xi/h,\\
 k\ge3,\quad c=k/2,\quad e=1+k(m-1),\quad q=e(k+1)^2,\\
 \lambda_{a,b}=c+(2a-k)\delta+i(2b-k)\gamma,\quad
 \chi(S)=\prod_{a,b=0}^k(S-\lambda_{a,b})^e,\quad E=\mathbb C[S]/(\chi),\\
 w_h(y)=|v_h(1/2+iy)|^2/(2\pi),\quad m_{h,k}=w_h^{*k},\quad
 d\sigma(y)=|\Gamma(1/4+iy/2)|^2dy/(2\pi).
\end{gathered}\tag{AMT1}
\]
The original source inner products, conjugate-linear in their first
argument, are
\[
 \langle P,Q\rangle_{\rm ar}
  =\int_{\mathbb R}\overline{P(c+iy)}Q(c+iy)m_{h,k}(y)\,dy,
 \qquad
 \langle P,Q\rangle_\sigma
  =\int_{\mathbb R}\overline{P(c+iy)}Q(c+iy)\,d\sigma(y).
 \tag{AMT2}
\]
No measure mass or complex-coordinate phase is changed. In particular
$\int d\sigma=\sqrt{2\pi}$. Write $\mathcal P_N$ for polynomials
in the original $S$ of degree at most $N$. In its coefficient frame
$(1,S,\ldots,S^N)$ let the two Grams be $H_N^{\rm ar},H_N^\sigma$.
For $N\ge q-1$, the literal remainder map
$J_N:\mathcal P_N\to E$ is onto and has kernel
$\chi\mathcal P_{N-q}$; at $N=q-1$ this is the zero vector space.
The frame on $E$ is $(1,S,\ldots,S^{q-1})$.

Here is the original observation being retained. Put $n=4m$,
$E_h=\mathbb C[z]/(h)$, $U=M_{j_hv_h}$, and
$\iota[P]=[P(z_1+\cdots+z_k)]\in E_h^{\otimes k}$.
The map $j_h$ records every divided Taylor coefficient through order
$m-1$ at every original root. Its unit $j_hv_h$ is invertible because
the stated order is complete. Fix the original $u\ne0$ and its
logarithm branch. With $\Phi_h'=h$ and $\Phi_h(0)=0$, let
$\ell_j$ be the outward ray of angle
$(\arg u+\pi+2\pi j)/(n+1)$ and $\Gamma_j=-\ell_0+\ell_j$.
The period map and projector are exactly
\[
\begin{gathered}
 (\Pi[P])_j=\int_{\Gamma_j}e^{\Phi_h(z)/u}
                           \operatorname{rem}_hP(z)\,dz,\quad1\le j\le n,\\
 C:\Gamma_j\mapsto\Gamma_{j+1}-\Gamma_1\ (j<n),\quad
        \Gamma_n\mapsto-\Gamma_1,\qquad
 P_k=\frac1{n+1}\sum_{a=0}^{n}(C^{-a})^{\otimes k},\\
 L=P_k\Pi^{\otimes k}U^{\otimes k}\iota,\quad
 B=\operatorname{im}L,\quad\jmath:B\hookrightarrow(\mathbb C^n)^{\otimes k},
 \quad\jmath\Lambda=L,\quad\Lambda:E\twoheadrightarrow B,\quad K=\ker\Lambda.
\end{gathered}\tag{AMT3}
\]
These are the convergent original representative integrals; their
existence and full determinant are proved in PDE and OPG. Fix the
original coefficient inclusion $I:K\hookrightarrow E$ and section
$S_*:B\to E$ with $\Lambda S_*=I_B$. Let
$r=\dim K$ and $b=\dim B=q-r$. All matrices below use these same
frames and this same full unit and period observation. Empty Grams
have determinant one, and maps to or from a zero space have their
unique empty values.

\subsection{All constants of the actual full-polynomial inequality}
Retain exactly the constants of TW and TWA:
\[
\begin{gathered}
 \alpha=\pi/2,\quad p=8m-\tfrac92,\quad B_h=42+8m,\quad M=21+4m,\\
 \vartheta_h=\int_{-1}^1w_h(y)\,dy>0,\quad
 M_\alpha=\int_{\mathbb R}e^{\alpha y}w_h(y)\,dy\in(0,\infty),\\
 c_\Gamma=\sqrt2e^{-7/6},\quad C_\Gamma=\sqrt{2\sqrt5}e^{2/3},\\
 a_k=\frac{c_h\vartheta_h^{k-3}e^{-\alpha(k-3)}(k-2)^{-B_h}}
                   {C_\Gamma2^{B_h/2}},\qquad
 A_k=\frac{C_h^{\rm tilt}}{c_\Gamma}k^{p+1}M_\alpha^{k-1},\qquad
 D_N=[1+4(N+M)^2]^M.
\end{gathered}\tag{AMT4}
\]
Here $c_h$ is the proved original convolution lower-envelope constant
TW7 and $C_h^{\rm tilt}$ is the original compact/tail maximum BT12.
Their full constructions occur in the preserved arithmetic endpoint
and BT providers; they depend on this fixed $h$ and not on $k,N,u$.
Their defining inequalities, proved there with the original numerator,
are
\[
\begin{split}
 m_{h,k}(y)&\ge c_h\vartheta_h^{k-3}
    e^{-\alpha(|y|+k-3)}(k-2+|y|)^{-B_h},\\
 m_{h,k}(y)&\le k C_h^{\rm tilt}M_\alpha^{k-1}
              e^{-\alpha|y|}(1+|y|/k)^{-p},\\
 c_\Gamma e^{-\alpha|y|}(1+|y|)^{-1/2}
 &\le d\sigma/dy\le
 C_\Gamma e^{-\alpha|y|}(1+|y|)^{-1/2}.
\end{split}\tag{AMT5}
\]
For clarity the full polynomial step is given here. The first and
third bounds, $k-2+|y|\le(k-2)(1+|y|)$, and
$(1+|y|)^{B_h}\le2^{B_h/2}(1+y^2)^M$, give
$m_{h,k}\ge a_k(1+y^2)^{-M}d\sigma/dy$. The second and third,
using $1+|y|/k\ge(1+|y|)/k$ and $p>1/2$, give
$m_{h,k}\le A_kd\sigma/dy$.
The fixed Gamma orthonormal polynomials obey
\[
 ye_j=\sqrt{(j+1)(j+1/2)}e_{j+1}
                         +\sqrt{j(j-1/2)}e_{j-1}.
\]
The two shifts on degree at most $N$ have norms at most $N+1$ and
$N$, respectively. Thus $\|yP\|_\sigma\le2(N+1)\|P\|_\sigma$.
Iterating at the successively increased degrees and expanding
$(1+y^2)^M$ proves
$\int|P|^2(1+y^2)^M d\sigma\le D_N\int|P|^2d\sigma$.
Cauchy--Schwarz applied to $|P|(1+y^2)^{-M/2}$ and
$|P|(1+y^2)^{M/2}$ then gives the reciprocal lower bound. The
coordinate map $P(S)\mapsto P(c+iy)$ has inverse
$f(y)\mapsto f((S-c)/i)$, preserves the degree and retains every
coefficient through this explicitly invertible substitution. It
therefore proves on the complete original source
\[
 \boxed{\quad \frac{a_k}{D_N}H_N^\sigma
                  \preceq H_N^{\rm ar}\preceq A_kH_N^\sigma.\quad}
 \tag{AMT6}
\]
All coefficients can be complex in this proof. In particular AMT6
holds on every affine remainder fibre, not only on monic vectors.
The two Grams are positive: the Gamma density is positive and the
lower bound in AMT5 is positive everywhere. Their finite moments
exist by the exponential upper bound.

\subsection{Exact remainder, minimum-section, kernel and boundary maps}
For $a\in\{\sigma,\mathrm{ar}\}$ put
\[
\begin{gathered}
 G_N^a=(J_N(H_N^a)^{-1}J_N^*)^{-1},\qquad
 V_N^a=(H_N^a)^{-1}J_N^*G_N^a:E\longrightarrow\mathcal P_N,\\
 H_{K,N}^a=I^*G_N^aI,\qquad
 Q_N^a=(\Lambda(G_N^a)^{-1}\Lambda^*)^{-1},\qquad
 S_N^a=(G_N^a)^{-1}\Lambda^*Q_N^a:B\longrightarrow E.
\end{gathered}\tag{AMT7}
\]
Inverses are on the indicated nonzero spaces. For an onto matrix
$J$ and positive $H$, direct multiplication gives $JV=I$ for
$V=H^{-1}J^*(JH^{-1}J^*)^{-1}$. Moreover
$Z^*HV=Z^*J^*G=0$ for $JZ=0$. Therefore every lift of $z$ is
$Vz+Z$, with squared norm
$z^*Gz+Z^*HZ$, proving its unique minimum and its Gram. Applying
this proof first to $J_N$, then to $\Lambda$, proves every formula
AMT7 and $\Lambda S_N^a=I_B$, $I^*G_N^aS_N^a=0$.
It also proves that $V_N^aS_N^a$ is the minimum lift for the
composite $\Lambda J_N$.

The exact comparison map on $\mathcal P_N$ is the identity on
original coefficients; it commutes with $J_N$ and $\Lambda J_N$.
It induces the identity on $E,K,B$. Between the two spaces of
minimum lifts the exact bijection and its inverse are
\[
 \operatorname{im}V_N^\sigma\longrightarrow\operatorname{im}V_N^{\rm ar},
 \quad V_N^\sigma z\longmapsto V_N^{\rm ar}z,
 \qquad T_N=V_N^{\rm ar}J_N\big|_{\operatorname{im}V_N^\sigma},
 \quad T_N^{-1}=V_N^\sigma J_N\big|_{\operatorname{im}V_N^{\rm ar}}.
 \tag{AMT8}
\]
Their difference $V_N^{\rm ar}z-V_N^\sigma z$ lies in
$\chi\mathcal P_{N-q}$, with its unique coefficient obtained by
division by the unchanged monic $\chi$. This also gives the
minimum-lift bijection on the preimage of $K$ and the analogous
bijection $V_N^\sigma S_N^\sigma y\mapsto
V_N^{\rm ar}S_N^{\rm ar}y$ for boundary lifts. The latter
differences lie in $\ker(\Lambda J_N)$, exactly as multiplication
by $\Lambda J_N$ verifies.

Minimize both sides of AMT6 over $J_NP=z$. Taking infima preserves
each inequality because these are the same nonempty fibres. Restrict
the resulting inequality to $z=Ix$, and independently minimize it
over $\Lambda z=y$. This proves the three full-vector comparisons
\[
\boxed{\begin{gathered}
 (a_k/D_N)G_N^\sigma\preceq G_N^{\rm ar}\preceq A_kG_N^\sigma,\\
 (a_k/D_N)H_{K,N}^\sigma\preceq H_{K,N}^{\rm ar}
                                      \preceq A_kH_{K,N}^\sigma,\\
 (a_k/D_N)Q_N^\sigma\preceq Q_N^{\rm ar}\preceq A_kQ_N^\sigma.
\end{gathered}}\tag{AMT9}
\]
The source identity is an exact bounded bijection with the two
constants in AMT6. No invariance of $K$ under an additional
multiplier was used to obtain its induced identities.

The matrix $[I,S_*]$ is invertible: its inverse sends $z$ to the
unique kernel coordinate of $z-S_*\Lambda z$ and the boundary
coordinate $\Lambda z$. Each $S_N^a-S_*$ lies in $K$, so the
change from this frame to $[I,S_N^a]$ is triangular with identity
diagonal and determinant one. The latter frame is orthogonal in
$G_N^a$ by AMT7. Consequently
\[
 \boxed{\quad\det H_{K,N}^a\det Q_N^a
       =|\det[I,S_*]|^2\det G_N^a,\qquad a\in\{\sigma,\mathrm{ar}\}.\quad}
 \tag{AMT10}
\]
This includes $r=0$ and $b=0$ and retains the original frame factor.

\subsection{The exact nested deficits and their rank factors}
Define the actual finite source determinant deficit and its width
\[
 \Xi_{k,N}=(N+1)\log A_k-
                  \log\frac{\det H_N^{\rm ar}}{\det H_N^\sigma},
 \qquad L_N=\log(A_kD_N/a_k)\ge0.
 \tag{AMT11}
\]
The last sign follows by applying AMT6 to any nonzero vector.
Put
\[
\begin{split}
 d_{E,N}&=q\log A_k-\log\frac{\det G_N^{\rm ar}}{\det G_N^\sigma},\\
 d_{K,N}&=r\log A_k-\log\frac{\det H_{K,N}^{\rm ar}}{\det H_{K,N}^\sigma},\\
 d_{B,N}&=b\log A_k-\log\frac{\det Q_N^{\rm ar}}{\det Q_N^\sigma}.
\end{split}\tag{AMT12}
\]
AMT9, by congruence with the inverse square root of each comparison
Gram, places each eigenvalue of its relative Gram in
$[a_k/D_N,A_k]$. Multiplication of all eigenvalues gives
$0\le d_{K,N}\le rL_N$, $0\le d_{B,N}\le bL_N$ and
$0\le d_{E,N}\le qL_N$.

Here is the full source-to-quotient determinant calculation, including
all relation directions. Choose an $H_N^\sigma$-orthonormal
coefficient frame $R_N$ of $\ker J_N$ solely for this computation,
and put $W_N=[R_N,V_N^\sigma(G_N^\sigma)^{-1/2}]$.
AMT7 gives $W_N^*H_N^\sigma W_N=I_{N+1}$. In this frame the
positive contraction
\[
 \mathscr D_N=A_k^{-1}W_N^*H_N^{\rm ar}W_N
       =\begin{pmatrix}F_N&C_N'\\(C_N')^*&A_N'\end{pmatrix}
       \preceq I_{N+1}
\]
has Schur complement
\[
 T_N'=A_N'-(C_N')^*F_N^{-1}C_N'
       =A_k^{-1}(G_N^\sigma)^{-1/2}G_N^{\rm ar}(G_N^\sigma)^{-1/2}.
\]
Indeed minimization in its first block is exactly the original
remainder minimum at $z=(G_N^\sigma)^{-1/2}w$.
Both $F_N$ and $T_N'$ are positive contractions: positivity follows
from positive block elimination, and the upper bound for the Schur
complement follows by testing its minimum at zero first coordinate.
Since $|\det W_N|^2=(\det H_N^\sigma)^{-1}$, elimination gives
\[
 \Xi_{k,N}=d_{R,N}+d_{E,N},\quad d_{R,N}:=-\log\det F_N\ge0,
 \qquad d_{E,N}=-\log\det T_N'\ge0.
 \tag{AMT13}
\]
When $N=q-1$, $F_N$ is empty, $d_{R,N}=0$, and the same equalities
hold without a relation inverse. The determinant identities are
independent of the auxiliary orthonormal frame.
Dividing the two instances of AMT10 proves
$d_{K,N}+d_{B,N}=d_{E,N}$. Thus the complete nested statement is
\[
\boxed{\begin{gathered}
 \Xi_{k,N}=d_{R,N}+d_{K,N}+d_{B,N},\qquad d_{R,N},d_{K,N},d_{B,N}\ge0,\\
 0\le d_{K,N}\le\min\{rL_N,d_{E,N}\}\le\min\{rL_N,\Xi_{k,N}\},\\
 0\le d_{B,N}\le\min\{bL_N,d_{E,N}\}\le\min\{bL_N,\Xi_{k,N}\}.
\end{gathered}}\tag{AMT14}
\]
In particular each metric keeps its own rank factor, and the two
deficits retain their exact common sum. Any proved finite upper
bound for the original $\Xi_{k,N}$, including the degree-specific
HMR44/TWA15 bound, can be substituted in these displayed minima.
No new density or source determinant is introduced by this substitution.

\subsection{All four original endpoints and the correlated mixed return}
Let
\[
 (N_0,N_1,N_2,N_3)=(q-1,q,2q-1,2q),\qquad
 (\epsilon_0,\epsilon_1,\epsilon_2,\epsilon_3)=(1,1,-1,-1),
\]
and define, in their original signs,
\[
\begin{split}
 F_K^a&=\sum_{j=0}^3\epsilon_j\log\det H_{K,N_j}^a,\qquad
 F_B^a=\sum_{j=0}^3\epsilon_j\log\det Q_{N_j}^a,\\
 \mathcal B_k^a&=\sum_{j=0}^3\epsilon_j\log\det G_{N_j}^a,
 \quad\Delta_K^\sigma=F_K^{\rm ar}-F_K^\sigma,
 \quad\Delta_B^\sigma=F_B^{\rm ar}-F_B^\sigma,
 \quad\delta_k^\sigma=\mathcal B_k^{\rm ar}-\mathcal B_k^\sigma.
\end{split}\tag{AMT15}
\]
AMT10 cancels its four identical frame factors only after these
signs are applied. It gives
$F_K^a+F_B^a=\mathcal B_k^a$ and
$\Delta_K^\sigma+\Delta_B^\sigma=\delta_k^\sigma$ exactly.
The spaces $\mathcal P_N$ increase with the same source norm and
remainder map, so their quotient minima decrease. Restriction to
$K$ and minimization onto $B$ preserve that order. Pairing $N_0$
with $N_2$ and $N_1$ with $N_3$ consequently proves
$F_K^a,F_B^a\ge0$.

For $X=K,B$, put $r_K=r$, $r_B=b$, and retain the exact caps
\[
 c_{X,N}=\min\{r_XL_N,d_{E,N}\},\qquad
 C_X^-=c_{X,q-1}+c_{X,q},\qquad
 C_X^+=c_{X,2q-1}+c_{X,2q}.
 \tag{AMT16}
\]
The entirely source-evaluated caps
$\widehat c_{X,N}=\min\{r_XL_N,\Xi_{k,N}\}$ and the corresponding
$\widehat C_X^\pm$ are at least these values; all conclusions below
also hold with the hats. Substituting AMT12 in AMT15 records each
endpoint and its sign:
\[
\boxed{\begin{split}
 \Delta_K^\sigma&=-d_{K,q-1}-d_{K,q}+d_{K,2q-1}+d_{K,2q},\\
 \Delta_B^\sigma&=-d_{B,q-1}-d_{B,q}+d_{B,2q-1}+d_{B,2q},\\
 \delta_k^\sigma&=-d_{E,q-1}-d_{E,q}+d_{E,2q-1}+d_{E,2q},\\
 -C_X^-&\le\Delta_X^\sigma\le C_X^+,\qquad X=K,B.
\end{split}}\tag{AMT17}
\]
The four copies of $r_X\log A_k$ cancel here because their signs
sum to zero. Since the original total correction is retained, the
joint interval is stronger than two independent intervals:
\[
\boxed{\begin{split}
 \max\{-C_K^-,\delta_k^\sigma-C_B^+\}
   &\le\Delta_K^\sigma\le
          \min\{C_K^+,\delta_k^\sigma+C_B^-\},\\
 \max\{-C_B^-,\delta_k^\sigma-C_K^+\}
   &\le\Delta_B^\sigma\le
          \min\{C_B^+,\delta_k^\sigma+C_K^-\}.
\end{split}}\tag{AMT18}
\]
These follow by subtracting the interval for either component from
the exact common sum; they require no independence assertion.

Here are the explicit original four endpoint widths:
\[
\begin{split}
 E_k^+&=L_{q-1}+L_q
   =2\log(A_k/a_k)+M\log[1+4(q-1+M)^2]
                         +M\log[1+4(q+M)^2],\\
 E_k^-&=L_{2q-1}+L_{2q}
   =2\log(A_k/a_k)+M\log[1+4(2q-1+M)^2]
                         +M\log[1+4(2q+M)^2].
\end{split}\tag{AMT19}
\]
They are the unchanged TWA5 constants, both nonnegative. AMT17
therefore gives in particular
\[
 \boxed{-rE_k^+\le F_K^{\rm ar}-F_K^\sigma\le rE_k^-,\qquad
        -bE_k^+\le F_B^{\rm ar}-F_B^\sigma\le bE_k^-.}
 \tag{AMT20}
\]
The total correction retains TWA15 with all four individual source
deficits:
\[
\begin{split}
 \underline\delta_k^\sigma
   &=\max\{-qE_k^+,-\Xi_{k,q-1}-\Xi_{k,q}\},\\
 \overline\delta_k^\sigma
   &=\min\{qE_k^-,\Xi_{k,2q-1}+\Xi_{k,2q}\},\qquad
 \underline\delta_k^\sigma\le\delta_k^\sigma\le\overline\delta_k^\sigma.
\end{split}\tag{AMT21}
\]
This also follows immediately from AMT14 and AMT17. It can be
inserted into AMT18 when an interval, rather than the actual finite
total determinant, is wanted. Thus the lower and upper endpoints
are kept in their distinct places throughout the mixed transfer.

For later use, subtraction of the exact logarithms in AMT4 gives
\[
\begin{gathered}
 \mathfrak s_h=\alpha+\log(M_\alpha/\vartheta_h),\qquad
 \mathfrak c_h=-\log M_\alpha+3\log\vartheta_h-3\alpha
       +\log\frac{C_h^{\rm tilt}C_\Gamma2^{B_h/2}}{c_hc_\Gamma},\\
 \log(A_k/a_k)=k\mathfrak s_h+(p+1)\log k
                         +B_h\log(k-2)+\mathfrak c_h,\qquad
 E_k^\pm=O_h(k+\log q).
\end{gathered}\tag{AMT22}
\]
The last statement follows from the displayed exact expressions,
since $M,p,B_h$ and the two integral constants are fixed and
$\log D_N=O_h(1+\log(N+1))$. It is a consequence of the original
constants and contains no change of measure.

\subsection{The full original multiplier and its transported observation}
On the existing tensor class $k=4l+1$, retain
\[
 t_j=2j+\tfrac12,\quad\xi_j=c+t_j,\quad
 f_l(S)=\prod_{j=0}^{l-1}(S-c-t_j),\quad
 \beta_k=\frac{(2\pi)^{k/2}}{\sqrt2\Gamma(k/2)},\quad
 dm_{0,k}(y)=\beta_k|f_l(c+iy)|^2d\sigma(y).
 \tag{AMT23}
\]
This is the exact Gamma recurrence identity TWA26 with masses
$(2\pi)^{k/2}$ and $\sqrt{2\pi}$. For $l=0$ both $f_l$ and
$\beta_k$ are one. Write $\mathsf E_\chi:E\to\mathbb C^q$ for
the complete divided Taylor map
$[P]\mapsto(P^{(r)}(\lambda_{a,b})/r!)_{a,b,\,0\le r<e}$.
It is invertible: a vector in its kernel is divisible by the
coprime factors of $\chi$, and dimensions agree. The full
multiplication matrix $U_f$ has, in each primary block, entries
\[
 (U_f)_{r,a}=\begin{cases}
 f_l^{(r-a)}(\lambda)/(r-a)!,&0\le a\le r<e,\\
 0,&a>r.
 \end{cases}\qquad
 \det U_f=\prod_{a,b=0}^k\prod_{j=0}^{l-1}
       [(2a-k)\delta-t_j+i(2b-k)\gamma]^e\ne0.
 \tag{AMT24}
\]
Indeed $k$ is odd, so every $\lambda_{a,b}$ has nonzero imaginary
part, while all $\xi_j$ are real and distinct. The inverse in each
block is the finite Taylor inverse of its nonzero constant term.
The exact polynomial map and its inverse on its stated range are
\[
 \mathcal P_N\longrightarrow f_l\mathcal P_N\subseteq\mathcal P_{N+l},
 \quad P\longmapsto f_lP,
 \qquad Q\longmapsto Q/f_l.
 \tag{AMT25}
\]
It sends the original fibre $J_NP=z$ bijectively to the product
jet fibre $(\mathsf E_fQ,\mathsf E_\chi Q)=(0,U_f\mathsf E_\chi z)$.
Vanishing of the $l$ distinct $f$-values proves divisibility by
$f_l$, and the inverse of $U_f$ proves the converse fibre statement.
The original relation space maps to $f_l\chi\mathcal P_{N-q}$.
On the free $\chi$ jet coordinate $Y$ of this constrained product
fibre, the exact observed map, its kernel, and a section are
\[
\boxed{\begin{gathered}
 \widetilde\Lambda=\Lambda\mathsf E_\chi^{-1}U_f^{-1},\qquad
 \widetilde\Lambda U_f\mathsf E_\chi=\Lambda,\\
 \ker\widetilde\Lambda=U_f\mathsf E_\chi K,\qquad
 \widetilde I=U_f\mathsf E_\chi I,\qquad
 \widetilde S_*=U_f\mathsf E_\chi S_*,\quad
 \widetilde\Lambda\widetilde S_*=I_B.
\end{gathered}}\tag{AMT26}
\]
Both inclusions in the kernel equality follow by applying the
displayed invertible map and its inverse. This retains the full
original amplitude in $\Lambda$ and every derivative in $U_f$.
The resulting kernel is proved by transport; it has not been
assigned the coordinates of $\mathsf E_\chi K$.

For an explicit metric on this transported space, let $b_n(y)$ be
the fixed monic Gamma polynomials and put
$\pi_n(S)=i^n b_n((S-c)/i)$, whose leading coefficient is one
and squared norm is $\gamma_n=\sqrt{2\pi}(2n)!/4^n$.
The evaluation kernels are the full finite sums
\[
 \mathscr K_{\chi,L}=\sum_{n=0}^L
   \frac{\mathsf E_\chi\pi_n(\mathsf E_\chi\pi_n)^*}{\gamma_n},\quad
 \mathscr K_{f,L}=\sum_{n=0}^L
   \frac{\mathsf E_f\pi_n(\mathsf E_f\pi_n)^*}{\gamma_n},\quad
 \mathscr K_{\chi,f;L}=\sum_{n=0}^L
   \frac{\mathsf E_\chi\pi_n(\mathsf E_f\pi_n)^*}{\gamma_n}.
\]
The product evaluation is onto for $L=N+l\ge q+l-1$, because
a polynomial of degree below $q+l$ vanishing at all its divided
jets is divisible by the degree-$q+l$ product $f_l\chi$ and is
zero. Thus its positive kernel has positive Schur complement
\[
\begin{gathered}
 C_N=\mathscr K_{\chi,N+l}-\mathscr K_{\chi,f;N+l}
            \mathscr K_{f,N+l}^{-1}\mathscr K_{f,\chi;N+l}>0,\\
 U_N=(\mathsf E_\chi\pi_{N+j}/\sqrt{\gamma_{N+j}})_{j=1}^l,
 \quad V_N=\mathscr K_{\chi,f;N+l}\mathscr K_{f,N+l}^{-1/2},\quad
 C_N=\mathscr K_{\chi,N}+U_NU_N^*-V_NV_N^*.
\end{gathered}\tag{AMT27}
\]
Empty $f$ blocks have the identity convention. The inverse of the
joint positive kernel, evaluated on $(0,Y)$, gives the minimum
$Y^*C_N^{-1}Y$: block elimination of the $f$-block proves this
identity with every mixed entry present. Combining it with AMT25
and the coefficient $\beta_k$ proves the exact tensor-source metrics
\[
\boxed{\begin{gathered}
 G_N^0=\beta_k\mathsf E_\chi^*U_f^*C_N^{-1}U_f\mathsf E_\chi,\qquad
 H_{K,N}^0=\beta_k\widetilde I^*C_N^{-1}\widetilde I,\\
 Q_N^0=\beta_k(\widetilde\Lambda C_N\widetilde\Lambda^*)^{-1},\qquad
 S_N^0=\mathsf E_\chi^{-1}U_f^{-1}C_N\widetilde\Lambda^*
                   (\widetilde\Lambda C_N\widetilde\Lambda^*)^{-1}.
\end{gathered}}\tag{AMT28}
\]
Direct multiplication verifies the section and its orthogonality.
The formulas hold on the original $E,K,B$, with the transported
observation exactly as in AMT26. No block or mixed product was
dropped from $C_N$.

The same coefficient identity also supplies a quantitative comparison
directly with this original tensor source. Define the positive exact
finite factors
\[
 \Pi_l=\prod_{j=0}^{l-1}t_j^2,\qquad
 U_{l,N}^{\rm poly}=\prod_{j=0}^{l-1}[4(N+j+1)^2+t_j^2],\qquad
 \widetilde L_N=L_N+\log(U_{l,N}^{\rm poly}/\Pi_l).
 \tag{AMT29}
\]
Pointwise $|f_l(c+iy)|^2\ge\Pi_l$. Also
$\|(iy-t)P\|_\sigma^2=\|yP\|_\sigma^2+t^2\|P\|_\sigma^2$,
so the degree-specific multiplication estimate in AMT6, applied
successively, gives
$\|f_lP\|_\sigma^2\le U_{l,N}^{\rm poly}\|P\|_\sigma^2$.
Writing $H_N^0$ for the unchanged $m_{0,k}$ source Gram yields
\[
 \beta_k\Pi_lH_N^\sigma\preceq H_N^0
                  \preceq\beta_k U_{l,N}^{\rm poly}H_N^\sigma,
 \qquad
 \frac{a_k}{D_N\beta_k U_{l,N}^{\rm poly}}H_N^0
       \preceq H_N^{\rm ar}\preceq\frac{A_k}{\beta_k\Pi_l}H_N^0.
 \tag{AMT30}
\]
The second pair follows by using the first upper bound for its lower
inequality and the first lower bound for its upper inequality.
Minimizing and restricting with the same $J_N,I,\Lambda$ gives
AMT30 for all three metrics in AMT28. Its common upper scalar is
$A_k/(\beta_k\Pi_l)$; it is retained before its four signed copies
cancel. The same determinant proof as AMT12--17 therefore gives
\[
\begin{gathered}
 \widetilde E_k^+=E_k^+
   +\log(U_{l,q-1}^{\rm poly}U_{l,q}^{\rm poly}/\Pi_l^2),\qquad
 \widetilde E_k^-=E_k^-
   +\log(U_{l,2q-1}^{\rm poly}U_{l,2q}^{\rm poly}/\Pi_l^2),\\
 \boxed{-r\widetilde E_k^+\le F_K^{\rm ar}-F_K^0
                 \le r\widetilde E_k^-,\qquad
        -b\widetilde E_k^+\le F_B^{\rm ar}-F_B^0
                 \le b\widetilde E_k^-.}
\end{gathered}\tag{AMT31}
\]
For an explicit bound on the extra widths, $t_j\ge1/2$ gives
\[
 0\le\log(U_{l,N}^{\rm poly}/\Pi_l)
   =\sum_{j=0}^{l-1}\log\left(1+\frac{4(N+j+1)^2}{t_j^2}\right)
   \le l\log[1+16(N+l)^2].
 \tag{AMT32}
\]
Thus $\widetilde E_k^\pm=O_h(k\log(q+k))$ on the stated endpoints.
This is a proved polynomial-source estimate for the exact original
$f_l$, with no inference that its multiplication preserves $K$.

\subsection{The evaluated original simple-quartet arithmetic kernel}
Now take the original simple-quartet class $m=1$, $k=4l+1$,
$l\ge1$, $q=(k+1)^2$. The parameter $u$ is precisely the period
parameter in AMT3, not the real integration coordinate $y$.
RPC21 and RPC23 give an explicit finite radius $R_*$ for the
unchanged quartet and arithmetic unit. For reference its exact
constants and cases are recorded here. Set
\[
\begin{gathered}
 \beta=2(\gamma^2-\delta^2),\quad\eta=(\delta^2+\gamma^2)^2,
 \quad a=v_h(\tfrac12+\delta+i\gamma)\ne0,\\
 c_3=\frac{a^{-2}-\overline a^{-2}}{4i\delta\gamma},\qquad
 c_1=\frac{a^{-2}+\overline a^{-2}}2-(\delta^2-\gamma^2)c_3,\\
 g_j=5^{j/5-1}e^{\pi ij/5}\Gamma(j/5)\ (1\le j\le4),\quad
 g_+=\max_j|g_j|,\quad g_-^{-1}=\max_j|g_j|^{-1},\quad
 \kappa=\max_{1\le j<4}|g_{j+1}/g_j|,\\
 E_j^{\rm ray}=\frac85\int_0^\infty r^j(r^3/3+r)
                  e^{-r^5/5+r^3/3+r}\,dr\ (0\le j\le3),\quad
 C_T=2\left(\sum_{j=0}^3(E_j^{\rm ray})^2\right)^{1/2},\\
 C_K=4C_Tg_-^{-1}+2g_+g_-^{-1}+2g_+g_-^{-2}C_T,\quad
 C_{K,3}=C_K(3\kappa^2+3\kappa C_K+C_K^2),\\
 \epsilon_*=\min\{1,(2C_Tg_-^{-1})^{-1}\},\qquad B_*=\beta+\eta.
\end{gathered}\tag{AMT33}
\]
Here $g_-^{-2}=(g_-^{-1})^2$. All integrals converge by their
negative fifth-degree exponent, and $(c_1,c_3)\ne(0,0)$ by the
original unit values. The exact radius is
\[
 R_*=
 \begin{cases}
 \displaystyle\left[\max\left\{1,\frac{B_*}{\epsilon_*},
 \frac{2(|c_3|C_{K,3}B_*+|c_1|(\kappa+C_K))}{|c_3g_4/g_1|}\right\}\right]^{5/2},
                  &c_3\ne0,\\[6pt]
 \displaystyle\left[\max\left\{1,\frac{B_*}{\epsilon_*},
                     \frac{2C_KB_*}{|g_2/g_1|}\right\}\right]^{5/2},
                  &c_3=0.
 \end{cases}\tag{AMT34}
\]
The complete period calculation RPC16--25 proves, for every original
branch and every $|u|\ge R_*$, a nonzero original commutator entry:
its $(4,1)$ entry in the first case and its $(2,1)$ entry in the
second. RPC14--15 identifies those entries with the exact test for
the original rank-four quadric's cyclic invariance. Thus its actual
projective orbit has five elements throughout this domain. OQX1--3,
using the product of precisely these five original quadrics, proves
\[
 \boxed{r=8k-16,\qquad b=k^2-6k+17,
        \qquad m=1,\ k=4l+1,\ l\ge1,\ |u|\ge R_*.}
 \tag{AMT35}
\]
This is the evaluated original period outcome; no value of the
amplitude or of a period coefficient is chosen. The radius depends
on the fixed one-factor quartet and unit and is independent of $k$.

For the same original packet, OKV22 proves for both existing source
orders $s=1,k$, and for the actual $K$ of any rank,
\[
\begin{gathered}
 D_*=\max\{3,\sqrt{\delta^2+\gamma^2}\},\quad
 C(\delta,\gamma)=12+4\log(2D_*),\\
 Z_q=1+2(q+1)q^2 2^{6q}(2D_*q)^{4q},\qquad
 0\le F_K^{(s)}\le2r\log Z_q
                \le2r C(\delta,\gamma)q\log(q+1),
 \quad F_K^{(1)}=F_K^\sigma,\quad F_K^{(k)}=F_K^0.
\end{gathered}\tag{AMT36}
\]
Its proof includes the exact source coefficient map and the unchanged
kernel inclusion. Combining it with the already proved AMT17--20
gives the finite arithmetic conclusion
\[
\boxed{\begin{split}
 \max\{0,F_K^\sigma-C_K^-\}
       &\le F_K^{\rm ar}\le F_K^\sigma+C_K^+,\\
 0\le F_K^{\rm ar}
       &\le U_K:=2r\log Z_q+\widehat C_K^+\\
       &\le2r C(\delta,\gamma)q\log(q+1)+rE_k^-.
\end{split}}\tag{AMT37}
\]
The cap symbols $C_K^\pm,\widehat C_K^\pm$ in this formula mean
the endpoint sums AMT16; the constant $C_K$ without a superscript
in AMT33 belongs only to the explicit period radius. They have
different stated indices and are not identified.
The alternative finite tensor-source bound
$F_K^{\rm ar}\le2r\log Z_q+r\widetilde E_k^-$ also follows from
AMT31. Formula AMT37 retains the sharper fixed-source widths.

On AMT35 the exact ranks give
\[
 \boxed{\quad F_K^{\rm ar}-F_K^\sigma=O_h(k^2),\qquad
 F_K^{\rm ar}=O_h(kq\log(q+1)),\qquad
 \frac{F_K^{\rm ar}}{q^2}\longrightarrow0.\quad}
 \tag{AMT38}
\]
For the first estimate use $r=8k-16$ in AMT20 and AMT22,
with $\log q=2\log(k+1)$. The second follows from AMT37.
For the limit, divide its last line by $q^2$: its first term is
at most $16C(\delta,\gamma)k\log(q+1)/q\to0$, and its second
is $O_h(k^2/q^2)\to0$. Nonnegativity gives the limit by the two
bounds. These estimates are uniform over the entire original
domain $|u|\ge R_*$, since none of their right sides depends on $u$.
The tensor comparison likewise gives
$F_K^{\rm ar}-F_K^0=O_h(k^2\log q)=o(q^2)$ from AMT31--32.

\subsection{The original boundary receives the full baseline coefficient}
Before using a limit, the exact finite arithmetic identity and its
consequence are
\[
\boxed{\begin{gathered}
 F_B^{\rm ar}=\mathcal B_k^\sigma+\delta_k^\sigma-F_K^{\rm ar},\\
 \max\{0,\mathcal B_k^\sigma+\delta_k^\sigma-U_K\}
    \le F_B^{\rm ar}\le\mathcal B_k^\sigma+\delta_k^\sigma,\\
 \max\{0,\mathcal B_k^\sigma+\underline\delta_k^\sigma-U_K\}
    \le F_B^{\rm ar}\le\mathcal B_k^\sigma+\overline\delta_k^\sigma.
\end{gathered}}\tag{AMT39}
\]
The first identity follows from AMT10 and AMT15; AMT37 and
$F_K^{\rm ar},F_B^{\rm ar}\ge0$ prove the first interval, and
AMT21 proves the second. The correlated boundary return also has
the sharper exact interval AMT18, with every original signed
total correction present.

BSL13--19 supplies the complete original finite baseline and its
proved coefficient
\[
 \mathcal B_k^\sigma=V_q+\eta_q+\eta_{q+1}-\eta_1,\qquad
 \frac{\mathcal B_k^{\rm ar}}{q^2}\longrightarrow
 C_B^{\rm base}:=9-8\log2+\mathcal F(2,\pi)
                      >\frac{769}{17010}>0.
 \tag{AMT40}
\]
Here $V_q$, the three distinct original relation errors $\eta_r$,
and the energy $\mathcal F$ are exactly the BSL/HBL/GEL quantities,
whose complete definitions and proofs are in the retained closure.
The superscript on $C_B^{\rm base}$ only distinguishes that existing
baseline coefficient from the endpoint caps $C_B^\pm$ in AMT16.
Its finite BSL14 enclosure can be substituted directly for
$\mathcal B_k^\sigma$ in AMT39, keeping the signs of all three
$\eta$ terms; no new evaluation is needed for that substitution.
Subtract AMT38 from the exact original identity
$F_B^{\rm ar}=\mathcal B_k^{\rm ar}-F_K^{\rm ar}$. It proves
\[
 \boxed{\quad
 \frac{F_B^{\rm ar}}{q^2}\longrightarrow C_B^{\rm base}>0,
 \qquad
 F_K^{\rm ar}+F_B^{\rm ar}=\mathcal B_k^{\rm ar},
 \quad m=1,\ k=4l+1,\ l\to\infty,\ |u|\ge R_*.
 \quad}\tag{AMT41}
\]
The convergence is uniform over that original period domain by
AMT38 and because the full-source baseline is independent of $u$.
Thus the proved original kernel estimate now passes to the actual
arithmetic kernel, and the original arithmetic boundary carries the
entire positive $q^2$ coefficient. All earlier full-sum identities
remain exact, with their individual mixed contributions now receiving
AMT37--41 and their signed finite corrections AMT17--21.

\subsection{The same comparison evaluates the mixed Gamma return on the $kq$ scale}
The first inequality pair AMT30 directly compares the two existing
Gamma sources on the same original coefficient space. Apply its
proved minimum and restriction maps AMT7--9 to $X=E,K,B$, writing
$H_{E,N}^a=G_N^a$, $H_{B,N}^a=Q_N^a$, and
$r_E=q,r_K=r,r_B=b$. At every original endpoint it gives
\[
 \beta_k\Pi_l H_{X,N}^\sigma\preceq H_{X,N}^0
       \preceq\beta_k U_{l,N}^{\rm poly}H_{X,N}^\sigma.
 \tag{AMT42}
\]
Define the exact relative determinant and the positive degree width
\[
\begin{gathered}
 z_{X,N}=\log\frac{\det H_{X,N}^0}{\det H_{X,N}^\sigma}
                              -r_X\log(\beta_k\Pi_l),\qquad
 \omega_N^\Gamma=\log(U_{l,N}^{\rm poly}/\Pi_l),\\
 0\le z_{X,N}\le r_X\omega_N^\Gamma,\qquad
 z_{K,N}+z_{B,N}=z_{E,N},\\
 \Omega_k^{\rm low}=\omega_{q-1}^\Gamma+\omega_q^\Gamma,\qquad
 \Omega_k^{\rm high}=\omega_{2q-1}^\Gamma+\omega_{2q}^\Gamma.
\end{gathered}\tag{AMT43}
\]
Congruence and multiplication of the eigenvalues in AMT42 prove
the two inequalities. The determinant identity AMT10, applied to
the tensor metric AMT28 as well, proves the common sum. These
facts preserve the complete original kernel and boundary rather
than assigning separate arbitrary errors to their determinants.

Use the precise return orientation
\[
 \Delta_X^\Gamma=F_X^\sigma-F_X^0\quad(X=K,B),\qquad
 \mathcal H_k=\mathcal B_k^\sigma-\mathcal B_k^0=-\mathfrak T_k.
\]
The last equality is the original TWA30/BSL18 sign. Substituting
all four $z$ terms, and canceling $\log(\beta_k\Pi_l)$ only
after the signed sum, proves
\[
\boxed{\begin{split}
 \Delta_X^\Gamma
   &=-z_{X,q-1}-z_{X,q}+z_{X,2q-1}+z_{X,2q},\\
 -r_X\Omega_k^{\rm low}
   &\le\Delta_X^\Gamma\le r_X\Omega_k^{\rm high},\\
 \Delta_K^\Gamma+\Delta_B^\Gamma&=\mathcal H_k,\\
 \mathcal H_k-r\Omega_k^{\rm high}
   &\le\Delta_B^\Gamma\le\mathcal H_k+r\Omega_k^{\rm low}.
\end{split}}\tag{AMT44}
\]
AMT29 displays every factor of both $\Omega$ values explicitly;
AMT32 gives the finite bounds
\[
\begin{split}
 0\le\Omega_k^{\rm low}
 &\le l\log[1+16(q-1+l)^2]+l\log[1+16(q+l)^2],\\
 0\le\Omega_k^{\rm high}
 &\le l\log[1+16(2q-1+l)^2]+l\log[1+16(2q+l)^2].
\end{split}\tag{AMT45}
\]
In particular both are $O(k\log(q+1))$ on the original
$m=1$ class. Its actual rank $r=8k-16$ therefore proves directly
\[
 \boxed{\Delta_K^\Gamma=O_h(k^2\log(q+1))=o(kq),\qquad
 \frac{F_K^{\rm ar}-F_K^\sigma}{kq}\longrightarrow0,\qquad
 \frac{F_K^{\rm ar}-F_K^0}{kq}\longrightarrow0.}
 \tag{AMT46}
\]
For the first equality divide the AMT44 bounds by $kq$ and use
$q=(k+1)^2$, so $k\log(q+1)/q\to0$. The second limit uses
AMT38, since $k^2/(kq)=k/q\to0$. The last uses the exact
identity $F_K^{\rm ar}-F_K^0=\Delta_K^\sigma+\Delta_K^\Gamma$.
All bounds are uniform for $|u|\ge R_*$.

For completeness the entire finite arithmetic boundary return is
also evaluated by those same quantities:
\[
\boxed{\begin{split}
 F_B^{\rm ar}-F_B^0
   &=\mathcal H_k+\delta_k^\sigma
                         -(F_K^{\rm ar}-F_K^0),\\
 \mathcal H_k+\delta_k^\sigma-r\widetilde E_k^-
   &\le F_B^{\rm ar}-F_B^0
             \le\mathcal H_k+\delta_k^\sigma+r\widetilde E_k^+.
\end{split}}\tag{AMT47}
\]
This uses AMT31 with its exact endpoint constants
$\widetilde E_k^+=E_k^++\Omega_k^{\rm low}$ and
$\widetilde E_k^-=E_k^-+\Omega_k^{\rm high}$, rather than a
new source comparison. Both tensor and arithmetic source masses
and the full original resultant have already been retained in
AMT23--30 before these signed cancellations.

The proved GEL18--19 constant is
\[
 \mathcal C_\Gamma=2\int_{a_*^2}^{b_*^2}\log x\,
                       \rho_{2,\pi}(x)\,dx-4(2\log2-1)>0,
 \qquad \frac{\mathfrak T_k}{lq}\longrightarrow\mathcal C_\Gamma.
 \tag{AMT48}
\]
Here $a_*,b_*$ and $\rho_{2,\pi}$ are exactly the proved GEL8/EIQ
equilibrium endpoints and density; their complete definitions and
integral proof remain in that provider. The endpoint letters here
avoid the already fixed original period parameter $u$. GEL18a
gives $4\log(27/(8\pi))\le\mathcal C_\Gamma\le6$ and proves
the strict positivity. TWA15 also proves the actual source return
$\delta_k^\sigma/(kq)\to0$. As $l/k\to1/4$, AMT44, AMT46
and AMT47 now prove
\[
\boxed{\begin{gathered}
 \frac{F_B^\sigma-F_B^0}{kq}\longrightarrow-\frac{\mathcal C_\Gamma}{4},
 \qquad
 \frac{F_B^{\rm ar}-F_B^0}{kq}\longrightarrow-\frac{\mathcal C_\Gamma}{4},
 \qquad
 \frac{F_B^{\rm ar}-F_B^\sigma}{kq}\longrightarrow0,\\
 \frac{F_B^\sigma-F_B^0}{\mathcal H_k}\longrightarrow1,
 \qquad
 \frac{F_B^{\rm ar}-F_B^0}{\mathcal H_k}\longrightarrow1,
 \quad m=1,\ k=4l+1,\ l\to\infty,\ |u|\ge R_*.
\end{gathered}}\tag{AMT49}
\]
Indeed the first two limits subtract an $o(kq)$ kernel return
from $\mathcal H_k=-\mathfrak T_k$, whose quotient by $kq$
tends to $-\mathcal C_\Gamma/4$. The third follows from
$\Delta_B^\sigma=\delta_k^\sigma-\Delta_K^\sigma$.
The nonzero limit of $\mathcal H_k/(kq)$ makes its denominator
nonzero for all sufficiently large original $k$ and proves the
last two ratios by division. The uniformity in $u$ follows from
the uniform kernel bounds and the period-independent total returns.
Thus the full existing source comparison evaluates the mixed
Gamma return also at its $kq$ scale, with its finite signed
remainders AMT44 and AMT47 still present.

% END RX INLINE 61

\paragraph{The factorial denominator in the current absolute control.}
The complete proof that follows sharpens the earlier absolute
kernel estimates by retaining their exact original factorial
denominators. MRI4d--e and BRI6c propagate its general-multiplicity
finite bounds and its actual simple-quartet $kq$ bound above.
The preceding AMT signed intervals and sourcewise differences
remain valid, with every arithmetic and Gamma source term retained.

% BEGIN RX INLINE 62; SHA256 74fb2088f85b54899953cc035f19d79507aab3d7210e4e85266498a2ca1b6de7
% Complete original-source estimate; preceding editions are unchanged.
\section{Factorial control of the original common kernel at the $kq$ scale}
\label{sec:KAF}

The exact denominator in OKV15 contains the information needed to improve
its absolute volume estimate. This proof keeps that denominator and also
allows every original primary multiplicity. The original Gamma sources,
root coordinates, coefficient maps and observation are retained. The
simple-quartet conclusion uses the actual period and rank calculations
PQO/RPC and OQC/OQX, with their complete proofs supplied as dependencies.

\subsection{Original coordinates, all primary multiplicities, and source norms}
Fix the original full-order quartet with $0<\delta<1/2$, $\gamma>2$,
and a positive integer primary order $m$. On the original tensor class set
\[
 k=4l+1,\quad l\ge1,\quad e_k=1+k(m-1),\quad
 q=e_k(k+1)^2,\quad c=k/2,\quad s\in\{1,k\},\quad b_s=s/2.
 \tag{KAF1}
\]
The original polynomial and its exact real-coordinate representative are
\[
\begin{split}
 \chi(S)&=\prod_{a,b=0}^k
 [S-c-(2a-k)\delta-i(2b-k)\gamma]^{e_k},\\
 \omega_{ab}&=(2b-k)\gamma-i(2a-k)\delta,\\
 P(y)&=i^{-q}\chi(c+iy)
       =\prod_{a,b=0}^k(y-\omega_{ab})^{e_k}.
\end{split}\tag{KAF2}
\]
In the full list $\omega_1,\ldots,\omega_q$, each root $\omega_{ab}$
occurs $e_k$ times. The exact maximal modulus remains
$R=k\sqrt{\delta^2+\gamma^2}$. Substitution gives inverse algebra maps
\[
 \Phi:\mathbb C[y]/(P)\longrightarrow E=\mathbb C[S]/(\chi),
 \quad[r]\longmapsto[r((S-c)/i)],\qquad
 \Phi^{-1}[f]=[f(c+iy)].\tag{KAF3}
\]
Indeed $P((S-c)/i)=i^{-q}\chi(S)$, and the two substitutions are
inverse before quotienting. Thus every root, multiplicity and factor of
$i$ in the original remainder operation is present.

Retain the complete measures and their monic orthogonal polynomials:
\[
\begin{gathered}
 M_s=(2\pi)^{s/2},\qquad
 dm_{0,s}(y)=\frac{M_s2^{b_s-2}}{\pi\Gamma(b_s)}
       \left|\Gamma\left(\frac{b_s}{2}+\frac{iy}{2}\right)\right|^2dy,\\
 p_0^{(s)}=1,\quad p_1^{(s)}=y,\qquad
 p_{d+1}^{(s)}=yp_d^{(s)}-d(b_s+d-1)p_{d-1}^{(s)},\\
 h_d^{(s)}=\int|p_d^{(s)}|^2dm_{0,s}=M_s d!(b_s)_d,\qquad
 \varphi_d^{(s)}=p_d^{(s)}/\sqrt{h_d^{(s)}},\qquad
 u_d^{(s)}(S)=\varphi_d^{(s)}((S-c)/i).
\end{gathered}\tag{KAF4}
\]
These are precisely the existing orders $1,k$. Their complete source
transform and norm calculation are IFB2--5 and OKV4--5. In particular
$\int e^{ty}dm_{0,s}=M_s(\cos t)^{-b_s}$ for $|t|<\pi/2$.
The generating series
$(1+z^2)^{-b_s/2}\exp(y\arctan z)$ has coefficients
$p_d^{(s)}z^d/d!$. Integrating its product at two sufficiently small
real variables gives $M_s(1-zw)^{-b_s}$, proving the displayed norms
and orthogonality by coefficient comparison. The original exponential
bound justifies those finite coefficient differentiations under the
integral. None of these source statements depends on $m$.

In the original polynomial frame of $E$ let $C_s$ have columns
$[u_0^{(s)}],\ldots,[u_{q-1}^{(s)}]$. Its determinant is exactly
\[
 \det C_s=i^{-q(q-1)/2}\prod_{d=0}^{q-1}(h_d^{(s)})^{-1/2}\ne0,
 \quad v_j^{(s)}=C_s^{-1}[u_{q+j}^{(s)}],\quad
 K_j^{(s)}=I_q+\sum_{a=0}^{j-1}v_a^{(s)}(v_a^{(s)})^*,
 \quad0\le j\le q+1.\tag{KAF5}
\]
The last included vector has original degree $2q$.

\subsection{A degree-dependent constant in the unchanged remainder estimate}
Define the actual finite constants
\[
 A_{k,m}=\max\left\{2\sqrt{1+\frac{k}{2q}},\frac{R}{q}\right\},
 \qquad T=qA_{k,m},\qquad
 \mathcal A_{k,m}=6\log2+4+4\log A_{k,m}.
 \tag{KAF6}
\]
The letter $A_{k,m}$ in this section is a coefficient-estimate constant;
it is not the arithmetic upper source constant $A_k$ of TWA/AMT.
In particular $A_{k,m}\ge2$ and $R\le T$. For fixed original quartet
and fixed $m$, $q\ge(k+1)^2$ implies
\[
 A_{k,m}\longrightarrow2,\qquad
 \mathcal A_{k,m}\longrightarrow4+10\log2.
 \tag{KAF7}
\]
Indeed $k/q\to0$ and $R/q=\sqrt{\delta^2+\gamma^2}\,k/q\to0$.
The root-dependent term in the maximum is retained at every finite $k$.

Here is the complete remainder estimate needed below. For the full
root list define elementary and complete homogeneous polynomials
$e_h(\omega)$ and $h_t(\omega)$ by their finite sums, with
$e_0=h_0=1$ and $e_h=0$ outside $0\le h\le q$. The formal identity
\[
 \left(\sum_{h=0}^q(-1)^he_h(\omega)z^h\right)
 \left(\sum_{t\ge0}h_t(\omega)z^t\right)=1
 \tag{KAF8}
\]
follows by multiplying each of the $q$ factors $1-\omega_a z$
by its geometric series. Repeated roots do not affect this identity.
It proves, for $D=q+j$ and $0\le j\le q$, the exact signed formula
\[
 [y^a]\operatorname{rem}_P(y^{q+j})
 =-\sum_{t=0}^j h_t(\omega)(-1)^{q+j-t-a}e_{q+j-t-a}(\omega),
 \qquad0\le a<q.\tag{KAF9}
\]
In detail the quotient is $\sum_{t=0}^jh_t(\omega)y^{j-t}$;
KAF8 cancels all coefficients of degree at least $q$ in its product
with $P$, and the remaining coefficients are KAF9. This proves the
formula by uniqueness of monic division, including all multiplicities.

The defining sums give
$|e_h|\le\binom qh R^h$ and
$|h_t|\le\binom{q+t-1}{t}R^t$. In every nonzero summand of KAF9
the total root exponent is $q+j-a$. The binomial coefficients of
$h_t$ increase with $t$, and the distinct indices in the $e_h$ sum
contribute at most $2^q$. Hence
\[
 |[y^a]\operatorname{rem}_P(y^{q+j})|
 \le2^q\binom{q+j-1}{j}R^{q+j-a}
 \le2^{3q-1}R^{q+j-a}.\tag{KAF10}
\]
The last bound is the binomial expansion at one, since $j\le q$.

The exact orthonormal recurrence from KAF4 has upward and downward
weights $\sqrt{(a+1)(b_s+a)}$ and $\sqrt{a(b_s+a-1)}$.
On degree at most $q-2$ each shift has norm at most
$\sqrt{q(b_s+q)}$. Because $b_s\le k/2$, their sum has norm at
most $2\sqrt{q(k/2+q)}\le T$. Iterating from the constant, with all
intermediate degrees in that range, proves
\[
 \|y^a\|_{L^2(m_{0,s})}\le\sqrt{M_s}\,T^a,
 \qquad0\le a<q.\tag{KAF11}
\]
For the finite coefficient sum $c_T(f)=\sum_a|[y^a]f|T^a$ the
original recurrence gives
$c_T(p_{d+1}^{(s)})\le T c_T(p_d^{(s)})+
d(b_s+d-1)c_T(p_{d-1}^{(s)})$.
For $1\le d<2q$,
\[
 d(b_s+d-1)\le2q(k/2+2q)=4q^2+kq
       \le4q^2+2kq\le T^2.
\]
Starting with $c_T(p_0)=1,c_T(p_1)=T$, induction therefore proves
$c_T(p_d)\le(2T)^d$ for $0\le d\le2q$: the recurrence bound is
$(3/4)(2T)^{d+1}\le(2T)^{d+1}$.
Combining KAF10 with all $q$ monomial norms KAF11, then using the
linearity of the original remainder map, gives
\[
 \|\operatorname{rem}_P(p_d^{(s)})\|_{L^2(m_{0,s})}
       \le\sqrt{M_s}\,q2^{3q}(2T)^d,
       \qquad0\le d\le2q.\tag{KAF12}
\]
For a monomial of degree $D\ge q$, each coefficient term is at
most $\sqrt{M_s}2^{3q-1}R^{D-a}T^a\le
\sqrt{M_s}2^{3q-1}T^D$. For $D<q$ use KAF11 directly.
These two cases and the proved coefficient sum establish KAF12
without deleting a remainder coefficient or changing the source norm.

\subsection{Retaining the factorial denominator}
Expansion of the degree-below-$q$ remainder in the orthonormal
polynomials of KAF4, followed by the exact map KAF3, proves
\[
 \|v_j^{(s)}\|_2^2
 =\frac{\|\operatorname{rem}_P(p_{q+j}^{(s)})\|_{L^2(m_{0,s})}^2}
          {M_s(q+j)!(b_s)_{q+j}}
 \le\frac{q^22^{6q}(2T)^{2(q+j)}}{(q+j)!(b_s)_{q+j}}.
 \tag{KAF13}
\]
The factor $M_s$ occurs in the original numerator and denominator
before this cancellation. For every $d\ge0$,
\[
 d!(b_s)_d\ge d!(1/2)_d=\frac{(2d)!}{4^d},
 \qquad
 \|v_j^{(s)}\|_2^2\le q^22^{6q}\frac{(4T)^{2(q+j)}}{(2q+2j)!}.
 \tag{KAF14}
\]
The inequality is factor by factor because $b_s\ge1/2$.
The equality follows by separating even and odd factors of $(2d)!$.
This is precisely the denominator whose further retention sharpens
OKV16--17.

Define the positive exact quantities
\[
 \eta_{k,m}=\frac{1-1/(4q)}{A_{k,m}^2}<\frac14,\qquad
 Z_{k,m}=1+q^22^{6q}\frac{(4qA_{k,m})^{4q}}{(4q)!}
              \frac{1-\eta_{k,m}^{q+1}}{1-\eta_{k,m}}.
 \tag{KAF15}
\]
For $t_d=(4T)^{2d}/(2d)!$, $q+1\le d\le2q$ gives
$t_{d-1}/t_d=(2d)(2d-1)/(16T^2)\le\eta_{k,m}$.
Consequently $\sum_{d=q}^{2q}t_d\le
t_{2q}\sum_{a=0}^q\eta_{k,m}^a$. Adding the $q+1$ actual rank-one
updates in KAF5 and using KAF14 proves the finite operator bound
\[
 \boxed{1\le\|K_{q+1}^{(s)}\|_2\le Z_{k,m},\qquad s\in\{1,k\}.}
 \tag{KAF16}
\]
Here $\|vv^*\|_2=\|v\|_2^2$ follows by Cauchy--Schwarz with
equality on $v\ne0$; a zero vector contributes zero.

For an explicit linear bound, monotonicity of $\log x$ on each
interval $[j-1,j]$ gives
\[
 \log n!=\sum_{j=1}^n\log j
       \ge\int_1^n\log x\,dx=n\log n-n+1.
\]
Thus $(4qA_{k,m})^{4q}/(4q)!\le
\mathrm e^{-1}(\mathrm e A_{k,m})^{4q}$, where $\mathrm e$ is
the exponential constant and differs from the primary integer $e_k$.
The geometric sum in KAF15 is at most $4/3$. Therefore
$Z_{k,m}\le1+[4/(3\mathrm e)]X$, with
$X=q^2\exp(q\mathcal A_{k,m})$. The positive exponential series
gives $\mathrm e>1+1+1/2+1/6=8/3$, so $4/(3\mathrm e)<1/2$.
Since $A_{k,m}\ge2$ and $q\ge36$, $X>2$, and hence
$1+[4/(3\mathrm e)]X\le X$. We have proved
\[
 \boxed{\log\|K_{q+1}^{(s)}\|_2\le\log Z_{k,m}
       \le q\mathcal A_{k,m}+2\log q.}
 \tag{KAF17}
\]
The full factorial and finite geometric sum in KAF15 remain available
as the sharper finite bound.

\subsection{All four actual kernel and boundary Grams}
Let $\Lambda:E\twoheadrightarrow B$ be the original period and full-unit
observation, let $I:\mathcal K\hookrightarrow E$ be its actual kernel
inclusion, and put $r=\dim\mathcal K$. No rank formula is needed in
this subsection, so it applies to every original $m$. The original
quotient and restricted Grams are
\[
 G_{q+j-1}^{(s)}=C_s^{-*}(K_j^{(s)})^{-1}C_s^{-1},\qquad
 T_s=C_s^{-1}I,\qquad H_j^{(s)}=T_s^*(K_j^{(s)})^{-1}T_s.
 \tag{KAF18}
\]
Indeed the original orthonormal source-coordinate remainder matrix is
$[I_q,v_0,\ldots,v_{j-1}]$. Multiplication by its adjoint gives $K_j$;
its unique minimum lift of $x$ is its adjoint applied to $K_j^{-1}x$,
with squared norm $x^*K_j^{-1}x$. This proves both matrices KAF18
in the original coefficient frame. The positive update chain implies
\[
 Z_{k,m}^{-1}T_s^*T_s\preceq H_j^{(s)}\preceq T_s^*T_s,
 \qquad H_{j+1}^{(s)}\preceq H_j^{(s)}.\tag{KAF19}
\]
Inversion reverses positive order by congruence and inversion of
positive eigenvalues. Restricting this order with the same injective
$T_s$ proves KAF19. After congruence with $(T_s^*T_s)^{-1/2}$,
every eigenvalue lies between $Z_{k,m}^{-1}$ and one. Thus
$-r\log Z_{k,m}\le\log\det H_j^{(s)}-
\log\det(T_s^*T_s)\le0$.

The exact four original degrees are $q-1,q,2q-1,2q$. With their
unchanged signs define
\[
 F_K^{(s)}=\log\det H_0^{(s)}+\log\det H_1^{(s)}
              -\log\det H_q^{(s)}-\log\det H_{q+1}^{(s)}.
\]
Pairing degrees $q-1$ with $2q-1$, and $q$ with $2q$, proves
nonnegativity by KAF19. Each difference is at most $r\log Z_{k,m}$.
The four original frame factors cancel only after these four signs
have been recorded. It follows that
\[
 \boxed{0\le F_K^{(s)}\le2r\log Z_{k,m}
       \le2r\,[q\mathcal A_{k,m}+2\log q],\qquad s\in\{1,k\}.}
 \tag{KAF20}
\]
If $r=0$, every kernel determinant is one and the statement has its
same empty-space meaning.

At each original degree the boundary Gram is
$Q_j^{(s)}=(\Lambda(G_{q+j-1}^{(s)})^{-1}\Lambda^*)^{-1}$.
For any fixed section $S_*$ of $\Lambda$, its minimum section
$S_j^{(s)}=(G_{q+j-1}^{(s)})^{-1}\Lambda^*Q_j^{(s)}$ differs from
$S_*$ by a kernel column and is orthogonal to $I$ in that metric.
The change of frame has determinant one. Taking determinants in this
orthogonal frame proves, with its original coefficient factor,
\[
 \det H_j^{(s)}\det Q_j^{(s)}
       =|\det[I,S_*]|^2\det G_{q+j-1}^{(s)}.
\]
The same four signs now prove the exact splitting
$F_K^{(s)}+F_B^{(s)}=\mathcal B_k^{(s)}$. Boundary minimum norms
decrease with degree, so $F_B^{(s)}\ge0$. In particular
\[
 \boxed{\max\{0,\mathcal B_k^{(s)}-2r\log Z_{k,m}\}
       \le F_B^{(s)}\le\mathcal B_k^{(s)}.}\tag{KAF21}
\]
These are statements about the original kernel and boundary, with
all multiplicities and the same observation throughout.

\subsection{The actual simple quartet and the arithmetic source}
For $m=1$, the complete RPC calculation supplies an explicit finite
$R_*$ from the original unit and quartet. On every logarithm branch
and for $|u|\ge R_*$ its original period quadric has five distinct
cyclic orbit members. The full invariant restriction ideal calculation
OQX then gives $r=8k-16$. Substitution in the finite KAF20 gives
\[
 \boxed{0\le\frac{F_K^{(s)}}{kq}
 \le\left(16-\frac{32}{k}\right)
       \left(\mathcal A_{k,1}+\frac{2\log q}{q}\right),
 \qquad s\in\{1,k\},\quad |u|\ge R_*.}\tag{KAF22}
\]
The original arithmetic-source comparison AMT17--22 proves
$-rE_k^+\le F_K^{\rm ar}-F_K^{(1)}\le rE_k^-$, where the
unchanged full constants are
\[
 E_k^+=2\log(A_k/a_k)+\log D_{q-1}+\log D_q,\qquad
 E_k^-=2\log(A_k/a_k)+\log D_{2q-1}+\log D_{2q},
 \quad D_N=[1+4(N+21+4m)^2]^{21+4m}.
\]
Here $a_k,A_k$ are exactly the original arithmetic constants of AMT4,
whose full source inequality and minimum-section proof are included
there. They retain the actual numerator $2\xi/h$ and all source
masses; AMT22 proves $E_k^\pm=O_h(k+\log q)$.
Thus the actual arithmetic metric satisfies the finite bound
\[
 \boxed{0\le F_K^{\rm ar}
       \le(16k-32)\log Z_{k,1}+(8k-16)E_k^-.}\tag{KAF23}
\]
Combining KAF7, KAF22--23, and $q=(k+1)^2$ proves uniformly on
the entire stated original period domain, for
$a\in\{\sigma,0,\mathrm{ar}\}$,
\[
 \boxed{0\le\liminf_{k\to\infty}\frac{F_K^a}{kq}
       \le\limsup_{k\to\infty}\frac{F_K^a}{kq}
       \le16(4+10\log2).}\tag{KAF24}
\]
For precision, uniformity means the same finite upper bound holds for
every such $u$ and branch at each original $k$; therefore it also
bounds every sequence of those parameters. The arithmetic extra term
divided by $kq$ is at most $(8-16/k)E_k^-/q\to0$.
AMT46 further proves that the three ratios differ by quantities
tending uniformly to zero. Thus their limit points coincide for any
fixed allowed sequence of period parameters; KAF24 does not assert
that this bounded sequence already has a unique limit.

Finally the exact arithmetic identity gives
\[
 \max\{0,\mathcal B_k^{\rm ar}-(16k-32)\log Z_{k,1}
                              -(8k-16)E_k^-\}
       \le F_B^{\rm ar}\le\mathcal B_k^{\rm ar}.\tag{KAF25}
\]
Its full-source value remains
$\mathcal B_k^{\rm ar}=\mathcal B_k^\sigma+\delta_k^\sigma
=\mathcal B_k^0+\mathcal H_k+\delta_k^\sigma$.
All exact Gamma and arithmetic signed returns of AMT44--49 are
retained. The new calculation improves the absolute common kernel
bound to the $kq$ scale, while its exact four-endpoint determinant
and actual period dependence remain the objects for the next leading
term calculation.

% END RX INLINE 62



\clearpage
\part*{Complete current frame, residue and absolute kernel refinements}

\clearpage
% BEGIN CURRENT OCF
% Original coefficient reduction. All maps are complex-linear until a metric is supplied.
\section{An exact conductor frame and a fixed invariant-generator recursion}
\label{sec:OCF}

Retain the original simple-quartet data, $m=1$, $k=4l+1$,
$l\geq1$, the original root order $(++),(+-),(-+),(--)$, and
the complete original period and amplitude matrix of OQC1:
\[
                  \mathsf H=F^{-1}\Pi U\mathcal E_{\rm prim}.
\]
Thus $\det\mathsf H\ne0$ and
$Q(x)=Q_0(\mathsf H^{-1}x)$, where
$Q_0(t)=t_{++}t_{--}-t_{+-}t_{-+}$.
Use $g f(x)=f(D^{-1}x)$,
$D=\operatorname{diag}(\zeta,\zeta^2,\zeta^3,\zeta^4)$,
$\zeta=e^{2\pi i/5}$, and the original invariant algebra
$\mathscr I=\mathbb C[x_1,x_2,x_3,x_4]^{\langle g\rangle}$.
This calculation applies on the actual five-orbit period locus,
which includes every original logarithm branch with $|u|\geq R_*$
by PQO and RPC. On that locus the five prime polynomials
$Q_a=g^aQ$, $0\leq a<5$, are pairwise nonassociate. Keep their
actual products
\[
 \mathcal A=\prod_{a=1}^4Q_a,\qquad
 \mathcal H_Q=Q\mathcal A=\prod_{a=0}^4Q_a,
 \qquad \deg\mathcal A=8,\quad\deg\mathcal H_Q=10.
 \tag{OCF1}
\]
No coefficient in these polynomials is replaced. The exact orbit
decision and the period domain are those of the cited original
calculations. The construction below treats every possible nonzero
leading coefficient of their actual restriction.

\subsection{The original Segre coefficient isomorphism}

For $d\geq0$ let $\mathcal B_d$ be the space of biforms of
bidegree $(d,d)$ in $(z_0,z_1)$ and $(w_0,w_1)$. Its ordered basis is
\[
 e_{ab}^{(d)}=z_0^a z_1^{d-a}w_0^b w_1^{d-b},
 \qquad 0\leq a,b\leq d.
 \tag{OCF2}
\]
Pairs are ordered lexicographically, first by $a$, then by $b$.
The order is used only for the exact coefficient elimination.
Set $\mathcal B_d=0$ for $d<0$. Multiplication satisfies
$e_{ab}^{(d)}e_{ij}^{(e)}=e_{a+i,b+j}^{(d+e)}$.
The original substitution is
\[
 \begin{gathered}
 t=(z_0w_0,z_0w_1,z_1w_0,z_1w_1)^T,\qquad x=\mathsf Ht,\\
 \mathscr P_d=\mathbb C[x_1,x_2,x_3,x_4]_d,\qquad
 \psi_d:\mathscr P_d\longrightarrow\mathcal B_d,\quad
 f\longmapsto f(\mathsf Ht).
 \end{gathered}\tag{OCF3}
\]
The induced map
$R_d=\mathscr P_d/(Q\mathscr P_{d-2})\longrightarrow\mathcal B_d$
is an isomorphism. Here is a coefficient proof with the given
orientation. For each $(a,b)$, choose an integer
$h$ between $\max(0,a+b-d)$ and $\min(a,b)$. The monomial
$t_{++}^{h}t_{+-}^{a-h}t_{-+}^{b-h}t_{--}^{d-a-b+h}$ maps
to $e_{ab}^{(d)}$, proving surjectivity. Two monomials with
the same pair differ in their exponent vectors by an integer
multiple of $(1,-1,-1,1)$. After removing their common monomial,
their difference is divisible by $t_{++}t_{--}-t_{+-}t_{-+}$,
using $A^n-B^n=(A-B)\sum_{j=0}^{n-1}A^{n-1-j}B^j$.
A polynomial in the kernel has coefficient sum zero within every
pair, so it is a sum of such differences. The reverse inclusion
in the kernel follows by direct substitution. The invertible
substitution $x=\mathsf Ht$ then gives exactly the asserted kernel.
Thus these maps preserve multiplication as well as every grading.

Write the full conductor restriction as
\[
 A=\psi_8(\mathcal A)
     =\sum_{r=0}^8\sum_{s=0}^8 a_{rs}e_{rs}^{(8)}\ne0.
 \tag{OCF4}
\]
To verify its nonvanishing, $Q$ is a rank-four irreducible quadric,
hence prime in the polynomial ring, and none of $Q_1,\ldots,Q_4$
is its associate. Therefore $Q$ does not divide their product.
The kernel calculation OCF3 proves OCF4. The ring
$\mathcal B=\bigoplus_{d\geq0}\mathcal B_d$, with the indicated
diagonal grading, embeds in the four-variable polynomial ring;
it is an integral domain. Multiplication by $A$ is consequently
injective on every $\mathcal B_d$.

\subsection{Every leading-coefficient case gives an exact strip frame}

Let $(r_*,s_*)$ be the lexicographically largest pair for which
$a_{r_*s_*}\ne0$, and put $a_*=a_{r_*s_*}$. This is an
actual nonzero coefficient of OCF4. For $d\geq8$ define
\[
 \begin{split}
 P_d&=\{(r_*+i,s_*+j):0\leq i,j\leq d-8\},\\
 E_d^{\rm strip}&=\{0,\ldots,d\}^2\setminus P_d,\qquad
 W_d=\mathbb C^{E_d^{\rm strip}}.
 \end{split}\tag{OCF5}
\]
For $0\leq d<8$ put $P_d=\varnothing$,
$E_d^{\rm strip}=\{0,\ldots,d\}^2$ and use the same definition
of $W_d$; for $d<0$ put $W_d=0$. Let
$S_d:W_d\hookrightarrow\mathbb C^{(d+1)^2}$ insert zeros in
the positions $P_d$. All coefficient columns below use OCF2.

The following finite elimination defines
$N_d:\mathbb C^{(d+1)^2}\to W_d$ and, for $d\geq8$,
$V_d:\mathbb C^{(d+1)^2}\to\mathbb C^{(d-7)^2}$.
Starting with the complete input biform $f$, traverse all
$(i,j)\in\{0,\ldots,d-8\}^2$ in decreasing lexicographic order.
At the step $(i,j)$, let $c$ be the current coefficient at
$(r_*+i,s_*+j)$, record $b_{ij}=c/a_*$, and replace the current
biform by
\[
       f_{\rm current}-b_{ij}A e_{ij}^{(d-8)}.
 \tag{OCF6}
\]
At completion take the coefficients on $E_d^{\rm strip}$ as
$N_df$ and take $(b_{ij})$ as $V_df$. If $d<8$, $N_d$ is the
identity and $V_d$ is the zero map into the zero space.
The exact multiplication matrix $M_{A,d}$ gives
\[
 S_dN_d+M_{A,d}V_d=I,\qquad
 N_dS_d=I_{W_d},\qquad
 \ker N_d=A\mathcal B_{d-8}.
 \tag{OCF7}
\]
To prove these identities, the largest term of
$A e_{ij}^{(d-8)}$ is $a_*e_{r_*+i,s_*+j}^{(d)}$;
every other term is strictly smaller in the same lexicographic
order. Therefore each step zeros its designated position without
altering an earlier zero. Summing all subtractions gives the
first identity. Input supported on the strip produces zero
coefficients at every elimination step, giving the second.
Finally, the leading term of $Ab$, for a nonzero homogeneous
biform $b$, is the product of the two leading terms and belongs
to $P_d$. Indeed every other product is smaller in the additive
lexicographic order. Thus
$A\mathcal B_{d-8}\cap S_dW_d=0$. The first identity now proves
the kernel formula, the uniqueness of the decomposition, and all
three identities in OCF7. All operations are linear in the input;
their coefficients are explicit rational expressions in the
original $a_{rs}$ with powers of the displayed $a_*$ as denominators.
The polynomial $A$ itself remains the full polynomial OCF4.

In particular the quotient map induced by $N_d$ is the exact
isomorphism
\[
 \mathcal B_d/(A\mathcal B_{d-8})\xrightarrow{\ \sim\ }W_d,
 \qquad
 n_d:=\dim W_d=
 \begin{cases}(d+1)^2,&0\leq d<8,\\16d-48,&d\geq8.
 \end{cases}
 \tag{OCF8}
\]
The second formula is
$(d+1)^2-(d-7)^2=16d-48$. The possible choices
$(r_*,s_*)\in\{0,\ldots,8\}^2$ form an exhaustive finite
partition: the chart for a pair specifies that its coefficient
is nonzero and every larger coefficient is zero. Hence all actual
coefficient cases have been included. At the first original
degree $k=5$ this construction retains all $36$ coordinates;
at every original $k\geq9$ it retains $16k-48$ coordinates.

\subsection{Twenty-five fixed invariant monomials generate all degrees}

Define the following fixed finite sets of exponent vectors:
\[
 G_e=\{\nu\in\mathbb N^4:|\nu|=e,
                   \nu_1+2\nu_2+3\nu_3+4\nu_4\equiv0\pmod5\},
 \quad 1\leq e\leq5,\qquad G=\bigcup_{e=1}^5G_e.
 \tag{OCF9}
\]
Here $G_1$ is empty. The complete remaining lists, in the original
$x_1,x_2,x_3,x_4$ order, are
\[
\begin{array}{c|l}
e&G_e\\\hline
2&(0,1,1,0),(1,0,0,1)\\
3&(0,0,2,1),(0,1,0,2),(1,2,0,0),(2,0,1,0)\\
4&(0,0,1,3),(0,2,2,0),(0,3,0,1),(1,0,3,0),\\
 & (1,1,1,1),(2,0,0,2),(3,1,0,0)\\
5&(0,0,0,5),(0,0,5,0),(0,1,3,1),(0,2,1,2),\\
 & (0,5,0,0),(1,0,2,2),(1,1,0,3),(1,3,1,0),\\
 & (2,1,2,0),(2,2,0,1),(3,0,1,1),(5,0,0,0).
\end{array}\tag{OCF10}
\]
Thus $(|G_1|,\ldots,|G_5|)=(0,2,4,7,12)$ and $|G|=25$.
These lists can also be obtained directly by choosing
$0\leq\nu_1,\nu_2,\nu_3\leq e$ with their sum at most $e$,
putting $\nu_4=e-\nu_1-\nu_2-\nu_3$, and applying the single
displayed congruence, which verifies completeness of each list.

Every positive-degree invariant monomial is a product of members
$x^\nu$, $\nu\in G$. To prove this, list the weights of its
individual factors in any order. If its degree is at most five,
the entire list gives a member of $G$. If its degree exceeds five,
consider the six partial sums of the first five weights, including
the empty partial sum, in $\mathbb Z/5\mathbb Z$. Two coincide;
the nonempty block between them has length at most five and total
weight zero. Its factors give a divisor $x^\nu$ with $\nu\in G$.
The complementary monomial again has total weight zero and smaller
degree. Induction completes the factorization. Because the action
is diagonal, a polynomial is invariant exactly when all its
nonzero monomials have weight zero: comparison of each monomial
coefficient in $gf=f$ proves this statement. Thus
\[
                \mathscr I=\mathbb C[x^\nu:\nu\in G].
 \tag{OCF11}
\]
All intermediate degrees here are degrees in the original
coefficient algebra. This construction introduces no new source
distribution and leaves the original source orders $1$ and $k$ intact.

For each $\nu\in G_e$ the complete generator biform is
\[
 f_\nu=\psi_e(x^\nu)
  =\prod_{j=1}^4
   (H_{j1}z_0w_0+H_{j2}z_0w_1+
                  H_{j3}z_1w_0+H_{j4}z_1w_1)^{\nu_j}
  =\sum_{a,b=0}^e c_{\nu,ab}e_{ab}^{(e)}.
 \tag{OCF12}
\]
For an entirely explicit coefficient formula, sum over all
nonnegative $4$ by $4$ integer arrays $(n_{jt})$ satisfying
$\sum_t n_{jt}=\nu_j$,
$\sum_j(n_{j1}+n_{j2})=a$ and
$\sum_j(n_{j1}+n_{j3})=b$. Then
\[
 c_{\nu,ab}=
 \sum_{(n_{jt})}
     \prod_{j=1}^4\left(
        \frac{\nu_j!}{\prod_{t=1}^4n_{jt}!}
        \prod_{t=1}^4H_{jt}^{n_{jt}}\right).
 \tag{OCF13}
\]
The multinomial theorem proves the formula, including every
original period and amplitude entry. Every factor has degree at
most five and at most $36$ possible biform coefficients.

\subsection{Exact multiplication on the strip and the invariant recursion}

Let $\mathcal L=\mathcal B/(A\mathcal B)$, a graded algebra.
For $\nu\in G_e$ and $d\geq e$ define
\[
 m_{\nu,d}:W_{d-e}\longrightarrow W_d,
 \qquad m_{\nu,d}=N_d M_{f_\nu}S_{d-e}.
 \tag{OCF14}
\]
Here $M_{f_\nu}$ is the exact coefficient convolution:
\[
       (M_{f_\nu}S_{d-e}v)_{ab}
          =\sum_{i,j}c_{\nu,a-i,b-j}(S_{d-e}v)_{ij},
\]
with out-of-range indices giving zero.
Therefore every entry of OCF14 is supplied
by the finite formulas OCF6 and OCF13.
For a representative $h\in\mathcal B_{d-e}$, OCF7 writes
$h-S_{d-e}N_{d-e}h=A b$. Multiplication by $f_\nu$ preserves
this ideal. Applying $N_d$ yields the exact identity
\[
       N_d(f_\nu h)=m_{\nu,d}N_{d-e}h.
 \tag{OCF15}
\]
This proves that OCF14 is multiplication by the specified element
of $\mathcal L$, including its independence of representatives.

The original invariant image is
$J_d=\psi_d(\mathscr I_d)\subseteq\mathcal B_d$.
Set $\overline J_d=N_dJ_d\subseteq W_d$, and set
$\overline J_d=0$ for $d<0$. The exact recursion is
\[
 \overline J_0=\mathbb C\cdot1,\qquad
 \overline J_d=
   \sum_{\substack{\nu\in G_e\\2\leq e\leq5,\ e\leq d}}
                 m_{\nu,d}\overline J_{d-e}
                 \quad(d\geq1).
 \tag{OCF16}
\]
Indeed OCF11 factors every invariant monomial of degree $d>0$
as $x^\nu$ times an invariant monomial of degree $d-e$.
Applying $\psi$, then OCF15, gives the inclusion from left to
right. Conversely every summand is the image of a product of two
invariant polynomials, so it belongs to the left side. Linearity
proves equality for the complete invariant space. In particular
the recursion gives $\overline J_1=0$ and includes all small degrees.

The conductor ideal is contained in this invariant image with
its original trace coefficient:
\[
 A\mathcal B_{d-8}\subseteq J_d,\qquad
 [5\mathsf R(\mathcal A h)]_{(Q)}=[\mathcal A h]_{(Q)},
 \quad\mathsf R=\frac15\sum_{a=0}^4g^a.
 \tag{OCF17}
\]
To prove the identity, $g^a\mathcal A$ is the product of all
five orbit equations except $Q_a$; for $a\ne0$ it contains $Q$.
Thus all four corresponding terms vanish modulo $Q$, and the
$a=0$ term is the displayed product. Every biform in
$\mathcal B_{d-8}$ has a polynomial lift by OCF3, so the identity
proves the inclusion. It also shows exactly why the quotient
used in OCF16 retains the original invariant image.

\subsection{Exact ranks, a finite pivot atlas, and the original quotient map}

For $d\geq0$ put
\[
 d_0(d)=\frac{\binom{d+3}{3}+4\epsilon_d}{5},\qquad
 \epsilon_d=\mathbf1_{d\equiv0\ (5)}-\mathbf1_{d\equiv1\ (5)},
 \qquad d_0(d)=0\quad(d<0).
 \tag{OCF18}
\]
This is the number of degree-$d$ weight-zero monomials: the
identity character has generating function $(1-z)^{-4}$, each
of the four nonidentity characters has generating function
$\prod_{j=1}^4(1-z\zeta^j)^{-1}=(1-z)/(1-z^5)$, and averaging
their coefficients gives OCF18.
Every invariant polynomial divisible by $Q$ is divisible by all
five $Q_a$, hence by $\mathcal H_Q$, since they are pairwise
nonassociate primes. Its quotient by this product is invariant:
apply $g$ and cancel the nonzero invariant product in the integral
domain. Conversely each $\mathcal H_Q$ times an invariant
polynomial is invariant and vanishes modulo $Q$. Consequently
$\ker(\mathscr I_d\to J_d)=\mathcal H_Q\mathscr I_{d-10}$,
and injection of multiplication by the nonzero product gives
$\dim J_d=d_0(d)-d_0(d-10)$.
Injection of multiplication by $A$, OCF7 and OCF17 then prove
\[
 j_d:=\dim\overline J_d
     =d_0(d)-d_0(d-10)-\max(d-7,0)^2.
 \tag{OCF19}
\]
In particular $j_0=1$, $j_1=0$, $j_5=12$, and
$j_d=8d-32$ for $d\geq8$. For the last formula at $d\geq10$,
$\epsilon_d=\epsilon_{d-10}$ and direct expansion gives
$\binom{d+3}{3}-\binom{d-7}{3}=5(d^2-6d+17)$;
subtracting $(d-7)^2$ gives $8d-32$. At $d=8,9$ the ranks
before this subtraction are $33,44$, respectively, which gives
$32,40$ and verifies both remaining cases.

Here is a complete finite construction of all the columns used
by OCF16. Start with $Z_0=(1)$ and the invariant polynomial
certificate $F_{0,1}=1$. Suppose all earlier matrices $Z_b$ have
independent columns spanning $\overline J_b$, and their columns
are $N_b\psi_b(F_{b,h})$ for invariant polynomials $F_{b,h}$.
Form the candidate matrix
\[
 C_d=\bigl[\,m_{\nu,d}Z_{d-e}\,\bigr]_{
                    \nu\in G_e,\ 2\leq e\leq5,\ e\leq d}.
 \tag{OCF20}
\]
Each candidate column has the explicit invariant certificate
$x^\nu F_{d-e,h}$. OCF16 proves that $C_d$ has column image
$\overline J_d$ and rank $j_d$. Select $j_d$ columns and $j_d$
rows giving a nonzero determinant. Let $Z_d$ be those columns,
let $I_d$ be those row positions, and let $T_d$ be the complementary
row positions in $E_d^{\rm strip}$. Retain the corresponding
invariant certificates. This completes the induction. When
$j_d=0$, use the empty matrix with zero columns, empty $I_d$,
and $T_d=E_d^{\rm strip}$.

Every selected determinant and every matrix entry is an explicit
expression obtained from the original entries of $\mathsf H$ and
the full $a_{rs}$. To cover all cases, order the finite choices
of column and row subsets and select the first nonzero determinant;
the corresponding chart asserts its nonvanishing and the vanishing
of all preceding candidates. A nonzero determinant exists because
the rank was proved in OCF19. At rank zero no determinant is needed.
Together with the exhaustive leading-coefficient charts after
OCF8, this gives a finite atlas for each specified $k$. It requires
no generic coefficient premise and supplies no numerical zero test
for an unevaluated analytic coefficient: each chart is specified
by its exact minors.

Define $\rho_d:W_d\to\mathbb C^{T_d}$ and the full original
coefficient map $\pi_d$ by
\[
 \rho_d=\operatorname{row}_{T_d}
            -(Z_d)_{T_d}(Z_d)_{I_d}^{-1}
                                  \operatorname{row}_{I_d},
 \qquad \pi_d=\rho_dN_d.
 \tag{OCF21}
\]
At rank zero the second term is the zero map. If $E_{T_d}$
inserts a column into its $T_d$ positions in $W_d$, then
\[
 \begin{split}
 w&=Z_d(Z_d)_{I_d}^{-1}w_{I_d}+E_{T_d}\rho_dw,\\
 \rho_dE_{T_d}&=I,\qquad\ker\rho_d=\operatorname{im}Z_d
                                  =\overline J_d,\\
 \ker\pi_d&=J_d,\qquad
       \pi_d S_d E_{T_d}=I_{\mathbb C^{T_d}}.
 \end{split}\tag{OCF22}
\]
For the first line, the subtraction on its right leaves zero
$I_d$ rows and exactly $\rho_dw$ on the complementary rows.
This proves the first three claims. If $\pi_df=0$, then
$N_df\in\overline J_d$; choose $j\in J_d$ with $N_dj=N_df$.
Now $f-j\in\ker N_d=A\mathcal B_{d-8}\subseteq J_d$ by
OCF17, so $f\in J_d$. The reverse inclusion follows immediately
from $N_dJ_d=\overline J_d$. The last equality uses $N_dS_d=I$.
This proves every map and kernel in OCF22.

At the original degrees, these are the exact dimensions
\[
 \begin{array}{c|c|c|c}
 k&n_k&j_k&|T_k|=n_k-j_k\\\hline
 5&36&12&24\\
 k\geq9,\ k=4l+1&16k-48&8k-32&8k-16.
 \end{array}\tag{OCF23}
\]
The original map $\Theta_k:R_k\xrightarrow{\sim}E_k^\vee$
in OQC7 sends the coefficient basis OCF2 to the original dual
primary basis in the same order. Indeed its evaluation on
$e^{wY_k}[1]$ is
$\exp(((2b-k)\gamma-i(2a-k)\delta)w)$. The corresponding
eigenvalue of $A_k$ is
$\beta_{ab}=k/2+(2a-k)\delta+i(2b-k)\gamma$; the eigenvalue
of $Y_k=(A_k-kI/2)/i$ is
$\omega_{ab}=(\beta_{ab}-k/2)/i=(2b-k)\gamma-i(2a-k)\delta$.
These eigenvalues are distinct because
$\delta>0$ and $\gamma>0$. Their exponentials are independent
by the Vandermonde determinant, so the evaluation identifies the
covectors exactly. By the original observation identity OQC8,
$\Theta_kJ_k=\operatorname{im}\Lambda_k^\vee$. Therefore OCF22
supplies the specified isomorphism
\[
 E_k^\vee/\operatorname{im}\Lambda_k^\vee
       \xrightarrow{\ \pi_k\ }\mathbb C^{T_k}
       \cong (\ker\Lambda_k)^\vee.
 \tag{OCF24}
\]
The displayed $\pi_k$ uses exactly the original primary coefficient
coordinates under $\Theta_k$. Every dualization in this statement
is complex-linear. No conjugation is inserted in this coefficient
map; the original metric maps can subsequently be composed with it.

Finally, this construction is finite with explicitly bounded
compressed matrices. OCF20 has $n_d$ rows and
$\sum_{e=2}^5 |G_e|j_{d-e}$ columns, with negative subscripts zero.
Both dimensions grow at most linearly in $d$: for example
$n_d\leq16d+1$ and the column number is at most $25(16d+1)$.
For $d\geq13$ the exact column number is $200d-1632$, obtained
from $\sum_e|G_e|=25$, $\sum_e e|G_e|=104$, and
$j_{d-e}=8(d-e)-32$. Only the preceding five graded bases are
needed to form the next matrix. Each column is obtained from
the at-most-$36$-coefficient formula OCF12, one exact convolution,
and OCF6. A full original biform coefficient column, if used as
temporary storage, has $(d+1)^2$ entries; every invariant relation
matrix in OCF20 has the stated linear row and column counts.
The invariant certificates can be retained as their products and
parent-column references, so their expanded high-degree monomials
need never be enumerated. The proof of their completeness is
OCF11--16, and the exact original quotient and its right inverse
are OCF21--24.

\subsection{The same recursion with fourteen proved generators}

The complete fixed set OCF9 can be reduced to the following
fourteen specific members while preserving every map and image
just proved:
\[
\begin{array}{c|l}
e&\mathcal G_e\\\hline
2&x_1x_4,\ x_2x_3\\
3&x_1^2x_3,\ x_1x_2^2,\ x_2x_4^2,\ x_3^2x_4\\
4&x_1^3x_2,\ x_1x_3^3,\ x_2^3x_4,\ x_3x_4^3\\
5&x_1^5,\ x_2^5,\ x_3^5,\ x_4^5 .
\end{array}\tag{OCF25}
\]
Here is a complete proof. Call a positive-degree invariant
monomial indecomposable when it is not the product of two
positive-degree invariant monomials. The divisor proof preceding
OCF11 shows that each indecomposable has degree at most five.
If it contains both $x_1,x_4$, it is their product, since otherwise
division by this invariant quadratic leaves a nonconstant
invariant monomial. The same assertion holds for $x_2,x_3$.
After these cases are removed, at most two variables occur:
every three indices contain one of the two opposite pairs.
One variable gives precisely its fifth power. The remaining
two-element supports are $\{1,2\},\{1,3\},\{2,4\},\{3,4\}$.
For positive exponents $(a,b)$ with $a+b\leq5$, the equation
$a+2b\equiv0\pmod5$ has exactly $(a,b)=(3,1),(1,2)$, and
$a+3b\equiv0\pmod5$ has exactly $(a,b)=(2,1),(1,3)$.
For $\{2,4\}$ divide the weights by $2$ in the field
$\mathbb Z/5\mathbb Z$, giving the first list; for $\{3,4\}$
divide by $3$, giving the second. These are exactly the eight
mixed monomials of degrees three and four in OCF25. Every
listed monomial has no proper invariant divisor, as is also
seen from these lists and the excluded opposite pairs.
This proves the classification. Factoring any invariant
monomial repeatedly into indecomposables terminates because
positive total degree strictly decreases. Hence OCF25 generates
the entire invariant algebra.

Use $\mathcal G=\bigcup_{e=2}^5\mathcal G_e$ in place of
$\{x^\nu:\nu\in G\}$ in OCF12--16 and OCF20. The same
factorization argument proves the same exact invariant image;
all coefficient formulas, invariant certificates, pivot charts,
quotient maps and ranks remain valid. The new recursion has
$2,4,4,4$ generators in degrees $2,3,4,5$, respectively. Thus
for $d\geq13$ its precise matrix dimensions and rank are
\[
 C_d^{(14)}\in
 \operatorname{Mat}_{(16d-48)\,\times\,(112d-864)}(\mathbb C),
 \qquad \operatorname{rank}C_d^{(14)}=8d-32.
 \tag{OCF26}
\]
Indeed the number of generators is $14$ and the sum of their
degrees is $2\cdot2+4\cdot3+4\cdot4+4\cdot5=52$; substituting
the already proved $j_{d-e}=8(d-e)-32$ gives
$14(8d-32)-8\cdot52=112d-864$. At smaller degrees the exact
number of columns is the same sum of the appropriate $j_{d-e}$
from OCF19, with negative subscripts zero. All coefficients
in this reduced recursion are still the original coefficients
of OCF12--14.

% END CURRENT OCF

\clearpage
% BEGIN CURRENT OKM
% Complete independent metric/CD provider. Original source orders only.
\section{The exact common kernel frame and its original Gamma metrics}
\label{sec:OKM}

Retain the original simple quartet $m=1$, $k=4l+1$, $l\geq1$,
$q=(k+1)^2$, $c=k/2$, $0<\delta<1/2$ and $\gamma>2$.
The original period parameter and its branch are fixed, on the proved
original OCF rank domain $|u|\geq R_*$, with $R_*$ the retained
explicit period-orbit radius. The kernel
below is the kernel of that one fixed original observation
$\Lambda:E\to B$. It is common to the two source metrics and to
their four degree endpoints. No intersection over different period
parameters or different observations is made.

Order all pairs $(a,b)$, $0\leq a,b\leq k$, lexicographically once
and retain that order. Define
\[
 \begin{gathered}
 \beta_{ab}=c+(2a-k)\delta+i(2b-k)\gamma,\qquad
 \omega_{ab}=\frac{\beta_{ab}-c}{i}
                =(2b-k)\gamma-i(2a-k)\delta,\\
 \chi(S)=\prod_{a,b=0}^k(S-\beta_{ab}),\qquad
 E=\mathbb C[S]/(\chi),\qquad
 \varepsilon_{ab}=
 \left[\prod_{(a',b')\ne(a,b)}
       \frac{S-\beta_{a'b'}}{\beta_{ab}-\beta_{a'b'}}\right],\\
 V_\beta=((\beta_{ab})^d)_{(a,b),\,0\leq d<q}.
 \end{gathered}\tag{OKM1}
\]
The points are distinct: equality of their real and imaginary parts
forces both indices equal. Hence every denominator displayed above
is nonzero. Evaluation at all points gives
$\varepsilon_{ab}(\beta_{a'b'})=\mathbf1_{(a,b)=(a',b')}$.
Polynomial interpolation proves that these are a basis of $E$, and
the evaluation map from increasing polynomial coordinates to this
primary basis is exactly $V_\beta$. Its determinant is the ordered
Vandermonde product, so it is invertible.

\subsection{The exact dual and primal maps of the strip quotient}
Use the coefficient presentation of OQC7--8 in Segre coordinates.
A biform of bidegree $(k,k)$ has the full coefficient expression
\[
 F=\sum_{a,b=0}^k f_{ab}
       X_1^aX_0^{k-a}Y_1^bY_0^{k-b}.
 \tag{OKM2}
\]
The literal variable dictionary to the OCF biform coefficients is
$X_1=z_0$, $X_0=z_1$, $Y_1=w_0$, $Y_0=w_1$.
The Segre substitution has
$t_{++}=X_1Y_1$, $t_{+-}=X_1Y_0$,
$t_{-+}=X_0Y_1$, $t_{--}=X_0Y_0$.
At the original exponential vector take
$X_1=e^{-i\delta w}$, $X_0=e^{i\delta w}$,
$Y_1=e^{\gamma w}$, $Y_0=e^{-\gamma w}$.
Then the monomial with coefficient $f_{ab}$ is $e^{\omega_{ab}w}$.
The original coefficient identification includes the entire invertible
matrix $\mathsf H$ of OQC1 before this Segre presentation; its unit
and period entries remain in the matrices produced by the strip
calculation. OQC7 is characterized by evaluation on $e^{wY}[1]$.
Since
$e^{wY}[1]=\sum_{a,b}e^{\omega_{ab}w}\varepsilon_{ab}$,
the independent exponentials in OQC4 therefore give the precise
complex-linear dual pairing
\[
 \Theta(F)(\varepsilon_{ab})=f_{ab},\qquad
 \Theta(F)(x)=f^{\mathsf T}x_{\rm prim}.
 \tag{OKM3}
\]
There is no binomial coefficient or complex conjugation in this
complex-linear dual pairing. The statement follows by equating all
coefficients of the distinct exponentials, whose independence is
proved by their Vandermonde derivative matrix in OQC4.

Let $J$ be exactly the coefficient subspace of the original observed
dual image, as identified in OQC8. The original strip and invariant
elimination provider OCF constructs its literal surjection
\[
 \pi:\mathbb C^q\twoheadrightarrow\mathbb C^{r_K},
 \qquad\ker\pi=J,\qquad r_K=8k-16.
 \tag{OKM4}
\]
This input is the actual OCF matrix, including its stated period
coefficients and its valid pivot choice. The coefficient pairing
OKM3 and OQC8 identify
$K:=\ker\Lambda=J^\perp$ in primary coordinates, where
$J^\perp=\{x:f^{\mathsf T}x=0\text{ for every }f\in J\}$.
Consequently its exact primal frame and polynomial frame are
\[
 I_{\rm prim}=\pi^{\mathsf T}:\mathbb C^{r_K}\hookrightarrow E_{\rm prim},
 \qquad I_{\rm poly}=V_\beta^{-1}\pi^{\mathsf T}:
                        \mathbb C^{r_K}\hookrightarrow E.
 \tag{OKM5}
\]
Here is a direct proof, including inverse coordinates. Choose any of
the explicit right sections $L_\pi$ of the surjective matrix $\pi$,
so $\pi L_\pi=I_{r_K}$. For $x=\pi^{\mathsf T}v$ and $f\in J$,
$f^{\mathsf T}x=(\pi f)^{\mathsf T}v=0$. The transpose is
injective because $L_\pi^{\mathsf T}\pi^{\mathsf T}=I_{r_K}$.
Conversely, if $x\in J^\perp$, then
$(I-L_\pi\pi)f\in J$ for every $f$, so
$x^{\mathsf T}(I-L_\pi\pi)f=0$ for every $f$. Thus
$x=\pi^{\mathsf T}L_\pi^{\mathsf T}x$.
This proves equality of the image and $K$, and its inverse on that
image is $L_\pi^{\mathsf T}$. In polynomial coordinates the inverse
is $L_\pi^{\mathsf T}V_\beta$. This proves the full typed primal
arrow, rather than selecting a Hermitian transpose in place of it.

\subsection{The complete original source and quotient metric}
For either original source $s\in\{1,k\}$ write $b_s=s/2$ and
$M_s=(2\pi)^{s/2}$. The unchanged real-coordinate measure is
\[
 dm_{0,s}(y)=M_s\frac{2^{b_s}}{4\pi\Gamma(b_s)}
      |\Gamma(b_s/2+iy/2)|^2\,dy,
 \qquad S=c+iy.
 \tag{OKM6}
\]
The complete canonical OGN1--3 proof gives the monic real
orthogonal polynomials and all their squared norms:
\[
 \begin{gathered}
 p_0(y)=1,\quad p_1(y)=y,\quad
 p_{d+1}(y)=yp_d(y)-d(b_s+d-1)p_{d-1}(y),\\
 h_d=M_s d!(b_s)_d,\qquad
 u_d(S)=\frac{p_d((S-c)/i)}{\sqrt{h_d}}.
 \end{gathered}\tag{OKM7}
\]
The square root is positive. Thus $u_d$ has the original leading
coefficient $i^{-d}/\sqrt{h_d}$ in $S$ and is orthonormal on the
line in OKM6. No source mass or phase has been removed.

For a polynomial cutoff $D\geq q-1$ let
\[
 \begin{gathered}
 T_D=\left(\frac{p_d(\omega_{ab})}{\sqrt{h_d}}\right)_
                        {(a,b),\,0\leq d\leq D},
 \qquad R_D=T_DT_D^*,\\
 G_D^{\rm prim}=R_D^{-1},\qquad
 G_D^{\rm poly}=V_\beta^*R_D^{-1}V_\beta,\qquad
 H_D^K=\overline\pi\,R_D^{-1}\pi^{\mathsf T}.
 \end{gathered}\tag{OKM8}
\]
The star denotes conjugate transpose, while the superscript
$\mathsf T$ is ordinary transpose. In particular
$(\pi^{\mathsf T})^*=\overline\pi$, with precisely the indicated
matrix dimensions.

These are the actual minimum quotient metrics, as follows directly
from the original remainder map. In the orthonormal coefficient
frame $u_0,\ldots,u_D$, the remainder map followed by primary
evaluation is exactly $T_D$, column by column. Its first $q$
columns are invertible: the monic polynomials $p_0,\ldots,p_{q-1}$
evaluated at the distinct $\omega_{ab}$ have a nonzero Vandermonde
determinant, and all $\sqrt{h_d}$ are nonzero. Thus $R_D>0$.
For a prescribed primary vector $x$, the vector
$T_D^*R_D^{-1}x$ maps to $x$. It is orthogonal to $\ker T_D$,
because its inner product with a kernel vector contains $T_D$
applied to that vector. Every other lift differs by such a vector,
so its squared norm is
$x^*R_D^{-1}x$ plus the squared norm of that difference.
This proves the first quotient metric in OKM8.
The identity $x_{\rm prim}=V_\beta x_{\rm poly}$ proves the
second by direct pairing. Restriction to the exact frame OKM5
then gives
\[
 I_{\rm poly}^*G_D^{\rm poly}I_{\rm poly}
 = (\pi^{\mathsf T})^*R_D^{-1}\pi^{\mathsf T}
 =H_D^K>0.
 \tag{OKM9}
\]
The last positivity follows from injectivity of $\pi^{\mathsf T}$.
The map $K\hookrightarrow E$ is the fixed original coefficient
map. No invariance of that kernel under multiplication is used.

\subsection{Christoffel--Darboux with the exact source mass}
For arbitrary $z,w\in\mathbb C$ and $D\geq0$, the complete
polynomial identity is
\[
 \sum_{d=0}^D\frac{p_d(z)p_d(w)}{h_d}
 =\begin{cases}
 \displaystyle\frac{p_{D+1}(z)p_D(w)-p_D(z)p_{D+1}(w)}
                     {h_D(z-w)},&z\ne w,\\[6pt]
 \displaystyle\frac{p_{D+1}'(z)p_D(z)-p_D'(z)p_{D+1}(z)}{h_D},
                       &z=w.
 \end{cases}\tag{OKM10}
\]
To prove it, put $a_d=d(b_s+d-1)$, with $a_0=0$, and set
$p_{-1}=0$ for this recurrence expression. The recurrence
in OKM7 gives
\[
 \begin{split}
 (z-w)p_d(z)p_d(w)
 &=p_{d+1}(z)p_d(w)-p_d(z)p_{d+1}(w)\\
 &\quad-a_d\bigl(p_d(z)p_{d-1}(w)-p_{d-1}(z)p_d(w)\bigr).
 \end{split}\tag{OKM11}
\]
For $d\geq1$, the full norm identity is $a_d/h_d=1/h_{d-1}$.
Divide OKM11 by $h_d$ and sum. Every adjacent term cancels,
leaving just the degree-$D$ numerator divided by $h_D$; the
$d=0$ lower term is zero. This proves the first line of OKM10.
Both sides of the resulting equality after multiplication by
$z-w$ are polynomials. Differentiating with respect to $z$ and
then setting $z=w$ proves the second line with the displayed
order and sign. This also proves the value at every confluent
point without an assertion about a limiting singular quotient.
The denominator is the full $h_D=M_sD!(b_s)_D$, including $M_s$.

Since all coefficients of the recurrence are real, the matrix
entries in OKM8 are therefore exactly the values of OKM10 at
\[
 z=\omega_{ab},\qquad w=\overline{\omega_{a'b'}}.
 \tag{OKM12}
\]
Their indexed conjugation and confluent locus are
\[
 \begin{gathered}
 \overline{\omega_{a'b'}}=\omega_{k-a',b'},\qquad
 \omega_{ab}-\overline{\omega_{a'b'}}
       =2(b-b')\gamma-2i(a+a'-k)\delta,\\
 \omega_{ab}=\overline{\omega_{a'b'}}
       \quad\Longleftrightarrow\quad b=b'\text{ and }a+a'=k.
 \end{gathered}\tag{OKM13}
\]
These follow by substituting OKM1; positivity of $\delta,\gamma$
proves both implications in the last line. Thus the conjugation
map on the point indices is precisely $(a,b)\mapsto(k-a,b)$.
The confluent formula of OKM10 applies at these indexed entries.
For the retained odd $k$, its matrix diagonal has $a'=a,b'=b$,
so $2a=k$ cannot occur; the literal Hermitian diagonal is given
by the first line of OKM10. The conjugate-index entries use its
second line. All entries, including the positive diagonal, are
simultaneously defined by the original polynomial sum on the
left of OKM10.

\subsection{Every pivot change and the four original endpoints}
Let $\pi'$ be another valid strip/invariant elimination matrix for
the same original $J$, with the same surjective codomain dimension.
Choose right sections $L_\pi,L_{\pi'}$ and put
\[
 A_{\pi',\pi}=\pi'L_\pi,
 \qquad A_{\pi,\pi'}=\pi L_{\pi'}.
 \tag{OKM14}
\]
Because $\pi'$ vanishes on $J=\ker\pi$, one has
$\pi' L_\pi\pi=\pi'$. Thus
$\pi'=A_{\pi',\pi}\pi$. Reversing the roles gives the
opposite identity, and surjectivity of each matrix then proves
$A_{\pi',\pi}A_{\pi,\pi'}=I$ and
$A_{\pi,\pi'}A_{\pi',\pi}=I$. This proves invertibility and
also independence from the chosen right sections. Writing
$A=A_{\pi',\pi}$ for these identities gives the exact frame
and Gram changes
\[
 \begin{gathered}
 I'_{\rm prim}=I_{\rm prim}A^{\mathsf T},\qquad
 I'_{\rm poly}=I_{\rm poly}A^{\mathsf T},\\
 (H_D^K)'=\overline A H_D^K A^{\mathsf T},\qquad
 \det(H_D^K)'=|\det A|^2\det H_D^K.
 \end{gathered}\tag{OKM15}
\]
The first two formulas follow by transposing $\pi'=A\pi$.
Substitution into OKM9 proves the congruence and determinant.
This retains all actual pivot determinants. They have not been
assigned modulus one.

The kernel is the same fixed $K=\ker\Lambda$ used in MRI1--5.
If its original BRU inclusion is $I_0$ and its fixed-section kernel
coordinate map is $\kappa_0$, let
\[
 \mathsf C_K=\kappa_0I_{\rm poly},\qquad
 \mathsf C_K^{-1}=L_\pi^{\mathsf T}V_\beta I_0.
 \tag{OKM16}
\]
The equalities $I_0\kappa_0x=x$ for $x\in K$ and
$L_\pi^{\mathsf T}V_\beta I_{\rm poly}=I_{r_K}$ prove that
these are inverse matrices and $I_0\mathsf C_K=I_{\rm poly}$.
Consequently
$H_D^K=\mathsf C_K^*(I_0^*G_D^{\rm poly}I_0)\mathsf C_K$.
The matrix $\mathsf C_K$ is independent of $D$ and of the choice of source
order $s$. The source metrics change; the two specified coefficient
frames do not.

It follows that the complete original four-endpoint kernel scalar is
\[
 \boxed{\displaystyle
 F_K^{(s)}=
 \log\det H_{q-1}^K+\log\det H_q^K
       -\log\det H_{2q-1}^K-\log\det H_{2q}^K.}
 \tag{OKM17}
\]
Indeed MRI3 gives the actual consecutive kernel determinant ratio
\[
 \frac{\det(I_0^*G_{q+j-1}^{\rm poly}I_0)}
      {\det(I_0^*G_{q+j}^{\rm poly}I_0)}
 =\frac{1+t_j}{1+\xi_j}.
\]
Each factor $|\det\mathsf C_K|^2$ cancels in that ratio by OKM16.
The weighted sum with endpoint weights $1,2,\ldots,2,1$ therefore
has precisely the four signs in OKM17, as in MRI4. This verifies
the original degree cutoffs, rather than shifting the column count
in $T_D$. Under any valid pivot change OKM15, the added scalar
$2\log|\det A|$ has total coefficient $1+1-1-1=0$, proving
invariance of OKM17 with its original frames.

For emphasis, the full mass in OKM7 appears in every column of
$T_D$ and in the denominator of OKM10. If that common mass is
factored explicitly, write $T_D=M_s^{-1/2}\widehat T_D$, where
$\widehat T_D$ has exactly the same numerator polynomials and
the denominator $\sqrt{d!(b_s)_d}$. Then
$H_D^K=M_s\overline\pi(\widehat T_D\widehat T_D^*)^{-1}
\pi^{\mathsf T}$, so its determinant has factor $M_s^{r_K}$.
It cancels only in the displayed signed four-endpoint scalar,
whose four ranks are equal. The individual metrics in OKM8--9
retain it. No inserted Gamma order, discarded action defect, or
kernel-invariance assertion enters any of these calculations.

% END CURRENT OKM

\clearpage
% BEGIN CURRENT OKA
\section{The original kernel graph, boundary basis and complete action}
\label{sec:OKA}

This section uses the complete conductor coefficient construction OCF
and its exact original dual identification in OKM. Fix the original
simple quartet, its complete amplitude unit and periods, the original
parameter on any logarithm branch with $|u|\geq R_*$, and
$k=4l+1$, $l\geq1$. Put $q=(k+1)^2$, $r=8k-16$, and
$\mathcal I_k=\{(a,b):0\leq a,b\leq k\}$. The roots, centre and
centred eigenvalues remain
\[
 c=k/2,\qquad
 \beta_{ab}=c+(2a-k)\delta+i(2b-k)\gamma,\qquad
 \omega_{ab}=(\beta_{ab}-c)/i
             =(2b-k)\gamma-i(2a-k)\delta.
 \tag{OKA1}
\]
The original algebra is $E=\mathbb C[S]/(\chi)$, with
$\chi(S)=\prod_{(a,b)\in\mathcal I_k}(S-\beta_{ab})$.
Let $\varepsilon_{ab}$ be its original primary idempotents and
$\mathcal V:E\to\mathbb C^{\mathcal I_k}$ the evaluation
isomorphism $[p]\mapsto(p(\beta_{ab}))_{ab}$. Its inverse sends
$e_{ab}$ to
\[
 \varepsilon_{ab}(S)
 =\frac{\chi(S)}{(S-\beta_{ab})\chi'(\beta_{ab})}.
 \tag{OKA2}
\]
These denominators are nonzero: equality of two roots in OKA1,
on comparing real and imaginary parts and using $\delta>0$,
$\gamma>2$, implies equality of both indices. Evaluation proves
that the polynomials in OKA2 are the coordinate idempotents,
so their columns give the exact inverse, with no changed root
or polynomial coefficient convention.

\subsection{The coefficient quotient gives a graph in the original primary basis}

Let $\pi:\mathbb C^{\mathcal I_k}\to\mathbb C^r$ be the
surjective coefficient quotient of OCF, whose kernel is precisely
the image $J_k$ of the invariant polynomials under the original
restriction. The domain here consists of the coefficients of the
biforms
$z_0^a z_1^{k-a}w_0^b w_1^{k-b}$, and is the complex-linear
dual of the displayed primary coordinate space. Let
$T\subset\mathcal I_k$ be the final OCF complement of its
invariant-image pivot rows, viewed as a subset of the full
biform coefficient grid. Its ordered coordinate inclusion is
$E_T:\mathbb C^r\to\mathbb C^{\mathcal I_k}$. The explicit
OCF formula gives
\[
 \pi E_T=I_r,\qquad
 N=\pi^T:\mathbb C^r\longrightarrow\mathbb C^{\mathcal I_k},
 \qquad E_T^TN=I_r.
 \tag{OKA3}
\]
Indeed a retained strip monomial in $T$ already has zero
conductor remainder outside its own coordinate, zero coordinates
on the invariant pivot rows, and its indicated coordinate equal
to one. Both reductions in OCF therefore leave that quotient
coordinate equal to its standard unit vector. Transposing gives
the last identity. This uses ordinary transpose, because the
restriction quotient and its dual are complex-linear maps.

Write $J=\mathcal I_k\setminus T$ only for the complementary
index set in this section; the invariant image retains the name
$J_k$. Denote coordinate restrictions by $x_T,x_J$, and put
$N_J=E_J^TN$. The original kernel inclusion is consequently
\[
 I_K=\mathcal V^{-1}N:\mathbb C^r\hookrightarrow E,
 \qquad
 \operatorname{im}I_K=K:=\ker\Lambda,\qquad
 N=\begin{pmatrix}I_r\\N_J\end{pmatrix}_{T,J}.
 \tag{OKA4}
\]
For completeness, a primary vector $v$ pairs with a biform
coefficient vector $f$ by $f^Tv$. The annihilator of
$\ker\pi=J_k$ is $\operatorname{im}\pi^T$: containment
follows by multiplication, while surjectivity of $\pi$ gives
both subspaces dimension $r$. OQC8 identifies precisely this
annihilator with the primary coordinates of the fixed original
$\ker\Lambda$. This proves both the original image assertion
and its direction in OKA4. In polynomial coordinates the same
inclusion is $V_\beta^{-1}N$, where
$(V_\beta)_{(a,b),d}=\beta_{ab}^d$ for $0\leq d<q$;
this is the full coefficient matrix of OKA2.

\subsection{An explicit basis of the actual boundary}

For $j\in J$ define the actual boundary vector
$b_j=\Lambda\varepsilon_j$. These vectors form a basis of
the original image $B=\operatorname{im}\Lambda$. To prove
independence, if $\sum_{j\in J}c_j\varepsilon_j$ lies in
$K$, its primary coordinates equal $N\zeta$. Its $T$
coordinates vanish, so OKA3 gives $\zeta=0$ and all $c_j=0$.
To prove spanning, use the literal decomposition
\[
 x=N x_T+E_J(x_J-N_Jx_T),\qquad
 \Lambda\mathcal V^{-1}x
       =\sum_{j\in J}b_j(x_J-N_Jx_T)_j.
 \tag{OKA5}
\]
The first equality follows separately on its $T$ and $J$
coordinates. Apply $\Lambda$, which kills the first term by
OKA4, to obtain the second. Thus no observation has been
substituted: in this basis the specified original observation
has the exact coefficient matrix $(-N_J,I_{q-r})$.

Here are the complete original period and tensor coordinates of
these boundary vectors. Write the quartet roots as
$\alpha_{\epsilon,\eta}=1/2+\epsilon\delta+i\eta\gamma$,
where $\epsilon,\eta\in\{+1,-1\}$ and the signs $+,-$ in
subscripts denote these same two numbers. In a subscript
$\alpha$ below means that complex root, indexed by its sign
pair. Keep the original one-factor primary idempotents
$\varepsilon^h_\alpha\in E_h$. Set
\[
 v_\alpha=\Pi U\varepsilon^h_\alpha
       =v_h(\alpha)\Pi\varepsilon^h_\alpha,
 \qquad
 Z_+=z_0,\ Z_-=z_1,\ W_+=w_0,\ W_-=w_1,
 \quad
 \mathfrak v(z,w)=\sum_{\epsilon,\eta}
                         v_{\epsilon\eta}Z_\epsilon W_\eta.
 \tag{OKA6}
\]
The vectors belong to the original one-factor period space;
all entries of $\Pi$ and all four actual nonzero amplitude
values remain present. For every $(a,b)\in\mathcal I_k$,
\[
 \boxed{\quad
 \jmath\Lambda\varepsilon_{ab}
  =[z_0^a z_1^{k-a}w_0^b w_1^{k-b}]
       P_k\bigl(\mathfrak v(z,w)^{\otimes k}\bigr),\qquad
 P_k=\frac15\sum_{h=0}^4(C^{-h})^{\otimes k}.
 \quad}\tag{OKA7}
\]
The original sum inclusion $\iota:E\hookrightarrow E_h^{\otimes k}$
sends $\varepsilon_{ab}$ to the sum of the primary tensor
idempotents whose word has $a$ positive real signs and $b$
positive imaginary signs. This assertion follows by evaluation:
on a primary word, the sum variable takes exactly the value
$\beta_{ab}$ of its two counts, and OKA2 evaluates to one
for those counts and zero otherwise. Applying
$(\Pi U)^{\otimes k}$ gives precisely the corresponding sum
of the tensors of the $v_\alpha$ in OKA6. Expanding the tensor
power and taking the indicated biform coefficient selects
exactly the same words, each once. The original definition
$\jmath\Lambda=P_k(\Pi U)^{\otimes k}\iota$ proves OKA7,
including the literal inverse powers and $1/5$. Its specialization
to $j\in J$ therefore supplies the actual basis $b_j$ of
OKA5. Neither a different target nor a numerical period matrix
has been introduced.

The fixed, source-independent kernel coordinate and section are
now explicitly
\[
 \begin{gathered}
 \kappa=E_T^T\mathcal V:E\to\mathbb C^r,
 \qquad S_*:B\to E,\quad S_*b_j=\varepsilon_j,
 \\
 I_K\kappa+S_*\Lambda=I_E,
 \quad\kappa I_K=I_r,\quad\kappa S_*=0,\quad\Lambda S_*=I_B.
 \end{gathered}
 \tag{OKA8}
\]
Every equality follows by applying OKA5 to an arbitrary primary
vector. In particular this is an actual splitting of the fixed
original short exact sequence, used unchanged for both source
orders and for all polynomial degrees.

\subsection{The whole original action and its retained defect}

The original centred sum action is $Y=(S-c)/i$ on $E$.
Under $\mathcal V$ it is the diagonal matrix of the
$\omega_{ab}$ in OKA1. Write its two diagonal submatrices
as $\Omega_T$ and $\Omega_J$. In the exact coordinates
$e=I_K\zeta+S_*b$ of OKA8, it has the complete matrix
\[
 \boxed{\quad
 Y\ \longleftrightarrow\
 \begin{pmatrix}
   \Omega_T&0\\
   \Omega_JN_J-N_J\Omega_T&\Omega_J
 \end{pmatrix},\qquad
 (\Omega_JN_J-N_J\Omega_T)_{jt}
                 =(\omega_j-\omega_t)(N_J)_{jt}.
 \quad}\tag{OKA9}
\]
Indeed the primary vector of $e$ is
$(\zeta,N_J\zeta+b)$. Applying the original diagonal matrix
gives $(\Omega_T\zeta,\Omega_JN_J\zeta+\Omega_Jb)$.
Applying the coordinates in OKA5 subtracts
$N_J\Omega_T\zeta$ from its second component and proves
every block in OKA9. This proves, with its precise codomain,
\[
 YI_K-I_K\Omega_T
     =S_*\bigl(\Omega_JN_J-N_J\Omega_T\bigr)
                :\mathbb C^r\to E,\qquad
 \Lambda YI_K=\Omega_JN_J-N_J\Omega_T:\mathbb C^r\to B,
 \tag{OKA10}
\]
where the matrix on the right uses the basis $b_j$.
The chosen complementary span of the original primary idempotents
is stable under $Y$; no stability of $K$ has been used. Its
entire failure is the displayed coefficient map.

For each original source $s\in\{1,k\}$ retain the monic
Gamma polynomials $p_d^{(s)}$, the complete norms
$h_d^{(s)}=(2\pi)^{s/2}d!(s/2)_d$, and the original
$e_d^{(s)}=[p_d^{(s)}(Y)]/\sqrt{h_d^{(s)}}$. Then the
fixed-section kernel and boundary coordinates of every actual
remainder, including all degrees $d\geq q$, are
\[
 \begin{split}
 (\kappa e_d^{(s)})_t
      &=\frac{p_d^{(s)}(\omega_t)}{\sqrt{h_d^{(s)}}},\\
 (\Lambda e_d^{(s)})_j
      &=\frac{p_d^{(s)}(\omega_j)
                -\sum_{t\in T}(N_J)_{jt}p_d^{(s)}(\omega_t)}
              {\sqrt{h_d^{(s)}}}.
 \end{split}\tag{OKA11}
\]
Polynomial evaluation in the original primary algebra gives the
full vector $p_d^{(s)}(\omega)/\sqrt{h_d^{(s)}}$ even when
its polynomial representative must be reduced modulo $\chi$.
Applying OKA5 and OKA8 proves OKA11. Thus the formula includes
the original finite remainder operation and every phase and source
mass. Composing its second line with the actual columns OKA7
gives exactly the period-vector increments of OPG. Its first line
is the fixed-section kernel vector used by OPG and OPR; the
degree-dependent minimum-section correction in OPR remains
required for the mixed norm. The metric computation in OKM
applies to the common inclusion OKA4 and does not replace that
correction by the fixed-section vector alone.

% END CURRENT OKA

\clearpage
% BEGIN CURRENT OMG
% Full original multiplicity geometry. No frozen predecessor is edited.
\section{The complete homogeneous geometry of the original multiplicity curve}
\label{sec:OMG}

\subsection{The actual primary, unit, period and Fourier coefficient map}
Retain the original full-order quartet and every original period
parameter. Write
\[
\begin{gathered}
 \alpha_{\epsilon,\eta}=\tfrac12+\epsilon\delta+i\eta\gamma,
 \quad\epsilon,\eta\in\{-1,1\},\quad0<\delta<\tfrac12,\quad\gamma>2,\\
 m\ge1,\quad d_m=m-1,\quad n=4m,\quad N=4m+1,\quad
 h(z)=\prod_{\epsilon,\eta}(z-\alpha_{\epsilon,\eta})^m,\\
 E_h=\mathbb C[z]/(h),\quad A_h=M_z,\quad
 v_h=2\xi/h,\quad U=M_{j_hv_h},\quad v_h(\alpha_{\epsilon,\eta})\ne0.
\end{gathered}\tag{OMG1}
\]
The last inequalities follow from the stipulated complete original
zero order $m$. The map $j_h$ records all divided Taylor coefficients
of orders $0,\ldots,m-1$ at all four roots. Let
$\mathcal E_{\rm prim}:\mathbb C^{4m}\to E_h$ send its coordinate
$(\alpha,d)$ to the original primary element $\varepsilon_\alpha
(z-\alpha)^d$, $0\le d<m$. Explicitly its representative is
\[
 E_{\alpha,d}(z)=\operatorname{rem}_h\left[
 h_\alpha(z)\sum_{j=0}^{m-1}
       \frac{(h_\alpha^{-1})^{(j)}(\alpha)}{j!}(z-\alpha)^{j+d}
                                      \right],\quad
 h_\alpha=h/(z-\alpha)^m.
\]
The Taylor inverse is defined since $h_\alpha(\alpha)\ne0$.
It gives the idempotent one on its own primary factor and zero
on the others. The four coprime factors therefore prove that
$\mathcal E_{\rm prim}$ is an isomorphism and include every
nilpotent degree. In this basis the full matrix of $U$ is
\[
 U_{\rm prim;\alpha,a;\alpha,d}
 =\begin{cases}v_h^{(a-d)}(\alpha)/(a-d)!,&a\ge d,\\0,&a<d,
   \end{cases}\qquad
 \det U=\prod_\alpha v_h(\alpha)^m\ne0.
 \tag{OMG2}
\]
Other primary blocks are zero. Inverting its constant coefficient
plus its nilpotent part by the finite geometric series proves
invertibility with all derivatives retained.

Keep the original $u\ne0$ and its logarithm branch, the primitive
$\Phi_h'=h$, $\Phi_h(0)=0$, the outward rays $\ell_j$ of angles
$(\arg u+\pi+2\pi j)/N$, and $\Gamma_j=-\ell_0+\ell_j$.
The exact period and Fourier maps are
\[
\begin{gathered}
 (\Pi[P])_j=\int_{\Gamma_j}e^{\Phi_h(z)/u}
                     \operatorname{rem}_hP(z)\,dz,\quad1\le j\le n,\\
 \zeta=e^{2\pi i/N},\quad F_{jt}=\zeta^{jt}-1,\quad
 (F^{-1})_{tj}=\zeta^{-tj}/N\quad(1\le j,t\le n),\quad
 C=F\operatorname{diag}(\zeta,\ldots,\zeta^n)F^{-1}.
\end{gathered}\tag{OMG3}
\]
The full original period theorem PDE/OPG gives the convergence and
invertibility of $\Pi$ on these representatives. The Fourier inverse
follows directly from
$\sum_{j=1}^{N-1}\zeta^{-tj}(\zeta^{sj}-1)=N\mathbf1_{t=s}$.
It is the same literal cyclic matrix as OCS4, including its inverse
powers in the original average. Its inherited Gram is
$F^*F=N(I_n+\mathbf1\mathbf1^*)$, as expansion of its entries gives.

Define the exact diagonal map and coefficient matrix
\[
\begin{gathered}
 D_m^{\rm jet}=\operatorname{diag}\left(i^{-d}/d!\right)_{\alpha,\,0\le d<m},
 \qquad \mathsf H_m=F^{-1}\Pi U\mathcal E_{\rm prim}D_m^{\rm jet},\\
 (\mathsf H_m)_{t;\alpha,d}
  =\frac{i^{-d}}{d!}\sum_{a=d}^{m-1}
     \left[\frac1N\sum_{j=1}^{n}\zeta^{-tj}
              \int_{\Gamma_j}e^{\Phi_h(z)/u}E_{\alpha,a}(z)\,dz\right]
                  \frac{v_h^{(a-d)}(\alpha)}{(a-d)!}.
\end{gathered}\tag{OMG4}
\]
This is exactly OCS13, now regarded as one square matrix. All
factors in its first line are invertible. In particular
\[
 \det\mathsf H_m=(\det F)^{-1}\det\Pi\,
      \left(\prod_\alpha v_h(\alpha)^m\right)
      \det\mathcal E_{\rm prim}
      \prod_\alpha\prod_{d=0}^{m-1}\frac{i^{-d}}{d!}\ne0,
\]
with the determinant of $\mathcal E_{\rm prim}$ taken in the fixed
original coefficient frame. None of these factors is prescribed.

Set $\omega_{\epsilon,\eta}=\eta\gamma-i\epsilon\delta$ and let
\[
 z_{\epsilon,\eta,d}(w)=w^d e^{\omega_{\epsilon,\eta}w},\qquad
 x(w)=F^{-1}\Pi Ue^{w(A_h-1/2)/i}[1].
 \tag{OMG5}
\]
The terminating exponential on each primary factor has coefficient
$i^{-d}w^d/d!$ on $\varepsilon_\alpha(z-\alpha)^d$.
Multiplication by OMG2 and application of the literal period maps
therefore prove the exact identity and inverse coordinate map
\[
 \boxed{x(w)=\mathsf H_m z(w),\qquad z(w)=\mathsf H_m^{-1}x(w).}
 \tag{OMG6}
\]
For the original Euclidean period metric, its Gram on these primary
curve coordinates is $\mathsf H_m^*N(I_n+\mathbf1\mathbf1^*)
\mathsf H_m$. Thus the coefficient change also retains its full
metric and all off-diagonal entries, without claiming orthogonality.

\subsection{The complete homogeneous ideal, with quadratic generators}
Use binary indices $a,b\in\{0,1\}$ for
$\epsilon=2a-1$, $\eta=2b-1$. Let
$\mathscr P_Z=\mathbb C[Z_{a,b,d}:a,b\in\{0,1\},0\le d\le d_m]$
with each variable of degree one. Define the algebra map
\[
 \Phi_m:\mathscr P_Z\longrightarrow
        \mathbb C[X_0,X_1,Y_0,Y_1,T_0,T_1],\qquad
 Z_{a,b,d}\longmapsto X_aY_bT_0^{d_m-d}T_1^d.
 \tag{OMG7}
\]
Its image is the complete diagonal graded algebra
\[
 R_m=\bigoplus_{k\ge0}
  \mathbb C[X_0,X_1]_k\otimes\mathbb C[Y_0,Y_1]_k
                              \otimes\mathbb C[T_0,T_1]_{kd_m}.
 \tag{OMG8}
\]
Indeed a basis monomial of its degree-$k$ component is
\[
 X_0^{k-a}X_1^aY_0^{k-b}Y_1^bT_0^{kd_m-D}T_1^D,
       \quad0\le a,b\le k,\quad0\le D\le kd_m.
 \tag{OMG9}
\]
It is a product of $k$ images in OMG7: choose $k$ binary first
indices with sum $a$, independently $k$ binary second indices with
sum $b$, and $k$ integers in $[0,d_m]$ with sum $D$; the last
choice follows by filling successive slots up to $d_m$. Combining
the three lists into rows gives the required preimage monomial.
Distinct displayed target monomials are linearly independent.

Let $\mathfrak t_m$ be generated by every quadratic binomial
\[
 Z_{a,b,d}Z_{a',b',d'}-Z_{c,e,f}Z_{c',e',f'}
 \quad\text{with}\quad
 a+a'=c+c',\quad b+b'=e+e',\quad d+d'=f+f',
 \tag{OMG10}
\]
where every index belongs to its original range. Repeated variables
are allowed. Then
\[
 \boxed{\ker\Phi_m=\mathfrak t_m,\qquad
           \mathscr P_Z/\mathfrak t_m\simeq R_m.}
 \tag{OMG11}
\]
Here is a complete ideal proof. Each displayed quadratic has zero
image. Two degree-$k$ monomials have the same image exactly when
the sums of their first indices, second indices and third indices
agree. Regard their $k$ factors temporarily as labelled rows. Align
the first indices with the target list by pairwise swaps of first
indices, keeping each row's other indices fixed. Each swap is an
OMG10 relation. Next align the second indices by the same operation,
keeping the aligned first indices fixed. Finally, if the third
indices differ from their target list, choose a row below its target
and a row above its target. Increase the former third index by one
and decrease the latter by one. These entries stay in $[0,d_m]$
because their target entries do; the total sum is unchanged and the
sum of absolute deviations decreases by two. This terminates at
the target list. Each transfer is another OMG10 relation. Thus the
difference between any two monomials with the same image is a sum
of multiples of the quadrics, by telescoping the successive monomial
differences. Labels only specify the sequence of polynomial maps;
the products are the original commutative monomials throughout.
For any polynomial in $\ker\Phi_m$, collect its coefficients by
their target monomial. Each such coefficient sum is zero by target
monomial independence. Subtracting one chosen monomial in each
group expresses that whole group as a linear combination of the
proved monomial differences. This proves the reverse inclusion for
all degrees and hence the entire ideal equality.

Every degree-$k$ original curve monomial has the function
\[
 w^D\exp\{[(2b-k)\gamma-i(2a-k)\delta]w\}
       \quad(0\le a,b\le k,\ 0\le D\le kd_m).
 \tag{OMG12}
\]
The frequencies are distinct: equality of real parts gives equality
of $b$ and equality of imaginary parts gives equality of $a$.
To prove independence with all polynomial powers, suppose
$\sum_{a,b}P_{a,b}(w)e^{\omega_{a,b}w}=0$, with
$\deg P_{a,b}\le kd_m$. Fix a pair and apply
$\prod_{(a',b')\ne(a,b)}(\partial_w-\omega_{a',b'})^{kd_m+1}$.
This kills every other summand. On its surviving polynomial each
factor is $\partial_w+(\omega_{a,b}-\omega_{a',b'})$, invertible
on the space of degree at most $kd_m$ because differentiation is
nilpotent and the constant is nonzero. The surviving polynomial
is therefore zero. Repeating this for every pair proves independence.

Thus the full homogeneous ideal of the original primary curve is
exactly $\mathfrak t_m$: a homogeneous polynomial restricts to zero
if and only if its coefficient sums at every target monomial OMG9
are zero. In particular
\[
 \boxed{\dim R_{m,k}=(k+1)^2[1+k(m-1)]=q_k,
 \quad\ker\{\mathscr P_{Z,k}\to\operatorname{Hol}(\mathbb C),
                      f\mapsto f(z(w))\}=\mathfrak t_{m,k}.}
 \tag{OMG13}
\]
This statement concerns the homogeneous ideal, equivalently the
projective closure of $[z(w)]$. Every one of its degrees has been
computed, including degree zero. No assertion about a different
nonhomogeneous affine ideal is used.

\subsection{The projective model and its complete chart inverses}
For $m\ge2$, the morphism
\[
 (\mathbb P^1)^3\longrightarrow\mathbb P^{4m-1},\qquad
 ([X_0:X_1],[Y_0:Y_1],[T_0:T_1])
      \longmapsto[X_aY_bT_0^{d_m-d}T_1^d]_{a,b,d}
 \tag{OMG14}
\]
has multidegree $(1,1,m-1)$. At least one first, second and third
coordinate is nonzero; their corresponding extreme product is
nonzero, so the map is everywhere defined. Its image and scheme
structure are exactly $\operatorname{Proj}R_m$, with complete
closed ideal OMG10. Here are explicit inverses proving this claim.
The eight charts $Z_{a,b,0}\ne0$ and $Z_{a,b,d_m}\ne0$ cover
$\operatorname{Proj}R_m$. Indeed OMG11 gives
\[
 Z_{a,b,d}^{d_m}=Z_{a,b,0}^{d_m-d}Z_{a,b,d_m}^{d}
             \quad\text{in }R_m.
\]
A homogeneous prime containing every extreme variable contains
every variable by this equality, and hence gives no point of Proj.
On $Z_{a,b,0}\ne0$ the three inverse affine coordinates are
\[
 x=Z_{1-a,b,0}/Z_{a,b,0},\quad
 y=Z_{a,1-b,0}/Z_{a,b,0},\quad
 t=Z_{a,b,1}/Z_{a,b,0}.
\]
For every variable its ratio to $Z_{a,b,0}$ is
$x^{\mathbf1_{a'\ne a}}y^{\mathbf1_{b'\ne b}}t^d$, by OMG7
and its proved injective quotient. Consequently the degree-zero
localized ring is $\mathbb C[x,y,t]$. These three coordinates
are algebraically independent, since their images in the parameter
fraction field are $X_{1-a}/X_a,Y_{1-b}/Y_b,T_1/T_0$.
For the chart at $Z_{a,b,d_m}$ replace the last ratio by
$t=Z_{a,b,d_m-1}/Z_{a,b,d_m}$; the general ratio has exponent
$d_m-d$ in $t$. Its coordinate ring is again $\mathbb C[x,y,t]$.
On chart overlaps these formulas invert the same parameter ratios,
so they glue to the inverse of OMG14. This proves the isomorphism
of schemes and its full defining homogeneous ideal. For $m=1$
the third factor has degree zero and the same argument uses the
four first/second charts: the model is the Segre $(\mathbb P^1)^2$
with ideal $Z_{0,0,0}Z_{1,1,0}-Z_{0,1,0}Z_{1,0,0}$.

In the actual Fourier coordinates let
$\mathscr P_x=\mathbb C[x_1,\ldots,x_n]$ and define
\[
 \mathfrak a=\{f(x):f(\mathsf H_mZ)\in\mathfrak t_m\},\qquad
 R=\mathscr P_x/\mathfrak a.
 \tag{OMG15}
\]
The invertible substitution OMG6 proves that $\mathfrak a$ is the
complete homogeneous ideal of $x(w)$ and that $R\simeq R_m$.
Its explicit generators are all OMG10 quadrics with every $Z$
replaced by its component of $\mathsf H_m^{-1}x$. These are the
actual period and unit coefficients of OMG4, and $\mathfrak a$
is prime because the target ring in OMG7 is a domain. Thus the
original curve's projective closure is the precise linear image
of OMG14 under $\mathsf H_m$, with no prescribed substitute matrix.

\subsection{The exact graded source-dual algebra and every factorial}
For each $k\ge0$ retain the original sum algebra
\[
\begin{gathered}
 e_k=1+kd_m,\quad c_k=k/2,\quad
 \beta_{k,a,b}=c_k+(2a-k)\delta+i(2b-k)\gamma,\\
 \chi_k(S)=\prod_{a,b=0}^k(S-\beta_{k,a,b})^{e_k},\quad
 E_k=\mathbb C[S]/(\chi_k),\quad Y_k=(M_S-c_kI)/i.
\end{gathered}\tag{OMG16}
\]
For $k=0$ this is $E_0=\mathbb C[S]/(S)=\mathbb C$. Let
$H_{k,a,b,D}=\varepsilon_{k,a,b}(S-\beta_{k,a,b})^D$,
$0\le D<e_k$, be the complete original primary basis obtained
by the same Taylor inverse as in OMG2. It satisfies
\[
 e^{wY_k}[1]=\sum_{a,b=0}^k e^{\omega_{k,a,b}w}
             \sum_{D=0}^{e_k-1}\frac{i^{-D}w^D}{D!}H_{k,a,b,D},
 \quad\omega_{k,a,b}=(2b-k)\gamma-i(2a-k)\delta.
 \tag{OMG17}
\]
For $f\in\mathscr P_{x,k}$ expand its original function uniquely as
$f(x(w))=\sum c_{a,b,D}(f)w^De^{\omega_{k,a,b}w}$ using OMG12.
Define the complex linear dual map, without a conjugation,
\[
 \Theta_k(f)(H_{k,a,b,D})=i^DD!c_{a,b,D}(f).
 \tag{OMG18}
\]
Then OMG17 proves exactly
$\Theta_k(f)(e^{wY_k}[1])=f(x(w))$.
The independence and surjectivity in OMG9--13 prove that this
map is onto $E_k^\vee$, with kernel $\mathfrak a_k$.
Consequently it induces a specified vector-space isomorphism
\[
 \boxed{\Theta_k:R_k\xrightarrow{\ \sim\ }E_k^\vee.}
 \tag{OMG19}
\]
All phases and factorials are those of OCS17; for a monomial
$x^\nu$ its coefficients are precisely the finite multinomial
OCS15 rule with the entries OMG4.

The collection of these vector spaces has an exact graded algebra
structure coming from the original sum maps. For $k,j\ge0$ put
\[
 \Delta_{k,j}:E_{k+j}\longrightarrow E_k\otimes E_j,
               \qquad[P(S)]\longmapsto[P(S_1+S_2)].
 \tag{OMG20}
\]
On a primary pair the eigenvalue is the displayed sum
$\beta_{k,a,b}+\beta_{j,a',b'}=\beta_{k+j,a+a',b+b'}$.
The sum of the two nilpotents has index $e_k+e_j-1=e_{k+j}$:
its power $e_k+e_j-2$ has nonzero top coefficient
$\binom{e_k+e_j-2}{e_k-1}$ and its next power is zero. Every
pair of total sign counts occurs. Thus the full minimal polynomial
of $S_1+S_2$ is exactly $\chi_{k+j}$, proving both well-definedness
and injectivity of OMG20. It also gives the complete primary formula
\[
 \Delta_{k,j}H_{k+j,A,B,D}
 =\sum_{\substack{a+a'=A,\ b+b'=B}}
   \sum_{\substack{d+d'=D\\0\le d<e_k,\ 0\le d'<e_j}}
       \binom Dd H_{k,a,b,d}\otimes H_{j,a',b',d'}.
 \tag{OMG21}
\]
Indeed the total primary idempotent is one precisely on the
matching eigenvalue pairs, and the nilpotent power expands by the
ordinary binomial theorem there. All other components vanish.
Substitution of the sum of three coordinates proves coassociativity;
interchange of the two coordinates proves cocommutativity. The
$k=0$ maps give the counit. Hence
\[
 (\ell\star\ell')(v)=(\ell\otimes\ell')(\Delta_{k,j}v)
       \quad(\ell\in E_k^\vee,\ \ell'\in E_j^\vee)
 \tag{OMG22}
\]
is an associative commutative graded product with unit in $E_0^\vee$.
For the exact dual basis $\lambda_{k,a,b,D}=i^DD!H_{k,a,b,D}^\vee$,
OMG21 gives
\[
 \lambda_{k,a,b,D}\star\lambda_{j,a',b',D'}
       =\lambda_{k+j,a+a',b+b',D+D'}.
\]
The coefficient is exactly
$\binom{D+D'}D i^DD!\,i^{D'}D'!=i^{D+D'}(D+D')!$;
this records the cancellation before stating the basis product.
Equivalently, OMG20 sends $e^{wY_{k+j}}[1]$ to
$e^{wY_k}[1]\otimes e^{wY_j}[1]$, with
$c_{k+j}=c_k+c_j$. Evaluating OMG18 on this identity proves
\[
 \boxed{\Theta:\ R\xrightarrow{\ \sim\ }
                    \bigoplus_{k\ge0}(E_k^\vee,\star)
                  \text{ is an isomorphism of graded algebras}.}
 \tag{OMG23}
\]
Its source multiplication is the actual polynomial multiplication;
its target multiplication has been constructed from the original sum
coalgebra rather than assigning a multiplication between unrelated
finite-dual vector spaces.

\subsection{The full semisimple and nilpotent action in this geometry}
On $\mathscr P_Z$ define the degree-preserving derivation
\[
 \partial_Z Z_{a,b,d}=\omega_{2a-1,2b-1}Z_{a,b,d}
                   +\begin{cases}d Z_{a,b,d-1},&d>0,\\0,&d=0.
                     \end{cases}
 \tag{OMG24}
\]
On the parameter ring OMG7 it is induced by
$\partial X_a=-i(2a-1)\delta X_a$,
$\partial Y_b=(2b-1)\gamma Y_b$,
$\partial T_0=0$, $\partial T_1=T_0$.
The product rule verifies OMG24 and commutation with $\Phi_m$;
therefore the complete ideal $\mathfrak t_m$ is stable.
The corresponding exact flow on the parameters is
\[
 (X_a,Y_b,T_0,T_1)\longmapsto
 (e^{-i(2a-1)\delta w}X_a,e^{(2b-1)\gamma w}Y_b,T_0,T_1+wT_0).
\]
At $(X_0,X_1,Y_0,Y_1,T_0,T_1)=(1,1,1,1,1,0)$ it gives precisely
the original functions OMG5. This is the entire stated analytic
flow on the explicit algebraic variety, with all polynomial degrees.

If $B_Z$ is the matrix in OMG24 on the column of variables, its
rows have diagonal $\omega_\alpha$ and entry $d$ in column
$(\alpha,d-1)$ of row $(\alpha,d)$. In actual coordinates the
derivation is
\[
 \partial_x x=\mathcal A x,\qquad
 \mathcal A=\mathsf H_mB_Z\mathsf H_m^{-1}.
 \tag{OMG25}
\]
The inverse substitution proves $\partial_x\mathfrak a\subseteq
\mathfrak a$, and its flow differentiates the original $x(w)$.
On the exact degree-$k$ basis of OMG9/OMG18 the induced action is
\[
 \partial_x\lambda_{k,a,b,D}
   =\omega_{k,a,b}\lambda_{k,a,b,D}
             +\begin{cases}D\lambda_{k,a,b,D-1},&D>0,\\0,&D=0.
              \end{cases}\qquad
 \Theta_k\partial_x=Y_k^\vee\Theta_k.
 \tag{OMG26}
\]
The displayed action on the basis uses its identification OMG23.
For a direct verification, on the original primal primary basis
$Y_kH_D=\omega H_D+i^{-1}H_{D+1}$ with the last term zero at
$D=e_k-1$. Applying $i^DD!H_D^\vee$ proves the exact lowering
coefficient $D$ in its dual action. Each frequency block has all
$e_k$ terms, with nonzero top iterated derivative
$(e_k-1)!$. Thus every original primary nilpotent has exactly the
same index $e_k$ in this model. The full annihilating polynomial
is retained with its leading phase:
\[
 \chi_k(c_k+i\partial_x)=0\quad\text{on }R_k,\qquad
 \chi_k(c_k+iT)=i^{q_k}\prod_{a,b=0}^k(T-\omega_{k,a,b})^{e_k}.
 \tag{OMG27}
\]
Distinct frequencies and the nonzero top chain prove that this
degree $q_k$ is minimal. The original $S$ action corresponds to
$c_k+i\partial_x$ on degree $k$, or to
$\tfrac12\mathsf E+i\partial_x$ on the graded algebra, where
$\mathsf E$ is the degree derivation. These equalities specify
every constant, sign and nilpotent action under OMG23.

\subsection{The exact original cyclic observation as invariant restriction}
Let $D=\operatorname{diag}(\zeta,\ldots,\zeta^n)$ and let
$G=\mathbb Z/N$ act on $\mathscr P_x$ by
$(g f)(x)=f(D^{-1}x)$. Thus a degree-$k$ monomial $x^\nu$
has character $\zeta^{-\sum t\nu_t}$. Write
\[
 \mathscr I=\mathscr P_x^G,\qquad
 \mathscr J=\operatorname{im}(\mathscr I\longrightarrow R).
 \tag{OMG28}
\]
These are graded algebras. The coefficient image of the original
observation is exactly their restriction at every degree, as follows.
Let $f_t$ be column $t$ of $F$ and let $S_\nu$ be the sum of
all distinct ordered words in the $f_t$ having multiplicity vector
$\nu$. Disjoint word supports make the $S_\nu$ a basis of the
permutation-invariant tensor space. The expansion and projector are
\[
 (Fx)^{\otimes k}=\sum_{|\nu|=k}x^\nu S_\nu,\qquad
 P_k(Fx)^{\otimes k}=\sum_{\substack{|\nu|=k\\\sum t\nu_t=0\pmod N}}
                    x^\nu S_\nu,\quad
 P_k=\frac1N\sum_{a=0}^{N-1}(C^{-a})^{\otimes k}.
\]
Every ordered word occurs once; the factor $k!/\prod\nu_t!$ is
the number of terms in $S_\nu$, not a missing coefficient.
The original sum injection is the iterated map OMG20 into
$E_h^{\otimes k}$. Hence the original observation
$L_k=P_k\Pi^{\otimes k}U^{\otimes k}\iota_k=\jmath\Lambda_k$
sends $e^{wY_k}[1]$ to the displayed projected curve with $x=x(w)$.
The coefficient covectors of its surviving $S_\nu$ pull back to
exactly $\Theta_k(x^\nu)$. They span the full dual of the image,
since every covector on a subspace of a finite dimensional space
extends to its ambient space. Consequently
\[
\boxed{\begin{gathered}
 \Theta_k(\mathscr J_k)=\operatorname{im}\Lambda_k^\vee
                          =(\ker\Lambda_k)^\perp,\qquad
 \operatorname{rank}\Lambda_k=\dim\mathscr J_k,\\
 R_k/\mathscr J_k\xrightarrow{\ \Theta_k|_{K_k}\ }K_k^\vee
       \text{ is an isomorphism of vector spaces},\qquad
 K_k=\ker\Lambda_k.
\end{gathered}}\tag{OMG29}
\]
For the last statement, a covector on $E_k$ vanishes on $K_k$
exactly when it factors through $E_k/K_k\simeq\operatorname{im}
\Lambda_k$. Restriction to $K_k$ is onto by extension of a basis.
This proves the claimed quotient and its precise type: the
degree-$k$ quotient is a vector-space quotient by $\mathscr J_k$.

\subsection{The actual orbit ideal, including every stabilizer}
The whole ideal $\mathfrak a$ is generated by the explicit quadratic
space in OMG15. Let $B_2$ be its coefficient matrix in the ordered
degree-two monomial frame of $\mathscr P_x$. Repeated or dependent
generator rows are harmless. For $a\in\mathbb Z/N$ let
$D_{2,a}$ be diagonal with entry $\zeta^{-a\sum t\nu_t}$ at
monomial $x^\nu$. The actual stabilizer and its exact finite test are
\[
 H=\{a:g^a\mathfrak a=\mathfrak a\}
   =\left\{a:\operatorname{rank}\begin{bmatrix}B_2\\B_2D_{2,a}\end{bmatrix}
                      =\operatorname{rank}B_2\right\}.
 \tag{OMG30}
\]
Equality of these row spaces is equivalent to equality of their
quadratically generated ideals; their dimensions agree under the
invertible diagonal map. This proves the finite test using the
actual entries OMG4. It prescribes no generic coefficient value.
The set $H$ is a subgroup by composition and inverse of ideal
equalities. Put $s=|H|$ and $d=N/s$. Since $G$ is cyclic, its
cosets have representatives $0,\ldots,d-1$ and
$H=\{0,d,2d,\ldots,(s-1)d\}$. Define all distinct component ideals
and the union algebra by
\[
 \mathfrak a_a=g^a\mathfrak a\ (0\le a<d),\qquad
 \mathfrak u=\bigcap_{a=0}^{d-1}\mathfrak a_a,
 \qquad A=\mathscr P_x/\mathfrak u.
 \tag{OMG31}
\]
These are distinct homogeneous prime ideals with the same Hilbert
function OMG13. Their intersection is a homogeneous radical ideal
stable under the original cyclic action. It is the complete ideal
of their scheme-theoretic union by the definition of the union's
ideal. No product formula for this intersection is substituted.

For any invariant polynomial, membership in one component ideal
implies membership in all its translates, by applying every $g^a$.
Thus the exact invariant restriction and its inverse description are
\[
\boxed{\begin{gathered}
 \mathscr I\cap\mathfrak a=\mathscr I\cap\mathfrak u,\\
 \mathscr I/(\mathscr I\cap\mathfrak u)
      \xrightarrow{\ [f]\mapsto[f]_{\mathfrak u}\ } A^G
      \xrightarrow{\operatorname{res}_0} \mathscr J\subseteq R
      \quad\text{are isomorphisms of graded algebras}.
\end{gathered}}\tag{OMG32}
\]
To prove surjectivity of the first map, average any representative
of an invariant class by $N^{-1}\sum g^a$; its class is unchanged.
Its kernel is exactly the displayed intersection. For the second,
use an invariant representative: a zero restriction lies in
$\mathscr I\cap\mathfrak a=\mathscr I\cap\mathfrak u$ and hence
is already zero in $A^G$. Its image is $\mathscr J$ by definition.
This gives the invariant image for every actual $H$, including
nontrivial stabilizers of a composite group order $4m+1$.

\subsection{The complete component conductor and an explicit finite witness}
Set $R_a=\mathscr P_x/\mathfrak a_a$. The simultaneous restriction
\[
 A\hookrightarrow\widetilde A:=\prod_{a=0}^{d-1}R_a,
 \qquad[f]_{\mathfrak u}\longmapsto([f]_{\mathfrak a_a})_a
 \tag{OMG33}
\]
is injective precisely because its kernel before quotienting is the
intersection OMG31. Its component conductor ideals are
\[
 \mathfrak c_a=
 \frac{(\bigcap_{b\ne a}\mathfrak a_b)+\mathfrak a_a}{\mathfrak a_a}
       \subseteq R_a,
 \qquad\bigcap_{b\ne a}\mathfrak a_b=\mathscr P_x\quad(d=1).
 \tag{OMG34}
\]
The conductor of this specified inclusion is exactly
\[
 \boxed{\{t\in\widetilde A:t\widetilde A\subseteq A\}
                   =\prod_{a=0}^{d-1}\mathfrak c_a\subseteq A.}
 \tag{OMG35}
\]
Indeed a tuple supported solely on component $a$ lies in the image
of $A$ if and only if its value belongs to $\mathfrak c_a$, by
lifting its other zero values to the intersection of the other
ideals. If $t$ is in the conductor, multiplication by each component
idempotent of $\widetilde A$ gives such a supported tuple; thus each
component of $t$ lies in $\mathfrak c_a$. Conversely these are
ideals of $R_a$. Multiplying a tuple with all its components in
the indicated ideals by any element of $\widetilde A$ gives a sum
of such supported tuples, each lying in $A$. This proves both
inclusions and the full formula. In particular this conductor is
also the annihilator of $\widetilde A/A$ in $A$.

The action of $G$ permutes these components with its original
coefficient substitutions. Projection of an invariant tuple to
component zero gives an isomorphism
\[
 \widetilde A^G\xrightarrow{\ \sim\ }R^H,\qquad
 r\in R^H\longmapsto(g^a r)_{0\le a<d}
                       \text{ for its inverse}.
 \tag{OMG36}
\]
The inverse is independent of coset representatives because $r$
is fixed by $H$, and invariance follows by reindexing the components.
The same argument applies to the conductor, whose ideals are
permuted. Equations OMG32--36 consequently prove the exact
inclusions of graded subalgebras and conductor ideals
\[
 \boxed{\mathfrak c_0^H\subseteq\mathscr J\subseteq R^H,
            \qquad\mathfrak c_0^H\,R^H\subseteq\mathscr J.}
 \tag{OMG37}
\]
Here $\mathfrak c_0^H$ denotes invariant elements of the specified
component ideal and is an ideal of $R^H$. An invariant tuple in
the conductor already lies in $A$, so its projection belongs to
$\mathscr J$ by OMG32. This proves the first inclusion; the
ideal property gives the second.

There is a nonzero homogeneous element of this conductor with a
completely bounded degree. For every $a\ne0$ choose from the
explicit transformed quadratic list of $\mathfrak a_a$ a quadratic
$f_a\notin\mathfrak a_0$. Such a choice exists: containment of
the entire quadratic space in that of $\mathfrak a_0$, which has
the same dimension, would make the spaces equal and their
generated ideals equal, contrary to distinctness. Membership is
the finite row-space test in the same degree-two monomial frame.
The product $f=\prod_{a=1}^{d-1}f_a$ lies in all other component
ideals. It is nonzero modulo $\mathfrak a_0$ because that ideal is
prime and none of its factors belongs to it. Thus
$\bar f\in\mathfrak c_0$ is nonzero of degree $2(d-1)$.
The group $H$ preserves $\mathfrak c_0$ and acts on the domain
$R$. Its exact norm
\[
 c_H=\prod_{h\in H}g^h\bar f\in\mathfrak c_0^H,\qquad
 \deg c_H=D_H=2s(d-1)=2(N-s),\qquad c_H\ne0
 \tag{OMG38}
\]
is invariant by permutation of its factors and nonzero in the
domain. For $d=1$ take the empty product $c_H=1$, $D_H=0$;
all statements remain exact. Multiplication by this actual
conductor element is injective and gives the finite dimension
bounds, with negative graded pieces set to zero,
\[
\boxed{\begin{gathered}
 (R^H)_{k-D_H}\xrightarrow{\ \times c_H\ }\mathscr J_k
                    \hookrightarrow(R^H)_k,\\
 \dim(R^H)_{k-D_H}\le\operatorname{rank}\Lambda_k
                                  \le\dim(R^H)_k,\\
 q_k-\dim(R^H)_k\le\dim K_k
                                  \le q_k-\dim(R^H)_{k-D_H}.
\end{gathered}}\tag{OMG39}
\]
Every ideal, group and matrix in these expressions is defined by
the actual original coefficient map. In particular the stabilizer
has its finite equality test OMG30 throughout; no particular value
of it has been inserted into the dimensions. To compute the displayed
invariant dimensions directly, let the actual $H$ substitutions act
on the basis OMG9 through $Z\mapsto\mathsf H_m^{-1}D^{-h}\mathsf H_mZ$
and reduce with OMG10. This is an explicit finite matrix on $q_k$
coordinates. The idempotent average $s^{-1}\sum_{h\in H}g^h$ has
trace equal to its rank, proving
$\dim(R^H)_k=s^{-1}\sum_{h\in H}\operatorname{Tr}(g^h|R_k)$.
Thus the dimensions in OMG39 also have complete original finite
coefficient formulas for every actual stabilizer.

\subsection{The exact action defect and all subsequently observed directions}
The derivation $\partial_x$ of OMG25 acts on $R$ with all primary
directions OMG26. The possible failure of its restriction to the
observed image is the specified finite linear map
\[
 \mathscr J_k\longrightarrow R_k/\mathscr J_k,
                       \qquad j\longmapsto[\partial_xj].
 \tag{OMG40}
\]
Via OMG29 it is exactly the complex linear dual of the original
map $\Lambda_kY_kI_K:K_k\to B_k$. Indeed a $j$ corresponds to
the covector $\lambda\Lambda_k$ on $E_k$ for a unique boundary
covector $\lambda$; evaluating OMG26 on $I_Kv$ gives
$\lambda\Lambda_kY_kI_Kv$, proving the identification and all
domains and codomains. The action on ambient polynomials also has
the exact conjugation formula
\[
 g^a\partial_xg^{-a}=\partial_{D^a\mathcal A D^{-a}},\qquad
 g^a\partial_xg^{-a}-\partial_x
                   =\partial_{D^a\mathcal A D^{-a}-\mathcal A}.
 \tag{OMG41}
\]
For a linear vector field $\partial_M f=(Mx)\cdot\nabla f$,
the chain rule proves this by substituting $D^{-a}x$ after
differentiating $f(D^a x)$. This formula retains every original
period/unit coefficient in $\mathcal A$ and does not require the
original ideal or the observed image to be preserved by extra
ambient actions.

Finally the subspace invisible to all original iterates has the
complete model
\[
\boxed{\begin{gathered}
 \mathcal K_{\infty,k}=\bigcap_{r=0}^{q_k-1}\ker(\Lambda_kY_k^r),\\
 R_k\Big/\sum_{r=0}^{q_k-1}\partial_x^r\mathscr J_k
       \xrightarrow{\ \Theta_k|_{\mathcal K_{\infty,k}}\ }
                          \mathcal K_{\infty,k}^\vee
                          \quad\text{is a vector-space isomorphism}.
\end{gathered}}\tag{OMG42}
\]
In finite dimensional duality the annihilator of an intersection
of kernels is the sum of the images of their dual maps: this
follows by taking the dual of the single stacked map with all
those components, whose kernel is precisely the intersection.
Equations OMG26 and OMG29 identify these images with
$\partial_x^r\mathscr J_k$, proving OMG42. The exact original
polynomial OMG27 expresses every higher iterate as a linear
combination of the first $q_k$, so the same intersection is taken
over all nonnegative iterates and is $Y_k$-invariant. This proves
the full nilpotent and cyclic-action receiver in the new homogeneous
model for every original multiplicity, with its actual observation,
periods and amplitude jets throughout.

% END CURRENT OMG

\clearpage
% BEGIN CURRENT AKS
% AKS: new continuation, 14 September 2026.
% This body is self-contained at the level of its explicitly stated inherited inputs.
\section{The absolute common kernel: the new finite result}
\label{sec:AKS}

This complete incoming calculation sharpens the current absolute common
kernel estimate. Its original comparison was against the preceding
logarithmic bound; the intervening KAF chapter already established
the $O(kq)$ scale. The new finite envelope reduces its uniform limsup
constant from $16(4+10\log2)$ to $32\log(25\log3/(4\pi))$.
AKJ below takes the pointwise minimum of both finite bounds. No
limiting kernel coefficient is inferred from these upper bounds.

Retain the actual hypothetical complete quartet and its entire quotient
\[
 \rho=\tfrac12+\delta+i\gamma,\qquad 0<\delta<\tfrac12,\quad \gamma>2,
 \qquad h(z)=\prod_{\alpha\in\{\rho,\bar\rho,1-\rho,1-\bar\rho\}}
 (z-\alpha)^m,\quad v_h=2\xi/h.
\]
The full primary unit is $U=M_{j_hv_h}$. The original tensor order, centre,
polynomial, and coefficient quotient are
\[
\begin{gathered}
 k=4l+1,\quad l\ge1,\quad e=1+k(m-1),\quad q=e(k+1)^2,\quad c=k/2,\\
 \lambda_{a,b}=c+(2a-k)\delta+i(2b-k)\gamma,\qquad
 \chi(S)=\prod_{a,b=0}^k(S-\lambda_{a,b})^e,
 \quad E=\mathbb C[S]/(\chi).
\end{gathered}\tag{AKS1}
\]
For the principal conclusion $m=1$. The source-operator estimates below
also hold with the stated repeated-root list for $m>1$, but no simple-quartet
rank is assigned to those multiplicities.

Let $\Lambda_k(u):E\twoheadrightarrow B_k(u)$ be the original observation,
$K(u)=\ker\Lambda_k(u)$, and $I:K(u)\hookrightarrow E$ its original
coefficient inclusion. For a source $s\in\{1,k\}$, let $G_N^{(s)}$ be the
original minimum-remainder Gram, in the frame $(1,S,\ldots,S^{q-1})$.
Define $H_{K,N}^{(s)}=I^*G_N^{(s)}I$ and retain exactly
\[
 F_K^{(s)}=
 \log\det H_{K,q-1}^{(s)}+\log\det H_{K,q}^{(s)}
 -\log\det H_{K,2q-1}^{(s)}-\log\det H_{K,2q}^{(s)}.
 \tag{AKS2}
\]
Every empty determinant is one. All conjugate transposes refer to the
specified coordinate matrices, not to a bilinear residue pairing.

Here are entirely explicit positive numbers used in the bound:
\[
\begin{gathered}
 R=k\sqrt{\delta^2+\gamma^2},\quad b_s=s/2,\quad
 t_q=\frac\pi2-\frac1q,\quad \mathfrak a_q=\frac{2\log3}{t_q},\quad
 \mathfrak s_q=\frac{\mathfrak a_q^q-1}{\mathfrak a_q-1},\\
 \mathfrak b_{s,q}=\left[\frac9{25}\sin\frac1q\right]^{-b_s}
       \exp\left(2R\left(\log3+\frac{t_q}{2}\right)\right),\\
 \mathcal Z_{s,j}=1+\mathfrak b_{s,q}\mathfrak s_q^2
        \sum_{d=q}^{q+j-1}\frac{d!}{(b_s)_d}
                          \left(\frac{25}{16}\right)^d,
 \quad \mathcal L_{s,j}=\log \mathcal Z_{s,j}\quad (0\le j\le q+1).
\end{gathered}\tag{AKS3}
\]
The empty sum gives $\mathcal Z_{s,0}=1$. In the admitted family $q\ge36$, so
$t_q\in(0,\pi/2)$ and $\mathfrak a_q>1$. These parameters are radii in the existing
Laplace and generating functions; neither is a new source order.

\paragraph{Main theorem (absolute finite bound).}
For the original simple quartet, every original logarithm branch, and
every $|u|\ge R_*$, where $R_*$ is precisely RPC21 or RPC23 (recorded in
Section~\ref{sec:AKS-radius}), put $r=8k-16$. Then
\[
 \boxed{\quad
 0\le F_K^{(s)}(k,u)\le r\bigl(\mathcal L_{s,q}+\mathcal L_{s,q+1}\bigr),
 \qquad s\in\{1,k\}.\quad}
 \tag{AKS4}
\]
There is a stronger, fully period-dependent interval in AKS28 below.
In particular, with
\[
 \kappa_* =32\log\left(\frac{25\log3}{4\pi}\right),
 \qquad 25.02078097649<\kappa_*<25.02078097650,
 \tag{AKS5}
\]
one has the uniform statement
\[
 \boxed{\quad
 0\le\limsup_{k\to\infty\atop k\equiv1\ (4)}
       \ \sup_{\substack{|u|\ge R_*\\\text{original branches}}}
         \frac{F_K^{(s)}(k,u)}{kq}\le\kappa_*.
 \quad}\tag{AKS6}
\]
Here the source may be $s=1$ for every $k$, or $s=k$. The same bound holds
for the actual arithmetic kernel. Its exact finite allowances are in
AKS34--36. Thus the absolute kernel is $O_h(kq)$ on this domain, not only
$O_h(kq\log q)$. Neither AKS5 nor AKS6 calls $\kappa_*$ the actual leading
coefficient.

\section{An exact coefficient realization of the actual lost directions}
\label{sec:AKS-frame}

\subsection{The period matrix and invariant restriction, with their full coefficients}
For $m=1$, order the one-factor roots as $(++),(+-),(-+),(--)$, and let
$\mathcal E_{\rm prim}$ have their original Lagrange idempotents as columns.
Set $\zeta=\exp(2\pi i/5)$ and $F_{jt}=\zeta^{jt}-1$, $1\le j,t\le4$.
The exact inverse is $(F^{-1})_{tj}=\zeta^{-tj}/5$. Retain
\[
\begin{gathered}
 \Pi[P]_j=\int_{\Gamma_j}e^{\Phi_h(z)/u}\operatorname{rem}_hP(z)\,dz,
 \qquad\Phi_h'=h,\quad\Phi_h(0)=0,\\
 \Gamma_j=-\ell_0+\ell_j,\quad
 \arg\ell_j=(\arg u+\pi+2\pi j)/5,\\
 \mathsf H=F^{-1}\Pi U\mathcal E_{\rm prim},\quad
 D=\operatorname{diag}(\zeta,\zeta^2,\zeta^3,\zeta^4),\\
 Q(x)=Q_0(\mathsf H^{-1}x),\qquad
 Q_0(t)=t_1t_4-t_2t_3,\quad Q_a(x)=Q(D^{-a}x).
\end{gathered}\tag{AKS7}
\]
The integral always uses the degree-less-than-four representative.
Invertibility of $\Pi$, all contour phases, and its full determinant are
the PDE/PQO results in the supplied proofs. The four unit values are
$(a,\bar a,\bar a,a)$, with the actual $a=v_h(\rho)\ne0$; no value of $a$
is selected. Hence $\mathsf H$ is the original invertible matrix.

The following formulas make the coefficient realization directly
computable. Write
\[
 \mathcal I_{k,0}=\{\nu\in\mathbb Z_{\ge0}^4:|\nu|=k,
                         \ \nu_1+2\nu_2+3\nu_3+4\nu_4\equiv0\pmod5\}.
\]
For $\nu\in\mathcal I_{k,0}$ define $B_{\nu;(a,b)}$ by summing over
$4\times4$ arrays $n_{t,\alpha}\ge0$ with row sums $\nu_t$, total plus-real
count $a$, and total plus-imaginary count $b$:
\[
 B_{\nu;(a,b)}=
 \sum_{\substack{\sum_\alpha n_{t,\alpha}=\nu_t\\
       \sum_{t,\alpha:\epsilon_\alpha=+}n_{t,\alpha}=a\\
       \sum_{t,\alpha:\eta_\alpha=+}n_{t,\alpha}=b}}
 \left(\prod_{t=1}^4\frac{\nu_t!}{\prod_\alpha n_{t,\alpha}!}\right)
                  \prod_{t,\alpha}\mathsf H_{t,\alpha}^{n_{t,\alpha}}.
 \tag{AKS8}
\]
Let $\mathsf V_{(a,b),d}=\lambda_{a,b}^d$ for $0\le d<q$ and set
$\mathsf A=B\mathsf V$. All entries of $\mathsf A$ are now finite
expressions in the actual periods and unit. There is no generic period
matrix in this definition.

For completeness, put $t_\alpha(w)=\exp((\eta_\alpha\gamma-
 i\epsilon_\alpha\delta)w)$ and $x(w)=\mathsf Ht(w)$. Expansion gives
\[
 x(w)^\nu=\sum_{a,b}B_{\nu;(a,b)}e^{\omega_{a,b}w},\qquad
 \omega_{a,b}=(2b-k)\gamma-i(2a-k)\delta.
\]
In the original primary algebra,
$e^{wY}[1]=\sum_{a,b}e^{\omega_{a,b}w}\varepsilon_{a,b}$, where
$Y=(M_S-cI)/i$. The exponentials are independent by the Vandermonde of
the distinct $\omega_{a,b}$. OCS9's unscaled symmetric word basis
therefore identifies the original observation with
\[
       \jmath\Lambda_k=\mathcal F_0\mathsf A,
       \qquad\ker\mathsf A=K(u)
       \quad\text{in the original coefficient frame.}
       \tag{AKS9}
\]
Every multinomial in AKS8 belongs to its actual ordered words; the
word multiplicity is not inserted a second time. The exact inherited
row Gram of $\mathcal F_0$ is the restriction to $\mathcal I_{k,0}$ of
\[
 \mathcal G_{\nu,\mu}=
 \sum_{\substack{n_{tu}\ge0\\\sum_u n_{tu}=\nu_t\\
                              \sum_t n_{tu}=\mu_u}}
 \frac{k!}{\prod_{t,u}n_{tu}!}\prod_{t,u}
          [5(\delta_{tu}+1)]^{n_{tu}}.
 \tag{AKS10}
\]
This follows by taking the coefficient of $\bar x^\nu y^\mu$ in
$(\bar x^TF^*Fy)^k$, since $F^*F=5(I+\mathbf1\mathbf1^*)$. It is positive
definite and includes every row cross term.

Here is the precise role of the orbit ideal. The map sending a homogeneous
polynomial $f$ to the covector specified by
$\Theta_k(f)(e^{wY}[1])=f(x(w))$ has kernel
$Q\mathbb C[x]_{k-2}$ and is onto $E^\vee$. Indeed equal exponential
indices correspond exactly to monomial differences divisible by
$t_1t_4-t_2t_3$; the displayed Vandermonde proves that there are no other
relations. Its invariant image is $\operatorname{im}\Lambda_k^\vee$.
On the five-member domain, the five $Q_a$ are distinct prime quadrics.
An invariant polynomial divisible by $Q$ is divisible by each $Q_a$,
and therefore by their full product; its quotient is invariant by
cancellation in the polynomial ring. Consequently
\[
\begin{gathered}
 \mathbb C[x]^{\langle D\rangle}\cap (Q)
       =\left(\prod_{a=0}^4Q_a\right)\mathbb C[x]^{\langle D\rangle},\\
 K(u)^\vee\simeq
 \mathbb C[x]_k\big/(\mathbb C[x]_k^{\langle D\rangle}
                              +Q\mathbb C[x]_{k-2}),\\
 b:=\operatorname{rank}\mathsf A
     =d_{5,k,0}-d_{5,k-10,0}=k^2-6k+17,\qquad r=q-b=8k-16.
\end{gathered}\tag{AKS11}
\]
Negative degrees are zero. The count is
$d_{5,d,0}=(\binom{d+3}{3}+4\epsilon_d)/5$,
$\epsilon_d=\mathbf1_{d\equiv0(5)}-\mathbf1_{d\equiv1(5)}$;
subtracting the cubics gives the stated expression for $k\ge10$, and
$k=5,9$ give $12,44$ directly. This is the inherited OQX calculation,
repeated to specify which kernel AKS9 realizes. The source metric is
still to be applied to this kernel, not to the invariant ring itself.

\subsection{A determinant formula requiring no assumed nonzero pivot}
Let $\mathcal G_0$ be the full positive row Gram in AKS10 and set
\[
 W=\mathsf A^*\mathcal G_0\mathsf A,\qquad
 \sigma_j=[z^j]\det(I_q+zW),\qquad
 P_K=\frac1{\sigma_b}\sum_{j=0}^b(-1)^j\sigma_{b-j}W^j.
 \tag{AKS12}
\]
Since $W$ has exactly $b$ positive eigenvalues, $\sigma_b>0$. The scalar
polynomial in AKS12 is one at zero and zero at every nonzero eigenvalue:
it is the normalized product of the factors $\lambda_i-z$ for the $b$
positive eigenvalues. Hence $P_K=P_K^*=P_K^2$ and
$\operatorname{im}P_K=K(u)$. This is the projector derived from the
literal observation, not a freely chosen projection. Replacing
$\mathcal G_0$ by any other positive coordinate Gram leaves this same
coefficient-orthogonal kernel projector, since its image and orthogonal
complement are unchanged. No such replacement is needed here.

At every original endpoint define
\[
 D_{K,N}^{(s)}=[z^r]\det(I_q+zP_KG_N^{(s)}P_K).
 \tag{AKS13}
\]
These numbers are strictly positive for $r>0$ and are one for $r=0$.
To prove their exact metric meaning, choose an orthonormal coefficient
basis $U_K$ of $\operatorname{im}P_K$ just for the proof. In a unitary
extension, $P_KG_NP_K$ has its positive block $U_K^*G_NU_K$ and a zero
block. Thus $D_{K,N}=\det(U_K^*G_NU_K)$. Write the original inclusion
as $I=U_KT$; then $|\det T|^2=\det(I^*I)$ and
\[
 \det H_{K,N}^{(s)}=\det(I^*I)D_{K,N}^{(s)},\qquad
 \boxed{\ F_K^{(s)}=
 \log\frac{D_{K,q-1}^{(s)}D_{K,q}^{(s)}}
               {D_{K,2q-1}^{(s)}D_{K,2q}^{(s)}}.\ }
 \tag{AKS14}
\]
This supplies an equivalent exact coefficient determinant formula for
all the $8k-16$ lost directions, including the entire period dependence.
The original frame factor is present at each endpoint and cancels only
with its four actual signs. It is unnecessary to suppose any particular
minor nonzero. Alternatively, selecting the first independent $r$ columns
of $P_K$ gives an explicit coefficient frame; the determinant formula
avoids even that chart choice.

For $m>1$ replace AKS8--9 by OCS12--18's full confluent matrix
$\mathfrak A_{\nu;(a,b,D)}=i^DD!c_{\nu;a,b,D}$ and its primary isomorphism.
AKS12--14 are then identical with the actual rank, not with the rank in
AKS11. All primary nilpotents and all unit derivatives remain in that
matrix.

\section{A bounded exact remainder operator on the original source}
\label{sec:AKS-operator}

\subsection{The complete source norms and coefficient map}
For either existing source set $M_s=(2\pi)^{s/2}$ and retain
\[
 dm_{0,s}(y)=\frac{M_s2^{b_s}}{4\pi\Gamma(b_s)}
           |\Gamma(b_s/2+iy/2)|^2\,dy,\qquad
 \int e^{ty}dm_{0,s}=M_s(\cos t)^{-b_s},\quad |t|<\pi/2.
 \tag{AKS15}
\]
The supplied IGO1--3 proof derives this transform by the beta integral,
Fourier inversion and a contour shift within $|\operatorname{Im}x|<\pi/2$.
In particular it includes the exact masses $M_1=\sqrt{2\pi}$ and
$M_k=(2\pi)^{k/2}$ and exponential domination on every smaller strip.
The original real monic polynomials satisfy
\[
\begin{gathered}
 p_0=1,\quad p_1=y,\quad
 p_{d+1}=yp_d-d(b_s+d-1)p_{d-1},\\
 \sum_{d\ge0}\frac{p_d(y)}{d!}z^d
      =(1+z^2)^{-b_s/2}e^{y\arctan z},\qquad
 h_d^{(s)}=\|p_d\|^2=M_s d!(b_s)_d.
\end{gathered}\tag{AKS16}
\]
The generating formula follows by differentiating in $z$. To check the
normalization directly, integrate the product of two generating
functions at small real $z,w$. AKS15 and the cosine addition identity
give $M_s(1-zw)^{-b_s}$. Taking any finite pair of derivatives at zero,
justified by the stated exponential domination, gives orthogonality and
the norms in AKS16. Thus the leading phase in
$u_d(S)=p_d((S-c)/i)/\sqrt{h_d^{(s)}}$ is exactly $i^{-d}$.
The same generating function agrees with NIST DLMF 18.23.7 after the
original parameter rescaling; it is derived here from the supplied
recurrence, not substituted for it.

List the original roots
$\omega_\nu=(\lambda_{a,b}-c)/i$ in lexicographic order, repeating each
exactly $e$ times, and put
\[
 P(y)=i^{-q}\chi(c+iy)=\prod_{\nu=1}^q(y-\omega_\nu),\qquad
 \Phi:[f(y)]\mapsto[f((S-c)/i)].
 \tag{AKS17}
\]
The inverse is substitution $S=c+iy$, and
$P((S-c)/i)=i^{-q}\chi(S)$. This is an exact algebra isomorphism, including
the complete confluent orders; there is no monomial comparison divisor.
Let $C_s$ have columns $[u_0],\ldots,[u_{q-1}]$. Its nonzero determinant
is $i^{-q(q-1)/2}\prod_{d<q}(h_d^{(s)})^{-1/2}$. Define
\[
 v_j=C_s^{-1}[u_{q+j}],\qquad
 \mathsf K_j=I_q+\sum_{a=0}^{j-1}v_av_a^*,\qquad 0\le j\le q+1.
 \tag{AKS18}
\]
Parseval in the low-degree original orthonormal polynomial basis proves
\[
 \|v_j\|_2^2=
 \frac{\|\operatorname{rem}_P p_{q+j}\|_{m_{0,s}}^2}
      {M_s(q+j)!(b_s)_{q+j}}.
 \tag{AKS19}
\]
This is the full norm of the actual remainder, not a diagonal entry of
an approximated quotient Gram.

\subsection{Newton interpolation with all roots and their orders}
For $0\le a<q$ define $N_a(y)=\prod_{\nu=1}^a(y-\omega_\nu)$, with
$N_0=1$. For every complex $v$, the unique entire-function remainder is
\[
 \operatorname{rem}_P(e^{vy})
   =\sum_{a=0}^{q-1}[\omega_1,\ldots,\omega_{a+1}]e^{v\cdot}\ N_a(y),
 \qquad
 \left|[\omega_1,\ldots,\omega_{a+1}]e^{v\cdot}\right|
       \le\frac{|v|^a}{a!}e^{R|v|}.
 \tag{AKS20}
\]
Here divided differences at repeated nodes mean Hermite divided
differences, not evaluation with a zero denominator.

A proof covering the confluent case is useful. On the standard
$a$-simplex with barycentric coordinates $\theta_1,\ldots,\theta_{a+1}$,
\[
 [\omega_1,\ldots,\omega_{a+1}]e^{v\cdot}
       =v^a\int_{\Delta_a}e^{v\sum_\nu\theta_\nu\omega_\nu}\,d\theta.
\]
The integral of $\prod\theta_\nu^{n_\nu}$ is
$\prod n_\nu!/(a+\sum n_\nu)!$, as follows by iterated integration of
integer beta integrals. Expanding the exponential therefore gives
$\sum_{n\ge0}v^{a+n}h_n(\omega_1,\ldots,\omega_{a+1})/(a+n)!$.
For monomials, the divided difference of $y^{a+n}$ is exactly $h_n$:
this follows by multiplying the geometric series
$\prod(1-\omega_\nu z)^{-1}$ and the Newton division recurrence.
It is a polynomial identity in the nodes and so also holds at repeated
nodes. This proves the integral formula. Iterated Newton division for
polynomials, followed by the convergent exponential series, proves the
first equality in AKS20, including every required derivative.
Finally the simplex has volume $1/a!$ and
$|\sum\theta_\nu\omega_\nu|\le R$, proving its stated bound.

For any $0<t<\pi/2$, put $\tau=2/t$. Positivity of every term of the
cosh series in AKS15 gives
\[
 \|y^a\|_{m_{0,s}}
 \le\sqrt{M_s}(\cos t)^{-b_s/2}\tau^a a!.
\]
Indeed its square is at most
$M_s(\cos t)^{-b_s}(2a)!t^{-2a}$ and
$(2a)!\le4^a(a!)^2$. Expanding the \emph{whole} Newton polynomial, then
using $|e_j(\omega_1,\ldots,\omega_a)|\le\binom ajR^j$, proves
\[
\begin{split}
 \|N_a\|_{m_{0,s}}
 &\le\sqrt{M_s}(\cos t)^{-b_s/2}
        \sum_{j=0}^a\binom ajR^{a-j}\tau^j j!\\
 &=\sqrt{M_s}(\cos t)^{-b_s/2}\tau^a a!
           \sum_{n=0}^a\frac{(R/\tau)^n}{n!}\\
 &\le\sqrt{M_s}(\cos t)^{-b_s/2}\tau^a a!e^{R/\tau}.
\end{split}\tag{AKS21}
\]
There are no omitted polynomial coefficients. Combining AKS20--21 gives
for every complex $v$
\[
 \|\operatorname{rem}_P(e^{vy})\|_{m_{0,s}}
 \le\sqrt{M_s}(\cos t)^{-b_s/2}e^{R(|v|+t/2)}
                       \sum_{a=0}^{q-1}(2|v|/t)^a.
 \tag{AKS22}
\]

\subsection{Cauchy extraction in the actual Gamma generating function}
Fix $0<z<1$ and write $a_z=\operatorname{arctanh}z$.
For $|w|=z$ the defining power series gives
$|\arctan w|\le a_z$ and $|1+w^2|\ge1-z^2$.
Apply the literal finite-dimensional remainder map to AKS16 and use
Cauchy's formula for its $d$th coefficient. The norm triangle inequality
on that circle and AKS22 give
\[
\begin{split}
 \frac{\|\operatorname{rem}_P p_d\|_{m_{0,s}}}{d!}
 \le{}&z^{-d}(1-z^2)^{-b_s/2}\sqrt{M_s}(\cos t)^{-b_s/2}\\
 &\mathrel{}\times e^{R(a_z+t/2)}
                    \sum_{a=0}^{q-1}(2a_z/t)^a.
\end{split}\tag{AKS23}
\]
Cauchy's formula is legitimate for the finite remainder vector: its
entries are holomorphic for $|w|<1$ by the actual matrix generating
identity OPG7. The source norm is taken only after this finite-vector
identity; there is no unjustified integral of a nonintegrable generating
function on the large circle.

Substitution of AKS23 into AKS19 retains the factorial denominator and
cancels the two occurrences of the original mass, giving
\[
 \boxed{\ 
 \|v_j\|_2^2\le\frac{d!}{(b_s)_d}z^{-2d}
 [\cos t(1-z^2)]^{-b_s}e^{2R(a_z+t/2)}
 \left[\sum_{a=0}^{q-1}(2a_z/t)^a\right]^2,
 \quad d=q+j.\ }
 \tag{AKS24}
\]
Adding the \emph{complete} rank-one positive matrices in AKS18 now proves
$\mathsf K_j\preceq \mathcal Z_{s,j}I_q$ after setting $z=4/5$ and $t=t_q$.
Indeed $a_{4/5}=\log3$, $1-z^2=9/25$, and $\cos t_q=\sin(1/q)$, so
the exact scalar obtained is AKS3. The high index $j=q+1$ includes the
original polynomial of degree $2q$.

There is also a root-by-root finite refinement. Put
$\mu_{s,d}=\int y^{2d}dm_{0,s}$,
$R_a=\max_{\nu\le a+1}|\omega_\nu|$, and
\[
\begin{gathered}
 B_{s,a}^{\rm root}=\sum_{d=0}^a
     |e_{a-d}(\omega_1,\ldots,\omega_a)|\sqrt{\mu_{s,d}},\\
 J_s^{\rm root}(z)=\frac{(1-z^2)^{-b_s}}{M_s}
 \left[\sum_{a=0}^{q-1}\frac{a_z^a}{a!}e^{a_zR_a}
                      B_{s,a}^{\rm root}\right]^2.
\end{gathered}\tag{AKS25}
\]
Replace the scalar multiplying
$d!z^{-2d}/(b_s)_d$ in AKS24 by $J_s^{\rm root}(z)$ to get a valid
sharper bound, by the same proof before AKS21 is relaxed. Every moment
is exactly $M_s(2d)![w^{2d}](\cos w)^{-b_s}$. Thus this refinement keeps
every actual root coefficient and a finite original moment formula.
AKS21 bounds it by the explicit scalar in AKS24. The proof never replaces
the source determinant by a product of diagonal moments: AKS24 is a
full-vector operator bound, and the next step restricts the inverse of
the full covariance.

\section{The four original kernel Grams and the evaluated scale}
\label{sec:AKS-volume}

The original orthonormal source map at degree $q+j-1$ is
$C_s[I_q,v_0,\ldots,v_{j-1}]$. Its minimum lift of $C_sx$ is
$[I_q,v_0,\ldots,v_{j-1}]^*\mathsf K_j^{-1}x$. It has that image and is
orthogonal to the full relation kernel, so Pythagoras proves
\[
 G_{q+j-1}^{(s)}=C_s^{-*}\mathsf K_j^{-1}C_s^{-1},\qquad
 H_j:=H_{K,q+j-1}^{(s)}=T_s^*\mathsf K_j^{-1}T_s,\quad T_s=C_s^{-1}I.
 \tag{AKS26}
\]
This verifies the source-minimum step in the coefficient determinant
AKS13--14. Neither a cross block nor any original relation is removed.

Since $I_q\preceq\mathsf K_j\preceq \mathcal Z_{s,j}I_q$, inversion and compression
give $\mathcal Z_{s,j}^{-1}H_0\preceq H_j\preceq H_0$, where $H_0=T_s^*T_s$.
Define $d_j=-\log\det(H_0^{-1/2}H_jH_0^{-1/2})$. Congruence and all
$r$ eigenvalues give
\[
 0=d_0\le d_1\le\cdots\le d_{q+1},\qquad
 0\le d_j\le r\mathcal L_{s,j},\qquad
 F_K^{(s)}=d_q+d_{q+1}-d_1.
 \tag{AKS27}
\]
The identity uses precisely the four endpoints in AKS2.
In particular it proves AKS4, as well as the sharper interval
\[
 \boxed{\quad d_1\le F_K^{(s)}\le
             r(\mathcal L_{s,q}+\mathcal L_{s,q+1})-d_1.\quad}
 \tag{AKS28}
\]
For a directly calculable expression for its lower endpoint set
$V=\|v_0\|^2$ and
$\theta=v_0^*T_s(T_s^*T_s)^{-1}T_s^*v_0$. Then
\[
 0\le\theta\le V,\qquad
 d_1=\log\frac{1+V}{1+V-\theta}.
 \tag{AKS29}
\]
Indeed the inverse of $I+v_0v_0^*$ is
$I-v_0v_0^*/(1+V)$; take its full compressed determinant. For $r=0$
set $\theta=d_1=0$. In GMB/OPR notation $V=t_0$ and $V-\theta=\xi_0$,
so this is exactly the first original section-corrected kernel increment,
not an independently chosen kernel observation.

Here is an evaluated finite remainder for AKS6. Put
\[
\begin{gathered}
 a_\infty=\frac{4\log3}{\pi},\qquad
 c_0=\log\frac{25\log3}{4\pi}>0,\\
 E_{s,q}(R)=\log2+\frac4{\pi-2/q}
  +\frac{s}{2}\log\frac{25\pi q}{18}
  +2R\left(\log3+\frac\pi4\right)\\
 \hspace{23mm}+\log((q+1)(4q+1))-2\log(a_\infty-1).
\end{gathered}\tag{AKS30}
\]
Then for both source orders, every $q\ge36$, and the same actual kernel,
\[
 \boxed{\quad
 \mathcal L_{s,q},\mathcal L_{s,q+1}\le2c_0q+E_{s,q}(R),\qquad
 0\le F_K^{(s)}\le4rc_0q+2rE_{s,q}(R).
 \quad}\tag{AKS31}
\]
To verify every factor, $b_s\ge1/2$ gives
\[
 \frac{d!}{(b_s)_d}\le\frac{4^d}{\binom{2d}{d}}\le2d+1.
\]
The last inequality uses that the central binomial is the largest of the
$2d+1$ positive coefficients whose sum is $4^d$. Thus the sum in AKS3
is at most $(q+1)(4q+1)(5/4)^{4q}$. Also
$\mathfrak s_q\le \mathfrak a_q^q/(a_\infty-1)$ and
\[
 2q\log \mathfrak a_q\le2q\log a_\infty+\frac4{\pi-2/q},\qquad
 \sin(1/q)\ge\frac2{\pi q}.
\]
The first follows from $-\log(1-x)\le x/(1-x)$ at $x=2/(\pi q)$;
the second is the chord bound for sine on $[0,\pi/2]$.
Use $t_q\le\pi/2$ in the remaining exponential. The order-$q$ terms
combine to
$4q\log(5/4)+2q\log a_\infty=2c_0q$.
The upper bound for the logarithm of the added positive scalar is
positive, so $\log(1+e^x)\le\log2+\max(0,x)$ proves the first bound in
AKS31. AKS28 proves the second.

On the original simple family $q=(k+1)^2$, $R=O_h(k)$ and
$r=8k-16$. Consequently $E_{1,q}(R)=O_h(k+\log q)$ and
$E_{k,q}(R)=O_h(k\log q)$. Dividing AKS31 by $kq$ proves AKS6, with
$32c_0=\kappa_*$. In fact the fixed-source upper bound is
\[
 F_K^{(1)}\le(32k-64)c_0q+(16k-32)E_{1,q}(R),
 \tag{AKS32}
\]
whose explicit second term is $O_h(k^2)$. The tensor-source version has
an $O_h(k^2\log q)$ upper remainder. These are remainders of upper
bounds, not assertions about the actual signed error or leading value.
The proof is uniform in all the period parameters in the stated domain:
no norm or condition number of a period matrix occurs in AKS30.

\paragraph{The precise earlier loss.}
OKV15 had already retained
$\|v_j\|^2\le q^22^{6q}(2D_*q)^{2d}/(d!(b_s)_d)$, $d=q+j$.
Replacing its denominator by $1/2$ caused the old $q\log q$ exponent.
Even without AKS20--24 one obtains the valid repair
\[
 \|\mathsf K_{q+1}\|\le
 1+(q+1)q^22^{6q}\frac{(4D_*q)^{4q}}{(4q)!},\qquad D_*\ge3.
 \tag{AKS33}
\]
Here $d!(b_s)_d\ge(2d)!/4^d$ and the remaining factorial ratio increases
up to $d=2q$ for $D_*\ge3$. The elementary inequality
$\log n!\ge n\log n-n+1$ already makes the logarithm of AKS33 $O_h(q)$.
The Newton calculation strengthens this repair to the packet-independent
limsup constant AKS5, while retaining $R$ in its finite remainder.


\paragraph{Exact enclosure of the displayed numerical ceiling.}
For $x>1$ and $z=(x-1)/(x+1)$ the integrated geometric series gives
\[
 0\le \log x-2\sum_{j=0}^{39}\frac{z^{2j+1}}{2j+1}
      \le\frac{2z^{81}}{81(1-z^2)}.
\]
For $0<x<1$, the alternating sum through $j=39$ for $\arctan x$
has error between zero and $x^{81}/81$. The tangent addition formula,
with the angle in $(0,\pi/2)$, proves
$\pi=16\arctan(1/5)-4\arctan(1/239)$: the tangent of
$4\arctan(1/5)-\arctan(1/239)$ is exactly one. Use the resulting
rational lower and upper bounds for $\pi$ and for $\log3$, then apply
the first displayed bound to the two endpoints of
$25\log3/(4\pi)$. Multiplication by $32$ gives, with outward decimal
rounding, the interval
\[
 \kappa_*\in[25.0207809764994871,\ 25.0207809764994872].
\]
The accompanying rational-arithmetic implementation checks these
inequalities by integer cross multiplication, rather than a floating-point
logarithm. This proves the less finely rounded strict interval AKS5.

\section{The complete arithmetic and mixed receivers}
\label{sec:AKS-transfer}

\subsection{The unchanged finite arithmetic allowances}
Use exactly AMT4's constants. To prevent a loss hidden in order notation,
write them here:
\[
\begin{gathered}
 \alpha=\pi/2,\ p=8m-9/2,\ B_h=42+8m,\ M=21+4m,\\
 \vartheta_h=\int_{-1}^1w_h(y)dy,\quad
 M_\alpha=\int e^{\alpha y}w_h(y)dy,\quad
 w_h(y)=|v_h(1/2+iy)|^2/(2\pi),\\
 c_\Gamma=\sqrt2e^{-7/6},\quad C_\Gamma=\sqrt{2\sqrt5}e^{2/3},\\
 a_k=\frac{c_h\vartheta_h^{k-3}e^{-\alpha(k-3)}(k-2)^{-B_h}}
                    {C_\Gamma2^{B_h/2}},\quad
 A_k=\frac{C_h^{\rm tilt}}{c_\Gamma}k^{p+1}M_\alpha^{k-1},\\
 D_N=[1+4(N+M)^2]^M,\quad L_N^{\rm ar}=\log(A_kD_N/a_k),\\
 E_k^+=L_{q-1}^{\rm ar}+L_q^{\rm ar},\qquad
 E_k^-=L_{2q-1}^{\rm ar}+L_{2q}^{\rm ar}.
\end{gathered}\tag{AKS34}
\]
The constants $c_h,C_h^{\rm tilt}$ are the original TW/BT envelope
constants, not new choices. Their actual inequalities and complete
source construction are preserved in AMT4--6 and the current source
closure. Minimize AMT6 on the same remainder fibre and then restrict to
$I(K)$. AMT9--20 gives, with the original arithmetic source $w_h^{*k}$,
\[
 -rE_k^+\le F_K^{\rm ar}-F_K^{(1)}\le rE_k^-.
\]
Let $\underline U_s=d_1^{(s)}$ and
$\overline U_s=r(\mathcal L_{s,q}+\mathcal L_{s,q+1})-d_1^{(s)}$. The new finite conclusion is
\[
 \boxed{\quad
 \max\{0,\underline U_1-rE_k^+\}\le F_K^{\rm ar}
                  \le\overline U_1+rE_k^-.
 \quad}\tag{AKS35}
\]
This improves AMT37's absolute upper bound. For the sharper correlated
version retain $d_{E,N}=q\log A_k-
\log(\det G_N^{\rm ar}/\det G_N^{(1)})$ and
$c_{X,N}=\min\{r_XL_N^{\rm ar},d_{E,N}\}$, where $r_K=r$, $r_B=q-r$.
Put $C_X^-=c_{X,q-1}+c_{X,q}$,
$C_X^+=c_{X,2q-1}+c_{X,2q}$. AMT18 then gives
\[
\begin{split}
 \max\{0,\underline U_1-C_K^-,
                  \underline U_1+\delta_k^\sigma-C_B^+\}
 &\le F_K^{\rm ar}\\
 &\le\overline U_1+
            \min\{C_K^+,\delta_k^\sigma+C_B^-\}.
\end{split}\tag{AKS36}
\]
One may replace $d_{E,N}$ in the caps by the original full-source
$\Xi_{k,N}$ or its proved finite majorant. The scalar
$\delta_k^\sigma=\mathcal B_k^{\rm ar}-\mathcal B_k^\sigma$ retains
exactly AMT21:
\[
\begin{gathered}
 \max\{-qE_k^+,-\Xi_{k,q-1}-\Xi_{k,q}\}
       \le\delta_k^\sigma\\
       \le\min\{qE_k^-,\Xi_{k,2q-1}+\Xi_{k,2q}\}.
\end{gathered}\tag{AKS37}
\]
Since $E_k^\pm=O_h(k+\log q)$ by the exact AMT22 expression, AKS35 and
AKS32 prove the arithmetic version of AKS6, with an explicit
$O_h(k^2)$ upper remainder. Nonnegativity follows from the actual
nested source minima, as in AMT15, not from comparison with a possibly
negative lower endpoint.

\subsection{The original Gamma multiplier, not a replacement higher source}
Retain $t_j=2j+1/2$, $f_l(S)=\prod_{j=0}^{l-1}(S-c-t_j)$ and
\[
 dm_{0,k}=\beta_k|f_l(c+iy)|^2d\sigma,\quad
 \beta_k=\frac{(2\pi)^{k/2}}{\sqrt2\Gamma(k/2)},\quad
 T_l=\prod_{j=0}^{l-1}t_j^2,\quad
 U_{l,N}=\prod_{j=0}^{l-1}[4(N+j+1)^2+t_j^2].
 \tag{AKS38}
\]
The exact source map is $\mathcal P_N\to f_l\mathcal P_N\subset
\mathcal P_{N+l}$, $P\mapsto f_lP$, with inverse division on this
specified subspace. Let $\mathsf E_\chi$ be the complete divided-jet
map and $U_f$ its Leibniz multiplication matrix,
$(U_f)_{r,a}=f_l^{(r-a)}(\lambda)/(r-a)!$ for $r\ge a$.
AMT25--28 gives the precise transported observation and kernel
\[
\begin{gathered}
 \widetilde\Lambda=\Lambda\mathsf E_\chi^{-1}U_f^{-1},\quad
 \widetilde I=U_f\mathsf E_\chi I,\quad
 \ker\widetilde\Lambda=U_f\mathsf E_\chi K,\\
 C_N=\mathscr K_{\chi,N+l}
   -\mathscr K_{\chi,f;N+l}\mathscr K_{f,N+l}^{-1}
                               \mathscr K_{f,\chi;N+l},\\
 H_{K,N}^{(k)}=\beta_k\widetilde I^*C_N^{-1}\widetilde I.
\end{gathered}\tag{AKS39}
\]
Each covariance is the original finite sum of full polynomial jet outer
products divided by its original monic norm. These formulas retain the
extra $f$ constraints and their full cross covariance. In particular
$M_fK=K$ is not assumed. AKS26 calculates the same tensor metric directly
in its own orthonormal source; the two representations agree by their
identical source norm and identical constrained minimum.

Define
\[
 w_N=\log(U_{l,N}/T_l),\quad
 \Omega^{\rm low}=w_{q-1}+w_q,\quad
 \Omega^{\rm high}=w_{2q-1}+w_{2q}.
\]
The already proved same-coefficient comparison AMT42--44 retains all
mass factors before cancellation and gives
\[
 -r\Omega^{\rm low}\le\Delta_K^\Gamma:=F_K^{(1)}-F_K^{(k)}
                         \le r\Omega^{\rm high}.
 \tag{AKS40}
\]
It therefore sharpens the two independent absolute bounds by
\[
\begin{split}
 \max\{\underline U_k,\underline U_1-r\Omega^{\rm high},0\}
       &\le F_K^{(k)}\\
       &\le\min\{\overline U_k,\overline U_1+r\Omega^{\rm low}\}.
\end{split}\tag{AKS41}
\]
The original finite estimate
$0\le w_N\le l\log[1+16(N+l)^2]$ proves that AKS40 is $o(kq)$ on the
actual rank-$r$ domain. Along with AKS35 this preserves the completed
fact that all three source kernels differ by $o(kq)$. AKS4, not AKS40,
is the new absolute-scale calculation.

\subsection{Boundary, QGQ and arithmetic receiving identities}
At each endpoint retain the original minimum section
$S_N^{(s)}=(G_N^{(s)})^{-1}\Lambda^*
(\Lambda(G_N^{(s)})^{-1}\Lambda^*)^{-1}$. The full Schur identity is
\[
 \det(I^*G_N^{(s)}I)\det Q_N^{(s)}
       =|\det[I,S_*]|^2\det G_N^{(s)},\qquad
 Q_N^{(s)}=(\Lambda(G_N^{(s)})^{-1}\Lambda^*)^{-1}.
 \tag{AKS42}
\]
Thus the new finite absolute boundary interval is
\[
 \max\{0,\mathcal B_k^{(s)}-\overline U_s\}
       \le F_B^{(s)}\le\mathcal B_k^{(s)}-\underline U_s.
 \tag{AKS43}
\]
The baselines in this formula are the \emph{complete} original BSL/HBL
objects, with their full finite relation errors. No $q^2$ coefficient is
assigned to $K$; its proved coefficient there remains zero.

Write $\mathcal H_k=\mathcal B_k^\sigma-\mathcal B_k^0=-\mathfrak T_k$.
The exact full identities still read
\[
\begin{gathered}
 F_B^{(1)}-F_B^{(k)}=\mathcal H_k-\Delta_K^\Gamma,\\
 \boxed{\quad\mathcal B_k^{\rm ar}
   =F_K^{(k)}+F_B^{(k)}+\mathcal H_k+\delta_k^\sigma
   =F_K^{(1)}+F_B^{(1)}+\delta_k^\sigma.\quad}
\end{gathered}\tag{AKS44}
\]
The Gamma interval is not re-estimated in this argument. For either
complete QRI1 or QRI2 interval $[\omega^-,\omega^+]$ for its already
calculated centre $W_k$, retain its exact source errors
$e_0\in[-a_0,b_0]$ and $e_q+e_{q+1}\in[-b_k^+,b_k^-]$. Then
\[
 \underline{\mathcal H}_k=-\omega^+-b_k^--a_0,
 \qquad\overline{\mathcal H}_k=-\omega^-+b_k^++b_0,
 \qquad \mathcal H_k\in
       [\underline{\mathcal H}_k,\overline{\mathcal H}_k].
 \tag{AKS45}
\]
These are exactly QRI3, not new unspecified error choices. The complete
QRI1--9, QGT and QGQ definitions remain verbatim in each updated receiver;
the new interval is intersected with all their retained tighter finite
intervals. In particular the simple-quartet nonzero error radius stays
\[
 \limsup\left|
 \frac{\mathcal H_k+lq\mathcal C_\Gamma}{q}
       +\frac{\partial_\alpha\mathcal L(2,\pi)}8
       -\frac{\log2}{4}\right|\le\frac{A_{\rm QGT}}{\sqrt2}.
 \tag{AKS46}
\]
No zero-radius or $o(q)$ arithmetic assertion is introduced.

For example the complete finite absolute defect relative to the tensor
boundary, with the actual Gamma and arithmetic scalars retained, is
\[
 \underline U_k+\mathcal H_k+\delta_k^\sigma
 \le\mathcal B_k^{\rm ar}-F_B^{(k)}
 \le\overline U_k+\mathcal H_k+\delta_k^\sigma.
 \tag{AKS47}
\]
Inserting AKS41 makes this sharper. Replacing $\mathcal H_k$ by the
appropriate endpoints of AKS45 and $\delta_k^\sigma$ by AKS37 gives its
fully inherited finite scalar enclosure with the same signs.
Also $F_B^{\rm ar}=\mathcal B_k^\sigma+\delta_k^\sigma-F_K^{\rm ar}$;
subtract the interval AKS35 or AKS36 to get its new absolute boundary
interval. This retains the exact correlations rather than taking
independent hypothetical source errors.

The existing limit is
$\mathcal H_k/(kq)\to-\mathcal C_\Gamma/4$ and
$\delta_k^\sigma/(kq)\to0$. Therefore the new interval implies
\[
\begin{gathered}
 \liminf_k\inf_{|u|\ge R_*}
 \frac{\mathcal B_k^{\rm ar}-F_B^{(k)}}{kq}
          \ge-\mathcal C_\Gamma/4,\\
 \limsup_k\sup_{|u|\ge R_*}
 \frac{\mathcal B_k^{\rm ar}-F_B^{(k)}}{kq}
          \le\kappa_*-\mathcal C_\Gamma/4.
\end{gathered}\tag{AKS48}
\]
All original branches are included in the infima and suprema. This does
not decide its sign. The already proved boundary \emph{difference}
limits of AMT49 are unchanged. Every subsequential limit of the absolute
normalized kernel lies in $[0,\kappa_*]$ and is shared by its three source
versions when they are evaluated on the same sequence of period
parameters. Existence of a single limit is not proved by boundedness.

Finally the HC residual remains exactly
$\mathcal R_k=\mathcal B_k^{\rm ar}-4q\log(D_hk)$, with its original
$D_h$ and original nonnegative contraction, phase and trace terms.
AKS44 changes none of them. The supplied positive $q^2$ baseline and
its complete Gamma correction are not re-presented as new results here.

\section{Actions, support faces, and the actual arithmetic class}
\label{sec:AKS-actions}

The coefficient formula AKS12 supplies all four blocks of the same sum
action: $P_KYP_K$, $P_KY(I-P_K)$, $(I-P_K)YP_K$ and
$(I-P_K)Y(I-P_K)$. There is no claim that either off-diagonal block
vanishes. In the actual minimum-section coordinates the corresponding
formula remains
\[
 [I,S_N]^{-1}Y[I,S_N]=
 \begin{pmatrix}
 H_{K,N}^{-1}I^*G_NYI&H_{K,N}^{-1}I^*G_NYS_N\\
 \Lambda YI&\Lambda YS_N
 \end{pmatrix}.
 \tag{AKS49}
\]
Its proof is multiplication by the inverse coordinate maps
$H_{K,N}^{-1}I^*G_N$ and $\Lambda$. In the invariant-ring model, the same
action is the derivation
$\mathcal Df=(\mathsf H\Omega\mathsf H^{-1}x)\cdot\nabla f$,
$\Omega=\operatorname{diag}(\gamma-i\delta,-\gamma-i\delta,
\gamma+i\delta,-\gamma+i\delta)$, with $\mathcal DQ=0$.
For the actual conductor $\mathcal A_Q=\prod_{a=1}^4Q_a$, its full defect
is $\mathcal D\mathcal A_Q=\sum_{a=1}^4(\mathcal DQ_a)
\prod_{b\ne a,\,1\le b\le4}Q_b$, exactly OQC22. Thus this coefficient
realization preserves the original action leaving the observed dual.

In particular the section-corrected vector is still OPR2's
\[
 C_s r_j^{\rm ker}=e_{q+j}-S_{q+j-1}^{(s)}\Lambda e_{q+j}
   =I\bigl(\zeta_{q+j}-\kappa S_{q+j-1}^{(s)}\Lambda e_{q+j}\bigr),
\]
with its full denominator $1+\xi_j$ in the kernel increment. The operator
proof estimates the complete covariance and its actual restriction, so
none of these coupled terms is replaced by a raw fixed-section vector.
The all-iterate kernel stays
$\bigcap_{a=0}^{q-1}\ker(\Lambda Y^a)$; nothing in AKS4 declares it zero.

On the existing arithmetic coefficient realization retain
$\mathcal V_{h,k}P=P(D_1+\cdots+D_k)F_h^{\otimes k}$ and
$J^{(k)}\mathcal V_{h,k}=\eta\operatorname{rem}_\chi$, where
$\eta=U^{\otimes k}\iota$ is the original injection. AKS17 and the
Newton coefficients compute precisely this same remainder. Any two
polynomial lifts differ by $\chi Q$ in the original relation space,
so their original arithmetic cohomology classes agree. The specified
primitive of the difference is EP.48, with the full differential and
cochain signs proved in AKP3--5 below; it depends on the lift difference.
The source comparison to $w_h^{*k}$ is the identity on original
polynomial coefficients, followed by the actual minima as in AMT8;
it is not a replacement arithmetic source.

The scalar object stays $G(R)=\{\tau\}\sqcup R^\bullet$, with supported
zero $0_R^\bullet$, injective amplitude/support pair, and infinite
integer target. Every coefficient map above is applied on each original
present support label and outer leg. It commutes with the specified
coefficient inclusions because both compositions apply the same linear
map to the same coefficient; a present zero is not converted into
absence. No new base object, purity condition, or identification of the
sum action with a Frobenius action is used. The actual distinguished
class and its higher-jet directions remain in the full fibre, even
when the observation kills them.

\section{Exact period domain and the interior alternative}
\label{sec:AKS-radius}

For reference the inherited radius used above is completely specified
as follows; these are AMT33--34/RPC17--23, with no new constants. Put
\[
\begin{gathered}
 \beta=2(\gamma^2-\delta^2),\quad\eta=(\delta^2+\gamma^2)^2,\quad
 a=v_h(\rho)\ne0,\\
 c_3=(a^{-2}-\bar a^{-2})/(4i\delta\gamma),\quad
 c_1=(a^{-2}+\bar a^{-2})/2-(\delta^2-\gamma^2)c_3,\\
 g_j=5^{j/5-1}e^{\pi ij/5}\Gamma(j/5),\quad
 g_+=\max_j|g_j|,\quad g_-^{-1}=\max_j|g_j|^{-1},\quad
 \kappa=\max_{j<4}|g_{j+1}/g_j|,\\
 E_j^{\rm ray}=\frac85\int_0^\infty r^j(r^3/3+r)
       e^{-r^5/5+r^3/3+r}dr,\quad
 C_T=2\left(\sum_{j=0}^3(E_j^{\rm ray})^2\right)^{1/2},\\
 C_K=4C_Tg_-^{-1}+2g_+g_-^{-1}+2g_+g_-^{-2}C_T,\\
 C_{K,3}=C_K(3\kappa^2+3\kappa C_K+C_K^2),\quad
 \epsilon_* =\min\{1,(2C_Tg_-^{-1})^{-1}\},\quad B_* =\beta+\eta.
\end{gathered}\tag{AKS50}
\]
Then
\[
R_*=
\begin{cases}
\displaystyle\left[\max\left\{1,\frac{B_*}{\epsilon_*},
\frac{2(|c_3|C_{K,3}B_*+|c_1|(\kappa+C_K))}{|c_3g_4/g_1|}
\right\}\right]^{5/2},&c_3\ne0,\\[4pt]
\displaystyle\left[\max\left\{1,\frac{B_*}{\epsilon_*},
\frac{2C_KB_*}{|g_2/g_1|}\right\}\right]^{5/2},&c_3=0.
\end{cases}\tag{AKS51}
\]
The actual unit implies $(c_1,c_3)\ne(0,0)$ and all integrals converge.
RPC's full period comparison proves a nonzero $(4,1)$ or $(2,1)$
commutator entry, respectively, for $|u|\ge R_*$ on every branch. It
therefore proves the five-member orbit needed in AKS11.

Inside this radius the same coefficient determinant AKS12--14 and the
operator bound AKS24 remain valid. The orbit is decided by the original
noninvariant quadratic coefficients, not by the radius alone. Since
$Q$ is nondegenerate, its projective orbit has either one or five
members. A one-member orbit must have character one: its character has
both fourth and fifth powers equal to one, by the quadratic determinant
and $D^5=I$. In that case $Q=a'x_1x_4+b'x_2x_3$ with $a'b'\ne0$, and
\[
 b=d_{5,k,0}-d_{5,k-2,0}
   =\left\lfloor\frac{(k+1)^2}{5}\right\rfloor
           +\mathbf1_{k\equiv0\text{ or }3\ (5)},\qquad r=q-b.
 \tag{AKS52}
\]
If a noninvariant coefficient is nonzero, AKS11 applies instead. These
are the complete OQC alternatives. The bound AKS31 always uses that
actual $r$; it does not extend the $8k-16$ specialization to an
unevaluated exceptional point. This calculation does not assert that
the interior exceptional set is empty.

\section{Verification scope and the remaining absolute coefficient}
\label{sec:AKS-scope}

The new mathematical content is AKS20--32 and its application to the
actual coefficient determinant AKS12--14, followed by the finite
receiving improvements AKS35--48. The complete source and orbit
identifications used in that application are the supplied OCS/OQC/OQX,
RPC, OPG/OPR and AMT results; their exact predecessors are preserved.
AKS33 pinpoints a repair obtainable immediately from the earlier proof,
whereas the Newton bound supplies the sharper explicit constant.

The bound \emph{at} scale $kq$ is now finite and uniform. What is not
proved is a matching lower coefficient, a unique limit for
$F_K^{(1)}/(kq)$, or a sign for the interval AKS48. Those questions involve
the orientation of the original period kernel within the inverse
remainder covariance; replacing that compression by a scalar bound
loses precisely that information. This is the finite-bounds outcome
allowed by the current request, not a claim to have evaluated an
asymptotic coefficient. The source differences are already completed
inputs and are not renamed as the remaining problem.

The accompanying exact checks distinguish finite algebraic diagnostics
from the analytic proof for the full family. They do not assert that a
diagnostic root is a zero of $\xi$, that a diagnostic matrix equals an
actual period matrix, or that a new Lean theorem has been checked.
No RH conclusion is drawn from this kernel bound.

% END CURRENT AKS

\clearpage
% BEGIN CURRENT AKJ
% Root continuation: full original residue, conductor and metric maps.
\section{The new ceiling in the original coefficient frame}
\label{sec:AKJ}
The complete incoming AKS proof sharpens the previously proved KAF bound.
Both estimates apply to the same original remainder and the same kernel.
We retain all their finite constants. In particular, for $m=1$,
$k=4l+1$, $l\ge1$, $q=(k+1)^2$, $c=k/2$, $0<\delta<1/2$,
$\gamma>2$, and every original branch with $|u|\ge R_*$, put
$r=8k-16$. The AKS constant is
\[
 \kappa_*=32\log\frac{25\log3}{4\pi}
 \in(25.02078097649,25.02078097650).
 \tag{AKJ1}
\]
AKS20--32 proves the full finite operator inequality from confluent
Newton division and coefficient extraction. Its parameters $t$ and $z$
are extraction radii in the existing source, whose orders remain $1,k$.
For $s\in\{1,k\}$ define, using the exact KAF and AKS quantities,
\[
 \begin{split}
 L_s&=d_1^{(s)},\\
 U_s&=\min\{2r\log Z_{k,1},\,
 r(\mathcal L_{s,q}+\mathcal L_{s,q+1})-d_1^{(s)}\}.
 \end{split}\tag{AKJ2}
\]
The first entry of the minimum is KAF's factorial estimate; the second
is AKS28, with its original first compressed increment retained. Thus
$L_s\le F_K^{(s)}\le U_s$. This pointwise minimum retains whichever
finite estimate is sharper at the actual original parameters. It does
not assume that the asymptotically smaller ceiling wins at every degree.

The two exact original coordinate constructions are related as follows.
Use OKM's invertible ordered Vandermonde $V_\beta$ and its actual
quotient matrix $\pi$. Put $I_\pi=V_\beta^{-1}\pi^{\mathsf T}$.
AKS12 constructs the coefficient-orthogonal projection $P_K$ from the
actual observation and its complete ordered-word Gram. Directly,
\[
 P_K=I_\pi(I_\pi^*I_\pi)^{-1}I_\pi^*,\qquad
 D_{K,N}^{(s)}
 =\frac{\det(\overline\pi(R_N^{(s)})^{-1}\pi^{\mathsf T})}
        {\det(I_\pi^*I_\pi)}.
 \tag{AKJ3}
\]
Indeed $I_\pi$ is injective with image $K$ by OKM5. The first displayed
matrix is Hermitian and idempotent, fixes this image, and kills its
coefficient-orthogonal complement, so it is exactly AKS12's projection.
Apply AKS14 with inclusion $I_\pi$ and then OKM8--9 to obtain the second
identity. The positive denominator is independent of endpoint and
source. Its exponent is $1+1-1-1=0$ in the four-endpoint quotient.
This proves the exact morphism between the incoming projector
determinant and the existing conductor-frame determinant, retaining the
original coefficient metric and every frame factor.

With the original AMT correlated caps $C_X^\pm$ and scalar
$\delta_k^\sigma$, the following joint intervals hold:
\[
\begin{split}
 \max\{0,L_1-C_K^-,L_1+\delta_k^\sigma-C_B^+\}
 &\le F_K^{\rm ar}\\
 &\le U_1+\min\{C_K^+,\delta_k^\sigma+C_B^-\},\\
 \max\{L_k,L_1-r\Omega^{\rm high},0\}
 &\le F_K^{(k)}
 \le\min\{U_k,U_1+r\Omega^{\rm low}\}.
\end{split}\tag{AKJ4}
\]
For proof, add each retained AMT lower or upper source-difference bound
to the corresponding endpoint of AKJ2, and intersect the resulting
intervals. In the first line the upper and lower source differences
are the same AMT18 correlated caps used in AKS36; no independent errors
have replaced their actual values. Every boundary interval is obtained
by subtracting these kernel intervals from the exact original total
$\mathcal B_k^a=F_K^a+F_B^a$. Consequently all three original kernels
have uniform limsup at most $\kappa_*$ after division by $kq$.
The source differences remain $o(kq)$. A matching leading value is
still to be calculated from these same matrices.

\section{Residue detection expressed in every original primary coefficient}
Retain now every original multiplicity $m\ge1$, $e=1+k(m-1)$,
$q=e(k+1)^2$, and the full polynomial
\[
 \chi(S)=\prod_{a,b=0}^k(S-\beta_{ab})^e,
 \quad \beta_{ab}=c+(2a-k)\delta+i(2b-k)\gamma,
 \quad E=\mathbb C[S]/(\chi).
 \tag{AKJ5}
\]
All distinct pairs give distinct roots. For a root $\beta$ write
$z=S-\beta$, $h_\beta(S)=\chi(S)/(S-\beta)^e$, and let
$\varepsilon_\beta$ be its original CRT idempotent. Define exactly
$\ell([f])=[S^{q-1}]\operatorname{rem}_\chi f$.
For every polynomial class the primary expansion gives
\[
 \ell(f)=\sum_\beta[z^{e-1}]
       \frac{f(\beta+z)}{h_\beta(\beta+z)}.
 \tag{AKJ6}
\]
The quotient on the right is a truncated power series through degree
$e-1$. To prove the identity, partial-fraction division of $f/\chi$
gives at each finite pole the coefficient of $(S-\beta)^{-1}$ shown
on the right. Expanding those principal parts at infinity, only their
first-order terms contribute to $S^{-1}$. Polynomial terms contribute
nothing. For the unique representative of degree below $q$, the same
$S^{-1}$ coefficient is its coefficient of $S^{q-1}$. This proves AKJ6
without deleting any repeated-root term. Since $h_\beta(\beta)\ne0$,
the bilinear pairing $(f,g)\mapsto\ell(fg)$ is nondegenerate: if the
first nonzero local coefficient of $f$ is $f_tz^t$, take
$g=\varepsilon_\beta z^{e-1-t}h_\beta$ to obtain $\ell(fg)=f_t\ne0$.

Let $\lambda_\nu$ be the specified zero-charge coefficient rows of
the original OCS observation, so the injection $\mathcal F_0$ gives
$\jmath\Lambda=\mathcal F_0\mathfrak A\mathcal H^{-1}$.
Retain the exact entries
\[
 a_{\nu,\beta,D}:=\lambda_\nu(\varepsilon_\beta z^D)
       =i^DD!c_{\nu;abD},\qquad 0\le D<e.
 \tag{AKJ7}
\]
Here $c_{\nu;abD}$ is the full OCS multinomial in the original period
and unit jets; $\beta=\beta_{ab}$. It is not an assigned coefficient.
The residue representative of each row has the explicit formula
\[
 A_\nu=\left[\sum_\beta\varepsilon_\beta h_\beta(S)
           \sum_{D=0}^{e-1}a_{\nu,\beta,D}(S-\beta)^{e-1-D}\right],
 \qquad \lambda_\nu(x)=\ell(xA_\nu).
 \tag{AKJ8}
\]
Indeed, restrict to one primary factor and substitute AKJ8 into AKJ6.
For $x=\varepsilon_\beta z^D$, multiplication cancels $h_\beta$ and
extracts exactly $a_{\nu,\beta,D}$. These vectors are a basis, proving
the asserted equality for every $x$. This gives the concrete inverse
of the residue equivalence used in the incoming formalization, on
the actual OCS rows, with their factorials and phases.

Define the integer $t_\beta$ from these actual coefficients by
\[
 t_\beta=\begin{cases}
 -1,&a_{\nu,\beta,D}=0\text{ for every }\nu,D,\\
 \max\{D:a_{\nu,\beta,D}\ne0\text{ for some }\nu\},&\text{otherwise},
 \end{cases}\qquad n_\beta=e-1-t_\beta.
 \tag{AKJ9}
\]
Define $\mathcal O_q:E\longrightarrow B^q$ by
$\mathcal O_q(x)=(\Lambda M_S^jx)_{0\le j<q}$, where $B=\operatorname{im}\Lambda$
is the original observation image. Thus $0\le n_\beta\le e$.
The complete original permanently invisible
submodule and its exact dimension are
\[
 \begin{gathered}
 K_\infty=\bigcap_{j=0}^{q-1}\ker(\Lambda M_S^j)
 =\bigoplus_\beta
   \varepsilon_\beta(z^{e-n_\beta})\subset E,\\
 \dim K_\infty=\sum_\beta n_\beta,
 \qquad \operatorname{rank}\mathcal O_q=q-\sum_\beta n_\beta.
 \end{gathered}\tag{AKJ10}
\]
Here $(z^e)=0$ and $(z^0)$ is the entire local algebra. To prove this,
vanishing of the first $q$ observations says
$\ell(S^jxA_\nu)=0$ for all $j<q,\nu$. Nondegeneracy just proved and
the power basis imply $xA_\nu=0$ for every row. Conversely those
products being zero kills every iterate. At the factor $\beta$,
AKJ8 is a unit $h_\beta$ times its displayed reverse coefficient
polynomial. The minimum valuation of all $A_\nu$ is precisely
$n_\beta$. Its common annihilator is $(z^{e-n_\beta})$: multiplication
by a coefficient of minimal valuation is a unit times $z^{n_\beta}$,
and all other coefficients have at least that valuation. This proves
the module formula and its dimension; rank-nullity proves the last
equality. Affine substitution $M_S=c+iY$ gives an invertible triangular
change between the first $q$ powers of the two operators, so the same
kernel is obtained using the original $Y$-stack of OCS.

In particular all iterates are jointly injective exactly when every
original highest-jet column $a_{\nu,\beta,e-1}$ has a nonzero entry.
For $m=1$, AKJ8 reduces to
\[
 A_\nu(\beta)=\chi'(\beta)B_{\nu,\beta},\qquad
 K_\infty=\operatorname{span}\{\varepsilon_\beta:
                         B_{\nu,\beta}=0\text{ for all }\nu\}.
 \tag{AKJ11}
\]
This is an exact test on the AKS8 matrix, not an assumption that it
passes. It constructs the inclusion $K_\infty\hookrightarrow K$:
each listed primary ideal is killed by all rows, hence in particular
by their original one-step observation. The source-minimum metric on
this subspace is its restriction from OKM8, retaining the same original
Gram. Full-iterate injectivity alone therefore supplies no estimate of
the one-step kernel determinant AKJ3.

\section{The actual conductor confines the kernel's polynomial degrees}
Return to $m=1$, $k=4l+1\ge9$ and the original five-member period domain.
Let $\mathcal A_Q=\prod_{j=1}^4Q_j$ and write its full Segre restriction
from OCF as
\[
 A=\sum_{r,s=0}^8a_{rs}
 X_1^rX_0^{8-r}Y_1^sY_0^{8-s}\ne0.
 \tag{AKJ12}
\]
OCF proves nonvanishing and retains all actual $\mathsf H$ coefficients.
No particular $a_{rs}$ is presumed nonzero. Write $k'=k-8$,
$q'=(k-7)^2$ and
\[
 b_{rs}=4+(2r-8)\delta+i(2s-8)\gamma,\quad
 (\mathcal T_Ap)(S')=\sum_{r,s=0}^8a_{rs}p(S'+b_{rs}),\quad
 \mu_j=\sum_{r,s}a_{rs}b_{rs}^j.
 \tag{AKJ13}
\]
All $81$ numbers $b_{rs}$ are distinct. Their Vandermonde is invertible,
so some $\mu_j$ with $0\le j\le80$ is nonzero. Define $v$ to be the first
such index. The exact finite differential expression on degree-below-$q$
polynomials is
\[
 \mathcal T_A=\sum_{j=0}^{q-1}\frac{\mu_j}{j!}\partial_S^j
 =\frac{\mu_v}{v!}\partial_S^v B(\partial_S),\quad
 B(X)=1+\sum_{j=v+1}^{q-1}\frac{v!\mu_j}{j!\mu_v}X^{j-v}.
 \tag{AKJ14}
\]
This is Taylor's finite formula. Hence it kills polynomials of degree
below $v$ and lowers every higher degree by exactly $v$. The operator
$B(\partial_S)$ is invertible: with $U=B(\partial_S)-I$, its inverse
is the finite sum $\sum_{a=0}^{q-1}(-U)^a$. Terms beyond nilpotency are
zero. Thus $\mathcal T_A$ maps polynomials of degree at most $d$
onto those of degree at most $d-v$ when $d\ge v$.

Let $\chi_{k'}$ be the original polynomial for tensor degree $k'$.
The conductor's exact annihilator in the original coefficient algebra is
\[
 V_A=\{p\in\mathbb C[S]_{<q}:\chi_{k'}\mid\mathcal T_Ap\},
 \qquad K\subset V_A.
 \tag{AKJ15}
\]
To prove this map, pair the biform $Ag$, of bidegree $(k,k)$, with
the primary values of $p$. If $g$ has coefficient $g_{ab}$ at bidegree
$(k',k')$, convolution of all original coefficients gives the pairing
\[
 \sum_{a,b=0}^{k'}g_{ab}
       \sum_{r,s=0}^8a_{rs}
       p(\beta^{(k')}_{ab}+b_{rs})
 =\sum_{a,b=0}^{k'}g_{ab}
                      (\mathcal T_Ap)(\beta^{(k')}_{ab}).
 \tag{AKJ16}
\]
The addition identity for the roots follows directly from their retained
centres and both signs. Vanishing for every $g$ is exactly vanishing at
all $q'$ distinct roots, hence the divisibility in AKJ15. OCF's proved
inclusion $A\mathbb C[X,Y]_{(k',k')}\subset J_k$ then gives $K\subset V_A$
by its original complex-linear coefficient pairing.

The following is a full degree-adapted basis of $V_A$. Define
$\mathcal J_v S^d=d!S^{d+v}/(d+v)!$, so
$\partial_S^v\mathcal J_v=I$, and set
\[
 \begin{gathered}
 1,S,\ldots,S^{v-1},\qquad
 P_h=\frac{(q'+h+v)!}{(q'+h)!}
       B(\partial_S)^{-1}\mathcal J_v(\chi_{k'}S^h),\\
 0\le h\le q-q'-v-1.
 \end{gathered}\tag{AKJ17}
\]
The first list is empty when $v=0$. Since $q-q'=16k-48\ge96>v$,
the range of the second list is nonnegative. Each $P_h$ is monic of
degree $q'+h+v$ and maps under $\mathcal T_A$ to
$\mu_v(q'+h+v)!/(v!(q'+h)!)$ times $\chi_{k'}S^h$. These statements
follow from AKJ14, the inverse's constant term one, and the leading
coefficient of $\mathcal J_v$. The images span every multiple of
$\chi_{k'}$ of degree at most $q-v-1$, and the low list spans the exact
kernel of $\mathcal T_A$. This proves spanning and independence. In
particular $\dim V_A=q-q'=16k-48$.

Choose the original OCF basis of $J_k/(A\mathbb C[X,Y]_{(k',k')})$,
whose exact dimension is $j_k=8k-32$. Evaluate each chosen biform row
on the primary values of every polynomial in AKJ17. This gives an
explicit $j_k$ by $(16k-48)$ matrix $\Xi$, with the full original
period coefficients. Its kernel maps isomorphically onto $K$ by the
displayed basis. Its rank is $j_k$: the original $J_k$ and conductor
annihilator duality, or dimension $\dim V_A-\dim K$, gives this rank.
For the submatrix $\Xi_{\le d}$ containing only basis columns of degree
at most $d$, the exact formula is
\[
 \dim(K\cap\mathbb C[S]_{\le d})
 =\#\{\text{degrees in AKJ17 at most }d\}
                      -\operatorname{rank}\Xi_{\le d}.
 \tag{AKJ18}
\]
Distinct monic degrees ensure that no omitted higher column can cancel
to form a degree-at-most-$d$ polynomial. Restricting the coefficient
isomorphism to this prefix and applying rank-nullity proves AKJ18.
It uses actual minors and makes no assertion of a generic prefix rank.

Let $d_0<\cdots<d_{r-1}$ be the degree jumps of the actual kernel.
The basis degrees of $V_A$, in increasing order, are $i$ for $i<v$,
and $q'+i$ for $i\ge v$. Inclusion and dimension of each degree prefix,
together with $K\subset\mathbb C[S]_{<q}$, imply
\[
 i+q'\mathbf1_{i\ge v}\le d_i\le q-r+i.
 \tag{AKJ19}
\]
For the lower inequality the $i$th independent kernel vector cannot
appear earlier than the $i$th independent vector of $V_A$. For the
upper inequality, intersection of an $r$-dimensional subspace of
$\mathbb C^q$ with its $(q-r+i+1)$-dimensional degree prefix has dimension
at least $i+1$. Summing yields the exact exterior-degree bounds
\[
 \begin{split}
 \frac{r(r-1)}2+\max\{r-v,0\}q'
 &\le\sum_{i=0}^{r-1}d_i
 \le r(q-r)+\frac{r(r-1)}2,\\
 0\le rq-\sum_i d_i
 &\le\frac{r(r+1)}2+rj_k+\min\{r,v\}q'.
 \end{split}\tag{AKJ20}
\]
The difference between the first line's upper and lower endpoints is
$rj_k+\min\{r,v\}q'$, since $q-q'=r+j_k$. Therefore the original
kernel's exterior degree is $rq-O(k^2)$, with the exact bound above
and $v\le80$, uniformly in the proved period domain. The affine map
$S=c+iy$ preserves all degree prefixes and multiplies a degree-$d$
leading coefficient by $i^d$. Thus these are also the degrees in the
original Gamma coordinate, with all phases known. They specify the
kernel's position without assigning a condition number to its frame.

\section{The exact conductor translations in the original source norm}
Keep either source $s\in\{1,k\}$, $b=s/2$, $M=(2\pi)^{s/2}$ and
the original $p_n,h_n=M n!(b)_n$. On polynomials of degree at most $D$,
let $\varphi_n=p_n/\sqrt{h_n}$ and $w_n=n!/(b)_n$.
Differentiating the original generating identity AKS16 in $y$ gives
\[
 \partial_y\varphi_n=
 \sum_{\substack{1\le a\le n\\a\text{ odd}}}
 \frac{(-1)^{(a-1)/2}}{a}
       \sqrt{\frac{w_n}{w_{n-a}}}\,\varphi_{n-a}.
 \tag{AKJ21}
\]
Indeed differentiation multiplies the generating series by
$\arctan z=\sum_{a\text{ odd}}(-1)^{(a-1)/2}z^a/a$.
Equating its finite coefficient and substituting the full norms gives
exactly AKJ21; the two copies of $M$ are cancelled only at this point.

Write $H_D=\sum_{j=1}^D1/j$ and $H_0=0$. The original source norm obeys
\[
 \|\partial_y\|\le\sqrt{H_D(H_D+4)}\le H_D+2,
 \qquad
 \|f(y+z)\|\le e^{|z|(H_D+2)}\|f\|.
 \tag{AKJ22}
\]
Here is a proof retaining both source orders. For $n>m\ge1$ and
$b\ge1/2$,
\[
 \frac{w_n}{w_m}\le\prod_{j=m+1}^n\frac{j}{j-1/2}
 \le\sqrt{\frac nm},\qquad w_n\le2\sqrt n\quad(n\ge1).
 \tag{AKJ23}
\]
The factor inequality follows by squaring and using
$j(j-1)\le(j-1/2)^2$; the product telescopes. For $m=0$, separate
the first factor, at most two. Apply the weighted Schur estimate to
the absolute matrix in AKJ21 with positive weights $\vartheta_n=1/\sqrt{w_n}$.
Its weighted row sum is at most $H_D$. Its weighted column sum at $n$
is at most $H_n+4$: for $1\le m<n$ the excess over $1/(n-m)$ is at most
\[
 \frac{\sqrt{n/m}-1}{n-m}
 =\frac1{\sqrt m(\sqrt n+\sqrt m)}\le\frac1{\sqrt{nm}},
\]
whose sum is at most two by comparison with $\int_0^n t^{-1/2}dt$.
The $m=0$ term is at most $2/\sqrt n\le2$. Enlarging the odd-index
sums to all indices preserves these bounds. For completeness, if a
nonnegative matrix has weighted row bound $A$ and column bound $B$,
Cauchy--Schwarz gives
$|(Tx)_m|^2\le A\vartheta_m\sum_nT_{mn}|x_n|^2/\vartheta_n$;
summing and using the column bound gives $\|Tx\|^2\le AB\|x\|^2$.
This proves the first bound in AKJ22, including $D=0$ where the
derivative is zero. The exact finite Taylor identity
$f(y+z)=\exp(z\partial_y)f(y)$ and its norm series prove the second.
The inverse translation is obtained by $-z$ with the same bound.

There is consequently an exact original-source upper control for AKJ13.
For $f(y)=p(c_k+iy)$, evaluate its output at $c_{k-8}+iy$, retaining
$c_k=k/2$ and $c_{k-8}=k/2-4$. Then
\[
 (\mathcal T_Ap)(c_{k-8}+iy)
 =\sum_{r,s=0}^8a_{rs}
   f\left(y+\frac{b_{rs}-4}{i}\right),
 \tag{AKJ24}
\]
and therefore
\[
 \|\mathcal T_Ap\|_{s,c_{k-8}}
 \le\left(\sum_{r,s}|a_{rs}|
          e^{|b_{rs}-4|(H_D+2)}\right)\|p\|_{s,c_k}.
 \tag{AKJ25}
\]
The notation $\|p\|_{s,c}=\|p(c+iy)\|_{m_{0,s}}$ changes only the
explicit centre; the measure, mass and source order $s$ are unchanged.
AKJ24 proves the centre correction rather than discarding it. For a
fixed original period and quartet the coefficient in AKJ25 grows at
most polynomially in $D$, since $H_D\le1+\log D$ for $D\ge1$ and
$|b_{rs}-4|\le8\sqrt{\delta^2+\gamma^2}$. The entire actual coefficient
sum remains visible. The finite inverse in AKJ17 has the specified
moment denominators $\mu_v$; AKJ25 asserts no bound on that inverse.
Together AKJ18--25 provide exact degree information and a source-norm
estimate for the actual conductor map, to be used in the remaining
four-endpoint compression calculation. No leading kernel value or
kernel-invariance assumption has replaced that calculation.

% END CURRENT AKJ

\clearpage
% BEGIN CURRENT AKP
\section{The original theta class and the checked metric interface}
\label{sec:AKP}
The absolute base and original coefficient resolution are retained:
\[
 \mathbf F_{1,\tau}=\{\tau,1\},\qquad
 \mathsf A_\tau(\mathcal T)=
 [V\oplus V\xrightarrow{\Theta(\phi-\widehat\psi)}\mathscr B],
 \qquad Q=\mathscr B/\Theta V.
 \tag{AKP1}
\]
The four-point chart has $+,-<\eta<\sigma$ and coefficients
$(V,V,\mathscr B,0)$, with the original restrictions $\Theta$ and
$J\Theta=\Theta\mathcal F$. Its tensor resolution identifies the top
cohomology with $Q^{\otimes k}$. AT1--2a supplies the complete original
construction. In its notation, retaining the full unit
$\upsilon_h=j_h(2\xi/h)$, the actual cochain and arithmetic class maps are
\[
 \begin{gathered}
 \mathcal V_{h,k}(P)=P(D_1+\cdots+D_k)F_h^{\otimes k},\qquad
 E=\mathbb C[S]/(\chi),\\
 \eta_{h,k}[P]=\upsilon_h^{\otimes k}P(A_h^{(k)})1,\qquad
 q^{(k)}\mathcal V_{h,k}(P)
   =\sigma_h^{\otimes k}\eta_{h,k}[P].
 \end{gathered}\tag{AKP2}
\]
Here $\sigma_h:E_h\hookrightarrow Q$ is the original section and
$\eta_{h,k}:E\hookrightarrow E_h^{\otimes k}$ is the full-jet injection.
Both injections are part of the supplied source proof; the cochain map
has not been replaced by an arbitrary observation.

For completeness, the exact representative-change calculation is as
follows. If $P'=P+\chi\kappa$, perform successive monic divisions of
$\chi(s_1+\cdots+s_k)$ by $h(s_a)$. Its final remainder is zero because
$\chi$ is the original annihilator of the sum in
$\bigotimes_a\mathbb C[s_a]/(h)$. Thus the actual divisions produce
$Q_a(\mathbf s)$ satisfying
\[
 \chi(s_1+\cdots+s_k)=\sum_{a=1}^kh(s_a)Q_a(\mathbf s).
 \tag{AKP3}
\]
The chosen ordered division determines these polynomials; no uniqueness
of all such decompositions is asserted. On the original tensor complex set
\[
 \Psi_\kappa=\sum_{a=1}^k(-1)^{a-1}
 Q_a(\mathbf D)\kappa\left(\sum_bD_b\right)
 \bigl(F_h^{\otimes(a-1)}\otimes\phi_0
                  \otimes F_h^{\otimes(k-a)}\bigr).
 \tag{AKP4}
\]
Its differential is
\[
 d\Psi_\kappa=\mathcal V_{h,k}(\chi\kappa)
             =\mathcal V_{h,k}(P')-\mathcal V_{h,k}(P).
 \tag{AKP5}
\]
Indeed the only nonzero differential in each summand acts at its
$\phi_0$ position. The tensor sign is $(-1)^{a-1}$, cancelling the
displayed sign. The identity $\Theta\phi_0=h(D)F_h$ gives the $a$th
term of AKP3; all operators commute with this differential. Summing
proves AKP5 with the original signs. This is EP.48, expressed at the
current receiving map. Consequently the quotient classes in AKP2 agree.
The chosen primitive for their difference depends on $\kappa$.
At an existing support label $\ell$, this quotient map sends the
relation to $(\ell,0)$ and sends external absence $\tau$ to $\tau$.
The coordinate change has therefore preserved the arithmetic class
and its support through an explicit cochain homotopy.

This calculation specifies the correction to the sentence following
AKS49: equality of cohomology classes is retained, with AKP4 the
primitive of the change. Membership in an additional period-observation
kernel alone supplies no such primitive. Its proved map is the inclusion
$K_\infty\hookrightarrow K\hookrightarrow E$ of AKJ10--11, followed by
the original arithmetic injection in AKP2. In particular, a nonzero
coefficient in this additional kernel remains a nonzero original
arithmetic class under that injection. Its zero observed value and its
arithmetic class are related by these exact maps.

\subsection{The complete minimum-section calculation at each endpoint}
Let $G$ be one of the original positive Hermitian coefficient Grams,
write $K_G=G^{-1}$, and let $\Lambda:E\twoheadrightarrow B$ be the
actual observation onto its image. Define
\[
 Q_G=(\Lambda K_G\Lambda^*)^{-1},\qquad
 S_G=K_G\Lambda^*Q_G.
 \tag{AKP6}
\]
The inverse exists because $\Lambda^*$ is injective and $K_G$ is
positive definite. Direct multiplication, using Hermitian symmetry,
gives
\[
 \Lambda S_G=I_B,\quad S_G^*G=Q_G\Lambda,\quad
 S_G^*GS_G=Q_G.
 \tag{AKP7}
\]
For every rectangular column family $X$ put $Z=\Lambda X$ and
$R=X-S_GZ$. Then
\[
 \Lambda R=0,\qquad R^*GS_G=0,\qquad
 R^*GR=X^*GX-Z^*Q_GZ\succeq0.
 \tag{AKP8}
\]
The first two equalities follow from AKP7. In the expansion of the
last Gram, each cross term is $Z^*Q_GZ$ and the section Gram is the
same matrix. The resulting single subtraction proves the identity,
including every off-diagonal pairing. Any other lift is $S_GZ+R'$
with $\Lambda R'=0$. Its Gram is $Z^*Q_GZ+R'^*GR'$, proving the
minimum property by the same orthogonality calculation.

Retain a fixed section $S_*$, the actual inclusion $I_K$, and original
coordinate map $\kappa_K$ satisfying
$I_K\kappa_K=I_E-S_*\Lambda$. Applying this identity to $R$ gives
\[
 R=I_K\bigl(\kappa_KX-\kappa_KS_G\Lambda X\bigr).
 \tag{AKP9}
\]
Thus the fixed-section coordinates include their specified minimum
correction. For one original new column $e_D$, define
$t_D=e_D^*Ge_D$, $\lambda_D=\Lambda e_D$ and
$\xi_D=\lambda_D^*Q_G\lambda_D$. AKP8 yields
$t_D-\xi_D\ge0$, so the exact ordered factors satisfy
\[
 (1+\xi_D)\left(1+\frac{t_D-\xi_D}{1+\xi_D}\right)=1+t_D.
 \tag{AKP10}
\]
The denominator is positive and is retained. These are the full OPR
minimum-section identities, now proved at the original Grams of
AKS26 and OKM8. For an empty observation image, the section is the
zero map, the Gram on that image has determinant one, and the same
identities reduce to $R=X$. No empty present support is deleted.

The new formalization in draft PR32 checks these finite identities and
the original observation-iterate and residue-annihilator maps. Its
checked implementation is
\[\texttt{8e78bc7c240b04d297ade6afdadfd863e0c6db7b}.\]
The later source and status head is
\[\texttt{9cec8482f2ecf978cf8bdb67b4c080b9dd74f5d6}.\]
AKJ6--11 evaluates the residue inverse in every actual primary-jet
coefficient, giving the exact map to those checked declarations.
The analytic AKS estimate, the conductor metric estimates, and the
polynomial-gcd dimension corollary have the full written proofs stated
here or in their supplied sources; they are outside that Lean run.

% END CURRENT AKP

\clearpage
% BEGIN CURRENT AIE
% Full original conductor inverse estimate. Child calculation, 2026-09-14.
\section{The finite conductor inverse in the two original Gamma norms}
\label{sec:AIE}
\paragraph{Current refinement.}
The finite estimates in this chapter remain valid. RUI5--15 and
CEP29--49 below sharpen their exterior consequence: complementary
singular values and the exact source quotient give an $o(kq)$
conductor error at all four original endpoints, for a fixed actual
period. The full original constants, moments and low-degree graph
are retained in those proofs.


Retain AKJ12--25 and the complete OCF conductor coefficient polynomial.
Thus $m=1$, $k=4l+1\ge9$, $q=(k+1)^2$, $0<\delta<1/2$,
$\gamma>2$, and the original period belongs to the proved five-orbit
locus. Its fixed $4$ by $4$ matrix $\mathsf H$ determines the four
quadrics and all $81$ coefficients $a_{rs}$ of their full product.
They remain the same coefficients as the tensor degree $k$ varies.
Put
\[
 b_{rs}=4+(2r-8)\delta+i(2s-8)\gamma,\quad
 \mu_j=\sum_{r,s=0}^8a_{rs}b_{rs}^j,\quad
 v=\min\{j:\mu_j\ne0\}\le80.
 \tag{AIE1}
\]
The existence and bound for $v$ are the full distinct-node Vandermonde
argument of AKJ13. Throughout, $0\le D\le q-1$ denotes the polynomial
degree cutoff. Every norm below is the original Gamma norm of source
order $s\in\{1,k\}$, with
\[
 \begin{gathered}
 b=s/2,\quad M=(2\pi)^{s/2},\quad h_n=M n!(b)_n,\\
 \varphi_n=p_n/\sqrt{h_n},\quad w_n=n!/(b)_n,
 \quad \|p\|_{s,c}=\|p(c+iy)\|_{m_{0,s}}.
 \end{gathered}\tag{AIE2}
\]
The centre $c$ is displayed wherever it changes. No source order,
source mass, coefficient, or period is assigned a replacement value.

\subsection{A coefficient-controlled disk for the literal finite symbol}

The exact symbol used in AKJ14 and AKJ17 is
\[
 B_q(z)=1+\sum_{j=v+1}^{q-1}
                \frac{v!\mu_j}{j!\mu_v}z^{j-v}.
 \tag{AIE3}
\]
The sum is empty if its lower index exceeds its upper index. Define
the positive, original-data constants
\[
 \begin{gathered}
 A_*^{\rm abs}=\sum_{r,s}|a_{rs}|,\qquad
 B_*^{\rm sh}=\max_{r,s}|b_{rs}|,\qquad
 C_v=\frac{A_*^{\rm abs}(B_*^{\rm sh})^v}{|\mu_v|},\\
 R_0=\frac{1}{B_*^{\rm sh}}
       \min\left\{1,\frac{v+1}{2\mathrm e C_v}\right\},
 \quad \rho_0=\tanh R_0.
 \end{gathered}\tag{AIE4}
\]
Here $A_*^{\rm abs}>0$, $B_*^{\rm sh}\ge4$, $C_v\ge1$,
$R_0>0$, and $0<\rho_0<1$. The superscripts distinguish these
constants from the original conductor polynomial and from every
earlier radius or pivot coefficient. The exact moment $\mu_v$,
including any cancellation in it, remains in the denominator.

For every $|z|\le R_0$ one has
\[
 |B_q(z)-1|
 \le \frac{v!}{|\mu_v|}
       \sum_{j=v+1}^{q-1}\frac{|\mu_j|}{j!}|z|^{j-v}
 \le C_v\frac{B_*^{\rm sh}|z|}{v+1}
                    \exp(B_*^{\rm sh}|z|)
 \le\frac12.
 \tag{AIE5}
\]
To prove the middle inequality, use
$|\mu_j|\le A_*^{\rm abs}(B_*^{\rm sh})^j$, write
$j=v+1+h$, and apply
$v!/(v+1+h)!\le1/((v+1)h!)$. Summation of the exponential
series gives the displayed bound. Definition AIE4 gives
$B_*^{\rm sh}R_0\le1$ and
$C_v B_*^{\rm sh}R_0/(v+1)\le1/(2\mathrm e)$, proving
the last inequality. The first upper bound in AIE5 retains every
individual finite moment and may be used directly for a larger
verified radius. The explicit $R_0$ proves a positive disk for
every actual coefficient chart; no uniform lower bound for
$|\mu_v|$ over varying periods is asserted.

The analytic function
\[
 F_q(z)=\frac1{B_q(-i\arctan z)}
          =\sum_{a=0}^\infty f_{q,a}z^a
 \tag{AIE6}
\]
is well defined on a neighbourhood of the closed disk
$|z|\le\rho_0$, and
$|F_q(z)|\le2$ there. Indeed the power series of $\arctan$
gives $|\arctan z|\le\operatorname{arctanh}|z|\le R_0$;
AIE5 gives $|B_q(-i\arctan z)|\ge1/2$. Compactness gives
an open neighbourhood with nonzero denominator. Cauchy's coefficient
formula therefore proves
\[
 f_{q,0}=1,\qquad |f_{q,a}|\le2\rho_0^{-a}
                    \quad(a\ge1).
 \tag{AIE7}
\]
Although the literal symbol $B_q$ changes with $q$, the disk and its
bound depend only on the original fixed conductor and shifts.

\subsection{The exact coefficient map and its norm}

Let $\mathcal P_D$ denote the polynomials in $y$ of degree at most
$D$, with the orthonormal basis in AIE2. Define the lowering map
\[
 L_s\varphi_0=0,\qquad
 L_s\varphi_n=\sqrt{\frac{w_n}{w_{n-1}}}\,\varphi_{n-1}
                         \quad(1\le n\le D).
 \tag{AIE8}
\]
Its powers have exactly the entries
$L_s^a\varphi_n=\sqrt{w_n/w_{n-a}}\varphi_{n-a}$ for
$a\le n$ and vanish for $a>n$. Thus $L_s^{D+1}=0$.
AKJ21 gives the polynomial identity
$\partial_y=\arctan L_s$, with every power above degree $D$
zero. For the coordinate pullback $\mathcal C_c p(y)=p(c+iy)$,
the chain rule is
$\mathcal C_c\partial_S\mathcal C_c^{-1}=-i\partial_y$.
Consequently the inverse of the literal finite operator is
\[
 \begin{split}
 \mathcal C_c B_q(\partial_S)^{-1}\mathcal C_c^{-1}
   &=F_q(L_s),\\
 F_q(L_s)\varphi_n
   &=\sum_{a=0}^n f_{q,a}
                     \sqrt{\frac{w_n}{w_{n-a}}}\,\varphi_{n-a}.
 \end{split}\tag{AIE9}
\]
For proof, evaluate the analytic power-series identity
$B_q(-i\arctan z)F_q(z)=1$ in the nilpotent matrix $L_s$.
All relevant coefficients are finite, so multiplication is an exact
finite polynomial identity. This also proves that AIE9 is the same
inverse as the finite nilpotent Neumann sum in AKJ14. The factor $-i$
is forced by the original coordinate map and has been retained.

Define the exact weight ratio and the resulting bound
\[
 \begin{split}
 \alpha_{D,s}&=\max_{0\le m\le n\le D}
                         \sqrt{\frac{w_n}{w_m}},\\
 N_{D,s}&=\alpha_{D,s}
            \left(1+2\sum_{a=1}^{D}\rho_0^{-a}\right)
 \le \frac{2\alpha_{D,s}}{1-\rho_0}\rho_0^{-D}.
 \end{split}\tag{AIE10}
\]
Every absolute row sum and every absolute column sum of AIE9 is at
most $\alpha_{D,s}\sum_{a=0}^D|f_{q,a}|$. Cauchy--Schwarz in each
row followed by the column bound proves that the operator norm is
at most this number; AIE7 gives $N_{D,s}$. It follows that, at every
unchanged centre $c$,
\[
 \|B_q(\partial_S)^{-1}p\|_{s,c}
                \le N_{D,s}\|p\|_{s,c}.
 \tag{AIE11}
\]
When $s=k$, $b=k/2\ge1$ and
$w_n/w_{n-1}=n/(b+n-1)\le1$, so $\alpha_{D,k}=1$.
When $s=1$, the same ratio is greater than one, so
$\alpha_{D,1}=\sqrt{D!/(1/2)_D}$; for $D\ge1$, the exact
factor estimate of AKJ23 yields
$\alpha_{D,1}\le\sqrt2 D^{1/4}$. At $D=0$, both values are
one. In particular
\[
 \log\|B_q(\partial_S)^{-1}\|_{s,c}
 \le D\log(\coth R_0)
       +\log\frac{2\alpha_{D,s}}{1-\tanh R_0}
             =O_{\mathsf H,\delta,\gamma}(D+\log(D+1)).
 \tag{AIE12}
\]
For the fixed original quartet and period, $q=(k+1)^2$ and
$D\le q-1$ imply that this logarithm divided by $kq$ tends to
zero. This estimate holds for both original source orders at once;
its constants explicitly retain the actual $\mu_v$.

\subsection{The infinite symbol differs by an exact low-degree map}

To record precisely the truncation in AIE3, define the entire function
\[
 B_\infty(z)=\frac{v!}{\mu_v z^v}
                 \sum_{r,s}a_{rs}e^{b_{rs}z},\qquad B_\infty(0)=1.
 \tag{AIE13}
\]
Its apparent singularity is removable by the definition of $v$.
On polynomials of degree at most $D\le q-1$, set
$E_q=B_\infty(\partial_S)-B_q(\partial_S)$; its exact expression is
\[
 E_q=\sum_{n=q-v}^{D}
       \frac{v!\mu_{n+v}}{(n+v)!\mu_v}\partial_S^n.
 \tag{AIE14}
\]
An empty sum is zero. Its image has degree at most
$D-q+v\le v-1$. Both symbols have constant term one, hence both
are invertible in this finite polynomial algebra. Their inverses obey
\[
 B_q(\partial_S)^{-1}-B_\infty(\partial_S)^{-1}
 =B_q(\partial_S)^{-1}E_qB_\infty(\partial_S)^{-1},
 \tag{AIE15}
\]
whose image is contained in $\mathbb C[S]_{<v}$ and whose rank is
at most $v$. Equation AIE15 follows by multiplying out its right
side. The degree statement holds because every factor is a power
series in the degree-lowering derivative. Thus the infinite-symbol
formula can be related to the literal AKJ inverse by the displayed
exact correction; it cannot replace that inverse without this map.

\subsection{An explicit right inverse of the original conductor}

Assume now $D\ge v$. Define the raising map, on the degree ranges
where its image lies in $\mathcal P_D$, by
\[
 J_s\varphi_n=\sqrt{\frac{w_n}{w_{n+1}}}\varphi_{n+1}
       =\sqrt{\frac{b+n}{n+1}}\varphi_{n+1}.
 \tag{AIE16}
\]
Then $L_s^vJ_s^v=I$ on $\mathcal P_{D-v}$ and
$\|J_s^v\|\le\max\{1,b^{v/2}\}$. The latter follows from
the displayed coefficient, at each raising step, by
$(b+n)/(n+1)\le\max\{1,b\}$.

The analytic function $G(z)=z/\arctan z$, with $G(0)=1$,
satisfies $|G(z)|\le2$ on $|z|<1$. Indeed
\[
 \frac{\arctan z}{z}=\int_0^1\frac{dt}{1+t^2z^2},\qquad
 \operatorname{Re}\frac1{1+u}>\frac12\quad(|u|<1).
 \tag{AIE17}
\]
The real-part inequality follows from
$2\operatorname{Re}(1+\overline u)>|1+u|^2$, which reduces
to $1>|u|^2$. It proves that the integral is nonzero with
real part greater than $1/2$, including the removable value at zero.
Every Taylor coefficient of $G(z)^v$ is therefore bounded by $2^v$:
apply Cauchy's formula on radius $t<1$ and let $t$ increase to one
for each fixed coefficient.

Define a map between the specified degree spaces by
\[
 R_{v,s}=G(L_s)^v J_s^v:\mathcal P_{D-v}\longrightarrow\mathcal P_D.
 \tag{AIE18}
\]
Since all functions of $L_s$ commute, AIE8 and the power-series
identity $(\arctan z)^vG(z)^v=z^v$ give
$\partial_y^vR_{v,s}=L_s^vJ_s^v=I$. This proves the right-inverse
identity even at the top allowed degree. The same row and column
argument as for AIE10, now with every coefficient bounded by $2^v$,
gives
\[
 \|R_{v,s}\|\le
   2^v(D+1)\alpha_{D,s}\max\{1,b^{v/2}\}.
 \tag{AIE19}
\]
When $v=0$ the map is exactly the identity; AIE19 remains a valid
upper bound. Define the corresponding map on polynomials in the
original $S$ coordinate by
\[
 \mathcal R_{v,s,c}
       =\mathcal C_c^{-1}i^vR_{v,s}\mathcal C_c.
 \tag{AIE20}
\]
The chain rule in AIE9 proves
$\partial_S^v\mathcal R_{v,s,c}=I$, with exactly the factor
$(-i)^vi^v=1$. This is a map in the original polynomial algebra;
the source norm only specifies its coefficients through the explicitly
given original orthonormal basis.

Let $c_k=k/2$ and $c'=c_{k-8}=k/2-4$. The full conductor of AKJ13
has the following right inverse on $\mathbb C[S']_{\le D-v}$:
\[
 \mathcal R_{A,s}g
     =\frac{v!}{\mu_v}
             B_q(\partial_S)^{-1}\mathcal R_{v,s,c'}g,
 \qquad \mathcal T_A\mathcal R_{A,s}g=g.
 \tag{AIE21}
\]
To check it, use the exact finite factorization
$\mathcal T_A=(\mu_v/v!)\partial_S^vB_q(\partial_S)$ from
AKJ14 on degree at most $D$. The two inverse factors cancel and
AIE20 gives the stated identity. In the identity the output variable
is $S'$; the differential expression uses the same numerical
polynomial coefficients, so the domain and codomain are as displayed.

Set $H_D=\sum_{j=1}^D1/j$, $H_0=0$, and
\[
 C_{A,D,s}=\frac{v!}{|\mu_v|}\,
  e^{4(H_D+2)}N_{D,s}\,
  2^v(D+1)\alpha_{D,s}\max\{1,b^{v/2}\}.
 \tag{AIE22}
\]
Then the original-centre norm estimate is
\[
 \|\mathcal R_{A,s}g\|_{s,c_k}
                    \le C_{A,D,s}\|g\|_{s,c'}.
 \tag{AIE23}
\]
Indeed AIE11 and AIE19 first prove the bound without the exponential
factor, with both norms at $c'$. For any output polynomial $p$,
$p(c_k+iy)=p(c'+i(y-4i))$, so AKJ22 applied to the exact shift
$-4i$ supplies the remaining factor. Thus the four-unit centre
difference is retained in the typed norm map. For the fixed actual
quartet and period, $v\le80$ and $b\in\{1/2,k/2\}$ give
\[
 \log C_{A,D,s}
 \le D\log(\coth R_0)+O_{\mathsf H,\delta,\gamma}
                         (\log(D+1)+\log(k+1)),
 \qquad \frac{\log C_{A,D,s}}{kq}\longrightarrow0.
 \tag{AIE24}
\]
Every factor in the first assertion is written in AIE22; the
asymptotic statement follows from $D\le(k+1)^2-1$ and the fixed
finite $v$. No varying-period bound is inferred from this statement.

For comparison with the exact monomial lift $\mathcal J_v$ of AKJ17,
put
$L_g=\mathcal R_{v,s,c'}g-\mathcal J_vg$. Both terms have
$v$th derivative $g$, so $L_g\in\mathbb C[S]_{<v}$. More explicitly,
\[
 L_g=\sum_{a=0}^{v-1}
       [S^a](\mathcal R_{v,s,c'}g)S^a,
 \quad
 \mathcal R_{A,s}g-\frac{v!}{\mu_v}B_q(\partial_S)^{-1}\mathcal J_vg
        =\frac{v!}{\mu_v}B_q(\partial_S)^{-1}L_g.
 \tag{AIE25}
\]
The first identity uses the exact zero coefficients of $\mathcal J_vg$
in degrees below $v$. This proves the complete degree-below-$v$
correction between the two lifts, rather than discarding their
different choices of integration constants.

\subsection{The exact quotient and exterior transfer supplied by this bound}

The kernel of $\mathcal T_A$ on degree at most $D\ge v$ is exactly
$\mathbb C[S]_{<v}$ by AKJ14. Give the quotient
$\mathbb C[S]_{\le D}/\mathbb C[S]_{<v}$ its Hilbert quotient norm
from the original source at $c_k$, and give its target
$\mathbb C[S']_{\le D-v}$ the original source at $c'$.
The induced conductor map $\overline{\mathcal T}_A$ is an
isomorphism, because AIE21 proves surjectivity and the displayed
kernel proves injectivity. Define the full upper norm from AKJ25 by
\[
 N_{A,D}=\sum_{r,s}|a_{rs}|
            e^{|b_{rs}-4|(H_D+2)}.
 \tag{AIE26}
\]
For every target $g$, its minimum-norm preimage class obeys
\[
 \frac{\|g\|_{s,c'}}{N_{A,D}}
 \le\|\overline{\mathcal T}_A^{-1}g\|_{\rm quot}
 \le C_{A,D,s}\|g\|_{s,c'}.
 \tag{AIE27}
\]
For the left inequality, apply AKJ25 to every preimage and take the
infimum. For the right inequality, use the actual preimage AIE21.
The infimum is attained since the omitted subspace is a closed
finite-dimensional subspace. These arguments prove the norm statement
with the specified quotient map, rather than assigning the coefficient
norm to the quotient.

In particular AIE21 gives the exact decomposition
\[
 V_A=\mathbb C[S]_{<v}\ \mathbin{\oplus}\
   \mathcal R_{A,s}\bigl(
         \chi_{k-8}\mathbb C[S']_{\le D-v-q'}\bigr),
 \quad q'=(k-7)^2,
 \tag{AIE28}
\]
where a negative degree means the zero space and $V_A$ here denotes
the AKJ15 space restricted to degree at most $D$. Indeed the image
under $\mathcal T_A$ of a member of $V_A$ is precisely a multiple
of $\chi_{k-8}$ in the displayed degree range. Subtracting its lift
leaves the exact kernel. The directness follows by applying
$\mathcal T_A$. Equation AIE25 identifies this decomposition with
the earlier monic basis, including its low-degree correction.

For any $t$-dimensional target subspace and any basis column matrix
$X$ in its original orthonormal coordinates, AIE27 implies the
full Gram determinant inequalities
\[
 N_{A,D}^{-2t}\det(X^*X)
 \le\det\bigl(X^*(\overline{\mathcal T}_A^{-1})^*
                   \overline{\mathcal T}_A^{-1}X\bigr)
 \le C_{A,D,s}^{2t}\det(X^*X).
 \tag{AIE29}
\]
For proof, AIE27 is a two-sided positive Hermitian matrix inequality
after congruence by $X$. Conjugating by $(X^*X)^{-1/2}$ gives
eigenvalues between the two scalar squared bounds, whose product is
the displayed determinant ratio. This proves the exact exterior
transfer and all powers of the norm constants.

Similarly, for any $t$-plane $U\subset\mathbb C[S]_{\le D}$,
AIE11 gives
$\|\bigwedge^t(B_q(\partial_S)^{-1}|_U)\|\le N_{D,s}^t$.
Although $B_q(\partial_S)$ is triangular with every diagonal entry
one and hence determinant one on the full polynomial space, that
full-space determinant does not cancel its restriction to an
arbitrary subspace. The exact restricted determinant is
$\det(X^*(B_q^{-1})^*B_q^{-1}X)/\det(X^*X)$ for a basis matrix
$X$ of that subspace; this is the map to which AIE11 applies.

For the original kernel rank $r=8k-16$ and $D\asymp q$, the upper
logarithm in these exterior bounds is $O_{\mathsf H,\delta,\gamma}(rq)
=O_{\mathsf H,\delta,\gamma}(kq)$. Therefore AIE12 and AIE24 prove
a sub-$kq$ logarithmic bound for each operator, while AIE29 retains
the rank factor required for volumes. They provide the original
conductor inverse, its full low-degree comparison and its quotient
metric control. A sub-$kq$ error for the actual rank-$r$ four-endpoint
determinant is not a consequence of these scalar bounds; that
remaining calculation must use the actual restricted matrices and
the existing endpoint correlation.

% END CURRENT AIE


\clearpage
% BEGIN CURRENT RUI
\section{Subleading exterior control from the actual finite conductor}
\label{sec:RUI}
Retain the original conductor of AKJ12--25 and AIE1--29. In particular
the quartet, full unit, original period and branch are fixed in the
proved five-orbit domain, $m=1$, $k=4l+1\ge9$, and $q=(k+1)^2$.
All $81$ coefficients $a_{rs}$ are those of the actual product of four
period quadrics. They are fixed as $k$ varies. Retain
\[
 \begin{gathered}
 b_{rs}=4+(2r-8)\delta+i(2s-8)\gamma,\qquad
 \mu_j=\sum_{r,s=0}^8a_{rs}b_{rs}^j,\qquad
 v=\min\{j:\mu_j\ne0\}\le80,\\
 A_{\rm abs}=\sum_{r,s}|a_{rs}|,\quad
 B_{\rm sh}=\max_{r,s}|b_{rs}|,\quad
 C_v=\frac{A_{\rm abs}B_{\rm sh}^v}{|\mu_v|}\ge1.
 \end{gathered}\tag{RUI1}
\]
No coefficient or moment denominator is replaced by a generic value.
The finite bound for $v$ is AKJ13's distinct-node Vandermonde proof.
The inequality $C_v\ge1$ follows from the triangle inequality applied
to the actual $v$th moment. The source measure at order $s\in\{1,k\}$
is exactly
\[
 dm_{0,s}(y)=(2\pi)^{s/2}
 \frac{2^{s/2}}{4\pi\Gamma(s/2)}
 |\Gamma(s/4+iy/2)|^2\,dy,\qquad
 \|p\|_{s,c}=\|p(c+iy)\|_{L^2(m_{0,s})}.
 \tag{RUI2}
\]
Every centre appearing below is stated, and the whole mass remains in
the source norm. Put $H_D=\sum_{j=1}^D1/j$, $H_0=0$.

\subsection{The forward polynomial bound and exact determinant}
For an integer cutoff $L\ge v+1$, define the literal finite symbol
\[
 B_L(z)=1+\sum_{a=1}^{L-v-1}
       \frac{v!\mu_{v+a}}{(v+a)!\mu_v}z^a.
 \tag{RUI3}
\]
The empty sum is zero. This includes the actual coefficient symbol
$B_q$ of AKJ14 and the source symbol $B_{N+1}$ needed at a larger
source cutoff $N$. For every displayed coefficient, including the
constant coefficient, one has
\[
 \left|\frac{v!\mu_{v+a}}{(v+a)!\mu_v}\right|
 \le C_v\frac{v!B_{\rm sh}^a}{(v+a)!}
 \le C_v\frac{B_{\rm sh}^a}{a!}.
 \tag{RUI4}
\]
The first inequality is the full moment bound
$|\mu_{v+a}|\le A_{\rm abs}B_{\rm sh}^{v+a}$.
The second follows from $(v+a)!\ge v!a!$.

On $\mathbb C[S]_{\le D}$ with the norm RUI2, the original
coordinate pullback sends $\partial_S$ to $-i\partial_y$.
AKJ21--23 proves $\|\partial_S\|\le H_D+2$, uniformly in both
original source orders. Applying RUI4 term by term to the finite
operator polynomial gives
\[
 \|B_L(\partial_S)\|_{s,c}\le M_D,
 \qquad M_D=C_v\exp\bigl(B_{\rm sh}(H_D+2)\bigr)\ge1.
 \tag{RUI5}
\]
Indeed each power has norm at most $(H_D+2)^a$; extending this
nonnegative finite sum to the exponential series proves the displayed
bound. The constant term is one and every other term strictly lowers
degree. Thus, on this exact $(D+1)$-dimensional polynomial space,
\[
 \det B_L(\partial_S)=1.
 \tag{RUI6}
\]
This is its determinant in the original increasing-power basis, and
therefore in every basis of the same source space, including the
original orthonormal Gamma basis. The assertion preserves the source
measure; it is not a rescaling of this determinant.

\subsection{Complementary singular values retain the full exterior norm}
Let $n=D+1$ and $T=B_L(\partial_S)$ acting on this same Hilbert space.
Write its singular values as $\sigma_1\ge\cdots\ge\sigma_n>0$.
RUI6 gives $\prod_i\sigma_i=1$. For every $0\le t\le n$,
\[
 \begin{split}
 \|\bigwedge^tT^{-1}\|
 &=\prod_{j=n-t+1}^n\sigma_j^{-1}
   =\prod_{j=1}^{n-t}\sigma_j
   =\|\bigwedge^{n-t}T\|
   \le M_D^{n-t},\\
 \|\bigwedge^tT\|&\le M_D^t.
 \end{split}\tag{RUI7}
\]
For a proof of the exterior norm formula, the singular-value
decomposition expresses $T$ as two unitary maps and a diagonal map
with these positive entries. Exterior powers of the unitary maps
preserve the induced norm; the diagonal exterior map has entries
equal to the products over all $t$-subsets. Its largest entry is the
product of the largest $t$ singular values. Applying this to $T^{-1}$
and using the full determinant proves every equality in RUI7. Empty
products cover $t=0,n$.

Consequently, for any actual $t$-dimensional subspace and any full-rank
basis column matrix $X$ in the original orthonormal source coordinates,
\[
 M_D^{-2t}\det(X^*X)
 \le\det\bigl(X^*(T^{-1})^*T^{-1}X\bigr)
 \le M_D^{2(n-t)}\det(X^*X).
 \tag{RUI8}
\]
For the upper inequality apply the first exterior norm of RUI7 to
the decomposable exterior vector of the columns of $X$. Its squared
norm is their Gram determinant. For the lower inequality, put
$Y=T^{-1}X$ and apply the second bound of RUI7 to $TY=X$.
This proves the lower bound on this same subspace without requiring
that it be invariant. Every original off-diagonal Gram entry is
retained by the displayed matrices.

In particular, for fixed original period and quartet, uniformly in
$s\in\{1,k\}$, every $0\le D\le2q$ and every allowed $t$ and $X$,
\[
 \left|\log\frac{
 \det(X^*(T^{-1})^*T^{-1}X)}{\det(X^*X)}\right|
 \le2(D+1)\bigl(\log C_v+B_{\rm sh}(H_D+2)\bigr)
 =O_{\rm actual}(q\log(q+1))=o(kq).
 \tag{RUI9}
\]
The harmonic estimate $H_D\le1+\log D$ for $D\ge1$ follows by
integrating $1/x$ on consecutive intervals; $D=0$ is immediate.
Since $q=(k+1)^2$, the displayed upper bound divided by $kq$ tends
to zero, independently of $D,t,X$. This is a uniform estimate over
the actual subspaces, not a presumed favourable position of the
kernel. Its constants are fixed only after fixing the original period
and quartet; no uniform bound on $1/|\mu_v|$ over moving periods is
asserted. The same inequalities for $T$ follow by exchanging the
image subspace and its inverse image.

RUI9 improves the rank-multiplied estimate at the end of AIE29.
The earlier estimate is valid, but its factor $t$ on a single inverse
operator bound is unnecessary. The complementary exterior identity
uses the exact full determinant and supplies the sharper bound on
the same original matrices.

\subsection{The actual conductor quotient inherits the sharper volume control}
Let $D\ge v$ and use the literal source truncation $L=D+1$.
Taylor's finite formula gives, on degree at most $D$,
\[
 \mathcal T_A=\sum_{j=v}^{D}\frac{\mu_j}{j!}\partial_S^j
             =\frac{\mu_v}{v!}\partial_S^vB_{D+1}(\partial_S),
 \qquad\ker\mathcal T_A=\mathbb C[S]_{<v}.
 \tag{RUI10}
\]
The first equality is exactly the translation operator in AKJ13;
all lower moments vanish by their proved definition. The factor
$B_{D+1}(\partial_S)$ is invertible and preserves each degree
filtration, so the last equality follows from the derivative kernel.

Use the original source centre $c_k=k/2$ and target centre
$c'=c_{k-8}=k/2-4$, both with source order $s$. AIE16--20 gives
the full right inverse $\mathcal R_{v,s,c'}$ of $\partial_S^v$
with norm at most
\[
 P_{D,s}=2^v(D+1)\alpha_{D,s}\max\{1,(s/2)^{v/2}\},
 \quad \alpha_{D,s}=\max_{0\le m\le n\le D}
          \sqrt{\frac{n!/(s/2)_n}{m!/(s/2)_m}}.
 \tag{RUI11}
\]
In particular $\alpha_{D,k}=1$, while
$\alpha_{D,1}\le\sqrt2\max\{1,D^{1/4}\}$; the full derivation
and the phase $i^v$ are the original AIE16--20 construction.
The exact conductor right inverse is
\[
 \mathcal R_{A,s,D}=
 \frac{v!}{\mu_v}B_{D+1}(\partial_S)^{-1}\mathcal R_{v,s,c'}.
 \tag{RUI12}
\]
Multiplying by RUI10 proves $\mathcal T_A\mathcal R_{A,s,D}=I$.
The numerical polynomial coefficients use the original source and
target variables, as in AIE21. The centre change from $c'$ to $c_k$
is the original shift $y\mapsto y-4i$, of norm at most
$\exp(4(H_D+2))$ by AKJ22.

Give the quotient $\mathbb C[S]_{\le D}/\mathbb C[S]_{<v}$ its
minimum norm from the source at $c_k$. The induced conductor
$\overline{\mathcal T}_A$ is an isomorphism onto
$\mathbb C[S']_{\le D-v}$ with its source norm at $c'$. Define
\[
 L_{D,s}=\frac{v!}{|\mu_v|}
          e^{4(H_D+2)}P_{D,s},\qquad
 N_{A,D}=\sum_{r,s}|a_{rs}|
             e^{|b_{rs}-4|(H_D+2)}.
 \tag{RUI13}
\]
For any $t$-dimensional target subspace, with original basis matrix
$X$, its exact quotient Gram ratio satisfies
\[
 N_{A,D}^{-2t}\det(X^*X)
 \le\det\bigl(X^*(\overline{\mathcal T}_A^{-1})^*
                       \overline{\mathcal T}_A^{-1}X\bigr)
 \le L_{D,s}^{2t}M_D^{2(D+1-t)}\det(X^*X).
 \tag{RUI14}
\]
To prove the upper bound, use the actual lift RUI12. Apply
$\mathcal R_{v,s,c'}$ to the $t$ basis columns, retaining its
specified image subspace. RUI7 bounds the exterior power of the
inverse $B_{D+1}$ on that image by $M_D^{D+1-t}$. The scalar,
right-inverse norm, and centre translation give $L_{D,s}^t$.
Projection to the minimum lift in the quotient is an orthogonal
projection on the same source Hilbert space, so its exterior norm
is at most one. This proves the complete upper Gram inequality.
For the lower bound, apply the forward original conductor to the
minimum lifts. AKJ24--25 gives its norm $N_{A,D}$; taking its
$t$th exterior power proves the result. This proof retains the
coupled Gram and the full quotient metric.

Every factor in RUI13 is explicit in the fixed original data.
For the asymptotic assertion restrict now to $v\le D\le2q$.
Since $v\le80$, $s\in\{1,k\}$ and $t\le D+1-v\le2q+1$,
both logarithmic endpoints of RUI14 have absolute value at most
\[
 2t\bigl(|\log L_{D,s}|+|\log N_{A,D}|\bigr)
 +2(D+1)\log M_D
 =O_{\rm actual}(q\log(q+1)+q\log(k+1))=o(kq).
 \tag{RUI15}
\]
For the logarithm of $N_{A,D}$, use
$A_{\rm abs}\le N_{A,D}\le
A_{\rm abs}\exp(\max|b_{rs}-4|(H_D+2))$.
For $L_{D,s}$, its displayed factors and the positive lower bound
$P_{D,s}\ge1$ give the same logarithmic estimate, retaining
$v!/|\mu_v|$. This proves the absolute bounds even if either
constant is below one. The harmonic bound and $q=(k+1)^2$
then prove the final limit. No dimension-$q$ scalar comparison has
silently acquired an extra factor of $q$ in this argument.

At the original coefficient cutoff $D=q-1$, RUI14 is the sharpened
version of AIE29, on precisely the same conductor quotient and
original coefficient subspaces. At a larger source cutoff $D=N$,
its literal symbol is $B_{N+1}$, and RUI10 proves the exact extension.
The old symbol $B_q$ cannot be substituted without recording the
finite difference. Passing this estimate through the original
remainder square requires the actual image of the relation ideal;
that square, its image graph, and the endpoint quotient metrics
are retained in the succeeding calculation.

% END CURRENT RUI

\clearpage
% BEGIN CURRENT CEP
% Original conductor endpoint square, graph quotient and determinant transfer.
\section{The original conductor at every source endpoint}
\label{sec:CEP}

Retain the complete simple-quartet conductor of OCF and AKJ, its
actual original period matrix, and $k=4l+1\ge9$ on the proved
five-orbit locus. Write
\[
 \begin{gathered}
 q=(k+1)^2,\quad k'=k-8,\quad q'=(k-7)^2,\quad
 \Delta=q-q'=16k-48,\\
 c=k/2,\quad c'=k/2-4,\quad
 \beta_{ij}=c+(2i-k)\delta+i(2j-k)\gamma,\\
 \beta'_{ab}=c'+(2a-k')\delta+i(2b-k')\gamma,\\
 A=\sum_{r,s=0}^8a_{rs}e_{rs}^{(8)}\ne0,\quad
 b_{rs}=4+(2r-8)\delta+i(2s-8)\gamma,\\
 \mu_j=\sum_{r,s}a_{rs}b_{rs}^j,\qquad
 v=\min\{j:\mu_j\ne0\}\le80.
 \end{gathered}\tag{CEP1}
\]
The factor $i$ before an imaginary term is the imaginary unit;
subscript indices are integers in their displayed ranges. The
original polynomials are
$\chi_k(S)=\prod_{i,j=0}^k(S-\beta_{ij})$ and
$\chi_{k'}(S')=\prod_{a,b=0}^{k'}(S'-\beta'_{ab})$.
Both have distinct roots. Set
$E_k=\mathbb C[S]/(\chi_k)$ and
$E_{k'}=\mathbb C[S']/(\chi_{k'})$.
Every source below has the original order $s\in\{1,k\}$ and
full mass $M_s=(2\pi)^{s/2}$, with $b=s/2$ and
$h_n=M_s n!(b)_n$. In particular the lower grid does not change
source order $k$ to $k-8$.

\subsection{The full source, remainder and primary-coordinate square}

For every polynomial define the original finite-translation map
\[
 (\mathcal T_Ap)(S')=\sum_{r,s=0}^8a_{rs}p(S'+b_{rs}).
 \tag{CEP2}
\]
For $0\le a,b\le k'$ and $0\le r,s\le8$, direct substitution gives
$\beta'_{ab}+b_{rs}=\beta_{a+r,b+s}$. All indices on the right lie
between zero and $k$. Consequently
\[
 (\mathcal T_A(\chi_k f))(\beta'_{ab})=0
 \quad\hbox{for every polynomial }f.
 \tag{CEP3}
\]
Division by the monic $\chi_{k'}$ and evaluation at its distinct roots
then show $\chi_{k'}\mid\mathcal T_A(\chi_k f)$. Thus CEP2 induces
the specified complex-linear map
$\overline{\mathcal T}_A:E_k\longrightarrow E_{k'}$.

In the original ordered primary coordinates its matrix is
\[
 C_{(a,b),(i,j)}
 =\begin{cases}
 a_{i-a,j-b},&0\le i-a,j-b\le8,\\
 0,&\text{otherwise}.
 \end{cases}
 \qquad C\in\operatorname{Mat}_{q'\times q}(\mathbb C).
 \tag{CEP4}
\]
Evaluation of CEP2 proves this entry formula, including its absence
of conjugation or binomial coefficients. Its ordinary transpose
$C^{\mathsf T}$ is the exact biform coefficient multiplication
by the nonzero polynomial $A$. That multiplication is injective
because the biform ring embeds in a polynomial integral domain.
Therefore $C$ has rank $q'$ and is surjective. With the original
ordered Vandermonde matrices $V_\beta,V_{\beta'}$, the polynomial
coordinate matrix is exactly $V_{\beta'}^{-1}CV_\beta$.
The kernel is the AKJ15 space
\[
 V_A=\{[p]\in E_k:\chi_{k'}\mid\mathcal T_Ap\},
 \qquad \dim V_A=\Delta.
 \tag{CEP5}
\]
CEP3 proves independence of the representative, and CEP4 proves
the dimension. The original common kernel $K=\ker\Lambda$ is
contained in $V_A$ by the full conductor inclusion OCF17.

For every endpoint
$N\in\{q-1,q,2q-1,2q\}$, put $d=N-v$ and define
\[
 \begin{gathered}
 \mathcal P_N=\mathbb C[S]_{\le N},\qquad
 \mathcal P'_d=\mathbb C[S']_{\le d},\\
 U_N:\mathcal P_N\longrightarrow E_k,\quad p\longmapsto[p],
 \qquad
 V_d:\mathcal P'_d\longrightarrow E_{k'},\quad g\longmapsto[g].
 \end{gathered}
\]
The top horizontal map is precisely
$\mathcal T_{A,N}:\mathcal P_N\longrightarrow\mathcal P'_{N-v}$.
Taylor's finite identity and the first nonzero moment show that
CEP2 lowers every degree at least $v$ by exactly $v$ and kills
exactly $\mathbb C[S]_{<v}$. The exact commutative square is
\[
 \begin{array}{ccc}
 \mathcal P_N &\xrightarrow{\ \mathcal T_{A,N}\ }&\mathcal P'_{N-v}\\
 {\scriptstyle U_N}\downarrow&&\downarrow{\scriptstyle V_{N-v}}\\
 E_k&\xrightarrow{\ \overline{\mathcal T}_A\ }&E_{k'} .
 \end{array}\tag{CEP6}
\]
The lower remainder is surjective because
$N-v\ge q-1-v\ge q'-1$; the latter inequality is
$\Delta\ge v$, and $\Delta\ge96>80\ge v$.

The map is complex-linear with its following full action defect.
If $M_S,M_{S'}$ denote the original multiplication operators, then
\[
 CM_S-M_{S'}C=C_b,\qquad
 (C_b)_{(a,b),(i,j)}=
 \begin{cases}a_{i-a,j-b}b_{i-a,j-b},&0\le i-a,j-b\le8,\\
 0,&\text{otherwise}.
 \end{cases}
 \tag{CEP7}
\]
Indeed each entry is multiplied by
$\beta_{ij}-\beta'_{ab}=b_{i-a,j-b}$. In the original centred
coordinates $Y=(M_S-cI)/i$, $Y'=(M_{S'}-c'I)/i$, this becomes
$CY-Y'C=(C_b-4C)/i$. Thus the exact relation to multiplication
retains its coefficient defect.

\subsection{The literal finite inverse at the larger cutoffs}

On $\mathcal P_N$, define the literal finite symbol
\[
 B_{N+1}(z)=1+\sum_{j=v+1}^{N}
              \frac{v!\mu_j}{j!\mu_v}z^{j-v},\qquad
 \mathcal T_{A,N}=\frac{\mu_v}{v!}
                   \partial_S^vB_{N+1}(\partial_S).
 \tag{CEP8}
\]
The expression on the right is evaluated in the output variable
$S'$. Taylor expansion of each term in CEP2 through degree $N$
proves this identity coefficient by coefficient. Its constant term
one makes $B_{N+1}(\partial_S)$ invertible: its difference from
the identity lowers degree, so the inverse is its finite nilpotent
Neumann sum through power $N$. A right inverse of $\partial_S^v$
on the polynomial spaces exists by integration of each monomial;
composition proves surjectivity of the top map in CEP6.

The relation to the earlier literal $B_q$ of AKJ14 is
\[
 \begin{split}
 D_{N,q}&=B_{N+1}(\partial_S)-B_q(\partial_S)
       =\sum_{j=q}^{N}\frac{v!\mu_j}{j!\mu_v}\partial_S^{j-v},\\
 B_{N+1}(\partial_S)^{-1}-B_q(\partial_S)^{-1}
       &=-B_{N+1}(\partial_S)^{-1}D_{N,q}B_q(\partial_S)^{-1},\\
 \mathcal T_{A,N}-\frac{\mu_v}{v!}\partial_S^vB_q(\partial_S)
       &=\sum_{j=q}^{N}\frac{\mu_j}{j!}\partial_S^j.
 \end{split}\tag{CEP9}
\]
Empty sums are zero. The first two differences have image contained
in $\mathbb C[S]_{\le N-q+v}$; the last has image contained in
$\mathbb C[S']_{\le N-q}$. The inverse formula follows by
multiplication, and its degree assertion follows since both inverse
factors preserve every degree prefix. At $N=q-1$ all differences
vanish. At the other three endpoints these are the full retained
defects; their ranks are not assigned the cutoff-independent bound
$v$.

For comparison, the entire symbol
$B_\infty(z)=v!\sum_{r,s}a_{rs}e^{b_{rs}z}/(\mu_vz^v)$,
with removable value one at zero, obeys at each cutoff
\[
 B_\infty(\partial_S)-B_{N+1}(\partial_S)
 =\sum_{n=N+1-v}^{N}
       \frac{v!\mu_{n+v}}{(n+v)!\mu_v}\partial_S^n.
 \tag{CEP10}
\]
Its image and the image of the corresponding inverse difference
lie in $\mathbb C[S]_{<v}$. This follows from the first derivative
degree $N+1-v$, and from the inverse identity with the same order
as AIE15. Thus the rank-at-most-$v$ comparison is with the
correctly growing finite symbol in CEP8.

For completeness the original-source inverse estimate extends
directly to these literal symbols. Retain the exact constants
$A_*^{\rm abs},B_*^{\rm sh},C_v,R_0,\rho_0$ of AIE4.
For every $|z|\le R_0$ the finite coefficient sum in CEP8 is at most
the full absolute moment sum:
\[
 |B_{N+1}(z)-1|
 \le \frac{v!}{|\mu_v|}
          \sum_{j=v+1}^N\frac{|\mu_j|}{j!}|z|^{j-v}
 \le C_v\frac{B_*^{\rm sh}|z|}{v+1}e^{B_*^{\rm sh}|z|}
 \le\frac12.
 \tag{CEP11}
\]
The middle inequality follows by writing $j=v+1+h$ and using
$v!/(v+1+h)!\le1/((v+1)h!)$. Hence
$F_{N+1}(z)=1/B_{N+1}(-i\arctan z)$ has coefficient absolute
values at most $2\rho_0^{-a}$ on degree $a\ge1$, and constant
coefficient one. This is Cauchy's formula on $|z|=\rho_0$, since
$|\arctan z|\le\operatorname{arctanh}\rho_0=R_0$.

In the original orthonormal Gamma basis, the lowering matrix
$L_s\varphi_n=\sqrt{w_n/w_{n-1}}\varphi_{n-1}$,
$w_n=n!/(b)_n$, satisfies
$\partial_y=\arctan L_s$ and $L_s^{N+1}=0$.
Evaluation of the power-series inverse in this nilpotent matrix
therefore gives the exact finite operator
$F_{N+1}(L_s)$, with entries
$[z^{n-m}]F_{N+1}(z)\sqrt{w_n/w_m}$.
The chain rule $S=c+iy$ supplies the retained $-i$.
Bounding the absolute row and column sums gives
\[
 \|B_{N+1}(\partial_S)^{-1}\|_{s,c}
 \le \alpha_{N,s}\left(1+2\sum_{a=1}^N\rho_0^{-a}\right),
 \qquad
 \alpha_{N,s}=\max_{0\le m\le n\le N}\sqrt{w_n/w_m}.
 \tag{CEP12}
\]
This proves the larger-cutoff assertion for the actual finite
symbol rather than substituting $N$ into an identity for $B_q$.
Likewise the explicit AIE18 raising right inverse on these degree
spaces, followed by $i^v$, $B_{N+1}(\partial_S)^{-1}$ and
$v!/\mu_v$, is a right inverse of CEP8. Its comparison with the
monomial integral remains exactly the degree-below-$v$ correction
in AIE25, now with this literal inverse.

\subsection{The exact source operator matrix and its full determinant}

Use the unchanged original orthonormal polynomials
$u_n(S)=p_n((S-c)/i)/\sqrt{h_n}$ in $\mathcal P_N$ and
$u'_m(S')=p_m((S'-c')/i)/\sqrt{h_m}$ in $\mathcal P'_{N-v}$.
Let $A_N$ denote the matrix of $\mathcal T_{A,N}$ in these two
frames. Define the analytic coefficient series at zero
\[
 t(z)=\sum_{r,s}a_{rs}
       \exp\left(\frac{b_{rs}-4}{i}\arctan z\right)
       =\sum_{\ell\ge0}t_\ell z^\ell .
 \tag{CEP13}
\]
Then its exact entries are
\[
 (A_N)_{m,n}=
 \begin{cases}
 t_{n-m}\sqrt{w_n/w_m},&n\ge m,\\
 0,&n<m,
 \end{cases}
 \quad
 0\le m\le N-v,\quad0\le n\le N.
 \tag{CEP14}
\]
To prove this, the original generating identity
$\sum p_n(y)z^n/n!=(1+z^2)^{-b/2}e^{y\arctan z}$
shows that translation by $\eta$ multiplies the series by
$e^{\eta\arctan z}$. Extracting its finite coefficients and
dividing by the complete $\sqrt{h_n}$ gives
$\sqrt{w_n/w_m}$ after the two copies of $M_s$ cancel at this
specified step. Evaluation of CEP2 at $c'+iy$ uses
$p(c+i(y+(b_{rs}-4)/i))$, proving the exact centre shift in CEP13.

For $\ell<v$, the shifted moment
$\sum a_{rs}(b_{rs}-4)^\ell$ is zero by the binomial theorem and
the definition of $v$; at $\ell=v$ it equals $\mu_v$.
Since $\arctan z=z+O(z^3)$, this gives
\[
 t_0=\cdots=t_{v-1}=0,\qquad
 t_v=\frac{(-i)^v\mu_v}{v!}\ne0.
 \tag{CEP15}
\]
The first $v$ columns of $A_N$ vanish. Its remaining square block
$\widehat A_N$, with columns $v,\ldots,N$, is upper triangular
with diagonal $t_v\sqrt{w_{m+v}/w_m}$.
Thus its rank is $N-v+1$, its kernel in the source is exactly
the first $v$ orthonormal coordinates, and
\[
 \begin{split}
 \det\widehat A_N
   &=\left(\frac{(-i)^v\mu_v}{v!}\right)^{N-v+1}
                   \prod_{m=0}^{N-v}\sqrt{\frac{w_{m+v}}{w_m}},\\
 \det(A_NA_N^*)
   &=\left|\frac{\mu_v}{v!}\right|^{2(N-v+1)}
                   \prod_{m=0}^{N-v}\frac{w_{m+v}}{w_m}.
 \end{split}\tag{CEP16}
\]
This retains the complex determinant phase and the original moment.
The operator norm also has the exact original-source upper bound
\[
 \|A_N\|\le
 M_{A,N}:=\sum_{r,s}|a_{rs}|
               e^{|b_{rs}-4|(H_N+2)},\qquad
 H_N=\sum_{j=1}^N1/j.
 \tag{CEP17}
\]
Indeed AKJ22 bounds each exact translation in CEP13 by its
displayed factor, and the triangle inequality gives the sum.

There is a closed four-endpoint expression for CEP16. For the
signs $\sigma=(1,1,-1,-1)$ at the ordered endpoints, write
$\mathcal D_A=\sum_N\sigma_N\log\det(A_NA_N^*)$.
Cancellation of the full products gives
\[
 \begin{split}
 \mathcal D_A
  &=-4q\log\left|\frac{\mu_v}{v!}\right|+\mathcal W_{v,s},\\
 \mathcal W_{v,s}
  &=\sum_{a=0}^{v-1}\left(
      \log w_{q-1-a}+\log w_{q-a}
      -\log w_{2q-1-a}-\log w_{2q-a}\right).
 \end{split}\tag{CEP18}
\]
For $v=0$ the second sum is zero. The product identity used here is
$\prod_{m=0}^{N-v}w_{m+v}/w_m
 =\prod_{a=0}^{v-1}w_{N-a}/\prod_{a=0}^{v-1}w_a$,
obtained by cancellation of consecutive products; it remains
valid when the two index intervals overlap.

Put $c_b^{\log}=\max\{1,b-1\}$. For every $j\ge1$ and $b\ge1/2$,
\[
 \left|\log\frac{j}{b+j-1}\right|\le\frac{c_b^{\log}}j .
 \tag{CEP19}
\]
For $b\ge1$ use $\log(1+(b-1)/j)\le(b-1)/j$.
For $1/2\le b<1$, put $a=1-b\le1/2$ and use
$-\log(1-a/j)\le a/(j-a)\le1/j$; the last inequality is
$aj\le j-a$, which follows from $a\le1/2$ and $j\ge1$.
Summing CEP19 over the two intervals of length $q$ in each
summand of CEP18 proves the completely explicit estimate
\[
 |\mathcal W_{v,s}|
 \le \frac{2v\,c_b^{\log}q}{q-v},\qquad
 |\mathcal D_A|
 \le4q\left|\log\left|\frac{\mu_v}{v!}\right|\right|
             +\frac{2v\,c_b^{\log}q}{q-v}.
 \tag{CEP20}
\]
All denominators are positive since $q\ge100>v$.
For either original source and the fixed actual quartet and period,
the displayed bound is $O_{\rm actual}(q+k)=o(kq)$.

\subsection{The induced original quotient metric and its correction}

Let the matrices of the two remainder maps in the orthonormal source
frames and original primary target frames be denoted by the same
$U_N,V_d$ as their maps. Explicitly
\[
 (U_N)_{(i,j),n}=\frac{p_n((\beta_{ij}-c)/i)}{\sqrt{h_n}},
 \qquad
 (V_d)_{(a,b),m}=\frac{p_m((\beta'_{ab}-c')/i)}{\sqrt{h_m}} .
 \tag{CEP21}
\]
They have sizes $q$ by $N+1$ and $q'$ by $N-v+1$,
respectively. Their first $q$, respectively $q'$, columns have
nonzero evaluation Vandermonde determinant, so both have full row
rank. The square CEP6 is the precise matrix identity
\[
 C U_N=V_dA_N,\qquad
 R_N=U_NU_N^*,\quad R'_d=V_dV_d^*,\quad
 G_N=R_N^{-1},\quad G'_d=(R'_d)^{-1}.
 \tag{CEP22}
\]
Here $G_N$ is the original upper primary quotient metric of OKM8,
and $G'_d$ is the lower-grid metric in the same source order $s$.
The induced metric on $E_k/V_A$, in its specified target coordinates
$E_{k'}$, is
\[
 Q_N=(C G_N^{-1}C^*)^{-1}
       =(V_dA_NA_N^*V_d^*)^{-1}.
 \tag{CEP23}
\]
To prove the first assertion, the minimum of $x^*G_Nx$ under
$Cx=y$ is attained at
$x=G_N^{-1}C^*(CG_N^{-1}C^*)^{-1}y$.
Every other lift differs by an element of $\ker C$ orthogonal
to this lift in the $G_N$ metric, proving the minimum and its
Gram matrix. The second equality is CEP22.

The exact correction to the lower original Gamma quotient metric is
\[
 \begin{gathered}
 W_d=V_d^*(R'_d)^{-1/2},\quad W_d^*W_d=I_{q'},\\
 K_{A,N}=W_d^*A_NA_N^*W_d>0,\qquad
 Q_N=(R'_d)^{-1/2}K_{A,N}^{-1}(R'_d)^{-1/2},\\
 \log\det Q_N=\log\det G'_d-\log\det K_{A,N}.
 \end{gathered}\tag{CEP24}
\]
All square roots are the unique positive Hermitian square roots.
These formulas follow by multiplying out the definition of
$K_{A,N}$ and inverting. In particular the passage between the
induced metric and the lower original source metric retains this
actual compressed covariance.

If $Z_d$ is any orthonormal basis of $\ker V_d$, completing $W_d$
to a unitary matrix, the further exact determinant factorization is
\[
 \det K_{A,N}
   =\det(A_NA_N^*)\,
        \det\bigl(Z_d^*(A_NA_N^*)^{-1}Z_d\bigr).
 \tag{CEP25}
\]
For a proof put $D=A_NA_N^*>0$ and write its block matrix in
the unitary frame $[W_d,Z_d]$. The upper left block is $K_{A,N}$.
Block Gaussian elimination gives
$\det D=\det K_{A,N}\det(D_{22}-D_{21}K_{A,N}^{-1}D_{12})$.
The lower right block of $D^{-1}$ is the inverse of that Schur
complement, proving CEP25. The empty-complement case has determinant
one. The dimension of this retained complement is
$N-v+1-q'$.

Choose any fixed full frame $I_V:\mathbb C^\Delta\hookrightarrow E_k$
of $\ker C$ and a fixed right section $L_C$ of $C$.
Then $F_V=[I_V,L_C]$ is invertible and
\[
 \det(I_V^*G_NI_V)
   =|\det F_V|^2\,\frac{\det G_N}{\det Q_N}
   =|\det F_V|^2\,\frac{\det G_N}{\det G'_d}\det K_{A,N}.
 \tag{CEP26}
\]
Indeed $CF_V=[0,I]$, so the lower right block of
$(F_V^*G_NF_V)^{-1}$ is $CG_N^{-1}C^*=Q_N^{-1}$.
The Schur determinant argument just used gives CEP26.
The frame factor is independent of source cutoff and cancels
only after applying the four signs.

For the actual inclusion $K\subset V_A$, choose its original
coefficient frame, and let $\Xi:\mathbb C^\Delta\twoheadrightarrow
\mathbb C^{\Delta-r}$ have its coordinate kernel, $r=8k-16$,
as in AKJ18. With $H_{V,N}=I_V^*G_NI_V$, the quotient metric is
$Q_{\Xi,N}=(\Xi H_{V,N}^{-1}\Xi^*)^{-1}$.
The same proof with a frame of $\ker\Xi$ and a right section gives,
with all constant frame factors cancelled by the endpoint signs,
\[
 F_K^{(s)}
 =\sum_N\sigma_N\log\det G_N
   -\sum_N\sigma_N\log\det G'_{N-v}
   +\sum_N\sigma_N\log\det K_{A,N}
   -\sum_N\sigma_N\log\det Q_{\Xi,N}.
 \tag{CEP27}
\]
Thus the conductor quotient and the remaining original observation
quotient have both been retained in the determinant relation.
The lower source cutoffs in this equation are exactly
$q-1-v,q-v,2q-1-v,2q-v$. Relative to the lower grid's own
four cutoffs they have the respective shifts
$\Delta-v,\Delta-v,2\Delta-v,2\Delta-v$.

\subsection{An exact graph quotient with controlled exterior metric error}

In original orthonormal coordinates define
\[
 S_N:\mathcal P_N\longrightarrow
         \mathbb C^v\oplus\mathcal P'_{N-v},\qquad
 S_Np=(P_{<v}p,\mathcal T_{A,N}p).
 \tag{CEP28}
\]
Here $P_{<v}$ records the first $v$ original upper orthonormal
coefficients, and the target carries the orthogonal direct-sum norm
with those coordinates and the unchanged lower Gamma source.
The first $v$ source coordinates span the polynomials of degree
below $v$. Thus CEP14--16 give its exact block matrix
$S_N=\operatorname{diag}(I_v,\widehat A_N)$ in the ordered
source split at degree $v$. In particular it is invertible, and
\[
 \begin{gathered}
 \|S_N\|\le M_N:=\max\{1,M_{A,N}\},\\
 |\det S_N|=\left|\frac{\mu_v}{v!}\right|^{N-v+1}
              \prod_{m=0}^{N-v}\sqrt{\frac{w_{m+v}}{w_m}},\\
 |\log|\det S_N||
 \le L_N:=(N-v+1)\left|\log\left|\frac{\mu_v}{v!}\right|\right|
                      +v c_b^{\log}H_N .
 \end{gathered}\tag{CEP29}
\]
The last bound follows by the consecutive-product identity preceding
CEP19 and $|\log w_j|\le c_b^{\log}H_j$: the two lists of $v$
logarithms are each bounded by $v c_b^{\log}H_N$, followed by
the factor $1/2$.

The upper and lower source ideals at this cutoff are
\[
 I_N=\chi_k\mathbb C[S]_{\le N-q},\qquad
 I'_N=\chi_{k'}\mathbb C[S']_{\le N-v-q'} .
 \tag{CEP30}
\]
A negative degree specifies the zero subspace.
Let $\mathcal W_N=U_N^{-1}(V_A)$. The exact image and ideal graph
under CEP28 are
\[
 \begin{gathered}
 S_N\mathcal W_N=\mathbb C^v\oplus I'_N,\\
 S_N I_N=\{(P_{<v}h,\mathcal T_{A,N}h):h\in I_N\}
                    \subset\mathbb C^v\oplus I'_N,\\
 V_A\ \xrightarrow{\ \simeq\ }\
   (\mathbb C^v\oplus I'_N)/(S_N I_N),\qquad
 U_Np\longmapsto[S_Np].
 \end{gathered}\tag{CEP31}
\]
The first equality follows because $\mathcal W_N$ is the preimage
of $I'_N$ under the top map of CEP6, and $S_N$ is invertible.
CEP3 proves the ideal inclusion. If representatives differ by
$I_N$, their images differ by $S_NI_N$, so the last map is
well defined. Conversely every class in the target has a preimage
under $S_N$ in $\mathcal W_N$, and the kernel consists exactly
of $I_N$. This proves the isomorphism. Its norm on the left is
the restriction of $G_N$: all lifts of a member of $V_A$ already
belong to $\mathcal W_N$, so their minimum source norm is the
quotient norm of $\mathcal W_N/I_N$.
The low-degree graph component $P_{<v}h$ is explicitly retained.

For reference the ideal dimensions at the four original endpoints are
\[
 \begin{array}{c|c|c|c}
 N&\dim I_N&\dim I'_N&
       \dim((\mathbb C^v\oplus I'_N)/S_NI_N)\\ \hline
 q-1&0&\Delta-v&\Delta\\
 q&1&\Delta+1-v&\Delta\\
 2q-1&q&q+\Delta-v&\Delta\\
 2q&q+1&q+\Delta+1-v&\Delta .
 \end{array}\tag{CEP32}
\]
These follow by counting monomials in CEP30; $S_N$ preserves
dimensions. Moreover $\mathcal T_A|_{I_N}$ is injective because
a nonzero member has degree at least $q>v$ and its image has
degree exactly $v$ less. Hence the second coordinate of the
retained ideal graph already has dimension $\dim I_N$.

Here is the complete exterior norm estimate for this exact map.
For an invertible complex matrix $S$ of dimension $n$ and
$0\le t\le n$, its singular values
$\sigma_1\ge\cdots\ge\sigma_n>0$ give
\[
 \|\bigwedge^t S^{-1}\|
 =\frac{\|\bigwedge^{n-t}S\|}{|\det S|}.
 \tag{CEP33}
\]
Indeed the numerator is $\sigma_1\cdots\sigma_{n-t}$ and
the denominator is the product of all singular values; their
ratio is the product of the $t$ largest singular values of
$S^{-1}$. This includes $t=0,n$.
For $n=N+1$, CEP29 therefore gives the finite bounds
\[
 \begin{split}
 \log^+\|\bigwedge^t S_N\|&\le t\log M_N,\\
 \log^+\|\bigwedge^t S_N^{-1}\|
        &\le(N+1-t)\log M_N+L_N .
 \end{split}\tag{CEP34}
\]
Here $\log^+x=\max\{0,\log x\}$. The right sides are nonnegative.
Thus for all these ranks both bounds are at most
\[
 E_N=(N+1)\log M_N+L_N.
 \tag{CEP35}
\]
With $R_{\rm sh}=\max_{r,s}|b_{rs}-4|$ and
$A_*^{\rm abs}=\sum|a_{rs}|$, equation CEP17 gives
$\log M_N\le\log^+A_*^{\rm abs}+R_{\rm sh}(H_N+2)$.
For the fixed original quartet and period, $v,R_{\rm sh},\mu_v$
and $A_*^{\rm abs}$ are fixed while $c_b^{\log}=O(k)$.
Since $N\le2q$ and $H_N\le1+\log N$, all four explicit bounds
in CEP35 satisfy
\[
 E_N=O_{\rm actual}(q\log(q+1)+k\log(q+1))=o(kq).
 \tag{CEP36}
\]
This uses only the forward norm and the exact determinant of
the original conductor map.

We now prove the exact relation to the metric of the original
kernel, including the quotient projection. In a finite-dimensional
Hilbert space let $I\subset W$ and let $S$ be invertible.
The induced map $W/I\longrightarrow SW/SI$ sends a minimal lift
$x\in W\cap I^\perp$ to the orthogonal projection of $Sx$
off $SI$. Orthogonal projection is a contraction on every exterior
power: in an orthonormal basis its singular values are zeros and
ones. Therefore, on any $t$-dimensional subspace of $W/I$, the
exterior norm of the induced map is at most
$\|\bigwedge^t S\|$. Applying the identical argument to $S^{-1}$
gives the inverse bound. The minimal-lift identification is
isometric since every other lift differs by a perpendicular
element of $I$, whose squared norm adds. This proves the
quotient assertion with no choice of coefficient norm.

Apply this result with $W=\mathcal W_N$, $I=I_N$, and the
actual rank-$r$ subspace $K\subset V_A$. The target subspace is
the explicit image of $K$ under CEP31; denote its Gram matrix
in the transported original frame by $\widetilde H_{K,N}$.
Let $I_K:\mathbb C^r\hookrightarrow E_k$ be that original
primary frame, and write its original minimum source lift as
$Y_N=U_N^*G_N I_K$.
Let $Z_N$ be the matrix in the original orthonormal source
coordinates of the independent columns
$\chi_k(S)S^a$, $0\le a\le N-q$. If there are no columns,
all following products involving them are zero.
The exact target orthogonal projection and Gram matrix are
\[
 \begin{gathered}
 P_N^{\rm graph}
  =I-S_NZ_N(Z_N^*S_N^*S_NZ_N)^{-1}Z_N^*S_N^*,\\
 \widetilde H_{K,N}
       =Y_N^*S_N^*P_N^{\rm graph}S_NY_N,\qquad
 H_{K,N}=Y_N^*Y_N=I_K^*G_NI_K .
 \end{gathered}\tag{CEP37}
\]
The inverted Gram matrix is positive definite because its
columns are independent and $S_N$ is invertible. Multiplication
shows that $P_N^{\rm graph}$ is Hermitian, idempotent, kills
$S_NI_N$ and fixes its orthogonal complement. Thus CEP37 is
precisely the minimum quotient Gram of CEP31, and its kernel
frame has full rank because CEP31 is injective.
Replacing the lift $Y_N$ by any lift of the same primary frame
adds columns in $I_N$, which the projection kills; hence every
choice of lift gives this same matrix.

Applying the proved exterior quotient estimate and squaring
volumes gives
\[
 \begin{split}
 -2\log\|\bigwedge^r S_N^{-1}\|
 &\le \log\det\widetilde H_{K,N}-\log\det H_{K,N}\\
 &\le2\log\|\bigwedge^r S_N\|,\\
 \left|\log\det\widetilde H_{K,N}-\log\det H_{K,N}\right|
 &\le2E_N .
 \end{split}\tag{CEP38}
\]
In detail, divide the chosen frame by its positive Gram square
root so it is orthonormal in the source quotient. The squared
norm of its exterior image is exactly the determinant ratio.
The forward and inverse exterior estimates prove the two
inequalities, and CEP34--35 prove the last assertion.

Consequently the following reduction of the actual four-endpoint
kernel determinant is proved:
\[
 \begin{gathered}
 \widetilde F_K^{(s)}
       =\sum_N\sigma_N\log\det\widetilde H_{K,N},\\
 |F_K^{(s)}-\widetilde F_K^{(s)}|
       \le2\sum_{N\in\{q-1,q,2q-1,2q\}}E_N=o(kq).
 \end{gathered}\tag{CEP39}
\]
The frame is the original fixed frame of $K$ at every endpoint;
CEP37 gives all four transported matrices explicitly.
Their target spaces retain the original ideal graphs, the
different source cutoffs, all conductor coefficients, and the
full original Gamma masses. CEP39 proves that the leading
$kq$ coefficient is unchanged by this particular graph-quotient
transport. It supplies no value for that coefficient: its
remaining determinant is the displayed, completely specified
four-endpoint graph quotient.

\subsection{All ideal graph entries as original polynomial products}

Each individual shift in CEP2 induces the unital algebra map
$\sigma_{rs}:E_k\to E_{k'}$, $[p]\mapsto[p(S'+b_{rs})]$.
The divisibility in CEP3, applied to the individual shifted
polynomial, proves independence of the representative; substitution
preserves multiplication and one. Thus
$\overline{\mathcal T}_A=\sum a_{rs}\sigma_{rs}$ is its exact
linear combination. For any $x,y\in E_k$, the complete
multiplicative defect of that linear combination is
\[
 \overline{\mathcal T}_A(xy)
   -\overline{\mathcal T}_A(x)\overline{\mathcal T}_A(y)
 =\sum_{r,s}a_{rs}\sigma_{rs}(x)\sigma_{rs}(y)
   -\sum_{r,s,t,u}a_{rs}a_{tu}
                        \sigma_{rs}(x)\sigma_{tu}(y).
 \tag{CEP40}
\]
This follows by applying the proved individual multiplicativity
and expanding the product of the two complete sums.

There is a complete finite product for the map on the source ideal.
For each original shift put
\[
 d_{rs}(S')=
 \prod_{\substack{0\le i,j\le k\\
  (i,j)\notin\{r,\ldots,r+k'\}\times\{s,\ldots,s+k'\}}}
                   (S'+b_{rs}-\beta_{ij}).
 \tag{CEP41}
\]
This is monic of degree $\Delta$, and
$\chi_k(S'+b_{rs})=\chi_{k'}(S')d_{rs}(S')$.
Indeed CEP3's root addition identifies exactly the $q'$ factors
in the indicated rectangle with the factors of $\chi_{k'}$.
All remaining factors are precisely the product in CEP41.
Consequently for every $h\in\mathbb C[S]_{\le L}$,
$L=N-q$, the exact ideal map is
\[
 \begin{split}
 \mathcal T_{A,N}(\chi_kh)&=\chi_{k'}\,\mathcal U_{A,L}h,\\
 (\mathcal U_{A,L}h)(S')&=\sum_{r,s}a_{rs}
                     d_{rs}(S')h(S'+b_{rs})
                  \in\mathbb C[S']_{\le L+\Delta-v}.
 \end{split}\tag{CEP42}
\]
The equality follows by the product identity just proved. The
degree statement follows from CEP8: if $h$ has degree $t$ and
leading coefficient $h_t\ne0$, the leading coefficient of its
displayed image is
$h_t\binom{q+t}{v}\mu_v$ at degree $t+\Delta-v$.
This follows because $\chi_k,\chi_{k'}$ are monic and
$\mathcal T_A$ lowers degree exactly by $v$ with its stated
leading derivative. Thus $\mathcal U_{A,L}$ is injective.
At $L=-1$ the domain and image are zero.

This gives a further exact sequence, with all maps specified:
\[
 0\longrightarrow\mathbb C[S]_{<v}
 \longrightarrow V_A
 \longrightarrow
 \mathbb C[S']_{\le L+\Delta-v}/
                     \mathcal U_{A,L}\mathbb C[S]_{\le L}
 \longrightarrow0 .
 \tag{CEP43}
\]
The first map is the original quotient class. For the second,
lift a member of $V_A$ to $p\in\mathcal W_N$, divide
$\mathcal T_Ap$ by $\chi_{k'}$, and take its class in the last
space. A change of lift by $\chi_kh$ adds exactly
$\mathcal U_{A,L}h$ by CEP42, proving independence. Every target
polynomial times $\chi_{k'}$ has a preimage under the onto
map $\mathcal T_{A,N}$, proving surjectivity. If the image is
zero, subtract the corresponding $\chi_kh$; the remaining
polynomial is killed by $\mathcal T_A$ and hence has degree
below $v$. Its class injects because $v<q$.
This proves exactness. The last space has dimension $\Delta-v$
at every original endpoint, as is also obtained by subtracting
the two polynomial dimensions.

For a literal matrix version of the ideal graph, equip
$\mathbb C[S']_{\le L+\Delta-v}$ with the ordered lower
orthonormal Gamma coefficient basis, and write $M_{\chi',N}$
for multiplication by the full original $\chi_{k'}$ into
$\mathcal P'_{N-v}$. It is injective. The map
\[
 D_{\chi',N}=\operatorname{diag}(I_v,M_{\chi',N})
 :\mathbb C^v\oplus\mathbb C[S']_{\le L+\Delta-v}
       \longrightarrow\mathbb C^v\oplus I'_N
 \tag{CEP44}
\]
is an isomorphism onto the indicated subspace. Its pullback metric
is exactly
$G_{\chi',N}=\operatorname{diag}(I_v,M_{\chi',N}^*M_{\chi',N})$;
this follows by taking the squared norm of each component.
For the domain polynomial coefficient frame of $h$, let
$Z_{\chi,N}h$ be the original upper orthonormal coefficients
of $\chi_kh$. In these coordinates the complete graph relation
matrix is
\[
 \mathcal Z_{A,N}=
 \begin{pmatrix}
 P_{<v}Z_{\chi,N}\\[2pt]
 [\mathcal U_{A,L}]
 \end{pmatrix},\qquad
 D_{\chi',N}\mathcal Z_{A,N}=S_N Z_{\chi,N}.
 \tag{CEP45}
\]
The first component is the original low-degree projection,
and the second consists of the exact finite products and
translations in CEP41--42. The equality is CEP28 and CEP42.
It proves that the graph quotient in CEP31 is precisely the
quotient by $\operatorname{im}\mathcal Z_{A,N}$ in the
displayed metric $G_{\chi',N}$. All matrix inverses and
minimum projections in CEP37 can therefore be computed using
these full original polynomial products. No coefficient of
the ideal, source metric or low-degree graph component is
discarded by this additional coordinate isomorphism.

\subsection{The compressed conductor correction is subleading at every endpoint}

The exact exterior estimate also controls the large compressed
covariance in CEP24 without multiplying a scalar inverse bound
by $q'$. Write $n'=N-v+1$ and retain
$A_N=[0,\widehat A_N]$ with its exact square block in CEP16.
Since $W_d^*W_d=I_{q'}$, the coefficient matrix of its image under
$\widehat A_N^*$ gives
\[
 K_{A,N}
   =(\widehat A_N^*W_d)^*(\widehat A_N^*W_d).
 \tag{CEP46}
\]
The squared norm of the exterior product of the $q'$ columns
on the right is its determinant. The norm of the exterior
product of the columns of $W_d$ is one, because those columns
are orthonormal. The forward estimate therefore gives
$\det K_{A,N}\le\|\widehat A_N\|^{2q'}\le M_N^{2q'}$.
For the reverse inequality apply
$\bigwedge^{q'}\widehat A_N^{-*}$ to the same image columns;
their image is the orthonormal column wedge of $W_d$. Thus
\[
 \det K_{A,N}\ge
      \|\bigwedge^{q'}\widehat A_N^{-1}\|^{-2}.
 \tag{CEP47}
\]
The singular values of an adjoint are the same, so the adjoint
inverse in this argument has exactly the displayed exterior norm.
The singular complement identity CEP33 in dimension $n'$
gives
$\|\bigwedge^{q'}\widehat A_N^{-1}\|
 \le M_N^{n'-q'}/|\det\widehat A_N|$.
CEP29 bounds the absolute logarithm of the last determinant
by $L_N$. Consequently the complete two-sided estimate is
\[
 \begin{gathered}
 -2\bigl((N-v+1-q')\log M_N+L_N\bigr)
       \le\log\det K_{A,N}\le2q'\log M_N,\\
 |\log\det K_{A,N}|\le2E_N=o(kq)
 \quad\hbox{at every original endpoint and source order.}
 \end{gathered}\tag{CEP48}
\]
All ranks in this argument are their actual ranks,
$q'\le N-v+1$ as proved after CEP6. In particular the estimate
applies to the complete original lower quotient, whose rank
grows quadratically.

Define the fully specified total-source difference and the
remaining original observation quotient return by
\[
 \begin{split}
 \mathcal T_k^{(s)}
   &=\sum_N\sigma_N
          \bigl(\log\det G_N-\log\det G'_{N-v}\bigr),\\
 F_{\Xi,k}^{(s)}&=\sum_N\sigma_N\log\det Q_{\Xi,N},\\
 \mathcal E_{A,k}^{(s)}
   &=\sum_N\sigma_N\log\det K_{A,N}.
 \end{split}
\]
Equations CEP26--27 and CEP48 prove the exact identities and bounds
\[
 \begin{gathered}
 F_{V_A}^{(s)}
      =\mathcal T_k^{(s)}+\mathcal E_{A,k}^{(s)},\qquad
 F_K^{(s)}
      =\mathcal T_k^{(s)}-F_{\Xi,k}^{(s)}
                              +\mathcal E_{A,k}^{(s)},\\
 |\mathcal E_{A,k}^{(s)}|\le2\sum_N E_N=o(kq),\\
 \frac{F_K^{(s)}-\mathcal T_k^{(s)}+F_{\Xi,k}^{(s)}}{kq}
        \longrightarrow0 .
 \end{gathered}\tag{CEP49}
\]
Here $F_{V_A}^{(s)}$ uses its fixed original coefficient frame,
and $Q_{\Xi,N}$ is the exact quotient metric of CEP27.
The frame factors cancel with coefficients $1+1-1-1$ only in
these returns, as proved in CEP26. Thus the leading metric
effect of the full conductor source transport has now been
controlled. The remaining rank-$8k-32$ observation quotient
and the shifted lower-grid source determinants are still their
explicit original matrices; CEP49 carries them forward into
the next calculation.

% END CURRENT CEP


\clearpage
% BEGIN CURRENT CTV
% Exact total-source variation in the original root polynomials and Gamma norms.
\section{The total-source difference as an outer-strip quotient}
\label{sec:CTV}

Retain CEP1--49, with $m=1$, $k=4l+1\ge9$,
$q=(k+1)^2$, $k'=k-8$, $q'=(k-7)^2$,
$\Delta=q-q'=16k-48$, and the original fixed
$0<\delta<1/2$, $\gamma>2$. The conductor has its actual
first nonzero moment index $0\le v\le80$. At each of the four
original endpoints $N\in\{q-1,q,2q-1,2q\}$ put
$t=N-q+1$ and $d=N-v$. The source order remains
$s\in\{1,k\}$ at both grids, with
\[
 b=s/2,\quad M_s=(2\pi)^{s/2},\quad
 dm_s(y)=M_s\frac{2^b}{4\pi\Gamma(b)}
             |\Gamma(b/2+iy/2)|^2\,dy,\quad
 h_n=M_s n!(b)_n .
 \tag{CTV1}
\]
These are the complete original measure and monic norms of
OKM6--7. The source centres are $c=k/2$ and $c'=k/2-4$.
The total-source difference in CEP49 is
\[
 \mathcal T_k^{(s)}
   =\sum_N\sigma_N\bigl(\log\det G_N-\log\det G'_{N-v}\bigr),
 \qquad \sigma=(1,1,-1,-1),
 \tag{CTV2}
\]
where both matrices use their original primary frames.
All calculations below retain these four cutoffs and the same
two source orders.

\subsection{The exact nested root polynomials in the original real coordinate}

Define the full monic polynomials
\[
 \begin{gathered}
 P_k(y)=i^{-q}\chi_k(c+iy)
       =\prod_{a,b=0}^k(y-\omega_{ab}),\qquad
 \omega_{ab}=(2b-k)\gamma-i(2a-k)\delta,\\
 P_{k'}(y)=i^{-q'}\chi_{k'}(c'+iy)
       =\prod_{a,b=0}^{k'}(y-\omega'_{ab}),\qquad
 \omega'_{ab}=(2b-k')\gamma-i(2a-k')\delta .
 \end{gathered}\tag{CTV3}
\]
The original phases $i^{-q}$ and $i^{-q'}$ have modulus one;
they occur exactly as written in the coefficient maps.
For every lower pair one has
$\omega'_{ab}=\omega_{a+4,b+4}$. Hence, with the complete
outer index set
$\mathcal O_k=\{0,\ldots,k\}^2\setminus\{4,\ldots,k-4\}^2$,
the exact factorization is
\[
 P_k(y)=P_{k'}(y)P_{\partial,k}(y),\qquad
 P_{\partial,k}(y)=\prod_{(a,b)\in\mathcal O_k}(y-\omega_{ab}),
 \qquad \deg P_{\partial,k}=\Delta .
 \tag{CTV4}
\]
Direct substitution proves the root equality; counting the retained
central rectangle gives $q'$ roots, leaving $\Delta$ distinct
outer roots. The products are monic, so the root equality fixes
the factor without any additional constant.

Since $k'$ is odd, $|2a-k'|\ge1$ for every lower index.
For real $y$ this proves, factor by factor,
\[
 |y-\omega'_{ab}|\ge|\operatorname{Im}\omega'_{ab}|
        \ge\delta,\qquad
 |P_{k'}(y)|\ge\delta^{q'}.
 \tag{CTV5}
\]
Also $|\omega'_{ab}|\le
R'=k'\sqrt{\delta^2+\gamma^2}$. These estimates keep every
original factor; no root has been moved.

\subsection{The full monic source and relation determinant identities}

For a nonzero polynomial $P$ and an integer $a\ge0$ define the relation
moment determinant
\[
 D_a(P;s)=
 \det\left[\int_{\mathbb R}y^{i+j}|P(y)|^2\,dm_s(y)
                                    \right]_{0\le i,j<a},
 \qquad D_0(P;s)=1 .
 \tag{CTV6}
\]
Every nonempty matrix is positive definite: the measure is strictly
positive on the real line, a nonzero polynomial has only finitely
many zeros, and the squared modulus of its product with $P$
is positive on an open interval. All moments exist because the
original source has its finite exponential moments in CTV14 below.

In increasing monomial $y$ coordinates, the quotient by a monic
polynomial $P$ of degree $a$ from source degree $L\ge a-1$
has determinant
\[
 \det G^{y}_{P,L}
   =\frac{\prod_{n=0}^L h_n}{D_{L-a+1}(P;s)} .
 \tag{CTV7}
\]
Here the quotient frame is $1,y,\ldots,y^{a-1}$.
For proof use the source frame
\[(1,y,\ldots,y^{a-1},P,yP,\ldots,y^{L-a}P).\]
Its increasing-degree coefficient matrix has diagonal entries one
and no entries above the degree of a column, hence determinant
one. The last block is exactly the relation Gram in CTV6.
Minimizing over those last coordinates gives its Schur complement,
whose determinant is the full source determinant divided by the
relation determinant. Monic orthogonalization by the original
$p_n$ also has determinant one, so that full determinant is
$\prod_{n=0}^L h_n$. This proves CTV7 including the empty relation
case. It is the full-source analogue of BSL3, with the original
Gamma order retained.

Let $V_\beta,V_{\beta'}$ be the original ordered Vandermonde
matrices in the $S,S'$ variables. The exact coordinate pullback
$S=c+iy$ sends $S^j$ to $(c+iy)^j$ and has diagonal coefficients
$i^j$; its determinant is $i^{a(a-1)/2}$ in degree below $a$.
Its modulus is one, so its Gram congruence has determinant factor
one. The same statement holds with $c'$. Evaluation in primary
coordinates then gives
\[
 \begin{split}
 \det G_N
  &=\frac{\prod_{n=0}^N h_n}
          {D_t(P_k;s)\,|\det V_\beta|^2},\\
 \det G'_{N-v}
  &=\frac{\prod_{n=0}^{N-v}h_n}
     {D_{t+\Delta-v}(P_{k'};s)\,|\det V_{\beta'}|^2}.
 \end{split}\tag{CTV8}
\]
The relation phase $i^q$ or $i^{q'}$ multiplying each source
relation column has squared modulus one, as used in this equality.
The Vandermonde factors are nonzero by the original distinctness
of all roots. Their logarithmic difference is constant across
the four endpoints and has signed coefficient zero. Therefore
CTV8 proves the exact finite expression
\[
 \mathcal T_k^{(s)}
  =\sum_N\sigma_N\left[
      \sum_{n=N-v+1}^{N}\log h_n
      +\log D_{t+\Delta-v}(P_{k'};s)
      -\log D_t(P_k;s)\right].
 \tag{CTV9}
\]
The first sum is empty for $v=0$. In particular both low relation
determinants and both high relation determinants remain present.

\subsection{The exact quotient of the outer-strip polynomial}

Introduce the positive measure
\[
 d\zeta_{k',s}(y)=|P_{k'}(y)|^2\,dm_s(y).
 \tag{CTV10}
\]
This measure is the exact pullback of the original source metric
under the multiplication map
$\mathcal M_{P_{k'}}:f\mapsto P_{k'}f$.
Indeed $\|\mathcal M_{P_{k'}}f\|_{m_s}^2
=\int|f|^2d\zeta_{k',s}$, including the same full factor $M_s$.
On polynomials of degree at most $N-q'$ its image is the specified
subspace $P_{k'}\mathbb C[y]_{\le N-q'}$ of the original
degree-at-most-$N$ source. The map is injective because $P_{k'}$
is nonzero. Thus CTV10 records a proved metric map to the
original source, rather than an assigned new Gamma order.

For $L=N-q'=t+\Delta-1$, let
$Q_{\partial,L}^{(s)}$ be the minimum quotient Gram of
\[
 \mathbb C[y]_{\le L}\longrightarrow
       \mathbb C[y]/(P_{\partial,k})
 \tag{CTV11}
\]
in the source norm CTV10 and in the increasing monomial
remainder frame of dimension $\Delta$.
The source degree is at least $\Delta-1$ at every endpoint.
The same monic-frame Schur proof as in CTV7, now using CTV4,
gives
\[
 \det Q_{\partial,L}^{(s)}
     =\frac{D_{t+\Delta}(P_{k'};s)}{D_t(P_k;s)} .
 \tag{CTV12}
\]
In fact the ideal columns are $P_{\partial,k}y^j$, $0\le j<t$;
their Gram in CTV10 is exactly $D_t(P_k;s)$ by CTV4.
The numerator is the Gram of the complete monomial source
of dimension $L+1=t+\Delta$. This proves every source and
relation factor in CTV12.

Write $\lambda_j^{(k',s)}$ for the squared norm of the monic
orthogonal polynomial of degree $j$ in CTV10. This polynomial
and its norm are completely determined by the positive moment
matrices: subtract from $y^j$ its projection onto lower powers.
If $H_j$ is that lower Gram and $a_j$ its inner-product column
with $y^j$, the lower coefficients are $-H_j^{-1}a_j$.
Completing the square proves minimality among all monic
polynomials of degree $j$; its squared norm is the Schur
complement. Thus
$\lambda_j^{(k',s)}=D_{j+1}(P_{k'};s)/D_j(P_{k'};s)>0$.
Multiplying these identities and substituting CTV12 into CTV9
gives
\[
 \begin{gathered}
 \mathcal F_{\partial,k}^{(s)}
    =\log\det Q_{\partial,\Delta-1}^{(s)}
       +\log\det Q_{\partial,\Delta}^{(s)}
       -\log\det Q_{\partial,q+\Delta-1}^{(s)}
       -\log\det Q_{\partial,q+\Delta}^{(s)},\\
 \mathcal T_k^{(s)}
    =\mathcal F_{\partial,k}^{(s)}
                        +\sum_N\sigma_N\mathcal R_{N,v}^{(s)},\\
 \mathcal R_{N,v}^{(s)}
    =\sum_{a=0}^{v-1}
          \left(\log h_{N-a}
                -\log\lambda_{N-q'-a}^{(k',s)}\right).
 \end{gathered}\tag{CTV13}
\]
Every weighted norm index is nonnegative:
$N-q'-(v-1)\ge\Delta-v\ge16$.
The total difference therefore has an exact rank-$\Delta$
quotient presentation with precisely the displayed norm tails.

\subsection{Finite bounds for all norm tails in both original sources}

The original Gamma generating identity gives the full exponential
moment
$\int e^{t y}\,dm_s(y)=M_s(\cos t)^{-b}$ for
$|t|<\pi/2$. By the evenness of the measure this also equals
$\int\cosh(t y)\,dm_s(y)$.
The nonnegative Taylor terms of $\cosh$ imply
\[
 \begin{gathered}
 \int y^{2j}\,dm_s(y)
  \le M_s(\cos t_0)^{-b}\frac{(2j)!}{t_0^{2j}},\\
 \|y^j\|_{m_s}
  \le\sqrt{M_s}(\cos t_0)^{-b/2}\mathfrak c_{\rm mom}^j j!,\qquad
 t_0=\frac\pi4,\quad \mathfrak c_{\rm mom}=\frac2{t_0}=\frac8\pi .
 \end{gathered}\tag{CTV14}
\]
The second inequality uses
$(2j)!\le4^j(j!)^2$, proved by
$\binom{2j}{j}\le\sum_{a=0}^{2j}\binom{2j}{a}=4^j$.
Thus every constant in this moment estimate includes the
original source mass and order.

For a monic polynomial of degree $n$ whose roots have modulus
at most $R'$, its coefficient of $y^a$ has absolute value at most
$\binom na(R')^{n-a}$: it is the sum of those products of
$n-a$ distinct roots. The triangle inequality and CTV14 therefore
give the full bound
\[
 \|P\|_{m_s}
  \le\sqrt{M_s}(\cos t_0)^{-b/2}
          \mathfrak c_{\rm mom}^n n!\sum_{h=0}^n
                              \frac{(R'/\mathfrak c_{\rm mom})^h}{h!}
  \le\sqrt{M_s}(\cos t_0)^{-b/2}
          \mathfrak c_{\rm mom}^n n!\,e^{R'/\mathfrak c_{\rm mom}}.
 \tag{CTV15}
\]
The monic test polynomial $y^j$ in the minimization defining
$\lambda_j^{(k',s)}$ gives the polynomial
$P_{k'}(y)y^j$ in CTV15, of degree $n=q'+j$; its extra
$j$ roots are exactly zero and also satisfy the stated bound.
For the lower bound, CTV5 applies to every competing monic
polynomial, whose original Gamma squared norm is at least
$h_j$. Taking the minimum proves
\[
 \delta^{2q'}h_j
 \le\lambda_j^{(k',s)}
 \le M_s(\cos t_0)^{-b}\mathfrak c_{\rm mom}^{2(q'+j)}
                ((q'+j)!)^2e^{2R'/\mathfrak c_{\rm mom}}.
 \tag{CTV16}
\]
No lower bound on an actual period coefficient is used here.
The lower root polynomial itself supplies the strictly positive
real-line lower bound.

For each term of CTV13 put $n=N-a$ and $j=n-q'$.
Subtracting the logarithms in CTV16 from the exact
$\log h_n=\log M_s+\log n!+\log(b)_n$ yields
\[
 \begin{split}
 b\log(\cos t_0)+\log\frac{(b)_n}{n!}
          -2n\log\mathfrak c_{\rm mom}
                         -\frac{2R'}{\mathfrak c_{\rm mom}}
 &\le \log h_n-\log\lambda_j^{(k',s)}\\
 &\le
   \sum_{r=j+1}^{n}\bigl(\log r+\log(b+r-1)\bigr)
                          -2q'\log\delta .
 \end{split}\tag{CTV17}
\]
The two full $\log M_s$ contributions cancel in this displayed
comparison because both norms arise from exactly CTV1.
They have not been removed from either individual norm.

Let $c_b^{\log}=\max\{1,b-1\}$. As proved directly in CEP19,
$|\log(j/(b+j-1))|\le c_b^{\log}/j$ for $b\ge1/2$:
use $\log(1+x)\le x$ for $b\ge1$; for $1/2\le b<1$
use $-\log(1-a/j)\le a/(j-a)\le1/j$, $a=1-b$.
Summation gives
$|\log((b)_n/n!)|\le c_b^{\log}H_n$.
With all original endpoints $N\ge99$, define the explicit
nonnegative numbers
\[
 \begin{split}
 L_N^{\rm tail}
    &=c_b^{\log}H_N+\frac b2\log2
                        +2N\log(8/\pi)+\frac{\pi R'}4,\\
 U_N^{\rm tail}
    &=q'\bigl(\log N+\log(b+N-1)+2|\log\delta|\bigr),\\
 C_{N,v}^{\rm tail}
    &=v\max\{L_N^{\rm tail},U_N^{\rm tail}\}.
 \end{split}\tag{CTV18}
\]
Indeed $\cos(\pi/4)=2^{-1/2}$ and
$2R'/\mathfrak c_{\rm mom}=\pi R'/4$.
In the upper sum of CTV17 there are exactly $q'$ terms, and
each is at most $\log N+\log(b+N-1)$.
Thus CTV17 proves
\[
 -vL_N^{\rm tail}\le\mathcal R_{N,v}^{(s)}
               \le vU_N^{\rm tail},\qquad
 |\mathcal R_{N,v}^{(s)}|\le C_{N,v}^{\rm tail}.
 \tag{CTV19}
\]
When $v=0$ the equality $\mathcal R_{N,0}^{(s)}=0$ and the
zero bound follow from the empty sum.
These estimates apply for all actual endpoints and for
both unchanged source orders simultaneously.

\subsection{The proved first-variation reduction and its remaining matrix}

Since $N\le2q$, $q'\le q$, $b\le k/2$ and
$R'\le k\sqrt{\delta^2+\gamma^2}$, equations CTV18--19 give
$C_{N,v}^{\rm tail}=O_{\delta,\gamma}(q\log(q+1)+
k\log(q+1))$ uniformly for every integer $0\le v\le80$.
To see this directly, use $H_N\le1+\log N$ in the first
line of CTV18 and the displayed upper bounds in its second;
$v$ is bounded by the fixed integer 80.
Dividing by $kq$ tends to zero, since
$\log(q+1)/k\to0$ for $q=(k+1)^2$.
Consequently CTV13 and the finite bounds just proved yield
\[
 \left|\mathcal T_k^{(s)}-\mathcal F_{\partial,k}^{(s)}\right|
        \le\sum_N C_{N,v}^{\rm tail}=o(kq).
 \tag{CTV20}
\]
The estimate is uniform over the possible actual moment indices
$v$; it requires no bound on the remaining period coordinates.
Every polynomial and measure defining
$\mathcal F_{\partial,k}^{(s)}$ is independent of that index.

For an explicit matrix target, let
$H_{L}^{\zeta}$ be the moment Gram of
$1,y,\ldots,y^L$ in the full measure CTV10, and let
$\mathsf R_{\partial,L}$ be the literal remainder matrix
modulo the full polynomial $P_{\partial,k}$ in CTV4.
The exact metric is
\[
 Q_{\partial,L}^{(s)}
  =\bigl(\mathsf R_{\partial,L}
          (H_{L}^{\zeta})^{-1}
          \mathsf R_{\partial,L}^*\bigr)^{-1}.
 \tag{CTV21}
\]
The map is onto because its first $\Delta$ columns are the
identity in the remainder frame. The minimum source lift
of a coefficient vector $x$ is
$(H_L^\zeta)^{-1}\mathsf R_{\partial,L}^*
 Q_{\partial,L}^{(s)}x$; it maps to $x$ and is perpendicular
to the kernel. Adding any other lift adds a perpendicular
kernel vector, proving the stated Gram. This also proves
direct equality with the quotient used in CTV11--12.

Combining the complete exact CEP49 and CTV13 identities gives
\[
 \begin{gathered}
 F_K^{(s)}
  =\mathcal F_{\partial,k}^{(s)}
       -F_{\Xi,k}^{(s)}
       +\mathcal E_{A,k}^{(s)}
       +\sum_N\sigma_N\mathcal R_{N,v}^{(s)},\\
 \left|F_K^{(s)}-\mathcal F_{\partial,k}^{(s)}
                         +F_{\Xi,k}^{(s)}\right|
 \le2\sum_NE_N+\sum_NC_{N,v}^{\rm tail}=o(kq).
 \end{gathered}\tag{CTV22}
\]
Here the first error has its fixed-actual-period scope from CEP49;
the new tail bound has the stronger uniform-in-$v$ scope just
proved. Thus the remaining leading calculation now consists
of the explicit outer-strip quotient of rank $16k-48$ and the
original observation quotient of rank $8k-32$. Both have their
complete original metric maps and four cutoffs specified.
No value of either leading volume is assigned by this reduction,
and no derivative of the earlier $q^2$ asymptotic law was used.

% END CURRENT CTV

\clearpage
% BEGIN CURRENT ROQ
% ROQ: original residual observation quotient, with full coefficient maps.
\section{The actual observation remaining inside the conductor kernel}
\label{sec:ROQ}

Retain the original simple quartet, its actual unit and period matrix
$\mathsf H=F^{-1}\Pi U\mathcal E_{\rm prim}$, and its proved five-member
quadric orbit. Put $k=4l+1\ge9$, $q=(k+1)^2$, $q'=(k-7)^2$,
$\Delta=q-q'=16k-48$, $r=8k-16$, and $j=8k-32$.
The period parameter and logarithm branch are fixed at a point on this
proved locus. All constructions below use the original ordered primary
coordinates at
\[
 \beta_{ab}=k/2+(2a-k)\delta+i(2b-k)\gamma,
 \qquad \omega_{ab}=(2b-k)\gamma-i(2a-k)\delta,
 \quad 0\le a,b\le k.
 \tag{ROQ1}
\]
The full conductor is the actual biform
$A=\sum_{a,b=0}^8a_{ab}e_{ab}^{(8)}$ of OCF4. Its matrix
$C:\mathbb C^q\to\mathbb C^{q'}$ is exactly CEP4:
$C_{(a,b),(i,h)}=a_{i-a,h-b}$ when both differences belong to
$\{0,\ldots,8\}$, and zero otherwise. Thus $C^{\mathsf T}$
is multiplication by the entire $A$, and $V_A=\ker C$.
The original observation is $\Lambda$, and $K=\ker\Lambda\subset V_A$.
Ordinary transpose throughout the coefficient duality is denoted by
$\mathsf T$; a star denotes conjugate transpose in a displayed metric.

\subsection{An explicit residual map and its actual boundary basis}

Use the complete OCF6 elimination, with its actual largest nonzero
coefficient $a_{r_*s_*}$ and every remaining $a_{ab}$. To avoid confusing
its degree-$k$ matrices with the source cutoff, write them as
\[
 \mathsf N:\mathbb C^q\twoheadrightarrow\mathbb C^\Delta,
 \quad \mathsf S:\mathbb C^\Delta\hookrightarrow\mathbb C^q,
 \quad \mathsf V:\mathbb C^q\to\mathbb C^{q'},
 \qquad
 \mathsf S\mathsf N+C^{\mathsf T}\mathsf V=I_q,
 \quad \mathsf N\mathsf S=I_\Delta.
 \tag{ROQ2}
\]
These are the exact identities OCF7. Specifically, at each decreasing
lexicographic conductor position the elimination subtracts its current
coefficient divided by $a_{r_*s_*}$ times the entire translated biform
$A e_{ab}^{(k-8)}$. This specifies all entries of $\mathsf N,\mathsf V$;
$\mathsf S$ inserts the retained strip coordinates. It also proves
$\ker\mathsf N=\operatorname{im}C^{\mathsf T}$.
Taking the annihilator for the bilinear coefficient pairing therefore gives
\[
 I_V=\mathsf N^{\mathsf T}:\mathbb C^\Delta\xrightarrow{\ \sim\ }V_A,
 \qquad I_V^{-1}=\mathsf S^{\mathsf T}|_{V_A}.
 \tag{ROQ3}
\]
For a direct verification, $C\mathsf N^{\mathsf T}=0$ by transposing
$\mathsf NC^{\mathsf T}=0$, and $\mathsf S^{\mathsf T}\mathsf N^{\mathsf T}=I$.
If $Cx=0$, transposing the first identity of ROQ2 gives
$x=\mathsf N^{\mathsf T}\mathsf S^{\mathsf T}x$.

Retain the $j$ independent original invariant certificates $F_{k,h}(x)$
selected by OCF20, and their entire biform coefficient columns $f_h$:
\[
 \mathsf F=[f_1,\ldots,f_j],\qquad
 f_h=\operatorname{coeff}\bigl(F_{k,h}(\mathsf Ht)\bigr),\qquad
 Z=\mathsf N\mathsf F\in\operatorname{Mat}_{\Delta\times j}(\mathbb C),
 \qquad D=\mathsf F^{\mathsf T}.
 \tag{ROQ4}
\]
Here $t=(z_0w_0,z_0w_1,z_1w_0,z_1w_1)^{\mathsf T}$ and the coefficient
order is the full original biform order. Every certificate is supplied
by the fourteen-generator recursion OCF25--26, or equivalently the
complete twenty-five-generator recursion OCF9--20. At a recursion step
its certificate is the actual product $x^\nu F_{k-e,h}$, and its column
is $\mathsf N_k M_{f_\nu}\mathsf S_{k-e}Z_{k-e,h}$. The coefficients
of $f_\nu$ are the full multinomials OCF13 in the actual $\mathsf H$.
Thus ROQ4 specifies finite entries without choosing a substitute period
or introducing a freely chosen observation. The proved image and rank are
$\operatorname{im}Z=\mathsf N J_k$ and $\operatorname{rank}Z=j$.

Let $I$ be the $j$ retained pivot rows of $Z$, let $T$ be their
complement among the $\Delta$ strip rows, and put $Z_I=E_I^{\mathsf T}Z$,
$Z_T=E_T^{\mathsf T}Z$. The original OCF selection guarantees that
$Z_I$ is invertible on its exhaustive exact pivot chart. ROQ2 applied
to $\mathsf F$ gives the full coefficient identity
\[
 \mathsf F=\mathsf SZ+C^{\mathsf T}\mathsf V\mathsf F,
 \qquad
 D=Z^{\mathsf T}\mathsf S^{\mathsf T}
                      +(\mathsf V\mathsf F)^{\mathsf T}C.
 \tag{ROQ5}
\]
In particular the unreduced invariant rows and the reduced strip rows
agree on $V_A$, with precisely the displayed conductor multiple as their
ambient difference. The residual observation has the explicit forms
\[
 \Xi:V_A\longrightarrow\mathbb C^j,\quad x\longmapsto Dx,
 \qquad \Xi I_V=Z^{\mathsf T},\qquad
 L=\mathsf N^{\mathsf T}E_I(Z_I^{\mathsf T})^{-1},
 \quad CL=0,\quad DL=I_j.
 \tag{ROQ6}
\]
The last three equalities follow by multiplication and ROQ3--5.
Moreover $\ker\Xi=K$: the rows $C$ and $D$ together span the
coefficient space $J_k$. Indeed conductor multiplication spans
$\ker\mathsf N$, its image is contained in $J_k$ by OCF17, and
the $f_h$ project to the basis $Z$ of $\mathsf NJ_k$.
Thus a primary vector killed by both $C$ and $D$ is exactly the
annihilator of $J_k$, which is the original $K$ by OKM3--5.
This proves rank $j$, surjectivity, and the complete exact sequence
\[
 0\longrightarrow K\longrightarrow V_A
       \xrightarrow{\ \Xi\ }\mathbb C^j\longrightarrow0.
 \tag{ROQ7}
\]

For every $y\in\mathbb C^j$ define
$\mathfrak b(y)=\Lambda Ly$. These vectors give the literal isomorphism
$\mathfrak b:\mathbb C^j\to\Lambda(V_A)$. They span because for
$x\in V_A$, $x-LDx\in K$ by ROQ6--7. They are independent because
$\Lambda Ly=0$ implies $Ly\in K$, hence $y=DLy=0$.
The inverse on the displayed image is $\Lambda x\mapsto Dx$ for
$x\in V_A$: replacing $x$ by another such lift changes it by $K$,
which $D$ kills. In particular the boundary basis vectors are exactly
$\mathfrak b_h=\sum_{a,b}L_{(a,b),h}\Lambda\varepsilon_{ab}$.
Each $\Lambda\varepsilon_{ab}$ is the full original projected tensor
coefficient in OKA7, including its literal $1/5$ cyclic average.

The whole original observation can be recovered from these rows as well.
Take an actual right inverse $L_C$ of $C$, for example the inverse of
the first nonzero square column minor with its coordinate insertion,
and define
\[
 \mathsf L=\bigl[L_C-LDL_C,\ L\bigr],\qquad
 \mathsf R=\begin{pmatrix}C\\D\end{pmatrix},\qquad
 \mathsf R\mathsf L=I_{q'+j},\qquad
 \Lambda=(\Lambda\mathsf L)\mathsf R.
 \tag{ROQ8}
\]
The first identity follows by multiplying its four blocks. For the
last, $\mathsf R(x-\mathsf L\mathsf Rx)=0$, whose kernel is $K$
by ROQ7. Thus the $q'+j=q-r$ columns of $\Lambda\mathsf L$ form an
actual basis of the original boundary: applying its kernel and then
$\mathsf R\mathsf L=I$ proves independence, and ROQ8 proves spanning.
The last $j$ columns are exactly $\mathfrak b_h$.

For completeness the inherited period metric on this residual basis is
also explicit. Let $B$ be the complete matrix AKS8 in primary
coordinates, so $\jmath\Lambda=\mathcal F_0B$. If $\nu,\mu$ range over
the original weight-zero degree-$k$ monomials, set
\[
 (\mathcal G_0)_{\nu\mu}
 =\sum_{\substack{n_{tu}\ge0\\\sum_u n_{tu}=\nu_t\\
                                      \sum_t n_{tu}=\mu_u}}
       \frac{k!}{\prod_{t,u}n_{tu}!}
               \prod_{t,u}[5(\delta_{tu}+1)]^{n_{tu}},
 \qquad P_{\rm per}=L^*B^*\mathcal G_0BL>0.
 \tag{ROQ9}
\]
Expansion of the original word inner product gives
$\mathcal G_0=\mathcal F_0^*\mathcal F_0$ as in AKS10, so the displayed
matrix is exactly the Gram of the independent vectors
$\jmath\mathfrak b_h$. Its strict positivity follows from that
independence and injectivity of $\jmath$. All off-diagonal word terms
and every actual entry of $\mathsf H$ remain in ROQ9.

\subsection{All four induced Gamma metrics and their fixed determinant factor}

For either original source $s\in\{1,k\}$ put
$b_s=s/2$, $M_s=(2\pi)^{s/2}$ and retain
\[
 \begin{gathered}
 p_0=1,\quad p_1=y,\quad
 p_{n+1}=yp_n-n(b_s+n-1)p_{n-1},\\
 h_n=M_s n!(b_s)_n,
 \qquad (u_n)_{ab}=p_n(\omega_{ab})/\sqrt{h_n}.
 \end{gathered}
 \tag{ROQ10}
\]
Thus $U_N=[u_0,\ldots,u_N]$ is precisely CEP21 and
$R_N=U_NU_N^*>0$, $G_N=R_N^{-1}$ for every $N\ge q-1$.
The positivity follows from its first $q$ columns, whose determinant
is the original Vandermonde times their nonzero leading coefficients.
Define
\[
 H_{V,N}=\overline{\mathsf N}G_N\mathsf N^{\mathsf T},\qquad
 Q_{\Xi,N}=
       (Z^{\mathsf T}H_{V,N}^{-1}\overline Z)^{-1}>0.
 \tag{ROQ11}
\]
These are matrices of sizes $\Delta$ and $j$, respectively. The
minimum of $z^*H_{V,N}z$ under $Z^{\mathsf T}z=y$ is attained at
$z=H_{V,N}^{-1}\overline ZQ_{\Xi,N}y$. Its difference from any other
lift lies in $\ker Z^{\mathsf T}$ and is perpendicular in the
$H_{V,N}$ metric. Its squared norm is $y^*Q_{\Xi,N}y$.
Thus ROQ11 calculates exactly the induced metric of ROQ7 in the
actual boundary coordinates ROQ6. This proof also identifies it with
the quotient metric appearing in CEP27 and CEP49.

The actual kernel frame in the strip is
\[
 \rho=E_T^{\mathsf T}-Z_T Z_I^{-1}E_I^{\mathsf T},\qquad
 I_K=\mathsf N^{\mathsf T}\rho^{\mathsf T},\qquad
 F_{\rm str}=[\rho^{\mathsf T},\ E_I(Z_I^{\mathsf T})^{-1}].
 \tag{ROQ12}
\]
The first two matrices are precisely the OCF/OKM frame. In the ordered
rows $(T,I)$ the last one has blocks
$\left(\begin{smallmatrix}I_r&0\\
 -(Z_I^{\mathsf T})^{-1}Z_T^{\mathsf T}&(Z_I^{\mathsf T})^{-1}
 \end{smallmatrix}\right)$.
Consequently $|\det F_{\rm str}|^2=|\det Z_I|^{-2}$, including the
absolute value of the row-order permutation. Block elimination in
$F_{\rm str}^*H_{V,N}F_{\rm str}$ gives the exact identity
\[
 \det(I_K^*G_NI_K)\det Q_{\Xi,N}
                 =\frac{\det H_{V,N}}{|\det Z_I|^2}.
 \tag{ROQ13}
\]
Indeed its upper block is the kernel Gram and its Schur complement
is the minimum quotient metric of the previous paragraph, since
$Z^{\mathsf T}F_{\rm str}=[0,I_j]$. The determinant of a block
matrix is the determinant of its upper block times that Schur
complement, while congruence contributes $|\det F_{\rm str}|^2$.
This proves ROQ13 without assigning modulus one to the actual pivot.

Define $F_{\Xi,k}^{(s)}$ by the four ordered endpoints
$\mathcal N=(q-1,q,2q-1,2q)$ and signs $\sigma=(1,1,-1,-1)$:
\[
 F_{\Xi,k}^{(s)}=\sum_{N\in\mathcal N}\sigma_N\log\det Q_{\Xi,N}.
 \tag{ROQ14}
\]
The fixed factor $|\det Z_I|^{-2}$ in ROQ13 cancels with these signs.
If a different exact certificate basis is selected, its residual rows
on $V_A$ are $D'=MD$ for an invertible $j\times j$ matrix $M$.
This follows because both lists are bases of the dual of $V_A/K$.
For an explicit inverse formula, $M=D'L$, and the old and new right
sections give its inverse; composing the two surjections proves both
inverse identities. The new metric is
$Q'_{\Xi,N}=M^{-*}Q_{\Xi,N}M^{-1}$, directly from their identical
minimum problem. Its determinant changes by $|\det M|^{-2}$, which
also cancels in ROQ14. No pivot or source mass has been removed at
an individual endpoint. In fact ROQ10 gives $Q_{\Xi,N}=M_s\widehat Q_{\Xi,N}$
when the common mass is explicitly factored out of every column;
its determinant carries $M_s^j$. The metric relative to the retained
period Gram ROQ9 has determinant $\det Q_{\Xi,N}/\det P_{\rm per}$,
whose fixed denominator cancels only in the same four-endpoint sum.

\subsection{The exact constrained covariance in the original source}

The residual metric can be computed without inverting the strip frame.
Define, using the complete original covariance,
\[
 \begin{split}
 A_N^C&=CR_NC^*,\qquad B_N^D=DR_NC^*,\qquad T_N^D=DR_ND^*,\\
 W_N&=R_N-R_NC^*(A_N^C)^{-1}CR_N,\\
 \Omega_N&=DW_ND^*
          =T_N^D-B_N^D(A_N^C)^{-1}(B_N^D)^*.
 \end{split}\tag{ROQ15}
\]
Since $C$ has full row rank and $R_N>0$, $A_N^C>0$. There is the exact
identity
\[
 W_N=\mathsf N^{\mathsf T}H_{V,N}^{-1}\overline{\mathsf N},
 \qquad \Omega_N=Z^{\mathsf T}H_{V,N}^{-1}\overline Z>0,
 \qquad Q_{\Xi,N}=\Omega_N^{-1}.
 \tag{ROQ16}
\]
To prove its first equality, $CW_N=0$, $W_N$ is Hermitian, and for
$x\in\ker C$ direct multiplication gives $W_NG_Nx=x$.
The other displayed candidate is Hermitian, has image in $\ker C$
by ROQ3, and has the same property on $x=\mathsf N^{\mathsf T}z$.
Such a matrix is unique: for their Hermitian difference $H$, its image
lies in $V_A$ and $HG_Nx=0$ for $x\in V_A$. For arbitrary $y$ and
$x\in V_A$, this yields $x^*G_NHy=(HG_Nx)^*y=0$.
Putting $x=Hy\in V_A$ gives $(Hy)^*G_N(Hy)=0$, hence $Hy=0$.
This proves equality. Now use $D\mathsf N^{\mathsf T}=Z^{\mathsf T}$
from ROQ6 and take the adjoint to prove the other two identities.

Every entry in ROQ15 is the finite sum of original column products:
$A_N^C=\sum_{n=0}^N(Cu_n)(Cu_n)^*$,
$B_N^D=\sum_{n=0}^N(Du_n)(Cu_n)^*$, and
$T_N^D=\sum_{n=0}^N(Du_n)(Du_n)^*$.
The conductor column is exactly
$Cu_n=V_{N-v}A_Ne_n$ by CEP22 when $n\le N$; its vanishing low
columns and centre difference are precisely CEP13--16.
The residual column is explicitly
$(Du_n)_h=\sum_{a,b}(f_h)_{ab}p_n(\omega_{ab})/\sqrt{h_n}$.
In particular ROQ15 retains all conductor--observation cross terms.
Replacing $D$ by $D+TC$ leaves $\Omega_N$ exactly unchanged because
$CW_N=0=W_NC^*$. Applied to ROQ5, this proves that the full and
strip row presentations give exactly the same metric, rather than
assuming their ambient coefficient matrices coincide.

\subsection{Exact increments and a positive residual four-endpoint return}

For $N\ge q-1$ let $u=u_{N+1}$ be the next original column ROQ10 and
define
\[
 \begin{gathered}
 a_N=1+(Cu)^*(A_N^C)^{-1}(Cu)>0,\\
 w_N=u-R_NC^*(A_N^C)^{-1}Cu\in V_A,\qquad
 z_N=Dw_N=Du-B_N^D(A_N^C)^{-1}Cu,\\
 \eta_N=\frac{z_N^*Q_{\Xi,N}z_N}{a_N}\ge0.
 \end{gathered}\tag{ROQ17}
\]
The complete updates are
\[
 \begin{gathered}
 W_{N+1}=W_N+w_Nw_N^*/a_N,\qquad
 \Omega_{N+1}=\Omega_N+z_Nz_N^*/a_N,\\
 Q_{\Xi,N+1}=Q_{\Xi,N}
       -\frac{Q_{\Xi,N}z_Nz_N^*Q_{\Xi,N}}
                         {a_N+z_N^*Q_{\Xi,N}z_N},\\
 \log\det Q_{\Xi,N}-\log\det Q_{\Xi,N+1}=\log(1+\eta_N).
 \end{gathered}\tag{ROQ18}
\]
Here is a direct proof of the covariance update. Put $a=Cu$ and
$v=R_NC^*(A_N^C)^{-1}a$. Since $R_{N+1}=R_N+uu^*$,
$A_{N+1}^C=A_N^C+aa^*$, whose inverse is
$ (A_N^C)^{-1}-(A_N^C)^{-1}aa^*(A_N^C)^{-1}/a_N$;
multiplying by $A_N^C+aa^*$ verifies this inverse exactly.
For the full covariance calculation write
$\alpha=a^*(A_N^C)^{-1}a$.
The subtracted covariance product expands as
\[
 R_NC^*(A_N^C)^{-1}CR_N
       +vu^*+uv^*+\alpha uu^*
       -\frac{(v+\alpha u)(v+\alpha u)^*}{1+\alpha}.
\]
Subtracting this from $R_N+uu^*$ and the old $W_N$ leaves
$(u-v)(u-v)^*/(1+\alpha)$, proving the first update.
Multiplying by $D,D^*$ proves the second. The same inverse identity
proves the third, and the determinant identity
$\det(H+zz^*/a)=\det H(1+z^*H^{-1}z/a)$ proves the fourth.
For that identity conjugate by $H^{-1/2}$; the rank-one matrix
$H^{-1/2}zz^*H^{-1/2}/a$ has one eigenvalue $z^*H^{-1}z/a$
and all its others zero. This completes the update proof.

Telescoping with the four actual endpoints, rather than changing the
source degree, gives the exact finite sum
\[
 F_{\Xi,k}^{(s)}=\log(1+\eta_{q-1})
       +2\sum_{N=q}^{2q-2}\log(1+\eta_N)
       +\log(1+\eta_{2q-1})\ge0.
 \tag{ROQ19}
\]
Indeed $\log\det Q_{q-1}-\log\det Q_{2q-1}$ sums the differences
for $N=q-1,\ldots,2q-2$, while
$\log\det Q_q-\log\det Q_{2q}$ sums those for $N=q,\ldots,2q-1$.
This proves each endpoint weight. Its first increment
$\lambda_s=\log(1+\eta_{q-1})$ is a completely specified
period-dependent nonnegative lower bound for $F_{\Xi,k}^{(s)}$.
The formulas also characterize equality: $\lambda_s=0$ exactly when
$Du_q=B_{q-1}^D(A_{q-1}^C)^{-1}Cu_q$; the full return vanishes
exactly when this corresponding row residual is zero for every
column $u_q,\ldots,u_{2q}$. Positivity of $Q_{\Xi,N}$ proves both
directions, without asserting a nonzero generic residual.

\subsection{Finite scalar enclosures and propagation into CEP49}

Retain the exact AKS3 scalars $\mathcal Z_{s,t}$ and
$\mathcal L_{s,t}=\log\mathcal Z_{s,t}$, whose full original polynomial
covariance bound is proved in AKS20--27. In primary coordinates it says
\[
 \mathcal Z_{s,t}^{-1}G_{q-1}\preceq G_{q+t-1}\preceq G_{q-1},
                  \qquad 0\le t\le q+1.
 \tag{ROQ20}
\]
The unchanged coefficient map from AKS26 to primary coordinates is
the ordered $V_\beta$, so congruence by that invertible map proves
exactly the displayed statement. Restrict it to $V_A$ and take the
minimum on the fixed fibre $Dx=y$ in $V_A$. Inequalities of squared
norms hold for every such $x$, hence also for their attained minima.
Thus
$\mathcal Z_{s,t}^{-1}Q_{\Xi,q-1}\preceq Q_{\Xi,q+t-1}
\preceq Q_{\Xi,q-1}$.
If $d_t=\log\det Q_{\Xi,q-1}-\log\det Q_{\Xi,q+t-1}$,
the eigenvalues after congruence by $Q_{\Xi,q-1}^{-1/2}$ give
$0\le d_t\le j\mathcal L_{s,t}$, and ROQ18 gives their monotonicity.
Since $d_1=\lambda_s$, the exact four-endpoint identity
$F_{\Xi,k}^{(s)}=d_q+d_{q+1}-d_1$ proves
\[
 \lambda_s\le F_{\Xi,k}^{(s)}
       \le j(\mathcal L_{s,q}+\mathcal L_{s,q+1})-\lambda_s.
 \tag{ROQ21}
\]
The actual rank $j$ enters only through this proved eigenvalue bound;
the determinant and increment formulas ROQ11--19 retain its full
orientation. No leading volume is assigned from a dimension.

Let $\epsilon_{A,k}^{(s)}=2\sum_{N\in\mathcal N}E_N$ with the
complete original constants of CEP29--36. CEP49 gives
$F_K^{(s)}=\mathcal T_k^{(s)}-F_{\Xi,k}^{(s)}+
\mathcal E_{A,k}^{(s)}$ and
$|\mathcal E_{A,k}^{(s)}|\le\epsilon_{A,k}^{(s)}=o(kq)$
for a fixed actual period. Substitution of ROQ21 proves
\[
 \begin{split}
 \mathcal T_k^{(s)}-\epsilon_{A,k}^{(s)}
       -j(\mathcal L_{s,q}+\mathcal L_{s,q+1})+\lambda_s
       &\le F_K^{(s)},\\
 F_K^{(s)}&\le
          \mathcal T_k^{(s)}+\epsilon_{A,k}^{(s)}-\lambda_s.
 \end{split}\tag{ROQ22}
\]
Intersecting this interval with the full AKJ2 or AKJ4 intervals preserves
each of their finite constants and source correlations. More precisely,
one may use the exact value $\mathcal E_{A,k}^{(s)}$ in place of its
upper or lower envelope in the appropriate side. Formula ROQ19 supplies
every additional original positive subtraction, rather than only the
first one. These are finite results on the residual quotient itself;
they do not claim that its leading $kq$ coefficient has been evaluated.

\subsection{The complete residual quotient in the conductor graph}

At a fixed original endpoint $N$, let $\mathcal P_N$, $I_N$,
$\mathcal W_N=U_N^{-1}(V_A)$ and the actual source isomorphism
$S_Np=(P_{<v}p,\mathcal T_{A,N}p)$ be precisely CEP28--31.
The source carries its original orthonormal Gamma coordinates of
order $s=1$ or $s=k$. The induced linear isomorphism on the residual
observation quotient is
\[
 \mathcal W_N/U_N^{-1}(K)
       \xrightarrow{\ [p]\mapsto[S_Np]\ }
 S_N\mathcal W_N/S_NU_N^{-1}(K),
 \qquad [p]\longmapsto DU_Np\in\mathbb C^j.
 \tag{ROQ23}
\]
The last map identifies the left quotient with $V_A/K$ and its fixed
coordinates ROQ6. Its kernel and surjectivity follow from ROQ7 and
surjectivity of $U_N$; applying the invertible $S_N$ proves the other
isomorphism. Thus all relation directions, including the source ideal
and the entire actual $K$, are retained in the transported denominator.

Here are exact matrices of that denominator and the resulting metric.
Let $Z_N^{\chi}$ be the original source coefficient columns
$\chi_k S^a$, $0\le a\le N-q$, with no columns when $N=q-1$.
Take the actual primary frames $I_K$ of ROQ12 and $L$ of ROQ6 and set
\[
 Y_{K,N}=U_N^*G_NI_K,\qquad Y_{L,N}=U_N^*G_NL,
 \qquad J_N=[Z_N^{\chi},Y_{K,N}].
 \tag{ROQ24}
\]
The columns of $Y_{K,N}$ are the minimum lifts of $I_K$ by the same
orthogonality proof as OKM8. In particular they are perpendicular to
$I_N=\ker U_N$, so $J_N$ has independent columns and its image is
exactly $U_N^{-1}(K)$. It has $N+1-q+r$ columns, while
$\dim\mathcal W_N=N+1-q+\Delta$, leaving precisely $j$ quotient
dimensions. The transported complete relation projector and Gram are
\[
 \begin{gathered}
 P_N^{\Xi,\rm graph}
   =I-S_NJ_N(J_N^*S_N^*S_NJ_N)^{-1}J_N^*S_N^*,\\
 \widetilde Q_{\Xi,N}
       =Y_{L,N}^*S_N^*P_N^{\Xi,\rm graph}S_NY_{L,N}>0.
 \end{gathered}\tag{ROQ25}
\]
The inverse exists because $S_N$ is invertible and $J_N$ has independent
columns. Multiplication verifies that the projector is Hermitian,
idempotent and projects onto the orthogonal complement of the entire
$S_NU_N^{-1}(K)$. Thus its Gram is the minimum norm of the transported
quotient ROQ23. Replacing $Y_{L,N}$ by another source lift of $L$ adds
columns in $I_N$, which this projector kills. Replacing the section
$L$ by another section of ROQ6 adds columns in $K$; their lifts lie
in the span of $J_N$ and are killed. Therefore ROQ25 computes
the specified quotient metric independently of either choice.

Before applying $S_N$, the same projection gives
\[
 Q_{\Xi,N}
 =L^*G_NL-L^*G_NI_K(I_K^*G_NI_K)^{-1}I_K^*G_NL.
 \tag{ROQ26}
\]
To prove it, $Y_{L,N}$ and $Y_{K,N}$ are perpendicular to the ideal,
so projection off $J_N$ only subtracts
$Y_{K,N}(Y_{K,N}^*Y_{K,N})^{-1}Y_{K,N}^*Y_{L,N}$.
Using $U_NU_N^*=R_N$ in their Gram products yields exactly ROQ26.
It is also the minimum of the original $G_N$ norm on $Ly+K$, so
it agrees with the already proved ROQ11 and ROQ16.

The all-rank exterior estimate CEP33--35 applies to ROQ23: represent
a quotient class by its minimum lift in the denominator's orthogonal
complement, apply $S_N$, then project off the target denominator.
This projection has every exterior norm at most one. Consequently the
induced $j$th exterior norm is at most $\|\bigwedge^jS_N\|$, and its
inverse is at most $\|\bigwedge^jS_N^{-1}\|$ by the identical argument
for $S_N^{-1}$. CEP34 bounds both logarithms from above by $E_N$.
Applying them to the fixed frame ROQ6 and squaring exterior volumes gives
\[
 \left|\log\det\widetilde Q_{\Xi,N}
                -\log\det Q_{\Xi,N}\right|\le2E_N,\qquad
 \left|\sum_{N\in\mathcal N}\sigma_N\log\det\widetilde Q_{\Xi,N}
                       -F_{\Xi,k}^{(s)}\right|
                       \le\epsilon_{A,k}^{(s)}=o(kq).
 \tag{ROQ27}
\]
All cross terms in $J_N^*S_N^*S_NJ_N$ remain present in ROQ25,
including $P_{<v}(\chi_k S^a)$ and its pairings with the minimum
kernel lifts. Its ideal columns can equally be generated from the full
products $d_{rs}$ and maps $\mathcal U_{A,L}$ in CEP41--45, because
those formulas prove exactly $S_NZ_N^{\chi}$ column by column.
This supplies the typed relation between the explicit residual
observation and the conductor graph at all four original cutoffs.

% END CURRENT ROQ

\clearpage
% BEGIN CURRENT OSP
% Full original-root comparison; no sealed provider is modified.
\section{A quantitative power comparison for the actual outer-strip quotient}
\label{sec:OSP}

Retain the original real parameters and integers
\[
\begin{gathered}
0<\delta<\tfrac12,\quad \gamma>2,\quad k=4l+1\ge9,\\
q=(k+1)^2,\quad k'=k-8,\quad q'=(k-7)^2,\quad
\Delta=q-q'=16k-48.
\end{gathered}\tag{OSP1}
\]
The two source orders in every formula below are precisely $s=1,k$.
Set $b=s/2$, $M_s=(2\pi)^{s/2}$, and retain the full original measure
\[
dm_s(y)=M_s\frac{2^b}{4\pi\Gamma(b)}
       |\Gamma(b/2+iy/2)|^2\,dy.
\tag{OSP2}
\]
In the original coordinate $S=k/2+iy$, and respectively
$S'=k/2-4+iy$, the complete monic root polynomials are
\[
\begin{split}
P_k(y)&=i^{-q}\chi_k(k/2+iy)
 =\prod_{a,b_0=0}^{k}
 [y-(2b_0-k)\gamma+i(2a-k)\delta],\\
P_{k'}(y)&=i^{-q'}\chi_{k'}(k/2-4+iy)
 =\prod_{a,b_0=0}^{k'}
 [y-(2b_0-k')\gamma+i(2a-k')\delta].
\end{split}\tag{OSP3}
\]
The index $b_0$ in these products is distinct from the source parameter
$b=s/2$. Each root occurs with exactly its displayed multiplicity.
The map $(a,b_0)\mapsto(a+4,b_0+4)$ identifies the second root list
with the central subrectangle of the first, so
\[
P_k=P_{k'}P_{\partial,k},\qquad
P_{\partial,k}(y)=
\prod_{(a,b_0)\in\{0,\ldots,k\}^2\setminus\{4,\ldots,k-4\}^2}
[y-(2b_0-k)\gamma+i(2a-k)\delta].
\tag{OSP4}
\]
This proves the factorization and its monic constant; the outer degree
is $\Delta$. Both complete root lists are invariant under
$\omega\mapsto-\omega$. Every root has modulus at most
\[
R=k\sqrt{\delta^2+\gamma^2}.
\tag{OSP5}
\]

\subsection{The two original source masses and a proved interval lower bound}

We first record the source facts needed in the comparison, including
their constants and proofs. For $a>0$ and real $x$, the two Euler
integrals and the change of variables in the beta integral give
\[
\frac{\Gamma(a+ix)\Gamma(a-ix)}{\Gamma(2a)}
 =\int_{\mathbb R}e^{ixu}(2\cosh(u/2))^{-2a}\,du.
\tag{OSP6}
\]
For example the intermediate beta integral is
$\int_0^\infty t^{a+ix-1}(1+t)^{-2a}dt$, and $t=e^u$
gives the displayed expression. Absolute convergence follows from
$a>0$. Rotating Euler's integral through
$\operatorname{sgn}(y)\theta$, $0<\theta<\pi/2$, gives
\[
|\Gamma(a+iy/2)|
 \le e^{-\theta|y|/2}\Gamma(a)(\cos\theta)^{-a}.
\tag{OSP7}
\]
The rotation follows from Cauchy's theorem on the intervening
sector: the large circular arc tends to zero because its real
exponential has positive cosine, and the small one tends to zero
because $a>0$. The remaining radial integral has absolute value
at most $\Gamma(a)(\cos\theta)^{-a}$. In particular the Fourier
transform in OSP6 is integrable. Fourier inversion, followed by
$y=2x$, now gives
\[
\int_{\mathbb R}e^{i\xi y}\,dm_s(y)
  =M_s(\cosh\xi)^{-b},\qquad
\int_{\mathbb R}dm_s=M_s.
\tag{OSP8}
\]
The same identity holds analytically for $|\operatorname{Im}\xi|<\pi/2$:
OSP7 with any larger $\theta<\pi/2$ supplies an integrable dominating
function on each smaller strip. Substituting $\xi=-it$ proves
\[
\int_{\mathbb R}e^{ty}\,dm_s(y)=M_s(\cos t)^{-b},
\qquad |t|<\pi/2.
\tag{OSP9}
\]
This also proves finiteness of all moments used below.

Write
\[
d\sigma(y)=\frac{|\Gamma(1/4+iy/2)|^2}{2\pi}\,dy,\quad
C_L=\frac{\pi}{\Gamma(3/4)^2},\quad
\ell_s=\frac{s-1}{4},\quad
\beta_s=\frac{M_s}{\sqrt2\,\Gamma(s/2)}.
\tag{OSP10}
\]
The two values of $\ell_s$ are the integers $0,l$. The reflection
identity gives a useful explicit lower bound
\[
\frac{d\sigma}{dy}(y)\ge C_Le^{-\pi|y|}.
\tag{OSP11}
\]
Here is its proof on the required domain. For $0<\operatorname{Re}z<1$
the Euler integrals give
$\Gamma(z)\Gamma(1-z)=\int_0^\infty t^{z-1}/(1+t)\,dt$.
On a keyhole contour with argument in $(0,2\pi)$, the circular
terms vanish with bounds $O(r^{\operatorname{Re}z})$ at zero
and $O(R^{\operatorname{Re}z-1})$ at infinity. The residue at $-1$
is $e^{i\pi(z-1)}$; the two radial integrals therefore give
$(1-e^{2\pi iz})\int_0^\infty t^{z-1}/(1+t)\,dt
=2\pi i e^{i\pi(z-1)}$. Their ratio is $\pi/\sin(\pi z)$.
At $z=1/4+iy/2$ this proves
\[
|\Gamma(1/4+iy/2)|^2|\Gamma(3/4-iy/2)|^2
 =\frac{2\pi^2}{\cosh(\pi y)}.
\]
The unrotated Euler integral bounds the second squared factor by
$\Gamma(3/4)^2$, proving
$d\sigma/dy\ge C_L/\cosh(\pi y)\ge C_Le^{-\pi|y|}$.

Gamma recurrence, proved by one integration by parts in Euler's
integral, gives the exact original-source identity
\[
dm_s(y)=\beta_s
 \prod_{j=0}^{\ell_s-1}\bigl(y^2+(2j+\tfrac12)^2\bigr)\,d\sigma(y).
\tag{OSP12}
\]
Indeed $b/2=\ell_s+1/4$, and each recurrence factor before
squaring and clearing denominators is
$j+1/4+iy/2$; its squared modulus is
$[y^2+(2j+1/2)^2]/4$. The factor $4^{-\ell_s}$ together with
the prefactor in OSP2 is exactly $\beta_s/(2\pi)$, because
$b=2\ell_s+1/2$. This proves every factor in OSP12.
For $s=1$ its product is empty and $\beta_1=1$.
In particular, on the full interval $q\le y\le2q$,
\[
\frac{dm_s}{dy}(y)
 \ge \beta_s q^{2\ell_s}C_Le^{-2\pi q}.
\tag{OSP13}
\]
The source mass-to-lower-density factor obeys the uniform finite bound
\[
\frac{M_s}{\beta_s q^{2\ell_s}}
 =\sqrt2\,\Gamma(b)q^{-2\ell_s}
 \le\sqrt{2\pi}\qquad(s=1,k).
\tag{OSP14}
\]
To prove the inequality, use
$\Gamma(b)=\sqrt\pi\prod_{j=0}^{2\ell_s-1}(j+1/2)
\le\sqrt\pi\,b^{2\ell_s}$ and $b\le q$.
All source masses have thus been retained before taking this exact
ratio.

\subsection{A complete coefficient-space estimate for the inner interval}

Let $p\in\mathbb C[y]$ have degree at most $d$, and let
$0<T\le q/2$. Then
\[
\sup_{|y|\le T}|p(y)|^2
 \le \frac{(d+1)^2\,10^{2d}}q
                  \int_q^{2q}|p(u)|^2\,du.
\tag{OSP15}
\]
This controls every coefficient direction; no assumption on zeros,
leading coefficient, or extremizers of $p$ is made.

For proof define the Legendre polynomial
$L_j(x)=(2^j j!)^{-1}(d/dx)^j(x^2-1)^j$.
Its leading coefficient is $(2j)!/[2^j(j!)^2]$. Integrating by
parts $j$ times, with all lower boundary derivatives zero, proves
orthogonality to lower powers and
\[
\int_{-1}^1L_j(x)^2\,dx
=\frac{(2j)!}{2^{2j}(j!)^2}\int_{-1}^1(1-x^2)^j\,dx
=\frac2{2j+1}.
\]
For the middle integral set $x=2t-1$; repeated integration by parts
gives $\int_0^1t^j(1-t)^jdt=(j!)^2/(2j+1)!$.
Consequently $\sqrt{2j+1}L_j(2z-3)$, $0\le j\le d$,
form an orthonormal basis for degree-at-most-$d$ polynomials on
$[1,2]$. Leibniz's rule also gives the explicit identity
\[
L_j(x)=2^{-j}\sum_{h=0}^j\binom jh^2
                         (x-1)^{j-h}(x+1)^h.
\]
For $|x|\le4$ its modulus is at most
$(5/2)^j\sum_h\binom jh^2\le(5/2)^j4^j=10^j$;
the sum inequality follows by comparing with
$(\sum_h\binom jh)^2=4^j$.
Expanding $p(qz)$ in the displayed orthonormal basis and applying
Cauchy--Schwarz gives, for $|z|\le T/q\le1/2$,
\[
|p(qz)|^2
\le\sum_{j=0}^d(2j+1)10^{2j}
                  \int_1^2|p(qv)|^2\,dv
\le(d+1)^2\,10^{2d}\int_1^2|p(qv)|^2\,dv .
\]
Here $|2z-3|\le4$ and $\sum_{j=0}^d(2j+1)=(d+1)^2$.
Changing variables proves OSP15, including $d=0$.

\subsection{The exact original-root map and a finite full-norm comparison}

In this subsection $P$ is either of the complete polynomials
$P_k,P_{k'}$ in OSP3, and $Q=\deg P$ is respectively $q,q'$.
For $d\ge0$ retain the two injections into the same original
polynomial Hilbert space
\[
\begin{aligned}
\mathcal J_P:\mathbb C[y]_{\le d}&\longrightarrow
          \mathbb C[y]_{\le Q+d},&p&\longmapsto Pp,\\
\mathcal J_{0,Q}:\mathbb C[y]_{\le d}&\longrightarrow
          \mathbb C[y]_{\le Q+d},&p&\longmapsto y^Qp .
\end{aligned}\tag{OSP16}
\]
Their codomain norm is precisely $m_s$ in OSP2.
If $P(y)=\sum_{a=0}^QP_a y^a$, set $P_a=0$ outside this range.
In increasing coefficient frames their matrices are
\[
(\mathsf J_P)_{rj}=P_{r-j},\qquad
(\mathsf J_{0,Q})_{rj}=\mathbf1_{\{r=Q+j\}},
\quad 0\le r\le Q+d,\quad 0\le j\le d.
\tag{OSP17}
\]
Thus the exact isomorphism between their images is
$y^Qp\mapsto Pp$, with inverse $Pp\mapsto y^Qp$.
It preserves the original coefficient vector $p$. The first
matrix contains every actual root coefficient; none is replaced
inside that map. Returning to $S$ or $S'$ multiplies the relevant
relation columns by the actual phase $i^Q$ and applies the
triangular coordinate pullback with diagonal entries $i^j$.
Each has the stated exact inverse and determinant of modulus one.

Suppose
\[
R<T\le q/2.
\tag{OSP18}
\]
Define the finite positive constants
\[
\begin{split}
\eta_{Q,T}&=\frac{Q R^2}{T^2-R^2},\\
\kappa_{Q,d,T,s}
 &=\frac{M_s}{\beta_s q^{2\ell_s}C_L}
    \frac{(d+1)^2}{q}\,
    e^{2\pi q}\,10^{2d}
             \left(\frac{T^2+R^2}{q^2}\right)^Q,\\
e_{Q,d,T,s}&=\eta_{Q,T}+\log(1+\kappa_{Q,d,T,s}).
\end{split}\tag{OSP19}
\]
Then every $p\in\mathbb C[y]_{\le d}$ satisfies
\[
e^{-e_{Q,d,T,s}}\int_{\mathbb R}|y|^{2Q}|p(y)|^2\,dm_s(y)
\le\int_{\mathbb R}|P(y)p(y)|^2\,dm_s(y)
\le e^{e_{Q,d,T,s}}
      \int_{\mathbb R}|y|^{2Q}|p(y)|^2\,dm_s(y).
\tag{OSP20}
\]

Here is a proof retaining the entire inner region. Pair the actual
roots as $\omega,-\omega$. On $|y|\ge T$, each paired quotient is
$1-\omega^2/y^2$, so
\[
\left|\log\frac{|P(y)|^2}{|y|^{2Q}}\right|
\le -Q\log(1-R^2/T^2)
\le\frac{Q R^2}{T^2-R^2}=\eta_{Q,T}.
\tag{OSP21}
\]
The first inequality uses $1-u\le|1-z|\le1+u$ for $|z|\le u<1$
on each of $Q/2$ pairs. The second follows by integrating
$1/(1-u)\le1/(1-u_0)$ on $0\le u\le u_0=R^2/T^2$.
For every real $y$, the same pairing gives
$|P(y)|^2\le(y^2+R^2)^Q$. Write
\[
A_B=\int_q^{2q}|y|^{2Q}|p(y)|^2\,dm_s(y).
\]
For nonzero $p$ this number is positive, and OSP13 gives
$A_B\ge q^{2Q}\beta_s q^{2\ell_s}C_Le^{-2\pi q}
 \int_q^{2q}|p(y)|^2dy$.
OSP15 and the exact full mass OSP8 imply
\[
\begin{split}
\int_{|y|<T}|P(y)p(y)|^2\,dm_s(y)
 &\le M_s(T^2+R^2)^Q\sup_{|y|\le T}|p(y)|^2
 \le\kappa_{Q,d,T,s}A_B,\\
\int_{|y|<T}|y|^{2Q}|p(y)|^2\,dm_s(y)
 &\le M_sT^{2Q}\sup_{|y|\le T}|p(y)|^2
 \le\kappa_{Q,d,T,s}A_B .
\end{split}\tag{OSP22}
\]
Let $A_{\rm out}$ and $B_{\rm out}$ be the integrals outside
$(-T,T)$ with factors $|y|^{2Q}$ and $|P|^2$, respectively.
Since $[q,2q]$ is outside that interval, $A_B\le A_{\rm out}$.
OSP21 gives
$e^{-\eta}A_{\rm out}\le B_{\rm out}\le e^\eta A_{\rm out}$.
OSP22 bounds both inner contributions by $\kappa A_{\rm out}$.
Therefore the full reference norm is at most
$(1+\kappa)A_{\rm out}$ and at least $A_{\rm out}$, while the
full actual norm is at least $e^{-\eta}A_{\rm out}$ and at most
$(e^\eta+\kappa)A_{\rm out}\le e^\eta(1+\kappa)A_{\rm out}$.
These four inequalities prove OSP20. The zero polynomial satisfies
it directly.

\subsection{Uniform bounded logarithmic cost at every required dimension}

Choose the fixed number and cutoff
\[
\epsilon=2^{-32},\qquad T=\epsilon q,
\]
and impose the explicit finite conditions
\[
k\ge29,\qquad
k\sqrt{\delta^2+\gamma^2}\le\frac{\epsilon q}{2}.
\tag{OSP23}
\]
In particular $T\ge2R$, as required for a uniform margin in OSP21.
These conditions hold eventually for every fixed original
$\delta,\gamma$; the sufficient bound
$k\ge\max\{29,2\sqrt{\delta^2+\gamma^2}/\epsilon\}$
follows from $q\ge k^2$. The first condition implies $q'\ge q/2$:
$2(k-7)^2-(k+1)^2=k^2-30k+97$ is positive at $k=29$
and is increasing thereafter.

For $Q\in\{q',q\}$ and every $0\le d\le2q$, set
\[
\begin{split}
K_q&=\frac{\sqrt{2\pi}}{C_L}\frac{(2q+1)^2}{q}e^{-2q},\\
\mathcal E_k&=\frac{qR^2}{\epsilon^2q^2-R^2}+\log(1+K_q).
\end{split}\tag{OSP24}
\]
Then OSP19 satisfies $e_{Q,d,T,s}\le\mathcal E_k$ for both
original source orders. To verify this uniformly, OSP14,
$T^2+R^2\le(5/4)\epsilon^2q^2$, $Q\ge q/2$ and $d\le2q$
give
\[
\kappa_{Q,d,T,s}
\le\frac{\sqrt{2\pi}}{C_L}\frac{(2q+1)^2}{q}
\exp\left(q\{2\pi+4\log10+\tfrac12\log(5/4)+\log\epsilon\}\right).
\tag{OSP25}
\]
The expression in braces is less than $-2$. Indeed
$\pi<4$, $\log10<4\log2$, $\log(5/4)<1/4$, and
$\log2>2/3$ give an upper bound
$8-16\log2+1/8<-61/24<-2$.
For completeness, $\pi<4$ follows from
$\pi/4=\int_0^1(1+x^2)^{-1}dx<1$; and the substitution
$x=(1+t)/(1-t)$ in $\log2=\int_1^2dx/x$ gives
$\log2=2\int_0^{1/3}(1-t^2)^{-1}dt>2/3$.
The logarithm bound follows by integrating $1/(1+x)<1$.
This proves OSP25 with $K_q$ as stated. Also
\[
0\le\mathcal E_k
\le \frac{4(\delta^2+\gamma^2)}{3\epsilon^2}+\log(1+K_q).
\tag{OSP26}
\]
Here $R^2\le\epsilon^2q^2/4$ and $R^2/q\le\delta^2+\gamma^2$.
Thus the logarithmic norm cost is bounded at fixed original
parameters, uniformly over all these coefficient degrees and both
unchanged source orders.

\subsection{Every relation determinant and the exact outer quotient}

For a nonzero polynomial $P$ and $a\ge0$ define in the unchanged monomial frame
\[
\begin{split}
D_a(P;s)&=\det\left[\int_{\mathbb R}y^{i+j}|P(y)|^2\,dm_s(y)
                                      \right]_{0\le i,j<a},\\
\mathcal D_a(Q;s)&=\det\left[\int_{\mathbb R}y^{2Q+i+j}\,dm_s(y)
                                      \right]_{0\le i,j<a},
\qquad D_0=\mathcal D_0=1 .
\end{split}\tag{OSP27}
\]
All nonempty Grams are positive definite, because a nonzero
polynomial has positive squared modulus on an open subinterval
of the strictly positive original measure. In the notation of
OSP17 their exact matrices are
$\mathsf J_P^*H^{(s)}_{Q+a-1}\mathsf J_P$ and
$\mathsf J_{0,Q}^*H^{(s)}_{Q+a-1}\mathsf J_{0,Q}$,
where $H^{(s)}_{Q+a-1}$ is the complete original moment Gram.
Consequently OSP20--26 prove for $1\le a\le2q+1$
\[
\left|\log D_a(P;s)-\log\mathcal D_a(Q;s)\right|
\le a\mathcal E_k,\qquad P=P_k\ \hbox{or}\ P_{k'} .
\tag{OSP28}
\]
To see the determinant step exactly, conjugate the actual Gram by
the inverse positive square root of the reference Gram. OSP20
places every eigenvalue of the resulting positive Hermitian
matrix in $[e^{-\mathcal E_k},e^{\mathcal E_k}]$. The logarithm
of its determinant is the sum of its $a$ eigenvalue logarithms,
proving OSP28. This is the full exterior determinant estimate,
including every polynomial coefficient direction. The estimate
also holds at $a=0$ with both sides zero.

For $t\in\{0,1,q,q+1\}$, let $L=t+\Delta-1$. The actual quotient
has source $\mathbb C[y]_{\le L}$ in measure
$|P_{k'}(y)|^2dm_s(y)$, target the increasing remainder frame of
$\mathbb C[y]/(P_{\partial,k})$, and its minimum Gram
$Q_{\partial,L}^{(s)}$. Define a comparison quotient with source
$\mathbb C[y]_{\le L}$ in measure $|y|^{2q'}dm_s(y)$, target the
increasing remainder frame of $\mathbb C[y]/(y^\Delta)$, and
minimum Gram $Q_{0,L}^{(s)}$. Its source measure still has order
$s$ and the full mass factor in OSP2.

For precision, the exact map relating the coefficient presentations
is the degree-preserving isomorphism
\[
\mathcal U_L(y^j)=
\begin{cases}
y^j,&0\le j<\Delta,\\
P_{\partial,k}(y)y^{j-\Delta},&\Delta\le j\le L.
\end{cases}\tag{OSP29}
\]
Its coefficient matrix has diagonal entries one and no terms of
degree greater than that of its input column. It is therefore
invertible with determinant one. If $\mathsf R_{\partial,L}$
is division remainder modulo $P_{\partial,k}$ and
$\mathsf R_{0,L}$ is division remainder modulo $y^\Delta$,
then
$\mathsf R_{\partial,L}\mathcal U_L=\mathsf R_{0,L}$:
the first $\Delta$ basis vectors map to their stated remainders,
and every later vector maps into the actual ideal. This proves
the precise morphism between the quotient coefficient spaces;
no equality of their metrics is assumed.

Using the original monic relation frames on the two sources gives
the exact identities
\[
\det Q_{\partial,t+\Delta-1}^{(s)}
 =\frac{D_{t+\Delta}(P_{k'};s)}{D_t(P_k;s)},\qquad
\det Q_{0,t+\Delta-1}^{(s)}
 =\frac{\mathcal D_{t+\Delta}(q';s)}{\mathcal D_t(q;s)} .
\tag{OSP30}
\]
For proof, append to $1,y,\ldots,y^{\Delta-1}$ the ideal frame
$P_{\partial,k},yP_{\partial,k},\ldots,y^{t-1}P_{\partial,k}$.
Its source coefficient determinant is one by OSP29. The relation
Gram is $D_t(P_k;s)$ because $P_{k'}P_{\partial,k}=P_k$.
Completing the square in the ideal coordinates gives the minimum
quotient Gram as the Schur complement; the determinant of the
full source Gram is the product of that determinant and the
relation determinant. This proves the first identity.
For the second repeat the same calculation with $y^\Delta$:
$y^{q'}y^\Delta=y^q$, with the original source order $s$.
The same proof includes $t=0$ using the empty relation block.

OSP28 applies to every index here, since
$t+\Delta\le q+\Delta+1\le2q+1$. Define the exact endpoint
difference $\varepsilon_t^{(s)}$ by
\[
\log\det Q_{\partial,t+\Delta-1}^{(s)}
=\log\mathcal D_{t+\Delta}(q';s)-\log\mathcal D_t(q;s)
 +\varepsilon_t^{(s)} .
\]
Then, separately at each of the four endpoints,
\[
|\varepsilon_t^{(s)}|\le(2t+\Delta)\mathcal E_k.
\tag{OSP31}
\]
In particular set
\[
\begin{split}
\mathcal F_{0,k}^{(s)}
={}&\log\mathcal D_\Delta(q';s)
    +\log\mathcal D_{\Delta+1}(q';s)
    -\log\mathcal D_{q+\Delta}(q';s)
    -\log\mathcal D_{q+\Delta+1}(q';s)\\
 &-\log\mathcal D_1(q;s)
    +\log\mathcal D_q(q;s)+\log\mathcal D_{q+1}(q;s).
\end{split}\tag{OSP32}
\]
The actual four-endpoint quantity of CTV13 therefore obeys the
proved finite estimate
\[
\boxed{\displaystyle
\left|\mathcal F_{\partial,k}^{(s)}-\mathcal F_{0,k}^{(s)}\right|
\le4(q+\Delta+1)\mathcal E_k
=O_{\delta,\gamma}(q)=o(kq),\qquad s=1,k .}
\tag{OSP33}
\]
Indeed its exact error is
$\varepsilon_0+\varepsilon_1-\varepsilon_q-\varepsilon_{q+1}$;
summing OSP31 gives
$\Delta+(\Delta+2)+(2q+\Delta)+(2q+\Delta+2)
=4(q+\Delta+1)$. OSP26 and $\Delta<q$ prove the last two
estimates. This calculation uses all four source degrees
$\Delta-1,\Delta,q+\Delta-1,q+\Delta$ and every source and
relation determinant in OSP30.

\subsection{The resulting same-source moment target}

The moments in OSP27 are completely specified by the original
analytic identity OSP9:
\[
\mu_j^{(s)}=\int y^j\,dm_s(y)
=M_s\left.\frac{d^j}{dz^j}(\cos z)^{-s/2}\right|_{z=0},
\qquad
\mathcal D_a(Q;s)=
\det[\mu_{2Q+i+j}^{(s)}]_{0\le i,j<a}.
\tag{OSP34}
\]
The differentiated integral is justified by OSP7, or directly by
the finite exponential moments on a slightly larger interval.
Odd moments vanish because the density is even. Reordering rows
and columns together by parity therefore proves the exact further
factorization
\[
\mathcal D_a(Q;s)
=\det[\mu_{2(Q+i+j)}^{(s)}]_{0\le i,j<\lceil a/2\rceil}
 \det[\mu_{2(Q+1+i+j)}^{(s)}]_{0\le i,j<\lfloor a/2\rfloor}.
\tag{OSP35}
\]
There is no permutation sign, since the same permutation is applied
to both rows and columns. Empty factors are one.
In OSP32 the coefficient of $\log M_s$ is
$\Delta+(\Delta+1)-(q+\Delta)-(q+\Delta+1)-1+q+(q+1)=0$.
This is an exact cancellation only after all original moments
and determinant dimensions have been retained.

Equation OSP33 proves an error-controlled comparison with the exact
finite determinant difference OSP32. It does not obtain a first
variation by differentiating a leading $q^2$ law. In particular no
numerical value of the limit of $\mathcal F_{0,k}^{(s)}/(kq)$ is
asserted by this comparison.

% END CURRENT OSP

\clearpage
% BEGIN CURRENT OAR
\section{The outer-root and actual residual bounds in the original arithmetic return}
\label{sec:OAR}
Retain the complete original data of CEP, CTV and ROQ. In particular
$m=1$, $k=4l+1\ge9$, $q=(k+1)^2$, $q'=(k-7)^2$,
$\Delta=16k-48$, $r=8k-16$ and $j=8k-32$. The quartet, full
amplitude unit, period and branch are the specified original ones on
the proved five-orbit locus. The source orders remain $s\in\{1,k\}$.
Every matrix below is the actual matrix in those chapters, with its
original mass and coefficient frame.

\subsection{The two finite error endpoints retain their signs}
Write
\[
 \mathcal N_+=\{q-1,q\},\qquad
 \mathcal N_-=\{2q-1,2q\},\qquad
 \epsilon_{A,k}^{(s)}=2\sum_{N\in\mathcal N_+\cup\mathcal N_-}E_N^{(s)}.
 \tag{OAR1}
\]
Here $E_N^{(s)}$ is the explicit full conductor constant of CEP35.
The nonnegative constants $L_N^{\rm tail},U_N^{\rm tail}$ are
exactly CTV18 in the same source order and with the original
$q'$, $\delta$, $\gamma$. Define
\[
 \begin{split}
 a_{k,s}^-&=v\left(\sum_{N\in\mathcal N_+}L_N^{\rm tail}
                         +\sum_{N\in\mathcal N_-}U_N^{\rm tail}\right),\\
 a_{k,s}^+&=v\left(\sum_{N\in\mathcal N_+}U_N^{\rm tail}
                         +\sum_{N\in\mathcal N_-}L_N^{\rm tail}\right),\\
 e_{k,s}^-&=\epsilon_{A,k}^{(s)}+a_{k,s}^-,\qquad
 e_{k,s}^+=\epsilon_{A,k}^{(s)}+a_{k,s}^+.
 \end{split}\tag{OAR2}
\]
These are finite quantities for every original period in the stated
domain. For each fixed original period, both $e_{k,s}^\pm=o(kq)$.
Indeed CEP36 proves this for $\epsilon_{A,k}^{(s)}$ and CTV18--20
proves it for each of the four tail contributions, uniformly over
the possible indices $0\le v\le80$ and both original source orders.
When $v=0$, both $a_{k,s}^\pm$ are exactly zero.

Set $\Phi_{k,s}=F_{\Xi,k}^{(s)}$ and retain the full outer quotient
$\mathcal F_{\partial,k}^{(s)}$ of CTV11--13. The exact identity is
\[
 \begin{gathered}
 F_K^{(s)}=\mathcal F_{\partial,k}^{(s)}-\Phi_{k,s}+\mathfrak e_{k,s},\\
 \mathfrak e_{k,s}=\mathcal E_{A,k}^{(s)}
     +\sum_{N\in\mathcal N_+}\mathcal R_{N,v}^{(s)}
     -\sum_{N\in\mathcal N_-}\mathcal R_{N,v}^{(s)},\\
 -e_{k,s}^-\le\mathfrak e_{k,s}\le e_{k,s}^+.
 \end{gathered}\tag{OAR3}
\]
The equality is CTV22, keeping its original two errors. CTV19 says
$-vL_N^{\rm tail}\le\mathcal R_{N,v}^{(s)}\le vU_N^{\rm tail}$.
At a positive endpoint this contributes lower bound $-vL_N^{\rm tail}$
and upper bound $vU_N^{\rm tail}$; at a negative endpoint it
contributes lower bound $-vU_N^{\rm tail}$ and upper bound
$vL_N^{\rm tail}$. Adding these four inequalities and the bounds
$-\epsilon_{A,k}^{(s)}\le\mathcal E_{A,k}^{(s)}
\le\epsilon_{A,k}^{(s)}$ proves each sign of OAR3.

\subsection{Original positive increments give explicit kernel caps}
The full value of $\Phi_{k,s}$ is ROQ19's finite sum
\[
 \Phi_{k,s}=\log(1+\eta_{q-1})
       +2\sum_{N=q}^{2q-2}\log(1+\eta_N)
       +\log(1+\eta_{2q-1}),\qquad
 \lambda_s=\log(1+\eta_{q-1})\ge0.
 \tag{OAR4}
\]
Every $\eta_N$ here is ROQ17's original constrained column, including
the complete conductor--observation cross covariance and its
positive denominator. No new observation is selected.
ROQ21 gives, in the original AKS notation,
\[
 \lambda_s\le\Phi_{k,s}\le
 V_s:=j(\mathcal L_{s,q}+\mathcal L_{s,q+1})-\lambda_s.
 \tag{OAR5}
\]
Retain the complete earlier kernel interval $L_s\le F_K^{(s)}\le U_s$
of AKJ2/MRI4e, with its pointwise minimum of the KAF and AKS bounds.
Define two explicit quantities
\[
 \begin{split}
 L_s^{\partial}&=\max\left\{L_s,
       \mathcal F_{\partial,k}^{(s)}-e_{k,s}^--V_s\right\},\\
 U_s^{\partial}&=\min\left\{U_s,
       \mathcal F_{\partial,k}^{(s)}+e_{k,s}^+-\lambda_s\right\}.
 \end{split}\tag{OAR6}
\]
Then
\[
 \begin{gathered}
 L_s^{\partial}\le F_K^{(s)}\le U_s^{\partial},\\
 \mathcal F_{\partial,k}^{(s)}-\Phi_{k,s}-e_{k,s}^-
       \le F_K^{(s)}
       \le\mathcal F_{\partial,k}^{(s)}-\Phi_{k,s}+e_{k,s}^+.
 \end{gathered}\tag{OAR7}
\]
For the first line substitute both ends of OAR5 into OAR3 and
intersect with the retained interval. The second line substitutes
the full finite value OAR4 into OAR3. Intersecting these intervals
preserves every bound. Both contain the actual value, so their
intersection is nonempty. The dimension $j$ is used only in the
proved finite comparison OAR5. It assigns no leading determinant.
Each additional evaluated positive term of OAR4 can be included
in the upper subtraction in OAR6 by the same inequality.

\subsection{The correlated arithmetic kernel and boundary}
Keep the original arithmetic total $\mathcal B_k^{\rm ar}$,
the scalar $\delta_k^\sigma$ and the four source-difference constants
$C_K^\pm,C_B^\pm$ of MRI4f. They satisfy, for the actual
$d=F_K^{\rm ar}-F_K^{(1)}$,
\[
 \max\{-C_K^-,\delta_k^\sigma-C_B^+\}\le d
       \le\min\{C_K^+,\delta_k^\sigma+C_B^-\}.
 \tag{OAR8}
\]
This is the pair of original kernel bounds together with the
boundary bounds, transferred by
$d+(F_B^{\rm ar}-F_B^{(1)})=\delta_k^\sigma$.
It retains their common scalar and signs. Define
\[
 \begin{split}
 L_{\rm ar}^{\partial}
   &=\max\{L_{\rm ar},L_1^{\partial}-C_K^-,
                    L_1^{\partial}+\delta_k^\sigma-C_B^+\},\\
 U_{\rm ar}^{\partial}
   &=\min\{U_{\rm ar},U_1^{\partial}+C_K^+,
                    U_1^{\partial}+\delta_k^\sigma+C_B^-\}.
 \end{split}\tag{OAR9}
\]
Here $[L_{\rm ar},U_{\rm ar}]$ is the complete interval of MRI4f.
Adding OAR8 to the source-$1$ first line of OAR7 proves
\[
 \begin{gathered}
 L_{\rm ar}^{\partial}\le F_K^{\rm ar}\le U_{\rm ar}^{\partial},\\
 \mathcal B_k^{\rm ar}-U_{\rm ar}^{\partial}
       \le F_B^{\rm ar}
       \le\mathcal B_k^{\rm ar}-L_{\rm ar}^{\partial}.
 \end{gathered}\tag{OAR10}
\]
The second line subtracts the first from the exact original identity
$F_K^{\rm ar}+F_B^{\rm ar}=\mathcal B_k^{\rm ar}$.
It uses the same arithmetic total at both endpoints. If the second
line of OAR7 is used instead, exactly the same addition and
subtraction yield the further interval containing the full residual
sum. Thus every exact positive residual subtraction is available to
the original arithmetic kernel and its correlated boundary.

\subsection{The signed mixed return with its complete finite error}
Write
$\mathfrak S_k^{\rm mix}=F_K^{\rm ar}+F_B^{\rm ar}-F_B^{(k)}$.
The original BRI6d identity is
$\mathfrak S_k^{\rm mix}=F_K^{(k)}+\mathcal H_k+\delta_k^\sigma$.
Consequently OAR3 gives the exact current expression
\[
 \begin{gathered}
 \mathfrak S_k^{\rm mix}=\mathcal F_{\partial,k}^{(k)}-\Phi_{k,k}
              +\mathcal H_k+\delta_k^\sigma+\mathfrak e_{k,k},\\
 -e_{k,k}^-\le\mathfrak e_{k,k}\le e_{k,k}^+,
 \qquad e_{k,k}^\pm=o(kq).
 \end{gathered}\tag{OAR11}
\]
This substitution proves every sign, keeping the original Gamma
return and arithmetic scalar. Addition of their unchanged values
to the source-$k$ bounds OAR7 yields the corresponding finite
intervals for this same signed quantity. The full original
four-endpoint metrics, outer-root polynomial, weighted source,
period-dependent certificates and graph denominator remain in
OAR11 through their proved maps CTV and ROQ. The next leading-scale
calculation is the difference of these two explicit quotient returns.

% END CURRENT OAR

\clearpage
% BEGIN CURRENT PMR
\section{The same-source Gamma moment comparison in the arithmetic return}
\label{sec:PMR}
Retain every original parameter, matrix and error in OAR1--11.
The complete OSP comparison supplies an additional bound directly
on the actual outer-root quotient. Set $\epsilon=2^{-32}$ and
retain its explicit range
\[
 k\ge29,\qquad k\sqrt{\delta^2+\gamma^2}
           \le\epsilon(k+1)^2/2.
 \tag{PMR1}
\]
This range contains all sufficiently large original $k=4l+1$ for
each fixed original quartet: $k\ge\max\{29,
2\sqrt{\delta^2+\gamma^2}/\epsilon\}$ is sufficient because
$(k+1)^2\ge k^2$. For the remaining finite values the full earlier
OAR estimates remain valid without using this comparison.

Keep OSP24's nonnegative $\mathcal E_k$ and put
\[
 c_k^\circ=4(q+\Delta+1)\mathcal E_k,
 \qquad e_{k,s}^{0,-}=e_{k,s}^-+c_k^\circ,
 \qquad e_{k,s}^{0,+}=e_{k,s}^++c_k^\circ.
 \tag{PMR2}
\]
OSP26 and OSP33 prove $c_k^\circ=O_{\delta,\gamma}(q)=o(kq)$
on PMR1, uniformly for the unchanged source orders $s=1,k$.
The actual outer quotient differs from OSP32's exact seven
original Gamma moment determinants by at most $c_k^\circ$.
Consequently the full identity and signed bounds are
\[
 \begin{gathered}
 F_K^{(s)}=\mathcal F_{0,k}^{(s)}-\Phi_{k,s}+\mathfrak e^0_{k,s},\\
 \mathfrak e^0_{k,s}=\mathfrak e_{k,s}
       +\mathcal F_{\partial,k}^{(s)}-\mathcal F_{0,k}^{(s)},\\
 -e_{k,s}^{0,-}\le\mathfrak e^0_{k,s}\le e_{k,s}^{0,+},
 \qquad e_{k,s}^{0,\pm}=o(kq).
 \end{gathered}\tag{PMR3}
\]
The first two equalities are OAR3 with the displayed exact
difference added and subtracted. The two bounds follow by adding
the interval $[-c_k^\circ,c_k^\circ]$ of OSP33 to the unequal
OAR3 endpoints. No metric equality is substituted for the
proved OSP coefficient maps and determinant comparison.
The residual sum $\Phi_{k,s}$ retains all actual period data.

Let $V_s$ and $\lambda_s$ be OAR4--5's exact residual upper
bound and first positive increment. Define
\[
 \begin{split}
 L_s^0&=\max\{L_s^\partial,
                \mathcal F_{0,k}^{(s)}-e_{k,s}^{0,-}-V_s\},\\
 U_s^0&=\min\{U_s^\partial,
                \mathcal F_{0,k}^{(s)}+e_{k,s}^{0,+}-\lambda_s\}.
 \end{split}\tag{PMR4}
\]
Then $L_s^0\le F_K^{(s)}\le U_s^0$. To prove this, insert
$\lambda_s\le\Phi_{k,s}\le V_s$ into PMR3 and intersect
with OAR7. Using the full finite sum for $\Phi_{k,s}$ in PMR3
gives the sharper interval centred at
$\mathcal F_{0,k}^{(s)}-\Phi_{k,s}$. Both retain the same
unequal finite errors of PMR2.

The exact arithmetic correlations of OAR8 give
\[
 \begin{split}
 L_{\rm ar}^0&=\max\{L_{\rm ar}^\partial,L_1^0-C_K^-,
                                      L_1^0+\delta_k^\sigma-C_B^+\},\\
 U_{\rm ar}^0&=\min\{U_{\rm ar}^\partial,U_1^0+C_K^+,
                                      U_1^0+\delta_k^\sigma+C_B^-\}.
 \end{split}\tag{PMR5}
\]
Adding those correlations to PMR4 at $s=1$ proves
$L_{\rm ar}^0\le F_K^{\rm ar}\le U_{\rm ar}^0$.
Subtracting from the exact same arithmetic total gives
\[
 \mathcal B_k^{\rm ar}-U_{\rm ar}^0
      \le F_B^{\rm ar}\le\mathcal B_k^{\rm ar}-L_{\rm ar}^0.
 \tag{PMR6}
\]
Thus the actual kernel and boundary retain their common scalar
and both original source comparisons.

Finally the original signed-return identity BRI6d and PMR3 at
source order $k$ give
\[
 \begin{gathered}
 F_K^{\rm ar}+F_B^{\rm ar}-F_B^{(k)}
   =\mathcal F_{0,k}^{(k)}-\Phi_{k,k}
      +\mathcal H_k+\delta_k^\sigma+\mathfrak e^0_{k,k},\\
 -e_{k,k}^{0,-}\le\mathfrak e^0_{k,k}\le e_{k,k}^{0,+}.
 \end{gathered}\tag{PMR7}
\]
This follows by substituting into
$F_K^{\rm ar}+F_B^{\rm ar}-F_B^{(k)}=F_K^{(k)}+
\mathcal H_k+\delta_k^\sigma$; every sign is retained.
The entire original same-source moment expression is OSP32--35.
The next asymptotic calculation must evaluate it with the actual
residual sum and an error small compared with $kq$.

\paragraph{The exact finite endpoints under this comparison.}
The interval transport in PMR4--6 preserves the already retained
outer-root endpoints exactly:
\[
 L_s^0=L_s^\partial,\qquad U_s^0=U_s^\partial,\qquad
 L_{\rm ar}^0=L_{\rm ar}^\partial,\qquad
 U_{\rm ar}^0=U_{\rm ar}^\partial.
 \tag{PMR8}
\]
For proof, OSP33 gives
$\mathcal F_{0,k}^{(s)}-c_k^\circ\le\mathcal F_{\partial,k}^{(s)}
\le\mathcal F_{0,k}^{(s)}+c_k^\circ$.
The new lower candidate in PMR4 is therefore at most
$\mathcal F_{\partial,k}^{(s)}-e_{k,s}^--V_s$, which is at
most $L_s^\partial$ by OAR6. The new upper candidate is at
least $\mathcal F_{\partial,k}^{(s)}+e_{k,s}^+-\lambda_s$,
which is at least $U_s^\partial$. Taking the displayed maximum
and minimum proves the source equalities. Substitution in PMR5
and the defining maxima and minima OAR9 proves the arithmetic
equalities; subtracting from the same total proves equality of
the boundary endpoints as well. Thus this comparison supplies
the exact same-source moment target and controlled error; a
further numerical improvement must use sharper evaluated moment
or residual information. The preceding ROQ/OAR bounds remain
the proved finite intervals throughout this transport.

% END CURRENT PMR


\clearpage
% BEGIN CURRENT OFV
% Uses complete proved providers OSP1--35, EIQ1--33, and QLG1--32.
% No leading law is differentiated without its finite partition remainder.
\section{The leading first variation of the actual outer-strip quotient}
\label{sec:OFV}

Retain every original object and both original source orders of
OSP1--35 and CTV1--22. In addition to the explicit finite conditions
OSP23, take $k\ge257$. This is an eventual threshold on the same
fixed original parameters. Put
\[
\vartheta=\pi/2,\qquad
\ell_s=(s-1)/4,\qquad
\beta_s=\frac{(2\pi)^{s/2}}{\sqrt2\,\Gamma(s/2)}
\quad(s=1,k).
\tag{OFV1}
\]
The symbol $\vartheta$ is the fixed Gamma tail rate; it is not either
original source centre.

We use the equilibrium of EIQ1--33 on its original positive half-line:
\[
\begin{gathered}
V_{\alpha,\beta}(x)=\beta\sqrt{x}-\alpha\log x,\quad
\mu_{\alpha,\beta}(dx)=\rho_{\alpha,\beta}(x)\,dx,\\
F(\alpha,\beta)=
\iint\log|x-y|\,d\mu_{\alpha,\beta}(x)d\mu_{\alpha,\beta}(y)
-\int V_{\alpha,\beta}\,d\mu_{\alpha,\beta}.
\end{gathered}
\tag{OFV2}
\]
Here $\rho$, its two endpoints, the Euler equality and the entire
exterior inequality are exactly the explicit proved EIQ formulas.
There is no new equilibrium assumption. EIQ29--30 proves smoothness
and
$F_\alpha(\alpha,\beta)=\int\log x\,d\mu_{\alpha,\beta}(x)$.
QLG31, with its complete circle and quantile proofs, gives
\[
\left|\log J_n(\alpha,\beta)-n^2F(\alpha,\beta)\right|\le C_0n,
\quad
J_n=\frac1{n!}\int_{(0,\infty)^n}
\prod_{i<j}(x_i-x_j)^2
\prod_i e^{-nV_{\alpha,\beta}(x_i)}\,dx_i
\tag{OFV3}
\]
for every $n\ge1$ and every
$(\alpha,\beta)\in[3/2,5/2]\times[\pi/2,3\pi/2]$.
The constant $C_0$ is the explicit finite constant QLG30.
We use precisely this uniform finite error, including the factor
$1/n!$; the conclusion below is not a derivative of an earlier
unquantified $q^2$ limit.

\subsection{Exact original-order recurrence and a small metric error}

For $a\ge0$, define
\[
\mathcal D_a(Q;s)
=\det\left[\int_{\mathbb R}y^{2Q+i+j}\,dm_s(y)\right]_{0\le i,j<a},
\qquad
\mathcal A_a(Q;\vartheta)
=\det\left[\int_{\mathbb R}y^{2Q+i+j}e^{-\vartheta|y|}\,dy
                                      \right]_{0\le i,j<a},
\tag{OFV4}
\]
with empty determinants one. The $Q$ used here is an integer;
all the instances below have $Q\ge q/2$.
The original recurrence OSP12 gives exactly
\[
|y|^{2Q}dm_s(y)=
\beta_s |y|^{2Q}
\prod_{j=0}^{\ell_s-1}(y^2+t_j^2)\,d\sigma(y),
\qquad t_j=2j+\tfrac12.
\tag{OFV5}
\]
For each actual source centre $c_0\in\{k/2,k/2-4\}$, the
coordinate polynomial of this multiplier is
$f_{\ell_s,c_0}(S)=\prod_j(S-c_0-t_j)$.
Its exact pullback satisfies
$f_{\ell_s,c_0}(c_0+iy)=\prod_j(iy-t_j)
=i^{\ell_s}\prod_j(y+it_j)$.
Thus multiplication by
$y^Q\prod_{j=0}^{\ell_s-1}(y+it_j)$, together with the
retained phase $i^{\ell_s}$, is the exact coefficient map
into the original one-factor Gamma source. Its squared modulus
is the displayed density multiplier, with no change to any $t_j$
or to $\beta_s$.

Set $T=\epsilon q$, $\epsilon=2^{-32}$ as in OSP23, and
\[
S_{\ell_s}=\sum_{j=0}^{\ell_s-1}t_j^2
=\frac{\ell_s(16\ell_s^2-12\ell_s-1)}{12},\qquad
r_{\ell_s}=\frac{S_{\ell_s}}{T^2}+\log(1+K_q),
\tag{OFV6}
\]
where $K_q$ is exactly OSP24. At $\ell_s=0$ we may and do take
$r_0=0$ instead, since the source identity is exact.
For $Q\in\{q',q\}$, $0\le a\le2q+1$,
\[
0\le
\log\mathcal D_a(Q;s)-a\log\beta_s
-\log\mathcal D_a(Q+\ell_s;1)
\le ar_{\ell_s}.
\tag{OFV7}
\]

Here is a proof of the full-form comparison underlying OFV7.
For every real $y$,
$|y|^{2Q}\prod_j(y^2+t_j^2)\ge |y|^{2(Q+\ell_s)}$.
On $|y|\ge T$, the logarithm of the ratio is at most
$\sum_jt_j^2/T^2=S_{\ell_s}/T^2$, by $\log(1+u)\le u$.
For $|y|<T$, its numerator is at most
$(T^2+(k/2)^2)^{Q+\ell_s}$.
OSP15 on the complete coefficient space of degree at most $a-1$,
OSP13 with source $s=1$, and mass $M_1=\sqrt{2\pi}$ therefore
bound its inner quadratic form by $K_q$ times the reference
quadratic form on $[q,2q]$. This is exactly the calculation
OSP22--25 with exponent $Q+\ell_s\ge q/2$, degree $a-1\le2q$,
and root bound $k/2\le R\le T/2$.
No upper bound on that exponent is needed in OSP25, since its
base $(5/4)\epsilon^2$ is less than one.
The full upper ratio is consequently at most
$e^{S_{\ell_s}/T^2}+K_q\le e^{r_{\ell_s}}$.
The lower full ratio is at least one by the pointwise lower
inequality. The positive-Gram eigenvalue argument OSP28 now proves
OFV7 in the original coefficient frames.
Since $S_{\ell_s}\le\ell_s k^2/4$ and $q\ge k^2$,
$r_{\ell_s}=O(1/k)$, with its stated fixed cutoff constant.

\subsection{A uniform finite Gamma-to-exponential determinant estimate}

Put $h=1/q$, and define the actual finite constants
\[
L_h=C_L(\sin h)^{3/2},\qquad
U_h=\frac{\Gamma(1/4)^2}{2\pi}(\sin h)^{-1/2}.
\tag{OFV8}
\]
Then the original one-factor Gamma density satisfies on the entire
real line
\[
L_he^{-(\vartheta+h)|y|}
\le\frac{d\sigma}{dy}(y)
\le U_he^{-(\vartheta-h)|y|}.
\tag{OFV9}
\]
Indeed apply OSP7 to $\Gamma(1/4+iy/2)$ with
$\theta=\pi/2-h$ for the upper bound.
For the lower bound apply it to $\Gamma(3/4-iy/2)$ in the exact
reflection product proved before OSP11. The upper bound on that
squared factor is
$\Gamma(3/4)^2(\sin h)^{-3/2}e^{-(\vartheta-h)|y|}$.
After division, the identity and
$\cosh(\pi y)\le e^{\pi|y|}$ give exactly the lower side of OFV9.
In particular both inequalities hold near zero as well as in
the tails.

The exact change of variables $y=\vartheta^{-1}u$ in OFV4 gives
\[
\mathcal A_a(Q;\vartheta)
=\vartheta^{-a(2Q+a)}\mathcal A_a(Q;1).
\tag{OFV10}
\]
The exponent follows from $a$ moment Jacobians, $2Q$ powers in
each, and the row and column powers
$\sum_{i=0}^{a-1}i+\sum_{j=0}^{a-1}j=a(a-1)$.
Applying OFV9 as full quadratic-form inequalities, and then using
OFV10, proves the separate finite bounds
\[
\begin{split}
a\log L_h-a(2Q+a)\log(1+h/\vartheta)
&\le\log\mathcal D_a(Q;1)-\log\mathcal A_a(Q;\vartheta)\\
&\le a\log U_h-a(2Q+a)\log(1-h/\vartheta).
\end{split}\tag{OFV11}
\]
Thus the absolute difference is at most
\[
B_a(Q)=a\max\{|\log L_h|,|\log U_h|\}
       +a(2Q+a)\frac{h}{\vartheta-h}.
\tag{OFV12}
\]
Here $-\log(1-u)\le u/(1-u)$ and $\log(1+u)\le u$
justify the second summand. Since
$2h/\pi\le\sin h\le h$ for $0<h\le\pi/2$,
the first summand is $O(a\log q)$.
The sine inequalities follow from concavity and the endpoint chord,
and by integrating $\cos t\le1$.
For all the present $a\le2q+1$, $Q\le q+\ell_s\le2q$,
OFV12 is $O(q\log q)$ uniformly in $a,Q$.
Combining OFV7 and OFV11 gives the error-controlled identity
\[
\log\mathcal D_a(Q;s)
=a\log\beta_s+\log\mathcal A_a(Q+\ell_s;\vartheta)
  +R_{a,Q,s},\qquad
|R_{a,Q,s}|\le ar_{\ell_s}+B_a(Q+\ell_s).
\tag{OFV13}
\]
This is a comparison on complete original polynomial coefficient
spaces, with the order shift proved by the literal multiplier
OFV5; it is not an assigned replacement of the original source order.

\subsection{Direct low-dimension estimates, outside the compact equilibrium range}

Define
\[
g(Q)=2Q\log(2Q/\vartheta)-2Q.
\tag{OFV14}
\]
For $Q\ge q/2$ and $1\le a\le\Delta+1$ we prove
\[
\log\mathcal A_a(Q;\vartheta)
=a g(Q)+O\bigl(a^2\log(Q+a)+a\log(Q+a)\bigr),
\tag{OFV15}
\]
with constants independent of $k,Q,a$ in these ranges.
There is no use of OFV3 here; its compact parameter rectangle
does not contain these low-dimension ratios.

For the lower bound take the full interval
$[y_0,y_0+1]$, $y_0=2Q/\vartheta$. On it
$y^{2Q}e^{-\vartheta y}\ge e^{g(Q)-\vartheta}$.
Consequently its moment Gram is bounded below by
$e^{g(Q)-\vartheta}$ times the unweighted monomial Gram on that interval.
Translation has a triangular coefficient matrix with diagonal one,
so the latter Gram determinant equals the one on $[0,1]$.
The Rodrigues norms proved in OSP15 show that it is exactly
\[
\prod_{j=0}^{a-1}
\frac1{(2j+1)\binom{2j}{j}^{\,2}}.
\]
Indeed $\binom{2j}{j}$ is the leading coefficient of
$L_j(2x-1)$, whose squared norm is $1/(2j+1)$.
Using $\binom{2j}{j}\le4^j$ therefore proves
\[
\log\mathcal A_a(Q;\vartheta)
\ge a g(Q)-a\vartheta-a\log(2a)-2a(a-1)\log2.
\tag{OFV16}
\]

For the upper bound, Gram--Schmidt shows that a positive Gram
determinant is at most the product of the squared norms of its
displayed vectors: each orthogonal component has norm no greater
than its original vector. The exact monomial moments give
\[
\log\mathcal A_a(Q;\vartheta)
\le\sum_{j=0}^{a-1}
\log\left(\frac{2(2(Q+j))!}{\vartheta^{2(Q+j)+1}}\right).
\tag{OFV17}
\]
For every positive integer $m$, integrating the increasing function
$\log x$ between successive integers proves
$m\log m-m+1\le\log(m!)\le m\log m-m+1+\log m$.
Each summand in OFV17 is consequently at most
$g(Q+j)+|\log(2/\vartheta)+1|+\log(2(Q+j))$.
Furthermore $g'(u)=2\log(2u/\vartheta)$, so
$g(Q+j)-g(Q)\le2j\log(2(Q+a)/\vartheta)$ in the present range.
This proves the explicit upper bound
\[
\begin{split}
\log\mathcal A_a(Q;\vartheta)\le{}&
a g(Q)+a\{|\log(2/\vartheta)+1|+\log(2(Q+a))\}\\
&+a(a-1)\log(2(Q+a)/\vartheta).
\end{split}\tag{OFV18}
\]
OFV16 and OFV18 prove OFV15 with its complete lower and upper
errors. They also apply at $a=1$.

\subsection{High-dimension partitions and their controlled first variation}

For an integer $a$ write $n_e=\lceil a/2\rceil$ and
$n_o=\lfloor a/2\rfloor$. The two exact changes of variables
$y=\sqrt{x}$ and $y=-\sqrt{x}$ give, with their two positive
Jacobians, the parity identity
\[
\mathcal A_a(Q;\vartheta)
=Z_{n_e}(Q-\tfrac12;\vartheta)\,
 Z_{n_o}(Q+\tfrac12;\vartheta),
\tag{OFV19}
\]
where
\[
Z_n(A;\vartheta)
=\det\left[\int_0^\infty x^{A+i+j}e^{-\vartheta\sqrt{x}}\,dx
                                      \right]_{0\le i,j<n}.
\]
The two branches add to the density $x^{-1/2}e^{-\vartheta\sqrt{x}}$;
thus no factor two is missing from OFV19. The odd block has the
additional factor $x$. Expanding the two Vandermonde determinants
and integrating each finite permutation sum proves the exact
Heine identity for this moment determinant:
\[
Z_n(A;\vartheta)=\frac1{n!}\int_{(0,\infty)^n}
\prod_{i<j}(x_i-x_j)^2\prod_i x_i^Ae^{-\vartheta\sqrt{x_i}}\,dx_i.
\]
Changing variables $x_i=n^2u_i$ now gives
\[
Z_n(A;\vartheta)=n^{2n(A+n)}J_n(A/n,\vartheta).
\tag{OFV20}
\]
There are $2n(A+1)$ powers from the weights and Jacobians
and $2n(n-1)$ from the squared Vandermonde, proving the exponent.

For the high instances we use exactly
\[
(a,Q)\in
\{(q,q+\ell_s),(q+1,q+\ell_s),
  (q+\Delta,q-\Delta+\ell_s),
  (q+\Delta+1,q-\Delta+\ell_s)\}.
\tag{OFV21}
\]
Since $k\ge257$, $\Delta/q\le16/k<1/16$ and
$\ell_s/q\le1/(4k)<1/64$.
For both parity blocks, elementary rounding therefore gives
\[
\frac{15q}{32}\le n\le\frac{9q}{16},\qquad
\frac{29q}{32}\le A\le\frac{33q}{32}.
\]
The first inequalities use $q\ge32$ to absorb the possible
rounding by $1/2$ and the extra endpoint $1$.
Hence $29/18\le A/n\le11/5$, contained strictly in $[3/2,5/2]$.
Every partition in OFV20 is thus in the proved parameter rectangle
OFV3, with $\beta=\vartheta=\pi/2$.

Define the smooth explicit function
\[
\mathcal H(a,Q)=
a(2Q+a)\log(a/2)+\frac{a^2}{2}F(2Q/a,\vartheta)
\quad(a>0,\ Q>0).
\tag{OFV22}
\]
Equations OFV3 and OFV19--20 give uniformly over OFV21
\[
\log\mathcal A_a(Q;\vartheta)=\mathcal H(a,Q)+O(q\log q).
\tag{OFV23}
\]
To verify the rounding error rather than omit it, set
$G(n,A)=2n(A+n)\log n+n^2F(A/n,\vartheta)$.
Each actual pair $(n,A)$ differs from $(a/2,Q)$ by at most $1/2$
in each coordinate. On all connecting segments $n$ is bounded
above and below by positive constant multiples of $q$, while
$A/n$ remains in $[3/2,5/2]$.
Its derivatives are
\[
G_A=2n\log n+nF_\alpha,\qquad
G_n=2(A+2n)\log n+2(A+n)+2nF-AF_\alpha .
\]
They are $O(q\log q)$ there, since the exact EIQ functions
$F,F_\alpha$ are bounded on that compact interval. The mean value
formula bounds each rounding error by $O(q\log q)$.
The two QLG errors sum to at most $C_0a=O(q)$, and
$2G(a/2,Q)=\mathcal H(a,Q)$, proving OFV23.

The derivative needed for the first variation is computed without
altering the direction of the two original degree changes:
\[
(\partial_a-\partial_Q)\mathcal H(a,Q)
=2Q\log(a/2)+(2Q+a)+aF(2Q/a,\vartheta)
 -(Q+a)F_\alpha(2Q/a,\vartheta).
\tag{OFV24}
\]
EIQ29--30 proves that $F_\alpha$ is smooth; in particular
$F,F_\alpha,F_{\alpha\alpha}$ are bounded on the compact interval
just used. Differentiating the right side of OFV24 shows that
its two first partial derivatives are $O(\log q)$ on the segment
$a=q+u$, $Q=q+\ell_s-u$, $0\le u\le\Delta$ and on its
connections to $(q,q)$. Every term follows by the product rule;
the derivatives of $2Q/a$ are $2/a$ and $-2Q/a^2$, both $O(1/q)$.
The mean value formula and then integration in $u$ consequently give
\[
\begin{split}
&\mathcal H(q+\Delta,q-\Delta+\ell_s)
             -\mathcal H(q,q+\ell_s)\\
&\quad=
\Delta q\{2\log(q/2)+3+F(2,\vartheta)-2F_\alpha(2,\vartheta)\}
+O\bigl(\Delta(\Delta+\ell_s)\log q\bigr).
\end{split}\tag{OFV25}
\]
This remainder is $O(q\log q)$ because
$\Delta=O(k)$, $\ell_s=O(k)$ and $q=(k+1)^2$.
The additional endpoint $a+1$ in OFV21 changes $\mathcal H$
by $O(q\log q)$, directly from its displayed derivatives.
Thus all rounding and endpoint changes have been quantified.

\subsection{All seven determinants and the resulting leading constant}

Replace the seven moment determinants in OSP32 by the exact
exponential determinants in OFV13, with their respective
$Q+\ell_s$. Their $a\log\beta_s$ terms cancel because their
signed dimensions are exactly
\[
\Delta+(\Delta+1)-(q+\Delta)-(q+\Delta+1)-1+q+(q+1)=0.
\tag{OFV26}
\]
The sum of their absolute finite errors in OFV13 is
$O(q\log q)$, since the sum of their dimensions is
$4(q+\Delta+1)$, $r_{\ell_s}=O(1/k)$, and OFV12 applies
uniformly with $Q+\ell_s\le2q$.
The low terms, using OFV15--18, are
\[
\begin{split}
&\log\mathcal A_\Delta(q-\Delta+\ell_s;\vartheta)
+\log\mathcal A_{\Delta+1}(q-\Delta+\ell_s;\vartheta)
-\log\mathcal A_1(q+\ell_s;\vartheta)\\
&\qquad=2\Delta g(q)+O(q\log q)
=4\Delta q\{\log(2q/\vartheta)-1\}+O(q\log q).
\end{split}\tag{OFV27}
\]
Indeed replacing each $g(q-\Delta+\ell_s)$ or $g(q+\ell_s)$
by $g(q)$ costs at most
$O(((2\Delta+1)(\Delta+\ell_s)+\ell_s)\log q)=O(q\log q)$,
using $g'(u)=2\log(2u/\vartheta)$ and the actual degree ranges.
The individual low determinant errors are
$O((\Delta+1)^2\log q)=O(q\log q)$.

The remaining four high terms are, by OFV23--25,
\[
\begin{split}
&\log\mathcal A_q(q+\ell_s;\vartheta)
+\log\mathcal A_{q+1}(q+\ell_s;\vartheta)\\
&\quad-\log\mathcal A_{q+\Delta}(q-\Delta+\ell_s;\vartheta)
-\log\mathcal A_{q+\Delta+1}(q-\Delta+\ell_s;\vartheta)\\
&=-2\Delta q
\{2\log(q/2)+3+F(2,\vartheta)-2F_\alpha(2,\vartheta)\}
+O(q\log q).
\end{split}\tag{OFV28}
\]
Combining OFV27 and OFV28 gives the complete constant
\[
\mathfrak C_{\partial}
=4\log(4/\vartheta)-10-2F(2,\vartheta)+4F_\alpha(2,\vartheta),
\qquad \vartheta=\pi/2,
\tag{OFV29}
\]
with every logarithmic scale term canceled by direct algebra.
OSP33 now proves the actual original-root asymptotic
\[
\boxed{\displaystyle
\mathcal F_{\partial,k}^{(s)}
=\Delta q\,\mathfrak C_{\partial}
   +O_{\delta,\gamma}(q\log q),\qquad
\frac{\mathcal F_{\partial,k}^{(s)}}{kq}
\longrightarrow16\,\mathfrak C_{\partial},
\qquad s=1,k.}
\tag{OFV30}
\]
All constants in the error depend only on the fixed original
parameters, the fixed cutoff, and the proved compact equilibrium
constants. The two original source orders have been treated
through their exact recurrence maps and full mass factors.
Since $\Delta/k=16-48/k$ and $\log q/k\to0$, the displayed
limit follows from its quantified remainder.

An equivalent expression uses the logarithmic energy of precisely
the original EIQ limiting potential, $V(x)=\pi\sqrt{x}-2\log x$:
\[
I_{2,\pi}=\iint\log|x-y|\,d\mu_{2,\pi}(x)d\mu_{2,\pi}(y),
\qquad
\mathfrak C_{\partial}=4\log(4/\pi)+2-2I_{2,\pi}.
\tag{OFV31}
\]
To prove it, EIQ31 gives
$\vartheta\int\sqrt{x}\,d\mu_{2,\vartheta}=6$, so
$F(2,\vartheta)=I_{2,\vartheta}-6+2F_\alpha(2,\vartheta)$.
Substituting in OFV29 yields
$4\log(4/\vartheta)+2-2I_{2,\vartheta}$.
The proved EIQ32 dilation $x\mapsto4x$ carries $\mu_{2,\pi}$
to $\mu_{2,\pi/2}$. Its exact logarithmic energy change is
$I_{2,\pi/2}=I_{2,\pi}+\log4$, proving OFV31.
This energy is a finite explicit integral of the fully specified
EIQ density; its logarithmic singularity is integrable because
that density is bounded and compactly supported.

Finally, the exact quotient metrics decrease when the source degree
increases: each smaller-degree lift remains an allowed lift in the
larger source with exactly the same norm. Hence each low-degree
quotient Gram in CTV13 dominates the corresponding high-degree
Gram on the identical target frame. Positive conjugation shows
their determinant ratios are at least one. It follows that
$\mathcal F_{\partial,k}^{(s)}\ge0$, and OFV30 therefore also proves
$\mathfrak C_{\partial}\ge0$.
The already proved finite tail bounds CTV18--20, with
$0\le v\le80$, further give
\[
\mathcal T_k^{(s)}
=\Delta q\,\mathfrak C_{\partial}
+O_{\delta,\gamma}(q\log q),\qquad s=1,k,
\tag{OFV32}
\]
uniformly over the possible actual first-moment indices.
No leading value of the separate original observation quotient
is assigned here.

% END CURRENT OFV

\clearpage
% BEGIN CURRENT FMC
% FMC: a fixed original module presentation and controlled coefficient frames.
\section{Fixed coefficient control of the original residual module}
\label{sec:FMC}

Retain the complete actual simple-quartet period matrix
$\mathsf H=F^{-1}\Pi U\mathcal E_{\rm prim}$, the original substitution
$t=(z_0w_0,z_0w_1,z_1w_0,z_1w_1)^{\mathsf T}$, and
$Q(x)=Q_0(\mathsf H^{-1}x)$, where
$Q_0(t)=t_{++}t_{--}-t_{+-}t_{-+}$.
The original action is $g f(x)=f(D^{-1}x)$ with
$D=\operatorname{diag}(\zeta,\zeta^2,\zeta^3,\zeta^4)$ and
$\zeta=e^{2\pi i/5}$. At the fixed actual period on the proved
five-member orbit locus retain the entire polynomials
$Q_a=g^aQ$ and $\mathcal A=\prod_{a=1}^4Q_a$.
Their coefficients and all phases are unchanged.
The degree-$k$ calculations below use $k=4l+1\ge9$ and
\[
 q=(k+1)^2,\quad q'=(k-7)^2,\quad
 \Delta=16k-48,\quad j=8k-32,\quad r=8k-16.
 \tag{FMC1}
\]
All constants introduced below are determined by this one fixed
actual coefficient datum. They are independent of $k$ and of the
original Gamma cutoff. A fixed presentation will supply the residual
frame in place of dividing new selected minors at each degree.

\subsection{A fixed free module containing the entire invariant image}

Let $P=\mathbb C[x_1,x_2,x_3,x_4]$ and introduce the specified
graded ring injection
\[
 T=\mathbb C[y_1,y_2,y_3,y_4]\longrightarrow P,
 \qquad y_i\longmapsto x_i^5,\qquad \deg y_i=5.
 \tag{FMC2}
\]
It is injective because distinct monomials have distinct exponent
vectors after multiplying each exponent by five. Put
$\mathcal E=\{a\in\mathbb Z^4:0\le a_i\le4\}$ and let $e_a$
be a free $T$-basis element of degree $|a|$. The exact isomorphism is
\[
 \mathscr P:=\bigoplus_{a\in\mathcal E}T e_a
       \xrightarrow{\ \iota\ }P,
 \qquad y^\nu e_a\longmapsto x^{5\nu+a}.
 \tag{FMC3}
\]
Every nonnegative exponent vector has a unique coordinatewise quotient
and remainder on division by five. This proves surjectivity, injectivity,
preservation of degree, and the explicit inverse. In particular this
is a rank-$625$ free module, with coefficient $\ell^1$ norm preserved
exactly by $\iota$. The ring and monomial coefficients in $P$ have
not been rescaled.

Let $\mathcal E_0=\{a\in\mathcal E:
a_1+2a_2+3a_3+4a_4\equiv0\pmod5\}$. The full original invariant
algebra, regarded here as a $T$-module, is
\[
 \mathscr I=P^{\langle g\rangle}
       =\iota\left(\bigoplus_{a\in\mathcal E_0}T e_a\right),
 \qquad |\mathcal E_0|=125.
 \tag{FMC4}
\]
Indeed the monomial $x^{5\nu+a}$ has character
$\zeta^{-(a_1+2a_2+3a_3+4a_4)}$. The diagonal action fixes a
polynomial precisely when each nonzero monomial has character one.
For each of the $5^3$ choices of the first three residues there is
exactly one $a_4$ since $4$ is invertible modulo five. This proves
the count and the displayed equality. These pure fifth powers are
four of the fourteen original generators OCF25; this is a new exact
module presentation of the same entire invariant algebra.

Define two homogeneous $T$-submodules of $\mathscr P$ by
\[
 \begin{split}
 \mathscr H&=\iota^{-1}(QP+\mathcal A P),\\
 \mathscr S&=\mathscr H+
                       \bigoplus_{a\in\mathcal E_0}T e_a.
 \end{split}\tag{FMC5}
\]
They have the fixed finite generating columns
$\iota^{-1}(Qx^a)$ and $\iota^{-1}(\mathcal A x^a)$ for
$a\in\mathcal E$, followed for $\mathscr S$ by the $125$ columns
$e_a$, $a\in\mathcal E_0$. The first two lists have respectively
homogeneous degrees $|a|+2$ and $|a|+8$, all at most $24$.
To obtain every matrix entry explicitly, expand either of the full
original polynomials as $f(x)=\sum_b f_bx^b$ and use
\[
 \iota^{-1}(f x^a)=
       \sum_b f_b\,y^{\lfloor(a+b)/5\rfloor}
                         e_{(a+b)\bmod5}.
 \tag{FMC6}
\]
Both quotient and remainder are coordinatewise. Formula FMC3 verifies
this identity term by term. Since $P$ is free on the $x^a$, these
columns generate precisely the two modules in FMC5; no later degree
is needed to define a new relation matrix.

The full coefficient substitution and conductor quotient give the exact
graded maps
\[
 \mathscr P/\mathscr H\xrightarrow{\ \iota,\,\psi\ }
          \mathcal B/(A\mathcal B)=:\mathcal L,
 \qquad
 \mathscr S/\mathscr H\xrightarrow{\ \sim\ }\overline{\mathscr J},
 \quad A=\psi_8(\mathcal A).
 \tag{FMC7}
\]
Here $\psi f=f(\mathsf Ht)$, and $\overline{\mathscr J}$ is the
image of the original invariant algebra in $\mathcal L$, exactly
OCF16. OCF3 proves $\ker\psi=QP$ and its surjectivity in every
degree. If $\psi f=A b$, lift $b$ to a polynomial $h$ using that
surjectivity; then $f-\mathcal A h\in QP$. Conversely every
$QP+\mathcal AP$ element maps to zero. This proves the first
isomorphism. FMC4--5 show that its restriction has precisely the
second image and kernel. Thus the quotient in FMC7 is the complete
original conductor quotient, and its indicated submodule is the
actual residual certificate image of ROQ4, with no new invariant
premise.

\subsection{Finite reduction data with only fixed actual denominators}

The next construction specifies fixed data for both modules in FMC5.
Order their module monomials $y^\nu e_a$ first lexicographically by
$\nu_1,\nu_2,\nu_3,\nu_4$, and then by a fixed index
$\kappa_a\in\{0,\ldots,624\}$ listing $a$ lexicographically.
Multiplication by a $y$-monomial preserves this order. It is a
well-order: in a decreasing infinite sequence the first coordinate
can decrease only finitely often; after it is fixed the same applies
successively to the other three coordinates and then the finite index.
For a nonzero column $f$, retain its full leading coefficient
$\lambda_f\ne0$ and its leading module monomial $m_f$.
No leading coefficient is deleted from the stored column.

Start with the finite nonzero homogeneous generators of FMC5--6.
Division of a column repeatedly cancels its largest divisible module
monomial $c y^\nu e_a$ using the first available column $f$ for which
$m_f=y^\mu e_a$ divides it: subtract
$(c/\lambda_f)y^{\nu-\mu}f$ and record that exact multiple.
A largest term divisible by none of the current leading monomials is
put into the remainder and removed from the current dividend.
Each operation strictly decreases the largest term still being
processed, so this division terminates. Remainders have no monomial
divisible by the stored leading monomials.

For every pair $f,h$ whose leading positions agree, put
$m=\operatorname{lcm}(m_f,m_h)$ in that position and form the full
column
\[
 \operatorname{Sp}(f,h)=
      \frac{m/m_f}{\lambda_f}f
                  -\frac{m/m_h}{\lambda_h}h.
 \tag{FMC8}
\]
Divide it as just described and append its nonzero remainder to the
list. Reconsider all new pairs. Track every appended column as a
$T$-linear combination of the original generators throughout these
operations. Every column and recorded certificate remains homogeneous:
the two terms in FMC8 have the same weighted degree, because their
leading position agrees and every input column is homogeneous in the
shifted grading of FMC3.

This algorithm stops after finitely many appended columns. Here is the
needed finiteness argument. Every infinite sequence in $\mathbb N^4$
has two terms in their original order with all coordinates weakly
increasing. To prove it, an infinite sequence of nonnegative integers
has either an infinitely repeated value or an increasing infinite
subsequence; use this successively for all four coordinates. For
module monomials first take a subsequence with one common position.
Consequently every infinite sequence of module monomials has an earlier
one dividing a later one. Each appended remainder's leading monomial
is divisible by none of the previous leading monomials. An infinite
sequence of appendages would contradict the proved property.

At termination every element of the generated module has leading
monomial divisible by one of the stored leading monomials. For
completeness, write an element as a finite sum of monomial multiples
of the stored columns. Among all such representations choose one for
which the largest leading module monomial of a summand is smallest,
and among these choose one with the fewest summands at that monomial.
These choices exist by the well-order and the positive integers.
If cancellation of those largest terms occurs, two summands at the
same module position can be combined using a multiple of FMC8,
with its exact two leading coefficients. Explicitly, if their
common leading monomial is $M$ and their leading coefficients are
$a,b$, their sum is
\[
 a\frac{M}{\operatorname{lcm}(m_f,m_h)}\operatorname{Sp}(f,h)
       +(a+b)\frac{M/m_h}{\lambda_h}h.
\]
This follows by substituting FMC8; it retains exactly both original
summands and every coefficient. The recorded division of
this S-column, together with its appended remainder when one was
needed, expresses it in stored columns with all
leading multiples strictly below its canceled common monomial.
Replacing the two summands by this expression either decreases the
largest monomial or decreases the number attaining it; combining
their coefficients if necessary gives the latter alternative.
Both contradict the chosen representation. Thus the largest
remaining monomial is an actual leading monomial of the element,
and it is divisible by one of the stored ones, as claimed.
This also proves that division of a module element has zero remainder.

Write $\mathcal G_H$ and $\mathcal G_S$ for the two resulting fixed
finite lists for FMC5. The preceding argument proves their full
reduction property. A column supported outside the respective leading
module is called a standard remainder. Such a remainder represents
each quotient class uniquely: two representatives of one class differ
by a module element, and a nonzero standard difference would violate
the proved leading-monomial property. Denote these exact linear
remainder maps by $\operatorname{NF}_H$ and $\operatorname{NF}_S$.
Their linearity follows from uniqueness and linearity of quotienting.

All data constructed here depend only on the actual matrices in FMC6.
Every operation is an addition, multiplication, or division by one of
the nonzero actual leading coefficients encountered in this finite
calculation. Thus the complete list of denominators is finite and
independent of $k$. Its precise coefficient chart records each tested
coefficient's vanishing or nonvanishing and every selected leading
monomial. For every actual period exactly its corresponding decisions
give the above terminating calculation. This does not posit a
numerical test for equality of unevaluated periods, nor does it assert
a lower bound for a later sequence of minors. The exact stored
coefficients and their nonzero leading entries define the constants
that follow.

\subsection{A fixed weight bounds every homogeneous reduction path}

The finite reductions admit a linear-in-degree depth bound. Let
$D_0\ge1$ be the largest of one and all absolute differences between
corresponding $y$-exponents in a leading and a nonleading term of
any column in $\mathcal G_H\cup\mathcal G_S$. Empty tail lists
contribute zero. Set
\[
 B_0=D_0+626,\qquad
 \mathfrak w(y^\nu e_a)
   =B_0^4\nu_1+B_0^3\nu_2+B_0^2\nu_3+B_0\nu_4+\kappa_a,
 \qquad L_k=B_0^4\lfloor k/5\rfloor+624.
 \tag{FMC9}
\]
Every nonleading term of either fixed list has weight at least one
less than its leading term. If the $y$-exponents agree, this follows
from the position order. Otherwise let their first differing
coordinate have weight $B_0^m$, $1\le m\le4$. Its exponent drops by
at least one, while the later differences have absolute value at
most $D_0$, and the position difference has absolute value at most
$624$. The total weight drop is at least
\[
 B_0^m-D_0\sum_{h=1}^{m-1}B_0^h-624
 \;\ge\;\frac{625B_0^m}{D_0+625}-624\;>\;1.
 \tag{FMC10}
\]
The middle inequality follows by bounding the geometric sum by
$B_0^m/(B_0-1)$; the last follows from $B_0^m\ge B_0=D_0+626$.
Multiplication by a monomial adds the same weight to the leading
and all its replacement terms, so the same drop holds in every
division step. A homogeneous degree-$k$ monomial has
$5|\nu|+|a|=k$, hence weight between zero and $L_k$. Every branch
of its reduction tree consequently has length at most $L_k$.

For each of the two lists define the complete coefficient constant
\[
 C_H=\max\left\{1,\max_{g\in\mathcal G_H}
                 \sum_{m\ne m_g}|[m]g/\lambda_g|\right\},
 \qquad
 C_S=\max\left\{1,\max_{g\in\mathcal G_S}
                 \sum_{m\ne m_g}|[m]g/\lambda_g|\right\}.
 \tag{FMC11}
\]
The maximum over an empty list is zero. They are finite actual
numbers because the lists are finite and their displayed leading
coefficients are nonzero. For every homogeneous degree-$k$ column,
\[
 \|\operatorname{NF}_H f\|_1\le C_H^{L_k}\|f\|_1,
 \qquad
 \|\operatorname{NF}_S f\|_1\le C_S^{L_k}\|f\|_1.
 \tag{FMC12}
\]
To prove these inequalities start with one monomial of coefficient
one. A reduction replaces its leading term by the complete tail with
total absolute coefficient at most the indicated constant. Along
each branch the nonnegative integer weight drops at least one.
Induction on that weight bounds the total absolute coefficient of
all final leaves by $C^{\mathfrak w(m)}$: a standard leaf contributes
one, and an internal node contributes at most $C$ times the bound at
weight one lower. Cancellation between branches can only decrease the
coefficient $\ell^1$ norm. Linearity, the triangle inequality, and
$\mathfrak w(m)\le L_k$ prove FMC12 for every input.

There is equally explicit control of the recorded original-generator
certificates. For each stored $g\in\mathcal G_S$ let
$c_g$ be its complete certificate in the fixed original generating
list of $\mathscr S$ in FMC5--6. Define
\[
 D_S=\max\left\{1,\max_{g\in\mathcal G_S}
                             \|c_g\|_1/|\lambda_g|\right\}.
 \tag{FMC13}
\]
At a reduction node with input coefficient $c$, the certificate
contribution has norm at most $D_S|c|$. At depth $h$ the total
absolute coefficient entering all nodes is at most $C_S^h$ times
the input norm, by the same tree proof. Summing through depth $L_k$
therefore proves that a certificate of
$f-\operatorname{NF}_S f$ in the original generators has norm at most
$D_S(L_k+1)C_S^{L_k}\|f\|_1$. All certificate entries retain the
original homogeneous grading. Their denominators are among the same
finite fixed reduction data; no degree-$k$ linear system is inverted.

\subsection{The two exact maps back to the original conductor strip}

Let $\mathsf N_k$ and $\mathsf S_k$ be the original OCF strip
remainder and its coordinate insertion, with leading conductor
coefficient $a_*=a_{r_*s_*}$. Define
\[
 A_0=\max\left\{1,
                \sum_{(a,b)\ne(r_*,s_*)}|a_{ab}/a_*|\right\}.
 \tag{FMC14}
\]
Their exact coefficient bounds are
$\|\mathsf N_k\|_1\le A_0^{10k}$ and
$\|\mathsf S_k\|_1=1$. To prove the first, give the biform monomial
$e_{ab}^{(k)}$ the weight $9a+b\in[0,10k]$. Every other nonzero
term of the full $A$ has smaller weight than its lexicographically
largest term: a decrease in the first index outweighs at most eight
increases in the second index, and a decrease in only the second
index also drops the weight. Every translated OCF6 subtraction
therefore has weight drop at least one. Its complete tail has
coefficient norm at most $A_0$. The reduction-tree proof of FMC12
now applies with depth $10k$. The insertion bound is immediate
from its unchanged coordinate coefficients.

Put
\[
 h_+=\max\{1,\max_i\sum_a|\mathsf H_{ia}|\},\qquad
 h_-=\max\{1,\max_a\sum_i|(\mathsf H^{-1})_{ai}|\},
 \quad h_* =\max\{h_+,h_-\}.
 \tag{FMC15}
\]
Substitution $x=\mathsf Ht$ on degree-$k$ polynomial coefficients
has $\ell^1$ norm at most $h_+^k$. Expanding each product of linear
forms and using the multinomial theorem proves this for each
monomial; the triangle inequality proves it for their sum. The Segre
substitution adds coefficients on identical biform monomials and so
has norm at most one. Conversely choose the following exact lift
of $e_{ab}^{(k)}$: with $h=\min(a,b)$ take
\[
 t_{++}^{h}t_{+-}^{a-h}t_{-+}^{b-h}t_{--}^{k-a-b+h},
 \qquad t=\mathsf H^{-1}x.
 \tag{FMC16}
\]
All exponents are nonnegative since $0\le a,b\le k$, and their
sum is $k$. Direct substitution gives the required biform monomial.
This lift, extended linearly and denoted $\mathscr L_k$, has norm
at most $h_-^k$ by the same complete multinomial estimate.

Let $\mathscr M_{H,k}$ be the set of degree-$k$ standard module
monomials for $\mathcal G_H$. Its coefficient space is the exact
degree-$k$ quotient model of FMC7. Define
\[
 \begin{split}
 U_k f&=\mathsf N_k\,\psi\,\iota f,\qquad
          f\in\mathbb C^{\mathscr M_{H,k}},\\
 V_k w&=\operatorname{NF}_H\,\iota^{-1}
                        \mathscr L_k\mathsf S_k w,
          \qquad w\in W_k.
 \end{split}\tag{FMC17}
\]
These are inverse maps. For the first composition, $\operatorname{NF}_H$
changes a lift only by $\mathscr H$, which maps to zero in the
original strip by FMC7, and $\psi\mathscr L_k=I$ together with
$\mathsf N_k\mathsf S_k=I$ proves $U_kV_k=I$.
For the opposite composition, $\mathscr L_k\mathsf S_k
\mathsf N_k\psi\iota f-\iota f$ maps to zero in the strip, so
it belongs to $QP+\mathcal AP$ by FMC7. Its free-module difference
belongs to $\mathscr H$ and is killed by the unique remainder.
The original $f$ is already standard. Hence $V_kU_k=I$.
The bounds already proved give
\[
 \|U_k\|_1\le h_+^k A_0^{10k},\qquad
 \|V_k\|_1\le h_-^k C_H^{L_k}.
 \tag{FMC18}
\]
These bounds use only the displayed original coefficients and the
two fixed homogeneous relation lists.

\subsection{An identity pivot block for the entire residual image}

Since $\mathscr H\subset\mathscr S$, every leading monomial of
$\mathscr H$ is a leading monomial of an element of $\mathscr S$.
The proved reduction property for $\mathcal G_S$ then implies
inclusion of their leading modules. Consequently
\[
 \mathscr M_{S,k}\subset\mathscr M_{H,k},\qquad
 \mathscr U_k=\mathscr M_{H,k}\setminus\mathscr M_{S,k},\qquad
 |\mathscr M_{H,k}|=\Delta,\quad
 |\mathscr M_{S,k}|=r,\quad |\mathscr U_k|=j.
 \tag{FMC19}
\]
The first count is FMC7 and OCF8. The quotient by $\mathscr S$
is exactly the original conductor quotient modulo its invariant
image, which has dimension $r$ by OCF22--24 and OKM3--5.
Subtracting proves the last count. These are counts of actual
standard monomials for the fixed relation lists, not independent
generic ranks.

Order $\mathscr M_{H,k}$ by listing $\mathscr U_k$ first and
$\mathscr M_{S,k}$ second, with the specified monomial order within
each list. For every $u\in\mathscr U_k$ put
$z_u=u-\operatorname{NF}_S(u)$. Both terms are $H$-standard,
and their difference belongs to $\mathscr S$. The $j$ resulting
columns form a basis of $\mathscr S_k/\mathscr H_k$ in this
quotient model. Indeed their coordinates on $\mathscr U_k$ are
the distinct unit vectors. For a standard column in $\mathscr S$,
the remainder under $\operatorname{NF}_S$ is zero. Its coordinates
on $\mathscr M_{S,k}$ must therefore be the negatives of the
remainders of its $\mathscr U_k$ terms. It is exactly the stated
linear combination of the $z_u$. This proves spanning and
independence, including all tails of every column.

Write $B_k$ for the matrix with columns $\operatorname{NF}_S(u)$
in the second standard set. The full exact graph frame and inverse are
\[
 P_k=\begin{pmatrix}I_j&0\\-B_k&I_r\end{pmatrix},\qquad
 P_k^{-1}=\begin{pmatrix}I_j&0\\B_k&I_r\end{pmatrix},\qquad
 \det P_k=1,\quad
 \|P_k\|_1,\|P_k^{-1}\|_1\le1+C_S^{L_k}.
 \tag{FMC20}
\]
The two matrix identities follow by multiplication and block
triangularity, and the norm bound follows from FMC12 applied to
each entire column. All entries of $B_k$ come from fixed reductions;
no minor is selected or divided at degree $k$.

Each $z_u$ has an actual invariant certificate. The recorded reduction
certificate of $u-\operatorname{NF}_S(u)$ from FMC13 expresses its
polynomial representative as
\[
 \iota z_u=Qh_u+\mathcal A a_u+f_u,
 \qquad f_u\in\mathscr I_k,
 \qquad \|f_u\|_{1,x}
                  \le D_S(L_k+1)C_S^{L_k}.
 \tag{FMC21}
\]
To obtain the three displayed pieces, collect the certificate terms
corresponding respectively to $Qx^a$, $\mathcal Ax^a$, and the
$125$ invariant basis elements $x^a$. Every coefficient is a
polynomial in the $y_i=x_i^5$. The last terms therefore belong
to $\mathscr I$, and homogeneity gives exactly degree $k$.
FMC3 preserves coefficient norm, and discarding the first two
certificate lists when bounding only the last cannot increase
their total $\ell^1$ norm. FMC13 proves the bound. Thus $U_kz_u$
is exactly $\mathsf N_k\psi f_u$, an original invariant restriction,
with the entire original periods in its coefficients.

Define the full original strip frame, its residual columns and their
explicit left inverse by
\[
 T_k=U_kP_k,\qquad
 Z_k^{\rm fix}=T_kE_j,\qquad
 L_{Z,k}=E_j^{\mathsf T}P_k^{-1}V_k,\qquad
 L_{Z,k}Z_k^{\rm fix}=I_j.
 \tag{FMC22}
\]
Here $E_j$ inserts the first $j$ coordinates into the new ordered
standard basis, while all output strip coordinates retain their
original OCF order. FMC17 and FMC20 prove the inverse formula
$T_k^{-1}=P_k^{-1}V_k$ and the last identity. By FMC21 these
$Z_k^{\rm fix}$ columns span precisely the original $\overline J_k$.
They provide all actual residual coefficients with a left inverse
obtained from fixed relations.

For the original coefficient Euclidean norms set the explicit number
\[
 M_k=2\sqrt{\Delta}\,h_*^k A_0^{10k}(C_HC_S)^{L_k}\ge1.
 \tag{FMC23}
\]
Then
\[
 \|T_k\|_2,\|T_k^{-1}\|_2,
       \|Z_k^{\rm fix}\|_2,\|L_{Z,k}\|_2\le M_k,
 \quad
 M_k^{-1}\|v\|_2\le\|Z_k^{\rm fix}v\|_2\le M_k\|v\|_2.
 \tag{FMC24}
\]
For the first two bounds, FMC18--20 give respectively the column-sum
bounds $h_+^kA_0^{10k}(1+C_S^{L_k})$ and
$(1+C_S^{L_k})h_-^kC_H^{L_k}$. For a matrix with $\Delta$
input coordinates, its Euclidean norm is at most $\sqrt\Delta$
times its coefficient column-sum norm: use
$\|Av\|_2\le\|Av\|_1\le\|A\|_1\|v\|_1
\le\sqrt\Delta\|A\|_1\|v\|_2$.
All constants are at least one, proving the two stated bounds.
The coordinate insertion and restriction have Euclidean norm one,
so FMC22 proves the remaining operator bounds. Its left-inverse
identity then proves the lower bound on $Z_k^{\rm fix}$.

For every $0\le t\le\Delta$, singular values now give
\[
 \log^+\|\bigwedge^t T_k\|_2,
 \log^+\|\bigwedge^t T_k^{-1}\|_2
       \le t\log M_k\le\Delta\log M_k
       =O_{\rm actual}(k^2)=o(kq).
 \tag{FMC25}
\]
Indeed every singular value of both square matrices is at most
$M_k$, and exterior singular values are their products. This last
fact follows by applying a singular-value decomposition: unitary
factors preserve exterior norms, and the diagonal factor multiplies
each coordinate wedge by the corresponding product. Explicitly
\[
 \log M_k=\log2+\tfrac12\log\Delta+k\log h_*
              +10k\log A_0+L_k\log(C_HC_S).
 \tag{FMC26}
\]
All data except $k,\Delta,L_k$ are fixed actual numbers,
$L_k\le B_0^4 k/5+624$, $\Delta=16k-48$, and
$kq=k(k+1)^2$. These formulas prove the last two statements of
FMC25 directly. The same bounds hold for transpose inverse maps,
whose singular values are unchanged. For $0\le t\le j$, FMC24
also puts every singular value of $Z_k^{\rm fix}$ between
$M_k^{-1}$ and $M_k$, so the same exterior bounds apply to its
map onto its image and its inverse there.

\subsection{Exact relation to every earlier certificate frame}

Let $Z_k^{\rm old}$ be the full actual OCF/ROQ certificate frame,
with its selected invertible $j\times j$ row block
$(Z_k^{\rm old})_I$. Both frames have exactly the same column
space, as proved in FMC21--22. Their exact transition and inverse are
\[
 A_k=((Z_k^{\rm old})_I)^{-1}(Z_k^{\rm fix})_I,
 \qquad A_k^{-1}=L_{Z,k}Z_k^{\rm old},\qquad
 Z_k^{\rm fix}=Z_k^{\rm old}A_k.
 \tag{FMC27}
\]
Each new column has a unique expansion in the old independent columns;
restriction to the old pivot rows gives the first formula. Applying
$L_{Z,k}$ gives the second and both inverse identities. Thus this
is the exact coefficient morphism to the earlier frame, including
its actual pivot determinant. No bound for that old determinant
is needed or asserted by FMC24--26. The new frame and its proved
bounds use only the fixed data of FMC5--23.

The two residual maps on $V_A$ therefore obey
$\Xi_k^{\rm fix}=A_k^{\mathsf T}\Xi_k^{\rm old}$ in their
respective coordinates, since the maps in the original strip dual
are the ordinary transposes of their certificate matrices.
For every original source and cutoff their minimum quotient metrics
satisfy the exact congruence and determinant formulas
\[
 Q_{\Xi,N}^{\rm fix}
        =\overline{A_k}^{-1}Q_{\Xi,N}^{\rm old}A_k^{-\mathsf T},
 \qquad
 \log\det Q_{\Xi,N}^{\rm fix}
        =\log\det Q_{\Xi,N}^{\rm old}-2\log|\det A_k|.
 \tag{FMC28}
\]
The constraints on a lift are transformed by precisely
$A_k^{\mathsf T}$; substituting the inverse coordinates in the
same attained minimum proves the congruence, and determinants give
the second formula. The scalar is independent of $N$ and of source
order. Its coefficient is exactly $1+1-1-1=0$ in the original
four-endpoint return. Thus this return agrees identically with the
entire earlier ROQ/RCA return, even if an earlier selected pivot is
small at some degrees.

The complete inherited period Gram is transported by the same exact
map at each endpoint. Let $L_{\rm old}$ be the primary right
section ROQ6 and let $L_{\rm fix}$ be any primary right section
of $\Xi_k^{\rm fix}$ inside $V_A$, for example the first $j$
columns of FMC30 below. Both $L_{\rm fix}$ and
$L_{\rm old}A_k^{-\mathsf T}$ have residual image $I_j$ under
$\Xi_k^{\rm fix}$, so their difference has columns in the
original $K=\ker\Lambda$. Consequently
\[
 \Lambda L_{\rm fix}=\Lambda L_{\rm old}A_k^{-\mathsf T},\qquad
 P_{\rm per}^{\rm fix}
     =\overline{A_k}^{-1}P_{\rm per}^{\rm old}A_k^{-\mathsf T},
 \qquad
 \frac{\det Q_{\Xi,N}^{\rm fix}}{\det P_{\rm per}^{\rm fix}}
     =\frac{\det Q_{\Xi,N}^{\rm old}}{\det P_{\rm per}^{\rm old}}.
\]
The first equality preserves the literal boundary vectors and their
full cyclic coefficient map. Taking their Gram with the unchanged
inherited metric ROQ9 proves the second equality. Taking both
determinants together with FMC28 proves the last equality pointwise,
before taking any endpoint difference.

\subsection{The unchanged Gamma source in the fixed degree model}

Retain $c=k/2$, the original nodes
$\beta_{ab}=c+(2a-k)\delta+i(2b-k)\gamma$, and
$\omega_{ab}=(\beta_{ab}-c)/i$. For each original source
$s\in\{1,k\}$ retain its complete recurrence and masses
\[
 \begin{gathered}
 p_0=1,\ p_1=y,\quad
 p_{n+1}=yp_n-n(s/2+n-1)p_{n-1},\\
 h_n=(2\pi)^{s/2}n!(s/2)_n,\qquad
 (U_N)_{(a,b),n}=p_n(\omega_{ab})/\sqrt{h_n},\qquad
 G_N=(U_NU_N^*)^{-1}.
 \end{gathered}\tag{FMC29}
\]
Here $N$ remains one of $q-1,q,2q-1,2q$. These are exactly the
original primary quotient metrics of CEP21--22 and ROQ10, with
every source mass and coordinate phase present.

The complete primary frame of $V_A$, in residual coordinates first
and common-kernel coordinates second, is
\[
 J_k=\mathsf N_k^{\mathsf T}T_k^{-\mathsf T},\qquad
 \Xi_k^{\rm fix}J_k=E_j^{\mathsf T},\qquad
 K=\operatorname{im}(J_kE_r).
 \tag{FMC30}
\]
Here $E_r$ inserts the last $r$ coordinates. For a primary vector
$x\in V_A$, ROQ3 gives its strip dual coordinate
$z=\mathsf S_k^{\mathsf T}x$. The new residual map is
$(Z_k^{\rm fix})^{\mathsf T}z=E_j^{\mathsf T}T_k^{\mathsf T}z$.
Substitution of $z=T_k^{-\mathsf T}y$ proves the second identity,
and its kernel is the last coordinate block, proving the third.
Thus this full frame has precisely the original common kernel;
it changes no observation. Its left inverse on $V_A$ is
$T_k^{\mathsf T}\mathsf S_k^{\mathsf T}$.
The original coefficient bounds additionally give
\[
 \|J_k\|_2\le\sqrt q A_0^{10k}M_k=:\widehat M_k,
 \qquad
 \|T_k^{\mathsf T}\mathsf S_k^{\mathsf T}\|_2\le M_k.
 \tag{FMC31}
\]
Indeed FMC14 bounds the norm of $\mathsf N_k$ by its column-sum
norm times $\sqrt q$, transpose preserves Euclidean norm, and
FMC24 applies to $T_k^{-\mathsf T}$. The insertion transpose is
a coordinate restriction of norm one. Both logarithms in FMC31
are $O_{\rm actual}(k+\log k)$ with their full displayed constants.

The exact source Gram in this specified degree model, and its actual
residual quotient, are
\[
 \begin{gathered}
 \mathcal H_N=J_k^*G_NJ_k
       =\overline{T_k}^{-1}
          \overline{\mathsf N_k}G_N\mathsf N_k^{\mathsf T}
                                               T_k^{-\mathsf T},\\
 \mathcal H_N=
       \begin{pmatrix}H_{11,N}&H_{12,N}\\H_{21,N}&H_{22,N}\end{pmatrix},
 \qquad
 Q_{\Xi,N}^{\rm fix}=H_{11,N}-H_{12,N}H_{22,N}^{-1}H_{21,N}.
 \end{gathered}\tag{FMC32}
\]
The first formula is direct substitution of the actual primary
frame. Since $G_N>0$ and $J_k$ is injective, $\mathcal H_N>0$;
therefore $H_{22,N}>0$. For prescribed first coordinate $y$, the
second coordinate minimizing the full quadratic form is
$-H_{22,N}^{-1}H_{21,N}y$. Completing the square, or multiplying
out this substitution, gives exactly the displayed Schur
complement. This is the original minimum on $\Xi_k^{\rm fix}x=y$
inside $V_A$, so it equals FMC28 and the full ROQ metric.
All four blocks, including every off-diagonal entry, are retained.
The original minimum source lift of the full frame is explicitly
$U_N^*G_NJ_k$, whose Gram is $\mathcal H_N$ by $U_NU_N^*=G_N^{-1}$.
It contains the entire original polynomial recurrence and remainder
at each stated cutoff.

There is a useful quantitative determinant comparison which does not
replace that Gamma form by a coefficient form. In strip dual coordinates
write $H_{V,N}=\overline{\mathsf N_k}G_N\mathsf N_k^{\mathsf T}$ and
$Z=Z_k^{\rm fix}$. Define the actual coefficient quotient form and
the relative determinant by
\[
 Q_{\rm coeff}=(Z^{\mathsf T}\overline Z)^{-1},\qquad
 \Phi_{\Xi,N}=\log\det Q_{\Xi,N}^{\rm fix}
                                  -\log\det Q_{\rm coeff}.
 \tag{FMC33}
\]
The first is the minimum quotient of the original coefficient
Euclidean norm in the fixed OCF strip, with the identical constraint
$Z^{\mathsf T}z=y$. The elementary minimum proof of ROQ11 gives
both this formula and
$Q_{\Xi,N}^{\rm fix}=(Z^{\mathsf T}H_{V,N}^{-1}\overline Z)^{-1}$.
Hence the full, directly calculable expression is
\[
 \begin{gathered}
 \Phi_{\Xi,N}=
       \log\det(Z^{\mathsf T}\overline Z)
       -\log\det(Z^{\mathsf T}H_{V,N}^{-1}\overline Z),\\
 |\log\det Q_{\rm coeff}|\le2j\log M_k
                       =O_{\rm actual}(k^2)=o(kq),\\
 F_{\Xi,k}^{(s)}
       =\Phi_{\Xi,q-1}+\Phi_{\Xi,q}
                              -\Phi_{\Xi,2q-1}-\Phi_{\Xi,2q}.
 \end{gathered}\tag{FMC34}
\]
FMC24 puts the singular values of $Z$ between $M_k^{-1}$ and
$M_k$. The positive eigenvalues of $Z^{\mathsf T}\overline Z$
are their squares, proving the determinant bound by summing $j$
logarithms. The first identity follows by the two exact inverse
Gram formulas. The final identity follows by canceling the fixed
$\log\det Q_{\rm coeff}$ with the four original signs and then
using FMC28. Thus all coefficient denominators and frame-volume
effects in this specified construction are controlled below the
$kq$ scale, while the original Gamma compression remains fully
visible in its exact determinant.

For the kernel frame $I_k=J_kE_r$, FMC31 likewise gives the concrete
coefficient-Gram bounds
\[
 M_k^{-2r}\le\det(I_k^*I_k)\le\widehat M_k^{2r},\qquad
 |\log\det(I_k^*I_k)|
             \le2r\log\widehat M_k=O_{\rm actual}(k^2)=o(kq).
 \tag{FMC35}
\]
Its left inverse on its image is
$E_r^{\mathsf T}T_k^{\mathsf T}\mathsf S_k^{\mathsf T}$, with norm
at most $M_k$; FMC31 bounds its forward norm. The singular-value
argument just given proves FMC35. Define the exact primary-coordinate
projection and determinant by
\[
 P_{K,\rm prim}=I_k(I_k^*I_k)^{-1}I_k^*,\qquad
 D_{K,N}^{\rm prim}=[z^r]\det(I_q+zP_{K,\rm prim}G_NP_{K,\rm prim}).
\]
Multiplication proves that this projection is Hermitian, idempotent,
and has exactly the original primary-coordinate $K$ as its image.
For the determinant proof take $U=I_k(I_k^*I_k)^{-1/2}$, so
$U^*U=I_r$, and extend its columns to an orthonormal basis of
$\mathbb C^q$. In that basis the projected metric has block
$U^*G_NU$ and a zero block. Its indicated determinant coefficient
is $\det(U^*G_NU)$, and substitution of the full Gram square root
therefore proves
\[
 \det(I_k^*G_NI_k)=\det(I_k^*I_k)D_{K,N}^{\rm prim}.
\]
The precise map to AKS12--14's original polynomial-coefficient
projector is the ordered Vandermonde: put
$I_{k,\rm poly}=V_\beta^{-1}I_k$ and
$G_{N,\rm poly}=V_\beta^*G_NV_\beta$.
Their compressed Gram is identically $I_k^*G_NI_k$ by
multiplication. Applying the determinant proof just given in
these polynomial coordinates proves the exact bridge
\[
 D_{K,N}^{\rm AKS}
    =\frac{\det(I_k^*I_k)}
            {\det(I_{k,\rm poly}^*I_{k,\rm poly})}
                         D_{K,N}^{\rm prim}.
\]
Here the polynomial projector has image the actual coefficient
kernel and thus agrees with AKS12's projector. The displayed
positive factor is independent of $N$ and source order and
cancels exactly with the four endpoint signs. FMC35 bounds its
primary-coordinate numerator; the polynomial denominator is
retained as the full displayed Vandermonde Gram. No bound for
that different denominator follows from FMC35 alone.

FMC25 is a coefficient-space estimate for specified exact maps.
Its relation to the original Gamma norms is the complete congruence
and Schur calculation FMC29--34, not an assertion that those maps
are Gamma isometries. In particular no upper or lower leading
value for the remaining Gamma compression is inferred from their
coefficient condition numbers. The new result supplies a fixed
degree presentation, actual invariant certificates, bounded forward
and inverse exterior maps, and the exact original determinant on
that model, with every fixed denominator specified and every
source order and endpoint retained.

\subsection{The exact map to the original polynomial degree flag}

The homogeneous degree $k$ in FMC3 is the original tensor degree.
Its precise map to the original polynomial-in-$S$ model is the
ordered evaluation isomorphism, not an identification of the two
degree indices. Retain
$\chi_k(S)=\prod_{a,b=0}^k(S-\beta_{ab})$ and the original
Vandermonde $(V_\beta)_{(a,b),d}=\beta_{ab}^d$, $0\le d<q$.
For a fixed-model primal coordinate column $y\in\mathbb C^\Delta$
the exact polynomial representative is
\[
 p_y(S)=\sum_{a,b=0}^k(J_ky)_{ab}
          \frac{\chi_k(S)}{(S-\beta_{ab})\chi_k'(\beta_{ab})},
 \qquad [p_y]_{\rm coeff}=V_\beta^{-1}J_ky.
 \tag{FMC36}
\]
All denominators are nonzero because the original nodes are
distinct. Evaluation at each node proves the displayed inverse
and hence the equality of the two formulas. It includes the
complete roots, their order, and their actual centre. Thus the
source degree-$N$ metric in FMC32 is equivalently
$(V_\beta^{-1}J_k)^*(V_\beta^*G_NV_\beta)
(V_\beta^{-1}J_k)$, identically $J_k^*G_NJ_k$ by multiplication.
This proves the full typed relation to the original coefficient
Gamma model.

There is an exact further map to the degree-adapted conductor
basis AKJ17. Retain all actual conductor moments $\mu_d$, their
first nonzero index $v\le80$, the literal $B_q(\partial_S)$ of
AKJ14, and $\mathcal J_v S^d=d!S^{d+v}/(d+v)!$. The ordered
polynomial columns of $P_{A,k}$ are
\[
 \begin{gathered}
 1,S,\ldots,S^{v-1},\\
 P_h=\frac{(q'+h+v)!}{(q'+h)!}
       B_q(\partial_S)^{-1}\mathcal J_v(\chi_{k'}S^h),
               \qquad 0\le h\le\Delta-v-1.
 \end{gathered}\tag{FMC37}
\]
Their distinct monic degrees are $0,\ldots,v-1$ and
$q'+v,\ldots,q-1$. AKJ14--17 proves directly that they form a
basis of this same $V_A$, using the complete conductor derivative
and its finite nilpotent inverse. Define $B_{A,k}(p)$ on $V_A$
by successively subtracting the coefficient of the highest
remaining degree times its matching column in FMC37, down
through the high list and then through the low monomials.
Each subtraction uses leading coefficient one. Spanning and the
distinct degree pivots show that it ends at zero and records the
unique coordinate column of $p$; this also proves its linearity
and both inverse identities on $V_A$.
Consequently the precise transition is
\[
 \begin{gathered}
 R_{A,k}=\mathsf S_k^{\mathsf T}V_\beta P_{A,k},\qquad
 R_{A,k}^{-1}z=B_{A,k}\bigl(V_\beta^{-1}\mathsf N_k^{\mathsf T}z\bigr),\\
 V_\beta^{-1}J_k=P_{A,k}R_{A,k}^{-1}T_k^{-\mathsf T},\qquad
 \Xi_{A,k}^{\rm fix}=E_j^{\mathsf T}T_k^{\mathsf T}R_{A,k}.
 \end{gathered}\tag{FMC38}
\]
The columns $V_\beta P_{A,k}$ lie in $V_A$. Therefore ROQ3 gives
$\mathsf N_k^{\mathsf T}R_{A,k}=V_\beta P_{A,k}$, which proves
the first inverse formula by applying the unique degree
coordinates just constructed. Both sides are invertible maps
between the displayed $\Delta$-dimensional coefficient spaces.
Substituting FMC30 gives the third equality, and evaluating
the new residual map on these same primary columns gives the
last. Thus its degree-prefix ranks give exactly
$\dim(K\cap\mathbb C[S]_{\le d})$ by subtracting that prefix
rank from the number of columns of FMC37 of degree at most $d$,
as in AKJ18. Every original factorial, finite inverse, primary
denominator and degree pivot is visible in these maps. The
subleading coefficient-frame estimate proved in FMC25 pertains
to the fixed strip module maps; FMC36--38 retain the separate
polynomial interpolation and its full Gamma metric exactly.

% END CURRENT FMC

\clearpage
% BEGIN CURRENT OVR
\section{The evaluated outer contribution in the original mixed return}
\label{sec:OVR}

Retain the actual simple quartet, original period on the proved five-orbit
locus, full amplitude unit, branch, and original source orders $s=1,k$ of
OAR and OFV. All arithmetic statements retain the original domain
$|u|\ge R_*$ of AMT; the actual period is fixed as $k=4l+1$ grows.
Here $q=(k+1)^2$, $\Delta=16k-48$, and every finite use of OFV retains
its stated threshold, including OSP23. Write
\[
 C_\partial=\mathfrak C_\partial
 =4\log(4/\pi)+2-2I_{2,\pi},\qquad
 I_{2,\pi}=\iint\log|x-y|\,d\mu_{2,\pi}(x)d\mu_{2,\pi}(y).
 \tag{OVR1}
\]
The measure here is the full EIQ density with its uniquely specified
endpoints. OFV proves this exact expression and $C_\partial\ge0$.
No value of the original residual observation is inserted into it.

\subsection{The exact remainder and the sourcewise leading identity}
Define the actual outer remainder and retain OAR3's full conductor and
monic-tail remainder by
\[
 b_{k,s}=\mathcal F_{\partial,k}^{(s)}-\Delta q C_\partial,
 \qquad
 R_{k,s}=b_{k,s}+\mathfrak e_{k,s}.
 \tag{OVR2}
\]
Thus $b_{k,s}=O_{\delta,\gamma}(q\log q)$ by OFV30. The second
summand is precisely the signed finite expression OAR3, with
$-e_{k,s}^-\le\mathfrak e_{k,s}\le e_{k,s}^+$ and
$e_{k,s}^\pm=o(kq)$. In particular $R_{k,s}=o(kq)$, because
$\log q/k\to0$. Its two exact finite error endpoints remain
$b_{k,s}-e_{k,s}^-$ and $b_{k,s}+e_{k,s}^+$; their signs are
not replaced by a symmetry assumption. Substitution into OAR3 proves
\[
 \begin{gathered}
 F_K^{(s)}+\Phi_{k,s}=\Delta q C_\partial+R_{k,s},\\
 \frac{F_K^{(s)}}{kq}+\frac{\Phi_{k,s}}{kq}
       =16C_\partial-\frac{48C_\partial}{k}
                         +\frac{R_{k,s}}{kq}
       \longrightarrow16C_\partial,\qquad s=1,k.
 \end{gathered}\tag{OVR3}
\]
Here $\Phi_{k,s}$ is exactly ROQ19's positive sum of constrained
covariance increments, built from the original period certificate
matrix $D$ and the complete conductor cross covariance. Its coefficient
maps, minimum sections, inherited period metric and graph denominator
are those already proved in ROQ; OVR3 changes none of them.

\subsection{Transfer between the two original Gamma sources}
Let $d_k=F_K^{(k)}-F_K^{(1)}$ denote the actual difference.
AMT44--46 proves $d_k=o(kq)$, with its original finite two-sided
source-comparison bounds. Subtracting the two exact lines OVR3 gives
\[
 \Phi_{k,k}-\Phi_{k,1}=R_{k,k}-R_{k,1}-d_k,
 \qquad
 \frac{\Phi_{k,k}-\Phi_{k,1}}{kq}\longrightarrow0.
 \tag{OVR4}
\]
This proves the transfer using the actual two residual observations,
whose individual metrics retain their respective original sources.
There is no asserted equality between those metrics. The equality in
OVR4 is the exact scalar consequence of their common original-kernel
restriction maps and the proved determinant identities.

\subsection{Arithmetic kernel and its correlated boundary}
Put $d_k^{\rm ar}=F_K^{\rm ar}-F_K^{(1)}$. This is the actual
quantity bounded in OAR8; AMT38 and AMT46 prove
$d_k^{\rm ar}=o(kq)$. Keeping its full value first gives
\[
 \begin{split}
 F_K^{\rm ar}
    &=\Delta q C_\partial-\Phi_{k,1}+R_{k,1}+d_k^{\rm ar},\\
 F_B^{\rm ar}
    &=\mathcal B_k^{\rm ar}-\Delta q C_\partial
          +\Phi_{k,1}-R_{k,1}-d_k^{\rm ar}.
 \end{split}\tag{OVR5}
\]
The first equality adds the literal definition of $d_k^{\rm ar}$
to OVR3 at $s=1$. The second subtracts the first from
$F_K^{\rm ar}+F_B^{\rm ar}=\mathcal B_k^{\rm ar}$.
In particular the same arithmetic total and the same two finite
errors occur with opposite signs in the boundary. OAR8 supplies
the unchanged correlated finite bounds on $d_k^{\rm ar}$.
Dividing these exact equalities gives
\[
 \frac{F_K^{\rm ar}+\Phi_{k,1}}{kq}\longrightarrow16C_\partial,
 \qquad
 \frac{\mathcal B_k^{\rm ar}-F_B^{\rm ar}+\Phi_{k,1}}{kq}
                  \longrightarrow16C_\partial.
 \tag{OVR6}
\]
Thus evaluation of the residual sum in either original Gamma order
passes through to the same original arithmetic kernel and boundary.

\subsection{The signed return retains its original Gamma correction}
The exact signed quantity of OAR11 is
$\mathfrak S_k^{\rm mix}=F_K^{\rm ar}+F_B^{\rm ar}-F_B^{(k)}$.
Applying OVR3 at $s=k$ to that identity yields
\[
 \mathfrak S_k^{\rm mix}+\Phi_{k,k}
     =\Delta q C_\partial+\mathcal H_k+\delta_k^\sigma+R_{k,k}.
 \tag{OVR7}
\]
All finite terms are the original ones. AMT48--49, with the retained
total correction estimates used there, proves
$(\mathcal H_k+\delta_k^\sigma)/(kq)\to-\mathcal C_\Gamma/4$,
where
\[
 \mathcal C_\Gamma
   =2\int\log x\,d\mu_{2,\pi}(x)-4(2\log2-1)>0.
\]
The equilibrium and source maps are exactly the EIQ/GEL providers
used in OFV and AMT. Hence OVR7 and OVR4 prove both limits
\[
 \frac{\mathfrak S_k^{\rm mix}+\Phi_{k,k}}{kq}
       \longrightarrow16C_\partial-\frac{\mathcal C_\Gamma}{4},
 \qquad
 \frac{\mathfrak S_k^{\rm mix}+\Phi_{k,1}}{kq}
       \longrightarrow16C_\partial-\frac{\mathcal C_\Gamma}{4}.
 \tag{OVR8}
\]
The two constants are explicit integrals of the same proved density,
respectively its full logarithmic pair energy and its logarithmic
moment. No identity between these distinct integrals is presumed.

\subsection{Bounds on the remaining residual coefficient}
Each $F_K^{(s)}$ is nonnegative: the original quotient metrics decrease
with increasing source degree, and restriction to the identical kernel
preserves the two Gram inequalities at the paired endpoints. Their
determinant ratios are consequently at least one, as proved in AMT15.
ROQ19 gives $\Phi_{k,s}\ge0$. Thus OVR3 bounds both nonnegative
sequences after division by $kq$. Let
$\kappa_*=32\log(25\log3/(4\pi))$ be the retained AKS asymptotic
kernel ceiling on this actual locus, applied in both original orders.
The actual residual cap OAR5, proved by ROQ21, is
$\Phi_{k,s}\le j(\mathcal L_{s,q}+\mathcal L_{s,q+1})-\lambda_s$,
where $j=8k-32$ and $\lambda_s\ge0$ is the original first increment.
AKS31 gives
$\limsup j(\mathcal L_{s,q}+\mathcal L_{s,q+1})/(kq)\le\kappa_*$
in both original source orders: its same explicit logarithmic
comparison constants apply, and $j/k\to8$ just as $r/k\to8$.
Thus both the original kernel and residual sequences have the same
proved upper ceiling. Combining their two ceilings with OVR3 gives
\[
 \begin{gathered}
 \max\{0,16C_\partial-\kappa_*\}
       \le\liminf\frac{F_K^{(s)}}{kq}
       \le\limsup\frac{F_K^{(s)}}{kq}
       \le\min\{\kappa_*,16C_\partial\},\\
 \max\{0,16C_\partial-\kappa_*\}
       \le\liminf\frac{\Phi_{k,s}}{kq}
       \le\limsup\frac{\Phi_{k,s}}{kq}
       \le\min\{\kappa_*,16C_\partial\}.
 \end{gathered}\tag{OVR9}
\]
For each line, subtract the other bounded nonnegative sequence from
the convergent sum in OVR3. More explicitly, that sum differs from
$16C_\partial$ by a sequence tending to zero; taking its upper and
lower eventual bounds proves
$\liminf(\Phi/(kq))=16C_\partial-\limsup(F_K/(kq))$ and
$\limsup(\Phi/(kq))=16C_\partial-\liminf(F_K/(kq))$.
The same argument with the two sequences interchanged proves the
first line, including its lower bound. The two first positive
increments also survive in the finite common upper bound. In the
original AKS notation let $d_1$ be the first section-corrected kernel
increment of AKS29, in the same source order $s$. AKS28 and OAR5 give
\[
 \begin{split}
 \Delta q C_\partial+R_{k,s}
   &=F_K^{(s)}+\Phi_{k,s}\\
   &\le(r+j)(\mathcal L_{s,q}+\mathcal L_{s,q+1})-d_1-\lambda_s\\
   &=\Delta(\mathcal L_{s,q}+\mathcal L_{s,q+1})-d_1-\lambda_s.
 \end{split}\tag{OVR10}
\]
The last equality is the exact rank identity
$(8k-16)+(8k-32)=16k-48$; both subtractions retain their actual
minimum-section definitions. Adding the eventual upper bounds
already proved and using OVR3 also proves the explicit consequence
\[
 0\le16C_\partial\le2\kappa_*,\qquad
 0\le C_\partial\le\frac{\kappa_*}{8}.
 \tag{OVR11}
\]
Indeed for every positive $\varepsilon$, both sequences are eventually
at most $\kappa_*+\varepsilon$, so their limiting sum is at most
$2\kappa_*+2\varepsilon$; let $\varepsilon$ decrease to zero.
Their nonnegativity proves the lower bound independently.
OVR4 and AMT46 carry the same limiting bounds to the two original
sources and the arithmetic kernel. The outstanding coefficient is
therefore the explicit original positive residual sum; no convergence
or numerical value for that sequence has been assumed.

% END CURRENT OVR

\end{document}
